% LOCKED SCAFFOLD - Volume I FINAL TOC - Chapter 2 only
% TOC authority: TOC/volume-I-ch01-04-FINAL.txt (Ch.2 block locked 2026-05-29)
% Flow audit: TOC/CH02-TOC-AUDIT-FLOW-LOCK.md
% Do not add Ch 3+ content to this file
\chapter{Maxwell Operators and Radiation Mappings}
\label{ch:02}
\setcounter{section}{-1}

\section{Introduction}
\label{sec:ch20}

Chapter~\ref{ch:01} closed with an inheritance ledger: which field-law results,
transport screens, gauge lemmas, and measurement contracts later chapters may
invoke without proof \cite{stratton1941}. Chapter~\ref{ch:02} consumes that
ledger and changes the working idiom. The governing balances of
Eq.~\eqref{eq:ch1.maxwell-local} are treated as fixed; the new objects are
\emph{maps}---from impressed source and boundary data to Maxwell-admissible field
states, and from those states to the outward-propagating content that alone can
carry far-zone flux compatible with the transport criteria of
Sections~\ref{sec:ch01.2.3}--\ref{sec:ch01.2.4}. Antenna practice speaks of
feed currents and radiation patterns; mathematical physics speaks of resolvents
and spectral projectors. Both vocabularies describe the same composition chain once
sources, closures, and far-zone selection are named as operators on declared
function spaces \cite{harrington2001,stratton1941}.

The shift is not a weakening of rigor. Chapter~\ref{ch:01} already showed that
angular descriptors attached to unspecified field ans\"atze can fail transport
admissibility even when local Maxwell residuals vanish. Chapter~\ref{ch:02}
therefore never postulates a modal field shape; it asks which \emph{linear} maps,
under which boundary classes, produce radiative outputs that downstream angular
analysis may safely inherit. Impressed data enter through the specification tuple
of Eq.~\eqref{eq:ch1.data-specification-tuple}; closure is encoded by the boundary
operator of Eq.~\eqref{eq:ch1.boundary-data-operator}; reactive versus radiative
partition is imported from Section~\ref{sec:ch01.2.4}. None of those structures
is re-derived here---they are the inputs on which operator domains and radiation
integrals are built \cite{stratton1941}.

Section~\ref{sec:ch20} states that program before the technical subsections
begin. Subsection~\ref{sec:ch201} gives the first HOME for the source-to-field
and radiation maps in frequency domain; Subsection~\ref{sec:ch202} packages the
harmonic Maxwell system as a closed boundary-value operator without reopening the
local derivations of Section~\ref{sec:ch01.1.1}; and Subsection~\ref{sec:ch203}
records narrative continuity with the Chapter~\ref{ch:01} bridge of
Section~\ref{sec:ch01.7.3}. Sections~\ref{sec:ch21}--\ref{sec:ch29} then develop
function-space domains, time--frequency equivalences, dyadic Green kernels,
Sommerfeld outgoing selection, radiation integrals, spectral/Fredholm structure,
non-radiating null spaces, reciprocity, and angular-spectrum action---each in its
assigned TOC home, each importing prior chapters only by reference.

Readers should expect a monograph pace comparable to Chapter~\ref{ch:01}:
every major relation carries explicit assumptions, domain qualifiers, and a
physical reading tied to measurable flux rather than interpretive labels alone.
Chapter~\ref{ch:02} supplies the operator spine on which Chapters~\ref{ch:03} and
\ref{ch:04} mount angular eigenstructure and realizability limits; it does not
anticipate vector spherical harmonics, Debye potentials, or gauge transformations,
which remain under their locked ownership elsewhere in Volume~I
\cite{stratton1941,harrington2001}.

\subsection{Radiation as a source-to-field mapping}
\label{sec:ch201}
% TOC tag: [HOME]

Laboratory radiation measurements always presuppose an excitation model even when
that model is left implicit. A network analyzer drives a port; a probe samples a
field component on an arc; a chamber termination fixes what fraction of energy
returns versus escapes. Subsection~\ref{sec:ch201} makes the implied chain
explicit: impressed data specified in the sense of
Section~\ref{sec:ch01.1.2} are mapped, through Maxwell closure, to a field pair
\((\E,\H)\), and only thereafter is the outward-propagating fraction identified as
the object that may carry far-zone power and angular flux
(Eq.~\eqref{eq:ch1.maxwell-local} and Sections~\ref{sec:ch01.1.4}--\ref{sec:ch01.1.5}
remain the inherited metering laws) \cite{stratton1941,harrington2001}. The
subsection establishes the first HOME in Volume~I for that two-step composition:
a harmonic \emph{solution map} and a \emph{radiation map} acting on admissible
sources at fixed angular frequency \(\omega>0\) under the \(\exp(-\jj\omega t)\)
convention locked in Chapter~\ref{ch:01}.

Let \(\Omega\subset\mathbb{R}^{3}\) be a piecewise smooth domain with boundary
\(\Gamma=\partial\Omega\). Phasor fields \(\E(\mathbf{r})\) and \(\H(\mathbf{r})\)
are suppressed in time unless causality is restored below. Impressed electric
current \(\Jimp\), charge density \(\rho\), and the data tuple \(\mathfrak{D}\) of
Eq.~\eqref{eq:ch1.data-specification-tuple} are not redefined; they are the
inputs on which the maps act. Admissibility is assumed in the sense of the
Chapter~\ref{ch:01} inheritance screen
(Eq.~\eqref{eq:ch1.inheritance-admissibility}): only configurations with
\(\mathcal{D}_{J}=1\) enter the domain of the radiation pipeline developed here
\cite{stratton1941}.

\begin{figure}[t]
\centering
\ChOneFigBlock{%
\begin{tikzpicture}[ch1fig, node distance=7mm and 11mm]
  \node[ch1box] (ex) {Excitation\\$\mathfrak{D}$};
  \node[ch1box, right=13mm of ex] (sw) {Solution map\\$\mathcal{S}_\omega$};
  \node[ch1box, right=13mm of sw] (st) {Full state\\$(\E,\H)$};
  \node[ch1box, right=13mm of st] (sel) {Outgoing\\selector $\Pi_{\mathrm{rad}}$};
  \node[ch1box, right=13mm of sel] (obs) {Far-zone\\observables};
  \draw[ch1flow] (ex) -- (sw) -- (st) -- (sel) -- (obs);
\end{tikzpicture}%
}
\caption{Excitation-to-observable chain: \(\mathcal{S}_\omega\) enforces Maxwell
closure with inherited boundary data; \(\Pi_{\mathrm{rad}}\) isolates the
radiative component on which far-zone flux integrals are meaningful.}
\label{fig:ch2.source-radiation-pipeline}
\end{figure}

Figure~\ref{fig:ch2.source-radiation-pipeline} organizes the logic used throughout
Chapter~\ref{ch:02}. The excitation block fixes \(\mathfrak{D}\), including
\(\Jimp\) and the boundary operator class
Eq.~\eqref{eq:ch1.boundary-data-operator}. The solution block implements
Eq.~\eqref{eq:ch1.maxwell-operator-block} together with the prescribed traces on
\(\Gamma\) \cite{maxwell1873,stratton1941}. The selector block is not a spatial
mask applied post hoc; it encodes outgoing-wave uniqueness compatible with
Section~\ref{sec:ch01.1.3} in the exterior limit treated in
Section~\ref{sec:ch243}. Observables in the final block are flux and pattern
quantities obtained only after selection.

Function spaces are declared before the maps. Let \(\mathcal{F}_\Omega\) be the
space of phasor pairs \((\E,\H)\) on \(\Omega\) whose components have
\(L^{2}\)-integrable curls and divergences in the sense required by
Eq.~\eqref{eq:ch1.maxwell-operator-block}, and whose boundary traces lie in the
domain of \(\mathcal{B}_\Gamma\). Let \(\mathcal{J}_\Omega\subset L^{2}(\Omega)^{3}\)
be the impressed-current space compatible with finite-energy admissibility
(Section~\ref{sec:ch01.0.3}), charge continuity
(Eq.~\eqref{eq:ch1.charge-continuity-harmonic}), and the same boundary class.
Refined Sobolev and distributional extensions of these spaces are supplied in
Section~\ref{sec:ch211}; here the declaration is minimal but precise enough to
state the maps \cite{stratton1941}.

The harmonic Maxwell \emph{solution map} sends admissible impressed data to a
field state:
\begin{equation}
\label{eq:ch2.maxwell-solution-map}
\mathcal{S}_\omega:\mathcal{J}_\Omega\longrightarrow\mathcal{F}_\Omega,
\qquad
\Jimp\longmapsto(\E,\H),
\end{equation}
defined by \(\mathcal{S}_\omega[\Jimp]=(\E,\H)\) satisfying the Maxwell operator
residuals of Eq.~\eqref{eq:ch1.maxwell-operator-block} with sources \(\mathfrak{D}\)
and boundary data \(\mathbf{g}_\Gamma\) \cite{stratton1941}.
Equation~\eqref{eq:ch2.maxwell-solution-map} is the Chapter~\ref{ch:02} packaging
of the initial-boundary value problem proven in Section~\ref{sec:ch01.1.1}; it
does not alter the law, only the input--output viewpoint.

Injectivity fails in general: distinct \(\Jimp\) can produce identical \((\E,\H)\)
on observational subdomains, and equivalent source distributions related by
divergence-free reparametrizations inside \(\Omega\) can leave exterior patterns
unchanged when \(\mathcal{B}_\Gamma\) permits reactive storage
(Section~\ref{sec:ch01.2.4}). Those equivalences are not nuisances to be removed
by definition; they foreshadow null-space and realizability analyses in
Sections~\ref{sec:ch27} and Chapter~\ref{ch:04}. Radiation bookkeeping therefore
never begins from a guessed modal field; it begins from \(\mathcal{S}_\omega[\Jimp]\)
with explicit closure.

The full state typically contains both stored and exported electromagnetic energy.
Reactive fields exchange power with sources and terminations at distances small
compared with \(k_0^{-1}\); radiative fields carry net outward flux at infinity
in open configurations. Let \(\mathcal{F}_{\mathrm{rad}}\subset\mathcal{F}_\Omega\)
denote pairs whose far-zone asymptotics satisfy outgoing selection (made analytic in
Section~\ref{sec:ch24}). Define the \emph{outgoing selector}
\(\Pi_{\mathrm{rad}}:\mathcal{F}_\Omega\to\mathcal{F}_{\mathrm{rad}}\) and the
\emph{radiation operator}
\begin{equation}
\label{eq:ch2.radiation-operator-def}
\mathcal{R}_\omega:\mathcal{J}_\Omega\longrightarrow\mathcal{F}_{\mathrm{rad}},
\qquad
\mathcal{R}_\omega[\Jimp]=\Pi_{\mathrm{rad}}\,\mathcal{S}_\omega[\Jimp].
\end{equation}
Equation~\eqref{eq:ch2.radiation-operator-def} is the first formal Volume~I home
for \(\mathcal{R}_\omega\) \cite{stratton1941}. Until Subsection~\ref{sec:ch224}
refines \(\Pi_{\mathrm{rad}}=\mathcal{O}_{\mathrm{out}}\circ\Pi_{\mathrm{prop}}\) in
Eq.~\eqref{eq:ch2.rad-projection-composition}, the symbol \(\Pi_{\mathrm{rad}}\) denotes
the outgoing radiative selector abstractly: every far-zone or transport claim before
§2.2.4 must be read as provisional pending spectral refinement
\cite{stratton1941,harrington2001}. Transport admissibility in
Sections~\ref{sec:ch01.2.3}--\ref{sec:ch01.2.4} applies to
\(\mathcal{R}_\omega[\Jimp]\), not to arbitrary components of
\(\mathcal{S}_\omega[\Jimp]\). Angular labels attached before selection can
therefore describe near-zone storage rather than exported flux.

Three composition stages organize the remainder of the chapter: specification
(\(\Jimp\in\mathcal{J}_\Omega\) with \(\mathcal{D}_{J}=1\)), closure
(\(\mathcal{S}_\omega\) enforcing Eq.~\eqref{eq:ch1.maxwell-local} on \(\Omega\)),
and export (\(\mathcal{R}_\omega\) retaining only outgoing radiative content).
Stages (i)--(ii) inherit law consistency from Chapter~\ref{ch:01}; stage (iii)
inherits the reactive/radiative partition of Section~\ref{sec:ch01.2.4}. Skipping
stage (iii) while asserting far-zone angular transport is the most common logical
failure mode in structured-beam interpretations flagged in
Section~\ref{sec:ch01.6.2} \cite{allen1992,stratton1941}.

When Green operators are available (Section~\ref{sec:ch23}), \(\mathcal{R}_\omega\)
admits a retarded integral representation. For \(\mathbf{r}\notin\Omega_s\) with
\(\operatorname{supp}\Jimp=\Omega_s\Subset\Omega\),
\begin{equation}
\label{eq:ch2.radiation-integral-preview}
\E_{\mathrm{rad}}(\mathbf{r})
=
-\jj\omega\mu_0
\int_{\Omega_s}
\boldsymbol{\Gamma}(\mathbf{r},\mathbf{r}')\cdot\Jimp(\mathbf{r}')\,d^{3}r',
\end{equation}
where \(\boldsymbol{\Gamma}\) is the outgoing dyadic Green function
\cite{stratton1941}. Equation~\eqref{eq:ch2.radiation-integral-preview} is a
forward pointer, not a duplicate definition: kernel construction, singularities,
and convergence domains belong to Sections~\ref{sec:ch231}--\ref{sec:ch234}; the
radiation integrals of Section~\ref{sec:ch25} realize the same map under explicit
far-zone asymptotics.

Every invocation of
Eqs.~\eqref{eq:ch2.maxwell-solution-map}--\eqref{eq:ch2.radiation-operator-def}
must state its domain class. Unless noted otherwise, \(\omega>0\) is fixed,
materials are linear isotropic passive regions, and \(\Omega\) is either bounded
with declared \(\mathcal{B}_\Gamma\) or an exterior domain with compact
\(\operatorname{supp}\Jimp\) and outgoing closure at infinity. Maps are defined
only on \(\mathcal{J}_\Omega\) satisfying Eq.~\eqref{eq:ch1.charge-continuity-harmonic};
generalized sources appear in Section~\ref{sec:ch232}.

Maxwell linearity on \(\Omega\) induces linearity of both maps on fixed boundary
data: for \(\Jimp^{(1)},\Jimp^{(2)}\in\mathcal{J}_\Omega\) and scalars \(a_{1},a_{2}\),
\begin{equation}
\label{eq:ch2.radiation-linearity}
\mathcal{R}_\omega\!\left[a_{1}\Jimp^{(1)}+a_{2}\Jimp^{(2)}\right]
=
a_{1}\mathcal{R}_\omega[\Jimp^{(1)}]
+
a_{2}\mathcal{R}_\omega[\Jimp^{(2)}],
\end{equation}
when \(\mathbf{g}_\Gamma\) and the material layout are held constant
\cite{maxwell1873,stratton1941}. Superposition of far-zone patterns is therefore
legitimate only after \(\mathcal{R}_\omega\), not after raw near-field addition.
Composite-source angular analysis decomposes \(\Jimp\) into admissible constituents,
maps each through \(\mathcal{R}_\omega\), and superposes radiative outputs---a
discipline that prevents reactive interference from masquerading as transport
(Section~\ref{sec:ch01.6.4}).

Write \(\operatorname{ran}\mathcal{S}_\omega\subset\mathcal{F}_\Omega\) for all
field pairs reachable from admissible sources and
\(\operatorname{ran}\mathcal{R}_\omega\subset\mathcal{F}_{\mathrm{rad}}\) for
far-zone radiative outputs. Strict inclusion
\(\operatorname{ran}\mathcal{R}_\omega\subsetneq\operatorname{ran}\mathcal{S}_\omega\)
occurs whenever reactive storage is nontrivial: distinct near-zone states can share
identical radiative images (Section~\ref{sec:ch01.2.4}). Pattern measurements
probe \(\operatorname{ran}\mathcal{R}_\omega\) and cannot invert
\(\mathcal{S}_\omega\) without additional near-zone data---a limitation revisited in
Section~\ref{sec:ch244}.

Radiative power is metered with inherited Poynting machinery. On a sphere
\(S_R=\{|\mathbf{r}|=R\}\) enclosing \(\operatorname{supp}\Jimp\), with outward
unit vector \(\hat{\mathbf{r}}\) and complex Poynting vector
\(\mathbf{S}=\tfrac{1}{2}\E\times\H^*\) from Section~\ref{sec:ch01.1.4},
\begin{equation}
\label{eq:ch2.radiated-power-screen}
P_{\mathrm{rad}}(R)
=
\Re\oint_{S_R}\mathbf{S}\cdot\hat{\mathbf{r}}\,dS.
\end{equation}
Equation~\eqref{eq:ch2.radiated-power-screen} does not introduce a new flux law; it
specifies which field pair enters the surface integral when the observable is
radiative. For \(\mathcal{R}_\omega[\Jimp]\), \(P_{\mathrm{rad}}(R)\) becomes
\(R\)-independent beyond the reactive shell and equals the power carried by
\(\operatorname{ran}\mathcal{R}_\omega\) \cite{poynting1884,stratton1941}. Angular
momentum budgets from Section~\ref{sec:ch01.1.5} attach to the same radiative
component only.

\begin{figure}[t]
\centering
\ChOneFigBlock{%
\begin{tikzpicture}[ch1fig]
  \node[ch1box, minimum width=22mm] (core) at (0,0) {Source\\$\Omega_s$};
  \node[ch1box, minimum width=28mm] (react) at (0,-12mm) {Reactive shell\\$1/r^2$, $1/r^3$};
  \node[ch1box, minimum width=32mm] (sel) at (0,-24mm) {Selector\\$\Pi_{\mathrm{rad}}$};
  \node[ch1box, minimum width=34mm] (far) at (0,-36mm) {Far zone\\$1/r$ outgoing};
  \draw[ch1flow] (core) -- (react) -- (sel) -- (far);
\end{tikzpicture}%
}
\caption{Spatial hierarchy seen by the maps: \(\mathcal{S}_\omega\) populates all
regions; \(\Pi_{\mathrm{rad}}\) passes only the outgoing transverse content
compatible with Section~\ref{sec:ch01.2.4}.}
\label{fig:ch2.reactive-radiative-split}
\end{figure}

Figure~\ref{fig:ch2.reactive-radiative-split} emphasizes geometry rather than
algebra. The solution map energizes reactive and radiative regions alike; the
radiation map discards non-outgoing storage before pattern or transport claims are
made. Misidentifying reactive circulation as exported angular structure is a
recurring error in high-gradient near-zone beams \cite{allen1992,stratton1941}.

Finite measurement channels (Section~\ref{sec:ch01.5}) implement linear functionals
on field space. Composing a declared far-zone observable with \(\mathcal{R}_\omega\)
yields a linear functional of \(\Jimp\); composing the same observable with
\(\mathcal{S}_\omega\) without selection generally includes reactive contamination.
No new measurement axioms are introduced---only the correct map for the declared
observable class \cite{stratton1941}.

Non-radiating admissible currents illustrate kernel visibility. If
\(\Jimp^{\mathrm{(NR)}}\in\mathcal{J}_\Omega\) has vanishing far-zone transverse
fields while maintaining near-zone storage (Section~\ref{sec:ch01.2.4}), then
\begin{equation}
\label{eq:ch2.null-radiation-preview}
\mathcal{R}_\omega[\Jimp^{\mathrm{(NR)}}]=\mathbf{0}
\quad\text{in}\ \mathcal{F}_{\mathrm{rad}},
\end{equation}
yet \(\mathcal{S}_\omega[\Jimp^{\mathrm{(NR)}}]\neq\mathbf{0}\) in general.
Equation~\eqref{eq:ch2.null-radiation-preview} targets the null space analyzed in
Section~\ref{sec:ch272}; ignoring \(\ker\mathcal{R}_\omega\) invites spurious modal
assignments to invisible sources.

Membership in \(\mathcal{J}_\Omega\) requires more than square integrability:
\begin{equation}
\label{eq:ch2.source-admissibility}
\mathcal{D}_{J}[\Jimp]=1
\iff
\left(
\Jimp\in L^{2}(\Omega)^{3},\
\divg\Jimp=-\jj\omega\rho,\
\Jimp\ \text{satisfies BC compatibility on }\Gamma
\right),
\end{equation}
with \(\mathcal{D}_{J}\) as in Eq.~\eqref{eq:ch1.inheritance-admissibility}.
Violating \(\divg\Jimp=-\jj\omega\rho\) contradicts
Eq.~\eqref{eq:ch1.charge-continuity-harmonic} before Maxwell closure is attempted.
Boundary compatibility ensures induced surface equivalents do not contradict
Eq.~\eqref{eq:ch1.boundary-data-operator} \cite{stratton1941}.

Causality remains an admissibility constraint even in harmonic form. With time-domain
current \(j(\mathbf{r},t)\) and retarded dyadic impulse
\(\mathcal{G}(\mathbf{r},\mathbf{r}',t-t')\),
\begin{equation}
\label{eq:ch2.causality-preview}
\E(\mathbf{r},t)
=
-\mu_0\,\frac{\partial}{\partial t}
\int
\mathcal{G}(\mathbf{r},\mathbf{r}',t-t')\cdot j(\mathbf{r}',t')\,d^{3}r'\,dt',
\end{equation}
integrated only over \(t'\le t\) \cite{stratton1941}. The steady-state operator
\(\mathcal{R}_\omega\) is the frequency specialization of this retarded chain;
equivalence conditions are proved in Section~\ref{sec:ch223}.

Vacuum scales from Eq.~\eqref{eq:ch1.vacuum-impedance-wavenumber} fix the carrier
on which selection acts: \(k_0=\omega\sqrt{\mu_0\varepsilon_0}\),
\(\eta_0=\sqrt{\mu_0/\varepsilon_0}\). Outgoing components retained by
\(\Pi_{\mathrm{rad}}\) carry \(\exp(-\jj k_0 r)\) phase fronts at leading order;
components whose phase does not lock to outgoing \(k_0 r\) remain in
\(\mathcal{S}_\omega[\Jimp]\) but outside \(\mathcal{F}_{\mathrm{rad}}\). Power
metered by Eq.~\eqref{eq:ch2.radiated-power-screen} stabilizes only on the latter
\cite{stratton1941}.

Distance labels become meaningful through operator action. For compact
\(\operatorname{supp}\Jimp\), \(\mathcal{S}_\omega[\Jimp]\) contains
\(1/r^{2}\) and \(1/r^{3}\) reactive terms; \(\mathcal{R}_\omega[\Jimp]\) retains
leading \(1/r\) transverse fields with \(kr\) phase in the far zone
(Section~\ref{sec:ch01.2.4}). Asymptotic partitions are sharpened in
Section~\ref{sec:ch252}.

Observability enters as functional composition rather than as a new postulate.
Let \(\mathcal{O}_{\mathrm{far}}\) denote a linear functional that extracts a
declared far-zone scalar or vector observable---pattern voltage, axial ratio
proxy, or a single spherical-harmonic coefficient---from a field pair in
\(\mathcal{F}_{\mathrm{rad}}\). The measured quantity associated with impressed
current \(\Jimp\) is the scalar
\(\mathcal{O}_{\mathrm{far}}\big[\mathcal{R}_\omega[\Jimp]\big]\), not
\(\mathcal{O}_{\mathrm{far}}\big[\mathcal{S}_\omega[\Jimp]\big]\) unless the
observable has been shown to vanish on reactive content. Because both
\(\mathcal{R}_\omega\) and \(\mathcal{O}_{\mathrm{far}}\) are linear on their
domains, the laboratory reading is a linear functional
\(\Lambda_{\mathrm{rad}}=\mathcal{O}_{\mathrm{far}}\circ\mathcal{R}_\omega\)
on \(\mathcal{J}_\Omega\). That composition is the operator-theoretic form of
the measurement screening in Section~\ref{sec:ch01.5}: finite apertures and
finite modal sets appear as rank-deficient \(\mathcal{O}_{\mathrm{far}}\), not
as modifications of Maxwell law \cite{stratton1941}. Angular-transport claims
require further that \(\Lambda_{\mathrm{rad}}\) be compatible with the flux
budgets of Section~\ref{sec:ch01.1.5} when the observable is interpreted as
momentum or power export.

The outgoing selector \(\Pi_{\mathrm{rad}}\) is idempotent on
\(\mathcal{F}_{\mathrm{rad}}\): applying selection twice leaves a radiative state
unchanged. On the larger space \(\mathcal{F}_\Omega\), \(\Pi_{\mathrm{rad}}\) is
a projection only after the Sommerfeld-class uniqueness of Section~\ref{sec:ch243}
is imposed; without that closure, multiple outgoing decompositions can coexist
for the same truncated near-zone data. The idempotence condition is therefore a
\emph{joint} property of field space and boundary class, not a algebraic trick
independent of physics \cite{harrington2001,stratton1941}. Numerical pipelines that
window fields in space and call the result ``radiative'' without outgoing closure
are approximating \(\Pi_{\mathrm{rad}}\) by a non-idempotent mask; such outputs
must not be fed into transport criteria without an explicit error budget.

The subsection closes with the composition law used downstream:
\begin{equation}
\label{eq:ch2.pipeline-composition}
\mathcal{R}_\omega
=
\Pi_{\mathrm{rad}}\circ\mathcal{S}_\omega,
\qquad
\mathcal{D}_{J}=1
\Longrightarrow
\mathcal{S}_\omega:\mathcal{J}_\Omega\to\mathcal{F}_\Omega,
\quad
\mathcal{R}_\omega:\mathcal{J}_\Omega\to\mathcal{F}_{\mathrm{rad}}.
\end{equation}
Green kernels realize \(\mathcal{S}_\omega\); Sommerfeld conditions define
\(\Pi_{\mathrm{rad}}\); spectral theory in Section~\ref{sec:ch26} studies
eigenstructure on spaces refined in Section~\ref{sec:ch21}
\cite{harrington2001,stratton1941}. Nothing in
Eqs.~\eqref{eq:ch2.maxwell-solution-map}--\eqref{eq:ch2.pipeline-composition}
re-proves results listed in Section~\ref{sec:ch01.7.1}; the bridge
Eq.~\eqref{eq:ch1.bridge-inheritance-map} is honored by citation-only inheritance
of curl--curl algebra, stress flux theorems, and gauge lemmas.

\paragraph{Worked illustration (conceptual).}
Take \(\Jimp(\mathbf{r})=\hat{\mathbf{x}}\,I_0\,h(\mathbf{r})\) with
\(h\in C_c^\infty(\Omega_s)\) supported in a ball of radius \(a\ll \lambda_0=2\pi/k_0\)
in free space with outgoing closure. \(\mathcal{S}_\omega[\Jimp]\) includes strong
near-zone reactive gradients scaling as \(1/R^{2}\) and \(1/R^{3}\) on a sphere
\(S_R\) with \(a\ll R\ll \lambda_0\). The selected radiative part carries dominant
\(1/R\) transverse fields whose normal Poynting flux through \(S_R\) approaches the
\(R\)-independent value of Eq.~\eqref{eq:ch2.radiated-power-screen}. Only
\(\mathcal{R}_\omega[\Jimp]\) may be tested against the transport criteria of
Section~\ref{sec:ch01.2.3}; numeric flux evaluation awaits Section~\ref{sec:ch251}.

In a closed cavity with perfect-conductor walls, the same \(\Jimp\) may excite
standing modes whose energy remains in \(\operatorname{ran}\mathcal{S}_\omega\) while
\(\Pi_{\mathrm{rad}}\) annihilates exterior export---not because
Eq.~\eqref{eq:ch2.maxwell-solution-map} fails, but because
Eq.~\eqref{eq:ch1.outer-boundary-classes} fixes a non-outgoing closure
\cite{harrington2001,stratton1941}. Open-region and cavity responses are comparable
only when \(\mathcal{B}_\Gamma\) and exterior termination are declared identical.

\paragraph{Closing summary.}
Radiation in Maxwell theory is the image of admissible impressed sources under
\(\Pi_{\mathrm{rad}}\circ\mathcal{S}_\omega\). Volume~I attaches every angular
descriptor downstream to that image, meters it with inherited flux laws, and screens
it with the transport and observability constraints established in
Chapter~\ref{ch:01}.

\subsection{Maxwell equations in operator form}
\label{sec:ch202}
% TOC tag: [HOME]
The solution map \(\mathcal{S}_\omega\) of Eq.~\eqref{eq:ch2.maxwell-solution-map} names an
input--output relation; Subsection~\ref{sec:ch202} identifies the \emph{operator}
whose inversion (under declared closure) produces that map. Harmonic Maxwell theory
on \(\Omega\) is packaged as a first-order residual system on a six-component state,
augmented by boundary traces through the operator
Eq.~\eqref{eq:ch1.boundary-data-operator} imported from
Section~\ref{sec:ch01.1.2}. Local curl and divergence balances remain those of
Eq.~\eqref{eq:ch1.maxwell-operator-block}; nothing in this subsection re-derives
them \cite{stratton1941,harrington2001}.

Collect phasor fields into the Maxwell state
\begin{equation}
\label{eq:ch2.maxwell-state-vector}
\mathbf{u}_\omega=
\begin{bmatrix}\E\\ \H\end{bmatrix}
\in\mathcal{F}_\Omega,
\end{equation}
where \(\mathcal{F}_\Omega\) is the admissible field space fixed in
Subsection~\ref{sec:ch201}. The pair \((\E,\H)\) is constrained jointly: divergence
rows of Eq.~\eqref{eq:ch1.maxwell-operator-block} must hold, and tangential traces on
\(\Gamma=\partial\Omega\) must lie in the domain of the boundary class declared with
\(\mathfrak{D}\). Treating \(\mathbf{u}_\omega\) as a single unknown is what permits
resolvent notation, variational forms, and the spectral compression developed in
Section~\ref{sec:ch26} \cite{harrington2001}.

Impressed data enter through a forcing vector aligned with the Maxwell residual
rows. For the tuple Eq.~\eqref{eq:ch1.data-specification-tuple},
\begin{equation}
\label{eq:ch2.maxwell-forcing-vector}
\mathbf{s}_\omega(\mathfrak{D})=
\begin{bmatrix}\mathbf{0}\\ -\Jimp\\ -\rho\\ \mathbf{0}\end{bmatrix},
\end{equation}
where row order matches Eq.~\eqref{eq:ch1.maxwell-operator-block}. Equation
\eqref{eq:ch2.maxwell-forcing-vector} records drive on the Amp\`ere and Gauss rows
only for the impressed-current class treated here; sheet and magnetic excitations
extend the vector through Eq.~\eqref{eq:ch1.interface-conditions} when trace spaces
are enriched in Section~\ref{sec:ch211} \cite{stratton1941}. The forcing is \emph{data},
not a field variable: it specifies how \(\mathcal{S}_\omega\) is sourced while
\(\mathbf{u}_\omega\) remains the unknown.

Let \(\mathcal{M}_\omega\) denote the harmonic Maxwell operator of
Eq.~\eqref{eq:ch1.maxwell-operator-block}. The volume subsystem reads
\begin{equation}
\label{eq:ch2.maxwell-driven-operator}
\mathcal{M}_\omega[\mathbf{u}_\omega]=\mathbf{s}_\omega(\mathfrak{D}),
\end{equation}
and is satisfied when each residual row vanishes. Linearity in \(\mathbf{u}_\omega\)
follows from linearity of Eq.~\eqref{eq:ch1.maxwell-local}; spatially varying
\(\varepsilon(\mathbf{r})\) and \(\mu(\mathbf{r})\) modify coefficients inside
\(\mathcal{M}_\omega\) without changing the first-order block shape
(Section~\ref{sec:ch222}) \cite{stratton1941}. In uniform vacuum regions,
\(\mathcal{M}_\omega\) has translation-invariant symbol, the property exploited by
dyadic Green kernels in Section~\ref{sec:ch23}.

\begin{figure}[t]
\centering
\ChOneFigBlock{%
\begin{tikzpicture}[ch1fig, node distance=8mm and 8mm]
  \node[ch1box] (vol) {Volume\\$\mathcal{M}_\omega\mathbf{u}=\mathbf{s}$};
  \node[ch1box, below=12mm of vol] (bnd) {Boundary\\$\mathcal{B}_\Gamma\gamma_\Gamma\mathbf{u}=\mathbf{g}$};
  \node[ch1box, left=18mm of vol] (u) {State $\mathbf{u}_\omega$};
  \draw[ch1flow] (u) -- (vol);
  \draw[ch1flow] (u) -- (bnd);
\end{tikzpicture}%
}
\caption{Parallel volume and boundary channels for the closed harmonic problem:
\(\mathbf{u}_\omega\) simultaneously satisfies Maxwell residuals and prescribed
traces.}
\label{fig:ch2.maxwell-boundary-stack}
\end{figure}

Figure~\ref{fig:ch2.maxwell-boundary-stack} is the structural diagram for
Eqs.~\eqref{eq:ch2.maxwell-driven-operator} and
\eqref{eq:ch1.boundary-data-operator}. A solver that enforces only volume rows
produces a formal solution compatible with Maxwell \emph{sources} but not necessarily
with termination physics; conversely, boundary data without volume balance violate
conservation. Subsection~\ref{sec:ch201} applies \(\mathcal{R}_\omega\) only after
both channels are closed \cite{harrington2001,stratton1941}.

Tangential traces are extracted by the restriction
\begin{equation}
\label{eq:ch2.trace-restriction}
\gamma_\Gamma:\mathcal{F}_\Omega\to\mathcal{T}_\Gamma,
\qquad
\gamma_\Gamma[\mathbf{u}_\omega]=
\begin{bmatrix}
\left.\hat{\mathbf{n}}\times\E\right|_{\Gamma}\\
\left.\hat{\mathbf{n}}\times\H\right|_{\Gamma}
\end{bmatrix},
\end{equation}
with outward unit normal \(\hat{\mathbf{n}}\) on \(\Gamma\) and trace space
\(\mathcal{T}_\Gamma\) induced by admissible interior states. The map
\(\gamma_\Gamma\) is not defined on arbitrary \(L^{2}\) six-tuples; Sobolev and
distribution refinements in Section~\ref{sec:ch211} specify its domain
\cite{stratton1941}.

The closed operator stacks volume and boundary residuals:
\begin{equation}
\label{eq:ch2.maxwell-boundary-composite}
\mathcal{L}_\omega[\mathbf{u}_\omega]=
\begin{bmatrix}
\mathcal{M}_\omega[\mathbf{u}_\omega]\\
\mathcal{B}_\Gamma[\gamma_\Gamma[\mathbf{u}_\omega]]
\end{bmatrix}.
\end{equation}
The harmonic boundary-value problem is to find \(\mathbf{u}_\omega\in\mathcal{F}_\Omega\)
such that
\begin{equation}
\label{eq:ch2.closed-maxwell-bvp}
\mathcal{L}_\omega[\mathbf{u}_\omega]=
\begin{bmatrix}\mathbf{s}_\omega(\mathfrak{D})\\ \mathbf{g}_\Gamma\end{bmatrix}.
\end{equation}
Perfect conductors impose \(\hat{\mathbf{n}}\times\E=\mathbf{0}\); radiation and
Leontovich closures encode outgoing or dissipative relations on
\(\gamma_\Gamma[\mathbf{u}_\omega]\). The operator \(\mathcal{L}_\omega\) is the
Chapter~\ref{ch:02} home for this coupled system \cite{harrington2001}.

Equation~\eqref{eq:ch2.maxwell-solution-map} is the resolvent relation
\begin{equation}
\label{eq:ch2.solution-map-resolvent}
\mathcal{S}_\omega[\Jimp]=\mathbf{u}_\omega
\quad\text{with}\quad
\mathcal{L}_\omega[\mathbf{u}_\omega]=
\begin{bmatrix}\mathbf{s}_\omega(\Jimp,\rho)\\ \mathbf{g}_\Gamma\end{bmatrix},
\end{equation}
where \(\rho\) is tied to \(\Jimp\) through
Eq.~\eqref{eq:ch1.charge-continuity-harmonic} when only electric impressed currents
drive the system. Invertibility of \(\mathcal{L}_\omega\) is not universal: existence
and uniqueness depend on \(\mathcal{B}_\Gamma\) and on whether outgoing selection is
imposed (Section~\ref{sec:ch243} for exterior domains; Section~\ref{sec:ch261} for
cavities) \cite{stratton1941}. \(\mathcal{S}_\omega\) is therefore a \emph{physical
branch} of \(\mathcal{L}_\omega^{-1}\), not an abstract algebraic inverse on all of
\(\mathcal{F}_\Omega\).

\begin{figure}[t]
\centering
\ChOneFigBlock{%
\begin{tikzpicture}[ch1fig, node distance=10mm]
  \node[ch1box] (E) at (0,6mm) {$\curl\E$};
  \node[ch1box] (H) at (0,-6mm) {$\curl\H$};
  \node[ch1box, right=20mm of E] (r1) {$-\jj\omega\mu_0\H$};
  \node[ch1box, right=20mm of H] (r2) {$\Jimp+\jj\omega\varepsilon_0\E$};
  \draw[ch1flow] (E) -- (r1);
  \draw[ch1flow] (H) -- (r2);
  \draw[ch1flow, dashed] (r1.south) to[bend left=25] (H.east);
  \draw[ch1flow, dashed] (r2.north) to[bend left=25] (E.east);
\end{tikzpicture}%
}
\caption{Circulation coupling in \(\mathcal{M}_\omega\): each curl row feeds the
partner field; \(\Jimp\) appears as external drive on the magnetic row.}
\label{fig:ch2.first-order-curl-coupling}
\end{figure}

Figure~\ref{fig:ch2.first-order-curl-coupling} shows why Maxwell structure is
first-order in both fields. Single-field Helmholtz reductions equivalent to
Eqs.~\eqref{eq:ch1.vacuum-wave-source-driven}--\eqref{eq:ch1.vacuum-curl-second-order}
are used later when a scalar symbol simplifies kernels; the elimination belongs to
Chapter~\ref{ch:01} and is not repeated. The electric-wave symbol used in
Chapter~\ref{ch:02} is recorded as
\begin{equation}
\label{eq:ch2.electric-wave-operator}
\mathcal{W}_{E,\omega}\E\equiv
\curl\!\left(\mu_0^{-1}\curl\E\right)-\omega^2\varepsilon_0\E
= -\jj\omega\mu_0\Jimp+\mu_0^{-1}\curl\mathbf{J}^{\mathrm{(mag)}}_{\mathrm{eff}},
\end{equation}
with \(\mathbf{J}^{\mathrm{(mag)}}_{\mathrm{eff}}\) collecting equivalent magnetic
drive when present \cite{stratton1941}. Subsection~\ref{sec:ch222} supplies the
magnetic counterpart and material-weighted generalizations.

Variational structure is inherited from Section~\ref{sec:ch01.1.4}. For test state
\(\mathbf{v}_\omega=(\E',\H')\),
\begin{equation}
\label{eq:ch2.maxwell-energy-form}
\mathfrak{a}_\omega(\mathbf{u}_\omega,\mathbf{v}_\omega)=
\int_\Omega\Big[\mu_0\,\H\!\cdot\!\curl\E'-\varepsilon_0\,\E\!\cdot\!\curl\H'\Big]d^{3}r.
\end{equation}
Integration by parts connects Eq.~\eqref{eq:ch2.maxwell-energy-form} to surface
fluxes and to the weak formulation
\begin{equation}
\label{eq:ch2.maxwell-weak-form}
\mathfrak{a}_\omega(\mathbf{u}_\omega,\mathbf{v}_\omega)=\mathfrak{f}_\omega(\mathbf{v}_\omega)
\quad\forall\,\mathbf{v}_\omega\in\mathcal{F}_\Omega^{0},
\end{equation}
where \(\mathcal{F}_\Omega^{0}\) carries homogeneous boundary traces compatible with
\(\mathcal{B}_\Gamma\) and \(\mathfrak{f}_\omega\) is induced by
\(\mathbf{s}_\omega(\mathfrak{D})\) \cite{poynting1884,stratton1941}. Equation
\eqref{eq:ch2.maxwell-weak-form} is the variational face of
Eq.~\eqref{eq:ch2.closed-maxwell-bvp}; finite-element discretizations favor this
route when material coefficients are piecewise constant
\cite{harrington2001}.

Well-posedness of Eq.~\eqref{eq:ch2.maxwell-weak-form} is not automatic on
unbounded domains or on multiply connected cavities. On a bounded region with
passive \(\mathcal{B}_\Gamma\) and homogeneous traces in
\(\mathcal{F}_\Omega^{0}\), coercivity fails only along divergence-free
cavity modes that satisfy both curl-curl balance and homogeneous tangential
data---precisely the null space that Section~\ref{sec:ch214} isolates before
imposing gauge or polarization fixes \cite{harrington2001,stratton1941}. For
exterior problems, coercivity is recovered only after restricting to the
outgoing subspace selected by Sommerfeld-type closure; without that restriction,
the bilinear form is indefinite and the weak problem admits non-physical
growing solutions. The practical rule for antenna analysis is therefore
threefold: declare \(\mathcal{D}(\mathcal{L}_\omega)\), specify
\(\mathcal{B}_\Gamma\) on every connected piece of \(\Gamma\), and verify that
the impressed tuple lies in the compatible range of
Eq.~\eqref{eq:ch2.source-compatibility}. Numerical codes that enforce volume
Maxwell rows on truncated domains but leave exterior faces open implement an
\emph{implicit} \(\mathcal{B}_\Gamma\); unless that implicit closure matches
the laboratory termination, \(\mathcal{S}_\omega\) in the code is a different
operator from the one measured on the bench \cite{harrington2001}.

Domain specification cannot remain at \(C^\infty\) when sources are distributional.
The minimal energy class declared here is
\begin{equation}
\label{eq:ch2.operator-domain-declaration}
\mathcal{D}(\mathcal{L}_\omega)\subset
\big\{\mathbf{u}_\omega=(\E,\H):
\E,\H\in L^{2}(\Omega)^{3},\ \curl\E,\curl\H\in L^{2}(\Omega)^{3}\big\},
\end{equation}
augmented by trace regularity required by \(\mathcal{B}_\Gamma\).
Section~\ref{sec:ch211} refines embeddings and compactness; evaluating
Eq.~\eqref{eq:ch2.maxwell-driven-operator} outside
\(\mathcal{D}(\mathcal{L}_\omega)\) generates spurious boundary terms that corrupt
flux identities \cite{stratton1941}.

Piecewise smooth \(\varepsilon(\mathbf{r})\) and \(\mu(\mathbf{r})\) modify
coefficients inside \(\mathcal{M}_\omega\) while preserving the block template of
Eq.~\eqref{eq:ch1.maxwell-operator-block}. Interface jumps satisfy
Eq.~\eqref{eq:ch1.interface-conditions}; multi-region Green operators in
Section~\ref{sec:ch234} realize \(\mathcal{L}_\omega^{-1}\) on composite domains
\cite{stratton1941}. The coupling pattern of
Figure~\ref{fig:ch2.first-order-curl-coupling} is unchanged---only symbols and trace
maps vary with geometry.

Finite terminals often impose impedance relations between tangential components:
\begin{equation}
\label{eq:ch2.boundary-impedance-operator}
\mathcal{Z}_\Gamma\!
\begin{bmatrix}
\left.\hat{\mathbf{n}}\times\E\right|_{\Gamma}\\
\left.\hat{\mathbf{n}}\times\H\right|_{\Gamma}
\end{bmatrix}=\mathbf{0},
\end{equation}
encoding Leontovich, absorbing, or mode-matched closures \cite{harrington2001}.
When Eq.~\eqref{eq:ch2.boundary-impedance-operator} replaces explicit
\(\mathbf{g}_\Gamma\) in Eq.~\eqref{eq:ch2.closed-maxwell-bvp}, admissible states
satisfy the impedance manifold on \(\Gamma\). Exterior outgoing problems correspond
to a limit in which \(\mathcal{Z}_\Gamma\) selects the Sommerfeld branch of
Section~\ref{sec:ch243}; conflating a finite-chamber termination with that limit
without normalization corrupts far-zone patterns \cite{stratton1941}.

Forcing must lie in the range implied by charge conservation. Divergence of the
Amp\`ere row together with Eq.~\eqref{eq:ch1.charge-continuity-harmonic} yields
\begin{equation}
\label{eq:ch2.source-compatibility}
\divg\mathbf{s}_{\omega,2}+\jj\omega\,\mathbf{s}_{\omega,3}=0
\end{equation}
on the second and third blocks of Eq.~\eqref{eq:ch2.maxwell-forcing-vector}.
Inconsistent \((\Jimp,\rho)\) pairs lie outside physically realizable drive;
numerical solvers may still return vectors that violate
Eq.~\eqref{eq:ch2.source-compatibility} \cite{stratton1941}. A compact
current filament with vanishing \(\rho\) automatically satisfies
Eq.~\eqref{eq:ch2.source-compatibility} because \(\mathbf{s}_{\omega,3}=0\).
When a volumetric \(\rho\) is prescribed independently of \(\Jimp\), the pair
must be adjusted so that their divergence balance matches the Amp\`ere--Gauss
rows; otherwise the closed problem is mathematically posed but physically
inconsistent. Feed models that inject \(\Jimp\) without updating \(\rho\) are
common in method-of-moments pipelines; they remain valid only when the
discretization enforces continuity on the impressed grid or when the problem
is magnetostatic in the relevant band \cite{harrington2001,stratton1941}.

Linking to Subsection~\ref{sec:ch201}, radiation acts on closed solutions:
\begin{equation}
\label{eq:ch2.maxwell-radiation-composition}
\mathcal{R}_\omega=\Pi_{\mathrm{rad}}\circ\mathcal{S}_\omega
=\Pi_{\mathrm{rad}}\circ\mathcal{L}_\omega^{-1}\circ
\begin{bmatrix}\mathbf{s}_\omega(\cdot)\\ \mathbf{g}_\Gamma\end{bmatrix},
\end{equation}
with \(\mathcal{L}_\omega^{-1}\) understood on the outgoing-admissible branch fixed by
\(\mathcal{B}_\Gamma\) and with \(\Pi_{\mathrm{rad}}\) from
Eq.~\eqref{eq:ch2.radiation-operator-def} \cite{stratton1941}. Closure must precede
radiative selection: applying \(\Pi_{\mathrm{rad}}\) to non-closed states mixes
reactive content into far-zone observables \cite{harrington2001}.

The energy form supports an integration-by-parts identity
\begin{equation}
\label{eq:ch2.maxwell-adjoint-pairing}
\mathfrak{a}_\omega(\mathbf{u}_\omega,\mathbf{v}_\omega)=
\langle\mathcal{M}_\omega[\mathbf{u}_\omega],\mathbf{v}_\omega\rangle_{\Omega}
+\int_\Gamma\big(\hat{\mathbf{n}}\times\E\cdot\H'-\hat{\mathbf{n}}\times\E'\cdot\H\big)dS,
\end{equation}
where \(\langle\cdot,\cdot\rangle_{\Omega}\) is the componentwise \(L^{2}\) pairing
\cite{stratton1941}. Self-adjoint truncations on bounded domains require that
\(\mathcal{B}_\Gamma\) cancel the boundary integrand on admissible traces; criteria
are supplied in Section~\ref{sec:ch214}. Equation
\eqref{eq:ch2.maxwell-adjoint-pairing} is the local spine for those arguments.

Fredholm structure preview: for bounded \(\Omega\) with passive
\(\mathcal{B}_\Gamma\), \(\mathcal{L}_\omega\) is a Fredholm operator on
\(\mathcal{D}(\mathcal{L}_\omega)\) modulo cavity modes, with index tied to
topology and boundary class. Exterior domains replace discrete spectra with
continuous radiation modes, and \(\mathcal{L}_\omega^{-1}\) becomes a resolvent on
the outgoing subspace rather than an ordinary inverse. Section~\ref{sec:ch263}
develops the spectral/Fredholm home; here the point is that invertibility questions
for \(\mathcal{S}_\omega\) are questions about \(\mathcal{L}_\omega\) on the
\emph{declared} domain, not about formal elimination of \(\H\) alone
\cite{harrington2001,stratton1941}.

Linearity of \(\mathcal{L}_\omega\) in \(\mathbf{u}_\omega\) implies that
superposition holds for fixed boundary class: if
\(\mathcal{L}_\omega[\mathbf{u}_\omega^{(a)}]=[\mathbf{s}^{(a)},\mathbf{g}^{(a)}]\)
and \(\mathcal{L}_\omega[\mathbf{u}_\omega^{(b)}]=[\mathbf{s}^{(b)},\mathbf{g}^{(b)}]\)
with identical \(\mathcal{B}_\Gamma\), then
\(\mathcal{L}_\omega[\mathbf{u}_\omega^{(a)}+\mathbf{u}_\omega^{(b)}]=
[\mathbf{s}^{(a)}+\mathbf{s}^{(b)},\mathbf{g}^{(a)}+\mathbf{g}^{(b)}]\).
Changing \(\mathcal{B}_\Gamma\) while holding \(\Jimp\) fixed breaks this
principle because boundary rows are part of the operator, not of the drive.
Laboratory comparisons that swap terminations without re-solving the coupled
system therefore compare different \(\mathcal{L}_\omega\), not different
``modes'' of the same map \cite{stratton1941,harrington2001}. The radiation
operator \(\mathcal{R}_\omega\) inherits the same sensitivity: composition
Eq.~\eqref{eq:ch2.maxwell-radiation-composition} is linear in impressed data
only after \(\mathcal{B}_\Gamma\) and the outgoing branch are held fixed.

\paragraph{Worked illustration (port termination).}
Consider a coaxial feed segment \(\Omega_f\) terminating on a perfectly conducting
aperture plane \(\Gamma_a\) that opens into an exterior region \(\Omega_{\mathrm{ext}}\)
with outgoing closure on a spherical surface \(\Gamma_\infty\). Let
\(\Omega=\Omega_f\cup\Omega_{\mathrm{ext}}\) with \(\Gamma=\partial\Omega_f\cup\Gamma_\infty\),
excluding the aperture from the outer boundary because \(\Gamma_a\) is an interior
interface between feed and free space. Impressed current \(\Jimp\) is supported on
the inner conductor \(C_{\mathrm{in}}\subset\Omega_f\) only, with \(\rho=0\) on the
feed segment so Eq.~\eqref{eq:ch2.source-compatibility} holds identically. On the
outer shield \(\Gamma_{\mathrm{sh}}\subset\partial\Omega_f\), set
\(\mathbf{g}_\Gamma=[\mathbf{0},\hat{\mathbf{n}}\times\H_{\mathrm{shield}}]^{\mathsf{T}}\)
with \(\hat{\mathbf{n}}\times\E=\mathbf{0}\) (perfect conductor). Across
\(\Gamma_a\), continuity of tangential \(\E\) and \(\H\) replaces a Dirichlet
row with interface matching; in the composite notation of
Eq.~\eqref{eq:ch2.maxwell-boundary-composite}, those jumps appear inside
\(\mathcal{B}_\Gamma\) as continuity constraints rather than as impressed data.

With outgoing data on \(\Gamma_\infty\), Eq.~\eqref{eq:ch2.closed-maxwell-bvp}
selects a unique branch of \(\mathcal{S}_\omega[\Jimp]\) in
\(\Omega_{\mathrm{ext}}\); the near-zone reactive fields in \(\Omega_f\) store
energy exchanged with the aperture admittance. Replacing \(\Gamma_a\) by a
short-circuit plane \(\Gamma_{\mathrm{sc}}\) changes the boundary stack while
\(\mathbf{s}_\omega(\Jimp)\) is unchanged, so \(\mathcal{L}_\omega\) itself
is altered. The same \(\Jimp\) then yields a different \(\mathbf{u}_\omega\),
different stored reactive energy in \(\Omega_f\), and different
\(\mathcal{R}_\omega[\Jimp]\) on \(\Gamma_\infty\). Pattern and flux reports
are therefore properties of the closed operator and termination geometry, not
of \(\Jimp\) alone \cite{stratton1941,harrington2001}. If the laboratory
measurement fixes \(\Gamma_{\mathrm{sc}}\) but the simulation leaves
\(\Gamma_a\) open with an approximate absorbing boundary, the computed
\(\mathcal{R}_\omega\) targets a different physical branch and no amount of
post-processing on \(\E\) alone can restore agreement without revisiting
\(\mathcal{B}_\Gamma\).

\paragraph{Closing summary.}
Harmonic Maxwell theory in operator form is the closed map
Eq.~\eqref{eq:ch2.maxwell-boundary-composite} driven by
Eq.~\eqref{eq:ch2.maxwell-forcing-vector} and solved through
Eq.~\eqref{eq:ch2.solution-map-resolvent}. State
Eq.~\eqref{eq:ch2.maxwell-state-vector} preserves first-order curl coupling;
Eqs.~\eqref{eq:ch2.maxwell-energy-form}--\eqref{eq:ch2.maxwell-weak-form} supply
variational access; and Eq.~\eqref{eq:ch2.maxwell-radiation-composition} places
Eq.~\eqref{eq:ch2.radiation-operator-def} on the correct branch of
\(\mathcal{L}_\omega^{-1}\). Function-space, Green, and time--frequency refinements
follow in Sections~\ref{sec:ch21}--\ref{sec:ch23}.

\subsection{Continuity with Chapter 1}
\label{sec:ch203}
% TOC tag: narrative bridge (imports §1.7 by reference; no [USE] tag)

Chapter~\ref{ch:02} does not reopen the Maxwell, transport, covariance, gauge, or
measurement proofs assembled in Chapter~\ref{ch:01}. Subsections~\ref{sec:ch201}
and~\ref{sec:ch202} introduced the harmonic maps \(\mathcal{S}_\omega\),
\(\mathcal{R}_\omega\), and the closed operator \(\mathcal{L}_\omega\); the present
subsection records how those maps inherit the Chapter~\ref{ch:01} admissibility
ledger of Section~\ref{sec:ch01.7} without recursive re-derivation
\cite{stratton1941}. The task is bookkeeping at the operator level: which
Chapter~\ref{ch:01} screens must remain satisfied when impressed data traverse
Eqs.~\eqref{eq:ch2.pipeline-composition} and
\eqref{eq:ch2.maxwell-radiation-composition}, and which Chapter~\ref{ch:02}
constructions are permitted to extend the inherited base.

Let \(\mathcal{R}_{\mathrm{ch1}}\) denote the result set established in
Sections~\ref{sec:ch01.1}--\ref{sec:ch01.6} under the chapter criterion
Eq.~\eqref{eq:ch1.inheritance-admissibility}. Section~\ref{sec:ch01.7.1} packages
that criterion as \(\mathcal{D}_{\mathrm{ch1}}=1\) if and only if
\(\mathcal{D}_{J}=1\), \(\mathcal{D}_{\mathrm{org}}=1\), \(\mathcal{D}_{\mathrm{g}}=1\),
and \(\mathcal{D}_{\mathrm{op}}=1\) simultaneously; the definitions of those four
flags, the master chain
Eq.~\eqref{eq:ch1.master-admissibility-chain}, and the corollary table
Table~\ref{tab:ch1.admissibility-corollaries} remain authoritative in
Chapter~\ref{ch:01} and are invoked here only by reference
\cite{stratton1941}. Chapter~\ref{ch:02} may add operator domains, resolvents,
Green kernels, and spectral projectors, but any quantity presented as a
transport-valid or measurement-valid radiation report must remain compatible with
\(\mathcal{R}_{\mathrm{ch1}}\) in the sense of
Eqs.~\eqref{eq:ch1.bridge-inheritance-map} and
\eqref{eq:ch1.bridge-compatibility-criteria}.

\begin{figure}[t]
\centering
\ChOneFigBlock{%
\begin{tikzpicture}[ch1fig, node distance=7mm and 10mm]
  \node[ch1box] (r) {$\mathcal{R}_{\mathrm{ch1}}$\\{\scriptsize §1.7.1}};
  \node[ch1box, below=9mm of r] (d) {$\mathcal{D}_{\mathrm{ch1}}=1$};
  \node[ch1box, right=16mm of d] (s) {$\mathcal{S}_\omega$};
  \node[ch1box, right=14mm of s] (l) {$\mathcal{L}_\omega$};
  \node[ch1box, right=14mm of l] (rad) {$\mathcal{R}_\omega$};
  \node[ch1box, below=11mm of rad] (obs) {Qualified observables};
  \draw[ch1flow] (r) -- (d);
  \draw[ch1flow] (d) -- (s) -- (l) -- (rad) -- (obs);
\end{tikzpicture}%
}
\caption{Continuity stack: Chapter~\ref{ch:01} results screen impressed data before
operator evaluation; \(\mathcal{R}_\omega\) outputs enter measurement and transport
reports only after the inherited admissibility chain remains satisfied.}
\label{fig:ch2.ch1-continuity-stack}
\end{figure}

Figure~\ref{fig:ch2.ch1-continuity-stack} is the reference architecture for the
remainder of Volume~I. The left block is not re-proved in Chapter~\ref{ch:02}; it
supplies the logical license under which \(\mathcal{S}_\omega\) and
\(\mathcal{L}_\omega\) may be evaluated. The right block is the new operator
machinery; its outputs are physically interpretable only when the left block's
screens remain unviolated \cite{harrington2001,stratton1941}. Confusing the two
blocks---treating an operator output as automatically transport-qualified because
Maxwell residuals vanish---reproduces the category error already diagnosed in
Section~\ref{sec:ch01.0.1} through its three-screen criterion.

Operator domains inherit the source screen without redefining it. Membership in
\(\mathcal{J}_\Omega\) was tied to \(\mathcal{D}_{J}=1\) in
Eq.~\eqref{eq:ch2.source-admissibility}; that packaging restates charge continuity,
finite-energy support, and boundary compatibility already fixed in
Section~\ref{sec:ch01.1.2} and Eq.~\eqref{eq:ch1.source-admissibility}. The
Chapter~\ref{ch:02} continuity relation is therefore
\begin{equation}
\label{eq:ch2.inheritance-operator-domain}
\mathcal{D}_{\mathrm{ch1}}=1
\ \Longrightarrow\
\Jimp\in\mathcal{J}_\Omega
\ \wedge\
\mathcal{S}_\omega[\Jimp]\in\mathcal{F}_\Omega,
\end{equation}
with \(\mathcal{F}_\Omega\) as declared in Subsection~\ref{sec:ch201}. Equation
\eqref{eq:ch2.inheritance-operator-domain} does not extend admissibility; it
records that the operator domains introduced in Chapter~\ref{ch:02} are
\emph{subordinate} to the Chapter~\ref{ch:01} criterion. Evaluating
\(\mathcal{L}_\omega\) on tuples that fail Eq.~\eqref{eq:ch2.source-admissibility}
produces formal vectors that need not correspond to any laboratory excitation, even
when a numerical solver returns a finite residual norm \cite{stratton1941}.

Transport qualification applies to radiative outputs, not to arbitrary components of
\(\mathcal{S}_\omega[\Jimp]\). Section~\ref{sec:ch01.2.3} fixed the three-part
transport test and the regime gate of Section~\ref{sec:ch01.2.4}; importing those
results at the operator level yields
\begin{equation}
\label{eq:ch2.transport-qualified-radiation}
\mathcal{D}_{J,\mathrm{reg}}=1
\ \Longrightarrow\
\text{angular-flux reports attach to}\ \mathcal{R}_\omega[\Jimp],
\ \text{not to}\ \mathcal{S}_\omega[\Jimp]\ \text{by default},
\end{equation}
where \(\mathcal{D}_{J,\mathrm{reg}}\) is the regime-qualified flag from
Eq.~\eqref{eq:ch1.regime-qualified-decision} and
Table~\ref{tab:ch1.admissibility-corollaries}. Equation
\eqref{eq:ch2.transport-qualified-radiation} is the operator restatement of the
reactive-radiative partition imported in Subsection~\ref{sec:ch201}: power and
angular-momentum integrals taken on \(\mathcal{S}_\omega[\Jimp]\) without outgoing
selection can include stored exchange with terminations and near-zone reactive
structure. Only after \(\Pi_{\mathrm{rad}}\) in
Eq.~\eqref{eq:ch2.radiation-operator-def} does the field pair enter the subspace on
which the transport index of Eq.~\eqref{eq:ch1.transport-admissibility-index} was
designed \cite{stratton1941,harrington2001}.

\begin{figure}[t]
\centering
\ChOneFigBlock{%
\begin{tikzpicture}[ch1fig, node distance=6mm and 8mm]
  \node[ch1box] (dj) {$\mathcal{D}_{J}$};
  \node[ch1box, right=10mm of dj] (do) {$\mathcal{D}_{\mathrm{org}}$};
  \node[ch1box, right=10mm of do] (dg) {$\mathcal{D}_{\mathrm{g}}$};
  \node[ch1box, right=10mm of dg] (dop) {$\mathcal{D}_{\mathrm{op}}$};
  \node[ch1box, below=12mm of do] (gate) {Operator evaluation gate};
  \draw[ch1flow] (dj) -- (gate);
  \draw[ch1flow] (do) -- (gate);
  \draw[ch1flow] (dg) -- (gate);
  \draw[ch1flow] (dop) -- (gate);
\end{tikzpicture}%
}
\caption{Four inherited screens from Eq.~\eqref{eq:ch1.inheritance-admissibility}
must remain satisfied before Chapter~\ref{ch:02} operator outputs enter transport,
covariance, gauge, or measurement reports.}
\label{fig:ch2.admissibility-operator-gates}
\end{figure}

Figure~\ref{fig:ch2.admissibility-operator-gates} separates the four inherited
flags that Section~\ref{sec:ch01.7.3} requires to remain unviolated. Each flag
couples to a distinct operator-level test in Chapter~\ref{ch:02}.

For \(\mathcal{D}_{\mathrm{org}}\), Section~\ref{sec:ch01.3.3} fixed origin-shift
rules for torque and angular-flux reports. Under a rotation \(R\) and origin shift
\(\mathbf{a}\), let \(\mathcal{T}_{R,\mathbf{a}}\) denote the induced action on
field pairs. Covariance of the radiation pipeline requires
\begin{equation}
\label{eq:ch2.covariance-operator-test}
\mathcal{R}_\omega\big[\mathcal{T}_{R,\mathbf{a}}\Jimp\big]
=
\mathcal{T}_{R,\mathbf{a}}\,\mathcal{R}_\omega[\Jimp]
\quad\text{on admissible data,}
\end{equation}
with residual \(\mathcal{R}_{\mathcal{F}}\) from
Eq.~\eqref{eq:ch1.covariance-residual} serving as the numerical diagnostic when
discretization breaks Eq.~\eqref{eq:ch2.covariance-operator-test}. Applying a
rotation to \(\Jimp\) but reporting angular momentum from an unrotated
\(\mathcal{S}_\omega[\Jimp]\) violates \(\mathcal{D}_{\mathrm{org}}\) even if
far-zone power from \(\mathcal{R}_\omega[\Jimp]\) appears correct
\cite{stratton1941}. Chapter~\ref{ch:02} supplies the field maps; Chapter~\ref{ch:01}
retains ownership of the covariance criterion.

For \(\mathcal{D}_{\mathrm{g}}\), Section~\ref{sec:ch01.4.1} showed that local
phase and longitudinal content can be gauge-sensitive while far-zone radiative
observables remain invariant on the outgoing subspace. At the operator level,
\begin{equation}
\label{eq:ch2.gauge-observable-composition}
\mathcal{Q}_{\mathrm{gauge-safe}}
=
\mathcal{O}_{\mathrm{far}}\circ\mathcal{R}_\omega,
\qquad
\mathcal{Q}_{\mathrm{gauge-safe}}\ \text{independent of}\ \chi
\ \text{on}\ \mathcal{F}_{\mathrm{rad}},
\end{equation}
where \(\mathcal{O}_{\mathrm{far}}\) is the far-zone readout functional introduced in
Subsection~\ref{sec:ch201} and \(\chi\) denotes gauge potentials used only in
intermediate calculations. Applying \(\mathcal{O}_{\mathrm{far}}\) directly to
\(\mathcal{S}_\omega[\Jimp]\) without radiative selection can reproduce
gauge-dependent near-zone structure in quantities advertised as observables. The
helicity screen \(\mathcal{D}_{h}\) of Section~\ref{sec:ch01.4.3} further restricts
which spin-resolved readouts remain meaningful when polarization is decomposed
\cite{stratton1941}.

For \(\mathcal{D}_{\mathrm{op}}\), Section~\ref{sec:ch01.5} classified reports as
direct observables, constrained surrogates, or inference-only quantities through
Eq.~\eqref{eq:ch1.report-class-lock}. The operator continuation is
\begin{equation}
\label{eq:ch2.operational-measurement-chain}
\Lambda_{\mathrm{lab}}
=
\mathcal{O}_{\mathrm{meas}}\circ\mathcal{O}_{\mathrm{far}}\circ\mathcal{R}_\omega,
\qquad
\mathcal{Q}\in\{\text{direct},\ \text{surrogate},\ \text{inference}\},
\end{equation}
where \(\mathcal{O}_{\mathrm{meas}}\) encodes finite aperture, calibration, and
modal truncation of Section~\ref{sec:ch01.5}. Equation
\eqref{eq:ch2.operational-measurement-chain} makes explicit that a pattern or
torque reading is a composition of maps, not a property of \(\Jimp\) alone. Rank
deficiency in \(\mathcal{O}_{\mathrm{meas}}\) produces surrogate reports even when
\(\mathcal{D}_{\mathrm{op}}=1\) for the underlying continuous observable
\cite{stratton1941,harrington2001}.

Section~\ref{sec:ch01.7.2} declared which topics require operator-domain proof in
Chapter~\ref{ch:02}, representation closure in Chapter~\ref{ch:03}, or realizability
limits in Chapter~\ref{ch:04}. The defer rule of
Eqs.~\eqref{eq:ch1.defer-topic-rule}--\eqref{eq:ch1.defer-reporting-lock} carries
forward unchanged:
\begin{equation}
\label{eq:ch2.defer-boundary-declaration}
\mathcal{X}\in\mathcal{D}_{\mathrm{defer}}
\ \Longrightarrow\
\text{no local proof of}\ \mathcal{X}\ \text{in Chapter~\ref{ch:02} unless TOC-owned},
\end{equation}
so Green-function construction, Sommerfeld uniqueness, Fredholm index computation,
and non-radiating null-space characterization are developed in their assigned
Sections~\ref{sec:ch23}--\ref{sec:ch27}, not by replaying Chapter~\ref{ch:01}
arguments \cite{stratton1941}. Conversely, Chapter~\ref{ch:02} must not import
deferred topics from later chapters as if already established; vector spherical
harmonics, Debye potentials, and gauge transformations remain outside the present
file boundary \cite{harrington2001}.

Non-recursion is the central compatibility constraint of
Eq.~\eqref{eq:ch1.bridge-inheritance-map}. Chapter~\ref{ch:02} constructions
\(\mathcal{Y}\) are admissible extensions when
\begin{equation}
\label{eq:ch2.bridge-nonrecursion}
\mathcal{D}_{\mathrm{bridge}}(\mathcal{Y})=1
\iff
\mathcal{Y}\prec\mathcal{R}_{\mathrm{ch1}}
\ \wedge\
\mathcal{Y}\ \text{adds no proof of}\ \mathcal{R}_{\mathrm{ch1}},
\end{equation}
with \(\prec\) as in Eq.~\eqref{eq:ch1.bridge-compatibility-criteria}. Packaging
Maxwell residuals as \(\mathcal{L}_\omega[\mathbf{u}_\omega]=\mathbf{s}_\omega\)
satisfies Eq.~\eqref{eq:ch2.bridge-nonrecursion} because it cites
Eq.~\eqref{eq:ch1.maxwell-operator-block} rather than re-deriving
Eq.~\eqref{eq:ch1.maxwell-local}. Re-displaying the full Maxwell derivation,
re-proving angular-momentum conservation, or introducing gauge freedom as a new
Chapter~\ref{ch:02} topic would violate the bridge lock and the Volume~I
no-recursion policy for Maxwell, gauge, and conservation content
\cite{stratton1941}.

The master admissibility chain of Eq.~\eqref{eq:ch1.master-admissibility-chain}
composes with the radiation pipeline as
\begin{equation}
\label{eq:ch2.continuity-composition-law}
\mathfrak{A}_{\mathrm{ch2}}^{\mathrm{(report)}}
=
\mathfrak{A}_{\mathrm{meas}}
\circ
\mathfrak{A}_{\mathrm{g}}
\circ
\mathfrak{A}_{\mathrm{cov}}
\circ
\mathfrak{A}_{\mathrm{tr}}
\circ
\Pi_{\mathrm{rad}}
\circ
\mathcal{S}_\omega,
\end{equation}
where \(\mathfrak{A}_{\mathrm{tr}}\) implements the transport and regime screens
\(\mathcal{D}_{J}\) and \(\mathcal{D}_{J,\mathrm{reg}}\), and the final four
\(\mathfrak{A}\)-maps are the inherited layers from
Table~\ref{tab:ch1.admissibility-corollaries}. Equation
\eqref{eq:ch2.continuity-composition-law} is the Chapter~\ref{ch:02} statement of
what ``continuity with Chapter~\ref{ch:01}'' means operationally: every report
downstream of an antenna model is a composition ending in \(\mathcal{S}_\omega\), but
only after \(\Pi_{\mathrm{rad}}\) and the inherited screens does the composition
return a quantity licensed for angular-transport interpretation
\cite{stratton1941,harrington2001}. Omitting any factor---evaluating
\(\mathfrak{A}_{\mathrm{tr}}\) on reactive fields, skipping
\(\mathfrak{A}_{\mathrm{g}}\) on a gauge-sensitive surrogate, or applying
\(\mathfrak{A}_{\mathrm{meas}}\) without declaring aperture rank deficiency---breaks
the chain even when \(\mathcal{L}_\omega[\mathbf{u}_\omega]=\mathbf{0}\) in the
interior \cite{stratton1941}.

Report classes from Eq.~\eqref{eq:ch1.report-class-lock} attach to the \emph{final}
output of Eq.~\eqref{eq:ch2.continuity-composition-law}, not to intermediate operator
values. A direct observable at the measurement stage can become inference-only after
\(\mathcal{O}_{\mathrm{meas}}\) if the modal set does not resolve the targeted
angular harmonic \cite{stratton1941}. Chapter~\ref{ch:02} therefore treats
\(\mathcal{R}_\omega[\Jimp]\) as the last field-theoretic object before
measurement composition; Sections~\ref{sec:ch25}--\ref{sec:ch29} develop flux
integrals and angular-spectrum actions on that object, always subject to the inherited
report class \cite{harrington2001}.

\paragraph{Worked illustration (admissibility failure before operators).}
Suppose \(\Jimp\) is a numerically specified filament current with
\(\divg\Jimp\neq-\jj\omega\rho\) because charge density was set to zero
independently on a volumetric grid. Then \(\mathcal{D}_{J}=0\) in the sense of
Eq.~\eqref{eq:ch2.source-admissibility}, even though a finite-element solve of
Eq.~\eqref{eq:ch2.closed-maxwell-bvp} may return a smooth \((\E,\H)\) pair. Applying
\(\mathcal{R}_\omega\) and then an angular-momentum flux functional from
Section~\ref{sec:ch01.1.5} produces a number that cannot inherit conservation
interpretation because the source tuple failed before Maxwell closure was physically
posed. Correcting \(\rho\) to restore
Eq.~\eqref{eq:ch1.charge-continuity-harmonic} changes \(\mathbf{s}_\omega\) in
Eq.~\eqref{eq:ch2.maxwell-forcing-vector} and therefore changes
\(\mathcal{S}_\omega[\Jimp]\), but leaves the \emph{form} of
Eqs.~\eqref{eq:ch2.continuity-composition-law} and
\eqref{eq:ch2.transport-qualified-radiation} unchanged. The illustration shows that
continuity with Chapter~\ref{ch:01} is enforced at the input screen, not by
post-hoc adjustment of operator outputs \cite{stratton1941,harrington2001}.

\paragraph{Worked illustration (regime error after operators).}
Next suppose \(\mathcal{D}_{J}=1\) but the device operates in a storage-dominated
near zone with \(kr\ll 1\) on the measurement surface while an angular-momentum
report applies far-zone transport formulas. Then \(\mathcal{D}_{J,\mathrm{reg}}=0\)
(Eq.~\eqref{eq:ch1.regime-qualified-decision}), and
Eq.~\eqref{eq:ch2.transport-qualified-radiation} forbids attaching transport meaning
to \(\mathcal{R}_\omega[\Jimp]\) on that surface even though Maxwell residuals vanish
and \(\mathcal{R}_\omega\) is well defined. Moving the evaluation surface outward
until Eq.~\eqref{eq:ch1.radiative-regime-condition} holds restores regime
qualification without altering \(\mathcal{L}_\omega\). Regime and operator closure
are independent gates: satisfying one does not substitute for the other
\cite{stratton1941}.

\paragraph{Operator symbol survival (Chapter~\ref{ch:02} ledger).}
Table~\ref{tab:ch2.operator-symbol-survival} records symbols that survive the full chapter
pipeline; downstream chapters relabel coordinates but must not alter the composition order
\cite{stratton1941,harrington2001}. Rigged-Hilbert and measure normalization for continuous
spectra are deferred to Chapter~\ref{ch:03} Subsections~\ref{sec:ch313} and~\ref{sec:ch352}
\cite{stratton1941}.

\begin{table}[t]
\centering
\caption{Operator symbol survival table for Chapter~\ref{ch:02}: domain, range, and report stage.
Symbols in the left column are authoritative in \texttt{00-NOTATION.md} and must not be overloaded.}
\label{tab:ch2.operator-symbol-survival}
\small
\begin{tabular}{@{}lll@{}}
\toprule
Symbol & Domain $\to$ range & Report role \\
\midrule
$\mathcal{S}_\omega$ & $\mathcal{J}_\Omega\to\mathcal{F}_\Omega$ & Maxwell closure; not transport-default \\
$\Pi_{\mathrm{prop}},\Pi_{\mathrm{ev}},\Pi_{\mathrm{non}}$ & $\mathcal{F}_\Omega\to\mathcal{F}_\Omega$ & Spectral split (formal until §3.1.3) \\
$\Pi_{\mathrm{rad}}=\mathcal{O}_{\mathrm{out}}\circ\Pi_{\mathrm{prop}}$ & $\mathcal{F}_\Omega\to\mathcal{F}_{\mathrm{rad}}$ & Outgoing export selector \\
$\mathcal{R}_\omega$ & $\mathcal{J}_\Omega\to\mathcal{F}_{\mathrm{rad}}$ & Radiation map; infinite $\ker$ \\
$\Lambda_{\mathrm{nf}}$ & $\mathcal{J}_\Omega\to a(\widehat{\mathbf{k}})$ & Near-to-far angular amplitude \\
$\Lambda_{\mathrm{rad}},\mathcal{M}_\omega$ & $\mathcal{J}_\Omega\to$ observables & Far readout; finite rank \\
$a,a_{\mathrm{ret}},a_{\mathrm{sup}},a_{\mathrm{acc}},a_{\mathrm{obs}}$ & on $S^{2}$ & Angular content stages \\
$\mathcal{D}_{J,\mathrm{reg}},\mathcal{D}_{\mathrm{ch2}}$ & flags & Regime and chapter gates \\
\bottomrule
\end{tabular}
\end{table}

\paragraph{Closing summary.}
Continuity with Chapter~\ref{ch:01} means that
Eqs.~\eqref{eq:ch2.inheritance-operator-domain},
\eqref{eq:ch2.transport-qualified-radiation}, and
\eqref{eq:ch2.continuity-composition-law} screen every use of
\(\mathcal{S}_\omega\), \(\mathcal{L}_\omega\), and \(\mathcal{R}_\omega\) against
\(\mathcal{R}_{\mathrm{ch1}}\) imported from
Sections~\ref{sec:ch01.7.1}--\ref{sec:ch01.7.3}. Chapter~\ref{ch:02} adds operator
machinery; it does not relax transport, covariance, gauge, or measurement discipline.
Sections~\ref{sec:ch21} onward refine domains and kernels while preserving the
composition order of Figure~\ref{fig:ch2.ch1-continuity-stack}.

With Subsections~\ref{sec:ch201}--\ref{sec:ch203} in place, Section~\ref{sec:ch20}
has fixed the radiation-map viewpoint, packaged harmonic Maxwell theory as a closed
operator \(\mathcal{L}_\omega\), and recorded the Chapter~\ref{ch:01} inheritance
screens that govern every subsequent map in the chapter. The technical development
now proceeds from function-space foundations through Green kernels, outgoing
selection, radiation integrals, spectral structure, null spaces, reciprocity, and
angular-spectrum operators---each subsection importing prior results only by
reference and honoring the composition law of
Eq.~\eqref{eq:ch2.continuity-composition-law}.

\section{Function Spaces and Maxwell Operators}
\label{sec:ch21}

Section~\ref{sec:ch20} introduced the maps \(\mathcal{S}_\omega\),
\(\mathcal{R}_\omega\), and the closed operator \(\mathcal{L}_\omega\) on
minimally declared spaces \(\mathcal{F}_\Omega\) and \(\mathcal{J}_\Omega\).
Those declarations were sufficient to state input--output relations and to record
continuity with Chapter~\ref{ch:01}, but they are not yet the function-space
apparatus required for trace theorems, variational stability, spectral
compression, or rank-deficient measurement models. Section~\ref{sec:ch21} supplies
that apparatus: Hilbert and Sobolev classes for fields and sources, distributional
extensions for impressed currents and interface data, and the operator-theoretic
vocabulary---domain, range, adjoint, boundedness, compactness, closure---through
which Maxwell maps become objects amenable to Fredholm analysis and angular
spectral theory in Sections~\ref{sec:ch26}--\ref{sec:ch29}
\cite{stratton1941,harrington2001}.

The finite-energy criterion of Eq.~\eqref{eq:ch1.energy-class} and the admissible
state class Eq.~\eqref{eq:ch1.admissible-space} from
Subsection~\ref{sec:ch01.0.3} remain the physical floor. Section~\ref{sec:ch21}
does not weaken or replace them; it embeds them in spaces where curl, divergence,
and tangential trace maps are closed operations. Subsection~\ref{sec:ch01.1.2}
fixed source tuples, boundary operators, and charge continuity for impressed data;
Subsection~\ref{sec:ch211} refines the corresponding distribution and trace spaces
without re-opening those proofs \cite{stratton1941}. The guiding constraint from
Subsection~\ref{sec:ch203} persists: every space declaration must remain compatible
with Eq.~\eqref{eq:ch2.inheritance-operator-domain} and must not introduce
recursive proofs of \(\mathcal{R}_{\mathrm{ch1}}\).

\begin{figure}[t]
\centering
\ChOneFigBlock{%
\begin{tikzpicture}[ch1fig, node distance=7mm and 9mm]
  \node[ch1box] (adm) {Admissible tuples\\{\scriptsize §1.0.3, §1.1.2}};
  \node[ch1box, right=14mm of adm] (sp) {Field/source\\Hilbert spaces};
  \node[ch1box, right=14mm of sp] (op) {Operator\\domain/range};
  \node[ch1box, right=14mm of op] (spec) {Compactness\\spectral tools};
  \node[ch1box, below=12mm of op] (max) {$\mathcal{L}_\omega,\ \mathcal{S}_\omega$};
  \draw[ch1flow] (adm) -- (sp) -- (op) -- (spec);
  \draw[ch1flow] (sp) -- (max);
  \draw[ch1flow] (op) -- (max);
\end{tikzpicture}%
}
\caption{Section~\ref{sec:ch21} development path: admissibility classes from
Chapter~\ref{ch:01} are embedded in Hilbert/Sobolev spaces, operator domains are
sharpened, and compactness/spectral structure is prepared for
Sections~\ref{sec:ch23}--\ref{sec:ch26}.}
\label{fig:ch2.function-space-roadmap}
\end{figure}

Figure~\ref{fig:ch2.function-space-roadmap} orders the four subsections that follow.
Subsection~\ref{sec:ch211} defines source, field, and dual spaces for phasor
Maxwell theory, including distributional sources and trace spaces on
\(\Gamma=\partial\Omega\), importing finite-energy admissibility from
Section~\ref{sec:ch01.0.3} and source--boundary coupling from
Section~\ref{sec:ch01.1.2}. Subsection~\ref{sec:ch212} treats linear operators on
those spaces: domain, range, graph norm, and adjoint pairing compatible with
Eq.~\eqref{eq:ch2.maxwell-adjoint-pairing}. Subsection~\ref{sec:ch213} develops
boundedness and compactness criteria, rank deficiency of observation maps, and the
spectral compression viewpoint previewed for radiation operators in
Section~\ref{sec:ch26}. Subsection~\ref{sec:ch214} addresses self-adjoint
realizations, Friedrichs and other closures, and the admissibility constraints
under which truncated Maxwell operators inherit coercivity or passivity on bounded
domains \cite{harrington2001,stratton1941}.

Several constructions deliberately remain outside Section~\ref{sec:ch21}. Dyadic
Green kernels and retarded inverses belong to Section~\ref{sec:ch23}; Sommerfeld
outgoing selection and exterior uniqueness to Section~\ref{sec:ch24}; radiation
integrals and flux operators at the radiative boundary to Section~\ref{sec:ch25}.
Section~\ref{sec:ch21} prepares the domains on which those maps act. Attempting to
evaluate \(\mathcal{L}_\omega^{-1}\) through a Green function before trace spaces
and dual pairings are fixed produces kernel formulas that ignore boundary rows and
corrupt flux identities \cite{stratton1941}. Conversely, staying at the
\(L^{2}\)-level of Eq.~\eqref{eq:ch2.operator-domain-declaration} without
Sobolev refinement is adequate for uniqueness statements on smooth domains but
insufficient for sheet sources, impedance boundaries, and finite-element
discretizations where \(\curl\) and \(\hat{\mathbf{n}}\times\) must be closed
simultaneously \cite{harrington2001}.

The energy bilinear form Eq.~\eqref{eq:ch2.maxwell-energy-form} and weak
formulation Eq.~\eqref{eq:ch2.maxwell-weak-form} from Subsection~\ref{sec:ch202}
anchor the Hilbert-space choices. Section~\ref{sec:ch21} identifies the minimal
completion in which \(\mathfrak{a}_\omega\) is continuous and coercive modulo
cavity modes, and in which the forcing functional \(\mathfrak{f}_\omega\) induced
by \(\mathbf{s}_\omega(\mathfrak{D})\) is bounded on test states with homogeneous
traces. Well-posedness of the weak problem on bounded domains then follows from
Lax--Milgram arguments once null spaces aligned with divergence-free resonances
are quotiented; the quotient construction is owned by Subsection~\ref{sec:ch214}
and feeds the cavity-mode analysis of Section~\ref{sec:ch263}
\cite{stratton1941}. Exterior domains require outgoing subspaces rather
than full Hilbert self-adjointness; that distinction is preserved throughout so
that interior cavity spectra are not conflated with continuous radiation spectra
\cite{harrington2001}.

Rank deficiency enters at two levels that Section~\ref{sec:ch21} must keep separate.
Observation maps---finite apertures, modal truncations, probe coupling---are
finite-rank operators on field space even when \(\mathcal{L}_\omega\) is
injective on the physical branch. Equivalent source distributions in
\(\ker\mathcal{R}_\omega\) produce identical far-zone patterns while differing in
reactive near fields; that null-space structure is analyzed in
Section~\ref{sec:ch27}, but the compact embedding lemmas used there are developed
in Subsection~\ref{sec:ch213}. Measurement surrogates from
Section~\ref{sec:ch01.5} therefore appear as rank-deficient compositions on the
right of Eq.~\eqref{eq:ch2.continuity-composition-law}, not as modifications of
Maxwell law \cite{stratton1941}.

Readers approaching from numerical methods should expect explicit domain
declarations for every operator symbol. A method-of-moments impedance matrix is a
Galerkin truncation of a boundary operator; a finite-element solution is a
Ritz--Galerkin approximation to Eq.~\eqref{eq:ch2.maxwell-weak-form}. Both are
valid only when the discrete spaces inherit the continuity and trace structure
defined in Subsections~\ref{sec:ch211}--\ref{sec:ch212}. Section~\ref{sec:ch21}
states the continuous targets those discretizations approximate; it does not
substitute for mesh convergence analysis \cite{harrington2001}.

Material inhomogeneity and interface conditions from
Eq.~\eqref{eq:ch1.interface-conditions} enter through piecewise constant
coefficients in \(\mathcal{M}_\omega\) without changing the space template:
fields remain in product Sobolev spaces with jump conditions enforced on traces.
Multi-region formulations in Section~\ref{sec:ch234} rely on the interface trace
classes introduced in Subsection~\ref{sec:ch211}. Perfect conductors impose
homogeneous tangential \(\E\); impedance and Leontovich boundaries modify the
boundary pairing in Eq.~\eqref{eq:ch2.maxwell-adjoint-pairing} and are treated as
closure operators in Subsection~\ref{sec:ch214} \cite{stratton1941}.

Subsection~\ref{sec:ch214} completes the function-space foundation begun here with
self-adjoint truncations and admissibility closures that connect directly to the
spectral/Fredholm program of Section~\ref{sec:ch263}. The subsection sequence
begins with the space definitions that make \(\mathcal{D}(\mathcal{L}_\omega)\) in
Eq.~\eqref{eq:ch2.operator-domain-declaration} precise enough for the Green,
radiation, and spectral sections that follow.


\subsection{Spaces for sources, fields, and distributions}
\label{sec:ch211}
% TOC tag: [HOME][USE SS1.0.3][USE SS1.1.2]

Subsections~\ref{sec:ch201}--\ref{sec:ch202} declared minimal spaces
\(\mathcal{F}_\Omega\) and \(\mathcal{J}_\Omega\) adequate to \emph{name} the maps
\(\mathcal{S}_\omega\), \(\mathcal{R}_\omega\), and \(\mathcal{L}_\omega\).
Subsection~\ref{sec:ch211} supplies the Sobolev, trace, and distributional
structure those declarations presuppose. Finite-energy admissibility from
Subsection~\ref{sec:ch01.0.3} and the source--boundary tuple of
Subsection~\ref{sec:ch01.1.2} remain authoritative; the present subsection embeds
them in spaces where \(\curl\), \(\divg\), and tangential restriction are
continuous maps with explicit domains \cite{stratton1941,harrington2001}. Nothing
here re-proves charge continuity, interface jumps, or outer-boundary classes already
fixed in Chapter~\ref{ch:01}.

Let \(\Omega\subset\mathbb{R}^{3}\) be a bounded Lipschitz domain or a finite
union of Lipschitz patches \(\{\Omega_a\}\) with outward boundary
\(\Gamma=\partial\Omega\). Phasor fields use the \(\exp(-\jj\omega t)\) convention
throughout. Single-vector Sobolev classes are
\begin{equation}
\label{eq:ch2.hcurl-space}
\HCurl(\Omega)
=
\left\{
\mathbf{u}\in L^{2}(\Omega)^{3}:
\curl\mathbf{u}\in L^{2}(\Omega)^{3}
\right\},
\end{equation}
\begin{equation}
\label{eq:ch2.hdiv-space}
\Hdiv(\Omega)
=
\left\{
\mathbf{u}\in L^{2}(\Omega)^{3}:
\divg\mathbf{u}\in L^{2}(\Omega)
\right\},
\end{equation}
each equipped with the graph norm
\(\|\mathbf{u}\|_{\HCurl}^{2}=\|\mathbf{u}\|_{L^{2}}^{2}+\|\curl\mathbf{u}\|_{L^{2}}^{2}\)
and the analogous \(\Hdiv\) norm \cite{stratton1941}. Equations
\eqref{eq:ch2.hcurl-space}--\eqref{eq:ch2.hdiv-space} are the vector analogues of
the finite-energy floor Eq.~\eqref{eq:ch1.energy-class}: they permit Maxwell
curl-coupling without assuming pointwise smoothness of \(\E\) or \(\H\). For
piecewise constant \(\varepsilon,\mu\), fields are understood patchwise in
Eqs.~\eqref{eq:ch2.hcurl-space}--\eqref{eq:ch2.hdiv-space} with interface traces
constrained by Eq.~\eqref{eq:ch1.interface-conditions} \cite{harrington2001}.

The Maxwell state space is the product
\begin{equation}
\label{eq:ch2.maxwell-product-space}
\mathcal{U}_\omega
=
\HCurl(\Omega)\times\HCurl(\Omega),
\qquad
\mathbf{u}_\omega=(\E,\H)\in\mathcal{U}_\omega.
\end{equation}
Equation~\eqref{eq:ch2.maxwell-product-space} refines the six-component packaging
Eq.~\eqref{eq:ch2.maxwell-state-vector}: both components carry square-integrable
curls, matching the first-order structure of
Eq.~\eqref{eq:ch1.maxwell-operator-block}. A weighted energy norm compatible with
Eq.~\eqref{eq:ch2.maxwell-energy-form} is
\begin{equation}
\label{eq:ch2.field-energy-norm}
\|\mathbf{u}_\omega\|_{\mathcal{U}}^{2}
=
\|\E\|_{L^{2}}^{2}+\|\H\|_{L^{2}}^{2}
+\alpha\,\|\curl\E\|_{L^{2}}^{2}+\beta\,\|\curl\H\|_{L^{2}}^{2},
\end{equation}
with positive weights \(\alpha,\beta\) fixed by constitutive scaling in each patch.
Norm~\eqref{eq:ch2.field-energy-norm} controls the bilinear form
\(\mathfrak{a}_\omega\) of Eq.~\eqref{eq:ch2.maxwell-weak-form} on bounded domains
and makes \(\mathcal{U}_\omega\) a Hilbert space \cite{stratton1941}.

\begin{figure}[t]
\centering
\ChOneFigBlock{%
\begin{tikzpicture}[ch1fig, node distance=7mm and 8mm]
  \node[ch1box] (l2) {$L^{2}(\Omega)^{3}$};
  \node[ch1box, above=10mm of l2] (hc) {$\HCurl(\Omega)$};
  \node[ch1box, below=10mm of l2] (hd) {$\Hdiv(\Omega)$};
  \node[ch1box, right=18mm of hc] (tr) {$\mathcal{T}_\Gamma$};
  \draw[ch1flow] (l2) -- (hc);
  \draw[ch1flow] (l2) -- (hd);
  \draw[ch1flow] (hc) -- node[above,font=\scriptsize] {$\gamma_\tau$} (tr);
\end{tikzpicture}%
}
\caption{Embedding ladder for Maxwell fields: \(\HCurl\) and \(\Hdiv\) sit above
\(L^{2}\); tangential traces map \(\HCurl\) fields to boundary data compatible with
Eq.~\eqref{eq:ch1.boundary-data-operator}.}
\label{fig:ch2.sobolev-embedding-ladder}
\end{figure}

Figure~\ref{fig:ch2.sobolev-embedding-ladder} summarizes the analytic hierarchy.
Divergence control enters through source spaces and through Gauss constraints rather
than through requiring both components to lie in \(\Hdiv\); forcing \(\E\in\Hdiv\)
and \(\H\in\Hdiv\) simultaneously is stronger than Maxwell's first-order system
requires and is reserved for scalar-potential reductions deferred to later chapters
\cite{stratton1941}.

Impressed sources refine the declaration of Subsection~\ref{sec:ch201}. Let
\(\Omega_s\subseteq\Omega\) denote compact source support. The admissible impressed
current space is
\begin{equation}
\label{eq:ch2.source-space-refined}
\mathcal{J}_\Omega
=
\left\{
\Jimp\in \HCurl(\Omega):
\operatorname{supp}\Jimp\subseteq\Omega_s,\
\mathcal{D}_{J}[\Jimp]=1
\right\},
\end{equation}
where \(\mathcal{D}_{J}\) is the screen of
Eq.~\eqref{eq:ch2.source-admissibility} inherited from
Eq.~\eqref{eq:ch1.inheritance-admissibility} and
Eq.~\eqref{eq:ch1.source-admissibility} \cite{stratton1941}. Equation
\eqref{eq:ch2.source-space-refined} strengthens square integrability to
\(\HCurl\) regularity so that \(\curl\Jimp\) is defined in \(L^{2}\) whenever
numerical or analytical pipelines differentiate sources. Charge-neutral subclasses
used in radiation null-space analysis satisfy
\begin{equation}
\label{eq:ch2.charge-neutral-source}
\mathcal{J}_\Omega^{(0)}
=
\left\{
\Jimp\in\mathcal{J}_\Omega:
\divg\Jimp=0
\right\},
\end{equation}
with \(\rho=0\) implied on \(\Omega_s\). Equation~\eqref{eq:ch2.charge-neutral-source}
is the subspace on which divergence-free reparametrizations act without altering
\(\mathcal{S}_\omega[\Jimp]\) in simply connected regions; null-space structure is
developed in Section~\ref{sec:ch272} \cite{harrington2001}.

Boundary data from the tuple Eq.~\eqref{eq:ch1.data-specification-tuple} live in
trace spaces rather than in volume Sobolev classes. Define
\begin{equation}
\label{eq:ch2.trace-space-definition}
\mathcal{T}_\Gamma
=
H^{-1/2}(\Gamma)^{3}\times H^{-1/2}(\Gamma)^{3},
\end{equation}
as the product space for tangential electric and magnetic traces induced by
\(\HCurl(\Omega)\) fields on Lipschitz \(\Gamma\) \cite{stratton1941}. The
tangential trace operators extend Eq.~\eqref{eq:ch2.trace-restriction}:
\begin{equation}
\label{eq:ch2.tangential-trace-maps}
\gamma_{\tau,E}:\HCurl(\Omega)\to H^{-1/2}(\Gamma)^{3},
\quad
\gamma_{\tau,E}[\E]=\hat{\mathbf{n}}\times\E|_{\Gamma},
\end{equation}
\begin{equation}
\label{eq:ch2.tangential-trace-magnetic}
\gamma_{\tau,H}:\HCurl(\Omega)\to H^{-1/2}(\Gamma)^{3},
\quad
\gamma_{\tau,H}[\H]=\hat{\mathbf{n}}\times\H|_{\Gamma},
\end{equation}
with outward unit normal \(\hat{\mathbf{n}}\). Equations
\eqref{eq:ch2.tangential-trace-maps}--\eqref{eq:ch2.tangential-trace-magnetic}
make Eq.~\eqref{eq:ch1.boundary-data-operator} an operator between declared
Banach spaces: \(\mathcal{B}_\Gamma:\mathcal{T}_\Gamma\to\mathcal{Y}_\Gamma\) for
data space \(\mathcal{Y}_\Gamma\) fixed by the boundary class of
Eq.~\eqref{eq:ch1.outer-boundary-classes} or outgoing traces of
Eq.~\eqref{eq:ch1.radiation-trace-condition} \cite{stratton1941,harrington2001}.

\begin{figure}[t]
\centering
\ChOneFigBlock{%
\begin{tikzpicture}[ch1fig, node distance=8mm and 10mm]
  \node[ch1box] (vol) {Volume\\$\HCurl(\Omega)$};
  \node[ch1box, right=16mm of vol] (tr) {Traces\\$\mathcal{T}_\Gamma$};
  \node[ch1box, right=16mm of tr] (dat) {Data\\$\mathcal{Y}_\Gamma$};
  \draw[ch1flow] (vol) -- node[above,font=\scriptsize] {$\gamma_\tau$} (tr);
  \draw[ch1flow] (tr) -- node[above,font=\scriptsize] {$\mathcal{B}_\Gamma$} (dat);
\end{tikzpicture}%
}
\caption{Trace--data chain for boundary closure: volume fields map to
\(\mathcal{T}_\Gamma\); the boundary operator from
Eq.~\eqref{eq:ch1.boundary-data-operator} selects admissible data.}
\label{fig:ch2.trace-boundary-chain}
\end{figure}

Figure~\ref{fig:ch2.trace-boundary-chain} is the function-space reading of
Figure~\ref{fig:ch2.maxwell-boundary-stack}. Surjectivity of \(\gamma_\tau\) onto
compatible trace subclasses fails globally---only fields in specified
\(\HCurl\) subspaces with prescribed divergence pairings map to physically realizable
\(\mathbf{g}_\Gamma\) \cite{stratton1941}. Subsection~\ref{sec:ch214} identifies
self-adjoint boundary realizations; here the point is that \(\mathcal{L}_\omega\) of
Eq.~\eqref{eq:ch2.maxwell-boundary-composite} acts on \(\mathcal{U}_\omega\) only when
\(\gamma_\tau[\mathbf{u}_\omega]\) lies in the domain of \(\mathcal{B}_\Gamma\).

Sheet currents and interface channels from Eq.~\eqref{eq:ch1.interface-conditions}
require distributional extensions. Let \(\mathcal{D}'(\Omega)^{3}\) denote
temporal-harmonic distributional three-vector fields. Impressed volume data remain
\(\Jimp\in\mathcal{J}_\Omega\), while surface equivalents \(\mathbf{J}_{s,ab}\) on
\(\Sigma_{ab}\) embed as
\begin{equation}
\label{eq:ch2.distributional-source-extension}
\langle\mathbf{J}_{s},\boldsymbol{\phi}\rangle
=
\sum_{ab}\int_{\Sigma_{ab}}\mathbf{J}_{s,ab}\cdot\boldsymbol{\phi}\,dS,
\qquad
\boldsymbol{\phi}\in C_c^\infty(\Omega)^{3},
\end{equation}
linking the data tuple Eq.~\eqref{eq:ch1.data-specification-tuple} to a continuous
linear functional on smooth test fields \cite{stratton1941}. Equation
\eqref{eq:ch2.distributional-source-extension} does not alter the jump laws of
Eq.~\eqref{eq:ch1.interface-conditions}; it specifies the functional-analytic class
in which those jumps appear in weak formulations. Interface trace spaces on
\(\Sigma_{ab}\) are denoted
\begin{equation}
\label{eq:ch2.interface-trace-space}
\mathcal{T}_{\Sigma_{ab}}
=
H^{-1/2}(\Sigma_{ab})^{3},
\end{equation}
with tangential continuity of \(\E\) and jump of \(\hat{\mathbf{n}}\times\H\) encoded
as constraints on \(\gamma_\tau\) taken from each patch \cite{harrington2001}.

Dual pairings for source and test fields use the \(L^{2}\) inner product
\begin{equation}
\label{eq:ch2.dual-source-pairing}
\langle\Jimp,\boldsymbol{\psi}\rangle_{\Omega}
=
\int_\Omega \Jimp\cdot\boldsymbol{\psi}^{*}\,d^{3}r,
\qquad
\boldsymbol{\psi}\in \HCurl(\Omega),
\end{equation}
which induces a bounded functional on test states when \(\Jimp\in L^{2}(\Omega)^{3}\)
and extends to Eq.~\eqref{eq:ch2.distributional-source-extension} for sheet data.
Forcing in Eq.~\eqref{eq:ch2.maxwell-forcing-vector} enters the weak form
Eq.~\eqref{eq:ch2.maxwell-weak-form} through
\(\mathfrak{f}_\omega(\mathbf{v}_\omega)=\langle\mathbf{s}_\omega,\mathbf{v}_\omega\rangle\)
once \(\mathbf{v}_\omega\in\mathcal{U}_\omega\) carries homogeneous boundary traces
\cite{stratton1941}.

The admissible Maxwell field space compatible with Chapter~\ref{ch:01} is
\begin{equation}
\label{eq:ch2.admissible-field-space}
\mathcal{F}_\Omega
=
\left\{
\mathbf{u}_\omega=(\E,\H)\in\mathcal{U}_\omega:
\gamma_\tau[\mathbf{u}_\omega]\in\mathcal{D}\big(\mathcal{B}_\Gamma\big),\
(\E,\H)\ \text{satisfy patchwise Gauss constraints}
\right\},
\end{equation}
where \(\mathcal{D}(\mathcal{B}_\Gamma)\) is the domain of boundary data compatible
with Eq.~\eqref{eq:ch1.admissible-data-class} imported from
Section~\ref{sec:ch01.1.2}. Equation~\eqref{eq:ch2.admissible-field-space} refines
the minimal declaration of Subsection~\ref{sec:ch201}: same symbol
\(\mathcal{F}_\Omega\), sharpened topology. Membership implies
Eq.~\eqref{eq:ch1.energy-class} but adds curl and trace control needed for
integration by parts in Eq.~\eqref{eq:ch2.maxwell-adjoint-pairing}
\cite{stratton1941,harrington2001}.

The closed Maxwell operator domain becomes
\begin{equation}
\label{eq:ch2.domain-refinement-L}
\mathcal{D}(\mathcal{L}_\omega)
=
\left\{
\mathbf{u}_\omega\in\mathcal{F}_\Omega:
\mathcal{M}_\omega[\mathbf{u}_\omega]\in L^{2}(\Omega)^{4}
\right\},
\end{equation}
which specializes Eq.~\eqref{eq:ch2.operator-domain-declaration} to Lipschitz
domains with patchwise smooth coefficients \cite{stratton1941}. On
\(\mathcal{D}(\mathcal{L}_\omega)\), the volume residuals of
Eq.~\eqref{eq:ch2.maxwell-driven-operator} are square integrable, and the boundary
row of Eq.~\eqref{eq:ch2.maxwell-boundary-composite} is evaluated in
\(\mathcal{T}_\Gamma\). Exterior domains append outgoing subspaces rather than
changing Eqs.~\eqref{eq:ch2.hcurl-space}--\eqref{eq:ch2.admissible-field-space};
that selection is function-space compatible with weighted norms that decay at
infinity and is detailed in Section~\ref{sec:ch243} \cite{harrington2001}.

Free-charge density paired with \(\Jimp\) in
Eq.~\eqref{eq:ch1.data-specification-tuple} is taken in the scalar class
\(\rho\in L^{2}(\Omega)\) whenever volumetric charge is prescribed, with
distributional interpretation \(\divg\Jimp+\jj\omega\rho=0\) understood in
\(\mathcal{D}'(\Omega)\) as in Eq.~\eqref{eq:ch1.source-admissibility}. Scalar
\(L^{2}\) control is sufficient because Gauss laws in
Eq.~\eqref{eq:ch1.maxwell-operator-block} involve \(\divg\) on \(\varepsilon_0\E\)
and \(\divg\) on \(\mu_0\H\), not gradients of \(\rho\). When \(\rho\) is absent and
\(\Jimp\in\mathcal{J}_\Omega^{(0)}\), the electric Gauss row reduces to a
divergence constraint on \(\E\) enforced weakly through test functions in
\(\HCurl(\Omega)\) \cite{stratton1941}.

Approximation properties matter for numerical projection. Smooth compactly supported
fields are dense in \(\HCurl(\Omega)\) and in \(\mathcal{U}_\omega\) equipped with
norm~\eqref{eq:ch2.field-energy-norm}; every admissible state is a limit of
\(C_c^\infty\) test pairs modified to satisfy boundary class on \(\Gamma\)
\cite{stratton1941}. Density underpins Galerkin convergence: edge-element spaces
approximate \(\mathcal{F}_\Omega\) when their degrees of freedom match tangential
trace continuity on \(\Gamma_{\mathrm{PEC}}\), where
\begin{equation}
\label{eq:ch2.pec-subspace}
\mathcal{U}_\omega^{\mathrm{PEC}}
=
\left\{
(\E,\H)\in\mathcal{U}_\omega:
\hat{\mathbf{n}}\times\E|_{\Gamma_{\mathrm{PEC}}}=\mathbf{0}
\right\},
\end{equation}
encodes perfect-conductor tangential data from
Eq.~\eqref{eq:ch1.outer-boundary-classes} without assuming \(\E=\mathbf{0}\) in
\(\Omega\) \cite{harrington2001}. Integration by parts on
\(\mathcal{U}_\omega^{\mathrm{PEC}}\) yields the boundary-term cancellation used in
Eq.~\eqref{eq:ch2.maxwell-adjoint-pairing} when \(\mathcal{B}_\Gamma\) selects
homogeneous tangential electric trace.

Compact embeddings used later for Fredholm arguments are recorded as
\begin{equation}
\label{eq:ch2.embedding-chain}
\HCurl(\Omega)\hookrightarrow L^{2}(\Omega)^{3}
\hookleftarrow H^{-1/2}(\Gamma)^{3},
\end{equation}
with the first inclusion compact on bounded \(\Omega\) and the trace map continuous
but not compact \cite{stratton1941}. Equation~\eqref{eq:ch2.embedding-chain}
explains why interior cavity spectra are discrete while boundary observation maps are
rank deficient: compact embeddings produce accumulating eigenvalues in volume,
whereas non-compact traces map infinite-dimensional field spaces to finite
observation vectors \cite{harrington2001}. Subsection~\ref{sec:ch213} develops the
operator-theoretic consequences.

\paragraph{Worked illustration (patchwise spaces).}
Consider \(\Omega=\Omega_1\cup\Omega_2\) with interface \(\Sigma_{12}\), vacuum
constants in each patch, and impressed \(\Jimp\in\mathcal{J}_\Omega\) supported in
\(\Omega_1\) only. Fields \(\E|_{\Omega_a},\H|_{\Omega_a}\in\HCurl(\Omega_a)\) with
traces glued on \(\Sigma_{12}\) through Eq.~\eqref{eq:ch1.interface-conditions}
(with \(\mathbf{J}_{s,12}=\mathbf{0}\)). A finite-element space that uses
\(H^{1}(\Omega)^{3}\) elements instead of \(\HCurl\) conforming elements violates
Eq.~\eqref{eq:ch2.hcurl-space}: tangential continuity may hold while
\(\curl\E\) fails to be square integrable across elements, breaking
Eq.~\eqref{eq:ch2.maxwell-weak-form}. Upgrading to edge elements restores
membership in \(\mathcal{U}_\omega\) and aligns the discrete problem with
\(\mathcal{D}(\mathcal{L}_\omega)\) \cite{harrington2001,stratton1941}.

\paragraph{Worked illustration (trace data mismatch).}
Suppose \(\mathbf{g}_\Gamma\) prescribes tangential \(\E\) on \(\Gamma\) with
\(H^{1/2}\) regularity but is not the trace of any \(\HCurl\) field in \(\Omega\).
Then no \(\mathbf{u}_\omega\in\mathcal{F}_\Omega\) satisfies
\(\gamma_{\tau,E}[\E]=\mathbf{g}_\Gamma\), regardless of how accurately a volume
solver minimizes Maxwell residuals. The failure is a trace-space mismatch, not a
Green-function error; correcting it requires either lifting \(\mathbf{g}_\Gamma\) into
\(\mathcal{T}_\Gamma\) as an image under \(\gamma_\tau\) or relaxing the boundary
class through impedance closure in Subsection~\ref{sec:ch214}
\cite{stratton1941,harrington2001}.

\paragraph{Closing summary.}
Subsection~\ref{sec:ch211} establishes the Chapter~\ref{ch:02} HOME for Sobolev
field spaces Eqs.~\eqref{eq:ch2.hcurl-space}--\eqref{eq:ch2.maxwell-product-space},
refined source space Eq.~\eqref{eq:ch2.source-space-refined}, trace and interface
classes Eqs.~\eqref{eq:ch2.trace-space-definition}--\eqref{eq:ch2.interface-trace-space},
and the admissible field domain Eq.~\eqref{eq:ch2.admissible-field-space} that
specializes \(\mathcal{D}(\mathcal{L}_\omega)\) in
Eq.~\eqref{eq:ch2.domain-refinement-L}. Chapter~\ref{ch:01} admissibility is
embedded, not re-proved. Subsection~\ref{sec:ch212} treats abstract operator domains,
ranges, and adjoints on these spaces.

\subsection{Linear operators: domain, range, and adjoint}
\label{sec:ch212}
% TOC tag: [HOME]

Subsection~\ref{sec:ch211} fixed the Hilbert and trace spaces on which Maxwell maps
act. Subsection~\ref{sec:ch212} treats those maps as linear operators with explicit
domains, ranges, and adjoints \cite{stratton1941,harrington2001}. The volume
operator \(\mathcal{M}_\omega\), boundary operator \(\mathcal{B}_\Gamma\), composite
\(\mathcal{L}_\omega\), and solution map \(\mathcal{S}_\omega\) introduced in
Subsections~\ref{sec:ch201}--\ref{sec:ch202} are not re-defined; their
function-analytic structure is recorded here as the Chapter~\ref{ch:02} HOME for
domain--range--adjoint formalism. Weak and energy forms of
Eqs.~\eqref{eq:ch2.maxwell-weak-form} and
\eqref{eq:ch2.maxwell-adjoint-pairing} supply the duality pairings below and are
cited, not re-derived \cite{stratton1941}.

Let \(\mathcal{X}\) and \(\mathcal{Y}\) be Hilbert spaces. A linear operator
\(T:\mathcal{D}(T)\to\mathcal{Y}\) is specified by its domain
\(\mathcal{D}(T)\subseteq\mathcal{X}\) and range
\(\mathcal{R}(T)=T\mathcal{D}(T)\subseteq\mathcal{Y}\). For Maxwell operators,
\(\mathcal{X}=\mathcal{U}_\omega\) from Eq.~\eqref{eq:ch2.maxwell-product-space}
and \(\mathcal{Y}\) is a product of volume and boundary data spaces. The graph
norm
\begin{equation}
\label{eq:ch2.graph-norm}
\|\mathbf{u}\|_{\mathrm{gr}(T)}^{2}
=
\|\mathbf{u}\|_{\mathcal{X}}^{2}+\|T\mathbf{u}\|_{\mathcal{Y}}^{2},
\qquad
\mathbf{u}\in\mathcal{D}(T),
\end{equation}
turns \(\mathcal{D}(T)\) into a Hilbert space when \(T\) is closed
\cite{stratton1941}. Equation~\eqref{eq:ch2.graph-norm} is the abstract tool used in
Subsection~\ref{sec:ch214} to construct self-adjoint closures; here it organizes
domain control for \(\mathcal{L}_\omega\).

\begin{figure}[t]
\centering
\ChOneFigBlock{%
\begin{tikzpicture}[ch1fig, node distance=8mm and 12mm]
  \node[ch1box] (d) {$\mathcal{D}(T)$};
  \node[ch1box, right=16mm of d] (r) {$\mathcal{R}(T)$};
  \node[ch1box, below=11mm of d] (x) {$\mathcal{X}$};
  \node[ch1box, below=11mm of r] (y) {$\mathcal{Y}$};
  \draw[ch1flow] (d) -- node[above,font=\scriptsize] {$T$} (r);
  \draw[ch1flow] (d) -- (x);
  \draw[ch1flow] (r) -- (y);
\end{tikzpicture}%
}
\caption{Operator domain and range: \(\mathcal{D}(T)\) embeds in state space
\(\mathcal{X}\); \(T\) maps into data space \(\mathcal{Y}\).}
\label{fig:ch2.operator-domain-range}
\end{figure}

Figure~\ref{fig:ch2.operator-domain-range} applies to each Maxwell map once
\(\mathcal{X}\) and \(\mathcal{Y}\) are chosen. Confusing \(\mathcal{D}(T)\) with
all of \(\mathcal{X}\) is a common source of spurious boundary terms in flux
identities \cite{stratton1941}.

The volume Maxwell operator of Eq.~\eqref{eq:ch2.maxwell-driven-operator} acts as
\begin{equation}
\label{eq:ch2.volume-operator-domain-range}
\mathcal{M}_\omega:\mathcal{D}(\mathcal{M}_\omega)\longrightarrow
\mathcal{Y}_{\mathrm{vol}},
\qquad
\mathcal{Y}_{\mathrm{vol}}=L^{2}(\Omega)^{3}\times L^{2}(\Omega)^{3}
\times L^{2}(\Omega)\times L^{2}(\Omega)^{3},
\end{equation}
with \(\mathcal{D}(\mathcal{M}_\omega)\subset\mathcal{F}_\Omega\) the subspace on
which all four residual rows of Eq.~\eqref{eq:ch1.maxwell-operator-block} lie in
\(\mathcal{Y}_{\mathrm{vol}}\) \cite{stratton1941}. Operand interpretation:
\(\mathcal{M}_\omega\) evaluates curl-coupled Maxwell residuals pointwise in the
\(L^{2}\) sense; physically, it returns whether a candidate \((\E,\H)\) satisfies
local Maxwell balance with impressed data removed to the forcing vector
Eq.~\eqref{eq:ch2.maxwell-forcing-vector}. The map is linear in
\(\mathbf{u}_\omega\) and bounded on bounded subsets of
\(\mathcal{D}(\mathcal{M}_\omega)\) but is not bounded on all of
\(\mathcal{U}_\omega\) because curl norms are not controlled by \(L^{2}\) norms alone
\cite{harrington2001}.

Tangential traces form bounded boundary operators on \(\HCurl(\Omega)\):
\begin{equation}
\label{eq:ch2.trace-operator-bounded}
\gamma_\tau:\HCurl(\Omega)\longrightarrow H^{-1/2}(\Gamma)^{3},
\qquad
\gamma_\tau[\mathbf{u}]=\hat{\mathbf{n}}\times\mathbf{u}|_{\Gamma},
\end{equation}
with \(\|\gamma_\tau\|_{\mathrm{op}}<\infty\) on Lipschitz \(\Gamma\)
\cite{stratton1941}. Equation~\eqref{eq:ch2.trace-operator-bounded} specializes
Eqs.~\eqref{eq:ch2.tangential-trace-maps}--\eqref{eq:ch2.tangential-trace-magnetic}
to a single template. The boundary operator of
Eq.~\eqref{eq:ch1.boundary-data-operator} is
\begin{equation}
\label{eq:ch2.boundary-operator-domain-range}
\mathcal{B}_\Gamma:\mathcal{D}(\mathcal{B}_\Gamma)\longrightarrow\mathcal{Y}_\Gamma,
\qquad
\mathcal{D}(\mathcal{B}_\Gamma)\subseteq\mathcal{T}_\Gamma,
\end{equation}
where \(\mathcal{T}_\Gamma\) and \(\mathcal{Y}_\Gamma\) are from
Subsection~\ref{sec:ch211}. Impedance and outgoing classes from
Eqs.~\eqref{eq:ch1.outer-boundary-classes}--\eqref{eq:ch1.radiation-trace-condition}
restrict \(\mathcal{D}(\mathcal{B}_\Gamma)\) to admissible trace pairs compatible
with Eq.~\eqref{eq:ch1.admissible-data-class} \cite{stratton1941,harrington2001}.

The composite operator of Eq.~\eqref{eq:ch2.maxwell-boundary-composite} is
\begin{equation}
\label{eq:ch2.composite-operator-domain}
\mathcal{L}_\omega:\mathcal{D}(\mathcal{L}_\omega)\longrightarrow
\mathcal{Y}_{\mathrm{vol}}\times\mathcal{Y}_\Gamma,
\end{equation}
with \(\mathcal{D}(\mathcal{L}_\omega)\) given by
Eq.~\eqref{eq:ch2.domain-refinement-L}. Range membership is \emph{coupled}: a state
cannot satisfy volume rows while violating boundary rows in the product sense of
Eq.~\eqref{eq:ch2.closed-maxwell-bvp}. The graph norm
\eqref{eq:ch2.graph-norm} on \(\mathcal{L}_\omega\) controls both residuals
simultaneously and is the natural norm for convergence tests in numerical Maxwell
solvers \cite{harrington2001}.

The harmonic solution map is an operator on source space:
\begin{equation}
\label{eq:ch2.solution-operator-domain}
\mathcal{S}_\omega:\mathcal{D}(\mathcal{S}_\omega)\longrightarrow\mathcal{F}_\Omega,
\qquad
\mathcal{D}(\mathcal{S}_\omega)=\mathcal{J}_\Omega,
\end{equation}
defined by Eq.~\eqref{eq:ch2.solution-map-resolvent} as
\(\mathcal{L}_\omega[\mathcal{S}_\omega[\Jimp]]=[\mathbf{s}_\omega(\Jimp),\mathbf{g}_\Gamma]\).
Well-posedness requires invertibility of \(\mathcal{L}_\omega\) on the physical
branch selected by \(\mathcal{B}_\Gamma\); then
\begin{equation}
\label{eq:ch2.resolvent-operator}
\mathcal{S}_\omega=\mathcal{L}_\omega^{-1}\circ
\big(\mathbf{s}_\omega(\cdot),\mathbf{g}_\Gamma\big),
\end{equation}
with \(\mathcal{L}_\omega^{-1}\) understood as the outgoing or cavity-truncated
resolvent of Sections~\ref{sec:ch243} and~\ref{sec:ch263}, not as an algebraic
inverse on all of \(\mathcal{U}_\omega\) \cite{stratton1941}. Radiation and
observation operators compose accordingly:
\begin{equation}
\label{eq:ch2.radiation-operator-composition}
\mathcal{R}_\omega=\Pi_{\mathrm{rad}}\circ\mathcal{S}_\omega,
\qquad
\Lambda_{\mathrm{rad}}=\mathcal{O}_{\mathrm{far}}\circ\mathcal{R}_\omega,
\end{equation}
with \(\Pi_{\mathrm{rad}}\) from Eq.~\eqref{eq:ch2.radiation-operator-def} and
\(\mathcal{O}_{\mathrm{far}}\) from Subsection~\ref{sec:ch201}
\cite{harrington2001}.

\begin{figure}[t]
\centering
\ChOneFigBlock{%
\begin{tikzpicture}[ch1fig, node distance=7mm and 11mm]
  \node[ch1box] (j) {$\mathcal{J}_\Omega$};
  \node[ch1box, right=13mm of j] (s) {$\mathcal{S}_\omega$};
  \node[ch1box, right=13mm of s] (f) {$\mathcal{F}_\Omega$};
  \node[ch1box, right=13mm of f] (r) {$\mathcal{R}_\omega$};
  \draw[ch1flow] (j) -- (s) -- (f) -- node[above,font=\scriptsize] {$\Pi_{\mathrm{rad}}$} (r);
\end{tikzpicture}%
}
\caption{Operator chain on declared domains:
\(\mathcal{S}_\omega:\mathcal{J}_\Omega\to\mathcal{F}_\Omega\),
\(\mathcal{R}_\omega:\mathcal{J}_\Omega\to\mathcal{F}_{\mathrm{rad}}\).}
\label{fig:ch2.solution-radiation-operator-chain}
\end{figure}

Figure~\ref{fig:ch2.solution-radiation-operator-chain} restates
Eq.~\eqref{eq:ch2.pipeline-composition} in domain--range language. Each arrow is
meaningful only on the declared domain; extending \(\mathcal{S}_\omega\) to
non-admissible sources produces elements of \(\mathcal{U}_\omega\) that need not
lie in \(\mathcal{F}_\Omega\) \cite{stratton1941}.

Adjoints are defined through duality pairings compatible with
Eq.~\eqref{eq:ch2.dual-source-pairing}. For the volume operator, the formal adjoint
\( \mathcal{M}_\omega^{\dagger}\) satisfies
\begin{equation}
\label{eq:ch2.volume-adjoint-definition}
\langle\mathcal{M}_\omega[\mathbf{u}_\omega],\mathbf{v}_\omega\rangle_{\mathrm{vol}}
=
\langle\mathbf{u}_\omega,\mathcal{M}_\omega^{\dagger}[\mathbf{v}_\omega]\rangle_{\mathcal{U}}
+\mathcal{B}_{\mathrm{surf}}(\mathbf{u}_\omega,\mathbf{v}_\omega),
\end{equation}
where \(\langle\cdot,\cdot\rangle_{\mathrm{vol}}\) is the componentwise
\(L^{2}\) pairing on \(\mathcal{Y}_{\mathrm{vol}}\),
\(\langle\cdot,\cdot\rangle_{\mathcal{U}}\) is induced by
Eq.~\eqref{eq:ch2.field-energy-norm}, and
\(\mathcal{B}_{\mathrm{surf}}\) collects boundary terms from integration by parts
\cite{stratton1941}. On \(\mathcal{D}(\mathcal{M}_\omega)\cap\mathcal{U}_\omega^{\mathrm{PEC}}\)
from Eq.~\eqref{eq:ch2.pec-subspace}, Eq.~\eqref{eq:ch2.maxwell-adjoint-pairing}
identifies \(\mathcal{B}_{\mathrm{surf}}\) with the tangential trace pairing of
Eq.~\eqref{eq:ch2.maxwell-adjoint-pairing}; that identity is imported, not
re-proved \cite{stratton1941,harrington2001}.

The composite adjoint satisfies the block structure
\begin{equation}
\label{eq:ch2.composite-adjoint-block}
\mathcal{L}_\omega^{\dagger}
=
\begin{bmatrix}
\mathcal{M}_\omega^{\dagger} & 0\\
0 & \mathcal{B}_\Gamma^{\dagger}
\end{bmatrix}
\quad\text{on compatible pairs,}
\end{equation}
with cross-coupling entering only through shared trace data when boundary and
volume domains are not synchronized. Self-adjoint \emph{realizations} of truncated
Maxwell problems require choosing \(\mathcal{B}_\Gamma\) so that
\(\mathcal{B}_{\mathrm{surf}}\) cancels on \(\mathcal{D}(\mathcal{L}_\omega)\);
Subsection~\ref{sec:ch214} owns that cancellation \cite{stratton1941}.

Forcing functionals induced by impressed data are bounded linear operators on test
spaces. The map
\begin{equation}
\label{eq:ch2.forcing-functional-operator}
\mathfrak{f}_\omega:\mathcal{U}_\omega^{0}\longrightarrow\mathbb{C},
\qquad
\mathfrak{f}_\omega[\mathbf{v}_\omega]=
\langle\mathbf{s}_\omega(\mathfrak{D}),\mathbf{v}_\omega\rangle_{\mathrm{vol}},
\end{equation}
is continuous when \(\mathbf{s}_\omega\in\mathcal{Y}_{\mathrm{vol}}\) and
\(\mathbf{v}_\omega\in\mathcal{U}_\omega^{0}\) carries homogeneous boundary traces
compatible with \(\mathcal{B}_\Gamma\) \cite{stratton1941}. Equation
\eqref{eq:ch2.forcing-functional-operator} is the operator reading of the right-hand
side of Eq.~\eqref{eq:ch2.maxwell-weak-form}. Source-to-field maps inherit
continuity on bounded subsets: if \(\|\Jimp\|_{\mathcal{J}}\le C\) and
\(\mathcal{L}_\omega\) is boundedly invertible on a fixed branch, then
\begin{equation}
\label{eq:ch2.solution-map-stability}
\|\mathcal{S}_\omega[\Jimp]\|_{\mathcal{U}}
\le
C_{\mathrm{res}}\,\|\mathbf{s}_\omega(\Jimp)\|_{\mathrm{vol}},
\end{equation}
with \(C_{\mathrm{res}}=\|\mathcal{L}_\omega^{-1}\|\) on the selected subspace
\cite{harrington2001,stratton1941}. Inequality~\eqref{eq:ch2.solution-map-stability}
is the operator form of solver stability under source perturbations.

Kernel and range spaces are recorded for later null-space analysis:
\begin{equation}
\label{eq:ch2.kernel-range-L}
\ker\mathcal{L}_\omega
=
\left\{
\mathbf{u}_\omega\in\mathcal{D}(\mathcal{L}_\omega):
\mathcal{L}_\omega[\mathbf{u}_\omega]=\mathbf{0}
\right\},
\end{equation}
\begin{equation}
\label{eq:ch2.range-L}
\mathcal{R}(\mathcal{L}_\omega)
=
\left\{
[\mathbf{s},\mathbf{g}]^{\mathsf{T}}\in
\mathcal{Y}_{\mathrm{vol}}\times\mathcal{Y}_\Gamma:
\exists\,\mathbf{u}_\omega\in\mathcal{D}(\mathcal{L}_\omega):\
\mathcal{L}_\omega[\mathbf{u}_\omega]=[\mathbf{s},\mathbf{g}]^{\mathsf{T}}
\right\}.
\end{equation}
On bounded domains with passive \(\mathcal{B}_\Gamma\), \(\ker\mathcal{L}_\omega\)
contains cavity modes; on exterior domains with outgoing closure,
\(\ker\mathcal{L}_\omega\) is trivial on the outgoing branch
\cite{stratton1941,harrington2001}. Section~\ref{sec:ch272} analyzes
\(\ker\mathcal{R}_\omega\); Subsection~\ref{sec:ch213} treats compactness and
Fredholm index on \(\mathcal{R}(\mathcal{L}_\omega)\).

Extension operators lift boundary data into volume fields. A right inverse
\(E_\Gamma:\mathcal{Y}_\Gamma\to\HCurl(\Omega)\) satisfying
\(\gamma_\tau\circ E_\Gamma=\mathrm{id}\) on compatible subclasses exists only for
\(\mathbf{g}_\Gamma\) in the range of \(\gamma_\tau\); otherwise no
\(\mathbf{u}_\omega\in\mathcal{F}_\Omega\) meets the boundary row
\cite{stratton1941}. Extension operators enter Calder\'on preconditioners and
boundary-element formulations developed in later sections; here they clarify that
\(\mathcal{B}_\Gamma\) acts on traces, not on arbitrary \(L^{2}\) boundary vectors
\cite{harrington2001}.

\paragraph{Worked illustration (domain mismatch).}
Let \(\mathbf{u}_\omega\in\mathcal{U}_\omega\setminus\mathcal{D}(\mathcal{L}_\omega)\)
have \(\curl\E\in L^{2}\) but \(\mathcal{M}_\omega[\mathbf{u}_\omega]\notin
\mathcal{Y}_{\mathrm{vol}}\) because Gauss constraints fail in the \(L^{2}\) sense.
Evaluating \(\mathfrak{a}_\omega(\mathbf{u}_\omega,\mathbf{v}_\omega)\) by parts
then produces boundary terms not accounted for in
Eq.~\eqref{eq:ch2.maxwell-weak-form}, and a numerical residual can appear small in
\(L^{2}\) while the operator adjoint
Eq.~\eqref{eq:ch2.volume-adjoint-definition} is undefined. Restricting to
\(\mathcal{D}(\mathcal{L}_\omega)\) restores synchronized volume and boundary rows
\cite{stratton1941,harrington2001}.

\paragraph{Worked illustration (resolvent branch).}
Fix \(\Jimp\in\mathcal{J}_\Omega\) and two distinct outgoing closures
\(\mathcal{B}_\Gamma^{(a)}\) and \(\mathcal{B}_\Gamma^{(b)}\) on the same exterior
domain. The operators \(\mathcal{L}_\omega^{(a)}\) and
\(\mathcal{L}_\omega^{(b)}\) have the same volume domain but different boundary
domains; Eq.~\eqref{eq:ch2.resolvent-operator} yields two resolvents
\(\mathcal{S}_\omega^{(a)}\neq\mathcal{S}_\omega^{(b)}\) in general. Comparing
simulation to measurement without matching \(\mathcal{D}(\mathcal{B}_\Gamma)\) is
therefore a comparison of different operators, not of different discretizations of
one operator \cite{stratton1941}.

\paragraph{Closing summary.}
Subsection~\ref{sec:ch212} establishes the Chapter~\ref{ch:02} HOME for operator
domains and ranges Eqs.~\eqref{eq:ch2.volume-operator-domain-range}--\eqref{eq:ch2.composite-operator-domain},
the solution and radiation chain
Eqs.~\eqref{eq:ch2.solution-operator-domain}--\eqref{eq:ch2.radiation-operator-composition},
adjoint templates Eqs.~\eqref{eq:ch2.volume-adjoint-definition}--\eqref{eq:ch2.composite-adjoint-block},
and kernel/range notation Eqs.~\eqref{eq:ch2.kernel-range-L}--\eqref{eq:ch2.range-L}.
Spaces from Subsection~\ref{sec:ch211} are presupposed; boundedness, compactness, and
spectral compression follow in Subsection~\ref{sec:ch213}.

\subsection{Boundedness, compactness, rank deficiency, spectral compression}
\label{sec:ch213}
% TOC tag: [HOME]

Subsections~\ref{sec:ch211}--\ref{sec:ch212} fixed spaces and operator domains.
Subsection~\ref{sec:ch213} records boundedness and compactness properties, rank
deficiency of observation maps, and the spectral-compression viewpoint that
Section~\ref{sec:ch26} develops for radiation operators
\cite{stratton1941,harrington2001}. Fredholm index computation and complete
spectral decomposition remain in Sections~\ref{sec:ch263} and~\ref{sec:ch261};
here the Chapter~\ref{ch:02} HOME establishes the operator-theoretic lemmas those
sections invoke by reference.

A linear operator \(T:\mathcal{X}\to\mathcal{Y}\) between Hilbert spaces is
\emph{bounded} if
\begin{equation}
\label{eq:ch2.operator-boundedness}
\|T\|_{\mathrm{op}}
=
\sup_{\|\mathbf{x}\|_{\mathcal{X}}=1}\|T\mathbf{x}\|_{\mathcal{Y}}
<\infty.
\end{equation}
Boundedness is equivalent to continuity at the origin and is the default
hypothesis for composition with measurement maps \cite{stratton1941}. On the
Maxwell chain, operators split into bounded and unbounded classes.

Trace and observation maps are bounded on declared spaces. The tangential trace
Eq.~\eqref{eq:ch2.trace-operator-bounded} satisfies
\(\|\gamma_\tau\|_{\mathrm{op}}<\infty\) from \(\HCurl(\Omega)\) to
\(H^{-1/2}(\Gamma)^{3}\) \cite{stratton1941}. Far-zone readouts from
Subsection~\ref{sec:ch201} define bounded functionals when restricted to
\(\mathcal{F}_{\mathrm{rad}}\):
\begin{equation}
\label{eq:ch2.far-observation-bounded}
\|\mathcal{O}_{\mathrm{far}}\|_{\mathrm{op}}
=
\sup_{\|\mathbf{u}\|_{\mathcal{U}}=1}
\big|\mathcal{O}_{\mathrm{far}}[\mathbf{u}]\big|
<\infty
\quad\text{on}\ \mathcal{F}_{\mathrm{rad}}.
\end{equation}
Laboratory maps \(\mathcal{O}_{\mathrm{meas}}\circ\mathcal{O}_{\mathrm{far}}\) in
Eq.~\eqref{eq:ch2.operational-measurement-chain} are bounded on the same subspace
because composition of bounded operators is bounded \cite{harrington2001}.

Volume and composite Maxwell operators are unbounded on full
\(\mathcal{U}_\omega\). The map \(\mathcal{M}_\omega\) of
Eq.~\eqref{eq:ch2.volume-operator-domain-range} has
\(\|\mathcal{M}_\omega\|_{\mathrm{op}}=\infty\) when
\(\mathcal{D}(\mathcal{M}_\omega)=\mathcal{U}_\omega\) because controlling
\(\|\mathcal{M}_\omega[\mathbf{u}]\|_{\mathrm{vol}}\) would require uniform curl
bounds beyond \(L^{2}\) control \cite{stratton1941}. Restriction to
\(\mathcal{D}(\mathcal{L}_\omega)\) with graph norm
Eq.~\eqref{eq:ch2.graph-norm} restores boundedness of the \emph{resolvent}
\(\mathcal{L}_\omega^{-1}\) on the physical branch where invertibility holds:
\begin{equation}
\label{eq:ch2.resolvent-bounded}
\|\mathcal{L}_\omega^{-1}\|_{\mathrm{op}}
<\infty
\quad\Longleftrightarrow\quad
\mathcal{L}_\omega\ \text{bijective on}\ \mathcal{D}(\mathcal{L}_\omega)/\ker\mathcal{L}_\omega.
\end{equation}
Equation~\eqref{eq:ch2.resolvent-bounded} specializes
Eq.~\eqref{eq:ch2.solution-map-stability}; exterior outgoing problems use
resolvent bounds on weighted subspaces rather than on all of
\(\mathcal{U}_\omega\) \cite{harrington2001,stratton1941}.

\begin{figure}[t]
\centering
\ChOneFigBlock{%
\begin{tikzpicture}[ch1fig, node distance=8mm and 10mm]
  \node[ch1box] (hc) {$\HCurl(\Omega)$};
  \node[ch1box, right=16mm of hc] (l2) {$L^{2}(\Omega)^{3}$};
  \node[ch1box, below=12mm of l2] (tr) {$H^{-1/2}(\Gamma)^{3}$};
  \draw[ch1flow] (hc) -- node[above,font=\scriptsize] {compact} (l2);
  \draw[ch1flow] (hc) -- node[right,font=\scriptsize] {$\gamma_\tau$} (tr);
\end{tikzpicture}%
}
\caption{Compact embedding \(\HCurl(\Omega)\hookrightarrow L^{2}(\Omega)^{3}\) on
bounded \(\Omega\); trace map is bounded but not compact
(Eq.~\eqref{eq:ch2.embedding-chain}).}
\label{fig:ch2.compact-embedding-schematic}
\end{figure}

Figure~\ref{fig:ch2.compact-embedding-schematic} restates
Eq.~\eqref{eq:ch2.embedding-chain}. Rellich-type compactness for
\(\HCurl(\Omega)\hookrightarrow L^{2}(\Omega)^{3}\) on bounded domains implies
that bounded sequences in \(\HCurl\) have \(L^{2}\)-convergent subsequences
\cite{stratton1941}. An operator \(K:\mathcal{X}\to\mathcal{Y}\) is
\emph{compact} if bounded sets in \(\mathcal{X}\) map to relatively compact sets
in \(\mathcal{Y}\):
\begin{equation}
\label{eq:ch2.compact-operator-def}
K\ \text{compact}
\iff
K(B_{\mathcal{X}}(0,1))\ \text{has compact closure in}\ \mathcal{Y}.
\end{equation}
Compact perturbations of identity produce Fredholm structure used in
Section~\ref{sec:ch263}; Subsection~\ref{sec:ch213} records the embedding lemmas
without computing index \cite{harrington2001}.

Radiation and Green maps on fixed bounded domains often decompose as
\begin{equation}
\label{eq:ch2.green-compact-split}
\mathcal{S}_\omega
=
G_{\mathrm{sing}}+K_{\mathrm{cpt}},
\end{equation}
with \(G_{\mathrm{sing}}\) capturing singular near-field kernels and
\(K_{\mathrm{cpt}}\) compact on \(\mathcal{J}_\Omega\) when sources have compact
support and \(\Omega\) is bounded \cite{stratton1941}. Equation
\eqref{eq:ch2.green-compact-split} is the operator template for
Section~\ref{sec:ch23}; kernel construction is not repeated here. The split
separates non-compact singular behavior at coincident points from compact
smoothing at finite separation \cite{harrington2001}.

Rank deficiency enters through finite observation. An aperture-limited measurement
map \(\mathcal{O}_{\mathrm{meas}}:\mathcal{F}_{\mathrm{rad}}\to\mathbb{C}^{N}\)
has finite-dimensional range with \(N<\infty\):
\begin{equation}
\label{eq:ch2.finite-rank-observation}
\mathcal{O}_{\mathrm{meas}}
=
\sum_{n=1}^{N}\langle\cdot,\mathbf{f}_n\rangle_{\mathcal{U}}\,\mathbf{e}_n,
\end{equation}
for dual vectors \(\mathbf{f}_n\) and coordinate directions \(\mathbf{e}_n\).
Therefore \(\mathrm{rank}(\mathcal{O}_{\mathrm{meas}})\le N\) and
\begin{equation}
\label{eq:ch2.observation-null-space}
\ker\mathcal{O}_{\mathrm{meas}}
=
\left\{
\mathbf{u}\in\mathcal{F}_{\mathrm{rad}}:
\mathcal{O}_{\mathrm{meas}}[\mathbf{u}]=\mathbf{0}
\right\}
\end{equation}
is infinite-dimensional whenever \(\mathcal{F}_{\mathrm{rad}}\) is
infinite-dimensional \cite{stratton1941,harrington2001}. Equation
\eqref{eq:ch2.observation-null-space} is the operator form of the measurement
screening in Section~\ref{sec:ch01.5}: finite modal sets cannot resolve all
angular content. Rank deficiency is structural, not a discretization artifact alone.

\begin{figure}[t]
\centering
\ChOneFigBlock{%
\begin{tikzpicture}[ch1fig, node distance=7mm and 9mm]
  \node[ch1box] (f) {$\mathcal{F}_{\mathrm{rad}}$};
  \node[ch1box, right=14mm of f] (o) {$\mathcal{O}_{\mathrm{meas}}$};
  \node[ch1box, right=14mm of o] (cn) {$\mathbb{C}^{N}$};
  \node[ch1box, below=11mm of f] (ker) {$\ker\mathcal{O}_{\mathrm{meas}}$};
  \draw[ch1flow] (f) -- (o) -- (cn);
  \draw[ch1flow, dashed] (ker) -- (f);
\end{tikzpicture}%
}
\caption{Finite-rank observation: infinite-dimensional field space maps to
\(\mathbb{C}^{N}\); \(\ker\mathcal{O}_{\mathrm{meas}}\) is infinite-dimensional for
\(N<\infty\).}
\label{fig:ch2.rank-deficient-observation}
\end{figure}

Figure~\ref{fig:ch2.rank-deficient-observation} applies to
\(\Lambda_{\mathrm{rad}}=\mathcal{O}_{\mathrm{meas}}\circ\mathcal{O}_{\mathrm{far}}
\circ\mathcal{R}_\omega\) from Eq.~\eqref{eq:ch2.radiation-operator-composition}.
Angular-spectrum reports that omit unresolved harmonics lie in the range of
\(\mathcal{O}_{\mathrm{meas}}\) but not in its injective sector
\cite{stratton1941}. Subsection~\ref{sec:ch213} does not classify report types;
Section~\ref{sec:ch01.5} remains authoritative.

Singular values quantify rank deficiency. For compact
\(K:\mathcal{X}\to\mathcal{Y}\), the singular value sequence
\(\{\sigma_n(K)\}_{n\ge 1}\) decreases to zero, and
\begin{equation}
\label{eq:ch2.singular-value-decay}
\|K-R_N K\|_{\mathrm{op}}
=
\inf_{\mathrm{rank}\,R_N\le N}\|K-R_N\|_{\mathrm{op}}
\sim \sigma_{N+1}(K),
\end{equation}
where \(R_N\) ranges over rank-\(N\) operators \cite{stratton1941}. For radiation
maps truncated to \(N\) spherical or plane-wave modes, Eq.~\eqref{eq:ch2.singular-value-decay}
controls how much field content lies outside the retained modal subspace; explicit
VSH eigenbases are deferred to Section~\ref{sec:ch261} and Chapter~\ref{ch:03}
\cite{harrington2001}.

\emph{Spectral compression} names the reduction of infinite-dimensional radiation
dynamics to a finite or countable spectral model on the physical branch. Let
\(\{\phi_n\}\) be a spectral family in \(\mathcal{F}_\Omega\) and
\(\{P_n\}\) the associated one-dimensional projectors. A truncated radiation operator
\begin{equation}
\label{eq:ch2.spectral-compression-operator}
\mathcal{R}_\omega^{(N)}
=
\sum_{n=1}^{N}\lambda_n\,P_n\circ\mathcal{S}_\omega
\end{equation}
retains only modes with eigenvalues \(\lambda_n\) above a declared threshold
\cite{stratton1941,harrington2001}. Equation
\eqref{eq:ch2.spectral-compression-operator} is the abstract template for
Section~\ref{sec:ch26}; concrete eigenfunctions are not constructed here. On
bounded cavities, compact embeddings yield accumulating eigenvalues
\(\lambda_n\to\infty\) for interior resonances; on exterior domains, continuous
spectra require outgoing projectors rather than discrete sums
\cite{stratton1941}.

Compression error is bounded in operator norm when \(K\) is compact:
\begin{equation}
\label{eq:ch2.compression-error-bound}
\|\mathcal{R}_\omega-\mathcal{R}_\omega^{(N)}\|_{\mathrm{op}}
\le
C\,\sigma_{N+1}(K_{\mathrm{cpt}}),
\end{equation}
with \(K_{\mathrm{cpt}}\) from Eq.~\eqref{eq:ch2.green-compact-split} and
\(C\) independent of impressed data on bounded \(\mathcal{J}_\Omega\)
\cite{harrington2001}. Physically, Eq.~\eqref{eq:ch2.compression-error-bound}
states that truncating modal content beyond the \((N+1)\)st singular value cannot
be recovered by any rank-\(N\) surrogate operator.

Rank deficiency of \(\mathcal{R}_\omega\) itself is distinct from observation rank
deficiency. Non-radiating sources satisfy
\(\mathcal{R}_\omega[\Jimp]=\mathbf{0}\) while
\(\mathcal{S}_\omega[\Jimp]\neq\mathbf{0}\) (Eq.~\eqref{eq:ch2.null-radiation-preview});
Section~\ref{sec:ch272} owns that null space. Here the point is compositional:
\begin{equation}
\label{eq:ch2.rank-composition-law}
\mathrm{rank}(\mathcal{O}_{\mathrm{meas}}\mathcal{R}_\omega)
\le
\min\big\{\mathrm{rank}(\mathcal{O}_{\mathrm{meas}}),
\dim\mathcal{R}(\mathcal{R}_\omega)\big\},
\end{equation}
so measurement rank caps reported angular content even when
\(\mathcal{R}_\omega\) is injective on admissible sources
\cite{stratton1941,harrington2001}. Both deficits must appear in uncertainty
budgets for angular-transport claims imported from Section~\ref{sec:ch01.2.3}.

Fredholm structure preview: on bounded \(\Omega\) with passive
\(\mathcal{B}_\Gamma\), \(\mathcal{L}_\omega\) is Fredholm on
\(\mathcal{D}(\mathcal{L}_\omega)\) modulo cavity kernels when written as
\begin{equation}
\label{eq:ch2.fredholm-split}
\mathcal{L}_\omega
=
\mathcal{L}_0+K,
\qquad
K\ \text{compact on graph domain},
\end{equation}
with \(\mathcal{L}_0\) an invertible reference operator on the quotient by
\(\ker\mathcal{L}_\omega\) \cite{stratton1941}. Index and solvability criteria are
developed in Section~\ref{sec:ch263}; Eq.~\eqref{eq:ch2.fredholm-split} supplies the
split Subsection~\ref{sec:ch212} anticipated without reopening the preview in
Subsection~\ref{sec:ch202}.

\paragraph{Worked illustration (modal truncation).}
A spherical near-field scanner records \(N=50\) tangential field samples on a
closed surface. The map \(\mathcal{O}_{\mathrm{meas}}\) has rank at most \(50\), so
angular harmonics beyond the resolved set lie in
\(\ker\mathcal{O}_{\mathrm{meas}}\). Applying Eq.~\eqref{eq:ch2.singular-value-decay}
to the compact smoothing operator from sample interpolation shows that recovered
far-zone coefficients error is controlled by \(\sigma_{51}\), not by Maxwell
residual alone \cite{stratton1941,harrington2001}. Increasing frequency without
increasing \(N\) lowers \(\sigma_{N+1}\) slowly for fixed aperture, worsening
compression error even when \(\mathcal{L}_\omega[\mathbf{u}_\omega]=\mathbf{0}\)
in the volume.

\paragraph{Worked illustration (interior versus exterior spectrum).}
In a closed cavity, compact embedding yields discrete resonances and
Eq.~\eqref{eq:ch2.spectral-compression-operator} sums over cavity modes. In an
exterior problem with outgoing closure, the same formula requires a continuous
spectral integral rather than a finite sum; using cavity compression on an
open antenna model conflates discrete storage modes with continuous radiation
modes \cite{stratton1941}. Section~\ref{sec:ch243} selects the outgoing branch;
Subsection~\ref{sec:ch213} only records that compression templates differ by
domain class.

\paragraph{Closing summary.}
Subsection~\ref{sec:ch213} establishes the Chapter~\ref{ch:02} HOME for
boundedness Eq.~\eqref{eq:ch2.operator-boundedness} and resolvent bounds
Eq.~\eqref{eq:ch2.resolvent-bounded}, compact operators and Green splits
Eqs.~\eqref{eq:ch2.compact-operator-def}--\eqref{eq:ch2.green-compact-split},
finite-rank observation Eqs.~\eqref{eq:ch2.finite-rank-observation}--\eqref{eq:ch2.observation-null-space},
spectral compression Eqs.~\eqref{eq:ch2.spectral-compression-operator}--\eqref{eq:ch2.compression-error-bound},
and rank composition Eq.~\eqref{eq:ch2.rank-composition-law}. Section~\ref{sec:ch263}
develops Fredholm consequences; Subsection~\ref{sec:ch214} treats self-adjoint
closures and admissibility.

\subsection{Self-adjointness, closures, and admissibility}
\label{sec:ch214}
% TOC tag: [HOME][USE SS1.0.3]

Subsections~\ref{sec:ch211}--\ref{sec:ch213} supplied spaces, operator domains,
compactness, and spectral compression templates. Subsection~\ref{sec:ch214}
identifies self-adjoint and coercive realizations of truncated Maxwell operators
on bounded domains, records closure and quotient constructions that remove cavity
kernels, and states the admissibility embedding that links Chapter~\ref{ch:01}
finite-energy criteria to the operator domains declared here
\cite{stratton1941,harrington2001}. Weak and energy forms
Eqs.~\eqref{eq:ch2.maxwell-energy-form}--\eqref{eq:ch2.maxwell-weak-form} and the
adjoint pairing Eq.~\eqref{eq:ch2.maxwell-adjoint-pairing} are presupposed; they
are not re-derived \cite{stratton1941}.

A densely defined operator \(T\) on Hilbert space \(\mathcal{H}\) is
\emph{symmetric} if
\(\langle T\mathbf{u},\mathbf{v}\rangle=\langle\mathbf{u},T\mathbf{v}\rangle\)
for all \(\mathbf{u},\mathbf{v}\in\mathcal{D}(T)\), and \emph{self-adjoint} if
\(\mathcal{D}(T^{\dagger})=\mathcal{D}(T)\) with equality of adjoints
\cite{stratton1941}. Maxwell operators are symmetric on suitable domains because
Eq.~\eqref{eq:ch2.maxwell-adjoint-pairing} cancels boundary pairings when
\(\mathcal{B}_\Gamma\) is passive and homogeneous on test states. Self-adjointness
requires matching domain and range of adjoints; that match fails on exterior
domains without restricting to outgoing subspaces \cite{harrington2001}.

On bounded \(\Omega\) with passive \(\mathcal{B}_\Gamma\), the sesquilinear form
\(\mathfrak{a}_\omega\) of Eq.~\eqref{eq:ch2.maxwell-energy-form} is coercive on
the quotient by cavity modes:
\begin{equation}
\label{eq:ch2.coercivity-quotient}
\mathfrak{a}_\omega(\mathbf{u},\mathbf{u})
\ge
\alpha\,\|\mathbf{u}\|_{\mathcal{U}}^{2},
\qquad
\mathbf{u}\in\mathcal{U}_\omega^{0}/\ker\mathcal{L}_\omega,
\end{equation}
for some \(\alpha>0\) when \(\ker\mathcal{L}_\omega\) is finite-dimensional and
consists of divergence-free resonances compatible with boundary class
\cite{stratton1941,harrington2001}. Equation~\eqref{eq:ch2.coercivity-quotient}
is the Lax--Milgram hypothesis that makes
Eq.~\eqref{eq:ch2.maxwell-weak-form} well posed on the quotient. Coercivity fails
on the unquotiented space because cavity modes satisfy
\(\mathfrak{a}_\omega(\mathbf{u},\mathbf{u})=0\) while
\(\|\mathbf{u}\|_{\mathcal{U}}>0\).

\begin{figure}[t]
\centering
\ChOneFigBlock{%
\begin{tikzpicture}[ch1fig, node distance=8mm and 11mm]
  \node[ch1box] (u) {$\mathcal{U}_\omega^{0}$};
  \node[ch1box, right=15mm of u] (q) {$\mathcal{U}_\omega^{0}/\ker\mathcal{L}_\omega$};
  \node[ch1box, below=12mm of q] (sol) {Unique weak solution};
  \draw[ch1flow] (u) -- node[above,font=\scriptsize] {quotient} (q);
  \draw[ch1flow] (q) -- (sol);
\end{tikzpicture}%
}
\caption{Cavity-mode quotient: coercivity and uniqueness hold on the quotient space
after removing \(\ker\mathcal{L}_\omega\).}
\label{fig:ch2.cavity-quotient-schematic}
\end{figure}

Figure~\ref{fig:ch2.cavity-quotient-schematic} is the functional-analytic form of
resonance removal in closed cavities. Physical excitation at a cavity frequency
requires impedance or source coupling that breaks pure standing-wave balance;
Subsection~\ref{sec:ch214} records the linear-algebraic quotient; Section~\ref{sec:ch263}
develops spectral consequences \cite{stratton1941}.

Impedance and Leontovich closures from
Eq.~\eqref{eq:ch2.boundary-impedance-operator} yield self-adjoint boundary
realizations when surface impedance is passive:
\begin{equation}
\label{eq:ch2.passive-impedance-class}
\Re\{Z_s\}\ge 0\ \text{on}\ \Gamma_Z
\quad\Longrightarrow\quad
\mathcal{B}_\Gamma\ \text{symmetric on}\ \mathcal{T}_\Gamma.
\end{equation}
Passive \(\Re\{Z_s\}\ge 0\) dissipates power through the boundary rather than
injecting net energy; active or negative-real parts require revisiting admissibility
as noted in Section~\ref{sec:ch01.1.2} \cite{stratton1941}. Perfect conductors
select Eq.~\eqref{eq:ch2.pec-subspace}; impedance boundaries enlarge
\(\mathcal{D}(\mathcal{B}_\Gamma)\) while preserving symmetry when
Eq.~\eqref{eq:ch2.passive-impedance-class} holds \cite{harrington2001}.

When \(\mathcal{B}_\Gamma\) is symmetric but the Maxwell operator is only densely
defined, the Friedrichs extension supplies a self-adjoint realization:
\begin{equation}
\label{eq:ch2.friedrichs-extension}
\mathcal{L}_\omega^{\mathrm{F}}
=
\overline{\mathcal{L}_\omega}\big|_{\mathcal{D}(\mathfrak{a}_\omega)},
\end{equation}
the closure of \(\mathcal{L}_\omega\) on the energy domain of
\(\mathfrak{a}_\omega\) \cite{stratton1941}. Equation
\eqref{eq:ch2.friedrichs-extension} is the Chapter~\ref{ch:02} HOME for
self-adjoint Maxwell truncations on bounded domains. Exterior outgoing problems
use a different closure template: restrict to the outgoing subspace selected in
Section~\ref{sec:ch243} rather than applying Eq.~\eqref{eq:ch2.friedrichs-extension}
to all of \(\mathcal{U}_\omega\) \cite{harrington2001,stratton1941}.

Well-posedness on the quotient combines coercivity and bounded forcing:
\begin{equation}
\label{eq:ch2.lax-milgram-wellposed}
\exists!\,\mathbf{u}_\omega\in\mathcal{U}_\omega^{0}/\ker\mathcal{L}_\omega:
\mathfrak{a}_\omega(\mathbf{u}_\omega,\mathbf{v})
=
\mathfrak{f}_\omega(\mathbf{v})
\quad\forall\,\mathbf{v}\in\mathcal{U}_\omega^{0}/\ker\mathcal{L}_\omega,
\end{equation}
with \(\mathfrak{f}_\omega\) from Eq.~\eqref{eq:ch2.forcing-functional-operator}.
Stability follows from Eqs.~\eqref{eq:ch2.coercivity-quotient} and
\eqref{eq:ch2.solution-map-stability} \cite{stratton1941}. Resolvent bound
Eq.~\eqref{eq:ch2.resolvent-bounded} holds on the quotient when
\(\mathcal{L}_\omega^{\mathrm{F}}\) is bijective there.

\begin{figure}[t]
\centering
\ChOneFigBlock{%
\begin{tikzpicture}[ch1fig, node distance=7mm and 10mm]
  \node[ch1box] (a) {$\mathfrak{a}_\omega$};
  \node[ch1box, right=14mm of a] (lm) {Lax--Milgram};
  \node[ch1box, right=14mm of lm] (sw) {$\mathcal{S}_\omega$};
  \node[ch1box, below=11mm of a] (b) {$\mathcal{B}_\Gamma$ passive};
  \draw[ch1flow] (a) -- (lm) -- (sw);
  \draw[ch1flow] (b) -- (a);
\end{tikzpicture}%
}
\caption{Self-adjoint closure pathway on bounded domains: passive
\(\mathcal{B}_\Gamma\) supports coercive \(\mathfrak{a}_\omega\) and stable
\(\mathcal{S}_\omega\) on the cavity quotient.}
\label{fig:ch2.self-adjoint-closure-pathway}
\end{figure}

Figure~\ref{fig:ch2.self-adjoint-closure-pathway} connects boundary physics to
solver stability. Changing \(\mathcal{B}_\Gamma\) while leaving
\(\mathfrak{a}_\omega\) formally unchanged can destroy coercivity or symmetry;
that sensitivity was already visible in the resolvent-branch illustration of
Subsection~\ref{sec:ch212} \cite{stratton1941,harrington2001}.

Admissibility from Subsection~\ref{sec:ch01.0.3} embeds into operator domains.
The Chapter~\ref{ch:01} finite-energy class
Eq.~\eqref{eq:ch1.admissible-space} requires
\((\E,\H)\in L^{2}(\Omega)^{3}\times L^{2}(\Omega)^{3}\) with compatible boundary
and source data \cite{stratton1941}. The Chapter~\ref{ch:02} refinement is
\begin{equation}
\label{eq:ch2.admissibility-embedding}
\mathcal{X}_{\mathrm{ch1}}(\Omega)\cap\mathcal{U}_\omega
\subseteq
\mathcal{F}_\Omega
\subseteq
\mathcal{D}(\mathcal{L}_\omega),
\end{equation}
with strict inclusions in general because \(\HCurl\) control and trace membership
strengthen \(L^{2}\) admissibility without replacing
Eq.~\eqref{eq:ch1.inheritance-admissibility} \cite{stratton1941}. Equation
\eqref{eq:ch2.admissibility-embedding} specializes
Eq.~\eqref{eq:ch2.inheritance-operator-domain}: when
\(\mathcal{D}_{\mathrm{ch1}}=1\), impressed data lie in \(\mathcal{J}_\Omega\) and
closed solutions lie in \(\mathcal{F}_\Omega\). Operator admissibility adds
graph-norm control:
\begin{equation}
\label{eq:ch2.graph-admissibility-criterion}
\mathbf{u}_\omega\in\mathcal{F}_\Omega
\ \Longrightarrow\
\|\mathbf{u}_\omega\|_{\mathrm{gr}(\mathcal{L}_\omega)}<\infty
\ \text{when}\ \mathcal{L}_\omega[\mathbf{u}_\omega]\ \text{is defined}.
\end{equation}
Inequality~\eqref{eq:ch2.graph-admissibility-criterion} is the quantitative form of
``physically realizable'' field pairs for variational solvers
\cite{harrington2001}.

Source-side admissibility remains coupled to
Eq.~\eqref{eq:ch2.source-space-refined} and
Eq.~\eqref{eq:ch1.source-admissibility} by reference. Closure operators must not
introduce incompatible traces:
\begin{equation}
\label{eq:ch2.closure-admissibility-lock}
\mathcal{B}_\Gamma\ \text{admissible}
\iff
\mathcal{R}(\gamma_\tau|_{\mathcal{F}_\Omega})\supseteq
\mathcal{D}(\mathcal{B}_\Gamma).
\end{equation}
Failure of Eq.~\eqref{eq:ch2.closure-admissibility-lock} means prescribed
\(\mathbf{g}_\Gamma\) is not the trace of any \(\HCurl\) field, reproducing the
trace-mismatch illustration of Subsection~\ref{sec:ch211} \cite{stratton1941}.

Exterior domains: outgoing closure from Section~\ref{sec:ch243} defines a
\emph{physically self-adjoint} branch rather than a Hilbert self-adjoint operator
on all of \(\mathcal{U}_\omega\):
\begin{equation}
\label{eq:ch2.outgoing-nonselfadjoint-split}
\mathcal{U}_\omega
=
\mathcal{U}_{\mathrm{out}}\oplus\mathcal{U}_{\mathrm{re}},
\qquad
\mathcal{L}_\omega\ \text{coercive on}\ \mathcal{U}_{\mathrm{out}},
\end{equation}
with reactive storage content in \(\mathcal{U}_{\mathrm{re}}\) controlled by
regime criteria of Section~\ref{sec:ch01.2.4} rather than by full coercivity on
\(\mathcal{U}_\omega\) \cite{stratton1941,harrington2001}. Conflating
Eq.~\eqref{eq:ch2.friedrichs-extension} with exterior radiation problems produces
incorrect cavity-like spectra in open regions.

\paragraph{Worked illustration (cavity at resonance).}
In a perfectly conducting cavity, \(\ker\mathcal{L}_\omega\) contains a standing
mode at \(\omega=\omega_c\). Attempting to solve
Eq.~\eqref{eq:ch2.maxwell-weak-form} without quotienting yields infinite
\(\|\mathbf{u}_\omega\|_{\mathcal{U}}\) for finite forcing. Quotienting by
\(\ker\mathcal{L}_\omega\) or adding loss through passive
Eq.~\eqref{eq:ch2.passive-impedance-class} restores unique
Eq.~\eqref{eq:ch2.lax-milgram-wellposed} solutions; feed coupling at \(\omega_c\)
without loss is physically singular \cite{stratton1941,harrington2001}.

\paragraph{Worked illustration (active impedance).}
If \(\Re\{Z_s\}<0\) on a portion of \(\Gamma_Z\), Eq.~\eqref{eq:ch2.passive-impedance-class}
fails and \(\mathfrak{a}_\omega\) is no longer coercive in the energy norm. A
simulation may still converge numerically while violating the passivity screen of
Section~\ref{sec:ch01.1.2}; power budgets imported from
Section~\ref{sec:ch01.1.4} then lose their sign-definite interpretation
\cite{stratton1941}. Admissibility requires revisiting \(\mathcal{B}_\Gamma\), not
post-processing fields alone.

\paragraph{Closing summary.}
Subsection~\ref{sec:ch214} establishes the Chapter~\ref{ch:02} HOME for coercivity
on the cavity quotient Eq.~\eqref{eq:ch2.coercivity-quotient}, passive impedance
symmetry Eq.~\eqref{eq:ch2.passive-impedance-class}, Friedrichs closure
Eq.~\eqref{eq:ch2.friedrichs-extension}, Lax--Milgram well-posedness
Eq.~\eqref{eq:ch2.lax-milgram-wellposed}, and admissibility embedding
Eqs.~\eqref{eq:ch2.admissibility-embedding}--\eqref{eq:ch2.closure-admissibility-lock}.
Chapter~\ref{ch:01} finite-energy criteria are embedded, not re-proved. Section~\ref{sec:ch21}
now supplies the complete function-space and operator foundation for
Sections~\ref{sec:ch22}--\ref{sec:ch29}.

With Subsections~\ref{sec:ch211}--\ref{sec:ch214} complete, Section~\ref{sec:ch21}
has moved from minimal declarations in Section~\ref{sec:ch20} to a closed operator
calculus on Sobolev spaces: traces, adjoints, compactness, rank-deficient readouts,
and self-adjoint truncations compatible with Chapter~\ref{ch:01} admissibility.
Section~\ref{sec:ch22} applies the same spaces to time--frequency Maxwell
equivalences and propagating-versus-evanescent splits; Green-function inverses in
Section~\ref{sec:ch23} and outgoing selection in Section~\ref{sec:ch24} inherit
\(\mathcal{D}(\mathcal{L}_\omega)\) and \(\mathcal{F}_\Omega\) as fixed input
classes \cite{stratton1941,harrington2001}.

\section{Time- and Frequency-Domain Maxwell Operators}
\label{sec:ch22}

Section~\ref{sec:ch21} fixed the Sobolev, trace, and self-adjoint domains on which
\(\mathcal{L}_\omega\) and \(\mathcal{S}_\omega\) act. Section~\ref{sec:ch22}
develops the time--frequency equivalences that connect those harmonic operators to
causal evolution, Helmholtz symbols, and propagating-versus-evanescent splits
\cite{stratton1941,harrington2001}. The \(\exp(-\jj\omega t)\) phasor convention
locked in Chapter~\ref{ch:01} remains in force; Section~\ref{sec:ch22} records when
that reduction is valid, how retarded time-domain maps relate to
\(\mathcal{R}_\omega\), and which spectral components carry far-zone transport
content importable from Section~\ref{sec:ch01.2.4}.

Harmonic Maxwell theory in Section~\ref{sec:ch20} treated \(\omega>0\) as a fixed
parameter. Time-domain formulations restore the full operator pencil
\(\partial/\partial t\) and clarify why frequency-domain solvers implement causal
limits when inversion is performed correctly. Subsection~\ref{sec:ch221} gives the
Chapter~\ref{ch:02} HOME for time-domain evolution on the same admissible classes
as Eq.~\eqref{eq:ch2.admissible-field-space}. Subsection~\ref{sec:ch222} packages
frequency-domain Maxwell and curl--curl Helmholtz operators without re-displaying
Eq.~\eqref{eq:ch1.maxwell-local}; full modal separation of vector fields is deferred
to Section~\ref{sec:ch33} in Chapter~\ref{ch:03} \cite{stratton1941}. Subsection
~\ref{sec:ch223} proves the equivalence between retarded Green propagation and
\(\mathcal{S}_\omega\) previewed in Eq.~\eqref{eq:ch2.causality-preview}.
Subsection~\ref{sec:ch224} applies the reactive-radiative partition of
Section~\ref{sec:ch01.2.4} to propagating, evanescent, and non-propagating
subspaces in operator form \cite{harrington2001}.

\begin{figure}[t]
\centering
\ChOneFigBlock{%
\begin{tikzpicture}[ch1fig, node distance=7mm and 10mm]
  \node[ch1box] (td) {Time domain\\$\partial_t$ evolution};
  \node[ch1box, right=14mm of td] (fd) {Frequency\\$\omega$-domain};
  \node[ch1box, right=14mm of fd] (sp) {Spectral split\\prop./evan.};
  \node[ch1box, below=12mm of fd] (op) {$\mathcal{L}_\omega,\ \mathcal{S}_\omega$};
  \draw[ch1flow] (td) -- node[above,font=\scriptsize] {FT} (fd);
  \draw[ch1flow] (fd) -- (sp);
  \draw[ch1flow] (fd) -- (op);
\end{tikzpicture}%
}
\caption{Section~\ref{sec:ch22} flow: time-domain evolution maps to harmonic
operators from Sections~\ref{sec:ch20}--\ref{sec:ch21}; spectral subspaces in
Subsection~\ref{sec:ch224} gate transport-qualified radiation content.}
\label{fig:ch2.time-frequency-roadmap}
\end{figure}

Figure~\ref{fig:ch2.time-frequency-roadmap} orders the four subsections. Time-domain
evolution supplies causality constraints that frequency-domain solvers must respect;
harmonic operators supply the symbols used in Green kernels of Section~\ref{sec:ch23}
and in radiation integrals of Section~\ref{sec:ch25}. Propagating and evanescent
labels are meaningful only after both Maxwell closure and regime qualification from
Chapter~\ref{ch:01} are satisfied \cite{stratton1941}.

Several constructions remain outside Section~\ref{sec:ch22}. Dyadic Green function
assembly belongs to Section~\ref{sec:ch23}; Sommerfeld outgoing selection to
Section~\ref{sec:ch24}; vector spherical harmonic separation to Chapter~\ref{ch:03}.
Section~\ref{sec:ch22} does not re-open function-space declarations from
Section~\ref{sec:ch21} or local Maxwell derivations from Chapter~\ref{ch:01}.
Helmholtz and curl--curl reductions developed in Subsection~\ref{sec:ch222} are
single-field tools: they simplify kernel construction but do not replace the
first-order \(\mathcal{L}_\omega\) packaging of
Eq.~\eqref{eq:ch2.maxwell-boundary-composite} when boundary rows and impedance
closures must be enforced jointly \cite{harrington2001,stratton1941}.

Material dispersion and broadband pulses can require frequency-dependent
\(\varepsilon(\omega)\), \(\mu(\omega)\) beyond the single-\(\omega\) harmonic
class treated here; within the locked harmonic band, time--frequency equivalence
states that \(\mathcal{S}_\omega\) is the \(\omega\)-slice of a causal time-domain
map whenever impressed data and \(\mathcal{B}_\Gamma\) are held fixed across the
transform \cite{stratton1941}. Changing termination or source switching in time
changes the operator pencil itself, not merely its frequency symbol.

Readers approaching from numerical codes should expect explicit statements of which
time--frequency route a solver implements. Finite-difference time-domain schemes
approximate Subsection~\ref{sec:ch221}; frequency-domain finite-element solvers
approximate Subsection~\ref{sec:ch222} on the domains of
Eq.~\eqref{eq:ch2.domain-refinement-L}. Modal and plane-wave decompositions used in
Subsection~\ref{sec:ch224} filter which components of
\(\mathcal{S}_\omega[\Jimp]\) may enter \(\mathcal{R}_\omega\) under the transport
criteria of Eq.~\eqref{eq:ch2.transport-qualified-radiation}
\cite{harrington2001}. The subsection sequence begins with time-domain evolution,
then harmonic reduction, then causality equivalence, then propagating and evanescent
subspaces as the spectral gate before Green-function inversion in
Section~\ref{sec:ch23}.


\subsection{Time-domain evolution}
\label{sec:ch221}
% TOC tag: [HOME]

Section~\ref{sec:ch21} fixed harmonic operator domains for
\(\mathcal{L}_\omega\) and \(\mathcal{S}_\omega\). Subsection~\ref{sec:ch221}
gives the Chapter~\ref{ch:02} HOME for \emph{time-domain} Maxwell evolution: the
first-order hyperbolic system, the evolution operator that maps initial and impressed
data to field trajectories, and the causal support constraints that harmonic solvers
must inherit \cite{stratton1941,harrington2001}. Local Maxwell balances remain those
of Eq.~\eqref{eq:ch1.ibvp-time-domain}; they are cited, not re-derived. Function
spaces from Section~\ref{sec:ch21} enter through time-dependent embeddings developed
below; Sobolev and trace declarations are not repeated \cite{stratton1941}.

Let \(\Omega\subset\mathbb{R}^{3}\) be a Lipschitz domain with boundary
\(\Gamma=\partial\Omega\). Real time-domain fields
\(\E(\mathbf{r},t)\), \(\H(\mathbf{r},t)\) use SI units and the same constitutive
class as Chapter~\ref{ch:01}. Collect the electromagnetic state at time \(t\) as
\begin{equation}
\label{eq:ch2.time-state-vector}
\mathbf{U}(t)
=
\begin{bmatrix}\E(\cdot,t)\\ \H(\cdot,t)\end{bmatrix}
\in\mathcal{V}(\Omega),
\end{equation}
where \(\mathcal{V}(\Omega)\) is the time-dependent product space of square-integrable
field pairs compatible with Eq.~\eqref{eq:ch1.admissible-space} at each frozen time
slice \cite{stratton1941}. Equation~\eqref{eq:ch2.time-state-vector} is the
time-domain analogue of Eq.~\eqref{eq:ch2.maxwell-state-vector}; phasor symbols are
reserved for the \(\exp(-\jj\omega t)\) reduction in Subsection~\ref{sec:ch222}.

The time-domain Maxwell system imported from Eq.~\eqref{eq:ch1.ibvp-time-domain} reads
in vacuum constitutive form as
\begin{equation}
\label{eq:ch2.time-maxwell-evolution}
\begin{aligned}
\curl\E(\mathbf{r},t)&=-\mu_0\,\partial_t\H(\mathbf{r},t),\\
\curl\H(\mathbf{r},t)&=\mathbf{j}_{\mathrm{imp}}(\mathbf{r},t)+\varepsilon_0\,\partial_t\E(\mathbf{r},t),\\
\divg\E(\mathbf{r},t)&=\frac{\rho(\mathbf{r},t)}{\varepsilon_0},\qquad
\divg\H(\mathbf{r},t)=0,
\end{aligned}
\end{equation}
with impressed current density \(\mathbf{j}_{\mathrm{imp}}(\mathbf{r},t)\) and charge
density \(\rho(\mathbf{r},t)\). Equation~\eqref{eq:ch2.time-maxwell-evolution}
packages local dynamics as a first-order system in \((\E,\H)\) and \(\partial_t\);
the hyperbolic coupling is the time-domain face of the curl structure in
Eq.~\eqref{eq:ch1.maxwell-operator-block} \cite{stratton1941}. Charge continuity
requires
\begin{equation}
\label{eq:ch2.charge-continuity-time}
\divg\mathbf{j}_{\mathrm{imp}}(\mathbf{r},t)+\partial_t\rho(\mathbf{r},t)=0
\end{equation}
in the sense of distributions when \(\rho\) is present \cite{stratton1941}. Pairs
\((\mathbf{j}_{\mathrm{imp}},\rho)\) violating
Eq.~\eqref{eq:ch2.charge-continuity-time} lie outside the time-domain admissible
class parallel to Eq.~\eqref{eq:ch2.source-admissibility}.

\begin{figure}[t]
\centering
\ChOneFigBlock{%
\begin{tikzpicture}[ch1fig, node distance=7mm and 10mm]
  \node[ch1box] (ic) {Initial data\\$(\E_0,\H_0)$};
  \node[ch1box, right=12mm of ic] (src) {Impressed\\$\mathbf{j}_{\mathrm{imp}}(\cdot,t)$};
  \node[ch1box, right=12mm of src] (ev) {Evolution\\$\mathcal{T}(t)$};
  \node[ch1box, right=12mm of ev] (st) {State $\mathbf{U}(t)$};
  \draw[ch1flow] (ic) -- (ev);
  \draw[ch1flow] (src) -- (ev);
  \draw[ch1flow] (ev) -- (st);
\end{tikzpicture}%
}
\caption{Time-domain evolution: initial data and impressed sources drive the
semigroup \(\mathcal{T}(t)\) to field state \(\mathbf{U}(t)\) under boundary class
\(\mathcal{B}_\Gamma\).}
\label{fig:ch2.time-evolution-pipeline}
\end{figure}

Figure~\ref{fig:ch2.time-evolution-pipeline} is the causal input--output chain for
Subsection~\ref{sec:ch221}. Boundary data enter through the same
\(\mathcal{B}_\Gamma\) of Eq.~\eqref{eq:ch1.boundary-data-operator}, now evaluated
on time-dependent traces \(\gamma_\tau[\mathbf{U}(t)]\in\mathcal{T}_\Gamma\) for
each \(t\ge 0\) \cite{stratton1941,harrington2001}. The closed time-domain problem
is to find \(\mathbf{U}(t)\) such that
\begin{equation}
\label{eq:ch2.time-ibvp-package}
\begin{cases}
\partial_t\mathbf{U}=\mathcal{A}\,\mathbf{U}+\mathbf{F}(t), & \mathbf{r}\in\Omega,\ t>0,\\
\mathbf{U}(0)=\mathbf{U}_0, & \mathbf{r}\in\Omega,\\
\mathcal{B}_\Gamma[\mathbf{U}(t)]=\mathbf{g}_\Gamma(t), & \mathbf{r}\in\Gamma,\ t\ge 0,
\end{cases}
\end{equation}
where \(\mathcal{A}\) is the Maxwell generator induced by
Eq.~\eqref{eq:ch2.time-maxwell-evolution} on
\(\mathcal{D}(\mathcal{A})\subset\mathcal{V}(\Omega)\) and
\(\mathbf{F}(t)\) collects impressed-source rows \cite{stratton1941}. Equation
\eqref{eq:ch2.time-ibvp-package} is the time-domain counterpart of
Eq.~\eqref{eq:ch2.closed-maxwell-bvp}; uniqueness on a fixed branch is inherited
from Section~\ref{sec:ch01.1.3} when data classes coincide with the harmonic limit
\cite{stratton1941}.

Define the evolution operator (semigroup) \(\mathcal{T}(t)\) by
\begin{equation}
\label{eq:ch2.time-semigroup-definition}
\mathcal{T}(t):\mathbf{U}_0\longmapsto \mathbf{U}(t),
\qquad
\mathcal{T}(0)=\mathrm{id},\qquad
\mathcal{T}(t+s)=\mathcal{T}(t)\mathcal{T}(s),
\end{equation}
when \(\mathbf{F}\equiv\mathbf{0}\) and \(\mathcal{B}_\Gamma\) is homogeneous. With
forcing,
\begin{equation}
\label{eq:ch2.time-evolution-forced}
\mathbf{U}(t)
=
\mathcal{T}(t)\mathbf{U}_0
+
\int_0^t \mathcal{T}(t-s)\,\mathbf{F}(s)\,ds,
\end{equation}
the Duhamel representation on \(\mathcal{D}(\mathcal{A})\)
\cite{stratton1941,harrington2001}. Equation~\eqref{eq:ch2.time-evolution-forced}
makes causality explicit: \(\mathbf{U}(t)\) depends only on
\(\mathbf{F}(s)\) for \(s\le t\). The time-domain \emph{solution map} sends impressed
data and initial state to trajectories:
\begin{equation}
\label{eq:ch2.time-domain-solution-map}
\mathcal{S}_{t}:\big(\mathbf{U}_0,\mathbf{j}_{\mathrm{imp}}\big)\longmapsto
\{\mathbf{U}(t)\}_{t\ge 0}.
\end{equation}
Harmonic \(\mathcal{S}_\omega\) of Eq.~\eqref{eq:ch2.maxwell-solution-map} is a
single-frequency slice of Eq.~\eqref{eq:ch2.time-domain-solution-map} when all data
are proportional to \(\exp(-\jj\omega t)\) and \(\mathcal{B}_\Gamma\) is fixed
\cite{stratton1941}; the equivalence is proved in Subsection~\ref{sec:ch223}.

Retarded representation specializes Eq.~\eqref{eq:ch2.time-evolution-forced} when
Green kernels exist:
\begin{equation}
\label{eq:ch2.retarded-convolution-form}
\E(\mathbf{r},t)
=
\E_{\mathrm{hom}}(\mathbf{r},t)
-\mu_0\,\partial_t
\int_0^t\!\!\int_\Omega
\mathcal{G}(\mathbf{r},\mathbf{r}',t-t')\cdot
\mathbf{j}_{\mathrm{imp}}(\mathbf{r}',t')\,d^{3}r'\,dt',
\end{equation}
with \(\mathcal{G}(\mathbf{r},\mathbf{r}',t-t')=\mathbf{0}\) for \(t-t'<0\)
\cite{stratton1941}. Equation~\eqref{eq:ch2.retarded-convolution-form} refines the
preview Eq.~\eqref{eq:ch2.causality-preview} without constructing \(\mathcal{G}\)
here; kernel ownership is Section~\ref{sec:ch23}. The support condition
\begin{equation}
\label{eq:ch2.causal-support-constraint}
\operatorname{supp}_{(\mathbf{r},t)}\E
\subseteq
\big\{(\mathbf{r},t): t\ge t_{\mathrm{ret}}(\mathbf{r})\big\}
\end{equation}
expresses finite propagation speed \(c_0=1/\sqrt{\mu_0\varepsilon_0}\) on vacuum
regions \cite{stratton1941,harrington2001}. Figure~\ref{fig:ch2.causal-light-cone}
schematizes Eq.~\eqref{eq:ch2.causal-support-constraint}.

\begin{figure}[t]
\centering
\ChOneFigBlock{%
\begin{tikzpicture}[ch1fig]
  \draw[ch1flow] (0,0) -- (2.2,2.2);
  \draw[ch1flow] (0,0) -- (2.2,-2.2);
  \node[ch1oncanvas, right=2mm] at (2.2,2.2) {$t=t_{\mathrm{ret}}$};
  \node[ch1box] at (0,0) {Source event};
\end{tikzpicture}%
}
\caption{Causal support: field at \((\mathbf{r},t)\) depends on sources inside the
backward light cone; retarded integration uses \(t'\le t\) only.}
\label{fig:ch2.causal-light-cone}
\end{figure}

Energy balance in time domain imports the Poynting and storage laws of
Section~\ref{sec:ch01.1.4}. With volume energy density
\begin{equation}
\label{eq:ch2.time-energy-density}
u(\mathbf{r},t)
=
\frac{\varepsilon_0}{2}\|\E(\mathbf{r},t)\|^2
+\frac{\mu_0}{2}\|\H(\mathbf{r},t)\|^2,
\end{equation}
integration over a control volume \(V\subset\Omega\) yields
\begin{equation}
\label{eq:ch2.time-poynting-balance}
\frac{d}{dt}\int_V u\,dV
+\int_{\partial V}\mathbf{S}\cdot\hat{\mathbf{n}}\,dS
=
-\int_V \E\cdot\mathbf{j}_{\mathrm{imp}}\,dV,
\end{equation}
with \(\mathbf{S}=\E\times\H\) \cite{stratton1941}. Equation
\eqref{eq:ch2.time-poynting-balance} is not re-derived; it supplies the physical
reading of Eq.~\eqref{eq:ch2.time-evolution-forced}: impressed work on the right
drives stored and exported energy on the left. Reactive exchange with boundaries
appears in the surface term and motivates the regime split of
Subsection~\ref{sec:ch224} \cite{harrington2001}.

The bridge to harmonic operators is the substitution
\(\E(\mathbf{r},t)=\Re\{\underline{\E}(\mathbf{r})\eim\}\) with
\(\underline{\H}\) analogous, yielding
\begin{equation}
\label{eq:ch2.time-to-frequency-bridge}
\partial_t \longrightarrow -\jj\omega,
\qquad
\mathcal{A}\longrightarrow \mathcal{M}_\omega,
\qquad
\mathcal{S}_t\big|_{\omega}\longrightarrow \mathcal{S}_\omega,
\end{equation}
on admissible data of Eq.~\eqref{eq:ch2.admissibility-embedding}
\cite{stratton1941}. Equation~\eqref{eq:ch2.time-to-frequency-bridge} is symbolic;
Subsection~\ref{sec:ch223} states equivalence conditions under which frequency-domain
solvers implement causal time-domain maps. Exterior domains require outgoing
boundary conditions in time domain compatible with
Eq.~\eqref{eq:ch1.radiation-trace-condition}; those closures are sharpened in
Section~\ref{sec:ch243} \cite{stratton1941,harrington2001}.

Time-domain admissibility couples Eq.~\eqref{eq:ch1.admissible-space} at each time
slice to Eq.~\eqref{eq:ch2.charge-continuity-time}:
\begin{equation}
\label{eq:ch2.time-admissibility-coupling}
\mathbf{U}(t)\in\mathcal{F}_\Omega\ \text{a.e.\ }t
\ \wedge\
(\mathbf{j}_{\mathrm{imp}},\rho)\ \text{satisfy charge continuity}
\ \Longrightarrow\
\mathbf{U}\in\mathcal{D}(\mathcal{S}_{t}).
\end{equation}
Finite-difference time-domain schemes discretize
Eq.~\eqref{eq:ch2.time-maxwell-evolution} on staggered grids; stability requires
discrete analogues of Eq.~\eqref{eq:ch2.causal-support-constraint} and consistent
boundary enforcement of \(\mathcal{B}_\Gamma\) \cite{harrington2001}. Subsection
~\ref{sec:ch221} states the continuous targets those schemes approximate.

\paragraph{Worked illustration (switch-on source).}
Let \(\mathbf{j}_{\mathrm{imp}}(\mathbf{r},t)=\mathbf{j}_0(\mathbf{r})H(t)\) with
\(H\) the Heaviside step and \(\mathbf{j}_0\in\mathcal{J}_\Omega\) at the frozen-time
slice. Then Eq.~\eqref{eq:ch2.time-evolution-forced} produces a transient
\(\mathbf{U}(t)\) that is not proportional to \(\exp(-\jj\omega t)\); evaluating a
single harmonic \(\mathcal{S}_\omega\) at one \(\omega\) cannot reproduce the full
trajectory without Fourier synthesis. Causality requires
Eq.~\eqref{eq:ch2.retarded-convolution-form} with \(t'\le t\); using advanced Green
functions violates Eq.~\eqref{eq:ch2.causal-support-constraint} even when
Eq.~\eqref{eq:ch2.time-maxwell-evolution} holds in the interior
\cite{stratton1941,harrington2001}.

\paragraph{Worked illustration (boundary switching).}
Suppose \(\mathcal{B}_\Gamma\) changes from open to short-circuit at \(t=t_s\).
Then \(\mathcal{A}\) itself changes at \(t_s\); Eq.~\eqref{eq:ch2.time-semigroup-definition}
applies piecewise on intervals separated by \(t_s\). A frequency-domain solve at fixed
\(\mathcal{B}_\Gamma\) cannot represent the switched problem without re-solving or
harmonic superposition over the entire time axis \cite{stratton1941}. Operator
identity, not mesh resolution, is the failure mode.

\paragraph{Closing summary.}
Subsection~\ref{sec:ch221} establishes the Chapter~\ref{ch:02} HOME for the
time-domain state Eq.~\eqref{eq:ch2.time-state-vector}, evolution system
Eq.~\eqref{eq:ch2.time-maxwell-evolution}, IBVP packaging
Eq.~\eqref{eq:ch2.time-ibvp-package}, semigroup and Duhamel forms
Eqs.~\eqref{eq:ch2.time-semigroup-definition}--\eqref{eq:ch2.time-evolution-forced},
solution map Eq.~\eqref{eq:ch2.time-domain-solution-map}, retarded convolution
Eq.~\eqref{eq:ch2.retarded-convolution-form}, and causal support
Eq.~\eqref{eq:ch2.causal-support-constraint}. Subsection~\ref{sec:ch222} develops
harmonic Maxwell and Helmholtz operators as the \(\omega\)-domain image of this
evolution under the bridge Eq.~\eqref{eq:ch2.time-to-frequency-bridge}.

\subsection{Frequency-domain Maxwell operators}
\label{sec:ch222}
% TOC tag: [HOME]

Subsection~\ref{sec:ch221} fixed time-domain evolution and the bridge
Eq.~\eqref{eq:ch2.time-to-frequency-bridge}. Subsection~\ref{sec:ch222} gives the
Chapter~\ref{ch:02} HOME for harmonic Maxwell operators and curl--curl Helmholtz
symbols on the spaces of Section~\ref{sec:ch21} \cite{stratton1941,harrington2001}.
The first-order system Eq.~\eqref{eq:ch1.maxwell-operator-block} and composite
\(\mathcal{L}_\omega\) of Eq.~\eqref{eq:ch2.maxwell-boundary-composite} are
presupposed; local Maxwell derivations from Chapter~\ref{ch:01} are not repeated.
The electric-wave operator Eq.~\eqref{eq:ch2.electric-wave-operator} was recorded in
Subsection~\ref{sec:ch202}; here it is embedded in domain declarations and extended
by magnetic and material-weighted counterparts. Full vector-field separation is
deferred to Section~\ref{sec:ch33} \cite{stratton1941}.

Under the locked \(\exp(-\jj\omega t)\) convention, time derivatives map to
\begin{equation}
\label{eq:ch2.frequency-substitution-rule}
\partial_t \longrightarrow -\jj\omega,
\qquad
\frac{\partial^2}{\partial t^2}\longrightarrow -\omega^2,
\end{equation}
on phasor fields \(\E(\mathbf{r})\), \(\H(\mathbf{r})\) and impressed data
\(\Jimp(\mathbf{r})\), \(\rho(\mathbf{r})\) \cite{stratton1941}. Equation
\eqref{eq:ch2.frequency-substitution-rule} specializes
Eq.~\eqref{eq:ch2.time-to-frequency-bridge} and converts
Eq.~\eqref{eq:ch1.ibvp-time-domain} to the harmonic BVP of
Eq.~\eqref{eq:ch1.ibvp-harmonic-bvp} by reference. The harmonic Maxwell residual
operator \(\mathcal{M}_\omega\) of Eq.~\eqref{eq:ch2.maxwell-driven-operator}
acts on \(\mathbf{u}_\omega=(\E,\H)\in\mathcal{U}_\omega\):
\begin{equation}
\label{eq:ch2.harmonic-maxwell-residual}
\mathcal{M}_\omega[\mathbf{u}_\omega]
=
\begin{bmatrix}
\curl\E+\jj\omega\mu_0\H\\
\curl\H-\jj\omega\varepsilon_0\E-\Jimp\\
\divg(\varepsilon_0\E)-\rho\\
\divg(\mu_0\H)
\end{bmatrix},
\end{equation}
row-aligned with Eq.~\eqref{eq:ch1.maxwell-operator-block} \cite{stratton1941}.
Equation~\eqref{eq:ch2.harmonic-maxwell-residual} is the frequency-domain packaging
of Eq.~\eqref{eq:ch2.time-maxwell-evolution}; closure and traces remain in
\(\mathcal{L}_\omega[\mathbf{u}_\omega]=[\mathbf{s}_\omega,\mathbf{g}_\Gamma]^{\mathsf{T}}\)
from Eq.~\eqref{eq:ch2.closed-maxwell-bvp}.

\begin{figure}[t]
\centering
\ChOneFigBlock{%
\begin{tikzpicture}[ch1fig, node distance=7mm and 10mm]
  \node[ch1box] (fo) {First-order\\$\mathcal{M}_\omega$};
  \node[ch1box, right=14mm of fo] (el) {$\mathcal{W}_{E,\omega}$};
  \node[ch1box, below=10mm of fo] (mag) {$\mathcal{W}_{H,\omega}$};
  \node[ch1box, right=14mm of el] (sym) {Symbol $k(\mathbf{r})$};
  \draw[ch1flow] (fo) -- node[above,font=\scriptsize] {eliminate $\H$} (el);
  \draw[ch1flow] (fo) -- node[left,font=\scriptsize] {eliminate $\E$} (mag);
  \draw[ch1flow] (el) -- (sym);
\end{tikzpicture}%
}
\caption{Frequency-domain reduction ladder: first-order \(\mathcal{M}_\omega\)
yields electric and magnetic Helmholtz operators when partner fields are eliminated
under stated domain and gauge constraints.}
\label{fig:ch2.frequency-operator-ladder}
\end{figure}

Figure~\ref{fig:ch2.frequency-operator-ladder} orders the reductions below.
Single-field Helmholtz forms equivalent to
Eqs.~\eqref{eq:ch1.vacuum-wave-source-driven}--\eqref{eq:ch1.vacuum-curl-second-order}
were established in Chapter~\ref{ch:01} and are invoked by reference; elimination
steps are not repeated \cite{stratton1941}. The Chapter~\ref{ch:02} electric-wave
operator Eq.~\eqref{eq:ch2.electric-wave-operator} is the vacuum electric Helmholtz
map on \(\mathcal{D}(\mathcal{W}_{E,\omega})\subset\HCurl(\Omega)\):
\begin{equation}
\label{eq:ch2.electric-helmholtz-domain}
\mathcal{D}(\mathcal{W}_{E,\omega})
=
\big\{\E\in\HCurl(\Omega):
\mathcal{W}_{E,\omega}\E\in L^{2}(\Omega)^{3}\big\}.
\end{equation}
The magnetic counterpart is
\begin{equation}
\label{eq:ch2.magnetic-wave-operator}
\mathcal{W}_{H,\omega}\H\equiv
\curl\!\left(\varepsilon_0^{-1}\curl\H\right)-\omega^2\mu_0\H
= \jj\omega\mu_0\Jimp-\varepsilon_0^{-1}\curl\Jimp,
\end{equation}
on \(\mathcal{D}(\mathcal{W}_{H,\omega})\subset\HCurl(\Omega)\) analogously
\cite{stratton1941,harrington2001}. Equations
\eqref{eq:ch2.electric-wave-operator} and~\eqref{eq:ch2.magnetic-wave-operator}
are dual single-field symbols; using one without the other is valid only when the
eliminated partner is recoverable from curl relations without ambiguity.

In piecewise smooth media with \(\varepsilon(\mathbf{r})\), \(\mu(\mathbf{r})\),
\begin{equation}
\label{eq:ch2.material-helmholtz-electric}
\mathcal{W}_{E,\omega}^{(\varepsilon,\mu)}\E\equiv
\curl\!\left(\mu^{-1}\curl\E\right)-\omega^2\varepsilon\E
= -\jj\omega\mu\Jimp+\mu^{-1}\curl\mathbf{J}^{\mathrm{(mag)}}_{\mathrm{eff}},
\end{equation}
\begin{equation}
\label{eq:ch2.material-helmholtz-magnetic}
\mathcal{W}_{H,\omega}^{(\varepsilon,\mu)}\H\equiv
\curl\!\left(\varepsilon^{-1}\curl\H\right)-\omega^2\mu\H
= \jj\omega\mu\Jimp-\varepsilon^{-1}\curl\Jimp,
\end{equation}
with interface jumps on \(\mathcal{D}(\mathcal{W}_{\cdot,\omega}^{(\varepsilon,\mu})\)
fixed by Eq.~\eqref{eq:ch1.interface-conditions} \cite{stratton1941}. Material
weights modify symbols but not the first-order coupling pattern of
Figure~\ref{fig:ch2.first-order-curl-coupling}; multi-region Green operators in
Section~\ref{sec:ch234} invert Eqs.~\eqref{eq:ch2.material-helmholtz-electric}--\eqref{eq:ch2.material-helmholtz-magnetic}
on composite domains \cite{harrington2001}.

The local Helmholtz symbol is
\begin{equation}
\label{eq:ch2.helmholtz-symbol}
k(\mathbf{r})=\omega\sqrt{\mu(\mathbf{r})\varepsilon(\mathbf{r})},
\qquad
k_0=\omega\sqrt{\mu_0\varepsilon_0},
\end{equation}
from Eq.~\eqref{eq:ch1.vacuum-impedance-wavenumber} in uniform vacuum
\cite{stratton1941}. In source-free regions,
Eqs.~\eqref{eq:ch1.vacuum-helmholtz-pair} apply with
\(k(\mathbf{r})\) replacing \(k_0\). The symbol \(k(\mathbf{r})\) labels
propagating versus evanescent behavior in Subsection~\ref{sec:ch224}; it is not a
substitute for outgoing selection on exterior domains \cite{harrington2001}.

Single-field reduction is admissible only under explicit constraints:
\begin{equation}
\label{eq:ch2.single-field-reduction-validity}
\text{eliminate }\H
\ \Longrightarrow\
\Jimp\in\HCurl(\Omega),\ \rho\ \text{compatible},\ 
\mathcal{B}_\Gamma\ \text{expressible in }\E\text{-traces alone},
\end{equation}
with the magnetic elimination requiring the parallel conditions on \(\H\)-traces
\cite{stratton1941}. Failure of Eq.~\eqref{eq:ch2.single-field-reduction-validity}
---for example impedance boundaries coupling both tangential components---forces
retention of the first-order \(\mathcal{L}_\omega\) system
Eq.~\eqref{eq:ch2.maxwell-boundary-composite} \cite{harrington2001}.

Forcing maps between impressed data and Helmholtz right-hand sides are linear:
\begin{equation}
\label{eq:ch2.helmholtz-forcing-map}
\mathbf{f}_{E,\omega}
=
-\jj\omega\mu\Jimp+\mu^{-1}\curl\mathbf{J}^{\mathrm{(mag)}}_{\mathrm{eff}},
\qquad
\mathcal{W}_{E,\omega}^{(\varepsilon,\mu)}\E=\mathbf{f}_{E,\omega},
\end{equation}
with \(\mathbf{f}_{E,\omega}\in L^{2}(\Omega)^{3}\) when
\(\Jimp\in\mathcal{J}_\Omega\) from Eq.~\eqref{eq:ch2.source-space-refined}
\cite{stratton1941}. The harmonic solution operator on the electric Helmholtz branch
is
\begin{equation}
\label{eq:ch2.helmholtz-solution-operator}
\mathcal{G}_{E,\omega}:\mathbf{f}_{E,\omega}\longmapsto\E,
\qquad
\mathcal{W}_{E,\omega}^{(\varepsilon,\mu)}\mathcal{G}_{E,\omega}=\mathrm{id}
\end{equation}
on the physical resolvent branch selected by \(\mathcal{B}_\Gamma\); explicit
\(\mathcal{G}_{E,\omega}\) construction is Section~\ref{sec:ch23} HOME
\cite{harrington2001,stratton1941}. Composition with the full map yields
\begin{equation}
\label{eq:ch2.frequency-solution-composition}
\mathcal{S}_\omega[\Jimp]
=
\big(\mathcal{G}_{E,\omega}\mathbf{f}_{E,\omega},\,
-\frac{1}{\jj\omega\mu}\curl\E\big)
\quad\text{when }\H\ \text{recovered from Amp\`ere row},
\end{equation}
linking Eq.~\eqref{eq:ch2.maxwell-solution-map} to Helmholtz inversion without
re-displaying Eq.~\eqref{eq:ch2.solution-map-resolvent} \cite{stratton1941}.

Spectral viewpoint preview: on bounded domains with discrete cavity spectra,
\(\mathcal{W}_{E,\omega}^{(\varepsilon,\mu)}\) admits eigenpairs
\((\lambda_n,\mathbf{e}_n)\) accumulating at infinity because
\(\HCurl(\Omega)\hookrightarrow L^{2}(\Omega)^{3}\) compactly
(Eq.~\eqref{eq:ch2.embedding-chain}) \cite{stratton1941}. On exterior domains,
continuous spectrum replaces discrete eigenvalues; Section~\ref{sec:ch263} develops
Fredholm and spectral consequences. Vector spherical harmonic separation of
eigenmodes is deferred to Section~\ref{sec:ch33} \cite{harrington2001}.

\begin{figure}[t]
\centering
\ChOneFigBlock{%
\begin{tikzpicture}[ch1fig, node distance=8mm and 12mm]
  \node[ch1box] (l) {$\mathcal{L}_\omega$\\first-order};
  \node[ch1box, right=16mm of l] (w) {$\mathcal{W}_{E,\omega}$\\curl--curl};
  \node[ch1box, below=11mm of w] (b) {Boundary row\\$\mathcal{B}_\Gamma$};
  \draw[ch1flow] (l) -- node[above,font=\scriptsize] {eliminate} (w);
  \draw[ch1flow, dashed] (b) -- (l);
  \draw[ch1flow, dashed] (b) -- (w);
\end{tikzpicture}%
}
\caption{First-order \(\mathcal{L}_\omega\) retains coupled boundary rows; single-field
\(\mathcal{W}_{E,\omega}\) is valid only when \(\mathcal{B}_\Gamma\) can be
projected without losing impedance coupling.}
\label{fig:ch2.first-order-vs-helmholtz}
\end{figure}

Figure~\ref{fig:ch2.first-order-vs-helmholtz} warns against replacing boundary-value
problems for \(\mathcal{L}_\omega\) with Helmholtz solves that enforce volume rows
only. Impedance and outgoing closures modify both volume and trace operators
simultaneously \cite{stratton1941,harrington2001}.

Weak form of the electric Helmholtz operator uses the same sesquilinear structure as
Eq.~\eqref{eq:ch2.maxwell-weak-form} restricted to \(\E\)-only test spaces when
reduction is valid \cite{stratton1941}. Coercivity on the quotient by cavity modes
from Eq.~\eqref{eq:ch2.coercivity-quotient} applies to
\(\mathcal{W}_{E,\omega}^{(\varepsilon,\mu)}\) on bounded domains with passive
\(\mathcal{B}_\Gamma\); Subsection~\ref{sec:ch214} owns the self-adjoint realizations
\cite{harrington2001}.

\paragraph{Worked illustration (vacuum dipole).}
For compact \(\Jimp\) in uniform vacuum with outgoing exterior closure,
Eq.~\eqref{eq:ch2.electric-wave-operator} with
\(\mathbf{f}_{E,\omega}=-\jj\omega\mu_0\Jimp\) yields \(\E\) whose far zone scales
as \(e^{-\jj k_0 r}/r\). Recovering
\(\H=-(\jj\omega\mu_0)^{-1}\curl\E\) closes the first-order pair without solving
Eq.~\eqref{eq:ch2.magnetic-wave-operator} separately; the two Helmholtz routes must
agree on \(\mathcal{D}(\mathcal{L}_\omega)\) \cite{stratton1941,harrington2001}.
Green-function evaluation is deferred to Section~\ref{sec:ch23}.

\paragraph{Worked illustration (impedance boundary).}
On a domain with passive Eq.~\eqref{eq:ch2.passive-impedance-class} closure,
eliminating \(\H\) to obtain Eq.~\eqref{eq:ch2.material-helmholtz-electric} alone
requires absorbing impedance into an equivalent surface operator on \(\E\)-traces.
If that projection is omitted, volume Helmholtz solutions satisfy
\(\mathcal{W}_{E,\omega}\E=\mathbf{f}_{E,\omega}\) but fail the coupled
\(\mathcal{B}_\Gamma\) row of \(\mathcal{L}_\omega\) \cite{stratton1941}. The
failure is structural: first-order closure cannot be dropped when impedance couples
both tangential components.

\paragraph{Closing summary.}
Subsection~\ref{sec:ch222} establishes the Chapter~\ref{ch:02} HOME for harmonic
residual packaging Eq.~\eqref{eq:ch2.harmonic-maxwell-residual}, electric and magnetic
Helmholtz operators Eqs.~\eqref{eq:ch2.electric-helmholtz-domain},
\eqref{eq:ch2.magnetic-wave-operator}, material-weighted symbols
Eqs.~\eqref{eq:ch2.material-helmholtz-electric}--\eqref{eq:ch2.material-helmholtz-magnetic},
reduction validity Eq.~\eqref{eq:ch2.single-field-reduction-validity}, and Helmholtz
solution composition Eqs.~\eqref{eq:ch2.helmholtz-forcing-map}--\eqref{eq:ch2.frequency-solution-composition}.
Subsection~\ref{sec:ch223} proves causality and time--frequency equivalence for
\(\mathcal{S}_\omega\) relative to \(\mathcal{S}_t\).

\subsection{Causality and time-frequency equivalence}
\label{sec:ch223}
% TOC tag: [HOME]

Subsections~\ref{sec:ch221}--\ref{sec:ch222} fixed time-domain evolution and harmonic
operators. Subsection~\ref{sec:ch223} gives the Chapter~\ref{ch:02} HOME proof that
\(\mathcal{S}_\omega\) and \(\mathcal{R}_\omega\) are the frequency-domain images of
causal time-domain maps when impressed data, media, and boundary class are held fixed
\cite{stratton1941,harrington2001}. The retarded preview
Eq.~\eqref{eq:ch2.causality-preview} and symbolic bridge
Eq.~\eqref{eq:ch2.time-to-frequency-bridge} are upgraded to an equivalence theorem;
Green-function construction remains in Section~\ref{sec:ch23} \cite{stratton1941}.

\paragraph{Causal support.}
A time-domain operator \(\mathcal{K}\) is \emph{causal} if
\begin{equation}
\label{eq:ch2.causal-operator-support}
\mathcal{K}[f](t)=0\ \text{for}\ t<t_0
\quad\text{whenever}\ 
\operatorname{supp}_t f\subseteq[t_0,\infty),
\end{equation}
for admissible test functions \(f\) \cite{stratton1941}. The retarded dyadic kernel
in Eq.~\eqref{eq:ch2.retarded-convolution-form} satisfies
\begin{equation}
\label{eq:ch2.retarded-kernel-support}
\mathcal{G}(\mathbf{r},\mathbf{r}',t-t')=\mathbf{0}\quad\text{for}\ t-t'<0,
\end{equation}
which enforces Eq.~\eqref{eq:ch2.causal-support-constraint} on \(\E\)
\cite{stratton1941,harrington2001}. Advanced kernels with
\(\mathcal{G}(\mathbf{r},\mathbf{r}',t-t')\neq\mathbf{0}\) for \(t-t'<0\) violate
Eq.~\eqref{eq:ch2.causal-operator-support} and are excluded from the equivalence
class treated here.

\begin{figure}[t]
\centering
\ChOneFigBlock{%
\begin{tikzpicture}[ch1fig, node distance=7mm and 11mm]
  \node[ch1box] (td) {$\mathcal{S}_t$\\time domain};
  \node[ch1box, right=14mm of td] (lt) {Laplace\\$s=-\jj\omega$};
  \node[ch1box, right=14mm of lt] (fd) {$\mathcal{S}_\omega$\\harmonic};
  \draw[ch1flow] (td) -- (lt) -- (fd);
\end{tikzpicture}%
}
\caption{Causality bridge: causal \(\mathcal{S}_t\) maps to harmonic
\(\mathcal{S}_\omega\) under Laplace/Fourier evaluation on the locked phasor axis
\(s=-\jj\omega\), \(\omega>0\).}
\label{fig:ch2.laplace-resolvent-bridge}
\end{figure}

Figure~\ref{fig:ch2.laplace-resolvent-bridge} organizes the proof template. Let
\(\mathbf{j}_{\mathrm{imp}}(\mathbf{r},t)\) be an impressed current with compact
support in \(\Omega_s\times[0,\infty)\) and let \(\mathbf{U}_0=\mathbf{0}\) for
simplicity. The Laplace transform in \(s\) with \(\Re\{s\}\) sufficiently large is
\begin{equation}
\label{eq:ch2.laplace-source-transform}
\widetilde{\mathbf{j}}(\mathbf{r},s)
=
\int_0^\infty \mathbf{j}_{\mathrm{imp}}(\mathbf{r},t)\,e^{-st}\,dt.
\end{equation}
On the phasor axis \(s=-\jj\omega\) with \(\omega>0\),
\begin{equation}
\label{eq:ch2.fourier-phasor-identification}
\widetilde{\mathbf{j}}(\mathbf{r},-\jj\omega)
=
\int_0^\infty \mathbf{j}_{\mathrm{imp}}(\mathbf{r},t)\,e^{\jj\omega t}\,dt,
\end{equation}
which is the Fourier transform restricted to causal support
\cite{stratton1941}. For \(\mathbf{j}_{\mathrm{imp}}(\mathbf{r},t)=\Re\{\underline{\mathbf{j}}(\mathbf{r})e^{-\jj\omega_0 t}\}\)
with fixed \(\omega_0>0\), Eq.~\eqref{eq:ch2.fourier-phasor-identification}
selects the phasor \(\underline{\mathbf{j}}\) under the locked convention
\(\exp(-\jj\omega t)\) \cite{stratton1941,harrington2001}.

Applying the Laplace transform to Eq.~\eqref{eq:ch2.time-ibvp-package} with fixed
\(\mathcal{B}_\Gamma\) and homogeneous initial data yields, formally,
\begin{equation}
\label{eq:ch2.laplace-maxwell-resolvent}
\widetilde{\mathbf{U}}(\mathbf{r},s)
=
\mathcal{R}(s)\,\widetilde{\mathbf{F}}(s),
\qquad
\mathcal{R}(s)=\mathcal{L}_s^{-1},
\end{equation}
where \(\mathcal{L}_s\) is the Laplace-domain Maxwell operator obtained from
\(\mathcal{A}\) by Eq.~\eqref{eq:ch2.frequency-substitution-rule} with
\(\partial_t\to s\) \cite{stratton1941}. On the axis \(s=-\jj\omega\),
\(\mathcal{L}_s=\mathcal{L}_\omega\) from Eq.~\eqref{eq:ch2.composite-operator-domain},
so
\begin{equation}
\label{eq:ch2.resolvent-green-equivalence}
\mathcal{R}(-\jj\omega)=\mathcal{L}_\omega^{-1}
\quad\text{on the outgoing/resolvent branch,}
\end{equation}
and the retarded convolution Eq.~\eqref{eq:ch2.retarded-convolution-form} is the
inverse Laplace transform of Eq.~\eqref{eq:ch2.laplace-maxwell-resolvent} when
\(\mathcal{G}\) is the causal Green operator constructed in Section~\ref{sec:ch23}
\cite{stratton1941,harrington2001}. Equation~\eqref{eq:ch2.resolvent-green-equivalence}
is the operator form of the preview Eq.~\eqref{eq:ch2.causality-preview}: harmonic
inversion is not an independent postulate but the \(s=-\jj\omega\) evaluation of a
causal resolvent.

The equivalence theorem requires explicit assumptions:
\begin{equation}
\label{eq:ch2.equivalence-assumptions}
\begin{aligned}
&\text{(i)}\ \mathcal{B}_\Gamma\ \text{and}\ \varepsilon,\mu\ \text{fixed in time},\\
&\text{(ii)}\ \mathbf{j}_{\mathrm{imp}}\ \text{causal with}\ \operatorname{supp}_t\subseteq[0,\infty),\\
&\text{(iii)}\ \text{outgoing/resolvent branch selected for exterior domains},\\
&\text{(iv)}\ \mathbf{U}_0=\mathbf{0}\ \text{or compatible with harmonic steady limit}.
\end{aligned}
\end{equation}
Under Eqs.~\eqref{eq:ch2.equivalence-assumptions},
\begin{equation}
\label{eq:ch2.harmonic-equivalence-theorem}
\mathcal{S}_\omega[\underline{\mathbf{j}}]
=
\eval{\widetilde{\mathbf{U}}(\mathbf{r},s)}_{s=-\jj\omega}
=
\eval{\mathcal{L}_\omega^{-1}\widetilde{\mathbf{F}}(s)}_{s=-\jj\omega},
\end{equation}
where \(\underline{\mathbf{j}}\) is the phasor of
\(\mathbf{j}_{\mathrm{imp}}(\mathbf{r},t)=\Re\{\underline{\mathbf{j}}(\mathbf{r})e^{-\jj\omega t}\}\)
and \(\widetilde{\mathbf{F}}\) is the Laplace forcing induced by
\(\widetilde{\mathbf{j}}\) \cite{stratton1941}. Equation
\eqref{eq:ch2.harmonic-equivalence-theorem} is the Chapter~\ref{ch:02} HOME
statement that frequency-domain solvers implement causal time-domain maps when
Eqs.~\eqref{eq:ch2.equivalence-assumptions} hold. Violating (i)--(iii) breaks
Eq.~\eqref{eq:ch2.harmonic-equivalence-theorem} even if
\(\mathcal{M}_\omega[\mathbf{u}_\omega]=\mathbf{0}\) in the interior
\cite{harrington2001}.

Radiation maps inherit the same equivalence. Because \(\Pi_{\mathrm{rad}}\) is linear and
commutes with harmonic extraction on the outgoing subspace when selection precedes
frequency evaluation,
\begin{equation}
\label{eq:ch2.radiation-time-frequency-equiv}
\mathcal{R}_\omega[\underline{\mathbf{j}}]
=
\Pi_{\mathrm{rad}}\,\mathcal{S}_\omega[\underline{\mathbf{j}}],
\end{equation}
and the time-domain radiation map \(\mathcal{R}_t=\Pi_{\mathrm{rad}}\circ\mathcal{S}_t\)
yields Eq.~\eqref{eq:ch2.radiation-time-frequency-equiv} under
Eq.~\eqref{eq:ch2.harmonic-equivalence-theorem} \cite{stratton1941,harrington2001}. Transport
and regime qualifications from Eq.~\eqref{eq:ch2.transport-qualified-radiation}
apply to Eq.~\eqref{eq:ch2.radiation-time-frequency-equiv} unchanged; causality does
not replace Section~\ref{sec:ch01.2.4} screening \cite{stratton1941}.

Analytic structure: causal retarded responses admit continuation to
\(\Im\{s\}>0\) (upper half-plane) under standard dispersive assumptions, while
\(\mathcal{L}_\omega\) is evaluated on \(s=-\jj\omega\) with \(\omega>0\)
\cite{stratton1941}. Poles of \(\mathcal{R}(s)\) in the lower half-plane correspond
to growing modes excluded by outgoing closure; discrete cavity poles on the axis
require loss or quotienting as in Eq.~\eqref{eq:ch2.coercivity-quotient}
\cite{harrington2001}. The steady-state limit after transients is
\begin{equation}
\label{eq:ch2.steady-state-harmonic-limit}
\lim_{t\to\infty}\mathbf{U}(t)
=
\Re\big\{\mathcal{S}_\omega[\underline{\mathbf{j}}]\,e^{-\jj\omega t}\big\}
\end{equation}
for single-frequency excitation when damping or outgoing radiation removes memory of
\(\mathbf{U}_0\) \cite{stratton1941}. Equation~\eqref{eq:ch2.steady-state-harmonic-limit}
justifies steady-state phasor measurements after switch-on transients decay.

\begin{figure}[t]
\centering
\ChOneFigBlock{%
\begin{tikzpicture}[ch1fig, node distance=8mm and 10mm]
  \node[ch1box] (ca) {Causal\\$\mathcal{G}(t-t')$};
  \node[ch1box, right=14mm of ca] (lt) {$\mathcal{R}(s)$};
  \node[ch1box, right=14mm of lt] (hw) {$\mathcal{L}_\omega^{-1}$};
  \node[ch1box, below=11mm of lt] (br) {Branch\\$\mathcal{B}_\Gamma$};
  \draw[ch1flow] (ca) -- (lt) -- (hw);
  \draw[ch1flow] (br) -- (lt);
\end{tikzpicture}%
}
\caption{Resolvent--Green equivalence: causal retarded \(\mathcal{G}\) Laplace-transforms
to resolvent \(\mathcal{R}(s)\); on \(s=-\jj\omega\), branch selection fixes
\(\mathcal{L}_\omega^{-1}\).}
\label{fig:ch2.causality-time-frequency-diagram}
\end{figure}

Figure~\ref{fig:ch2.causality-time-frequency-diagram} emphasizes that
Eq.~\eqref{eq:ch2.resolvent-green-equivalence} is branch-dependent. Frequency-domain
codes that invert \(\mathcal{L}_\omega\) without outgoing or passive closure implement
a different \(\mathcal{R}(s)\) than the causal retarded \(\mathcal{G}\) of
Eq.~\eqref{eq:ch2.retarded-convolution-form} \cite{stratton1941,harrington2001}.

Failure modes are structural. Advanced Green functions satisfy
\begin{equation}
\label{eq:ch2.advanced-kernel-violation}
\mathcal{G}_{\mathrm{adv}}(\mathbf{r},\mathbf{r}',t-t')\neq\mathbf{0}
\ \text{for some}\ t-t'<0,
\end{equation}
and produce harmonic symbols that cannot equal
Eq.~\eqref{eq:ch2.harmonic-equivalence-theorem} for any causal
\(\mathbf{j}_{\mathrm{imp}}\) \cite{stratton1941}. Broadband synthesis requires
integrating Eq.~\eqref{eq:ch2.harmonic-equivalence-theorem} over \(\omega\); a single
\(\mathcal{S}_\omega\) slice cannot represent arbitrary \(\mathbf{j}_{\mathrm{imp}}(t)\)
without superposition \cite{harrington2001}. Time-varying \(\mathcal{B}_\Gamma(t)\)
breaks assumption (i) and requires piecewise \(\mathcal{S}_t\) as in the boundary-switch
illustration of Subsection~\ref{sec:ch221}.

\paragraph{Worked illustration (switch-on and steady state).}
Let \(\mathbf{j}_{\mathrm{imp}}(\mathbf{r},t)=\mathbf{j}_0(\mathbf{r})H(t)\) with
\(\mathbf{j}_0\in\mathcal{J}_\Omega\) and outgoing exterior closure. The causal
\(\mathcal{S}_t\) produces a transient \(\mathbf{U}(t)\) approaching
Eq.~\eqref{eq:ch2.steady-state-harmonic-limit} for large \(t\). A harmonic solve at
\(\mathcal{S}_\omega\) with \(\Jimp=\mathbf{j}_0\) returns the phasor of that steady
limit, not the transient at finite \(t\); equivalence is between frequency symbols and
causal resolvents, not between arbitrary time slices and one phasor
\cite{stratton1941,harrington2001}.

\paragraph{Worked illustration (wrong branch).}
A finite-domain solve with perfectly absorbing layers approximates an outgoing
\(\mathcal{B}_\Gamma\); a closed-cavity solve at the same \(\omega\) selects a
different branch of \(\mathcal{L}_\omega^{-1}\). Both satisfy volume Maxwell rows but
only one matches Eq.~\eqref{eq:ch2.resolvent-green-equivalence} for the laboratory
open configuration. Causality cannot repair branch mismatch
\cite{stratton1941,harrington2001}.

\paragraph{Closing summary.}
Subsection~\ref{sec:ch223} establishes the Chapter~\ref{ch:02} HOME for causal operator
support Eqs.~\eqref{eq:ch2.causal-operator-support}--\eqref{eq:ch2.retarded-kernel-support},
Laplace--resolvent equivalence Eqs.~\eqref{eq:ch2.laplace-maxwell-resolvent}--\eqref{eq:ch2.resolvent-green-equivalence},
equivalence assumptions Eq.~\eqref{eq:ch2.equivalence-assumptions}, harmonic theorem
Eq.~\eqref{eq:ch2.harmonic-equivalence-theorem}, radiation inheritance
Eq.~\eqref{eq:ch2.radiation-time-frequency-equiv}, and steady-state limit
Eq.~\eqref{eq:ch2.steady-state-harmonic-limit}. Subsection~\ref{sec:ch224} applies
spectral splits to propagating and evanescent subspaces under the qualified regime of
Section~\ref{sec:ch01.2.4}.
\subsection{Propagating, evanescent, and non-propagating subspaces}
\label{sec:ch224}
% TOC tag: [HOME][USE SS1.2.4]

Subsections~\ref{sec:ch221}--\ref{sec:ch223} fixed evolution, harmonic operators,
and causality. Subsection~\ref{sec:ch224} gives the Chapter~\ref{ch:02} HOME for
spectral partitions into propagating, evanescent, and non-propagating subspaces in
operator form \cite{stratton1941,harrington2001}. Reactive versus radiative regime
criteria, \(\chi(r)\), and \(\mathcal{D}_{J,\mathrm{reg}}\) remain those of
Section~\ref{sec:ch01.2.4} and are invoked by reference through
Eqs.~\eqref{eq:ch1.radiative-regime-condition}--\eqref{eq:ch1.regime-qualified-decision};
flux integrals and Chu bounds are not re-derived \cite{stratton1941}. The outgoing
selector \(\Pi_{\mathrm{rad}}\) of Eq.~\eqref{eq:ch2.radiation-operator-def} is
refined spectrally, not replaced.

In a homogeneous region with symbol \(k_0\) from Eq.~\eqref{eq:ch2.helmholtz-symbol},
plane-wave trials \(\mathbf{u}_{\mathbf{k}}(\mathbf{r})=\mathbf{u}_0\,e^{-\jj\mathbf{k}\cdot\mathbf{r}}\)
satisfy the Helmholtz condition when
\begin{equation}
\label{eq:ch2.dispersion-relation}
|\mathbf{k}|=k_0,\qquad
k_z^2=k_0^2-k_x^2-k_y^2,
\end{equation}
with lateral wavenumber \(k_t=\sqrt{k_x^2+k_y^2}\) \cite{stratton1941}. Equation
\eqref{eq:ch2.dispersion-relation} splits spectral content by longitudinal behavior:
\begin{equation}
\label{eq:ch2.longitudinal-wavenumber-split}
k_z=
\begin{cases}
\pm\sqrt{k_0^2-k_t^2}, & k_t\le k_0\ \text{(propagating)},\\[4pt]
-\jj\,\alpha(k_t), & k_t>k_0\ \text{(evanescent)},\ \alpha>0,
\end{cases}
\end{equation}
where the evanescent root is chosen with decay in the \(+z\) half-space
\cite{stratton1941,harrington2001}. For \(k_t>k_0\) the decay rate is
\begin{equation}
\label{eq:ch2.evanescent-decay-rate}
\alpha(k_t)=\sqrt{k_t^2-k_0^2},
\end{equation}
so a pure evanescent trial on \(z>0\) scales as
\(e^{-\alpha(k_t)z}\) with penetration depth \(\delta_{\mathrm{ev}}=1/\alpha(k_t)\)
\cite{stratton1941}. Propagating modes carry real \(k_z=\pm\sqrt{k_0^2-k_t^2}\) and
oscillatory phase along the longitudinal axis; evanescent modes carry imaginary
\(k_z\) and store reactive energy in the near zone without far-zone power export
\cite{harrington2001}. At fixed \((k_x,k_y)\), the Helmholtz symbol
\(\mathcal{L}_\omega\) from Subsection~\ref{sec:ch222} acts on each longitudinal
branch as a scalar multiplier in the plane-wave basis, which is why spectral
partitions commute with harmonic extraction once outgoing closure is fixed on the
propagating sheet \cite{stratton1941}.

Transverse polarization completes the mode description. For a trial electric field
\(\widetilde{\E}(k_x,k_y)\,e^{-\jj(k_x x+k_y y+k_z z)}\), source-free admissibility
requires \(\mathbf{k}\cdot\widetilde{\E}=0\) in the homogeneous region
\cite{stratton1941,harrington2001}. Decompose \(\mathbf{k}=\mathbf{k}_t+k_z\hat{\mathbf{z}}\)
with \(\mathbf{k}_t=k_x\hat{\mathbf{x}}+k_y\hat{\mathbf{y}}\). Then
\begin{equation}
\label{eq:ch2.transverse-polarization-constraint}
k_x\widetilde{E}_x+k_y\widetilde{E}_y+k_z\widetilde{E}_z=0,
\end{equation}
which leaves two independent transverse amplitudes on each branch of
Eq.~\eqref{eq:ch2.longitudinal-wavenumber-split}. On the propagating sheet,
Eq.~\eqref{eq:ch2.transverse-polarization-constraint} supports traveling power
along \(\hat{\mathbf{z}}\) when outgoing \(k_z\) is selected; on the evanescent sheet
the same constraint forces predominantly tangential storage with longitudinal
\(E_z\) tied to \(\alpha(k_t)\) \cite{stratton1941}. Subsection~\ref{sec:ch224}
does not classify TE/TM guide nomenclature---that labeling is deferred to
waveguide sections---but Eq.~\eqref{eq:ch2.transverse-polarization-constraint}
is the local polarization gate paired with the dispersion gate of
Eq.~\eqref{eq:ch2.dispersion-relation} \cite{harrington2001}.

\begin{figure}[t]
\centering
\ChOneFigBlock{%
\begin{tikzpicture}[ch1fig, node distance=8mm and 10mm]
  \node[ch1box] (kt) {$k_t$};
  \node[ch1box, right=14mm of kt] (prop) {$k_t\le k_0$\\propagating};
  \node[ch1box, below=11mm of kt] (ev) {$k_t>k_0$\\evanescent};
  \draw[ch1flow] (kt) -- (prop);
  \draw[ch1flow] (kt) -- (ev);
\end{tikzpicture}%
}
\caption{Dispersion split Eq.~\eqref{eq:ch2.longitudinal-wavenumber-split}: lateral
wavenumber \(k_t\) selects propagating versus evanescent longitudinal behavior at
fixed \(k_0\).}
\label{fig:ch2.propagating-evanescent-split}
\end{figure}

Figure~\ref{fig:ch2.propagating-evanescent-split} is the local spectral gate used
throughout Section~\ref{sec:ch22}. Angular-spectrum representations in
Section~\ref{sec:ch29} expand fields in these modes; here the subspaces are defined
as operator domains \cite{stratton1941}.

Let \(\mathcal{F}_{\mathrm{prop}}\subset\mathcal{F}_\Omega\) be the closure of
finite superpositions of propagating plane waves with outgoing \(k_z\) selected by
exterior closure, let \(\mathcal{F}_{\mathrm{ev}}\subset\mathcal{F}_\Omega\) be the
closure of evanescent superpositions, and let
\(\mathcal{F}_{\mathrm{non}}\subset\mathcal{F}_\Omega\) collect modes that are
neither outgoing-propagating nor evanescent-decaying in the declared half-space---for
example cutoff guide modes and purely reactive standing content tied to terminals
\cite{stratton1941,harrington2001}. Orthogonal spectral projectors (formal) satisfy
\begin{equation}
\label{eq:ch2.spectral-projectors}
\Pi_{\mathrm{prop}}+\Pi_{\mathrm{ev}}+\Pi_{\mathrm{non}}=\mathrm{id}
\quad\text{on}\ \widehat{\mathcal{F}}_\Omega,
\end{equation}
with \(\widehat{\mathcal{F}}_\Omega\) the spectral completion of admissible states
\cite{stratton1941}. Idempotency and mutual annihilation read
\begin{equation}
\label{eq:ch2.spectral-projector-algebra}
\Pi_{\mathrm{prop}}^2=\Pi_{\mathrm{prop}},\quad
\Pi_{\mathrm{ev}}^2=\Pi_{\mathrm{ev}},\quad
\Pi_{\mathrm{non}}^2=\Pi_{\mathrm{non}},\quad
\Pi_{\mathrm{prop}}\Pi_{\mathrm{ev}}=\Pi_{\mathrm{ev}}\Pi_{\mathrm{prop}}=0,
\end{equation}
with analogous zero products for \(\Pi_{\mathrm{non}}\) \cite{stratton1941}. On
\(\mathcal{D}(\mathcal{L}_\omega)\) from Subsection~\ref{sec:ch222}, each projector
block is an invariant subspace of the scalar Helmholtz symbol in the plane-wave
basis,

\medskip
\noindent\textbf{Formal status.}
Equations~\eqref{eq:ch2.spectral-projectors}--\eqref{eq:ch2.spectral-projector-algebra} are
\emph{formal} on the spectral completion \(\widehat{\mathcal{F}}_\Omega\): idempotency and
orthogonality are justified in the rigged Hilbert framework of Chapter~\ref{ch:03}
Subsections~\ref{sec:ch313}--\ref{sec:ch352} \cite{stratton1941,kato1995}. Until that closure,
projector algebra is used only on admissible finite superpositions and on
\(\mathcal{R}_\omega[\Jimp]\), where outgoing selection is already enforced
\cite{harrington2001}.

\begin{equation}
\label{eq:ch2.helmholtz-subspace-commutation}
\mathcal{L}_\omega\big(\Pi_{\mathrm{prop}}\mathbf{u}\big)
=
\Pi_{\mathrm{prop}}\big(\mathcal{L}_\omega\mathbf{u}\big),
\end{equation}
and similarly for \(\Pi_{\mathrm{ev}}\) and \(\Pi_{\mathrm{non}}\) when boundary
closure does not mix sheets \cite{harrington2001}. Equation
\eqref{eq:ch2.helmholtz-subspace-commutation} is the operator reason harmonic
solvers can treat propagating and evanescent content separately before radiation
selection; Sommerfeld outgoing closure in Section~\ref{sec:ch24} acts only on
\(\Pi_{\mathrm{prop}}\) \cite{stratton1941}.

The decomposition
\begin{equation}
\label{eq:ch2.propagating-evanescent-decomposition}
\mathcal{S}_\omega[\Jimp]
=
\Pi_{\mathrm{prop}}\mathcal{S}_\omega[\Jimp]
+\Pi_{\mathrm{ev}}\mathcal{S}_\omega[\Jimp]
+\Pi_{\mathrm{non}}\mathcal{S}_\omega[\Jimp]
\end{equation}
holds when the impressed spectrum excites all three classes. Only
\(\Pi_{\mathrm{prop}}\mathcal{S}_\omega[\Jimp]\) is eligible for far-zone radiation
content after outgoing closure; evanescent and non-propagating components remain in
reactive storage exchange \cite{stratton1941,harrington2001}. Energy norms from
Eq.~\eqref{eq:ch2.field-energy-norm} decompose orthogonally in the spectral basis
when the impressed angular spectrum has sufficient decay in \((k_x,k_y)\)
\cite{stratton1941}.

The link to Eq.~\eqref{eq:ch2.radiation-operator-def} is
\begin{equation}
\label{eq:ch2.rad-projection-composition}
\Pi_{\mathrm{rad}}
=
\mathcal{O}_{\mathrm{out}}\circ\Pi_{\mathrm{prop}},
\end{equation}
where \(\mathcal{O}_{\mathrm{out}}\) enforces outgoing selection on propagating
content compatible with Eq.~\eqref{eq:ch1.radiation-trace-condition}
\cite{stratton1941}. Equation~\eqref{eq:ch2.rad-projection-composition} refines
\(\Pi_{\mathrm{rad}}\): radiative subspace \(\mathcal{F}_{\mathrm{rad}}\) is a
subset of \(\Pi_{\mathrm{prop}}\mathcal{F}_\Omega\), not of all of
\(\mathcal{F}_\Omega\) \cite{harrington2001}. Evanescent and non-propagating
components lie in the null space of far-zone radiation:
\begin{equation}
\label{eq:ch2.rad-null-on-evanescent}
\Pi_{\mathrm{rad}}\Pi_{\mathrm{ev}}=0,\qquad
\Pi_{\mathrm{rad}}\Pi_{\mathrm{non}}=0,
\end{equation}
while \(\Pi_{\mathrm{rad}}\Pi_{\mathrm{prop}}=\Pi_{\mathrm{rad}}\) after outgoing
closure \cite{stratton1941}. Equations~\eqref{eq:ch2.rad-null-on-evanescent} and
\eqref{eq:ch2.rad-projection-composition} restate, in projector form, the regime
logic of Section~\ref{sec:ch01.2.4}: reactive near fields may be large even when
\(\Pi_{\mathrm{rad}}\mathcal{S}_\omega[\Jimp]=\mathbf{0}\) \cite{harrington2001}.
Causality from Subsection~\ref{sec:ch223} ensures retarded support but does not,
by itself, promote evanescent content into \(\mathcal{F}_{\mathrm{rad}}\).

\medskip
\noindent\textbf{Effective definition (locked from §2.2.4 onward).}
For all \(r>r_{\min}(a,\lambda_0)\) and every transport, flux, near-to-far, and angular-spectrum
statement in Sections~\ref{sec:ch25}--\ref{sec:ch210},
\begin{equation*}
\Pi_{\mathrm{rad}}=\mathcal{O}_{\mathrm{out}}\circ\Pi_{\mathrm{prop}}
\quad\text{as in Eq.~\eqref{eq:ch2.rad-projection-composition}.}
\end{equation*}
Earlier uses of \(\Pi_{\mathrm{rad}}\) in Section~\ref{sec:ch20} are to be read through this
refinement \cite{stratton1941,harrington2001}.

Transport qualification from Section~\ref{sec:ch01.2.4} enters as an operator screen:
\begin{equation}
\label{eq:ch2.regime-operator-screen}
\mathcal{D}_{J,\mathrm{reg}}=1
\ \Longrightarrow\
\text{transport reports use}\ 
\Pi_{\mathrm{prop}}\mathcal{R}_\omega[\Jimp],
\ \text{not}\ \Pi_{\mathrm{ev}}\mathcal{S}_\omega[\Jimp],
\end{equation}
with \(\mathcal{D}_{J,\mathrm{reg}}\) from
Eq.~\eqref{eq:ch1.regime-qualified-decision} \cite{stratton1941}. Equation
\eqref{eq:ch2.regime-operator-screen} is the operator restatement of
Eq.~\eqref{eq:ch2.transport-qualified-radiation}; regime indices \(\chi(r)\) and
\(\chi_\Omega\) diagnose when Eq.~\eqref{eq:ch2.regime-operator-screen} fails on a
measurement shell \cite{stratton1941,harrington2001}.

Waveguide and cutoff structures illustrate non-propagating subspaces. For a uniform
guide with lateral eigenvalue \(\kappa_n\), mode \(n\) propagates only when
\begin{equation}
\label{eq:ch2.cutoff-condition}
k_0^2>\kappa_n^2
\quad\Longleftrightarrow\quad
\omega>\omega_{c,n}=\frac{c_0\kappa_n}{\sqrt{\varepsilon_r\mu_r}},
\end{equation}
with \(\omega_{c,n}\) the cutoff angular frequency \cite{stratton1941}. Below
cutoff, the mode lies in \(\mathcal{F}_{\mathrm{ev}}\) or
\(\mathcal{F}_{\mathrm{non}}\) depending on boundary termination; applying far-zone
transport formulas to such fields violates Eq.~\eqref{eq:ch2.regime-operator-screen}
as in the TE\(_{10}\) illustration of Section~\ref{sec:ch01.1.2}
\cite{harrington2001}. In a rectangular guide of width \(a\), lateral eigenvalues
\(\kappa_{m,n}=\sqrt{(m\pi/a)^2+(n\pi/b)^2}\) replace \(k_t\) in
Eq.~\eqref{eq:ch2.cutoff-condition}; each mode index \((m,n)\) switches from
\(\mathcal{F}_{\mathrm{non}}\) to \(\mathcal{F}_{\mathrm{prop}}\) only when
\(\omega\) exceeds its cutoff \cite{stratton1941}. Near cutoff the propagating
group velocity vanishes and stored reactive energy dominates even though the mode
has crossed into \(\Pi_{\mathrm{prop}}\); Eq.~\eqref{eq:ch1.radiative-regime-condition}
remains the far-zone transport gate, not Eq.~\eqref{eq:ch2.cutoff-condition} alone
\cite{harrington2001}.

Angular-spectrum preview: expand \(\E(\mathbf{r})\) on \(z=\mathrm{const}\) as
\begin{equation}
\label{eq:ch2.angular-spectrum-preview}
\E(\mathbf{r})
=
\int\!\!\int \widetilde{\E}(k_x,k_y)\,
e^{-\jj(k_x x+k_y y+k_z(k_t) z)}\,dk_x\,dk_y,
\end{equation}
with \(k_z(k_t)\) from Eq.~\eqref{eq:ch2.longitudinal-wavenumber-split}
\cite{stratton1941}. Equation~\eqref{eq:ch2.angular-spectrum-preview} is the
Chapter~\ref{ch:02} spectral template for Section~\ref{sec:ch29}; full radiation
operator action on angular spectra is HOME there. Subsection~\ref{sec:ch224} only
defines the propagating/evanescent partition induced by
Eq.~\eqref{eq:ch2.dispersion-relation} \cite{harrington2001}.

\begin{figure}[t]
\centering
\ChOneFigBlock{%
\begin{tikzpicture}[ch1fig, node distance=7mm and 9mm]
  \node[ch1box] (s) {$\mathcal{S}_\omega[\Jimp]$};
  \node[ch1box, right=13mm of s] (pp) {$\Pi_{\mathrm{prop}}$};
  \node[ch1box, right=13mm of pp] (pr) {$\Pi_{\mathrm{rad}}$};
  \node[ch1box, below=11mm of s] (pe) {$\Pi_{\mathrm{ev}}$};
  \node[ch1box, below=11mm of pe] (pn) {$\Pi_{\mathrm{non}}$};
  \draw[ch1flow] (s) -- (pp) -- (pr);
  \draw[ch1flow] (s) -- (pe);
  \draw[ch1flow] (s) -- (pn);
\end{tikzpicture}%
}
\caption{Spectral gate: only \(\Pi_{\mathrm{prop}}\) content passes to
\(\Pi_{\mathrm{rad}}\); evanescent and non-propagating components remain reactive.}
\label{fig:ch2.spectral-radiation-gate}
\end{figure}

Figure~\ref{fig:ch2.spectral-radiation-gate} connects Subsection~\ref{sec:ch224} to
Figure~\ref{fig:ch2.source-radiation-pipeline}. Causality from
Subsection~\ref{sec:ch223} does not automatically place evanescent energy in
\(\mathcal{F}_{\mathrm{rad}}\); spectral and regime screens are additional
\cite{stratton1941}.

Near-zone reactive fields decompose into large evanescent and non-propagating
fractions even when \(\mathcal{S}_\omega[\Jimp]\) is exact. As \(r\to\infty\),
\(\Pi_{\mathrm{prop}}\) content dominates \(P_{\mathrm{rad}}\) when
Eq.~\eqref{eq:ch1.radiative-regime-condition} holds \cite{stratton1941}. Operator
norms on subspaces satisfy
\begin{equation}
\label{eq:ch2.subspace-energy-inequality}
\|\Pi_{\mathrm{ev}}\mathcal{S}_\omega[\Jimp]\|_{\mathcal{U}}
\gg
\|\Pi_{\mathrm{prop}}\mathcal{S}_\omega[\Jimp]\|_{\mathcal{U}}
\quad\text{on near shells,}
\end{equation}
while the reverse inequality governs far-zone radiative observables
\cite{harrington2001,stratton1941}. Subsection~\ref{sec:ch224} does not compute
shell integrals; it supplies the spectral vocabulary for
Section~\ref{sec:ch252} and flux operators in Section~\ref{sec:ch25}.

\paragraph{Worked illustration (near-field scan).}
A probe at \(kr=1\) measures strong tangential fields dominated by
\(\Pi_{\mathrm{ev}}\mathcal{S}_\omega[\Jimp]\). Declaring a transport channel from
that scan without \(\mathcal{D}_{J,\mathrm{reg}}=1\) violates
Eq.~\eqref{eq:ch2.regime-operator-screen} even if
\(\mathcal{L}_\omega[\mathbf{u}_\omega]=\mathbf{0}\) \cite{stratton1941}. Moving
the evaluation surface until Eq.~\eqref{eq:ch1.radiative-regime-condition} holds
increases the \(\Pi_{\mathrm{prop}}\) fraction without changing impressed data
\cite{harrington2001}.

\paragraph{Worked illustration (cutoff waveguide).}
Below Eq.~\eqref{eq:ch2.cutoff-condition}, guide modes are evanescent in the
propagation direction. \(\Pi_{\mathrm{rad}}\) annihilates such content in the
exterior, yet \(\|\mathcal{S}_\omega[\Jimp]\|\) remains large in the feed region.
Confusing stored guide energy with radiative export is a regime error, not a gauge
artifact \cite{stratton1941,harrington2001}.

\paragraph{Worked illustration (total internal reflection).}
At a planar dielectric interface, an incident propagating mode with
\(k_{t,\mathrm{inc}}<k_0\) can generate a transmitted evanescent sheet when the
transmitted region has \(k_{0,\mathrm{trans}}<k_{t,\mathrm{inc}}\)
\cite{stratton1941}. The transmitted longitudinal wavenumber is then
\(-\jj\alpha(k_{t,\mathrm{inc}})\) from Eq.~\eqref{eq:ch2.evanescent-decay-rate},
so \(\Pi_{\mathrm{ev}}\) captures the tunneling field while \(\Pi_{\mathrm{prop}}\)
remains on the incident side. Declaring a radiative channel from the evanescent
sheet alone violates Eq.~\eqref{eq:ch2.rad-null-on-evanescent} even though
\(|\Pi_{\mathrm{ev}}\mathcal{S}_\omega[\Jimp]|\) is measurable at the interface
\cite{harrington2001}. Only after recombination with outgoing propagating content
in the exterior half-space does Eq.~\eqref{eq:ch2.regime-operator-screen} permit
transport reporting \cite{stratton1941}.

\paragraph{Closing summary.}
Subsection~\ref{sec:ch224} establishes the Chapter~\ref{ch:02} HOME for dispersion
split Eqs.~\eqref{eq:ch2.dispersion-relation}--\eqref{eq:ch2.longitudinal-wavenumber-split},
evanescent decay Eq.~\eqref{eq:ch2.evanescent-decay-rate}, transverse gate
Eq.~\eqref{eq:ch2.transverse-polarization-constraint}, spectral projectors
Eqs.~\eqref{eq:ch2.spectral-projectors}--\eqref{eq:ch2.spectral-projector-algebra},
Helmholtz commutation Eq.~\eqref{eq:ch2.helmholtz-subspace-commutation},
decomposition Eq.~\eqref{eq:ch2.propagating-evanescent-decomposition},
radiation composition Eqs.~\eqref{eq:ch2.rad-projection-composition}--\eqref{eq:ch2.rad-null-on-evanescent},
regime screen Eq.~\eqref{eq:ch2.regime-operator-screen}, cutoff condition
Eq.~\eqref{eq:ch2.cutoff-condition}, angular-spectrum preview
Eq.~\eqref{eq:ch2.angular-spectrum-preview}, and subspace energy inequality
Eq.~\eqref{eq:ch2.subspace-energy-inequality}. Section~\ref{sec:ch01.2.4} regime logic
is embedded by reference.

With Subsections~\ref{sec:ch221}--\ref{sec:ch224} complete, Section~\ref{sec:ch22}
has connected time-domain evolution, harmonic Maxwell and Helmholtz operators,
causal equivalence, and propagating--evanescent spectral gates on the function spaces
of Section~\ref{sec:ch21}. Green-function inversion in Section~\ref{sec:ch23}
realizes \(\mathcal{L}_\omega^{-1}\) and \(\mathcal{G}_{E,\omega}\) with retarded
causality; Sommerfeld selection in Section~\ref{sec:ch24} fixes outgoing branches for
\(\Pi_{\mathrm{prop}}\) and \(\Pi_{\mathrm{rad}}\); radiation integrals in
Section~\ref{sec:ch25} meter flux on the propagating subspace qualified by
Eq.~\eqref{eq:ch2.regime-operator-screen} \cite{stratton1941,harrington2001}.
\section{Green Functions and Electromagnetic Inversion}
\label{sec:ch23}

Section~\ref{sec:ch22} connected harmonic Maxwell operators to causal time-domain
evolution and to propagating--evanescent spectral partitions on the function spaces
of Section~\ref{sec:ch21}. Those results fixed the symbols \(\mathcal{L}_\omega\),
\(\mathcal{S}_\omega\), and the retarded equivalence previewed in
Eqs.~\eqref{eq:ch2.retarded-convolution-form} and~\eqref{eq:ch2.resolvent-green-equivalence},
but they did not construct the kernels that implement
\(\mathcal{L}_\omega^{-1}\) on admissible domains. Section~\ref{sec:ch23} supplies
that construction layer: scalar, vector, and dyadic Green operators for Maxwell
inversion, distributional regularization of singular source--field couplings,
retarded causal propagation, and solvability criteria on single- and multi-region
geometries \cite{stratton1941,harrington2001,jackson1999}. The phasor convention
\(\exp(-\jj\omega t)\) and SI units remain those locked in Chapter~\ref{ch:01};
Section~\ref{sec:ch23} does not reopen local Maxwell derivations or angular-momentum
balance.

The resolvent viewpoint from Subsections~\ref{sec:ch202} and~\ref{sec:ch212} treats
\(\mathcal{S}_\omega\) as a physical branch of \(\mathcal{L}_\omega^{-1}\) selected
by boundary class \(\mathcal{B}_\Gamma\) and outgoing closure
(Eq.~\eqref{eq:ch2.resolvent-operator}). Green functions are the integral
realizations of that resolvent: a scalar Helmholtz kernel inverts the Laplacian or
Helmholtz symbol on scalar potentials; vector and dyadic kernels invert coupled
curl--curl systems while respecting divergence constraints and tangential trace rows
on \(\Gamma=\partial\Omega\) \cite{stratton1941}. Equation
\eqref{eq:ch2.solution-map-resolvent} remains the operator statement; kernel formulas
are the concrete maps that numerical codes, radiation integrals, and reciprocity
identities evaluate once domains and branch selection are fixed
\cite{harrington2001,jackson1999}. \(\mathcal{S}_\omega\) is therefore not an
abstract matrix inverse but a retarded or outgoing-truncated resolvent on
\(\mathcal{D}(\mathcal{L}_\omega)\subset\mathcal{U}_\omega\).

Subsection~\ref{sec:ch231} owns the Chapter~\ref{ch:02} HOME for scalar, vector, and
dyadic Green function definitions in harmonic Maxwell theory. The scalar kernel
\(G(\mathbf{r},\mathbf{r}')\) satisfies \((\nabla^2+k_0^2)G=-\delta(\mathbf{r}-\mathbf{r}')\)
with boundary conditions inherited from \(\mathcal{B}_\Gamma\); the electric dyadic
\(\mathcal{G}_{E}(\mathbf{r},\mathbf{r}')\) inverts the electric curl--curl operator
of Subsection~\ref{sec:ch222} on \(\mathcal{D}(\mathcal{W}_{E,\omega})\) without
re-displaying Eq.~\eqref{eq:ch2.electric-wave-operator} \cite{stratton1941}. Vector
potentials and Lorenz-gauge reductions appear only as intermediate tools when they
shorten a derivation; gauge freedom and helicity limits fixed in Chapter~\ref{ch:01}
are not re-proved \cite{harrington2001}. Magnetic dyadics and Hertz potentials are
introduced when they reduce duplication between electric and magnetic source
couplings, always on the trace spaces of Subsection~\ref{sec:ch211}.

\begin{figure}[t]
\centering
\ChOneFigBlock{%
\begin{tikzpicture}[ch1fig, node distance=7mm and 10mm]
  \node[ch1box] (sc) {Scalar\\Helmholtz $G$};
  \node[ch1box, right=13mm of sc] (vec) {Vector\\$\mathbf{G}$};
  \node[ch1box, right=13mm of vec] (dy) {Dyadic\\$\mathcal{G}_E$};
  \node[ch1box, below=12mm of vec] (reg) {Regularization\\$\mathcal{D}'$};
  \node[ch1box, below=12mm of dy] (ret) {Retarded\\$\mathcal{G}(t-t')$};
  \node[ch1box, right=13mm of ret] (sol) {Solvability\\multi-region};
  \draw[ch1flow] (sc) -- (vec) -- (dy);
  \draw[ch1flow] (dy) -- (reg);
  \draw[ch1flow] (reg) -- (ret) -- (sol);
\end{tikzpicture}%
}
\caption{Section~\ref{sec:ch23} development path: scalar Helmholtz kernels support
vector and dyadic Maxwell inverses; distributional regularization and retarded
causality precede solvability on composite domains
(Subsections~\ref{sec:ch231}--\ref{sec:ch234}).}
\label{fig:ch2.green-function-roadmap}
\end{figure}

Figure~\ref{fig:ch2.green-function-roadmap} orders the four subsections that follow.
Subsection~\ref{sec:ch231} defines the kernel hierarchy and the operator identities
that connect \(\mathcal{G}_{E,\omega}\) to \(\mathcal{S}_\omega\) on impressed
currents \(\Jimp\in\mathcal{J}_\Omega\). Subsection~\ref{sec:ch232} treats
singularities at \(\mathbf{r}=\mathbf{r}'\), distributional extensions beyond
classical \(L^{2}\) fields, and the regularized self-coupling that method-of-moments
and finite-element codes implement implicitly \cite{jackson1999,harrington2001}.
Subsection~\ref{sec:ch233} upgrades the retarded preview of
Eq.~\eqref{eq:ch2.retarded-convolution-form} to a causal propagation theorem aligned
with Eq.~\eqref{eq:ch2.harmonic-equivalence-theorem} from Subsection~\ref{sec:ch223}.
Subsection~\ref{sec:ch234} addresses existence, uniqueness, and constructive inversion
on piecewise-smooth domains with interface traces from Subsection~\ref{sec:ch211} and
material Helmholtz symbols from Subsection~\ref{sec:ch222}
\cite{stratton1941}.

Singularities are structural, not numerical artifacts. The free-space dyadic Green
function blows up as \(|\mathbf{r}-\mathbf{r}'|^{-3}\) near coincidence; impressed
currents in \(\mathcal{J}_\Omega\) are often surface-supported or filamentary, so
field evaluation requires distributional interpretation of the convolution
\(\int \mathcal{G}_{E}(\mathbf{r},\mathbf{r}')\cdot\Jimp(\mathbf{r}')\,dV'\)
\cite{stratton1941,jackson1999}. Subsection~\ref{sec:ch232} supplies the Chapter~\ref{ch:02}
HOME for principal-value and finite-part prescriptions, for extraction of singular
static terms from dynamic kernels, and for the self-term corrections that distinguish
raw Green evaluation from Galerkin-tested impedance entries \cite{harrington2001}.
Section~\ref{sec:ch23} does not substitute mesh-refinement heuristics for those
distributional definitions: a code that omits self-term regularization may still
return a finite number while violating the operator identity that
Eq.~\eqref{eq:ch2.solution-map-resolvent} encodes.

Retardation ties harmonic kernels to the causal time-domain map
\(\mathcal{S}_t\) of Subsection~\ref{sec:ch221}. The retarded scalar kernel
\(G(\mathbf{r},\mathbf{r}',t-t')\) vanishes for \(t-t'<0\), matching
Eq.~\eqref{eq:ch2.retarded-kernel-support}; its Fourier transform on the locked
phasor axis yields the outgoing harmonic \(\mathcal{G}_{E,\omega}\) when branch
selection is consistent with Eq.~\eqref{eq:ch2.resolvent-green-equivalence}
\cite{stratton1941}. Advanced or acausal kernels produce harmonic symbols that cannot
equal \(\mathcal{S}_\omega\) for any physical impressed current
(Eq.~\eqref{eq:ch2.advanced-kernel-violation}). Subsection~\ref{sec:ch233} makes
that inheritance explicit for dyadic Maxwell systems and records the support
constraints that time-domain finite-difference schemes approximate when they march
\(\mathbf{U}(t)\) forward from Subsection~\ref{sec:ch221} \cite{harrington2001}.

Solvability is branch-dependent. On bounded domains with passive impedance or PEC
closure, \(\mathcal{L}_\omega\) may possess discrete cavity spectra requiring
quotienting as in Eq.~\eqref{eq:ch2.coercivity-quotient}; on exterior domains,
outgoing Green functions exist only after Sommerfeld selection in
Section~\ref{sec:ch24} \cite{stratton1941}. Subsection~\ref{sec:ch234} states
constructive inversion theorems for composite \(\Omega=\cup_j\Omega_j\) with
interface jumps from Eq.~\eqref{eq:ch1.interface-conditions}, importing trace classes
without re-proving jump laws from Chapter~\ref{ch:01}. Multi-region dyadic assembly
uses continuity of \(\hat{\mathbf{n}}\times\E\) and \(\hat{\mathbf{n}}\times\H\)
across interfaces while allowing piecewise \(k_j\) in each subdomain
\cite{harrington2001}. When solvability fails---for example, resonant cavities without
loss, or exterior problems with incoming rather than outgoing radiation---the failure
mode is detected at the resolvent level, not patched by redefining \(\Jimp\).

Several constructions deliberately remain outside Section~\ref{sec:ch23}. Sommerfeld
outgoing integrals and exterior uniqueness proofs belong to Section~\ref{sec:ch24};
far-zone radiation integrals and flux meters to Section~\ref{sec:ch25}; vector
spherical harmonic separation to Chapter~\ref{ch:03}. Section~\ref{sec:ch23} builds
kernels on the domains declared in Section~\ref{sec:ch21} and on the harmonic symbols
fixed in Section~\ref{sec:ch22}; it does not classify propagating versus evanescent
far-zone content beyond citing the spectral projectors of Subsection~\ref{sec:ch224}.
Green convolution returns the full admissible field
\(\mathcal{S}_\omega[\Jimp]=\Pi_{\mathrm{prop}}\mathcal{S}_\omega[\Jimp]
+\Pi_{\mathrm{ev}}\mathcal{S}_\omega[\Jimp]+\Pi_{\mathrm{non}}\mathcal{S}_\omega[\Jimp]\);
only the propagating fraction is eligible for \(\Pi_{\mathrm{rad}}\) under
Eq.~\eqref{eq:ch2.regime-operator-screen} \cite{stratton1941,harrington2001}. Kernel
construction therefore precedes radiation screening but must not be confused with it.

The compactness split of Eq.~\eqref{eq:ch2.green-compact-split} from
Subsection~\ref{sec:ch213} reappears in kernel form: singular static parts plus
smooth dynamic corrections define Fredholm structures used later in
Section~\ref{sec:ch263} \cite{stratton1941}. Radiation operators inherit that split
when observation maps are composed on the right of \(\mathcal{S}_\omega\); Green
function theory here supplies the integral operator whose compact remainder encodes
scattered fields relative to the primary source coupling \cite{harrington2001}.

Readers approaching from computational electromagnetics should map each subsection to
a familiar object. Subsection~\ref{sec:ch231} is the continuous target behind
free-space and layered-media dyads tabulated in references \cite{jackson1999}.
Subsection~\ref{sec:ch232} is the analytical basis for self-term and near-singular
integration in method-of-moments matrices. Subsection~\ref{sec:ch233} is the
continuum limit that FDTD retarded updates approximate. Subsection~\ref{sec:ch234}
is the well-posedness check that finite-element and boundary-element solvers assume
when they invert a stiffness or impedance system on a declared mesh
\cite{harrington2001}. None of those discretizations is derived here; Section~\ref{sec:ch23}
fixes the continuous operators they must converge toward.

Material dispersion, anisotropy, and moving media can modify kernel symbols beyond
the piecewise-scalar \(\varepsilon(\mathbf{r})\), \(\mu(\mathbf{r})\) class treated
in Subsection~\ref{sec:ch222}; within the locked harmonic band, kernels are evaluated
at fixed \(\omega>0\) with \(k(\mathbf{r})\) from
Eq.~\eqref{eq:ch2.helmholtz-symbol}. Loss enters through complex \(k\) or through
passive impedance boundaries in Eq.~\eqref{eq:ch2.boundary-impedance-operator},
shifting poles of \(\mathcal{L}_\omega^{-1}\) without changing the convolution
template \cite{stratton1941}. Perfectly conducting inclusions are sharp limits of
impedance closures; Subsection~\ref{sec:ch234} treats them through homogeneous
tangential trace conditions on \(\Gamma_{\mathrm{PEC}}\) rather than through
modified volume Green functions alone \cite{harrington2001}.

Charge neutrality and continuity constraints on \(\Jimp\) from
Eqs.~\eqref{eq:ch1.charge-continuity-harmonic} and~\eqref{eq:ch2.charge-neutral-source}
remain in force during inversion: a Green formula that ignores \(\divg\Jimp\)
when assembling \(\nabla\cdot\E\) produces source images that are not in
\(\mathcal{J}_\Omega\) \cite{stratton1941}. Subsection~\ref{sec:ch231} states dyadic
identities in divergence-compatible form; Subsection~\ref{sec:ch234} verifies that
interface charge accumulation is consistent with impressed data on each subdomain.

The subsection sequence begins with kernel definitions that realize
Eq.~\eqref{eq:ch2.resolvent-operator}, continues with distributional regularization
and retarded causality, and closes with solvability on composite domains before
Sommerfeld radiation conditions in Section~\ref{sec:ch24} select the exterior branch
of \(\mathcal{G}_{E,\omega}\) \cite{stratton1941,harrington2001,jackson1999}.


\subsection{Scalar, vector, and dyadic Green functions}
\label{sec:ch231}
% TOC tag: [HOME]

Section~\ref{sec:ch23} declared that Green kernels realize the resolvent
Eq.~\eqref{eq:ch2.resolvent-operator}. Subsection~\ref{sec:ch231} gives the
Chapter~\ref{ch:02} HOME definitions for scalar, vector, and dyadic Green functions
that implement \(\mathcal{G}_{E,\omega}\) previewed in
Eq.~\eqref{eq:ch2.helmholtz-solution-operator} and connect to the radiation integral
preview Eq.~\eqref{eq:ch2.radiation-integral-preview}
\cite{stratton1941,harrington2001,jackson1999}. The electric-wave operator
Eq.~\eqref{eq:ch2.electric-wave-operator} and material-weighted symbols
Eqs.~\eqref{eq:ch2.material-helmholtz-electric}--\eqref{eq:ch2.material-helmholtz-magnetic}
from Subsection~\ref{sec:ch222} are presupposed; their curl--curl forms are not
re-derived. Gauge structure from Chapter~\ref{ch:01} enters only through the Lorenz
reduction below; gauge freedom is not re-opened \cite{stratton1941}. Singular
coincidence limits, distributional convolutions, and self-term regularization belong
to Subsection~\ref{sec:ch232}; outgoing branch selection at infinity belongs to
Section~\ref{sec:ch24}.

\paragraph{Assumptions.}
Unless noted otherwise, \(\omega>0\) is fixed, \(\Jimp\in\mathcal{J}_\Omega\) from
Eq.~\eqref{eq:ch2.source-space-refined}, charge continuity
Eq.~\eqref{eq:ch1.charge-continuity-harmonic} holds, and \(\mathcal{B}_\Gamma\) is
declared on \(\Gamma=\partial\Omega\). In uniform vacuum regions the wavenumber is
\(k_0\) from Eq.~\eqref{eq:ch2.helmholtz-symbol}. Kernels are stated first in
operator form on \(\mathcal{D}(\mathcal{W}_{E,\omega})\subset\HCurl(\Omega)\) from
Eq.~\eqref{eq:ch2.electric-helmholtz-domain}; piecewise \(k(\mathbf{r})\) extensions
use the same template with patchwise coefficients
\cite{stratton1941,harrington2001}.

\paragraph{Scalar Helmholtz kernel.}
The scalar Green function \(G(\mathbf{r},\mathbf{r}')\) inverts the Helmholtz operator
\((\nabla^2+k_0^2)\) on scalar potentials and potentials-of-vector fields. It satisfies
\begin{equation}
\label{eq:ch2.scalar-helmholtz-green-pde}
\big(\nabla^{2}+k_0^{2}\big)\,G(\mathbf{r},\mathbf{r}')
=
-\delta(\mathbf{r}-\mathbf{r}'),
\end{equation}
with \(\delta\) the three-dimensional Dirac distribution and boundary conditions on
\(\Gamma\) inherited from \(\mathcal{B}_\Gamma\) \cite{stratton1941}. Equation
\eqref{eq:ch2.scalar-helmholtz-green-pde} is the local symbol of spectral inversion:
for fixed \(\mathbf{r}'\), \(G(\cdot,\mathbf{r}')\) is the response to a unit point
source at \(\mathbf{r}'\) under the same closure as \(\mathcal{L}_\omega\)
\cite{jackson1999}. The associated solution operator is
\begin{equation}
\label{eq:ch2.scalar-green-solution-operator}
\phi(\mathbf{r})
=
\int_{\Omega} G(\mathbf{r},\mathbf{r}')\,q(\mathbf{r}')\,d^{3}r',
\qquad
\big(\nabla^{2}+k_0^{2}\big)\phi=q,
\end{equation}
for scalar sources \(q\) in the dual class paired with \(\phi\)
\cite{stratton1941}. On exterior domains, \(G\) must satisfy outgoing decay compatible
with Section~\ref{sec:ch24}; on bounded domains, Dirichlet, Neumann, or impedance
data on \(\Gamma\) replace the free-space radiation template
\cite{harrington2001}.

In unbounded vacuum with outgoing selection at infinity,
\begin{equation}
\label{eq:ch2.free-space-scalar-green}
G_0(\mathbf{r},\mathbf{r}')
=
\frac{e^{-\jj k_0 R}}{4\pi R},
\qquad
R=|\mathbf{r}-\mathbf{r}'|,
\end{equation}
with \(R>0\) and the branch of \(k_0\) fixed by the locked phasor axis
\cite{stratton1941,jackson1999}. Equation~\eqref{eq:ch2.free-space-scalar-green}
satisfies Eq.~\eqref{eq:ch2.scalar-helmholtz-green-pde} for
\(\mathbf{r}\neq\mathbf{r}'\); the distributional definition at coincidence is
completed in Subsection~\ref{sec:ch232}. Physically,
Eq.~\eqref{eq:ch2.free-space-scalar-green} is the spherical outgoing wave produced
by a unit monopole source at \(\mathbf{r}'\); its phase gradient defines local
propagation direction on the propagating sheet of Subsection~\ref{sec:ch224}
\cite{stratton1941}.

\paragraph{Vector potential and vector kernels.}
Lorenz-gauge vector potentials \(\mathbf{A}\) reduce vector Maxwell inversion to a
vector Helmholtz system driven by \(\Jimp\). Under the gauge and continuity constraints
fixed in Chapter~\ref{ch:01}, each Cartesian component of \(\mathbf{A}\) satisfies
Eq.~\eqref{eq:ch2.scalar-helmholtz-green-pde} with \(k_0\) and source \(-\mu_0\Jimp\)
\cite{stratton1941}. The vector field recovered from \(\mathbf{A}\) is
\begin{equation}
\label{eq:ch2.lorenz-field-from-potential}
\E(\mathbf{r})
=
\jj\omega\mu_0\,\mathbf{A}(\mathbf{r})
+\frac{1}{\jj\omega\varepsilon_0}\,\grad\,\divg\mathbf{A}(\mathbf{r}),
\end{equation}
with \(\H\) obtained from \(\curl\E\) when the first-order system is required
\cite{harrington2001,jackson1999}. Equation~\eqref{eq:ch2.lorenz-field-from-potential}
is a vector reduction tool, not a substitute for the dyadic definition below when
\(\mathcal{B}_\Gamma\) couples \(\E\) and \(\H\) at boundaries: impedance and PEC
closures are enforced on \(\mathcal{L}_\omega\), not on \(\mathbf{A}\) alone
\cite{stratton1941}. A vector Green function \(\mathbf{G}^{A}(\mathbf{r},\mathbf{r}')\)
may be assembled from three scalar kernels; the electric dyadic is preferred when
\(\Jimp\) is the primary excitation because it maps \(\Jimp\) to \(\E\) directly
without intermediate gauge potentials \cite{harrington2001}.

\paragraph{Electric dyadic Green function.}
Let \(\overline{\mathbf{G}}_{E}(\mathbf{r},\mathbf{r}')\) denote a rank-two tensor
(dyad) acting on vector sources by dot product from the left. The electric dyadic
Green function inverts the electric-wave operator
Eq.~\eqref{eq:ch2.electric-wave-operator} in the sense
\begin{equation}
\label{eq:ch2.dyadic-green-defining-operator}
\mathcal{W}_{E,\omega}\,
\overline{\mathbf{G}}_{E}(\cdot,\mathbf{r}')
=
\overline{\mathbf{I}}\,\delta(\mathbf{r}-\mathbf{r}'),
\end{equation}
where \(\overline{\mathbf{I}}\) is the identity dyad and the operator
\(\mathcal{W}_{E,\omega}\) acts on the \(\mathbf{r}\)-dependence of
\(\overline{\mathbf{G}}_{E}\) \cite{stratton1941,harrington2001}. Equation
\eqref{eq:ch2.dyadic-green-defining-operator} is the dyadic analogue of
Eq.~\eqref{eq:ch2.scalar-helmholtz-green-pde}: each column of
\(\overline{\mathbf{G}}_{E}\) is the electric field produced by an impulsive vector
source aligned with a Cartesian unit vector at \(\mathbf{r}'\)
\cite{jackson1999}. The electric field due to impressed current is
\begin{equation}
\label{eq:ch2.electric-dyadic-convolution}
\E(\mathbf{r})
=
\jj\omega\mu_0
\int_{\Omega}
\overline{\mathbf{G}}_{E}(\mathbf{r},\mathbf{r}')\cdot\Jimp(\mathbf{r}')\,d^{3}r'
\;+\;\E_{\mathrm{hom}}(\mathbf{r}),
\end{equation}
where \(\E_{\mathrm{hom}}\) is the homogeneous solution fixed by \(\mathcal{B}_\Gamma\)
\cite{stratton1941}. Equation~\eqref{eq:ch2.electric-dyadic-convolution} is the
Chapter~\ref{ch:02} HOME realization of the preview
Eq.~\eqref{eq:ch2.radiation-integral-preview} with outgoing dyad
\(\boldsymbol{\Gamma}=\overline{\mathbf{G}}_{E}\) on the physical branch
\cite{harrington2001}. The prefactor \(\jj\omega\mu_0\) matches the locked phasor
convention and the forcing map Eq.~\eqref{eq:ch2.helmholtz-forcing-map}.

The integral operator \(\mathcal{G}_{E,\omega}\) of
Eq.~\eqref{eq:ch2.helmholtz-solution-operator} is realized by
\begin{equation}
\label{eq:ch2.dyadic-green-integral-operator}
\big(\mathcal{G}_{E,\omega}\mathbf{f}\big)(\mathbf{r})
=
\int_{\Omega}
\overline{\mathbf{G}}_{E}(\mathbf{r},\mathbf{r}')\cdot\mathbf{f}(\mathbf{r}')\,d^{3}r',
\end{equation}
for \(\mathbf{f}\) in the forcing class of Eq.~\eqref{eq:ch2.helmholtz-forcing-map}
\cite{stratton1941}. When \(\mathbf{f}_{E,\omega}\) is the electric Helmholtz drive
induced by \(\Jimp\), Eqs.~\eqref{eq:ch2.electric-dyadic-convolution}
and~\eqref{eq:ch2.dyadic-green-integral-operator} coincide up to the linear map from
\(\Jimp\) to \(\mathbf{f}_{E,\omega}\) \cite{harrington2001}. The solution map on
impressed currents is therefore
\begin{equation}
\label{eq:ch2.solution-map-dyadic-realization}
\mathcal{S}_\omega[\Jimp]
=
\mathbf{u}_\omega
\quad\text{with}\quad
\E\ \text{from Eq.~\eqref{eq:ch2.electric-dyadic-convolution}},\
\H\ \text{from}\ \curl\E\ \text{when single-field reduction is valid},
\end{equation}
on domains satisfying Eq.~\eqref{eq:ch2.single-field-reduction-validity}; otherwise
the first-order dyadic system of Section~\ref{sec:ch234} applies
\cite{stratton1941}.

\begin{figure}[t]
\centering
\ChOneFigBlock{%
\begin{tikzpicture}[ch1fig, node distance=7mm and 10mm]
  \node[ch1box] (j) {$\Jimp(\mathbf{r}')$};
  \node[ch1box, right=14mm of j] (g) {$\overline{\mathbf{G}}_{E}$};
  \node[ch1box, right=14mm of g] (e) {$\E(\mathbf{r})$};
  \node[ch1box, below=11mm of g] (bc) {$\mathcal{B}_\Gamma$};
  \node[ch1box, below=11mm of e] (h) {$\H=\mu_0^{-1}\jj\omega^{-1}\curl\E$};
  \draw[ch1flow] (j) -- node[above,font=\scriptsize] {conv.} (g) -- (e);
  \draw[ch1flow] (bc) -- (g);
  \draw[ch1flow] (e) -- (h);
\end{tikzpicture}%
}
\caption{Dyadic inversion chain Eq.~\eqref{eq:ch2.electric-dyadic-convolution}:
impressed current convolves with \(\overline{\mathbf{G}}_{E}\) fixed by
\(\mathcal{B}_\Gamma\); \(\H\) follows from curl relations when reduction is admissible.}
\label{fig:ch2.dyadic-convolution-pipeline}
\end{figure}

Figure~\ref{fig:ch2.dyadic-convolution-pipeline} is the local inversion template used
in Sections~\ref{sec:ch25}--\ref{sec:ch27}. Boundary data enter through the kernel
\(G\) or \(\overline{\mathbf{G}}_{E}\) satisfying \(\mathcal{B}_\Gamma\), not through
a separate post-multiplication on \(\E\) alone \cite{stratton1941,harrington2001}.

\paragraph{Free-space electric dyad.}
In uniform vacuum with outgoing condition at infinity, let
\(\mathbf{R}=\mathbf{r}-\mathbf{r}'\), \(R=|\mathbf{R}|\), and
\(\widehat{\mathbf{R}}=\mathbf{R}/R\). The standard outgoing electric dyad is
\begin{equation}
\label{eq:ch2.free-space-dyadic-green}
\overline{\mathbf{G}}_{E,0}(\mathbf{r},\mathbf{r}')
=
\frac{e^{-\jj k_0 R}}{4\pi R}
\left[
\overline{\mathbf{I}}
-\widehat{\mathbf{R}}\widehat{\mathbf{R}}
+\left(\frac{1}{\jj k_0 R}-\frac{1}{k_0^{2}R^{2}}\right)
\Big(\overline{\mathbf{I}}-3\,\widehat{\mathbf{R}}\widehat{\mathbf{R}}\Big)
\right],
\end{equation}
valid for \(R>0\) \cite{stratton1941,harrington2001,jackson1999}. The
\(\overline{\mathbf{I}}-\widehat{\mathbf{R}}\widehat{\mathbf{R}}\) term is transverse
to \(\mathbf{R}\) and dominates in the radiative far zone; the terms in parentheses
supply near-field reactive structure through inverse powers of \(k_0 R\)
\cite{stratton1941}. Equation~\eqref{eq:ch2.free-space-dyadic-green} satisfies
Eq.~\eqref{eq:ch2.dyadic-green-defining-operator} for
\(\mathbf{r}\neq\mathbf{r}'\); coincidence evaluation is deferred to
Subsection~\ref{sec:ch232}. The same kernel appears in
Eq.~\eqref{eq:ch2.radiation-integral-preview} once \(\mathbf{r}\) lies outside
\(\operatorname{supp}\Jimp\) and outgoing closure is active
\cite{harrington2001}.

\paragraph{Magnetic dyadic and duality.}
The magnetic field may be generated by the electric dyad through Maxwell curl relations
when \(\Jimp\) is divergence-free in the source region:
\begin{equation}
\label{eq:ch2.magnetic-from-electric-dyad}
\H(\mathbf{r})
=
\frac{1}{\jj\omega\mu_0}\,\curl\E(\mathbf{r}),
\qquad
\E\ \text{from Eq.~\eqref{eq:ch2.electric-dyadic-convolution}},
\end{equation}
on simply connected source-free observation regions \cite{stratton1941}. A magnetic
dyadic \(\overline{\mathbf{G}}_{H}\) inverting
Eq.~\eqref{eq:ch2.magnetic-wave-operator} is defined analogously to
Eq.~\eqref{eq:ch2.dyadic-green-defining-operator} with \(\mathcal{W}_{H,\omega}\) in
place of \(\mathcal{W}_{E,\omega}\) \cite{harrington2001}. Duality
\(\varepsilon_0\E\leftrightarrow\mu_0\H\), \(\Jimp\leftrightarrow\mathbf{J}^{\mathrm{(mag)}}_{\mathrm{eff}}\)
maps \(\overline{\mathbf{G}}_{E,0}\leftrightarrow\overline{\mathbf{G}}_{H,0}\) in
vacuum without re-deriving the magnetic-wave symbol. Subsection~\ref{sec:ch231} states
\(\overline{\mathbf{G}}_{H}\) only where magnetic impressed currents are primary;
otherwise Eq.~\eqref{eq:ch2.magnetic-from-electric-dyad} avoids duplicate kernel
tabulations \cite{stratton1941}.

\paragraph{Material weighting and reciprocity.}
In piecewise media, local symbols \(k(\mathbf{r})\), \(\varepsilon(\mathbf{r})\),
\(\mu(\mathbf{r})\) modify Eq.~\eqref{eq:ch2.dyadic-green-defining-operator} to
\begin{equation}
\label{eq:ch2.material-dyadic-defining-operator}
\mathcal{W}_{E,\omega}^{(\varepsilon,\mu)}\,
\overline{\mathbf{G}}_{E}^{(\varepsilon,\mu)}(\cdot,\mathbf{r}')
=
\overline{\mathbf{I}}\,\delta(\mathbf{r}-\mathbf{r}'),
\end{equation}
with \(\mathcal{W}_{E,\omega}^{(\varepsilon,\mu)}\) from
Eq.~\eqref{eq:ch2.material-helmholtz-electric} \cite{stratton1941}. Free-space
Eq.~\eqref{eq:ch2.free-space-dyadic-green} applies patchwise only when
\(k(\mathbf{r})\) is constant in each patch and interface jumps are enforced by
continuity of tangential \(\E,\H\) across shared boundaries
(Eq.~\eqref{eq:ch1.interface-conditions}); multi-patch assembly is
Subsection~\ref{sec:ch234} \cite{harrington2001}. Reciprocity for fixed \(\omega\)
reads
\begin{equation}
\label{eq:ch2.dyadic-reciprocity-swap}
\overline{\mathbf{G}}_{E}(\mathbf{r},\mathbf{r}')
=
\overline{\mathbf{G}}_{E}^{\mathsf{T}}(\mathbf{r}',\mathbf{r}),
\end{equation}
for passive media without magneto-electric coupling; full proof and lossless
restrictions are in Section~\ref{sec:ch28} \cite{stratton1941}. Equation
\eqref{eq:ch2.dyadic-reciprocity-swap} is recorded here because it constrains
probe-source interchange in Eq.~\eqref{eq:ch2.electric-dyadic-convolution}.

\paragraph{Boundary of validity.}
Equations~\eqref{eq:ch2.electric-dyadic-convolution}--\eqref{eq:ch2.solution-map-dyadic-realization}
hold when \(\Jimp\) lies in the domain of the forcing map, when
\(\mathcal{B}_\Gamma\) is time-independent and identical under phasor reduction, and
when the selected kernel branch matches the resolvent of
Eq.~\eqref{eq:ch2.resolvent-green-equivalence}. They fail for resonant cavities
without quotienting, for incoming rather than outgoing exterior conditions, and for
sources with \(\divg\Jimp\neq 0\) unless charge density \(\rho\) is included in the
complete source tuple of Section~\ref{sec:ch01.1.2} \cite{stratton1941,harrington2001}.
Near \(\mathbf{r}=\mathbf{r}'\), Eqs.~\eqref{eq:ch2.free-space-scalar-green}
and~\eqref{eq:ch2.free-space-dyadic-green} are pointwise undefined; convolution
meaning is distributional (Subsection~\ref{sec:ch232}).

\paragraph{Worked illustration (Hertzian dipole).}
Replace \(\Jimp(\mathbf{r}')\) by a localized current density with dipole moment
\(\mathbf{p}\) aligned with \(\hat{\mathbf{z}}\) in vacuum. Evaluating
Eq.~\eqref{eq:ch2.electric-dyadic-convolution} with
Eq.~\eqref{eq:ch2.free-space-dyadic-green} at \(k_0 r\gg 1\) recovers the familiar
\(\theta\)-dependent far pattern proportional to
\(\widehat{\boldsymbol{\theta}}\times(\widehat{\boldsymbol{\theta}}\times\mathbf{p})\);
the near zone is dominated by reactive terms in Eq.~\eqref{eq:ch2.free-space-dyadic-green}
and lies largely in \(\Pi_{\mathrm{ev}}\mathcal{S}_\omega[\Jimp]\) of
Subsection~\ref{sec:ch224} \cite{stratton1941,harrington2001}. Declaring a transport
channel from near-zone reactive scans without
Eq.~\eqref{eq:ch2.regime-operator-screen} repeats the regime error class of
Section~\ref{sec:ch01.2.4}.

\paragraph{Worked illustration (filamentary current).}
A thin wire carrying uniform phasor current \(I_0\) produces \(\Jimp\) supported on a
one-dimensional manifold. Equation~\eqref{eq:ch2.electric-dyadic-convolution} remains
formally valid, but the integral requires distributional reduction to line integrals
with principal-value treatment of the \(R^{-3}\) singularity when observation points
approach the wire \cite{stratton1941,jackson1999}. Subsection~\ref{sec:ch232} supplies
the finite-part prescription; Subsection~\ref{sec:ch231} only identifies the kernel
\(\overline{\mathbf{G}}_{E,0}\) as the continuum target.

\paragraph{Closing summary.}
Subsection~\ref{sec:ch231} establishes the Chapter~\ref{ch:02} HOME for scalar
Helmholtz Eqs.~\eqref{eq:ch2.scalar-helmholtz-green-pde}--\eqref{eq:ch2.free-space-scalar-green},
Lorenz reduction Eq.~\eqref{eq:ch2.lorenz-field-from-potential}, dyadic defining
operator Eq.~\eqref{eq:ch2.dyadic-green-defining-operator}, electric convolution
Eqs.~\eqref{eq:ch2.electric-dyadic-convolution}--\eqref{eq:ch2.solution-map-dyadic-realization},
free-space dyad Eq.~\eqref{eq:ch2.free-space-dyadic-green}, magnetic recovery
Eq.~\eqref{eq:ch2.magnetic-from-electric-dyad}, material extension
Eq.~\eqref{eq:ch2.material-dyadic-defining-operator}, and reciprocity swap
Eq.~\eqref{eq:ch2.dyadic-reciprocity-swap}. Subsection~\ref{sec:ch232} completes
coincidence and self-term theory; Subsection~\ref{sec:ch233} ties harmonic kernels to
retarded time-domain propagation.

\subsection{Singularities, distributions, and regularization}
\label{sec:ch232}
% TOC tag: [HOME]

Subsection~\ref{sec:ch231} defined scalar and dyadic Green kernels that satisfy
Eqs.~\eqref{eq:ch2.scalar-helmholtz-green-pde} and~\eqref{eq:ch2.dyadic-green-defining-operator}
for \(\mathbf{r}\neq\mathbf{r}'\) and deferred coincidence evaluation to the present
subsection \cite{stratton1941,harrington2001,jackson1999}. Subsection~\ref{sec:ch232}
gives the Chapter~\ref{ch:02} HOME for distributional convolutions, principal-value
and finite-part prescriptions, static--dynamic kernel splits, and self-term
regularization that make Eq.~\eqref{eq:ch2.electric-dyadic-convolution} well-defined
on \(\Jimp\in\mathcal{J}_\Omega\) including sheet and filamentary support. The
Sobolev and distributional source classes of Subsection~\ref{sec:ch211} are presupposed
through Eqs.~\eqref{eq:ch2.distributional-source-extension} and~\eqref{eq:ch2.dual-source-pairing};
their functional-analytic construction is not repeated. Retarded time-domain support
belongs to Subsection~\ref{sec:ch233}; outgoing branch selection at infinity belongs
to Section~\ref{sec:ch24}.

\paragraph{Assumptions.}
Let \(\Omega\Subset\mathbb{R}^{3}\) be Lipschitz with piecewise smooth internal
interfaces, \(\omega>0\) fixed, and \(\Jimp\) in the refined source class
Eq.~\eqref{eq:ch2.source-space-refined} or its distributional extension
Eq.~\eqref{eq:ch2.distributional-source-extension}. Kernels are the vacuum templates
Eqs.~\eqref{eq:ch2.free-space-scalar-green} and~\eqref{eq:ch2.free-space-dyadic-green}
unless \(\mathcal{B}_\Gamma\) or piecewise \(k(\mathbf{r})\) modify them patchwise.
Observation points \(\mathbf{r}\) may approach \(\operatorname{supp}\Jimp\) but
remain in the domain where \(\E\in\HCurl(\Omega)\) is sought
\cite{stratton1941,harrington2001}.

\paragraph{Near-field singular structure.}
For \(\mathbf{r}\notin\operatorname{supp}\Jimp\), the convolution
Eq.~\eqref{eq:ch2.electric-dyadic-convolution} is a classical integral. As
\(\mathbf{r}\to\mathbf{r}'\in\operatorname{supp}\Jimp\), the free-space dyad
Eq.~\eqref{eq:ch2.free-space-dyadic-green} expands in powers of \(k_0 R\) with
leading static singularity \cite{stratton1941,jackson1999}. With
\(\mathbf{R}=\mathbf{r}-\mathbf{r}'\), \(R=|\mathbf{R}|\),
\(\widehat{\mathbf{R}}=\mathbf{R}/R\),
\begin{equation}
\label{eq:ch2.dyadic-singular-expansion}
\overline{\mathbf{G}}_{E,0}(\mathbf{r},\mathbf{r}')
=
\frac{1}{4\pi R^{3}}\,\overline{\mathbf{P}}_{t}(\widehat{\mathbf{R}})
\;+\;
\frac{\jj k_0}{4\pi R^{2}}\,\overline{\mathbf{Q}}(\widehat{\mathbf{R}})
\;+\;
\frac{k_0^{2}}{4\pi R}\,\overline{\mathbf{S}}(\widehat{\mathbf{R}})
\;+\;\mathcal{O}(k_0^{3}),
\end{equation}
where \(\overline{\mathbf{P}}_{t}(\widehat{\mathbf{R}})=\overline{\mathbf{I}}-3\widehat{\mathbf{R}}\widehat{\mathbf{R}}\)
is transverse to \(\mathbf{R}\) and \(\overline{\mathbf{Q}}\), \(\overline{\mathbf{S}}\)
collect the next two anisotropic terms in the locked phasor convention
\cite{harrington2001}. Equation~\eqref{eq:ch2.dyadic-singular-expansion} matches the
\(R^{-3}\), \(R^{-2}\), and \(R^{-1}\) pieces of
Eq.~\eqref{eq:ch2.free-space-dyadic-green} and separates reactive near storage from
radiative \(R^{-1}\) decay. The \(R^{-3}\) term is non-integrable in a neighborhood
of coincidence; convolution with \(\Jimp\) therefore requires a distributional or
regularized definition, not pointwise Riemann integration
\cite{stratton1941}. The scalar kernel Eq.~\eqref{eq:ch2.free-space-scalar-green}
carries the parallel hierarchy \(R^{-1}\), \(R^{0}\), \ldots\ with
\(R^{-1}\) non-integrable on neighborhoods of the source.

\paragraph{Static limit and transverse projection.}
In the electrostatic limit \(k_0\to 0\) at fixed \(R>0\),
Eq.~\eqref{eq:ch2.dyadic-singular-expansion} reduces to
\begin{equation}
\label{eq:ch2.static-dyadic-limit}
\overline{\mathbf{G}}_{E,0}^{\mathrm{(st)}}(\mathbf{r},\mathbf{r}')
=
\frac{1}{4\pi R^{3}}\Big(\overline{\mathbf{I}}-3\,\widehat{\mathbf{R}}\widehat{\mathbf{R}}\Big),
\end{equation}
the field of a unit dipole aligned with each Cartesian source direction
\cite{jackson1999,stratton1941}. Equation~\eqref{eq:ch2.static-dyadic-limit} is the
leading singular part extracted in method-of-moments static--dynamic splits: dynamic
corrections are \(\mathcal{O}(k_0 R)\) relative to the static term when
\(k_0 R\ll 1\) on source patches \cite{harrington2001}. For divergence-free
\(\Jimp\), the \(\grad\,\divg\) structure of Lorenz reduction
Eq.~\eqref{eq:ch2.lorenz-field-from-potential} removes monopole artifacts; the
dominant self-coupling on compact patches is dipolar and transverse to the separation
vector between observation and source centroids \cite{stratton1941}.

\begin{figure}[t]
\centering
\ChOneFigBlock{%
\begin{tikzpicture}[ch1fig, node distance=7mm and 10mm]
  \node[ch1box] (src) {Source patch\\$\operatorname{supp}\Jimp$};
  \node[ch1box, right=16mm of src] (near) {Near zone\\$R\to 0$};
  \node[ch1box, right=16mm of near] (reg) {Regularized\\convolution};
  \node[ch1box, below=11mm of near] (sing) {$G_{\mathrm{sing}}$};
  \node[ch1box, below=11mm of reg] (cpt) {$K_{\mathrm{cpt}}$};
  \draw[ch1flow] (src) -- (near) -- (reg);
  \draw[ch1flow] (near) -- (sing);
  \draw[ch1flow] (reg) -- (cpt);
\end{tikzpicture}%
}
\caption{Singular--regular split at coincidence: non-integrable \(G_{\mathrm{sing}}\)
is extracted analytically; remainder \(K_{\mathrm{cpt}}\) is smooth for Galerkin or
point evaluation away from the extracted neighborhood.}
\label{fig:ch2.singular-regularization-split}
\end{figure}

Figure~\ref{fig:ch2.singular-regularization-split} implements the operator template
Eq.~\eqref{eq:ch2.green-compact-split} at kernel level. Write
\begin{equation}
\label{eq:ch2.kernel-singular-regular-split}
\overline{\mathbf{G}}_{E}(\mathbf{r},\mathbf{r}')
=
\overline{\mathbf{G}}_{E}^{\mathrm{(sing)}}(\mathbf{r},\mathbf{r}')
+
\overline{\mathbf{G}}_{E}^{\mathrm{(reg)}}(\mathbf{r},\mathbf{r}'),
\end{equation}
with \(\overline{\mathbf{G}}_{E}^{\mathrm{(sing)}}\) collecting the non-integrable
static and leading dynamic singularities of Eq.~\eqref{eq:ch2.dyadic-singular-expansion}
and \(\overline{\mathbf{G}}_{E}^{\mathrm{(reg)}}\) bounded on
\(\Omega\times\Omega\) \cite{stratton1941,harrington2001}. The corresponding
solution operator split is
\begin{equation}
\label{eq:ch2.solution-singular-compact-split}
\mathcal{S}_\omega
=
\mathcal{G}_{\mathrm{sing}}+\mathcal{K}_{\mathrm{cpt}},
\qquad
\mathcal{G}_{\mathrm{sing}}[\Jimp]
=
\jj\omega\mu_0\int_{\Omega}
\overline{\mathbf{G}}_{E}^{\mathrm{(sing)}}(\mathbf{r},\mathbf{r}')\cdot\Jimp(\mathbf{r}')\,d^{3}r',
\end{equation}
with \(\mathcal{K}_{\mathrm{cpt}}\) compact on bounded \(\Omega\) when
\(\overline{\mathbf{G}}_{E}^{\mathrm{(reg)}}\) is square-integrable in both variables
\cite{stratton1941}. Equations~\eqref{eq:ch2.kernel-singular-regular-split}
and~\eqref{eq:ch2.solution-singular-compact-split} realize
Eq.~\eqref{eq:ch2.green-compact-split} for Green-based inversion; Fredholm consequences
are developed in Section~\ref{sec:ch263} \cite{harrington2001}.

\paragraph{Distributional convolution.}
When \(\Jimp\) is a bounded function, Eq.~\eqref{eq:ch2.electric-dyadic-convolution}
defines \(\E\in L^{2}_{\mathrm{loc}}(\Omega)^{3}\) away from the source. For sheet
currents Eq.~\eqref{eq:ch2.distributional-source-extension}, the convolution is defined
by duality: for all smooth compactly supported test fields \(\boldsymbol{\phi}\),
\begin{equation}
\label{eq:ch2.distributional-dyadic-pairing}
\langle\E,\boldsymbol{\phi}\rangle_{\Omega}
=
\jj\omega\mu_0
\sum_{ab}\int_{\Sigma_{ab}}\int_{\Omega}
\mathbf{J}_{s,ab}(\mathbf{r}')
\cdot\overline{\mathbf{G}}_{E}^{\mathsf{T}}(\mathbf{r}',\mathbf{r})\cdot\boldsymbol{\phi}(\mathbf{r})\,d^{3}r\,dS',
\end{equation}
with appropriate Fubini interpretation and outgoing closure on exterior problems
\cite{stratton1941}. Equation~\eqref{eq:ch2.distributional-dyadic-pairing} is the
weak form of Eq.~\eqref{eq:ch2.electric-dyadic-convolution} when pointwise evaluation
fails; it pairs with Eq.~\eqref{eq:ch2.maxwell-weak-form} once \(\boldsymbol{\phi}\)
is replaced by admissible test states in \(\mathcal{F}_\Omega^{0}\)
\cite{harrington2001}. Physically, sheet sources produce jump conditions on
\(\hat{\mathbf{n}}\times\E\) and \(\hat{\mathbf{n}}\times\H\); the distributional
definition ensures those jumps arise from the same kernel as volume currents, not from
ad hoc post-factors \cite{stratton1941}.

\paragraph{Principal value and finite part.}
For volume \(\Jimp\) and observation points \(\mathbf{r}\notin\operatorname{supp}\Jimp\)
approaching the support, define the principal-value convolution
\begin{equation}
\label{eq:ch2.principal-value-convolution}
\E(\mathbf{r})
=
\jj\omega\mu_0\,\mathrm{p.v.}\int_{\Omega}
\overline{\mathbf{G}}_{E}(\mathbf{r},\mathbf{r}')\cdot\Jimp(\mathbf{r}')\,d^{3}r',
\end{equation}
as the limit of integrals over \(\Omega\setminus B_{\varepsilon}(\mathbf{r})\) as
\(\varepsilon\downarrow 0\) when that limit exists in \(\HCurl(\Omega)\)
\cite{jackson1999,stratton1941}. For observation on the source support---for example
\(\mathbf{r}\in\Sigma_{ab}\) for a surface patch---the principal value generally
diverges because of the \(R^{-3}\) term; the \emph{Hadamard finite part} subtracts
the singular expansion explicitly:
\begin{equation}
\label{eq:ch2.finite-part-convolution}
\E(\mathbf{r})
=
\jj\omega\mu_0\lim_{\varepsilon\to 0}
\left[
\int_{\Omega\setminus B_{\varepsilon}(\mathbf{r})}
\overline{\mathbf{G}}_{E}(\mathbf{r},\mathbf{r}')\cdot\Jimp(\mathbf{r}')\,d^{3}r'
\;+\;
\int_{B_{\varepsilon}(\mathbf{r})\cap\Omega}
\overline{\mathbf{G}}_{E}^{\mathrm{(sing)}}(\mathbf{r},\mathbf{r}')\cdot\Jimp(\mathbf{r}')\,d^{3}r'
\right],
\end{equation}
provided \(\Jimp\) is smooth enough on the patch that the limit exists
\cite{harrington2001}. Equation~\eqref{eq:ch2.finite-part-convolution} is the
Chapter~\ref{ch:02} HOME prescription for on-surface field evaluation; it differs
from Eq.~\eqref{eq:ch2.principal-value-convolution} by the explicit analytic
subtraction of \(\overline{\mathbf{G}}_{E}^{\mathrm{(sing)}}\), not by numerical
mesh refinement alone \cite{stratton1941}. Mathematically, the finite part is the
constant term in the Laurent expansion of the integral as \(\varepsilon\to 0\);
physically, it removes infinite self-energy of pointwise source models while retaining
finite reactive coupling between extended patches \cite{jackson1999}.

\paragraph{Self-term and Galerkin testing.}
Method-of-moments and boundary-element codes replace
Eq.~\eqref{eq:ch2.electric-dyadic-convolution} by Galerkin entries
\begin{equation}
\label{eq:ch2.mom-dyadic-entry}
Z_{mn}
=
\jj\omega\mu_0\int_{\Omega_m}\int_{\Omega_n}
\mathbf{f}_m(\mathbf{r})\cdot
\overline{\mathbf{G}}_{E}(\mathbf{r},\mathbf{r}')\cdot
\mathbf{f}_n(\mathbf{r}')\,d^{3}r\,d^{3}r',
\end{equation}
for basis functions \(\mathbf{f}_m,\mathbf{f}_n\) on patches \(\Omega_m,\Omega_n\)
\cite{harrington2001}. When \(m=n\), the integrand is non-integrable at
\(\mathbf{r}=\mathbf{r}'\); the Chapter~\ref{ch:02} regularized self-term is
\begin{equation}
\label{eq:ch2.mom-self-term-regularized}
Z_{mm}^{\mathrm{(reg)}}
=
\jj\omega\mu_0\int_{\Omega_m}\int_{\Omega_m}
\mathbf{f}_m(\mathbf{r})\cdot
\overline{\mathbf{G}}_{E}^{\mathrm{(reg)}}(\mathbf{r},\mathbf{r}')\cdot
\mathbf{f}_m(\mathbf{r}')\,d^{3}r\,d^{3}r'
\;+\;
Z_{mm}^{\mathrm{(sing)}},
\end{equation}
where \(Z_{mm}^{\mathrm{(sing)}}\) is the analytic integral of
\(\overline{\mathbf{G}}_{E}^{\mathrm{(sing)}}\) over the coincident patch, often
reduced to closed forms for flat triangles and rectangles \cite{stratton1941}. Point
matching---evaluating Eq.~\eqref{eq:ch2.electric-dyadic-convolution} at patch centers
with \(\overline{\mathbf{G}}_{E}\) replaced by
\(\overline{\mathbf{G}}_{E}^{\mathrm{(reg)}}\) only---is \emph{not} equivalent to
Eq.~\eqref{eq:ch2.mom-self-term-regularized} unless basis overlap and singular
subtraction match the same finite-part convention \cite{harrington2001}. The
distinction is operator-theoretic: Galerkin testing defines a bounded bilinear form on
\(\mathcal{J}_\Omega\times\mathcal{J}_\Omega\); raw point evaluation does not
\cite{stratton1941}.

\paragraph{Near-singular subtraction.}
When \(\mathbf{r}\) and \(\mathbf{r}'\) lie on distinct patches separated by small
distance \(d\ll \lambda_0=2\pi/k_0\), direct quadrature of
Eq.~\eqref{eq:ch2.free-space-dyadic-green} is ill-conditioned. The standard
subtraction template is
\begin{equation}
\label{eq:ch2.near-singular-subtraction}
\int_{\Omega_n}
\overline{\mathbf{G}}_{E}(\mathbf{r},\mathbf{r}')\cdot\Jimp(\mathbf{r}')\,d^{3}r'
=
\int_{\Omega_n}
\Big[\overline{\mathbf{G}}_{E}-\overline{\mathbf{G}}_{E}^{\mathrm{(near)}}(\mathbf{r},\mathbf{r}')\Big]
\cdot\Jimp(\mathbf{r}')\,d^{3}r'
\;+\;
\int_{\Omega_n}
\overline{\mathbf{G}}_{E}^{\mathrm{(near)}}(\mathbf{r},\mathbf{r}')\cdot\Jimp(\mathbf{r}')\,d^{3}r',
\end{equation}
where \(\overline{\mathbf{G}}_{E}^{\mathrm{(near)}}\) matches the first few terms of
Eq.~\eqref{eq:ch2.dyadic-singular-expansion} about \(\mathbf{r}\) and the integral of
\(\overline{\mathbf{G}}_{E}^{\mathrm{(near)}}\) is evaluated analytically
\cite{harrington2001,stratton1941}. Equation~\eqref{eq:ch2.near-singular-subtraction}
does not alter the continuum operator; it stabilizes numerical approximation of
\(\mathcal{K}_{\mathrm{cpt}}\) in Eq.~\eqref{eq:ch2.solution-singular-compact-split}.

\paragraph{Scalar coincidence limit.}
The scalar Green function satisfies, in distributional sense,
\begin{equation}
\label{eq:ch2.scalar-coincidence-limit}
\int_{B_{\varepsilon}(\mathbf{r})}
G_0(\mathbf{r},\mathbf{r}')\,d^{3}r'
=
\frac{\varepsilon}{2}
\;+\;\mathcal{O}(\varepsilon^{3}),
\end{equation}
for small balls \(B_{\varepsilon}(\mathbf{r})\) in \(\mathbb{R}^{3}\), reflecting the
\(R^{-1}\) integrability borderline in three dimensions \cite{jackson1999}. Equation
\eqref{eq:ch2.scalar-coincidence-limit} explains why monopole self-terms in scalar
potential formulations require either finite source extent or explicit regularization;
vector dyads inherit stronger \(R^{-3}\) divergence captured by
Eq.~\eqref{eq:ch2.static-dyadic-limit} \cite{stratton1941}.

\paragraph{Boundary of validity.}
Equations~\eqref{eq:ch2.finite-part-convolution}--\eqref{eq:ch2.mom-self-term-regularized}
require patchwise smooth \(\Jimp\), Lipschitz \(\Omega\), and fixed \(\omega\) in the
passive linear band. They fail for point dipoles modeled as \(\delta\)-supported
\(\Jimp\) without accompanying finite-part renormalization, and for non-Lipschitz
edges where \(\HCurl\) membership of \(\E\) requires additional edge singularities
not captured by Eq.~\eqref{eq:ch2.dyadic-singular-expansion}
\cite{stratton1941,harrington2001}. Lossy media with complex \(k_0\) shift singular
expansion coefficients but not the order of the \(R^{-3}\) pole; Subsection~\ref{sec:ch232}
does not re-tabulate complex-\(k\) subtractors \cite{stratton1941}.

\paragraph{Worked illustration (flat triangular patch).}
A PEC-backed triangular patch carries \(\Jimp=\hat{\mathbf{n}}\times\H_{\mathrm{exc}}\)
in the method-of-moments sense. The self impedance uses
Eq.~\eqref{eq:ch2.mom-self-term-regularized} with
\(\overline{\mathbf{G}}_{E}^{\mathrm{(sing)}}\) from
Eq.~\eqref{eq:ch2.static-dyadic-limit} integrated over the triangle and
\(\overline{\mathbf{G}}_{E}^{\mathrm{(reg)}}\) integrated numerically. Omitting
\(Z_{mm}^{\mathrm{(sing)}}\) while retaining the \(R^{-3}\) kernel in the numerical
part produces a finite but incorrect reactive component; the error is largest at
\(k_0 h\ll 1\) for patch characteristic size \(h\) \cite{harrington2001,stratton1941}.

\paragraph{Worked illustration (on-surface observation).}
For \(\mathbf{r}\in\Sigma_{ab}\), Eq.~\eqref{eq:ch2.principal-value-convolution}
diverges because the exclusion ball removes a set of measure zero in the surface
integral sense. Eq.~\eqref{eq:ch2.finite-part-convolution} with
\(\overline{\mathbf{G}}_{E}^{\mathrm{(sing)}}\) from
Eq.~\eqref{eq:ch2.dyadic-singular-expansion} yields finite tangential \(\E\) consistent
with the jump conditions imported from Eq.~\eqref{eq:ch1.interface-conditions}
\cite{stratton1941}. Confusing principal-value limits with finite-part values is a
common source of incorrect boundary-element tangential fields \cite{harrington2001}.

\paragraph{Closing summary.}
Subsection~\ref{sec:ch232} establishes the Chapter~\ref{ch:02} HOME for singular
expansion Eq.~\eqref{eq:ch2.dyadic-singular-expansion}, static limit
Eq.~\eqref{eq:ch2.static-dyadic-limit}, kernel split
Eqs.~\eqref{eq:ch2.kernel-singular-regular-split}--\eqref{eq:ch2.solution-singular-compact-split},
distributional pairing Eq.~\eqref{eq:ch2.distributional-dyadic-pairing},
principal-value Eq.~\eqref{eq:ch2.principal-value-convolution}, finite-part
Eq.~\eqref{eq:ch2.finite-part-convolution}, MoM entries
Eqs.~\eqref{eq:ch2.mom-dyadic-entry}--\eqref{eq:ch2.mom-self-term-regularized},
near-singular subtraction Eq.~\eqref{eq:ch2.near-singular-subtraction}, and scalar
coincidence limit Eq.~\eqref{eq:ch2.scalar-coincidence-limit}. Subsection~\ref{sec:ch233}
promotes harmonic kernels to retarded causal propagation; Subsection~\ref{sec:ch234}
uses the regularized convolution on composite domains.

\subsection{Retardation and causal propagation}
\label{sec:ch233}
% TOC tag: [HOME]

Subsections~\ref{sec:ch231}--\ref{sec:ch232} defined harmonic Green kernels and their
regularized convolutions on \(\Jimp\in\mathcal{J}_\Omega\). Subsection~\ref{sec:ch233}
gives the Chapter~\ref{ch:02} HOME for \emph{retarded} time-domain Green operators,
their causal support, and the kernel-level upgrade of the resolvent equivalence
Eq.~\eqref{eq:ch2.resolvent-green-equivalence} and harmonic theorem
Eq.~\eqref{eq:ch2.harmonic-equivalence-theorem} from Subsection~\ref{sec:ch223}
\cite{stratton1941,harrington2001,jackson1999}. Operator-level causality and
Laplace--Fourier bridges are presupposed through
Eqs.~\eqref{eq:ch2.causal-operator-support}--\eqref{eq:ch2.fourier-phasor-identification};
the present subsection constructs the Green kernels that realize those statements.
Harmonic formulas Eqs.~\eqref{eq:ch2.free-space-scalar-green} and~\eqref{eq:ch2.free-space-dyadic-green}
are not re-derived; they appear as \(s=-\jj\omega\) evaluations of causal Laplace
transforms. Sommerfeld outgoing selection at infinity belongs to Section~\ref{sec:ch24};
multi-region solvability belongs to Subsection~\ref{sec:ch234}.

\paragraph{Assumptions.}
Let \(\mathbf{j}_{\mathrm{imp}}(\mathbf{r},t)\) be causal with
\(\operatorname{supp}_{(\mathbf{r},t)}\subseteq\Omega_s\times[0,\infty)\),
\(\Omega_s\Subset\Omega\), and let \(\mathcal{B}_\Gamma\), \(\varepsilon\), and
\(\mu\) be time-independent. Initial data are homogeneous,
\(\mathbf{U}_0=\mathbf{0}\), unless noted for switch-on illustrations. Vacuum
regions use \(c_0=1/\sqrt{\mu_0\varepsilon_0}\) and \(k_0=\omega/c_0\) from
Eq.~\eqref{eq:ch2.helmholtz-symbol} \cite{stratton1941}. The equivalence class
excludes advanced kernels satisfying Eq.~\eqref{eq:ch2.advanced-kernel-violation}.

\paragraph{Retarded scalar kernel.}
The causal scalar Green function for the wave operator
\((\nabla^{2}-\varepsilon_0\mu_0\,\partial_t^{2})\) in vacuum satisfies
\begin{equation}
\label{eq:ch2.retarded-scalar-wave-operator}
\Big(\nabla^{2}-\varepsilon_0\mu_0\,\partial_t^{2}\Big)\,
G(\mathbf{r},\mathbf{r}',t-t')
=
-\delta(\mathbf{r}-\mathbf{r}')\,\delta(t-t'),
\end{equation}
with retarded causality \(G=\mathbf{0}\) for \(t-t'<0\) \cite{stratton1941,jackson1999}.
The free-space outgoing solution is
\begin{equation}
\label{eq:ch2.retarded-scalar-green-time}
G_0(\mathbf{r},\mathbf{r}',t-t')
=
\frac{\delta\!\big(t-t'-R/c_0\big)}{4\pi R},
\qquad
R=|\mathbf{r}-\mathbf{r}'|,
\end{equation}
which transports impulse support on the light cone \(t-t'=R/c_0\)
\cite{stratton1941}. Equation~\eqref{eq:ch2.retarded-scalar-green-time} is the
time-domain counterpart of Eq.~\eqref{eq:ch2.free-space-scalar-green}: the locked
phasor field \(\exp(-\jj\omega t)\) selects the \(s=-\jj\omega\) Laplace transform
of Eq.~\eqref{eq:ch2.retarded-scalar-green-time} on the outgoing branch
\cite{harrington2001}. Physically, Eq.~\eqref{eq:ch2.retarded-scalar-green-time}
states that a unit impulse at \((\mathbf{r}',t')\) influences \(\mathbf{r}\) only
after travel time \(R/c_0\), enforcing finite propagation speed and
Eq.~\eqref{eq:ch2.causal-support-constraint} on potentials generated by scalar
convolution \cite{stratton1941}.

\paragraph{Retarded dyadic kernel.}
The retarded electric dyadic \(\overline{\mathbf{G}}_{E}(\mathbf{r},\mathbf{r}',t-t')\)
upgrades Eq.~\eqref{eq:ch2.electric-dyadic-convolution} to the time domain. The
causal Maxwell field is
\begin{equation}
\label{eq:ch2.retarded-dyadic-convolution}
\E(\mathbf{r},t)
=
\E_{\mathrm{hom}}(\mathbf{r},t)
-\mu_0\,\partial_t\int_{0}^{t}\!\!\int_{\Omega}
\overline{\mathbf{G}}_{E}(\mathbf{r},\mathbf{r}',t-t')\cdot
\mathbf{j}_{\mathrm{imp}}(\mathbf{r}',t')\,d^{3}r'\,dt',
\end{equation}
which specializes Eq.~\eqref{eq:ch2.retarded-convolution-form} once the dyadic kernel
is identified \cite{stratton1941,harrington2001}. Retarded causality requires
\begin{equation}
\label{eq:ch2.retarded-dyadic-support}
\overline{\mathbf{G}}_{E}(\mathbf{r},\mathbf{r}',t-t')
=
\mathbf{0}\quad\text{for}\ t-t'<0,
\end{equation}
analogous to Eq.~\eqref{eq:ch2.retarded-kernel-support} \cite{stratton1941}. For
\(\mathbf{r}\notin\Omega_s\), Eq.~\eqref{eq:ch2.retarded-dyadic-convolution} produces
\(\E(\mathbf{r},t)=\mathbf{0}\) until the earliest retarded arrival time
\begin{equation}
\label{eq:ch2.light-cone-arrival-time}
t_{\mathrm{ret}}(\mathbf{r})
=
\inf_{(\mathbf{r}',t')\in\operatorname{supp}\mathbf{j}_{\mathrm{imp}}}
\big\{t'\,+\,R/c_0\big\},
\end{equation}
which is the geometric light-cone support condition underlying
Eq.~\eqref{eq:ch2.causal-support-constraint} \cite{jackson1999,stratton1941}. The
harmonic dyad Eq.~\eqref{eq:ch2.free-space-dyadic-green} is not repeated here; it is
the \(s=-\jj\omega\) symbol of \(\overline{\mathbf{G}}_{E,0}\) under
Eqs.~\eqref{eq:ch2.equivalence-assumptions} and outgoing closure
\cite{harrington2001}.

\begin{figure}[t]
\centering
\ChOneFigBlock{%
\begin{tikzpicture}[ch1fig, node distance=7mm and 10mm]
  \node[ch1box] (src) {$(\mathbf{r}',t')$\\source event};
  \node[ch1box, right=16mm of src] (cone) {Light cone\\$t-t'=R/c_0$};
  \node[ch1box, right=16mm of cone] (obs) {$(\mathbf{r},t)$\\observation};
  \node[ch1box, below=11mm of cone] (adv) {Advanced\\excluded};
  \draw[ch1flow] (src) -- (cone) -- (obs);
  \draw[ch1flow,dashed] (src) to[bend right=20] (adv);
\end{tikzpicture}%
}
\caption{Retarded support Eq.~\eqref{eq:ch2.retarded-dyadic-support}: field at
\((\mathbf{r},t)\) depends only on sources with \(t'\le t-R/c_0\); advanced
coupling is excluded.}
\label{fig:ch2.retarded-light-cone}
\end{figure}

Figure~\ref{fig:ch2.retarded-light-cone} is the space--time gate for
Eq.~\eqref{eq:ch2.retarded-dyadic-convolution}. Subsection~\ref{sec:ch232} treated
coincidence in \(\mathbf{r}\); the present subsection treats ordering in \(t\)
\cite{stratton1941}. Regularized self-terms from
Eq.~\eqref{eq:ch2.mom-self-term-regularized} apply to harmonic convolutions; time
stepping discretizes Eq.~\eqref{eq:ch2.retarded-dyadic-convolution} with the same
kernel class \cite{harrington2001}.

\paragraph{Laplace transform and harmonic extraction.}
Let \(\widetilde{\overline{\mathbf{G}}}_{E}(\mathbf{r},\mathbf{r}',s)\) denote the
Laplace transform of \(\overline{\mathbf{G}}_{E}(\mathbf{r},\mathbf{r}',t)\) in
\(t-t'\) with \(\Re\{s\}\) large enough for convergence. Causal support
Eq.~\eqref{eq:ch2.retarded-dyadic-support} implies
\(\widetilde{\overline{\mathbf{G}}}_{E}\) is analytic in the right half-plane and,
under outgoing closure,
\begin{equation}
\label{eq:ch2.retarded-laplace-harmonic-link}
\overline{\mathbf{G}}_{E,\omega}(\mathbf{r},\mathbf{r}')
=
\eval{\widetilde{\overline{\mathbf{G}}}_{E}(\mathbf{r},\mathbf{r}',s)}_{s=-\jj\omega},
\end{equation}
which is the kernel form of Eq.~\eqref{eq:ch2.resolvent-green-equivalence}
\cite{stratton1941,harrington2001}. Applying the Laplace transform to
Eq.~\eqref{eq:ch2.retarded-dyadic-convolution} with causal
\(\mathbf{j}_{\mathrm{imp}}\) yields
\begin{equation}
\label{eq:ch2.laplace-dyadic-convolution}
\widetilde{\E}(\mathbf{r},s)
=
\widetilde{\E}_{\mathrm{hom}}(\mathbf{r},s)
-\mu_0 s\int_{\Omega}
\widetilde{\overline{\mathbf{G}}}_{E}(\mathbf{r},\mathbf{r}',s)\cdot
\widetilde{\mathbf{j}}(\mathbf{r}',s)\,d^{3}r',
\end{equation}
with \(\widetilde{\mathbf{j}}\) from Eq.~\eqref{eq:ch2.laplace-source-transform}
\cite{stratton1941}. On \(s=-\jj\omega\) under Eqs.~\eqref{eq:ch2.equivalence-assumptions},
Eq.~\eqref{eq:ch2.laplace-dyadic-convolution} reduces to
Eq.~\eqref{eq:ch2.electric-dyadic-convolution} with
\(\overline{\mathbf{G}}_{E,\omega}\) from Eq.~\eqref{eq:ch2.retarded-laplace-harmonic-link}
and \(\Jimp=\eval{\widetilde{\mathbf{j}}(\mathbf{r},s)}_{s=-\jj\omega}\)
\cite{harrington2001}. The Chapter~\ref{ch:02} causal propagation theorem at kernel
level is
\begin{equation}
\label{eq:ch2.retarded-harmonic-kernel-theorem}
\mathcal{S}_\omega[\Jimp]
=
\eval{\mathcal{S}_t[\mathbf{j}_{\mathrm{imp}}](\mathbf{r},t)}_{\text{harmonic at }\omega}
\quad\Longleftrightarrow\quad
\overline{\mathbf{G}}_{E,\omega}
=
\eval{\widetilde{\overline{\mathbf{G}}}_{E}}_{s=-\jj\omega}
\end{equation}
for the same \(\mathcal{B}_\Gamma\) and outgoing branch, with \(\mathcal{S}_t\) the
time-domain map Eq.~\eqref{eq:ch2.time-domain-solution-map} \cite{stratton1941}.
Equation~\eqref{eq:ch2.retarded-harmonic-kernel-theorem} does not replace
Eq.~\eqref{eq:ch2.harmonic-equivalence-theorem}; it identifies the Green operators
whose resolvent relation Eq.~\eqref{eq:ch2.resolvent-green-equivalence} implements
\cite{harrington2001,jackson1999}.

\paragraph{Advanced kernels and branch failure.}
Advanced dyads satisfying Eq.~\eqref{eq:ch2.advanced-kernel-violation} produce
Laplace symbols that cannot equal
Eq.~\eqref{eq:ch2.retarded-laplace-harmonic-link} for any causal
\(\mathbf{j}_{\mathrm{imp}}\) \cite{stratton1941}. Time-harmonic solvers that
implicitly use incoming Green functions therefore fail
Eq.~\eqref{eq:ch2.retarded-harmonic-kernel-theorem} even when interior Maxwell rows
vanish \cite{harrington2001}. Exterior problems require outgoing closure consistent
with Section~\ref{sec:ch24}; interior cavity problems may require loss or quotienting
as in Eq.~\eqref{eq:ch2.coercivity-quotient} before a causal Laplace contour can be
closed \cite{stratton1941}.

\paragraph{Broadband synthesis and dispersion.}
For finite-bandwidth excitation, the harmonic slice at fixed \(\omega\) is insufficient.
Causal fields superpose as
\begin{equation}
\label{eq:ch2.broadband-retarded-synthesis}
\E(\mathbf{r},t)
=
\int_{-\infty}^{\infty}
\widetilde{\E}(\mathbf{r},-\jj\omega)\,e^{-\jj\omega t}\,\frac{d\omega}{2\pi},
\end{equation}
with \(\widetilde{\E}(\mathbf{r},-\jj\omega)\) the Fourier transform induced by
Eq.~\eqref{eq:ch2.laplace-dyadic-convolution} on causal data
\cite{stratton1941,jackson1999}. Equation~\eqref{eq:ch2.broadband-retarded-synthesis}
is the inverse of the single-frequency reduction used throughout Volume~I; dispersive
\(\varepsilon(\omega)\), \(\mu(\omega)\) modify
\(\widetilde{\overline{\mathbf{G}}}_{E}(\mathbf{r},\mathbf{r}',s)\) and are deferred
beyond the locked single-\(\omega\) harmonic band except where loss enters as
\(\Im\{k_0\}\neq 0\) \cite{harrington2001}. Within the harmonic band,
Eq.~\eqref{eq:ch2.steady-state-harmonic-limit} describes the long-time limit after
switch-on when outgoing radiation or damping removes transients
\cite{stratton1941}.

\paragraph{Radiation inheritance.}
Because \(\Pi_{\mathrm{rad}}\) commutes with harmonic extraction under the ordering
assumed in Subsection~\ref{sec:ch223},
\begin{equation}
\label{eq:ch2.retarded-radiation-kernel-composition}
\mathcal{R}_t[\mathbf{j}_{\mathrm{imp}}]
=
\Pi_{\mathrm{rad}}\circ\mathcal{S}_t[\mathbf{j}_{\mathrm{imp}}],
\qquad
\mathcal{R}_\omega[\Jimp]
=
\Pi_{\mathrm{rad}}\circ\mathcal{S}_\omega[\Jimp],
\end{equation}
with \(\mathcal{R}_\omega\) from Eq.~\eqref{eq:ch2.radiation-time-frequency-equiv}
\cite{stratton1941,harrington2001}. Causality does not promote evanescent content to
radiative export: Eq.~\eqref{eq:ch2.regime-operator-screen} remains the transport gate
after time-domain evolution \cite{stratton1941}. Retarded propagation ensures support
constraints only; spectral and regime screens from Subsection~\ref{sec:ch224} are
additional \cite{harrington2001}.

\paragraph{Boundary of validity.}
Equations~\eqref{eq:ch2.retarded-dyadic-convolution}--\eqref{eq:ch2.retarded-harmonic-kernel-theorem}
require Eqs.~\eqref{eq:ch2.equivalence-assumptions}, outgoing or passive closure on
exterior domains, and \(\mathbf{j}_{\mathrm{imp}}\) in the causal class of
Subsection~\ref{sec:ch221}. They fail for time-varying \(\mathcal{B}_\Gamma(t)\),
moving media, and nonlinear sources. Coincidence and self-term evaluation of
\(\overline{\mathbf{G}}_{E,\omega}\) remains governed by Subsection~\ref{sec:ch232}
even when Eq.~\eqref{eq:ch2.retarded-dyadic-convolution} is discretized in time
\cite{stratton1941,harrington2001}.

\paragraph{Worked illustration (step switch-on).}
Let \(\mathbf{j}_{\mathrm{imp}}(\mathbf{r},t)=\mathbf{j}_0(\mathbf{r})H(t)\) with
\(\mathbf{j}_0\in\mathcal{J}_\Omega\). Equation~\eqref{eq:ch2.retarded-dyadic-convolution}
produces a transient \(\E(\mathbf{r},t)\) vanishing for
\(t<t_{\mathrm{ret}}(\mathbf{r})\) from Eq.~\eqref{eq:ch2.light-cone-arrival-time}
and approaching Eq.~\eqref{eq:ch2.steady-state-harmonic-limit} when
\(\mathbf{j}_0(\mathbf{r})=\Re\{\Jimp(\mathbf{r})e^{-\jj\omega t}\}\) is not assumed;
for single-frequency phasor drive after long times, the harmonic field is
\(\mathcal{S}_\omega[\Jimp]\) from Eq.~\eqref{eq:ch2.retarded-harmonic-kernel-theorem}
\cite{stratton1941}. Evaluating Eq.~\eqref{eq:ch2.electric-dyadic-convolution} at
finite \(t\) before transient decay returns the steady phasor only if
\(t\gg t_{\mathrm{transient}}\), not because causality is optional
\cite{harrington2001}.

\paragraph{Worked illustration (delta-pulse excitation).}
A short pulse \(\mathbf{j}_{\mathrm{imp}}(\mathbf{r},t)=\mathbf{j}_0(\mathbf{r})\,\delta(t)\)
injects broadband content. Eq.~\eqref{eq:ch2.broadband-retarded-synthesis} rebuilds
\(\E(\mathbf{r},t)\) from the full \(\omega\) spectrum; a single
\(\mathcal{S}_\omega[\Jimp]\) slice cannot represent the pulse without superposition
as noted in Subsection~\ref{sec:ch223} \cite{stratton1941,jackson1999}. The retarded
support Eq.~\eqref{eq:ch2.retarded-dyadic-support} still confines influences to
\(t\ge R/c_0\), independent of bandwidth \cite{stratton1941}.

\paragraph{Closing summary.}
Subsection~\ref{sec:ch233} establishes the Chapter~\ref{ch:02} HOME for retarded scalar
wave operator Eq.~\eqref{eq:ch2.retarded-scalar-wave-operator}, causal scalar kernel
Eq.~\eqref{eq:ch2.retarded-scalar-green-time}, dyadic convolution
Eqs.~\eqref{eq:ch2.retarded-dyadic-convolution}--\eqref{eq:ch2.retarded-dyadic-support},
light-cone arrival Eq.~\eqref{eq:ch2.light-cone-arrival-time}, Laplace link
Eqs.~\eqref{eq:ch2.retarded-laplace-harmonic-link}--\eqref{eq:ch2.laplace-dyadic-convolution},
kernel theorem Eq.~\eqref{eq:ch2.retarded-harmonic-kernel-theorem}, broadband synthesis
Eq.~\eqref{eq:ch2.broadband-retarded-synthesis}, and radiation composition
Eq.~\eqref{eq:ch2.retarded-radiation-kernel-composition}. Subsection~\ref{sec:ch234}
addresses constructive inversion and solvability on composite domains.

\subsection{Inversion and solvability}
\label{sec:ch234}
% TOC tag: [HOME]

Subsections~\ref{sec:ch231}--\ref{sec:ch233} defined harmonic and retarded Green
kernels, their regularized convolutions, and causal Laplace links. Subsection~\ref{sec:ch234}
gives the Chapter~\ref{ch:02} HOME for constructive inversion and solvability on
single- and multi-region domains with interface data from
Eq.~\eqref{eq:ch1.interface-conditions} \cite{stratton1941,harrington2001}. Weak
well-posedness on bounded quotients through
Eqs.~\eqref{eq:ch2.lax-milgram-wellposed} and~\eqref{eq:ch2.coercivity-quotient}
from Subsection~\ref{sec:ch214} is presupposed; Lax--Milgram and Friedrichs arguments
are not repeated. Outgoing uniqueness on exterior domains is finalized in
Section~\ref{sec:ch243}; here the resolvent branch and its Green realization are
stated as hypotheses \cite{stratton1941}. Material Helmholtz symbols
Eqs.~\eqref{eq:ch2.material-helmholtz-electric}--\eqref{eq:ch2.material-helmholtz-magnetic}
and trace spaces Eqs.~\eqref{eq:ch2.interface-trace-space}--\eqref{eq:ch2.domain-refinement-L}
from Sections~\ref{sec:ch211}--\ref{sec:ch222} are fixed input classes.

\paragraph{Assumptions.}
Let \(\Omega=\cup_{j=1}^{N_p}\Omega_j\) be a finite union of Lipschitz patches with
pairwise interfaces \(\Sigma_{ab}=\overline{\Omega_a}\cap\overline{\Omega_b}\),
piecewise constants \(\varepsilon_j,\mu_j>0\) on \(\Omega_j\), and outer boundary
\(\Gamma=\partial\Omega\). Impressed data satisfy
\(\Jimp\in\mathcal{J}_\Omega\), charge continuity
Eq.~\eqref{eq:ch1.charge-continuity-harmonic}, and boundary class
\(\mathcal{B}_\Gamma\) compatible with Eq.~\eqref{eq:ch2.admissible-field-space}.
Regularized convolutions follow Subsection~\ref{sec:ch232}; retarded causality follows
Subsection~\ref{sec:ch233} \cite{stratton1941,harrington2001}.

\paragraph{Composite domain and patchwise inversion.}
Decompose the solution map as
\begin{equation}
\label{eq:ch2.composite-domain-decomposition}
\Omega
=
\bigsqcup_{j=1}^{N_p}\Omega_j
\ \cup\
\bigsqcup_{a<b}\Sigma_{ab},
\qquad
k_j=\omega\sqrt{\mu_j\varepsilon_j},
\end{equation}
with wavenumbers \(k_j\) from Eq.~\eqref{eq:ch2.helmholtz-symbol} on each patch
\cite{stratton1941}. On \(\Omega_j\), the electric field satisfies the patchwise
Helmholtz drive of Eq.~\eqref{eq:ch2.material-helmholtz-electric} and is represented
by the regularized dyadic convolution
\begin{equation}
\label{eq:ch2.patchwise-dyadic-inversion}
\E(\mathbf{r})
=
\E_{\mathrm{hom}}(\mathbf{r})
+\jj\omega\mu_j
\int_{\Omega_j}
\overline{\mathbf{G}}_{E}^{(j)}(\mathbf{r},\mathbf{r}')\cdot\Jimp(\mathbf{r}')\,d^{3}r',
\qquad
\mathbf{r}\in\Omega_j,
\end{equation}
where \(\overline{\mathbf{G}}_{E}^{(j)}\) satisfies
Eq.~\eqref{eq:ch2.material-dyadic-defining-operator} with patch coefficients and
finite-part evaluation when \(\mathbf{r}\in\Omega_j\) approaches
\(\operatorname{supp}\Jimp\) \cite{harrington2001}. Equation
\eqref{eq:ch2.patchwise-dyadic-inversion} is the constructive content of
Eq.~\eqref{eq:ch2.solution-map-dyadic-realization} on uniform subdomains; it does
not yet enforce interface jumps across \(\Sigma_{ab}\) \cite{stratton1941}.

\paragraph{Interface coupling.}
Fields in adjacent patches are linked by tangential continuity and normal-flux jumps
imported from Eq.~\eqref{eq:ch1.interface-conditions} \cite{stratton1941}. In
operator form,
\begin{equation}
\label{eq:ch2.interface-dyadic-continuity}
\hat{\mathbf{n}}_{ab}\times\E^{(a)}\big|_{\Sigma_{ab}}
=
\hat{\mathbf{n}}_{ab}\times\E^{(b)}\big|_{\Sigma_{ab}},
\qquad
\hat{\mathbf{n}}_{ab}\times\H^{(a)}\big|_{\Sigma_{ab}}
=
\hat{\mathbf{n}}_{ab}\times\H^{(b)}\big|_{\Sigma_{ab}},
\end{equation}
with \(\E^{(j)}\) from Eq.~\eqref{eq:ch2.patchwise-dyadic-inversion} and
\(\H^{(j)}\) from Eq.~\eqref{eq:ch2.magnetic-from-electric-dyad} when valid
\cite{harrington2001}. Equations~\eqref{eq:ch2.interface-dyadic-continuity} replace
free-space Eq.~\eqref{eq:ch2.free-space-dyadic-green} by patch Green functions
\(\overline{\mathbf{G}}_{E}^{(j)}\) plus interface unknowns on
\(\mathcal{T}_{\Sigma_{ab}}\) from Eq.~\eqref{eq:ch2.interface-trace-space}
\cite{stratton1941}. The assembled multi-region system is
\begin{equation}
\label{eq:ch2.multi-region-inversion-system}
\begin{bmatrix}
\mathcal{I}+\mathcal{K}_{11} & \cdots & \mathcal{K}_{1N_p}\\
\vdots & \ddots & \vdots\\
\mathcal{K}_{N_p1} & \cdots & \mathcal{I}+\mathcal{K}_{N_pN_p}
\end{bmatrix}
\begin{bmatrix}\mathbf{x}_1\\ \vdots\\ \mathbf{x}_{N_p}\end{bmatrix}
=
\begin{bmatrix}\mathbf{b}_1\\ \vdots\\ \mathbf{b}_{N_p}\end{bmatrix},
\end{equation}
where \(\mathbf{x}_j\) collects tangential field or equivalent surface currents on
patch boundaries, \(\mathcal{K}_{jk}\) are dyadic layer potentials generated by
\(\overline{\mathbf{G}}_{E}^{(k)}\), and \(\mathbf{b}_j\) encodes impressed sources
and outer \(\mathcal{B}_\Gamma\) data \cite{harrington2001,stratton1941}. Equation
\eqref{eq:ch2.multi-region-inversion-system} is the Chapter~\ref{ch:02} HOME
statement of multi-region solvability: invertibility of the block operator is the
Green-function analogue of Eq.~\eqref{eq:ch2.lax-milgram-wellposed} on composite
domains \cite{stratton1941}.

\begin{figure}[t]
\centering
\ChOneFigBlock{%
\begin{tikzpicture}[ch1fig, node distance=7mm and 9mm]
  \node[ch1box] (p1) {$\Omega_1$\\$k_1$};
  \node[ch1box, right=14mm of p1] (sig) {$\Sigma_{12}$};
  \node[ch1box, right=14mm of sig] (p2) {$\Omega_2$\\$k_2$};
  \node[ch1box, below=11mm of sig] (g) {$\overline{\mathbf{G}}_{E}^{(j)}$};
  \node[ch1box, below=11mm of p2] (bc) {$\mathcal{B}_\Gamma$};
  \draw[ch1flow] (p1) -- (sig) -- (p2);
  \draw[ch1flow] (g) -- (p1);
  \draw[ch1flow] (g) -- (p2);
  \draw[ch1flow] (bc) -- (p2);
\end{tikzpicture}%
}
\caption{Multi-region inversion Eq.~\eqref{eq:ch2.multi-region-inversion-system}:
patchwise dyadic kernels coupled by interface continuity
Eq.~\eqref{eq:ch2.interface-dyadic-continuity} and exterior
\(\mathcal{B}_\Gamma\).}
\label{fig:ch2.multi-region-green-assembly}
\end{figure}

Figure~\ref{fig:ch2.multi-region-green-assembly} is the constructive template for
boundary-element and layered-medium solvers. Subsection~\ref{sec:ch232} supplies
self-term regularization on each patch; Subsection~\ref{sec:ch233} supplies causal
time marching when Eq.~\eqref{eq:ch2.multi-region-inversion-system} is embedded in
\(\mathcal{S}_t\) \cite{harrington2001}.

\paragraph{Existence on bounded domains.}
On a bounded Lipschitz domain with passive \(\mathcal{B}_\Gamma\) satisfying
Eq.~\eqref{eq:ch2.passive-impedance-class}, coercivity on the quotient
\(\mathcal{U}_\omega^{0}/\ker\mathcal{L}_\omega\) from
Eq.~\eqref{eq:ch2.coercivity-quotient} implies existence of a unique weak solution
Eq.~\eqref{eq:ch2.lax-milgram-wellposed} \cite{stratton1941}. Green-function
inversion is equivalent when the regularized operator
\begin{equation}
\label{eq:ch2.green-bvp-equivalence}
\mathcal{S}_\omega[\Jimp]
=
\mathbf{u}_\omega
\quad\Longleftrightarrow\quad
\mathcal{L}_\omega[\mathbf{u}_\omega]
=
\begin{bmatrix}\mathbf{s}_\omega(\Jimp,\rho)\\ \mathbf{g}_\Gamma\end{bmatrix},
\end{equation}
with \(\mathbf{u}_\omega=(\E,\H)\) obtained from
Eqs.~\eqref{eq:ch2.patchwise-dyadic-inversion}--\eqref{eq:ch2.mom-self-term-regularized}
on simply connected patches and from
Eq.~\eqref{eq:ch2.multi-region-inversion-system} otherwise
\cite{harrington2001,stratton1941}. Equation~\eqref{eq:ch2.green-bvp-equivalence}
restates Eq.~\eqref{eq:ch2.solution-map-resolvent} with constructive kernels; stability
is controlled by Eq.~\eqref{eq:ch2.solution-map-stability} and resolvent bound
Eq.~\eqref{eq:ch2.resolvent-bounded} on the physical branch
\cite{stratton1941}.

\paragraph{Exterior domains and branch selection.}
On exterior \(\Omega\) with compact \(\operatorname{supp}\Jimp\), existence of
\(\mathcal{S}_\omega\) requires outgoing Green functions whose far zone matches
Section~\ref{sec:ch24} \cite{stratton1941}. Without that selection, the block system
Eq.~\eqref{eq:ch2.multi-region-inversion-system} admits incoming solutions violating
Eq.~\eqref{eq:ch2.retarded-harmonic-kernel-theorem} \cite{harrington2001}. Uniqueness
on the outgoing branch is
\begin{equation}
\label{eq:ch2.exterior-uniqueness-hypothesis}
\mathcal{L}_\omega[\mathbf{u}_\omega]=\mathbf{0}
\ \text{and outgoing}\ \mathcal{B}_\Gamma
\ \Longrightarrow\
\mathbf{u}_\omega=\mathbf{0},
\end{equation}
proved in Section~\ref{sec:ch243}; Subsection~\ref{sec:ch234} invokes
Eq.~\eqref{eq:ch2.exterior-uniqueness-hypothesis} as the solvability gate for
exterior Green inversion \cite{stratton1941}. The split
Eq.~\eqref{eq:ch2.outgoing-nonselfadjoint-split} reminds that exterior problems are
not Friedrichs self-adjoint on all of \(\mathcal{U}_\omega\)
\cite{harrington2001}.

\paragraph{Resonance and failure modes.}
Bounded cavities with PEC walls possess discrete spectra \(\{\omega_n\}\). At
\(\omega=\omega_n\),
\begin{equation}
\label{eq:ch2.cavity-resonance-failure}
\ker\mathcal{L}_\omega\neq\{\mathbf{0}\}
\quad\Longrightarrow\quad
\mathcal{L}_\omega^{-1}\ \text{unbounded on}\ \mathcal{U}_\omega,
\end{equation}
and Eq.~\eqref{eq:ch2.green-bvp-equivalence} fails unless quotienting by standing
modes as in Eq.~\eqref{eq:ch2.coercivity-quotient} or introducing loss
\cite{stratton1941,harrington2001}. Poles of the resolvent appear as
\begin{equation}
\label{eq:ch2.resolvent-pole-cavity}
\det\big(\mathcal{I}+\mathcal{K}\big)=0
\quad\Longleftrightarrow\quad
\omega=\omega_n,
\end{equation}
for the coupled operator \(\mathcal{K}\) in
Eq.~\eqref{eq:ch2.multi-region-inversion-system} on closed resonant structures
\cite{stratton1941}. Impressed \(\Jimp\) at resonance requires compatibility with
\(\ker\mathcal{L}_\omega\); otherwise amplitudes diverge even when Green formulas
are formally applied \cite{harrington2001}. Incoming exterior excitation produces a
distinct failure: Eq.~\eqref{eq:ch2.exterior-uniqueness-hypothesis} is violated
without additional incident-field subtraction in Section~\ref{sec:ch251}
\cite{stratton1941}.

\paragraph{Numerical equivalence.}
Finite-element stiffness systems and boundary-element impedance matrices are
Galerkin discretizations of Eq.~\eqref{eq:ch2.maxwell-weak-form} and
Eq.~\eqref{eq:ch2.multi-region-inversion-system}, respectively:
\begin{equation}
\label{eq:ch2.discrete-green-equivalence}
\mathbf{Z}\mathbf{I}=\mathbf{V}
\quad\Longleftrightarrow\quad
\mathcal{K}\mathbf{x}=\mathbf{b},
\end{equation}
when basis functions match trace spaces and self-terms implement
Eq.~\eqref{eq:ch2.mom-self-term-regularized} \cite{harrington2001,stratton1941}.
Equation~\eqref{eq:ch2.discrete-green-equivalence} is not a separate physics; it is
the matrix form of the same inversion problem. Ill-conditioning near
Eq.~\eqref{eq:ch2.resolvent-pole-cavity} mirrors practical difficulty at cavity
modes and near-spurious resonances in open-region formulations
\cite{harrington2001}.

\paragraph{Boundary of validity.}
Equations~\eqref{eq:ch2.patchwise-dyadic-inversion}--\eqref{eq:ch2.multi-region-inversion-system}
require Lipschitz interfaces, passive linear media on each patch, and fixed
\(\omega\) in the harmonic band. They fail for time-varying boundaries, nonlinear
constitutive laws, and anisotropic magneto-electric coupling without retensoring
\(\overline{\mathbf{G}}_{E}^{(j)}\) \cite{stratton1941}. Point sources without
Subsection~\ref{sec:ch232} regularization produce formal but non-admissible fields;
\(\Jimp\) must remain in \(\mathcal{J}_\Omega\) \cite{harrington2001}.

\paragraph{Worked illustration (coated sphere).}
A spherical core \(\Omega_1\) and coating \(\Omega_2\) with distinct \((\varepsilon,\mu)\)
assemble Eq.~\eqref{eq:ch2.multi-region-inversion-system} with a single interface
\(\Sigma_{12}\). Continuity Eq.~\eqref{eq:ch2.interface-dyadic-continuity} replaces
free-space dyads in the exterior region; Mie coefficients computed in
Section~\ref{sec:ch33} are the modal diagonalization of the same block system
\cite{stratton1941,harrington2001}. Subsection~\ref{sec:ch234} does not compute
Mie modes; it identifies the operator that modal analysis diagonalizes.

\paragraph{Worked illustration (PEC cavity at resonance).}
When \(\omega\) matches a cavity eigenfrequency, Eq.~\eqref{eq:ch2.cavity-resonance-failure}
blocks inversion. A finite-Q loss layer or quotienting
Eq.~\eqref{eq:ch2.coercivity-quotient} restores bounded
\(\mathcal{L}_\omega^{-1}\); applying Eq.~\eqref{eq:ch2.free-space-dyadic-green}
without boundary-compliant kernels on a closed PEC box returns a finite matrix solve
that is not Eq.~\eqref{eq:ch2.green-bvp-equivalence} for the closed problem
\cite{stratton1941,harrington2001}.

\paragraph{Closing summary.}
Subsection~\ref{sec:ch234} establishes the Chapter~\ref{ch:02} HOME for composite
decomposition Eq.~\eqref{eq:ch2.composite-domain-decomposition}, patch inversion
Eq.~\eqref{eq:ch2.patchwise-dyadic-inversion}, interface coupling
Eq.~\eqref{eq:ch2.interface-dyadic-continuity}, multi-region system
Eq.~\eqref{eq:ch2.multi-region-inversion-system}, Green--BVP equivalence
Eq.~\eqref{eq:ch2.green-bvp-equivalence}, exterior uniqueness hypothesis
Eq.~\eqref{eq:ch2.exterior-uniqueness-hypothesis}, resonance failure
Eqs.~\eqref{eq:ch2.cavity-resonance-failure}--\eqref{eq:ch2.resolvent-pole-cavity},
and discrete equivalence Eq.~\eqref{eq:ch2.discrete-green-equivalence}.

With Subsections~\ref{sec:ch231}--\ref{sec:ch234} complete, Section~\ref{sec:ch23}
has constructed scalar, vector, and dyadic Green kernels, regularized their
singular convolutions, tied them to retarded causal propagation, and stated
constructive inversion on composite domains compatible with
Eq.~\eqref{eq:ch2.resolvent-operator}. Sommerfeld outgoing selection in
Section~\ref{sec:ch24} fixes the exterior branch assumed in
Eq.~\eqref{eq:ch2.exterior-uniqueness-hypothesis}; radiation integrals in
Section~\ref{sec:ch25} evaluate far-zone observables from the propagating subspace
qualified by Eq.~\eqref{eq:ch2.regime-operator-screen}; spectral and Fredholm
consequences of Eq.~\eqref{eq:ch2.green-compact-split} continue in
Section~\ref{sec:ch263} \cite{stratton1941,harrington2001,jackson1999}.

\section{Radiation Conditions and Uniqueness}
\label{sec:ch24}

Section~\ref{sec:ch23} constructed harmonic and retarded Green kernels, regularized
their convolutions, and stated multi-region inversion subject to a branch hypothesis
Eq.~\eqref{eq:ch2.exterior-uniqueness-hypothesis}. That hypothesis is not a Maxwell
identity: exterior problems admit incoming as well as outgoing square-integrable
solutions unless a radiation condition selects the physical resolvent branch
\cite{stratton1941,harrington2001}. Section~\ref{sec:ch24} supplies the Chapter~\ref{ch:02}
closure for that selection---Sommerfeld-type outgoing conditions, operator-level
outgoing-wave gates compatible with retarded Green functions and
Eq.~\eqref{eq:ch2.retarded-harmonic-kernel-theorem}, uniqueness theorems for radiative
solutions, and the consequences for source reconstruction and equivalent-current
inversion \cite{stratton1941,jackson1999}. Outer-boundary taxonomy and local
uniqueness templates from Subsections~\ref{sec:ch01.1.2} and~\ref{sec:ch01.1.3} are
invoked through Eqs.~\eqref{eq:ch1.outer-boundary-classes}--\eqref{eq:ch1.radiation-trace-condition}
and Eqs.~\eqref{eq:ch1.uniqueness-difference-system}--\eqref{eq:ch1.uniqueness-estimate}
without re-opening those proofs \cite{stratton1941}. Local Maxwell derivations,
gauge structure, and flux integral definitions remain in Chapter~\ref{ch:01}.

The problem addressed here is branch degeneracy. On a bounded domain with passive
\(\mathcal{B}_\Gamma\), Eq.~\eqref{eq:ch2.lax-milgram-wellposed} and cavity quotients
Eq.~\eqref{eq:ch2.coercivity-quotient} from Subsection~\ref{sec:ch214} fix
\(\mathcal{S}_\omega\) up to discrete resonances Eq.~\eqref{eq:ch2.cavity-resonance-failure}.
On an exterior domain with compact sources, infinitely many harmonic fields satisfy
\(\mathcal{M}_\omega[\mathbf{u}_\omega]=\mathbf{0}\) in the volume unless far-zone
decay and phase are constrained \cite{stratton1941,harrington2001}. Free-space kernels
Eqs.~\eqref{eq:ch2.free-space-scalar-green} and~\eqref{eq:ch2.free-space-dyadic-green}
therefore require outgoing interpretation: the same formula can represent incoming or
outgoing waves depending on branch and sign convention on \(k_0\)
\cite{jackson1999}. Causality from Subsection~\ref{sec:ch233} narrows the class to
retarded kernels but does not, by itself, replace Sommerfeld selection when harmonic
solvers invert \(\mathcal{L}_\omega\) directly at fixed \(\omega\)
\cite{stratton1941}. Section~\ref{sec:ch24} makes the outgoing subspace explicit as
the domain on which \(\Pi_{\mathrm{rad}}\) from
Eq.~\eqref{eq:ch2.radiation-operator-def} is a projector rather than a formal
post-filter \cite{harrington2001}.

\begin{figure}[t]
\centering
\ChOneFigBlock{%
\begin{tikzpicture}[ch1fig, node distance=7mm and 10mm]
  \node[ch1box] (sf) {Sommerfeld\\far-zone};
  \node[ch1box, right=14mm of sf] (out) {Outgoing\\branch};
  \node[ch1box, right=14mm of out] (uniq) {Uniqueness\\$\mathcal{S}_\omega$};
  \node[ch1box, right=14mm of uniq] (src) {Source\\reconstruction};
  \node[ch1box, below=12mm of out] (gr) {$\overline{\mathbf{G}}_{E,\omega}$};
  \draw[ch1flow] (sf) -- (out) -- (uniq) -- (src);
  \draw[ch1flow] (gr) -- (out);
\end{tikzpicture}%
}
\caption{Section~\ref{sec:ch24} development path: Sommerfeld far-zone conditions
select the outgoing Green branch used in Section~\ref{sec:ch23}; uniqueness fixes
\(\mathcal{S}_\omega\); Subsection~\ref{sec:ch244} addresses reconstruction
identifiability.}
\label{fig:ch2.radiation-uniqueness-roadmap}
\end{figure}

Figure~\ref{fig:ch2.radiation-uniqueness-roadmap} orders the four subsections.
Subsection~\ref{sec:ch241} gives the Chapter~\ref{ch:02} HOME for the Sommerfeld
radiation condition and its link to Eqs.~\eqref{eq:ch1.outer-boundary-classes} and
\(\Pi_{\mathrm{prop}}\) from Subsection~\ref{sec:ch224} \cite{stratton1941}. Subsection
~\ref{sec:ch242} packages outgoing-wave selection as an operator gate on
\(\mathcal{F}_\Omega\), aligning Green-function inversion with the retarded/resolvent
branch of Eq.~\eqref{eq:ch2.resolvent-green-equivalence}
\cite{harrington2001}. Subsection~\ref{sec:ch243} proves uniqueness of radiative
solutions on the outgoing subspace, upgrading the preview used in
Eq.~\eqref{eq:ch2.exterior-uniqueness-hypothesis} and importing the energy method of
Subsection~\ref{sec:ch01.1.3} by reference \cite{stratton1941}. Subsection~\ref{sec:ch244}
records implications for inverse source problems, null-space structure of
\(\mathcal{R}_\omega\), and the limits of reconstruction without regime qualification
Eq.~\eqref{eq:ch2.regime-operator-screen} \cite{harrington2001}.

Several constructions deliberately remain outside Section~\ref{sec:ch24}. Far-zone
integral representations and flux meters belong to Section~\ref{sec:ch25}; vector
spherical harmonic separation of outgoing fields to Chapter~\ref{ch:03}; complete
Fredholm index computation to Section~\ref{sec:ch263} \cite{stratton1941}. Section~\ref{sec:ch24}
does not re-derive Green kernels or regularization from Section~\ref{sec:ch23}; it
selects which constructed branch is physically admissible in open regions
\cite{harrington2001}. Regime logic from Section~\ref{sec:ch01.2.4} is unchanged:
uniqueness of \(\mathcal{S}_\omega\) does not promote evanescent near fields to
transport channels; Eq.~\eqref{eq:ch2.regime-operator-screen} remains the measurement
gate after outgoing closure \cite{stratton1941}.

The split Eq.~\eqref{eq:ch2.outgoing-nonselfadjoint-split} from Subsection~\ref{sec:ch214}
warns against importing cavity self-adjointness into exterior problems. Outgoing
selection defines a physical subspace \(\mathcal{U}_{\mathrm{out}}\) on which
\(\mathcal{L}_\omega\) is injective; reactive content in \(\mathcal{U}_{\mathrm{re}}\)
may be large even when Eq.~\eqref{eq:ch2.exterior-uniqueness-hypothesis} holds for the
total outgoing component \cite{harrington2001,stratton1941}. Radiation operators
inherit idempotency of \(\Pi_{\mathrm{rad}}\) only after Subsection~\ref{sec:ch243};
until then, post hoc far-zone windowing of \(\mathcal{S}_\omega[\Jimp]\) risks
confusing stored energy with radiative export, the failure mode flagged in
Section~\ref{sec:ch01.2.4} \cite{stratton1941}.

Readers approaching from numerical electromagnetics should map Section~\ref{sec:ch24}
to the boundary artifacts they already use. Absorbing boundary conditions, perfectly
matched layers, and radiation boundary integrals are discrete surrogates for
Sommerfeld outgoing selection \cite{harrington2001}. A frequency-domain finite-element
solve that omits such closure may return a finite \(\mathbf{u}_\omega\) that is not
the retarded branch of Eq.~\eqref{eq:ch2.retarded-harmonic-kernel-theorem}; the
mismatch is branch error, not mesh error alone \cite{stratton1941}. Boundary-element
exterior formulations embed outgoing conditions in the Green function choice; interior
cavity formulations require quotienting Eq.~\eqref{eq:ch2.coercivity-quotient} instead
\cite{harrington2001}. Section~\ref{sec:ch24} states the continuous targets those
schemes approximate.

Incoming versus outgoing distinctions interact with propagation subspaces from
Subsection~\ref{sec:ch224}. Only \(\Pi_{\mathrm{prop}}\) content carries far-zone
power after outgoing closure; evanescent and non-propagating components may dominate
near measurements while satisfying the same Sommerfeld condition at infinity
\cite{stratton1941,harrington2001}. Section~\ref{sec:ch24} therefore closes
uniqueness of the \emph{radiative} exterior solution, not uniqueness of reactive
near-field parameterizations. Transport reporting waits for Section~\ref{sec:ch25} on
the qualified propagating subspace \cite{stratton1941}.

Source reconstruction and equivalent-current inversion previewed in
Section~\ref{sec:ch244} depend on both uniqueness and observability. Uniqueness of
\(\mathcal{S}_\omega\) on the outgoing branch does not imply uniqueness of
\(\Jimp\) given far-zone samples: \(\ker\mathcal{R}_\omega\) is generally
infinite-dimensional, as recorded in Eq.~\eqref{eq:ch2.observation-null-space}
from Subsection~\ref{sec:ch213} \cite{harrington2001,stratton1941}. Subsection~\ref{sec:ch244}
separates \emph{field} uniqueness from \emph{source} identifiability and connects
the latter to measurement surrogates in Chapter~\ref{ch:01} without re-opening
angular-momentum observables \cite{stratton1941}.

Loss and dissipation modify uniqueness without changing the outgoing template.
Passive impedance Eq.~\eqref{eq:ch2.passive-impedance-class} and complex \(k_0\) shift
resolvent poles off the real axis, suppressing the sharp divergence of
Eq.~\eqref{eq:ch2.cavity-resonance-failure} while preserving the need for explicit
outgoing selection in open regions \cite{stratton1941,harrington2001}. Section~\ref{sec:ch24}
treats loss as a regularization of the spectrum, not as a substitute for
Eq.~\eqref{eq:ch1.radiation-trace-condition} on exterior boundaries
\cite{stratton1941}.

The subsection sequence begins with Sommerfeld far-zone conditions, continues with
outgoing operator selection on Green kernels from Section~\ref{sec:ch23}, proves
uniqueness on the radiative branch, and closes with reconstruction consequences before
radiation integrals in Section~\ref{sec:ch25} evaluate flux and patterns on the
propagating subspace \cite{stratton1941,harrington2001,jackson1999}.


\subsection{Sommerfeld radiation condition}
\label{sec:ch241}
% TOC tag: [HOME][USE SS1.1.2]

Section~\ref{sec:ch24} identified outgoing branch selection as the missing closure
for exterior Green inversion Eq.~\eqref{eq:ch2.exterior-uniqueness-hypothesis}.
Subsection~\ref{sec:ch241} gives the Chapter~\ref{ch:02} HOME for the Sommerfeld
radiation condition and its operator embedding on \(\mathcal{F}_\Omega\)
\cite{stratton1941,harrington2001,jackson1999}. Outer-boundary taxonomy and the
outgoing-trace preview Eq.~\eqref{eq:ch1.radiation-trace-condition} from
Subsection~\ref{sec:ch01.1.2} are presupposed through
Eqs.~\eqref{eq:ch1.outer-boundary-classes} and~\eqref{eq:ch1.boundary-data-operator};
those boundary classes are not re-derived. Free-space kernels
Eqs.~\eqref{eq:ch2.free-space-scalar-green} and~\eqref{eq:ch2.free-space-dyadic-green}
from Subsection~\ref{sec:ch231} are cited as asymptotic templates, not re-displayed.
Uniqueness proofs and reconstruction limits belong to Subsections~\ref{sec:ch243}
and~\ref{sec:ch244}.

\paragraph{Assumptions.}
Let \(\Omega_{\mathrm{ext}}\) denote an exterior region surrounding compact
\(\operatorname{supp}\Jimp\subset\Omega_s\), with \(\omega>0\) fixed and vacuum
constants \(\varepsilon_0,\mu_0\) in the far zone. Spherical coordinates
\((r,\theta,\phi)\) are centered so that \(\Omega_s\subset B_R(0)\) for some
\(R_0>0\). Fields are phasors under \(\exp(-\jj\omega t)\); outgoing means
\(e^{-\jj k_0 r}/r\) decay with \(k_0=\omega\sqrt{\mu_0\varepsilon_0}\) from
Eq.~\eqref{eq:ch2.helmholtz-symbol} \cite{stratton1941}. Interior interfaces and
finite domains use Eq.~\eqref{eq:ch1.outer-boundary-classes}; Sommerfeld closure
applies on artificial or physical spheres at large \(r\) \cite{harrington2001}.

\paragraph{Scalar Sommerfeld template.}
For scalar potentials and Helmholtz kernels, outgoing radiation requires
\begin{equation}
\label{eq:ch2.sommerfeld-scalar-asymptotic}
G(\mathbf{r},\mathbf{r}')
=
A(\theta,\phi,\mathbf{r}')\,\frac{e^{-\jj k_0 r}}{r}
\;+\;
\mathcal{O}(r^{-2}),
\qquad
r=|\mathbf{r}|\to\infty,
\end{equation}
with amplitude \(A\) independent of \(r\) at leading order and branch fixed by the
locked phasor axis \cite{stratton1941,jackson1999}. Equation
\eqref{eq:ch2.sommerfeld-scalar-asymptotic} is the Chapter~\ref{ch:02} scalar
Sommerfeld condition: it excludes incoming \(e^{+\jj k_0 r}/r\) and standing
\(e^{\pm\jj k_0 r}/r\) superpositions that degenerate exterior uniqueness
\cite{harrington2001}. The free-space template Eq.~\eqref{eq:ch2.free-space-scalar-green}
satisfies Eq.~\eqref{eq:ch2.sommerfeld-scalar-asymptotic} for fixed \(\mathbf{r}'\)
\cite{stratton1941}. Physically, Eq.~\eqref{eq:ch2.sommerfeld-scalar-asymptotic}
states that monopole-like responses propagate outward with phase fronts
\(k_0 r=\mathrm{const}\) and amplitude decay \(r^{-1}\), the only radial
combination compatible with energy flux to infinity in three dimensions
\cite{jackson1999}.

\paragraph{Vector Sommerfeld condition.}
For harmonic Maxwell fields in vacuum far from sources,
\begin{equation}
\label{eq:ch2.sommerfeld-field-condition}
\mathbf{r}\times\E
\;+\;
\eta_0\,(\mathbf{r}\times\H)\times\hat{\mathbf{r}}
=
\mathbf{o}(1),
\qquad
r\to\infty,
\end{equation}
with \(\eta_0=\sqrt{\mu_0/\varepsilon_0}\) and \(\hat{\mathbf{r}}=\mathbf{r}/r\)
\cite{stratton1941,harrington2001}. Equation~\eqref{eq:ch2.sommerfeld-field-condition}
is the Chapter~\ref{ch:02} vector Sommerfeld radiation condition. It implies
transverse radiative fields with asymptotic ratio
\begin{equation}
\label{eq:ch2.sommerfeld-impedance-ratio}
\H_{\mathrm{rad}}
\sim
\eta_0^{-1}\,\hat{\mathbf{r}}\times\E_{\mathrm{rad}},
\qquad
r\to\infty,
\end{equation}
and matches the outgoing-trace preview Eq.~\eqref{eq:ch1.radiation-trace-condition}
from Subsection~\ref{sec:ch01.1.2} when tangential traces are taken on a sphere of
large radius \cite{stratton1941}. Mathematically,
Eq.~\eqref{eq:ch2.sommerfeld-field-condition} selects the subspace on which the
exterior Laplacian/Helmholtz resolvent coincides with retarded Green functions
Eq.~\eqref{eq:ch2.retarded-harmonic-kernel-theorem}; physically, it enforces
power flow outward rather than inward at infinity \cite{harrington2001,jackson1999}.

\paragraph{Far-zone transversality and propagation.}
Sommerfeld-admissible fields decompose at leading order as
\begin{equation}
\label{eq:ch2.far-zone-transverse-radiation}
\E(\mathbf{r})
=
\E_0(\theta,\phi)\,\frac{e^{-\jj k_0 r}}{r}
\;+\;\mathcal{O}(r^{-2}),
\qquad
\E_0\cdot\hat{\mathbf{r}}=0,
\end{equation}
with \(\E_0\) the radiation amplitude pattern and \(\E_0\cdot\hat{\mathbf{r}}=0\)
the transversality constraint \cite{stratton1941}. Equation
\eqref{eq:ch2.far-zone-transverse-radiation} connects Sommerfeld closure to
\(\Pi_{\mathrm{prop}}\) from Subsection~\ref{sec:ch224}: far-zone content lies on
the propagating sheet with real \(k_r=k_0\) and outgoing phase
\cite{harrington2001}. Evanescent and reactive components may still be large at
finite \(r\) while satisfying Eq.~\eqref{eq:ch2.sommerfeld-field-condition}; regime
qualification Eq.~\eqref{eq:ch2.regime-operator-screen} remains separate from
Sommerfeld selection \cite{stratton1941}.

\begin{figure}[t]
\centering
\ChOneFigBlock{%
\begin{tikzpicture}[ch1fig, node distance=7mm and 10mm]
  \node[ch1box] (src) {Compact\\$\Omega_s$};
  \node[ch1box, right=16mm of src] (near) {Near/reactive\\$kr\lesssim 1$};
  \node[ch1box, right=16mm of near] (far) {Sommerfeld\\$e^{-jk_0 r}/r$};
  \node[ch1box, below=11mm of near] (ev) {$\Pi_{\mathrm{ev}}$};
  \draw[ch1flow] (src) -- (near) -- (far);
  \draw[ch1flow] (near) -- (ev);
\end{tikzpicture}%
}
\caption{Sommerfeld far zone Eq.~\eqref{eq:ch2.sommerfeld-field-condition}: compact
sources drive near fields with evanescent content; at large \(r\) only outgoing
transverse \(e^{-jk_0 r}/r\) waves satisfy exterior closure.}
\label{fig:ch2.sommerfeld-far-zone}
\end{figure}

Figure~\ref{fig:ch2.sommerfeld-far-zone} separates Sommerfeld closure from regime
screening. Subsection~\ref{sec:ch241} fixes the asymptotic \emph{boundary at
infinity}; Subsection~\ref{sec:ch224} and Section~\ref{sec:ch01.2.4} govern when
far-zone samples may be interpreted as transport \cite{stratton1941,harrington2001}.

\paragraph{Green-function asymptotics.}
Dyadic kernels from Subsection~\ref{sec:ch231} inherit scalar Sommerfeld decay.
For \(r=|\mathbf{r}-\mathbf{r}'|\to\infty\) at fixed \(\mathbf{r}'\),
\begin{equation}
\label{eq:ch2.sommerfeld-dyadic-asymptotic}
\overline{\mathbf{G}}_{E,0}(\mathbf{r},\mathbf{r}')
=
\overline{\mathbf{G}}_{\infty}(\theta,\phi,\mathbf{r}')\,
\frac{e^{-\jj k_0 r}}{r}
\;+\;\mathcal{O}(r^{-2}),
\end{equation}
with \(\overline{\mathbf{G}}_{\infty}\) transverse in the \(\mathbf{r}\)-direction
\cite{stratton1941,harrington2001}. Equation~\eqref{eq:ch2.sommerfeld-dyadic-asymptotic}
is the kernel form of Eq.~\eqref{eq:ch2.sommerfeld-field-condition} under
Eq.~\eqref{eq:ch2.electric-dyadic-convolution}: impressed currents that satisfy
Eq.~\eqref{eq:ch1.source-admissibility} generate fields whose far zone is built from
outgoing dyadic radiation \cite{stratton1941}. The leading \(r^{-1}\) term in
Eq.~\eqref{eq:ch2.free-space-dyadic-green} is the explicit vacuum instance of
Eq.~\eqref{eq:ch2.sommerfeld-dyadic-asymptotic} \cite{jackson1999}. Subsection~\ref{sec:ch242}
promotes Eq.~\eqref{eq:ch2.sommerfeld-dyadic-asymptotic} to an operator selector on
\(\mathcal{S}_\omega\); Subsection~\ref{sec:ch243} uses it in uniqueness energy
integrals \cite{harrington2001}.

\paragraph{Outgoing function class.}
Define the Sommerfeld outgoing subspace on an exterior shell \(\Gamma_R=\{r=R\}\):
\begin{equation}
\label{eq:ch2.sommerfeld-outgoing-class}
\mathcal{F}_{\mathrm{SF}}
=
\left\{
(\E,\H)\in\HCurl(\Omega_{\mathrm{ext}}):
\text{Eqs.~\eqref{eq:ch2.sommerfeld-field-condition}--\eqref{eq:ch2.far-zone-transverse-radiation} hold}
\right\}.
\end{equation}
Equation~\eqref{eq:ch2.sommerfeld-outgoing-class} embeds Eq.~\eqref{eq:ch1.radiation-trace-condition}
into the Chapter~\ref{ch:02} function-space framework of
Eq.~\eqref{eq:ch2.admissible-field-space} \cite{stratton1941}. The physical
solution operator \(\mathcal{S}_\omega\) maps \(\mathcal{J}_\Omega\) into
\(\mathcal{F}_{\mathrm{SF}}\) when \(\mathcal{B}_\Gamma\) includes outgoing closure
\cite{harrington2001}. Incoming solutions lie outside
\(\mathcal{F}_{\mathrm{SF}}\) even if they satisfy volume Maxwell rows; that exclusion
is the operator content of Sommerfeld radiation \cite{stratton1941,jackson1999}.

\paragraph{Large-sphere impedance closure.}
On a finite truncation sphere \(\Gamma_R\), Eq.~\eqref{eq:ch2.sommerfeld-impedance-ratio}
motivates the first-order absorbing condition
\begin{equation}
\label{eq:ch2.silver-muller-trace-approx}
\hat{\mathbf{n}}\times\H
\;-\;
\eta_0^{-1}(\hat{\mathbf{n}}\times\E)\times\hat{\mathbf{n}}
=
\mathbf{0}\quad\text{on }\Gamma_R,
\end{equation}
which approximates Eq.~\eqref{eq:ch1.radiation-trace-condition} at finite \(R\)
\cite{stratton1941,harrington2001}. Equation~\eqref{eq:ch2.silver-muller-trace-approx}
is not identical to Sommerfeld closure at infinity: spurious reflections from
\(\Gamma_R\) contaminate near fields unless \(R\) is large compared with source
extent and wavelength \cite{harrington2001}. Section~\ref{sec:ch24} treats
Eq.~\eqref{eq:ch2.silver-muller-trace-approx} as a practical surrogate; the
continuous target remains Eq.~\eqref{eq:ch2.sommerfeld-field-condition}
\cite{stratton1941}. Perfectly matched layers and radiation boundary integrals are
discrete refinements of the same outgoing-trace logic \cite{harrington2001}.

\paragraph{Link to radiation operator.}
The outgoing selector \(\Pi_{\mathrm{rad}}\) of
Eq.~\eqref{eq:ch2.radiation-operator-def} is idempotent on fields satisfying
Eq.~\eqref{eq:ch2.sommerfeld-outgoing-class} once uniqueness is established in
Subsection~\ref{sec:ch243} \cite{stratton1941}. Prior to that proof, Sommerfeld
closure defines the domain on which \(\Pi_{\mathrm{rad}}\) acts as projection rather
than heuristic windowing \cite{harrington2001}. Composition with
Eq.~\eqref{eq:ch2.transport-qualified-radiation} yields
\begin{equation}
\label{eq:ch2.sommerfeld-radiation-composition}
\mathcal{R}_\omega[\Jimp]
=
\Pi_{\mathrm{rad}}\,\mathcal{S}_\omega[\Jimp],
\qquad
\mathcal{S}_\omega[\Jimp]\in\mathcal{F}_{\mathrm{SF}},
\end{equation}
with transport qualification from Section~\ref{sec:ch01.2.4} applied after
Eq.~\eqref{eq:ch2.sommerfeld-radiation-composition} \cite{stratton1941}. Equation
\eqref{eq:ch2.sommerfeld-radiation-composition} restates
Eq.~\eqref{eq:ch2.radiation-time-frequency-equiv} on the Sommerfeld domain without
re-proving causality from Subsection~\ref{sec:ch223} \cite{harrington2001}.

\paragraph{Boundary of validity.}
Equations~\eqref{eq:ch2.sommerfeld-scalar-asymptotic}--\eqref{eq:ch2.sommerfeld-outgoing-class}
require linear vacuum far zones, fixed \(\omega>0\), and compact \(\operatorname{supp}\Jimp\).
They fail for anisotropic or magneto-electric exterior media without modified
asymptotics, for lossless closed cavities where Eq.~\eqref{eq:ch1.outer-boundary-classes}
uses PEC rather than outgoing closure, and for time-harmonic fields with net incoming
power at infinity \cite{stratton1941,harrington2001}. Loss shifts
Eq.~\eqref{eq:ch2.sommerfeld-scalar-asymptotic} to complex \(k_0\) but preserves
outgoing \emph{selection} as distinct from incoming \cite{stratton1941}.

\paragraph{Worked illustration (outgoing spherical wave).}
In vacuum with \(\E,\H\propto e^{-\jj k_0 r}/r\) and \(\E_0\cdot\hat{\mathbf{r}}=0\),
Eqs.~\eqref{eq:ch2.sommerfeld-field-condition} and~\eqref{eq:ch2.sommerfeld-impedance-ratio}
hold exactly. Reversing \(k_0\to -k_0\) produces an incoming wave still satisfying
volume Maxwell rows but violating Eq.~\eqref{eq:ch2.sommerfeld-outgoing-class} and
Eq.~\eqref{eq:ch2.retarded-harmonic-kernel-theorem} \cite{stratton1941,jackson1999}.
Harmonic solvers that do not enforce Eq.~\eqref{eq:ch2.sommerfeld-field-condition}
may return such incoming components with equal legitimacy at the PDE level
\cite{harrington2001}.

\paragraph{Worked illustration (finite truncation sphere).}
Replace Eq.~\eqref{eq:ch2.sommerfeld-field-condition} by
Eq.~\eqref{eq:ch2.silver-muller-trace-approx} on \(\Gamma_R\) with \(R=5\lambda_0\).
Fields inside \(\Omega_s\) differ from the exact Sommerfeld solution by
\(\mathcal{O}((k_0 R)^{-1})\) in typical radiative components while reactive near
content may differ by order unity \cite{stratton1941,harrington2001}. Declaring a
transport channel from \(\Gamma_R\) probes without
Eq.~\eqref{eq:ch2.regime-operator-screen} conflates truncation error with angular
transport \cite{stratton1941}.

\paragraph{Closing summary.}
Subsection~\ref{sec:ch241} establishes the Chapter~\ref{ch:02} HOME for scalar
Sommerfeld asymptotic Eq.~\eqref{eq:ch2.sommerfeld-scalar-asymptotic}, vector
condition Eq.~\eqref{eq:ch2.sommerfeld-field-condition}, impedance ratio
Eq.~\eqref{eq:ch2.sommerfeld-impedance-ratio}, far-zone transverse radiation
Eq.~\eqref{eq:ch2.far-zone-transverse-radiation}, dyadic asymptotic
Eq.~\eqref{eq:ch2.sommerfeld-dyadic-asymptotic}, outgoing class
Eq.~\eqref{eq:ch2.sommerfeld-outgoing-class}, Silver--Müller surrogate
Eq.~\eqref{eq:ch2.silver-muller-trace-approx}, and radiation composition
Eq.~\eqref{eq:ch2.sommerfeld-radiation-composition}, embedding
Eq.~\eqref{eq:ch1.radiation-trace-condition} and
Eq.~\eqref{eq:ch1.outer-boundary-classes} by reference. Subsection~\ref{sec:ch242}
operatorizes outgoing selection on Green kernels; Subsection~\ref{sec:ch243} proves
uniqueness on \(\mathcal{F}_{\mathrm{SF}}\).

\subsection{Outgoing-wave selection}
\label{sec:ch242}
% TOC tag: [HOME]

Subsection~\ref{sec:ch241} fixed Sommerfeld asymptotics and the outgoing function class
Eq.~\eqref{eq:ch2.sommerfeld-outgoing-class}. Subsection~\ref{sec:ch242} gives the
Chapter~\ref{ch:02} HOME for the \emph{outgoing-wave selector}
\(\mathcal{O}_{\mathrm{out}}\) that implements that closure as an operator on
\(\mathcal{F}_\Omega\), on harmonic Green kernels, and on the resolvent branch of
\(\mathcal{L}_\omega^{-1}\) \cite{stratton1941,harrington2001,jackson1999}. The
composition preview Eq.~\eqref{eq:ch2.rad-projection-composition} from
Subsection~\ref{sec:ch224} named \(\mathcal{O}_{\mathrm{out}}\) without constructing
it; the present subsection supplies the operator definition and its links to
\(\Pi_{\mathrm{prop}}\), \(\Pi_{\mathrm{rad}}\), and retarded Green functions
Eq.~\eqref{eq:ch2.retarded-harmonic-kernel-theorem}. Sommerfeld field conditions
Eqs.~\eqref{eq:ch2.sommerfeld-field-condition}--\eqref{eq:ch2.sommerfeld-outgoing-class}
are presupposed; uniqueness proofs belong to Subsection~\ref{sec:ch243}.

\paragraph{Assumptions.}
Exterior regions use vacuum far-zone constants with \(k_0>0\) real in the locked
harmonic band, compact \(\operatorname{supp}\Jimp\), and time-independent
\(\mathcal{B}_\Gamma\) including outgoing closure at infinity or on
Eq.~\eqref{eq:ch2.silver-muller-trace-approx} at large finite \(R\)
\cite{stratton1941}. Propagating and evanescent projectors
Eqs.~\eqref{eq:ch2.spectral-projectors}--\eqref{eq:ch2.spectral-projector-algebra}
from Subsection~\ref{sec:ch224} are fixed on \(\widehat{\mathcal{F}}_\Omega\)
\cite{harrington2001}. Incoming and advanced branches satisfying
Eq.~\eqref{eq:ch2.advanced-kernel-violation} are excluded.

\paragraph{Outgoing selector on fields.}
Let \(\mathcal{F}_{\mathrm{in}}\) denote the incoming counterpart of
\(\mathcal{F}_{\mathrm{SF}}\) obtained by reversing outgoing phase in the far zone.
The outgoing-wave selector is the idempotent map
\begin{equation}
\label{eq:ch2.outgoing-selector-def}
\mathcal{O}_{\mathrm{out}}:\mathcal{F}_\Omega\to\mathcal{F}_{\mathrm{SF}},
\qquad
\mathcal{O}_{\mathrm{out}}[\mathbf{u}_\omega]
=
\mathbf{u}_{\mathrm{out}},
\end{equation}
where \(\mathbf{u}_{\mathrm{out}}\) is the unique member of
\(\mathcal{F}_{\mathrm{SF}}\) whose propagating spectral content matches the outgoing
\(k_z\) branch of Eq.~\eqref{eq:ch2.longitudinal-wavenumber-split} and whose
incoming content is annihilated \cite{stratton1941,harrington2001}. Equation
\eqref{eq:ch2.outgoing-selector-def} is the Chapter~\ref{ch:02} HOME operator
realization of Sommerfeld closure: it acts \emph{before} far-zone measurement and
\emph{after} volume closure by \(\mathcal{L}_\omega\), selecting the physical
resolvent branch assumed in Eq.~\eqref{eq:ch2.resolvent-green-equivalence}
\cite{stratton1941}. On \(\mathcal{F}_{\mathrm{SF}}\),
\begin{equation}
\label{eq:ch2.outgoing-idempotency}
\mathcal{O}_{\mathrm{out}}^2=\mathcal{O}_{\mathrm{out}},
\qquad
\mathcal{O}_{\mathrm{out}}\big|_{\mathcal{F}_{\mathrm{SF}}}=\mathrm{id},
\end{equation}
while incoming states satisfy \(\mathcal{O}_{\mathrm{out}}[\mathbf{u}_{\mathrm{in}}]=\mathbf{0}\)
when extended by continuity from the far zone \cite{harrington2001}. Physically,
\(\mathcal{O}_{\mathrm{out}}\) removes incoming spherical and plane-wave components
that satisfy Maxwell equations but violate retarded causality
\cite{jackson1999,stratton1941}.

\paragraph{Plane-wave and spectral gate.}
On homogeneous regions, outgoing selection in lateral wavenumber \((k_x,k_y)\) uses
the longitudinal root from Eq.~\eqref{eq:ch2.longitudinal-wavenumber-split}:
\begin{equation}
\label{eq:ch2.outgoing-plane-wave-gate}
k_z=
\begin{cases}
+\sqrt{k_0^2-k_t^2}, & k_t\le k_0\ \text{(outgoing propagating)},\\[4pt]
-\jj\,\alpha(k_t), & k_t>k_0\ \text{(decaying evanescent)},
\end{cases}
\end{equation}
with the sign of the propagating root fixed by exterior closure and the locked
\(\exp(-\jj\omega t)\) convention \cite{stratton1941,harrington2001}. Equation
\eqref{eq:ch2.outgoing-plane-wave-gate} is the spectral action of
\(\mathcal{O}_{\mathrm{out}}\) on each plane-wave mode; reversing
\(k_z\to -k_z\) on the propagating sheet produces incoming content annihilated by
\(\mathcal{O}_{\mathrm{out}}\) \cite{stratton1941}. Evanescent modes are unchanged
in decay direction but remain outside \(\Pi_{\mathrm{rad}}\) through
Eq.~\eqref{eq:ch2.rad-null-on-evanescent} \cite{harrington2001}. The propagating
projector \(\Pi_{\mathrm{prop}}\) from Subsection~\ref{sec:ch224} retains all
outgoing and incoming propagating modes; \(\mathcal{O}_{\mathrm{out}}\) acts only
after \(\Pi_{\mathrm{prop}}\) in the radiation chain
Eq.~\eqref{eq:ch2.rad-projection-composition} \cite{stratton1941}.

\begin{figure}[t]
\centering
\ChOneFigBlock{%
\begin{tikzpicture}[ch1fig, node distance=7mm and 10mm]
  \node[ch1box] (f) {$\mathcal{F}_\Omega$};
  \node[ch1box, right=14mm of f] (pp) {$\Pi_{\mathrm{prop}}$};
  \node[ch1box, right=14mm of pp] (oo) {$\mathcal{O}_{\mathrm{out}}$};
  \node[ch1box, right=14mm of oo] (pr) {$\Pi_{\mathrm{rad}}$};
  \node[ch1box, below=11mm of f] (ev) {$\Pi_{\mathrm{ev}}$};
  \draw[ch1flow] (f) -- (pp) -- (oo) -- (pr);
  \draw[ch1flow] (f) -- (ev);
\end{tikzpicture}%
}
\caption{Outgoing selection Eq.~\eqref{eq:ch2.outgoing-selector-def} sits between
propagating projection and radiation export:
\(\Pi_{\mathrm{rad}}=\mathcal{O}_{\mathrm{out}}\circ\Pi_{\mathrm{prop}}\).}
\label{fig:ch2.outgoing-selector-pipeline}
\end{figure}

Figure~\ref{fig:ch2.outgoing-selector-pipeline} refines
Eq.~\eqref{eq:ch2.rad-projection-composition} with explicit operator order. Regime
qualification Eq.~\eqref{eq:ch2.regime-operator-screen} applies after
\(\Pi_{\mathrm{rad}}\), not at the \(\mathcal{O}_{\mathrm{out}}\) stage
\cite{stratton1941,harrington2001}.

\paragraph{Green-function branch selection.}
Harmonic kernels from Section~\ref{sec:ch23} admit incoming and outgoing analytic
continuations. The outgoing dyadic branch is
\begin{equation}
\label{eq:ch2.outgoing-green-branch}
\overline{\mathbf{G}}_{E,\omega}^{\mathrm{(out)}}
=
\mathcal{O}_{\mathrm{out}}\,\overline{\mathbf{G}}_{E,\omega},
\qquad
\overline{\mathbf{G}}_{E,\omega}^{\mathrm{(out)}}
\ \text{satisfies}\ 
\text{Eq.~\eqref{eq:ch2.sommerfeld-dyadic-asymptotic}},
\end{equation}
with \(\mathcal{O}_{\mathrm{out}}\) acting on the observation-point dependence
\cite{stratton1941,harrington2001}. Equation~\eqref{eq:ch2.outgoing-green-branch}
identifies the kernel used in Eq.~\eqref{eq:ch2.electric-dyadic-convolution} on
exterior problems with the retarded symbol of
Eq.~\eqref{eq:ch2.retarded-harmonic-kernel-theorem} \cite{stratton1941}. Incoming
kernels \(\overline{\mathbf{G}}_{E,\omega}^{\mathrm{(in)}}\) are related by
\(k_0\to -k_0\) on the propagating Riemann sheet and produce fields outside
\(\mathcal{F}_{\mathrm{SF}}\) \cite{jackson1999}. Subsection~\ref{sec:ch242} does
not re-tabulate Eq.~\eqref{eq:ch2.free-space-dyadic-green}; it fixes the sheet on
which that formula is evaluated \cite{harrington2001}.

\paragraph{Resolvent and solution-operator selection.}
The physical resolvent branch is
\begin{equation}
\label{eq:ch2.outgoing-resolvent-selection}
\mathcal{L}_\omega^{-1}\big|_{\mathrm{phys}}
=
\mathcal{O}_{\mathrm{out}}\circ\mathcal{L}_\omega^{-1},
\end{equation}
understood on \(\mathcal{D}(\mathcal{L}_\omega)\) with outgoing \(\mathcal{B}_\Gamma\)
\cite{stratton1941}. The solution map on impressed currents becomes
\begin{equation}
\label{eq:ch2.outgoing-solution-operator}
\mathcal{S}_\omega[\Jimp]
=
\mathcal{O}_{\mathrm{out}}\,
\Big(\jj\omega\mu_0\int_{\Omega}
\overline{\mathbf{G}}_{E,\omega}^{\mathrm{(out)}}(\mathbf{r},\mathbf{r}')\cdot\Jimp(\mathbf{r}')\,d^{3}r'
\;+\;\E_{\mathrm{hom}}\Big),
\end{equation}
with regularization from Subsection~\ref{sec:ch232} when \(\mathbf{r}\) approaches
\(\operatorname{supp}\Jimp\) \cite{harrington2001,stratton1941}. Equations
\eqref{eq:ch2.outgoing-resolvent-selection} and~\eqref{eq:ch2.outgoing-solution-operator}
are the operator content of Eq.~\eqref{eq:ch2.exterior-uniqueness-hypothesis} once
Subsection~\ref{sec:ch243} proves uniqueness on the selected branch
\cite{stratton1941}. Causality from Subsection~\ref{sec:ch233} is consistent with
Eq.~\eqref{eq:ch2.outgoing-green-branch} but not sufficient when harmonic solvers
bypass time-domain construction \cite{harrington2001}.

\paragraph{Radiation export refinement.}
On propagating content,
\begin{equation}
\label{eq:ch2.outgoing-radiation-refinement}
\Pi_{\mathrm{rad}}
=
\mathcal{O}_{\mathrm{out}}\circ\Pi_{\mathrm{prop}},
\qquad
\mathcal{R}_\omega
=
\Pi_{\mathrm{rad}}\circ\mathcal{S}_\omega,
\end{equation}
restating Eqs.~\eqref{eq:ch2.rad-projection-composition} and~\eqref{eq:ch2.radiation-operator-def}
with the Chapter~\ref{ch:02} HOME definition of \(\mathcal{O}_{\mathrm{out}}\)
\cite{stratton1941,harrington2001}. Radiative subspace
\(\mathcal{F}_{\mathrm{rad}}\subset\mathcal{F}_{\mathrm{SF}}\) consists of fields
with nonzero far-zone transverse amplitude Eq.~\eqref{eq:ch2.far-zone-transverse-radiation};
\(\mathcal{O}_{\mathrm{out}}\) maps into \(\mathcal{F}_{\mathrm{SF}}\) but not
automatically into \(\mathcal{F}_{\mathrm{rad}}\) when evanescent content dominates
near sources \cite{stratton1941}. Transport reporting therefore requires
Eq.~\eqref{eq:ch2.regime-operator-screen} after
Eq.~\eqref{eq:ch2.outgoing-radiation-refinement} \cite{harrington2001}.

\paragraph{Incoming--outgoing split.}
Define \(\mathcal{O}_{\mathrm{in}}\) analogously on \(\mathcal{F}_{\mathrm{in}}\).
On the propagating subspace,
\begin{equation}
\label{eq:ch2.incoming-outgoing-complement}
\mathrm{id}
=
\mathcal{O}_{\mathrm{out}}+\mathcal{O}_{\mathrm{in}}
\quad\text{on}\ \Pi_{\mathrm{prop}}\mathcal{F}_\Omega,
\qquad
\mathcal{O}_{\mathrm{out}}\mathcal{O}_{\mathrm{in}}=0,
\end{equation}
with evanescent and non-propagating components orthogonal to both when decay
direction is fixed by Eq.~\eqref{eq:ch2.outgoing-plane-wave-gate}
\cite{stratton1941,harrington2001}. Equation~\eqref{eq:ch2.incoming-outgoing-complement}
explains branch degeneracy in unconstrained Helmholtz solvers: both
\(\mathcal{O}_{\mathrm{out}}[\mathbf{u}]\) and
\(\mathcal{O}_{\mathrm{in}}[\mathbf{u}]\) satisfy volume equations unless one is
selected by boundary or radiation closure \cite{jackson1999}.

\paragraph{Boundary of validity.}
Equations~\eqref{eq:ch2.outgoing-selector-def}--\eqref{eq:ch2.outgoing-solution-operator}
require linear media, fixed \(\omega\), and outgoing \(\mathcal{B}_\Gamma\) on
exterior problems. They fail on closed PEC cavities where
Eq.~\eqref{eq:ch1.outer-boundary-classes} replaces Sommerfeld closure, and on
lossless open resonators where discrete spectra require quotienting
Eq.~\eqref{eq:ch2.coercivity-quotient} \cite{stratton1941,harrington2001}.
Absorbing layers approximate \(\mathcal{O}_{\mathrm{out}}\) but introduce spurious
modes unless tuned to Eq.~\eqref{eq:ch2.sommerfeld-field-condition} at large distance
\cite{harrington2001}.

\paragraph{Worked illustration (plane-wave pair).}
For \(\E=\E_0 e^{-\jj(k_x x+k_y y+k_z z)}\) with \(k_z=+\sqrt{k_0^2-k_t^2}\),
\(\mathcal{O}_{\mathrm{out}}\) leaves the field unchanged; for
\(k_z=-\sqrt{k_0^2-k_t^2}\), \(\mathcal{O}_{\mathrm{out}}\) annihilates the mode on
exterior domains \cite{stratton1941,jackson1999}. Applying \(\Pi_{\mathrm{rad}}\) to
the incoming plane wave yields zero export even though \(|\E|\) is nonzero
\cite{harrington2001}.

\paragraph{Worked illustration (two-sheet Green evaluation).}
Evaluating Eq.~\eqref{eq:ch2.free-space-dyadic-green} with \(k_0\) and \(-k_0\) on
the propagating sheet produces distinct fields in the near zone. Only
Eq.~\eqref{eq:ch2.outgoing-green-branch} matches
Eq.~\eqref{eq:ch2.retarded-harmonic-kernel-theorem} for causal
\(\mathbf{j}_{\mathrm{imp}}\) \cite{stratton1941,harrington2001}. Frequency-domain
codes that default to principal-value Helmholtz inverses without
Eq.~\eqref{eq:ch2.outgoing-resolvent-selection} may superpose both sheets
\cite{stratton1941}.

\paragraph{Closing summary.}
Subsection~\ref{sec:ch242} establishes the Chapter~\ref{ch:02} HOME for outgoing
selector Eqs.~\eqref{eq:ch2.outgoing-selector-def}--\eqref{eq:ch2.outgoing-idempotency},
spectral gate Eq.~\eqref{eq:ch2.outgoing-plane-wave-gate}, outgoing Green branch
Eq.~\eqref{eq:ch2.outgoing-green-branch}, resolvent and solution selection
Eqs.~\eqref{eq:ch2.outgoing-resolvent-selection}--\eqref{eq:ch2.outgoing-solution-operator},
radiation refinement Eq.~\eqref{eq:ch2.outgoing-radiation-refinement}, and
incoming--outgoing complement Eq.~\eqref{eq:ch2.incoming-outgoing-complement},
completing Eq.~\eqref{eq:ch2.rad-projection-composition}. Subsection~\ref{sec:ch243}
proves uniqueness on \(\mathcal{F}_{\mathrm{SF}}\); Subsection~\ref{sec:ch244} treats
reconstruction identifiability.

\subsection{Uniqueness of radiative solutions}
\label{sec:ch243}
% TOC tag: [HOME][USE SS1.1.3]

Subsections~\ref{sec:ch241}--\ref{sec:ch242} fixed Sommerfeld closure and the outgoing
selector \(\mathcal{O}_{\mathrm{out}}\). Subsection~\ref{sec:ch243} gives the
Chapter~\ref{ch:02} HOME proof that exterior radiative solutions are unique on the
outgoing branch, upgrading the hypothesis
Eq.~\eqref{eq:ch2.exterior-uniqueness-hypothesis} to Theorem~\ref{thm:ch2.exterior-radiative-uniqueness}
\cite{stratton1941,harrington2001,jackson1999}. The difference-field reduction
Eqs.~\eqref{eq:ch1.uniqueness-difference-system}--\eqref{eq:ch1.uniqueness-null} and
energy identity Eq.~\eqref{eq:ch1.uniqueness-energy-identity} from
Subsection~\ref{sec:ch01.1.3} are presupposed; passive bounded-domain uniqueness
Proposition~\ref{prop:ch1.uniqueness-passive-closure} and coercive estimate
Eq.~\eqref{eq:ch1.uniqueness-estimate} are not re-proved. The present subsection
specializes the energy method to exterior domains with
Eqs.~\eqref{eq:ch2.sommerfeld-field-condition}--\eqref{eq:ch2.sommerfeld-outgoing-class},
outgoing resolvent selection
Eqs.~\eqref{eq:ch2.outgoing-resolvent-selection}--\eqref{eq:ch2.outgoing-solution-operator},
and radiation refinement Eq.~\eqref{eq:ch2.outgoing-radiation-refinement}. Source
reconstruction limits belong to Subsection~\ref{sec:ch244}.

\paragraph{Assumptions and function class.}
Let \(\Omega_{\mathrm{ext}}\) denote an exterior domain with compact
\(\operatorname{supp}\Jimp\subset B_R(0)\), vacuum constants \(\varepsilon_0,\mu_0\)
in the far zone, and fixed \(\omega>0\) in the harmonic band of
Subsection~\ref{sec:ch223}. Data are admissible per
Eq.~\eqref{eq:ch1.admissible-data-class} with outgoing closure
Eq.~\eqref{eq:ch2.sommerfeld-outgoing-class} at infinity, equivalently
Eq.~\eqref{eq:ch1.radiation-trace-condition} on spherical shells
\cite{stratton1941}. Constitutive response in the source region is passive:
\(\Im\{\varepsilon\}\ge 0\), \(\Im\{\mu\}\ge 0\) or equivalent loss so homogeneous
difference energy is dissipative \cite{stratton1941}. Incoming and advanced branches
are excluded by Eq.~\eqref{eq:ch2.outgoing-selector-def} and
Eq.~\eqref{eq:ch2.incoming-outgoing-complement} \cite{harrington2001}. Fields lie in
\(\mathcal{F}_{\mathrm{SF}}\) as defined in Subsection~\ref{sec:ch241}; radiative
export uses \(\mathcal{F}_{\mathrm{rad}}\subset\mathcal{F}_{\mathrm{SF}}\) from
Eq.~\eqref{eq:ch2.far-zone-transverse-radiation} \cite{stratton1941}.

\paragraph{Difference fields on the outgoing branch.}
Let \((\E_1,\H_1)\) and \((\E_2,\H_2)\) be two harmonic solutions on the outgoing
branch for the same impressed data and \(\mathcal{B}_\Gamma\). Define
\(\Delta\E=\E_1-\E_2\), \(\Delta\H=\H_1-\H_2\), and
\(\mathbf{u}_\Delta=(\Delta\E,\Delta\H)\). Linearity and
Eq.~\eqref{eq:ch1.uniqueness-difference-system} give the homogeneous exterior system
\begin{equation}
\label{eq:ch2.exterior-difference-system}
\curl\Delta\E=-\jj\omega\mu\,\Delta\H,\quad
\curl\Delta\H=\jj\omega\varepsilon\,\Delta\E,\quad
\mathcal{B}_\Gamma[\Delta\E,\Delta\H]=\mathbf{0},
\end{equation}
with \(\varepsilon,\mu\) matching the physical model in \(\Omega_s\) and
\(\varepsilon_0,\mu_0\) outside \cite{stratton1941}. Outgoing closure requires
\begin{equation}
\label{eq:ch2.exterior-difference-outgoing}
\mathcal{O}_{\mathrm{out}}[\mathbf{u}_\Delta]=\mathbf{u}_\Delta,
\qquad
\mathbf{u}_\Delta\in\mathcal{F}_{\mathrm{SF}},
\end{equation}
which is the exterior specialization of the third row of
Eq.~\eqref{eq:ch1.uniqueness-difference-system} when
\(\mathcal{B}_\Gamma\) includes Sommerfeld radiation \cite{harrington2001}. Equations
\eqref{eq:ch2.exterior-difference-system}--\eqref{eq:ch2.exterior-difference-outgoing}
state that uniqueness testing reduces to showing only the trivial pair satisfies the
homogeneous outgoing problem \cite{stratton1941}. Incoming difference modes would
violate Eq.~\eqref{eq:ch2.exterior-difference-outgoing} even if they satisfied the
curl rows \cite{jackson1999}.

\paragraph{Energy identity on truncated exterior shells.}
Form the complex Poynting vector for difference fields,
\(\mathbf{S}_\Delta=\tfrac12\Delta\E\times\Delta\H^{*}\), as in
Eq.~\eqref{eq:ch1.uniqueness-energy-identity} \cite{stratton1941}. For a large
sphere \(\Gamma_R=\{r=R\}\) enclosing \(\Omega_s\), let
\(\Omega_R=B_R\setminus\Omega_s\). The divergence theorem on \(\Omega_R\) yields
\begin{equation}
\label{eq:ch2.exterior-flux-identity}
\int_{\Gamma_R}\hat{\mathbf{r}}\cdot\mathbf{S}_\Delta\,dS
=
-\jj\omega\!\int_{\Omega_R}
\Big(u_{\Delta H}-u_{\Delta E}\Big)\,dV
\;+\;\mathcal{P}_{\Delta,s},
\end{equation}
where \(u_{\Delta E}=\tfrac14\varepsilon|\Delta\E|^{2}\),
\(u_{\Delta H}=\tfrac14\mu|\Delta\H|^{2}\), and \(\mathcal{P}_{\Delta,s}\) collects
source-region passive dissipation when \(\Im\{\varepsilon\},\Im\{\mu\}>0\)
\cite{stratton1941,harrington2001}. Equation~\eqref{eq:ch2.exterior-flux-identity}
is Eq.~\eqref{eq:ch1.uniqueness-energy-identity} with \(\Delta\Jimp=\mathbf{0}\) and
\(\Gamma=\Gamma_R\); its real part controls uniqueness through the sign of the
boundary flux \cite{stratton1941}. In lossless vacuum on
\(\Omega_R\setminus\Omega_s\), the volume term is purely imaginary at steady state,
so \(\Re\int_{\Gamma_R}\hat{\mathbf{r}}\cdot\mathbf{S}_\Delta\,dS\) is determined
entirely by far-zone radiation \cite{harrington2001}.

\paragraph{Sommerfeld flux limit.}
Sommerfeld condition Eq.~\eqref{eq:ch2.sommerfeld-field-condition} implies that on
\(\Gamma_R\) with \(R\gg a=\max|\mathbf{r}|\) over \(\operatorname{supp}\Jimp\), the
leading transverse contribution to \(\Delta\E,\Delta\H\) scales as
\(e^{-\jj k_0 R}/R\) with \(k_0=\omega\sqrt{\varepsilon_0\mu_0}\)
\cite{stratton1941,jackson1999}. The far-zone transverse radiation law
Eq.~\eqref{eq:ch2.far-zone-transverse-radiation} then gives
\begin{equation}
\label{eq:ch2.exterior-flux-asymptotic}
\Re\int_{\Gamma_R}\hat{\mathbf{r}}\cdot\mathbf{S}_\Delta\,dS
=
P_{\Delta,\mathrm{rad}}(R)
=
\frac{1}{2\eta_0}\oint_{\Gamma_R}
\big|\mathbf{E}_{\Delta,t}\big|^{2}\,dS
\;+\;\mathcal{O}(R^{-1}),
\end{equation}
with \(\eta_0=\sqrt{\mu_0/\varepsilon_0}\) and \(\mathbf{E}_{\Delta,t}\) the
transverse spherical component of \(\Delta\E\) \cite{stratton1941}. Taking
\(R\to\infty\) under outgoing closure,
\begin{equation}
\label{eq:ch2.exterior-sommerfeld-flux-limit}
\lim_{R\to\infty}P_{\Delta,\mathrm{rad}}(R)
=
P_{\Delta,\mathrm{rad}}\ge 0,
\end{equation}
where \(P_{\Delta,\mathrm{rad}}\) is the net outgoing power carried by the difference
field \cite{stratton1941,harrington2001}. Equation
\eqref{eq:ch2.exterior-sommerfeld-flux-limit} is the exterior boundary term induced
by Eq.~\eqref{eq:ch1.radiation-trace-condition}: outgoing difference fields radiate
nonnegative power to infinity, while incoming difference fields would produce negative
flux and are excluded by \(\mathcal{O}_{\mathrm{out}}\)
\cite{jackson1999,stratton1941}. Dyadic asymptotics
Eq.~\eqref{eq:ch2.sommerfeld-dyadic-asymptotic} are presupposed; they are not
re-tabulated here \cite{harrington2001}.

\begin{figure}[t]
\centering
\ChOneFigBlock{%
\begin{tikzpicture}[ch1fig, node distance=7mm and 10mm]
  \node[ch1box] (s) {Compact\\$\Omega_s$};
  \node[ch1box, right=16mm of s] (d) {$\Delta\E,\Delta\H$\\outgoing};
  \node[ch1box, right=16mm of d] (gr) {$\Gamma_R$\\flux};
  \node[ch1box, right=16mm of gr] (out) {$P_{\Delta,\mathrm{rad}}\ge 0$};
  \draw[ch1flow] (s) -- (d) -- (gr) -- (out);
\end{tikzpicture}%
}
\caption{Exterior uniqueness flux argument: homogeneous difference fields on the
outgoing branch satisfy Eqs.~\eqref{eq:ch2.exterior-flux-identity}--\eqref{eq:ch2.exterior-sommerfeld-flux-limit}
on spheres \(\Gamma_R\).}
\label{fig:ch2.exterior-uniqueness-flux}
\end{figure}

Figure~\ref{fig:ch2.exterior-uniqueness-flux} organizes the proof template imported
from Subsection~\ref{sec:ch01.1.3}: volume terms in
Eq.~\eqref{eq:ch2.exterior-flux-identity} are controlled by passivity when loss is
present, and vanish in lossless vacuum away from sources \cite{stratton1941}.

\paragraph{Lemma (vanishing outgoing transverse radiation).}
\begin{lemma}[Zero outgoing flux implies trivial difference field]
\label{lem:ch2.zero-flux-difference}
Let \(\mathbf{u}_\Delta=(\Delta\E,\Delta\H)\) satisfy
Eqs.~\eqref{eq:ch2.exterior-difference-system}--\eqref{eq:ch2.exterior-difference-outgoing}
on an exterior domain with vacuum far zone. If
\(P_{\Delta,\mathrm{rad}}=0\) in Eq.~\eqref{eq:ch2.exterior-sommerfeld-flux-limit},
then \(\mathbf{E}_{\Delta,t}=\mathbf{0}\) on every sphere \(\Gamma_R\) for
\(R\) large enough, and \(\mathbf{u}_\Delta\) has no propagating outgoing content in
\(\mathcal{F}_{\mathrm{rad}}\) \cite{stratton1941,harrington2001}.
\end{lemma}
\begin{proof}
Eq.~\eqref{eq:ch2.exterior-flux-asymptotic} with \(P_{\Delta,\mathrm{rad}}=0\) forces
\(\oint_{\Gamma_R}|\mathbf{E}_{\Delta,t}|^{2}\,dS=0\) for all sufficiently large
\(R\). Outgoing transverse fields with zero spherical flux cannot carry
\(1/R\) radiation; the only remaining homogeneous outgoing transverse modes are
identically zero on \(\mathcal{F}_{\mathrm{rad}}\) \cite{stratton1941,jackson1999}.
\end{proof}

\paragraph{Exterior uniqueness theorem.}
\begin{theorem}[Exterior radiative uniqueness]
\label{thm:ch2.exterior-radiative-uniqueness}
Let \(\Omega_{\mathrm{ext}}\) be an exterior domain with compact impressed sources,
vacuum far zone, and outgoing Sommerfeld closure
Eqs.~\eqref{eq:ch2.sommerfeld-outgoing-class}--\eqref{eq:ch2.outgoing-selector-def}.
If \((\E_1,\H_1)\) and \((\E_2,\H_2)\) are harmonic Maxwell solutions on the outgoing
branch for the same admissible data, then \(\E_1=\E_2\) and \(\H_1=\H_2\).
Equivalently,
\begin{equation}
\label{eq:ch2.exterior-uniqueness-theorem}
\mathcal{L}_\omega[\mathbf{u}_\omega]=\mathbf{0}
\ \text{and}\ 
\mathbf{u}_\omega\in\mathcal{F}_{\mathrm{SF}}
\ \Longrightarrow\
\mathbf{u}_\omega=\mathbf{0},
\end{equation}
which upgrades Eq.~\eqref{eq:ch2.exterior-uniqueness-hypothesis} to the Chapter~\ref{ch:02}
HOME uniqueness statement \cite{stratton1941,harrington2001}.
\end{theorem}
\begin{proof}
Let \(\mathbf{u}_\Delta\) be the difference of two outgoing solutions. Then
Eqs.~\eqref{eq:ch2.exterior-difference-system}--\eqref{eq:ch2.exterior-difference-outgoing}
hold. Apply Eq.~\eqref{eq:ch2.exterior-flux-identity} and take
\(\Re\{\cdot\}\). In lossless vacuum outside \(\Omega_s\), the volume integral is
purely imaginary, so \(\Re P_{\Delta,\mathrm{rad}}(R)\) equals the real boundary flux.
Outgoing closure forces \(\Re P_{\Delta,\mathrm{rad}}\ge 0\); passivity in
\(\Omega_s\) adds nonnegative \(\Re\mathcal{P}_{\Delta,s}\), so the only way to
satisfy Eq.~\eqref{eq:ch2.exterior-sommerfeld-flux-limit} with equality throughout is
\(P_{\Delta,\mathrm{rad}}=0\) \cite{stratton1941}. Lemma~\ref{lem:ch2.zero-flux-difference}
eliminates outgoing transverse content. Evanescent components remain in
\(\mathcal{F}_{\mathrm{SF}}\) but are tied to the same source-free outgoing solution
operator; with no radiative export and homogeneous forcing, the unique retarded
exterior state is trivial by Eq.~\eqref{eq:ch2.outgoing-solution-operator} with
\(\Jimp=\mathbf{0}\) \cite{harrington2001,stratton1941}. Hence
\(\mathbf{u}_\Delta=\mathbf{0}\), proving Eq.~\eqref{eq:ch2.exterior-uniqueness-theorem}.
\end{proof}

\paragraph{Relation to Chapter~\ref{ch:01} coercivity.}
Theorem~\ref{thm:ch2.exterior-radiative-uniqueness} is the exterior counterpart of
Proposition~\ref{prop:ch1.uniqueness-passive-closure}: instead of a bounded passive
\(\Gamma\) with coercivity Eq.~\eqref{eq:ch1.uniqueness-estimate}, the exterior uses
Sommerfeld flux Eq.~\eqref{eq:ch2.exterior-sommerfeld-flux-limit} as the coercive
boundary channel \cite{stratton1941}. The uniqueness index
Eq.~\eqref{eq:ch1.uniqueness-condition-index} is finite on exterior outgoing problems
away from embedded resonances in finite scatterers; it diverges on bounded cavities at
Eq.~\eqref{eq:ch2.cavity-resonance-failure} \cite{harrington2001}. Exterior uniqueness
therefore replaces bounded coercivity with an \emph{open} boundary condition that
selects decaying outgoing phase \cite{stratton1941,jackson1999}.

\paragraph{Consequences for resolvent and solution map.}
Theorem~\ref{thm:ch2.exterior-radiative-uniqueness} makes the physical resolvent
injective on the outgoing subspace:
\begin{equation}
\label{eq:ch2.exterior-resolvent-injective}
\mathcal{L}_\omega^{-1}\big|_{\mathrm{phys}}:
\mathcal{Y}_\omega\to\mathcal{F}_{\mathrm{SF}}
\ \text{is single-valued},
\qquad
\mathcal{S}_\omega[\Jimp]
=
\mathcal{L}_\omega^{-1}\big|_{\mathrm{phys}}
\begin{bmatrix}\mathbf{s}_\omega(\Jimp,\rho)\\ \mathbf{g}_\Gamma\end{bmatrix},
\end{equation}
consistent with Eq.~\eqref{eq:ch2.green-bvp-equivalence}, outgoing Green branch
Eq.~\eqref{eq:ch2.outgoing-green-branch}, and solution selection
Eq.~\eqref{eq:ch2.outgoing-solution-operator} \cite{stratton1941,harrington2001}.
\begin{corollary}[Uniqueness of outgoing Green convolution]
\label{cor:ch2.outgoing-green-unique}
For fixed \(\Jimp\) in the admissible class, the exterior field
\(\mathcal{S}_\omega[\Jimp]\) obtained from
Eq.~\eqref{eq:ch2.outgoing-solution-operator} is unique in
\(\mathcal{F}_{\mathrm{SF}}\) \cite{stratton1941}.
\end{corollary}
\begin{proof}
Two convolutions differing by a homogeneous outgoing state violate
Theorem~\ref{thm:ch2.exterior-radiative-uniqueness} \cite{harrington2001}.
\end{proof}

\paragraph{Radiation selector idempotency.}
Subsection~\ref{sec:ch201} declared \(\Pi_{\mathrm{rad}}\) idempotent on
\(\mathcal{F}_{\mathrm{rad}}\); the joint property with outgoing closure can now be
stated precisely. On radiative states,
\begin{equation}
\label{eq:ch2.rad-idempotency-radiative}
\Pi_{\mathrm{rad}}^2=\Pi_{\mathrm{rad}}
\quad\text{on}\ \mathcal{F}_{\mathrm{rad}},
\end{equation}
as previewed in Subsection~\ref{sec:ch201} \cite{stratton1941}. After
Theorem~\ref{thm:ch2.exterior-radiative-uniqueness}, the pipeline
Eq.~\eqref{eq:ch2.outgoing-radiation-refinement} defines a single-valued map
\(\mathcal{R}_\omega=\Pi_{\mathrm{rad}}\circ\mathcal{S}_\omega\) without branch
ambiguity from duplicate outgoing decompositions \cite{harrington2001}. Evanescent
components remain in \(\ker\Pi_{\mathrm{rad}}\) through
Eq.~\eqref{eq:ch2.rad-null-on-evanescent}; uniqueness of \(\mathcal{S}_\omega\) does
not imply uniqueness of \(\mathcal{R}_\omega[\Jimp]\) when near-zone evanescent
storage dominates export \cite{stratton1941}.

\paragraph{Selector and resolvent consistency.}
Outgoing selector idempotency Eq.~\eqref{eq:ch2.outgoing-idempotency} and resolvent
selection Eq.~\eqref{eq:ch2.outgoing-resolvent-selection} are consistent with
Theorem~\ref{thm:ch2.exterior-radiative-uniqueness}:
\begin{equation}
\label{eq:ch2.outgoing-resolvent-unique}
\mathcal{O}_{\mathrm{out}}\circ\mathcal{L}_\omega^{-1}
=
\mathcal{L}_\omega^{-1}\big|_{\mathrm{phys}}
\quad\text{is unique on}\ \mathcal{Y}_\omega,
\end{equation}
because any two selections differing by a homogeneous outgoing mode would violate
Eq.~\eqref{eq:ch2.exterior-uniqueness-theorem} \cite{stratton1941,harrington2001}.
Equation~\eqref{eq:ch2.outgoing-resolvent-unique} closes the operator chain of
Figure~\ref{fig:ch2.outgoing-selector-pipeline} on exterior domains: once data are
fixed, \(\mathcal{O}_{\mathrm{out}}\), \(\mathcal{S}_\omega\), and
\(\Pi_{\mathrm{rad}}\) act on a single field orbit in \(\mathcal{F}_{\mathrm{SF}}\)
\cite{stratton1941}.

\paragraph{Bounded versus exterior contrast.}
On bounded domains with passive \(\mathcal{B}_\Gamma\), uniqueness follows from
Eq.~\eqref{eq:ch2.lax-milgram-wellposed} and
Proposition~\ref{prop:ch1.uniqueness-passive-closure} after quotienting
Eq.~\eqref{eq:ch2.coercivity-quotient} \cite{stratton1941}. At cavity eigenfrequencies,
\begin{equation}
\label{eq:ch2.exterior-vs-cavity-uniqueness}
\text{bounded PEC cavity:}\ \ker\mathcal{L}_\omega\neq\{\mathbf{0}\}
\quad\text{vs.}\quad
\text{exterior outgoing:}\ \ker\mathcal{L}_\omega|_{\mathcal{F}_{\mathrm{SF}}}=\{\mathbf{0}\},
\end{equation}
so Eq.~\eqref{eq:ch2.exterior-uniqueness-theorem} does not apply to closed resonators
where Eq.~\eqref{eq:ch2.cavity-resonance-failure} holds \cite{harrington2001,stratton1941}.
Exterior uniqueness is therefore a property of \(\mathcal{F}_{\mathrm{SF}}\) with
Sommerfeld closure, not of arbitrary \(\mathcal{B}_\Gamma\) from
Eq.~\eqref{eq:ch1.outer-boundary-classes} \cite{stratton1941}. A finite scatterer
embedded in free space may support Mie resonances in the near field without breaking
exterior uniqueness: those modes are driven by \(\Jimp\) and do not appear as free
homogeneous outgoing states \cite{jackson1999,harrington2001}.

\paragraph{Failure modes and non-applicability.}
Theorem~\ref{thm:ch2.exterior-radiative-uniqueness} requires linear passive media,
fixed \(\omega\) away from pathological embedded spectra, and exact outgoing closure
rather than approximate ABC alone \cite{stratton1941,harrington2001}. First-order
absorbing boundaries that do not enforce Eq.~\eqref{eq:ch2.sommerfeld-field-condition}
at large \(R\) can admit spurious incoming content; uniqueness is then a property of
the \emph{augmented} truncated problem, not of the physical exterior
\cite{harrington2001}. Anisotropic, moving, or nonlinear media require modified flux
functionals; time-varying \(\mathcal{B}_\Gamma(t)\) breaks the harmonic reduction of
Subsection~\ref{sec:ch223} \cite{stratton1941}. Uniqueness of \(\mathcal{S}_\omega\)
does not imply uniqueness of \(\Jimp\) in \(\mathcal{R}_\omega\)---that is
Subsection~\ref{sec:ch244} \cite{harrington2001}. Regime qualification for transport
remains Eq.~\eqref{eq:ch2.regime-operator-screen}, independent of field uniqueness
\cite{stratton1941}.

\paragraph{Worked illustration (duplicate outgoing solvers).}
Suppose two frequency-domain solvers return \((\E_1,\H_1)\) and \((\E_2,\H_2)\) for
the same \(\Jimp\) with identical near fields but opposite far-zone phase on the
propagating sheet of Eq.~\eqref{eq:ch2.outgoing-plane-wave-gate}. Their difference
violates Eq.~\eqref{eq:ch2.exterior-sommerfeld-flux-limit} unless amplitudes cancel;
Theorem~\ref{thm:ch2.exterior-radiative-uniqueness} forbids distinct outgoing
solutions \cite{stratton1941,jackson1999}. One solver has selected an incoming sheet
despite matching volume residuals; only
Eq.~\eqref{eq:ch2.outgoing-resolvent-selection} restores
Corollary~\ref{cor:ch2.outgoing-green-unique} \cite{harrington2001}.

\paragraph{Worked illustration (lossy dielectric scatterer).}
With \(\Im\{\varepsilon\}>0\) in \(\Omega_s\), Eq.~\eqref{eq:ch2.exterior-flux-identity}
acquires a negative definite contribution from \(\Re\mathcal{P}_{\Delta,s}\) for
nontrivial difference fields, strengthening the flux argument: passive loss cannot
sustain a nonzero homogeneous difference state with outgoing closure
\cite{stratton1941}. Uniqueness then holds even when near-zone fields are dominated
by evanescent content; export reporting still requires
Eq.~\eqref{eq:ch2.regime-operator-screen} \cite{harrington2001}.

\paragraph{Worked illustration (short dipole exterior).}
For \(\Jimp=\hat{\mathbf{z}}\,I_0\,\delta^{(3)}(\mathbf{r})\) in free space,
Eq.~\eqref{eq:ch2.outgoing-solution-operator} produces the unique outgoing dipole
field in \(\mathcal{F}_{\mathrm{SF}}\) \cite{jackson1999,stratton1941}. Any
homogeneous outgoing correction would add \(1/R\) transverse radiation with
\(P_{\Delta,\mathrm{rad}}>0\), contradicting
Theorem~\ref{thm:ch2.exterior-radiative-uniqueness} for zero impressed data
\cite{harrington2001}. Reactive near-zone storage does not signal non-uniqueness;
only duplicate \emph{outgoing} closures do \cite{stratton1941}.

\paragraph{Worked illustration (ABC versus Sommerfeld).}
Truncating at finite \(R\) with a first-order ABC may return two numerically distinct
fields that agree on \(\Gamma_R\) but differ in spurious incoming content reflected
between \(\Omega_s\) and the ABC \cite{harrington2001}. Those differences are not
counterexamples to Theorem~\ref{thm:ch2.exterior-radiative-uniqueness} on
\(\Omega_{\mathrm{ext}}\); they signal that the truncated BVP lacks
Eq.~\eqref{eq:ch2.sommerfeld-outgoing-class}. Silver--Müller surrogates
Eq.~\eqref{eq:ch2.silver-muller-trace-approx} improve the closure but still require
\(R\gg\lambda_0\) for uniqueness in practice \cite{stratton1941}.

\paragraph{Closing summary.}
Subsection~\ref{sec:ch243} establishes the Chapter~\ref{ch:02} HOME proof of exterior
radiative uniqueness Theorem~\ref{thm:ch2.exterior-radiative-uniqueness} via
Eqs.~\eqref{eq:ch2.exterior-difference-system}--\eqref{eq:ch2.exterior-uniqueness-theorem},
flux identities Eqs.~\eqref{eq:ch2.exterior-flux-identity}--\eqref{eq:ch2.exterior-sommerfeld-flux-limit},
Lemma~\ref{lem:ch2.zero-flux-difference}, resolvent injectivity
Eq.~\eqref{eq:ch2.exterior-resolvent-injective}, Green uniqueness
Corollary~\ref{cor:ch2.outgoing-green-unique}, radiation idempotency
Eq.~\eqref{eq:ch2.rad-idempotency-radiative}, resolvent--selector consistency
Eq.~\eqref{eq:ch2.outgoing-resolvent-unique}, and bounded/exterior contrast
Eq.~\eqref{eq:ch2.exterior-vs-cavity-uniqueness}, importing the energy method of
Subsection~\ref{sec:ch01.1.3} by reference. Subsection~\ref{sec:ch244} addresses
source reconstruction, \(\ker\mathcal{R}_\omega\), and observability limits.

\subsection{Implications for source reconstruction}
\label{sec:ch244}
% TOC tag: [HOME]

Theorem~\ref{thm:ch2.exterior-radiative-uniqueness} fixes uniqueness of
\(\mathcal{S}_\omega[\Jimp]\) on the outgoing branch. Subsection~\ref{sec:ch244}
records the Chapter~\ref{ch:02} HOME consequences for \emph{inverse} source problems:
field uniqueness does not imply source identifiability, the forward radiation map
\(\mathcal{R}_\omega\) is generally non-injective, and finite observation channels
compress rank further \cite{stratton1941,harrington2001}. Observation null-space
Eq.~\eqref{eq:ch2.observation-null-space}, rank composition
Eq.~\eqref{eq:ch2.rank-composition-law}, and non-radiating preview
Eq.~\eqref{eq:ch2.null-radiation-preview} from Subsection~\ref{sec:ch213} are
presupposed; complete classification of \(\ker\mathcal{R}_\omega\) belongs to
Section~\ref{sec:ch272} \cite{harrington2001}. Measurement taxonomy from
Section~\ref{sec:ch01.5} is invoked by reference; angular-momentum observables from
Chapter~\ref{ch:01} are not re-opened \cite{stratton1941}.

\paragraph{Field uniqueness versus source identifiability.}
Forward exterior closure yields
\begin{equation}
\label{eq:ch2.field-source-uniqueness-split}
\underbrace{\mathcal{S}_\omega[\Jimp^{(1)}]=\mathcal{S}_\omega[\Jimp^{(2)}]}_{\text{same impressed data}}
\ \Longrightarrow\
\Jimp^{(1)}=\Jimp^{(2)}
\quad\text{false in general},
\end{equation}
whereas Theorem~\ref{thm:ch2.exterior-radiative-uniqueness} gives
\begin{equation}
\label{eq:ch2.field-uniqueness-forward}
\mathcal{S}_\omega[\Jimp^{(1)}]=\mathcal{S}_\omega[\Jimp^{(2)}]
\ \Longrightarrow\
(\E_1,\H_1)=(\E_2,\H_2)
\quad\text{on}\ \mathcal{F}_{\mathrm{SF}}
\end{equation}
for the same boundary class and outgoing branch \cite{stratton1941}. Equations
\eqref{eq:ch2.field-source-uniqueness-split}--\eqref{eq:ch2.field-uniqueness-forward}
separate \emph{state} uniqueness from \emph{excitation} uniqueness: many admissible
currents can produce the same radiative exterior field once reactive degrees of
freedom are quotiented \cite{harrington2001}. Inverse problems therefore target
equivalence classes of \(\Jimp\), not individual source distributions
\cite{stratton1941}.

\paragraph{Forward radiation map and its null space.}
The radiative forward map from Subsection~\ref{sec:ch201} is
\begin{equation}
\label{eq:ch2.radiation-forward-map}
\mathcal{R}_\omega:\mathcal{J}_\Omega\longrightarrow\mathcal{F}_{\mathrm{rad}},
\qquad
\mathcal{R}_\omega[\Jimp]=\Pi_{\mathrm{rad}}\,\mathcal{S}_\omega[\Jimp],
\end{equation}
restating Eq.~\eqref{eq:ch2.radiation-operator-def} after outgoing closure
Eqs.~\eqref{eq:ch2.outgoing-radiation-refinement}--\eqref{eq:ch2.outgoing-resolvent-unique}
\cite{stratton1941}. The source-side null space previewed in
Eq.~\eqref{eq:ch2.null-radiation-preview} is
\begin{equation}
\label{eq:ch2.ker-R-omega-preview}
\ker\mathcal{R}_\omega
=
\left\{
\Jimp\in\mathcal{J}_\Omega:
\Pi_{\mathrm{rad}}\,\mathcal{S}_\omega[\Jimp]=\mathbf{0}
\right\},
\end{equation}
generally infinite-dimensional on open domains \cite{harrington2001,stratton1941}.
Section~\ref{sec:ch272} owns the complete null-space analysis; Subsection~\ref{sec:ch244}
records the reconstruction consequences of Eq.~\eqref{eq:ch2.ker-R-omega-preview}
without constructing non-radiating bases \cite{stratton1941}. Elements of
\(\ker\mathcal{R}_\omega\) are invisible to far-zone radiative observables yet may
carry substantial near-zone reactive energy through \(\mathcal{S}_\omega\)
\cite{harrington2001}.

\begin{figure}[t]
\centering
\ChOneFigBlock{%
\begin{tikzpicture}[ch1fig, node distance=7mm and 10mm]
  \node[ch1box] (j) {$\mathcal{J}_\Omega$};
  \node[ch1box, right=14mm of j] (s) {$\mathcal{S}_\omega$};
  \node[ch1box, right=14mm of s] (f) {$\mathcal{F}_{\mathrm{SF}}$};
  \node[ch1box, right=14mm of f] (r) {$\mathcal{F}_{\mathrm{rad}}$};
  \node[ch1box, below=12mm of j] (ker) {$\ker\mathcal{R}_\omega$};
  \draw[ch1flow] (j) -- (s) -- (f) -- node[above,font=\scriptsize] {$\Pi_{\mathrm{rad}}$} (r);
  \draw[ch1flow, dashed] (ker) -- (j);
\end{tikzpicture}%
}
\caption{Forward chain for reconstruction: \(\mathcal{S}_\omega\) is single-valued on
\(\mathcal{F}_{\mathrm{SF}}\) after Theorem~\ref{thm:ch2.exterior-radiative-uniqueness};
\(\mathcal{R}_\omega\) collapses many sources to the same radiative output.}
\label{fig:ch2.reconstruction-forward-chain}
\end{figure}

Figure~\ref{fig:ch2.reconstruction-forward-chain} contrasts injectivity of
\(\mathcal{S}_\omega|_{\mathcal{F}_{\mathrm{SF}}}\) with non-injectivity of
\(\mathcal{R}_\omega\) on \(\mathcal{J}_\Omega\) \cite{stratton1941}.

\paragraph{Observation composition and inverse problem.}
Laboratory readouts compose with the forward map. From
Eq.~\eqref{eq:ch2.radiation-operator-composition} and finite-rank observation
Eq.~\eqref{eq:ch2.finite-rank-observation},
\begin{equation}
\label{eq:ch2.reconstruction-observation-map}
\mathcal{M}_\omega
=
\mathcal{O}_{\mathrm{meas}}\circ\mathcal{O}_{\mathrm{far}}\circ\mathcal{R}_\omega
=
\mathcal{O}_{\mathrm{meas}}\circ\Lambda_{\mathrm{rad}}:
\mathcal{J}_\Omega\longrightarrow\mathbb{C}^{N},
\end{equation}
with \(\Lambda_{\mathrm{rad}}=\mathcal{O}_{\mathrm{far}}\circ\mathcal{R}_\omega\)
\cite{stratton1941,harrington2001}. Given measurement vector
\(\mathbf{y}\in\mathbb{C}^{N}\), the inverse source problem seeks
\(\Jimp\in\mathcal{J}_\Omega\) such that
\begin{equation}
\label{eq:ch2.inverse-source-equation}
\mathcal{M}_\omega[\Jimp]=\mathbf{y},
\qquad
\mathcal{D}_{J}[\Jimp]=1,
\end{equation}
subject to admissibility Eq.~\eqref{eq:ch2.source-admissibility}
\cite{harrington2001}. Equation~\eqref{eq:ch2.inverse-source-equation} is ill-posed
whenever \(\ker\mathcal{M}_\omega\neq\{\mathbf{0}\}\), which occurs if either
\(\ker\mathcal{R}_\omega\neq\{\mathbf{0}\}\) or \(\ker\mathcal{O}_{\mathrm{meas}}\)
is nontrivial on \(\operatorname{ran}\mathcal{R}_\omega\)
\cite{stratton1941}. Rank cap Eq.~\eqref{eq:ch2.rank-composition-law} applies directly:
\begin{equation}
\label{eq:ch2.reconstruction-rank-cap}
\mathrm{rank}(\mathcal{M}_\omega)
\le
\min\big\{N,\,\dim\mathcal{R}(\mathcal{R}_\omega)\big\},
\end{equation}
so no measurement protocol with \(N\) channels resolves more than \(N\) independent
radiative degrees of freedom \cite{harrington2001,stratton1941}.

\paragraph{Reconstruction null space and equivalence classes.}
The composed null space satisfies
\begin{equation}
\label{eq:ch2.reconstruction-null-space}
\ker\mathcal{M}_\omega
\supseteq
\ker\mathcal{R}_\omega
\ \cup\
\mathcal{R}_\omega^{-1}\!\big(\ker\mathcal{O}_{\mathrm{meas}}\big),
\end{equation}
with set inclusion on admissible sources \cite{stratton1941}. Identifiable sources
modulo null space form the quotient
\begin{equation}
\label{eq:ch2.source-identifiability-quotient}
\mathcal{J}_\Omega\big/\ker\mathcal{M}_\omega
\ \cong\
\operatorname{ran}\mathcal{M}_\omega
\subset\mathbb{C}^{N},
\end{equation}
so reconstruction recovers at best a finite-dimensional image, not a unique
\(\Jimp(\mathbf{r})\) \cite{harrington2001}. Equivalent-current models that expand
\(\Jimp\) on auxiliary surfaces inherit the same non-uniqueness: distinct surface
distributions related by \(\ker\mathcal{R}_\omega\) produce identical far-zone
patterns \cite{stratton1941,harrington2001}. Boundary-element inversion therefore
requires regularization or additional constraints beyond Maxwell consistency alone
\cite{harrington2001}.

\begin{figure}[t]
\centering
\ChOneFigBlock{%
\begin{tikzpicture}[ch1fig, node distance=7mm and 9mm]
  \node[ch1box] (j1) {$\Jimp^{(1)}$};
  \node[ch1box, right=12mm of j1] (j2) {$\Jimp^{(2)}$};
  \node[ch1box, right=14mm of j2] (plus) {$\Jimp^{(1)}+\delta\Jimp$};
  \node[ch1box, below=11mm of plus] (y) {Same $\mathbf{y}$};
  \node[ch1box, below=11mm of j1] (delta) {$\delta\Jimp\in\ker\mathcal{M}_\omega$};
  \draw[ch1flow] (j1) -- (plus);
  \draw[ch1flow] (delta) -- (plus);
  \draw[ch1flow] (plus) -- (y);
\end{tikzpicture}%
}
\caption{Non-uniqueness: adding \(\delta\Jimp\in\ker\mathcal{M}_\omega\) leaves
measurement vector \(\mathbf{y}\) unchanged.}
\label{fig:ch2.reconstruction-null-class}
\end{figure}

Figure~\ref{fig:ch2.reconstruction-null-class} is the inverse image of
Eq.~\eqref{eq:ch2.observation-null-space} pulled back to source space through
Eq.~\eqref{eq:ch2.reconstruction-observation-map} \cite{stratton1941}.

\paragraph{Regularized reconstruction.}
When Eq.~\eqref{eq:ch2.inverse-source-equation} is ill-posed, stable inversion uses
a penalized least-squares or Tikhonov template on admissible sources:
\begin{equation}
\label{eq:ch2.tikhonov-source-reconstruction}
{\Jimp}_{\mathrm{reg}}
=
\arg\min_{\Jimp\in\mathcal{J}_\Omega}
\left\{
\|\mathcal{M}_\omega[\Jimp]-\mathbf{y}\|^{2}
+\alpha\,\mathcal{P}[\Jimp]
\right\},
\end{equation}
with regularization parameter \(\alpha>0\) and penalty
\(\mathcal{P}\) enforcing smoothness, minimum norm, or sparsity on an equivalent
surface \cite{harrington2001,stratton1941}. Equation
\eqref{eq:ch2.tikhonov-source-reconstruction} does not remove
\(\ker\mathcal{M}_\omega\); it selects one representative from the affine coset
\(\mathcal{M}_\omega^{-1}(\mathbf{y})\) \cite{stratton1941}. The choice of
\(\mathcal{P}\) is physics-adjacent modeling, not a Maxwell identity; different
penalties yield different \({\Jimp}_{\mathrm{reg}}\) with identical radiative
predictions when \(\delta\Jimp\in\ker\mathcal{R}_\omega\) \cite{harrington2001}.

\paragraph{Regime qualification for inverse claims.}
Transport and angular reports attach to \(\mathcal{R}_\omega[\Jimp]\), not to raw
\(\mathcal{S}_\omega[\Jimp]\), per Eq.~\eqref{eq:ch2.transport-qualified-radiation}
\cite{stratton1941}. Regime screen Eq.~\eqref{eq:ch2.regime-operator-screen} gates
whether \(\mathcal{R}_\omega[\Jimp]\) lies in the propagating sector on which
far-zone observables were designed:
\begin{equation}
\label{eq:ch2.reconstruction-regime-gate}
\mathcal{D}_{J,\mathrm{reg}}=1
\ \Longrightarrow\
\text{inverse claims target}\ \mathcal{R}_\omega[\Jimp],
\ \text{not reactive near fields},
\end{equation}
with \(\mathcal{D}_{J,\mathrm{reg}}\) from Eq.~\eqref{eq:ch1.regime-qualified-decision}
\cite{stratton1941,harrington2001}. Attempting to reconstruct \(\Jimp\) from
near-zone reactive samples without regime qualification conflates storage with export,
the failure mode flagged in Section~\ref{sec:ch01.2.4} \cite{stratton1941}. Outgoing
uniqueness from Theorem~\ref{thm:ch2.exterior-radiative-uniqueness} stabilizes the
\emph{field} side of Eq.~\eqref{eq:ch2.inverse-source-equation}; it does not promote
evanescent-dominated states to identifiable radiative sources
\cite{harrington2001}.

\paragraph{Adjoint and back-projection template.}
Formal adjoint back-projection provides a reconstruction \emph{template} without
claiming uniqueness:
\begin{equation}
\label{eq:ch2.adjoint-backprojection}
\widetilde{\Jimp}
=
\mathcal{M}_\omega^{\dagger}\,\mathbf{y},
\end{equation}
where \(\mathcal{M}_\omega^{\dagger}\) is the Hilbert-space adjoint on declared
domains \cite{stratton1941,harrington2001}. Equation
\eqref{eq:ch2.adjoint-backprojection} is useful for sensitivity analysis and
iterative inversion but generally satisfies
\(\mathcal{M}_\omega[\widetilde{\Jimp}]\neq\mathbf{y}\) when
\(\ker\mathcal{M}_\omega\neq\{\mathbf{0}\}\) \cite{harrington2001}. Consistent
reconstruction requires projecting onto \(\operatorname{ran}\mathcal{M}_\omega^{\dagger}\)
or applying Eq.~\eqref{eq:ch2.tikhonov-source-reconstruction}
\cite{stratton1941}. Full integral representations for
Eq.~\eqref{eq:ch2.adjoint-backprojection} belong to Section~\ref{sec:ch251}; they are
not repeated here \cite{harrington2001}.

\paragraph{Equivalent sources and Huygens surfaces.}
Equivalent-current inversion on a closed surface \(\Gamma_{\mathrm{eq}}\) seeks
\(\mathbf{J}_{\mathrm{eq}}\) on \(\Gamma_{\mathrm{eq}}\) such that exterior fields
match measured data \cite{stratton1941,harrington2001}. Uniqueness of exterior
\(\mathcal{S}_\omega\) from Theorem~\ref{thm:ch2.exterior-radiative-uniqueness}
fixes the \emph{field} on \(\mathcal{F}_{\mathrm{SF}}\) but not
\(\mathbf{J}_{\mathrm{eq}}\): any non-radiating tangential distribution on
\(\Gamma_{\mathrm{eq}}\) can be added without changing far-zone samples
\cite{harrington2001}. The operator form is
\begin{equation}
\label{eq:ch2.equivalent-source-ambiguity}
\mathbf{J}_{\mathrm{eq},1}-\mathbf{J}_{\mathrm{eq},2}
\in
\ker\mathcal{R}_\omega\big|_{\Gamma_{\mathrm{eq}}},
\end{equation}
with restriction to surface currents in the trace class of
Subsection~\ref{sec:ch211} \cite{stratton1941}. Numerical equivalence in
Eq.~\eqref{eq:ch2.discrete-green-equivalence} inherits the same null-space structure
at finite matrix rank \cite{harrington2001}.

\paragraph{Near-field versus far-field data.}
Near-field scanning samples \(\mathcal{S}_\omega[\Jimp]\) on a surface enclosing
\(\Omega_s\), not \(\mathcal{R}_\omega[\Jimp]\) alone \cite{stratton1941}. Because
\(\mathcal{S}_\omega\) is injective on outgoing fields but many sources share the
same \(\mathcal{R}_\omega\) image, near-field data can \emph{reduce} null-space
dimension relative to far-zone-only inversion without eliminating
\(\ker\mathcal{R}_\omega\) in general \cite{harrington2001}. Outgoing closure still
required: two incoming/outgoing superpositions with identical near traces on a finite
surface need not agree in the far zone unless Theorem~\ref{thm:ch2.exterior-radiative-uniqueness}
applies to the full exterior problem \cite{stratton1941,jackson1999}. Near-field
inversion therefore couples Subsection~\ref{sec:ch243} field uniqueness with
Subsection~\ref{sec:ch244} source non-uniqueness \cite{harrington2001}.

\paragraph{Boundary of validity.}
Equations~\eqref{eq:ch2.inverse-source-equation}--\eqref{eq:ch2.tikhonov-source-reconstruction}
require linear Maxwell theory, fixed \(\omega\), admissible
Eq.~\eqref{eq:ch2.source-admissibility}, and outgoing forward closure
Eq.~\eqref{eq:ch2.outgoing-resolvent-unique} \cite{stratton1941,harrington2001}.
Nonlinear, time-varying, or dispersive inversion breaks the harmonic reduction of
Subsection~\ref{sec:ch223}. Single-frequency inversion cannot resolve broadband
excitation without superposition across \(\omega\) \cite{stratton1941}. Loss improves
conditioning of Eq.~\eqref{eq:ch2.inverse-source-equation} by damping resonances but
does not make \(\mathcal{R}_\omega\) injective on all of \(\mathcal{J}_\Omega\)
\cite{harrington2001}. Complete null-space bases and NR source examples belong to
Section~\ref{sec:ch27} \cite{stratton1941}.

\paragraph{Worked illustration (pattern-matched equivalent currents).}
Two distinct surface-current expansions on a Huygens sphere produce identical
far-zone \(\mathcal{R}_\omega[\Jimp]\) when their difference lies in
Eq.~\eqref{eq:ch2.ker-R-omega-preview} \cite{harrington2001,stratton1941}. A
method-of-moments fit that minimizes \(\|\mathcal{M}_\omega[\Jimp]-\mathbf{y}\|\)
without penalty returns one of infinitely many equivalent solutions; changing
\(\mathcal{P}\) in Eq.~\eqref{eq:ch2.tikhonov-source-reconstruction} shifts
\({\Jimp}_{\mathrm{reg}}\) while leaving \(\mathbf{y}\) fixed \cite{harrington2001}.

\paragraph{Worked illustration (aperture-limited pattern inversion).}
With \(N\) probe channels, Eq.~\eqref{eq:ch2.reconstruction-rank-cap} limits
recoverable modal content to rank \(N\) regardless of mesh density
\cite{stratton1941}. Refining a boundary-element mesh without adding independent
measurements does not resolve components in
\(\ker\mathcal{O}_{\mathrm{meas}}\circ\Lambda_{\mathrm{rad}}\)
\cite{harrington2001}. This is the operator reading of
Eq.~\eqref{eq:ch2.observation-null-space}, distinct from discretization error alone
\cite{stratton1941}.

\paragraph{Worked illustration (reactive near-field trap).}
A source with \(\mathcal{D}_{J,\mathrm{reg}}=0\) may produce strong near-zone
\(\mathcal{S}_\omega[\Jimp]\) while \(\mathcal{R}_\omega[\Jimp]\approx\mathbf{0}\)
\cite{stratton1941}. Inverting near samples without
Eq.~\eqref{eq:ch2.reconstruction-regime-gate} assigns spurious radiative labels to
storage-dominated fields \cite{harrington2001}. Regime screen
Eq.~\eqref{eq:ch2.regime-operator-screen} must precede angular-transport claims on
reconstructed sources \cite{stratton1941}.

\paragraph{Closing summary.}
Subsection~\ref{sec:ch244} establishes the Chapter~\ref{ch:02} HOME reconstruction
record via field/source split
Eqs.~\eqref{eq:ch2.field-source-uniqueness-split}--\eqref{eq:ch2.field-uniqueness-forward},
forward map Eq.~\eqref{eq:ch2.radiation-forward-map}, null-space preview
Eq.~\eqref{eq:ch2.ker-R-omega-preview}, observation composition
Eq.~\eqref{eq:ch2.reconstruction-observation-map}, inverse equation
Eq.~\eqref{eq:ch2.inverse-source-equation}, rank cap
Eq.~\eqref{eq:ch2.reconstruction-rank-cap}, reconstruction null space
Eq.~\eqref{eq:ch2.reconstruction-null-space}, identifiability quotient
Eq.~\eqref{eq:ch2.source-identifiability-quotient}, Tikhonov template
Eq.~\eqref{eq:ch2.tikhonov-source-reconstruction}, regime gate
Eq.~\eqref{eq:ch2.reconstruction-regime-gate}, adjoint template
Eq.~\eqref{eq:ch2.adjoint-backprojection}, and equivalent-source ambiguity
Eq.~\eqref{eq:ch2.equivalent-source-ambiguity}, separating Theorem~\ref{thm:ch2.exterior-radiative-uniqueness}
field uniqueness from source non-identifiability \cite{stratton1941,harrington2001}.

\medskip
\noindent
Section~\ref{sec:ch24} closes with the operator chain fixed for open regions:
Sommerfeld asymptotics and outgoing class
Eqs.~\eqref{eq:ch2.sommerfeld-outgoing-class}--\eqref{eq:ch2.sommerfeld-field-condition}
(Subsection~\ref{sec:ch241}), outgoing selector and resolvent branch
Eqs.~\eqref{eq:ch2.outgoing-selector-def}--\eqref{eq:ch2.outgoing-solution-operator}
(Subsection~\ref{sec:ch242}), exterior radiative uniqueness
Theorem~\ref{thm:ch2.exterior-radiative-uniqueness} and resolvent injectivity
Eq.~\eqref{eq:ch2.exterior-resolvent-injective} (Subsection~\ref{sec:ch243}), and
inverse reconstruction limits
Eqs.~\eqref{eq:ch2.inverse-source-equation}--\eqref{eq:ch2.ker-R-omega-preview}
(Subsection~\ref{sec:ch244}) \cite{stratton1941,harrington2001,jackson1999}. Green
kernels from Section~\ref{sec:ch23} are now tied to a single physical branch;
\(\mathcal{S}_\omega\) is single-valued on \(\mathcal{F}_{\mathrm{SF}}\), while
\(\mathcal{R}_\omega\) and composed measurement maps remain rank-deficient on
\(\mathcal{J}_\Omega\) \cite{harrington2001}. Section~\ref{sec:ch25} evaluates
radiation integrals, flux meters, and near/intermediate/far partitions on the
qualified propagating subspace; Section~\ref{sec:ch27} completes non-radiating
null-space classification \cite{stratton1941}. Regime qualification
Eq.~\eqref{eq:ch2.regime-operator-screen} and transport gate
Eq.~\eqref{eq:ch2.transport-qualified-radiation} remain mandatory before
angular-transport claims on operator outputs \cite{stratton1941,harrington2001}.

\section{Radiation Integrals and Source-to-Field Mappings}
\label{sec:ch25}

Section~\ref{sec:ch24} fixed outgoing branch selection, exterior radiative uniqueness, and
the limits of source reconstruction on \(\mathcal{J}_\Omega\). The operator chain
\(\Jimp\mapsto\mathcal{S}_\omega[\Jimp]\mapsto\mathcal{R}_\omega[\Jimp]\) is therefore
single-valued on \(\mathcal{F}_{\mathrm{SF}}\) at the field level, yet rank-deficient at
the source level through Eq.~\eqref{eq:ch2.ker-R-omega-preview}
\cite{stratton1941,harrington2001}. Section~\ref{sec:ch25} supplies the Chapter~\ref{ch:02}
\emph{evaluation layer} for that chain: explicit source-to-field integral representations
on the outgoing Green branch, distance-regime asymptotics that separate reactive from
radiative spatial scales, flux and operator-norm meters on qualified radiative outputs, and
the radiative--non-radiative partition aligned with propagation projectors from
Subsection~\ref{sec:ch224} \cite{stratton1941,jackson1999}. Poynting flux, momentum
budgets, and transport indices from Subsections~\ref{sec:ch01.1.4},~\ref{sec:ch01.1.5},
and~\ref{sec:ch01.2.3} are presupposed; Section~\ref{sec:ch25} restates them as operator
evaluations on \(\mathcal{R}_\omega[\Jimp]\) without re-opening their HOME proofs in
Chapter~\ref{ch:01} \cite{stratton1941}. Dyadic kernel construction and Sommerfeld closure
remain in Sections~\ref{sec:ch23}--\ref{sec:ch24}; vector spherical harmonic separation
belongs to Chapter~\ref{ch:03}.

The problem addressed here is \emph{observable evaluation} after closure. Green convolution
Eq.~\eqref{eq:ch2.electric-dyadic-convolution} and outgoing solution operator
Eq.~\eqref{eq:ch2.outgoing-solution-operator} return field pairs on
\(\mathcal{F}_\Omega\); laboratory and far-zone claims attach to
\(\mathcal{R}_\omega[\Jimp]=\Pi_{\mathrm{rad}}\,\mathcal{S}_\omega[\Jimp]\) per
Eqs.~\eqref{eq:ch2.radiation-operator-def} and~\eqref{eq:ch2.outgoing-radiation-refinement}
\cite{stratton1941,harrington2001}. The preview integral
Eq.~\eqref{eq:ch2.radiation-integral-preview} from Subsection~\ref{sec:ch201} named the
retarded map without far-zone asymptotics; Section~\ref{sec:ch25} upgrades that preview to
evaluable radiation integrals on the outgoing dyad
\(\overline{\mathbf{G}}_{E,\omega}^{\mathrm{(out)}}\) from
Eq.~\eqref{eq:ch2.outgoing-green-branch} once \(\mathbf{r}\) lies in the intermediate or
far zone \cite{jackson1999,stratton1941}. Power screening through
Eq.~\eqref{eq:ch2.radiated-power-screen} stabilizes only on outputs that pass regime gate
Eq.~\eqref{eq:ch2.regime-operator-screen} and transport qualification
Eq.~\eqref{eq:ch2.transport-qualified-radiation} \cite{stratton1941}.

\paragraph{Evaluation chain.}
Let \(\Omega_s=\operatorname{supp}\Jimp\) be compact with characteristic scale \(a\),
vacuum wavenumber \(k_0=\omega\sqrt{\varepsilon_0\mu_0}\), and wavelength
\(\lambda_0=2\pi/k_0\). Section~\ref{sec:ch25} organizes evaluation in four stages:
\begin{equation}
\label{eq:ch2.section25-evaluation-chain}
\Jimp
\ \xrightarrow{\ \mathcal{S}_\omega\ }
(\E,\H)\in\mathcal{F}_{\mathrm{SF}}
\ \xrightarrow{\ \Pi_{\mathrm{rad}}\ }
\mathcal{R}_\omega[\Jimp]
\ \xrightarrow{\ \mathcal{F}_{\mathrm{flux}}\ }
\text{power, pattern, norm reports},
\end{equation}
where \(\mathcal{F}_{\mathrm{flux}}\) denotes the flux and operator-norm functionals
developed in Subsection~\ref{sec:ch253} by reference to Chapter~\ref{ch:01} flux identities
\cite{stratton1941}. Equation~\eqref{eq:ch2.section25-evaluation-chain} is the section-level
restatement of Eq.~\eqref{eq:ch2.pipeline-composition}: stage (iii) export must precede
stage (iv) metering \cite{harrington2001}. Applying \(\mathcal{F}_{\mathrm{flux}}\) to
\(\mathcal{S}_\omega[\Jimp]\) without \(\Pi_{\mathrm{rad}}\) imports reactive circulation
into quantities designed for outgoing transverse radiation
(Section~\ref{sec:ch01.2.4}) \cite{stratton1941}.

The outgoing integral realization of \(\mathcal{R}_\omega\) for
\(\mathbf{r}\notin\Omega_s\) is
\begin{equation}
\label{eq:ch2.radiation-integral-outgoing-upgrade}
\E_{\mathrm{rad}}(\mathbf{r})
=
-\jj\omega\mu_0
\int_{\Omega_s}
\overline{\mathbf{G}}_{E,\omega}^{\mathrm{(out)}}(\mathbf{r},\mathbf{r}')\cdot\Jimp(\mathbf{r}')\,d^{3}r',
\qquad
\H_{\mathrm{rad}}
=
-\frac{1}{\jj\omega\mu_0}\,\curl\E_{\mathrm{rad}},
\end{equation}
with \(\overline{\mathbf{G}}_{E,\omega}^{\mathrm{(out)}}\) from
Eq.~\eqref{eq:ch2.outgoing-green-branch} and curl taken in the observation variable
\(\mathbf{r}\) \cite{stratton1941,jackson1999}. Equation
\eqref{eq:ch2.radiation-integral-outgoing-upgrade} upgrades
Eq.~\eqref{eq:ch2.radiation-integral-preview} after Sections~\ref{sec:ch23}--\ref{sec:ch24};
Subsection~\ref{sec:ch251} gives the Chapter~\ref{ch:02} HOME for magnetic couplings,
scalar-potential routes, and equivalent surface formulations that reduce to the same
\(\mathcal{R}_\omega\) image \cite{harrington2001}. Near coincidence
\(\mathbf{r}\to\Omega_s\) requires regularization from Subsection~\ref{sec:ch232}; Section~\ref{sec:ch25}
evaluates Eq.~\eqref{eq:ch2.radiation-integral-outgoing-upgrade} only where the observation
point lies outside the source support or on declared equivalent surfaces with trace
matching \cite{stratton1941}.

\begin{figure}[t]
\centering
\ChOneFigBlock{%
\begin{tikzpicture}[ch1fig, node distance=7mm and 10mm]
  \node[ch1box] (int) {Integral\\representations};
  \node[ch1box, right=13mm of int] (reg) {Near / int / far\\regimes};
  \node[ch1box, right=13mm of reg] (flx) {Flux \&\\operator norms};
  \node[ch1box, right=13mm of flx] (split) {Rad / NR\\components};
  \node[ch1box, below=12mm of int] (gr) {$\overline{\mathbf{G}}_{E,\omega}^{\mathrm{(out)}}$};
  \node[ch1box, below=12mm of reg] (kr) {$k_0 r$ scales};
  \draw[ch1flow] (gr) -- (int) -- (reg) -- (flx) -- (split);
  \draw[ch1flow] (kr) -- (reg);
\end{tikzpicture}%
}
\caption{Section~\ref{sec:ch25} development path: outgoing integral representations
(Subsection~\ref{sec:ch251}), distance-regime asymptotics (Subsection~\ref{sec:ch252}),
flux and norm meters on \(\mathcal{R}_\omega[\Jimp]\) (Subsection~\ref{sec:ch253}), and
radiative--non-radiative split (Subsection~\ref{sec:ch254}).}
\label{fig:ch2.radiation-integrals-roadmap}
\end{figure}

Figure~\ref{fig:ch2.radiation-integrals-roadmap} orders the four subsections.
Subsection~\ref{sec:ch251} owns the Chapter~\ref{ch:02} HOME for source-to-field integral
representations that realize \(\mathcal{S}_\omega\) and \(\mathcal{R}_\omega\) on the
outgoing branch, including magnetic current couplings and Huygens-equivalent surface
forms when they avoid duplicating Subsection~\ref{sec:ch231} kernel tabulations
\cite{stratton1941,harrington2001}. Subsection~\ref{sec:ch252} supplies distance-regime
asymptotics: reactive \(1/r^{2}\), \(1/r^{3}\) near-zone dominance versus leading
\(1/r\) transverse radiation in the far zone (Eq.~\eqref{eq:ch2.far-zone-transverse-radiation}),
with explicit \(k_0 r\) thresholds that qualify Eq.~\eqref{eq:ch2.regime-operator-screen}
\cite{stratton1941,jackson1999}. Subsection~\ref{sec:ch253} imports flux and operator-norm
templates from Chapter~\ref{ch:01} as evaluations on \(\mathcal{R}_\omega[\Jimp]\), including
spherical-shell power integrals compatible with Eq.~\eqref{eq:ch2.radiated-power-screen}
\cite{stratton1941}. Subsection~\ref{sec:ch254} applies propagating and evanescent projectors
Eqs.~\eqref{eq:ch2.spectral-projectors}--\eqref{eq:ch2.rad-null-on-evanescent} to separate
radiative export from non-radiative storage, inheriting Section~\ref{sec:ch01.2.4} and
Subsection~\ref{sec:ch224} by reference \cite{harrington2001}.

\paragraph{Distance regimes (preview).}
For observation radius \(r=|\mathbf{r}|\) from a reference point in \(\Omega_s\), define
\begin{equation}
\label{eq:ch2.distance-regime-preview}
\begin{aligned}
\text{near:}\ & k_0 r\lesssim 1,\\
\text{intermediate:}\ & 1\lesssim k_0 r\lesssim (k_0 a)^{-1},\\
\text{far:}\ & k_0 r\gg k_0 a,\quad r\gg a,
\end{aligned}
\end{equation}
with implicit constants fixed in Subsection~\ref{sec:ch252} \cite{stratton1941,jackson1999}.
Equations~\eqref{eq:ch2.distance-regime-preview} are scale labels, not sharp cutoffs:
reactive and radiative terms coexist in the near zone, and intermediate zones require
explicit asymptotic bookkeeping before Eq.~\eqref{eq:ch2.radiation-integral-outgoing-upgrade}
reduces to transverse \(1/r\) radiation \cite{harrington2001}. Subsection~\ref{sec:ch252}
states the Chapter~\ref{ch:02} HOME for field magnitude scalings and for the observation
surfaces on which each regime is valid \cite{stratton1941}. Pattern and gain definitions
used in antenna theory attach only after the far-zone reduction in
Eq.~\eqref{eq:ch2.far-zone-transverse-radiation} \cite{harrington2001,jackson1999}.

\paragraph{Flux meters on the radiative branch.}
Complex Poynting flux from Eq.~\eqref{eq:ch1.poynting-complex} applies to any admissible
field pair, but power budgets for transport claims use the radiative branch:
\begin{equation}
\label{eq:ch2.flux-on-radiative-branch}
P_{\mathrm{rad}}(R)
=
\Re\oint_{\Gamma_R}\hat{\mathbf{r}}\cdot\mathbf{S}\big[\mathcal{R}_\omega[\Jimp]\big]\,dS,
\end{equation}
with \(\mathbf{S}[\cdot]\) the Poynting vector functional imported from
Subsection~\ref{sec:ch01.1.4} and \(\Gamma_R=\{r=R\}\) large enough for far-zone reduction
when reporting export \cite{stratton1941}. Equation~\eqref{eq:ch2.flux-on-radiative-branch}
specializes Eq.~\eqref{eq:ch2.radiated-power-screen} to the outgoing selector chain; it is
not equivalent to integrating \(\mathbf{S}[\mathcal{S}_\omega[\Jimp]]\) unless reactive
content has been shown negligible or orthogonal to the declared observable
\cite{harrington2001}. Subsection~\ref{sec:ch253} develops momentum and angular-momentum
flux operators on the same qualified subspace, citing
Eqs.~\eqref{eq:ch1.poynting-complex}--\eqref{eq:ch1.maxwell-stress} and transport index
Eq.~\eqref{eq:ch1.transport-admissibility-index} without re-proving conservation laws
\cite{stratton1941}.

Several constructions deliberately remain outside Section~\ref{sec:ch25}. Vector spherical
harmonic expansions of far fields belong to Chapter~\ref{ch:03}; spectral decomposition and
Fredholm index of radiation operators to Section~\ref{sec:ch26}; complete non-radiating
null-space classification to Section~\ref{sec:ch27}; reciprocity identities to
Section~\ref{sec:ch28} \cite{stratton1941,harrington2001}. Section~\ref{sec:ch25} does not
re-derive Green kernels, Sommerfeld conditions, or exterior uniqueness; it evaluates
observables on the branch fixed in Sections~\ref{sec:ch23}--\ref{sec:ch24}. Angular-spectrum
filters and plane-wave action operators belong to Section~\ref{sec:ch29}; here only the
integral representations they act upon are constructed \cite{stratton1941}.

Regime logic from Section~\ref{sec:ch01.2.4} is unchanged. Uniqueness of
\(\mathcal{S}_\omega\) on \(\mathcal{F}_{\mathrm{SF}}\) does not promote evanescent near
fields to transport channels; Eq.~\eqref{eq:ch2.regime-operator-screen} remains the
measurement gate after outgoing closure \cite{stratton1941,harrington2001}. Source
reconstruction limits from Subsection~\ref{sec:ch244} apply to any inversion built from
far-zone samples of Eq.~\eqref{eq:ch2.radiation-integral-outgoing-upgrade}: finite
rank of \(\mathcal{M}_\omega\) caps identifiable content even when forward integrals are
exact \cite{harrington2001}. Section~\ref{sec:ch25} therefore addresses \emph{forward}
metering; inverse identifiability remains rank-deficient as recorded in
Eqs.~\eqref{eq:ch2.reconstruction-rank-cap}--\eqref{eq:ch2.ker-R-omega-preview}
\cite{stratton1941}.

The split between \(\Pi_{\mathrm{prop}}\), \(\Pi_{\mathrm{ev}}\), and \(\Pi_{\mathrm{rad}}\)
from Subsection~\ref{sec:ch224} governs which integral components enter
Eq.~\eqref{eq:ch2.flux-on-radiative-branch}. Evanescent modes may satisfy Sommerfeld
closure at infinity yet carry no far-zone transverse export through
Eq.~\eqref{eq:ch2.rad-null-on-evanescent} \cite{stratton1941,harrington2001}. Subsection~\ref{sec:ch254}
operatorizes that split on integral representations without re-displaying spectral
projector algebra from Subsection~\ref{sec:ch224} \cite{harrington2001}. Reactive
near-zone circulation flagged in Section~\ref{sec:ch01.2.4} therefore remains in
\(\mathcal{S}_\omega[\Jimp]\) until explicitly projected; it is not removed by outgoing
selection alone when evanescent storage dominates local measurements
\cite{stratton1941}.

Readers approaching from numerical electromagnetics should map Section~\ref{sec:ch25} to
familiar far-field routines. Subsection~\ref{sec:ch251} is the continuous target behind
method-of-moments far-field post-processing and near-field-to-far-field transforms
\cite{harrington2001}. Subsection~\ref{sec:ch252} is the analytical basis for reactive
versus radiative distance labels in antenna range definitions \cite{stratton1941,jackson1999}.
Subsection~\ref{sec:ch253} is the operator form of integrated Poynting power and gain
metrics \cite{stratton1941}. Subsection~\ref{sec:ch254} is the frequency-domain counterpart
of propagating-mode filtering before pattern reports \cite{harrington2001}. None of those
post-processors is derived here; Section~\ref{sec:ch25} fixes the continuous maps they
must approximate on the outgoing branch \cite{stratton1941}.

Loss and dispersive materials modify kernel symbols at fixed \(\omega\) without changing
the evaluation template: Eq.~\eqref{eq:ch2.radiation-integral-outgoing-upgrade} uses
\(\overline{\mathbf{G}}_{E,\omega}^{\mathrm{(out)}}\) evaluated with complex
\(k(\mathbf{r})\) from Eq.~\eqref{eq:ch2.helmholtz-symbol} \cite{stratton1941}. Passive
dissipation shifts poles and damps resonances but does not replace regime qualification
when near-zone evanescent content dominates \cite{harrington2001}. Time-domain broadband
metering requires superposition across \(\omega\) and is deferred to the causal chain of
Subsection~\ref{sec:ch223}; Section~\ref{sec:ch25} works at fixed harmonic frequency
\cite{stratton1941}.

Charge continuity Eq.~\eqref{eq:ch1.charge-continuity-harmonic} and admissibility
Eq.~\eqref{eq:ch2.source-admissibility} remain in force during integral evaluation: source
distributions that violate \(\divg\Jimp=-\jj\omega\rho\) produce formal convolutions outside
\(\mathcal{J}_\Omega\) even when Eq.~\eqref{eq:ch2.radiation-integral-outgoing-upgrade}
returns finite numbers \cite{stratton1941}. Subsection~\ref{sec:ch251} states integral
identities in divergence-compatible form; equivalent surface currents must match tangential
fields through trace classes from Subsection~\ref{sec:ch211} \cite{harrington2001}.

\paragraph{Worked illustration (compact dipole ranges).}
For \(\Jimp\) supported in a ball of radius \(a\ll\lambda_0\), \(\mathcal{S}_\omega[\Jimp]\)
includes dominant reactive \(1/r^{2}\), \(1/r^{3}\) terms on spheres with \(r\sim a\),
while \(\mathcal{R}_\omega[\Jimp]\) carries leading \(1/r\) transverse radiation with
\(k_0 r\) phase for \(r\gg a\) \cite{jackson1999,stratton1941}. Evaluating
Eq.~\eqref{eq:ch2.flux-on-radiative-branch} on \(S_R\) with \(R\sim a\) conflates storage
with export; only when \(k_0 R\gg 1\) and \(R\gg a\) does \(P_{\mathrm{rad}}(R)\) approach
the \(R\)-independent value predicted by far-zone reduction \cite{harrington2001}. Regime
gate Eq.~\eqref{eq:ch2.regime-operator-screen} must be checked before angular-transport
claims \cite{stratton1941}.

\paragraph{Worked illustration (near-field scan versus pattern).}
Near-field probes sample \(\mathcal{S}_\omega[\Jimp]\) on a surface enclosing \(\Omega_s\)
\cite{stratton1941}. Transforming those samples to a far-zone pattern implements
Eq.~\eqref{eq:ch2.radiation-integral-outgoing-upgrade} only when outgoing closure and
equivalent-source uniqueness on the exterior side are satisfied
(Theorem~\ref{thm:ch2.exterior-radiative-uniqueness}) \cite{harrington2001,jackson1999}.
Distinct impressed currents in \(\ker\mathcal{R}_\omega\) produce identical patterns while
differing in reactive near fields; near scans alone do not resolve that null space without
additional constraints (Subsection~\ref{sec:ch244}) \cite{stratton1941,harrington2001}.

The subsection sequence begins with integral representations that realize
Eq.~\eqref{eq:ch2.radiation-integral-outgoing-upgrade}, continues with near, intermediate,
and far asymptotics, develops flux and operator-norm meters on the radiative branch, and
closes with the radiative--non-radiative component split before spectral theory in
Section~\ref{sec:ch26} and reciprocity in Section~\ref{sec:ch28} compose additional
structure on the same integral operators \cite{stratton1941,harrington2001,jackson1999}.

\subsection{Source-to-field integral representations}
\label{sec:ch251}
% TOC tag: [HOME]

Section~\ref{sec:ch25} upgraded the radiation preview
Eq.~\eqref{eq:ch2.radiation-integral-outgoing-upgrade} on the outgoing dyad branch.
Subsection~\ref{sec:ch251} gives the Chapter~\ref{ch:02} HOME for the full family of
source-to-field integral representations that realize \(\mathcal{S}_\omega\) and
\(\mathcal{R}_\omega\) on admissible domains after Sections~\ref{sec:ch23}--\ref{sec:ch24}
\cite{stratton1941,harrington2001,jackson1999}. Scalar, vector, and dyadic kernels from
Subsection~\ref{sec:ch231}, regularization from Subsection~\ref{sec:ch232}, and outgoing
selection Eqs.~\eqref{eq:ch2.outgoing-green-branch}--\eqref{eq:ch2.outgoing-solution-operator}
are presupposed; free-space dyad tabulations and defining operators are not re-displayed
\cite{stratton1941}. Distance-regime asymptotics belong to Subsection~\ref{sec:ch252}; flux
meters to Subsection~\ref{sec:ch253}.

\paragraph{Assumptions.}
Fix \(\omega>0\), compact \(\operatorname{supp}\Jimp=\Omega_s\Subset\Omega\), admissible
data Eq.~\eqref{eq:ch1.admissible-data-class}, charge continuity
Eq.~\eqref{eq:ch1.charge-continuity-harmonic}, and outgoing closure
Eqs.~\eqref{eq:ch2.sommerfeld-outgoing-class}--\eqref{eq:ch2.outgoing-selector-def} on
exterior problems \cite{stratton1941}. Observation points \(\mathbf{r}\) lie outside
\(\Omega_s\) or on equivalent surfaces with trace matching declared below; coincidence
limits use Subsection~\ref{sec:ch232} \cite{harrington2001}. Kernels
\(\overline{\mathbf{G}}_{E,\omega}^{\mathrm{(out)}}\),
\(\overline{\mathbf{G}}_{H,\omega}^{\mathrm{(out)}}\), and scalar
\(G_\omega^{\mathrm{(out)}}\) denote outgoing branches fixed by
Eq.~\eqref{eq:ch2.outgoing-green-branch} \cite{stratton1941,jackson1999}.

\paragraph{Electric impressed-current integral.}
The primary volume representation is
\begin{equation}
\label{eq:ch2.outgoing-electric-volume-integral}
\E(\mathbf{r})
=
\jj\omega\mu_0
\int_{\Omega_s}
\overline{\mathbf{G}}_{E,\omega}^{\mathrm{(out)}}(\mathbf{r},\mathbf{r}')\cdot\Jimp(\mathbf{r}')\,d^{3}r'
\;+\;\E_{\mathrm{hom}}(\mathbf{r}),
\end{equation}
\begin{equation}
\label{eq:ch2.outgoing-magnetic-from-curl-integral}
\H(\mathbf{r})
=
\frac{1}{\jj\omega\mu_0}\,\curl\E(\mathbf{r}),
\end{equation}
for \(\mathbf{r}\notin\Omega_s\) on simply connected observation regions
\cite{stratton1941,harrington2001}. Equations
\eqref{eq:ch2.outgoing-electric-volume-integral}--\eqref{eq:ch2.outgoing-magnetic-from-curl-integral}
are the Subsection~\ref{sec:ch251} HOME specialization of
Eqs.~\eqref{eq:ch2.electric-dyadic-convolution}--\eqref{eq:ch2.magnetic-from-electric-dyad}
with outgoing dyads and exterior observation \cite{stratton1941}. They coincide with
Eq.~\eqref{eq:ch2.radiation-integral-outgoing-upgrade} when
\(\E_{\mathrm{hom}}=\mathbf{0}\) and \(\curl\) is taken in \(\mathbf{r}\)
\cite{harrington2001}. The prefactor \(\jj\omega\mu_0\) matches
Eq.~\eqref{eq:ch2.helmholtz-forcing-map} and the locked phasor axis
\cite{stratton1941}.

\paragraph{Magnetic impressed-current integral.}
When magnetic impressed sources
\(\mathbf{J}^{\mathrm{(mag)}}_{\mathrm{imp}}\) are primary, duality maps
\(\varepsilon_0\E\leftrightarrow\mu_0\H\) and
\(\Jimp\leftrightarrow\mathbf{J}^{\mathrm{(mag)}}_{\mathrm{imp}}\) yields
\begin{equation}
\label{eq:ch2.outgoing-magnetic-volume-integral}
\H(\mathbf{r})
=
\jj\omega\varepsilon_0
\int_{\Omega_s}
\overline{\mathbf{G}}_{H,\omega}^{\mathrm{(out)}}(\mathbf{r},\mathbf{r}')\cdot
\mathbf{J}^{\mathrm{(mag)}}_{\mathrm{imp}}(\mathbf{r}')\,d^{3}r'
\;+\;\H_{\mathrm{hom}}(\mathbf{r}),
\end{equation}
\begin{equation}
\label{eq:ch2.outgoing-electric-from-curl-magnetic}
\E(\mathbf{r})
=
-\frac{1}{\jj\omega\varepsilon_0}\,\curl\H(\mathbf{r}),
\end{equation}
with \(\overline{\mathbf{G}}_{H,\omega}^{\mathrm{(out)}}\) the magnetic dyad defined
duality-wise from Subsection~\ref{sec:ch231} without re-tabulating its closed form
\cite{stratton1941,harrington2001}. Equations
\eqref{eq:ch2.outgoing-magnetic-volume-integral}--\eqref{eq:ch2.outgoing-electric-from-curl-magnetic}
complete the dual route to Eqs.~\eqref{eq:ch2.outgoing-electric-volume-integral}--\eqref{eq:ch2.outgoing-magnetic-from-curl-integral}
for aperture magnetic currents and dual-equivalent models \cite{stratton1941}.

\paragraph{Combined impressed sources.}
Superposition gives the coupled volume representation
\begin{equation}
\label{eq:ch2.combined-impressed-source-integral}
\begin{aligned}
\E(\mathbf{r})
&=
\jj\omega\mu_0\!\int_{\Omega_s}\!
\overline{\mathbf{G}}_{E,\omega}^{\mathrm{(out)}}(\mathbf{r},\mathbf{r}')\cdot\Jimp(\mathbf{r}')\,d^{3}r'
-\frac{1}{\jj\omega\varepsilon_0}\,\curl\H^{(\mathrm{mag})}(\mathbf{r})
+\E_{\mathrm{hom}}(\mathbf{r}),\\[4pt]
\H(\mathbf{r})
&=
\jj\omega\varepsilon_0\!\int_{\Omega_s}\!
\overline{\mathbf{G}}_{H,\omega}^{\mathrm{(out)}}(\mathbf{r},\mathbf{r}')\cdot
\mathbf{J}^{\mathrm{(mag)}}_{\mathrm{imp}}(\mathbf{r}')\,d^{3}r'
+\frac{1}{\jj\omega\mu_0}\,\curl\E^{(\mathrm{elec})}(\mathbf{r})
+\H_{\mathrm{hom}}(\mathbf{r}),
\end{aligned}
\end{equation}
where \(\E^{(\mathrm{elec})},\H^{(\mathrm{mag})}\) denote the single-source contributions
from Eqs.~\eqref{eq:ch2.outgoing-electric-volume-integral} and
\eqref{eq:ch2.outgoing-magnetic-volume-integral} \cite{stratton1941,harrington2001}.
Equation~\eqref{eq:ch2.combined-impressed-source-integral} is the first-order Maxwell
integral system on outgoing kernels; it reduces to either branch when one impressed
source vanishes \cite{stratton1941}. Linearity
Eq.~\eqref{eq:ch2.radiation-linearity} follows from Eq.~\eqref{eq:ch2.combined-impressed-source-integral}
on fixed \(\mathcal{B}_\Gamma\) \cite{harrington2001}.

\begin{figure}[t]
\centering
\ChOneFigBlock{%
\begin{tikzpicture}[ch1fig, node distance=7mm and 9mm]
  \node[ch1box] (jev) {$\Jimp$ volume};
  \node[ch1box, right=12mm of jev] (jmv) {$\mathbf{J}^{\mathrm{(mag)}}_{\mathrm{imp}}$};
  \node[ch1box, right=12mm of jmv] (surf) {Huygens\\$\Gamma_{\mathrm{eq}}$};
  \node[ch1box, below=12mm of jmv] (ge) {$\overline{\mathbf{G}}_{E}^{\mathrm{(out)}}$};
  \node[ch1box, below=12mm of surf] (gh) {$\overline{\mathbf{G}}_{H}^{\mathrm{(out)}}$};
  \node[ch1box, right=14mm of jmv] (ef) {$(\E,\H)$};
  \draw[ch1flow] (jev) -- (ge) -- (ef);
  \draw[ch1flow] (jmv) -- (gh) -- (ef);
  \draw[ch1flow] (surf) -- (ef);
\end{tikzpicture}%
}
\caption{Subsection~\ref{sec:ch251} integral routes: electric volume, magnetic volume,
and Huygens surface representations converge to the same \(\mathcal{S}_\omega\) image on
the outgoing branch when trace matching holds.}
\label{fig:ch2.source-to-field-integral-routes}
\end{figure}

Figure~\ref{fig:ch2.source-to-field-integral-routes} summarizes the representation
family; each route must use outgoing kernels and admissible sources
\cite{stratton1941,harrington2001}.

\paragraph{Charge and polarization augmentation.}
When \(\divg\Jimp=-\jj\omega\rho\) with nonzero \(\rho\), the complete electric field
includes a scalar-potential gradient driven by \(\rho\):
\begin{equation}
\label{eq:ch2.charge-augmented-electric-integral}
\E(\mathbf{r})
=
\jj\omega\mu_0\!\int_{\Omega_s}\!
\overline{\mathbf{G}}_{E,\omega}^{\mathrm{(out)}}(\mathbf{r},\mathbf{r}')\cdot\Jimp(\mathbf{r}')\,d^{3}r'
-\grad\!\int_{\Omega_s} G_\omega^{\mathrm{(out)}}(\mathbf{r},\mathbf{r}')\,
\frac{\rho(\mathbf{r}')}{\varepsilon_0}\,d^{3}r'
\;+\;\E_{\mathrm{hom}}(\mathbf{r}),
\end{equation}
with \(G_\omega^{\mathrm{(out)}}\) the outgoing scalar kernel from
Eq.~\eqref{eq:ch2.scalar-green-solution-operator} \cite{stratton1941,jackson1999}. Equation
\eqref{eq:ch2.charge-augmented-electric-integral} enforces Gauss coupling without leaving
\(\mathcal{J}_\Omega\); when Eq.~\eqref{eq:ch1.charge-continuity-harmonic} holds and
\(\rho\) is determined by \(\Jimp\), the gradient term is fixed by the same impressed data
\cite{stratton1941}. Lorenz-gauge vector-potential reduction
Eq.~\eqref{eq:ch2.lorenz-field-from-potential} collapses to
Eq.~\eqref{eq:ch2.charge-augmented-electric-integral} on uniform exterior domains; gauge
freedom is not re-opened \cite{harrington2001}.

\paragraph{Huygens equivalent-surface representation.}
Let \(\Gamma_{\mathrm{eq}}\) be a closed surface enclosing \(\Omega_s\), with outward
normal \(\hat{\mathbf{n}}\) and equivalent surface densities
\(\mathbf{J}_{s,\mathrm{eq}}\), \(\mathbf{K}_{s,\mathrm{eq}}\) on
\(\Gamma_{\mathrm{eq}}\) in the trace class of Subsection~\ref{sec:ch211}. The exterior
Huygens representation is
\begin{equation}
\label{eq:ch2.huygens-equivalent-integral}
\E(\mathbf{r})
=
\jj\omega\mu_0\!\int_{\Gamma_{\mathrm{eq}}}\!
\overline{\mathbf{G}}_{E,\omega}^{\mathrm{(out)}}(\mathbf{r},\mathbf{r}')\cdot
\mathbf{J}_{s,\mathrm{eq}}(\mathbf{r}')\,dS'
+\int_{\Gamma_{\mathrm{eq}}}\!
\nabla'\!G_\omega^{\mathrm{(out)}}(\mathbf{r},\mathbf{r}')\times
\mathbf{K}_{s,\mathrm{eq}}(\mathbf{r}')\,dS',
\end{equation}
\begin{equation}
\label{eq:ch2.huygens-magnetic-integral}
\H(\mathbf{r})
=
\jj\omega\varepsilon_0\!\int_{\Gamma_{\mathrm{eq}}}\!
\overline{\mathbf{G}}_{H,\omega}^{\mathrm{(out)}}(\mathbf{r},\mathbf{r}')\cdot
\mathbf{K}_{s,\mathrm{eq}}(\mathbf{r}')\,dS'
-\int_{\Gamma_{\mathrm{eq}}}\!
\nabla'\!G_\omega^{\mathrm{(out)}}(\mathbf{r},\mathbf{r}')\times
\mathbf{J}_{s,\mathrm{eq}}(\mathbf{r}')\,dS',
\end{equation}
for \(\mathbf{r}\) outside the closed surface bounded by \(\Gamma_{\mathrm{eq}}\)
\cite{stratton1941,harrington2001}. Equations
\eqref{eq:ch2.huygens-equivalent-integral}--\eqref{eq:ch2.huygens-magnetic-integral} are
the Subsection~\ref{sec:ch251} HOME surface form of the same outgoing propagation operator;
\(\nabla'\) acts on the source coordinate \(\mathbf{r}'\) \cite{stratton1941}. Tangential
field matching
\begin{equation}
\label{eq:ch2.huygens-tangential-match}
\hat{\mathbf{n}}\times\E\big|_{\Gamma_{\mathrm{eq}}^{+}}
=
\hat{\mathbf{n}}\times\E\big|_{\Gamma_{\mathrm{eq}}^{-}},
\qquad
\hat{\mathbf{n}}\times\H\big|_{\Gamma_{\mathrm{eq}}^{+}}
=
\hat{\mathbf{n}}\times\H\big|_{\Gamma_{\mathrm{eq}}^{-}}
\end{equation}
fixes \(\mathbf{J}_{s,\mathrm{eq}}=\hat{\mathbf{n}}\times(\H^{-}-\H^{+})\),
\(\mathbf{K}_{s,\mathrm{eq}}=\hat{\mathbf{n}}\times(\E^{-}-\E^{+})\) in the stratified
sense of Eq.~\eqref{eq:ch1.interface-conditions} \cite{stratton1941,harrington2001}.
Equivalent-source non-uniqueness Eq.~\eqref{eq:ch2.equivalent-source-ambiguity} applies:
distinct \((\mathbf{J}_{s,\mathrm{eq}},\mathbf{K}_{s,\mathrm{eq}})\) pairs related by
\(\ker\mathcal{R}_\omega\) yield identical far-zone radiation \cite{harrington2001}.

\paragraph{Integral realization of \(\mathcal{S}_\omega\).}
\begin{theorem}[Outgoing source-to-field integral representation]
\label{thm:ch2.source-to-field-integral}
Let \(\mathfrak{D}=(\Jimp,\mathbf{J}^{\mathrm{(mag)}}_{\mathrm{imp}},\rho,\mathcal{B}_\Gamma)\)
be admissible with outgoing closure on exterior problems. Then the physical solution
\(\mathcal{S}_\omega[\mathfrak{D}]=(\E,\H)\) is given by
Eqs.~\eqref{eq:ch2.combined-impressed-source-integral}--\eqref{eq:ch2.charge-augmented-electric-integral}
with outgoing kernels, or by Eqs.~\eqref{eq:ch2.huygens-equivalent-integral}--\eqref{eq:ch2.huygens-tangential-match}
when equivalent surface data replace volume sources \cite{stratton1941,harrington2001}.
The representation agrees with Eq.~\eqref{eq:ch2.outgoing-solution-operator} and
Green--BVP equivalence Eq.~\eqref{eq:ch2.green-bvp-equivalence} on the outgoing branch
fixed by Theorem~\ref{thm:ch2.exterior-radiative-uniqueness}.
\end{theorem}
\begin{proof}
Subsection~\ref{sec:ch231} defined dyadic kernels satisfying
Eq.~\eqref{eq:ch2.dyadic-green-defining-operator}; Subsection~\ref{sec:ch232} justified
convolution on \(\mathcal{J}_\Omega\); Section~\ref{sec:ch24} selected the outgoing
branch. Substituting \(\overline{\mathbf{G}}_{E,\omega}^{\mathrm{(out)}}\) into
Eq.~\eqref{eq:ch2.electric-dyadic-convolution} yields
Eq.~\eqref{eq:ch2.outgoing-electric-volume-integral}. Magnetic and coupled forms follow
from duality and superposition \cite{stratton1941}. Huygens formulas follow by replacing
volume sources with surface distributions that reproduce the same tangential traces
\cite{harrington2001}. Uniqueness on \(\mathcal{F}_{\mathrm{SF}}\) from
Theorem~\ref{thm:ch2.exterior-radiative-uniqueness} identifies the integral representation
with \(\mathcal{S}_\omega\) \cite{stratton1941}.
\end{proof}

\paragraph{Radiation map \(\mathcal{R}_\omega\).}
Integral representations return full admissible fields in \(\mathcal{F}_\Omega\). The
radiation map applies after evaluation:
\begin{equation}
\label{eq:ch2.integral-radiation-map}
\mathcal{R}_\omega[\mathfrak{D}]
=
\Pi_{\mathrm{rad}}\,
\Big(\E[\mathfrak{D}],\,\H[\mathfrak{D}]\Big),
\end{equation}
with \((\E,\H)\) from Theorem~\ref{thm:ch2.source-to-field-integral} and
\(\Pi_{\mathrm{rad}}\) from Eq.~\eqref{eq:ch2.outgoing-radiation-refinement}
\cite{stratton1941,harrington2001}. Equation~\eqref{eq:ch2.integral-radiation-map} is
the integral restatement of Eq.~\eqref{eq:ch2.radiation-forward-map}; far-zone pattern
and power integrals in Subsections~\ref{sec:ch252}--\ref{sec:ch253} act on
Eq.~\eqref{eq:ch2.integral-radiation-map}, not on raw
Eq.~\eqref{eq:ch2.outgoing-electric-volume-integral} alone \cite{stratton1941}. Regime
qualification Eq.~\eqref{eq:ch2.regime-operator-screen} remains mandatory before transport
claims \cite{harrington2001}.

\paragraph{Equivalence with operator resolvent.}
On the outgoing branch,
\begin{equation}
\label{eq:ch2.integral-resolvent-equivalence}
\mathcal{S}_\omega[\mathfrak{D}]
=
\mathcal{O}_{\mathrm{out}}\,
\mathcal{L}_\omega^{-1}\!
\begin{bmatrix}\mathbf{s}_\omega(\mathfrak{D})\\ \mathbf{g}_\Gamma\end{bmatrix}
=
\mathcal{O}_{\mathrm{out}}\,
\Big(\mathcal{I}_{\mathrm{vol}}[\mathfrak{D}]\Big),
\end{equation}
with \(\mathcal{I}_{\mathrm{vol}}[\mathfrak{D}]\) the combined volume integrals
Eqs.~\eqref{eq:ch2.combined-impressed-source-integral}--\eqref{eq:ch2.charge-augmented-electric-integral},
connecting integral and operator routes through
Eqs.~\eqref{eq:ch2.resolvent-green-equivalence}--\eqref{eq:ch2.outgoing-resolvent-selection}
\cite{stratton1941,harrington2001}. Equation~\eqref{eq:ch2.integral-resolvent-equivalence}
is the constructive content of Eq.~\eqref{eq:ch2.green-bvp-equivalence} for evaluation
purposes: numerical codes implement the left side via linear algebra and the right side
via quadrature on outgoing kernels \cite{harrington2001}.

\paragraph{Multi-region and boundary data.}
When \(\Omega=\cup_j\Omega_j\) with interface traces from Subsection~\ref{sec:ch234},
Eqs.~\eqref{eq:ch2.combined-impressed-source-integral}--\eqref{eq:ch2.charge-augmented-electric-integral}
hold patchwise with local \(k_j\) and outgoing exterior closure on the infinite exterior
region \cite{stratton1941}. Impressed \(\Jimp\) may vanish in passive exterior regions;
equivalent currents on material boundaries enter through
Eqs.~\eqref{eq:ch2.huygens-equivalent-integral}--\eqref{eq:ch2.huygens-tangential-match}
with \(\Gamma_{\mathrm{eq}}\) on scatterer surfaces \cite{harrington2001}. Interior
cavity problems replace outgoing \(G\) with standing-mode Green functions and require
quotienting Eq.~\eqref{eq:ch2.coercivity-quotient}; those representations are not
outgoing integrals and are excluded from Theorem~\ref{thm:ch2.source-to-field-integral}
\cite{stratton1941}.

\paragraph{Boundary of validity.}
Theorem~\ref{thm:ch2.source-to-field-integral} requires linear media, fixed \(\omega\),
admissible \(\mathfrak{D}\), outgoing closure on exterior problems, and observation
points outside singular coincidence unless regularization from Subsection~\ref{sec:ch232}
is applied \cite{stratton1941,harrington2001}. Incoming kernels produce different
integrals violating Eq.~\eqref{eq:ch2.retarded-harmonic-kernel-theorem}
\cite{jackson1999}. Anisotropic, moving, or nonlinear media require modified kernels
beyond the locked patchwise-scalar class \cite{stratton1941}. Near-zone reactive dominance
does not invalidate the integral; it signals that \(\mathcal{R}_\omega[\mathfrak{D}]\)
may differ sharply from \(\mathcal{S}_\omega[\mathfrak{D}]\) until
Eq.~\eqref{eq:ch2.regime-operator-screen} is satisfied \cite{harrington2001}.

\paragraph{Worked illustration (electric Hertzian dipole).}
Replace \(\Jimp\) by a localized dipole current in vacuum. Evaluating
Eq.~\eqref{eq:ch2.outgoing-electric-volume-integral} with
Eq.~\eqref{eq:ch2.free-space-dyadic-green} at \(k_0 r\gg 1\) reproduces the transverse
far pattern cited in Subsection~\ref{sec:ch231}; Eq.~\eqref{eq:ch2.integral-radiation-map}
returns the same transverse content as \(\Pi_{\mathrm{rad}}\mathcal{S}_\omega[\Jimp]\)
\cite{stratton1941,jackson1999}. Near-zone reactive terms remain in
\(\mathcal{S}_\omega[\Jimp]\) but are excluded from radiative export when
Eq.~\eqref{eq:ch2.rad-null-on-evanescent} applies \cite{harrington2001}.

\paragraph{Worked illustration (aperture Huygens).}
On a plane aperture in an infinite ground screen, equivalent magnetic and electric surface
densities on the aperture produce the exterior field through
Eqs.~\eqref{eq:ch2.huygens-equivalent-integral}--\eqref{eq:ch2.huygens-magnetic-integral}
with \(\Gamma_{\mathrm{eq}}\) the aperture \cite{stratton1941,harrington2001}. Distinct
equivalent pairs related by Eq.~\eqref{eq:ch2.equivalent-source-ambiguity} yield identical
\(\mathcal{R}_\omega\) images; Theorem~\ref{thm:ch2.exterior-radiative-uniqueness} fixes
the exterior \((\E,\H)\) once outgoing closure is imposed \cite{stratton1941}.

\paragraph{Worked illustration (method-of-moments far field).}
A boundary-element code evaluates surface integrals equivalent to
Eq.~\eqref{eq:ch2.huygens-equivalent-integral} with outgoing \(G\) embedded in the
impedance matrix \cite{harrington2001}. Far-field post-processing applies
Eq.~\eqref{eq:ch2.integral-radiation-map} through asymptotic reduction deferred to
Subsection~\ref{sec:ch252}; the continuous target is
Theorem~\ref{thm:ch2.source-to-field-integral}, not a mesh-specific heuristic
\cite{stratton1941}.

\paragraph{Closing summary.}
Subsection~\ref{sec:ch251} establishes the Chapter~\ref{ch:02} HOME for outgoing electric
volume integral Eqs.~\eqref{eq:ch2.outgoing-electric-volume-integral}--\eqref{eq:ch2.outgoing-magnetic-from-curl-integral},
magnetic dual Eqs.~\eqref{eq:ch2.outgoing-magnetic-volume-integral}--\eqref{eq:ch2.outgoing-electric-from-curl-magnetic},
combined system Eq.~\eqref{eq:ch2.combined-impressed-source-integral}, charge augmentation
Eq.~\eqref{eq:ch2.charge-augmented-electric-integral}, Huygens surface integrals
Eqs.~\eqref{eq:ch2.huygens-equivalent-integral}--\eqref{eq:ch2.huygens-tangential-match},
source-to-field theorem Theorem~\ref{thm:ch2.source-to-field-integral}, radiation map
Eq.~\eqref{eq:ch2.integral-radiation-map}, and resolvent equivalence
Eq.~\eqref{eq:ch2.integral-resolvent-equivalence}, realizing \(\mathcal{S}_\omega\) and
\(\mathcal{R}_\omega\) on the outgoing branch without re-tabulating Subsection~\ref{sec:ch231}
kernels \cite{stratton1941,harrington2001,jackson1999}. Subsection~\ref{sec:ch252} supplies
near, intermediate, and far asymptotics of these integrals.

\subsection{Near, intermediate, and far regimes}
\label{sec:ch252}
% TOC tag: [HOME]

Subsection~\ref{sec:ch251} fixed outgoing source-to-field integrals on the physical branch.
Subsection~\ref{sec:ch252} gives the Chapter~\ref{ch:02} HOME for near, intermediate, and
far distance regimes: field-magnitude scalings, far-zone reduction of
Eqs.~\eqref{eq:ch2.outgoing-electric-volume-integral}--\eqref{eq:ch2.combined-impressed-source-integral},
Fresnel--Fraunhofer transition scales, and power-flux convergence on large spheres
\cite{stratton1941,harrington2001,jackson1999}. The preview labels
Eq.~\eqref{eq:ch2.distance-regime-preview} from Section~\ref{sec:ch25} are upgraded here;
Sommerfeld asymptotics Eqs.~\eqref{eq:ch2.sommerfeld-dyadic-asymptotic} and
Eq.~\eqref{eq:ch2.far-zone-transverse-radiation} from Subsection~\ref{sec:ch241} are
presupposed \cite{stratton1941}. Regime indices
Eqs.~\eqref{eq:ch1.regime-index}--\eqref{eq:ch1.radiative-regime-condition} from
Section~\ref{sec:ch01.2.4} are invoked by reference; reactive-radiative proofs are not
re-opened \cite{stratton1941}. Flux meters belong to Subsection~\ref{sec:ch253}.

\paragraph{Assumptions and scales.}
Let \(\Omega_s=\operatorname{supp}\Jimp\) be compact with characteristic radius
\(a=\max_{\mathbf{r},\mathbf{r}'\in\Omega_s}|\mathbf{r}-\mathbf{r}'|/2\), vacuum
wavenumber \(k_0=\omega\sqrt{\varepsilon_0\mu_0}\), wavelength \(\lambda_0=2\pi/k_0\),
and electrical size \(k_0 a\) \cite{stratton1941,jackson1999}. Observation points
\(\mathbf{r}\) satisfy \(|\mathbf{r}|=r\notin\Omega_s\); reference origin is chosen in
\(\Omega_s\). Outgoing kernels from Subsection~\ref{sec:ch251} and dyadic structure from
Eq.~\eqref{eq:ch2.free-space-dyadic-green} are presupposed without re-tabulation
\cite{stratton1941}. Constants \(C_1,C_2,\ldots\) below depend on \(k_0,a\) and source
geometry but not on \(r\) in the asymptotic limits stated \cite{harrington2001}.

\paragraph{Regime definitions.}
The Chapter~\ref{ch:02} HOME regime partition is
\begin{equation}
\label{eq:ch2.distance-regime-definitions}
\begin{aligned}
\text{\emph{near (reactive) zone:}}\ & k_0 r\le C_{\mathrm{n}},\\[2pt]
\text{\emph{intermediate (transition) zone:}}\ &
C_{\mathrm{n}}<k_0 r\le C_{\mathrm{i}}(k_0 a),\\[2pt]
\text{\emph{far (radiative) zone:}}\ &
k_0 r> C_{\mathrm{i}}(k_0 a),\quad r> C_{\mathrm{f}}a^{2}/\lambda_0,
\end{aligned}
\end{equation}
with declared positive constants \(C_{\mathrm{n}}=\mathcal{O}(1)\),
\(C_{\mathrm{i}}=\mathcal{O}(1)\) when \(k_0 a=\mathcal{O}(1)\), and Fraunhofer scale
\(r_{\mathrm{F}}=2a^{2}/\lambda_0\) entering \(C_{\mathrm{f}}=\mathcal{O}(1)\)
\cite{stratton1941,jackson1999,harrington2001}. Equations
\eqref{eq:ch2.distance-regime-definitions} upgrade Eq.~\eqref{eq:ch2.distance-regime-preview}:
they are diagnostic bands, not disjoint sharp cutoffs \cite{stratton1941}. For
\(k_0 a\ll 1\) (electrically small sources), \(C_{\mathrm{i}}\) may be taken near unity
and the intermediate zone is narrow; for extended apertures with \(k_0 a\gg 1\), the
Fresnel--Fraunhofer condition \(r\gtrsim r_{\mathrm{F}}\) controls pattern formation
\cite{harrington2001,jackson1999}.

\paragraph{Near-zone field scalings.}
Evaluating Eq.~\eqref{eq:ch2.outgoing-electric-volume-integral} with
Eq.~\eqref{eq:ch2.free-space-dyadic-green} for \(k_0 r\le C_{\mathrm{n}}\) yields
\begin{equation}
\label{eq:ch2.near-zone-field-scaling}
|\E(\mathbf{r})|
=
\mathcal{O}\!\big(\mu_0\omega\|\Jimp\|_{\mathcal{J}}\,k_0^{-2} r^{-3}\big)
\;+\;
\mathcal{O}\!\big(\mu_0\omega\|\Jimp\|_{\mathcal{J}}\,k_0^{-1} r^{-2}\big),
\end{equation}
with \(\|\Jimp\|_{\mathcal{J}}\) a source amplitude scale on \(\Omega_s\)
\cite{stratton1941,jackson1999}. Equation~\eqref{eq:ch2.near-zone-field-scaling} reflects
the \(R^{-3}\) and \(R^{-2}\) terms in Eq.~\eqref{eq:ch2.free-space-dyadic-green}; the
leading \(r^{-3}\) contribution is reactive (stored) in the sense of
Section~\ref{sec:ch01.2.4} \cite{stratton1941}. Magnetic field follows from
Eq.~\eqref{eq:ch2.outgoing-magnetic-from-curl-integral} with the same near-zone
polynomial-in-\(k_0 r\) hierarchy \cite{harrington2001}. On shells with
Eq.~\eqref{eq:ch1.regime-index} \(\chi(r)\gg 1\), Poynting flux through
Eq.~\eqref{eq:ch2.radiated-power-screen} is not a credible export meter
\cite{stratton1941}.

\paragraph{Intermediate-zone transition.}
When \(C_{\mathrm{n}}<k_0 r\le C_{\mathrm{i}}(k_0 a)\), reactive and radiative terms in
Eq.~\eqref{eq:ch2.free-space-dyadic-green} are comparable:
\begin{equation}
\label{eq:ch2.intermediate-zone-balance}
|\E_{\mathrm{reactive}}(\mathbf{r})|
=
\mathcal{O}\!\big(|\E_{\mathrm{rad}}(\mathbf{r})|\,k_0 r\big),
\end{equation}
with \(\E_{\mathrm{rad}}\) the leading transverse \(e^{-\jj k_0 r}/r\) component and
\(\E_{\mathrm{reactive}}\) the remainder \cite{stratton1941,harrington2001}. Equation
\eqref{eq:ch2.intermediate-zone-balance} is the operator reading of regime mixing flagged
in Section~\ref{sec:ch01.2.4}: \(\mathcal{S}_\omega[\Jimp]\) contains both storage and
export, while \(\mathcal{R}_\omega[\Jimp]\) isolates export only after
Eq.~\eqref{eq:ch2.integral-radiation-map} \cite{stratton1941}. For extended sources, the
Fresnel parameter
\begin{equation}
\label{eq:ch2.fresnel-parameter}
F
=
\frac{\pi a^{2}}{\lambda_0 r}
=
\frac{k_0 a^{2}}{2r}
\end{equation}
separates Fresnel (\(F\gtrsim 1\)) from Fraunhofer (\(F\ll 1\)) pattern regions on
observation shells \cite{jackson1999,harrington2001}. Subsection~\ref{sec:ch252} records
\(F\) as a transition diagnostic; complete aperture diffraction theory is not re-developed
\cite{stratton1941}.

\begin{figure}[t]
\centering
\ChOneFigBlock{%
\begin{tikzpicture}[ch1fig, node distance=7mm and 10mm]
  \node[ch1box] (src) {Source\\$k_0 a$};
  \node[ch1box, right=14mm of src] (near) {Near\\$k_0 r\lesssim 1$};
  \node[ch1box, right=14mm of near] (mid) {Intermediate\\$F\sim 1$};
  \node[ch1box, right=14mm of mid] (far) {$1/r$\\radiative};
  \node[ch1box, below=11mm of near] (chi) {$\chi(r)\gg 1$};
  \node[ch1box, below=11mm of far] (chi2) {$\chi(r)\ll 1$};
  \draw[ch1flow] (src) -- (near) -- (mid) -- (far);
  \draw[ch1flow] (near) -- (chi);
  \draw[ch1flow] (far) -- (chi2);
\end{tikzpicture}%
}
\caption{Distance regimes Eq.~\eqref{eq:ch2.distance-regime-definitions}: reactive
dominance near the source, transition through Fresnel parameter
Eq.~\eqref{eq:ch2.fresnel-parameter}, radiative export in the far zone with
Eq.~\eqref{eq:ch1.regime-index} \(\chi(r)\ll 1\).}
\label{fig:ch2.distance-regime-schematic}
\end{figure}

Figure~\ref{fig:ch2.distance-regime-schematic} pairs spatial bands with regime index
\(\chi(r)\) from Section~\ref{sec:ch01.2.4} \cite{stratton1941}.

\paragraph{Far-zone reduction of integrals.}
When \(k_0 r\gg 1\) and \(r\gg r_{\mathrm{F}}\) for extended sources, outgoing dyadic
asymptotics Eq.~\eqref{eq:ch2.sommerfeld-dyadic-asymptotic} reduce
Eq.~\eqref{eq:ch2.outgoing-electric-volume-integral} to
\begin{equation}
\label{eq:ch2.far-field-integral-reduction}
\E(\mathbf{r})
=
\E_0(\theta,\phi)\,\frac{e^{-\jj k_0 r}}{r}
\;+\;\mathcal{O}(r^{-2}),
\qquad
\E_0\cdot\hat{\mathbf{r}}=0,
\end{equation}
\begin{equation}
\label{eq:ch2.far-field-amplitude-integral}
\E_0(\theta,\phi)
=
-\jj\omega\mu_0\,\frac{1}{4\pi}
\int_{\Omega_s}
\Jimp(\mathbf{r}')\,
e^{-\jj k_0\hat{\mathbf{r}}\cdot\mathbf{r}'}\,d^{3}r',
\end{equation}
with \(\hat{\mathbf{r}}=\mathbf{r}/r\) and the transversality constraint enforced by
\(\overline{\mathbf{G}}_{\infty}\) in Eq.~\eqref{eq:ch2.sommerfeld-dyadic-asymptotic}
\cite{stratton1941,jackson1999,harrington2001}. Equation
\eqref{eq:ch2.far-field-integral-reduction} is the Subsection~\ref{sec:ch252} HOME
specialization of Eq.~\eqref{eq:ch2.far-zone-transverse-radiation}; the amplitude integral
Eq.~\eqref{eq:ch2.far-field-amplitude-integral} is the vector radiation integral used in
antenna far-field post-processing \cite{harrington2001}. Magnetic field in the far zone
satisfies Eq.~\eqref{eq:ch2.sommerfeld-impedance-ratio} with the same \(\E_0\)
\cite{stratton1941}.

\begin{theorem}[Far-field approximation of source integrals]
\label{thm:ch2.far-field-approximation}
Let \(\Jimp\in\mathcal{J}_\Omega\) be supported in \(B_a(0)\) with outgoing closure.
For \(r> r_{\min}(a,\lambda_0)\) as in Eq.~\eqref{eq:ch2.distance-regime-definitions},
the field from Eq.~\eqref{eq:ch2.outgoing-electric-volume-integral} satisfies
Eqs.~\eqref{eq:ch2.far-field-integral-reduction}--\eqref{eq:ch2.far-field-amplitude-integral}
with error \(\mathcal{O}\big((k_0 a)^{2}/(k_0 r)\,|\E_0|\,r^{-1}\big)\) for extended
sources and \(\mathcal{O}\big((k_0 r)^{-1}\,|\E_0|\,r^{-1}\big)\) for \(k_0 a\ll 1\)
\cite{stratton1941,jackson1999}.
\end{theorem}
\begin{proof}
Substitute Eq.~\eqref{eq:ch2.sommerfeld-dyadic-asymptotic} into
Eq.~\eqref{eq:ch2.outgoing-electric-volume-integral} and expand
\(|\mathbf{r}-\mathbf{r}'|^{-1}\) for \(|\mathbf{r}'|\le a\ll r\). The leading
\(e^{-\jj k_0 r}/r\) term factorizes into \(\E_0(\theta,\phi)\) times phase
\(e^{-\jj k_0\hat{\mathbf{r}}\cdot\mathbf{r}'}\) over the source; transversality follows
from the dyadic structure Eq.~\eqref{eq:ch2.free-space-dyadic-green} at leading order
\cite{stratton1941,harrington2001}. Higher orders in \(a/r\) and \(1/(k_0 r)\) give the
stated error bounds \cite{jackson1999}.
\end{proof}

\paragraph{Power-flux convergence.}
On spheres \(S_r=\{|\mathbf{r}|=r\}\), define shell power
\(P(R)=\Re\oint_{S_R}\mathbf{S}\cdot\hat{\mathbf{r}}\,dS\) as in
Eq.~\eqref{eq:ch2.radiated-power-screen} \cite{stratton1941}. In the far zone,
\begin{equation}
\label{eq:ch2.power-flux-regime-convergence}
|P(r)-P_\infty|
=
\mathcal{O}(r^{-1}),
\qquad
P_\infty
=
\frac{1}{2\eta_0}\oint_{S_r}
\big|\E_0(\theta,\phi)\big|^{2}\,d\Omega,
\end{equation}
with \(\E_0\) from Eq.~\eqref{eq:ch2.far-field-amplitude-integral} and \(\eta_0\) from
Eq.~\eqref{eq:ch2.exterior-flux-asymptotic} \cite{stratton1941,harrington2001}. Equation
\eqref{eq:ch2.power-flux-regime-convergence} is the operator form of
Eq.~\eqref{eq:ch1.radiative-regime-condition}: \(P(r)\) stabilizes once reactive
\(r^{-2}\), \(r^{-3}\) contributions no longer dominate the shell integral
\cite{stratton1941}. Evaluating Eq.~\eqref{eq:ch2.radiated-power-screen} on
\(\mathcal{S}_\omega[\Jimp]\) before \(r\gtrsim r_c\) from
Eq.~\eqref{eq:ch1.crossover-radius} confuses storage with export
\cite{harrington2001}.

\paragraph{Operator regime map.}
Distance regimes induce an observation-shell map on operator outputs:
\begin{equation}
\label{eq:ch2.observation-shell-regime-map}
\begin{aligned}
r\in\text{near}
&\ \Longrightarrow\
\text{meter}\ \mathcal{S}_\omega[\Jimp],\\
r\in\text{far}
&\ \Longrightarrow\
\text{meter}\ \mathcal{R}_\omega[\Jimp]\ \text{via Eqs.~\eqref{eq:ch2.far-field-integral-reduction}--\eqref{eq:ch2.far-field-amplitude-integral}},
\end{aligned}
\end{equation}
with intermediate shells requiring explicit \(\chi(r)\) from
Eq.~\eqref{eq:ch1.regime-index} before transport claims \cite{stratton1941}. Regime screen
Eq.~\eqref{eq:ch2.regime-operator-screen} restates the same rule on projected subspaces:
only \(\Pi_{\mathrm{prop}}\mathcal{R}_\omega[\Jimp]\) enters transport diagnostics when
\(\mathcal{D}_{J,\mathrm{reg}}=1\) \cite{stratton1941,harrington2001}. Subsection~\ref{sec:ch252}
does not replace Eq.~\eqref{eq:ch2.regime-operator-screen}; it supplies the spatial
\(k_0 r\) and \(F\) thresholds that make the screen numerically checkable on shells
\cite{stratton1941}.

\paragraph{Electrical-size limits.}
When \(k_0 a\ll 1\), Eq.~\eqref{eq:ch2.far-field-amplitude-integral} reduces to the
electric-dipole pattern proportional to
\(\widehat{\boldsymbol{\theta}}\times(\widehat{\boldsymbol{\theta}}\times\mathbf{p})\)
with \(\mathbf{p}=\int \Jimp(\mathbf{r}')\,d^{3}r'\) \cite{jackson1999,stratton1941}. The
near zone extends to \(r\sim\lambda_0\); far-zone reduction applies for
\(k_0 r\gg 1\) \cite{harrington2001}. When \(k_0 a\gg 1\), multiple lobes appear in
\(\E_0(\theta,\phi)\) and Fraunhofer distance \(r_{\mathrm{F}}=2a^{2}/\lambda_0\) dominates
the far-zone threshold \cite{stratton1941,jackson1999}. Angular-spectrum action in
Section~\ref{sec:ch29} refines \(\E_0\); Subsection~\ref{sec:ch252} stops at the
\(e^{-\jj k_0 r}/r\) transverse template \cite{harrington2001}.

\paragraph{Boundary of validity.}
Equations~\eqref{eq:ch2.distance-regime-definitions}--\eqref{eq:ch2.power-flux-regime-convergence}
require linear vacuum or piecewise-constant media, fixed \(\omega\), outgoing closure,
and observation outside \(\Omega_s\) \cite{stratton1941,harrington2001}. They fail in
cutoff waveguides where propagating modes do not exist (Eq.~\eqref{eq:ch2.cutoff-condition}),
in cavity standing-wave fields, and in lossless open resonances without outgoing selection
\cite{stratton1941}. Anisotropic beams and surface waves require modified asymptotics beyond
the locked free-space dyad template \cite{harrington2001}. Theorem~\ref{thm:ch2.far-field-approximation}
does not assert uniqueness of sources---only field asymptotics on the outgoing branch
\cite{stratton1941}.

\paragraph{Worked illustration (small dipole shells).}
For \(k_0 a=0.5\), Eq.~\eqref{eq:ch2.near-zone-field-scaling} dominates on \(k_0 r=0.5\)
while Eq.~\eqref{eq:ch2.far-field-integral-reduction} applies on \(k_0 r=10\)
\cite{stratton1941,jackson1999}. Table-level regime indices from
Section~\ref{sec:ch01.2.4} show \(\chi(r)\gg 1\) on the inner shell and \(\chi(r)\ll 1\)
on the outer shell \cite{stratton1941}. \(P(r)\) from
Eq.~\eqref{eq:ch2.power-flux-regime-convergence} is not \(R\)-independent on the inner shell
\cite{harrington2001}.

\paragraph{Worked illustration (aperture Fraunhofer distance).}
For a circular aperture of radius \(a\) with \(k_0 a=10\), Eq.~\eqref{eq:ch2.fresnel-parameter}
gives \(F=1\) at \(r=r_{\mathrm{F}}=2a^{2}/\lambda_0\). Pattern measurements at
\(r<r_{\mathrm{F}}\) lie in the intermediate zone; at \(r\gg r_{\mathrm{F}}\),
Eq.~\eqref{eq:ch2.far-field-amplitude-integral} defines the stationary-phase far pattern
\cite{jackson1999,harrington2001,stratton1941}.

\paragraph{Worked illustration (NF/FF transform target).}
Near-field scanning implements Eq.~\eqref{eq:ch2.outgoing-electric-volume-integral} on a
surface in the intermediate zone and extrapolates to
Eq.~\eqref{eq:ch2.far-field-amplitude-integral} \cite{harrington2001}. The continuous target
is Theorem~\ref{thm:ch2.far-field-approximation} on the outgoing branch; transform error
combines measurement noise with violation of Eq.~\eqref{eq:ch2.distance-regime-definitions}
\cite{stratton1941}.

\paragraph{Closing summary.}
Subsection~\ref{sec:ch252} establishes the Chapter~\ref{ch:02} HOME for regime definitions
Eq.~\eqref{eq:ch2.distance-regime-definitions}, near-zone scaling
Eq.~\eqref{eq:ch2.near-zone-field-scaling}, intermediate balance
Eq.~\eqref{eq:ch2.intermediate-zone-balance}, Fresnel parameter
Eq.~\eqref{eq:ch2.fresnel-parameter}, far-field reduction
Eqs.~\eqref{eq:ch2.far-field-integral-reduction}--\eqref{eq:ch2.far-field-amplitude-integral},
Theorem~\ref{thm:ch2.far-field-approximation}, power convergence
Eq.~\eqref{eq:ch2.power-flux-regime-convergence}, and observation-shell map
Eq.~\eqref{eq:ch2.observation-shell-regime-map}, upgrading
Eq.~\eqref{eq:ch2.distance-regime-preview} and linking spatial scales to regime screen
Eq.~\eqref{eq:ch2.regime-operator-screen} \cite{stratton1941,harrington2001,jackson1999}.
Subsection~\ref{sec:ch253} develops flux and operator-norm meters on far-zone outputs.

\subsection{Energy, momentum, angular-momentum flux, power, and operator norms}
\label{sec:ch253}
% TOC tag: [USE SS1.1.4][USE SS1.1.5][USE SS1.2.3]

Subsections~\ref{sec:ch251}--\ref{sec:ch252} fixed outgoing integrals and distance regimes.
Subsection~\ref{sec:ch253} records the Chapter~\ref{ch:02} operator evaluation of energy,
momentum, and angular-momentum flux meters, power budgets, and radiation norms on
\(\mathcal{R}_\omega[\Jimp]\), importing flux identities from
Subsections~\ref{sec:ch01.1.4} and~\ref{sec:ch01.1.5} and the transport triad from
Subsection~\ref{sec:ch01.2.3} by reference \cite{stratton1941,poynting1884,belinfante1940}.
Poynting balance Eq.~\eqref{eq:ch1.poynting-complex}, Maxwell stress
Eq.~\eqref{eq:ch1.maxwell-stress}, angular balance
Eqs.~\eqref{eq:ch1.angular-momentum-balance}--\eqref{eq:ch1.angular-balance-integral}, and
transport indices Eqs.~\eqref{eq:ch1.transport-admissibility-index}--\eqref{eq:ch1.transport-qualified-decision}
are presupposed; their HOME proofs remain in Chapter~\ref{ch:01} \cite{stratton1941}. The
present subsection packages those functionals as operators on
\(\mathcal{F}_{\mathrm{rad}}\) reached through
Eq.~\eqref{eq:ch2.integral-radiation-map} after regime qualification
Eq.~\eqref{eq:ch2.regime-operator-screen} \cite{harrington2001}. Radiative--non-radiative
splitting belongs to Subsection~\ref{sec:ch254}.

\paragraph{Assumptions.}
Let \((\E,\H)=\mathcal{S}_\omega[\mathfrak{D}]\) be the outgoing solution from
Theorem~\ref{thm:ch2.source-to-field-integral} and
\((\E_{\mathrm{rad}},\H_{\mathrm{rad}})=\mathcal{R}_\omega[\mathfrak{D}]\) its radiative
image Eq.~\eqref{eq:ch2.integral-radiation-map} \cite{stratton1941}. Flux meters on
spherical shells \(S_r=\{|\mathbf{r}|=r\}\) assume \(r\) lies in the far zone of
Eq.~\eqref{eq:ch2.distance-regime-definitions} unless noted otherwise, with
\(\mathcal{D}_{J,\mathrm{reg}}=1\) from Eq.~\eqref{eq:ch1.regime-qualified-decision}
\cite{stratton1941,harrington2001}. Complex Poynting vector
\(\mathbf{S}=\tfrac12\E\times\H^*\) and stress tensor \(\mathbf{T}\) follow
Eqs.~\eqref{eq:ch1.poynting-complex} and~\eqref{eq:ch1.maxwell-stress} \cite{stratton1941}.
Angular-flux tensor \(\mathbf{M}=\mathbf{r}\times\mathbf{T}\) uses origin \(\mathbf{r}_0\)
declared in Subsection~\ref{sec:ch01.1.5} \cite{belinfante1940,stratton1941}.

\paragraph{Flux functionals on the radiative branch.}
Define the radiative Poynting functional
\begin{equation}
\label{eq:ch2.poynting-functional-radiative}
\mathbf{S}_{\mathrm{rad}}[\Jimp]
=
\frac{1}{2}\,\E_{\mathrm{rad}}\times\H_{\mathrm{rad}}^{*},
\qquad
(\E_{\mathrm{rad}},\H_{\mathrm{rad}})=\mathcal{R}_\omega[\Jimp],
\end{equation}
with \(\E_{\mathrm{rad}},\H_{\mathrm{rad}}\) the field components selected by
Eq.~\eqref{eq:ch2.outgoing-radiation-refinement} \cite{stratton1941}. Equation
\eqref{eq:ch2.poynting-functional-radiative} specializes
Eq.~\eqref{eq:ch1.poynting-complex} to the operator output on which export claims are
allowed per Eq.~\eqref{eq:ch2.transport-qualified-radiation} \cite{harrington2001}. The
reactive counterpart \(\mathbf{S}_{\mathrm{re}}=\tfrac12(\E-\E_{\mathrm{rad}})\times(\H-\H_{\mathrm{rad}})^*\)
is not a transport meter unless separately qualified \cite{stratton1941}. Stress and
angular-flux functionals on the radiative branch are
\begin{equation}
\label{eq:ch2.stress-functional-radiative}
\mathbf{T}_{\mathrm{rad}}[\Jimp]
=
\mathbf{T}(\E_{\mathrm{rad}},\H_{\mathrm{rad}}),
\qquad
\mathbf{M}_{\mathrm{rad}}[\Jimp]
=
\mathbf{r}\times\mathbf{T}_{\mathrm{rad}}[\Jimp],
\end{equation}
with \(\mathbf{T}(\cdot,\cdot)\) the Maxwell stress map Eq.~\eqref{eq:ch1.maxwell-stress}
\cite{stratton1941,belinfante1940}. Equations
\eqref{eq:ch2.poynting-functional-radiative}--\eqref{eq:ch2.stress-functional-radiative}
are the Subsection~\ref{sec:ch253} operator restatement of Chapter~\ref{ch:01} flux objects;
they introduce no new conservation law \cite{stratton1941}.

\paragraph{Power flux meter.}
Shell power on the radiative branch is the bounded linear functional
\begin{equation}
\label{eq:ch2.power-flux-functional}
\mathcal{P}_{\mathrm{rad}}(r;\Jimp)
=
\Re\oint_{S_r}\hat{\mathbf{r}}\cdot\mathbf{S}_{\mathrm{rad}}[\Jimp]\,dS,
\end{equation}
which coincides with Eq.~\eqref{eq:ch2.flux-on-radiative-branch} and
Eq.~\eqref{eq:ch2.radiated-power-screen} when \((\E,\H)\) in those equations is replaced
by \(\mathcal{R}_\omega[\Jimp]\) \cite{stratton1941,poynting1884}. In the far zone,
Eqs.~\eqref{eq:ch2.far-field-integral-reduction}--\eqref{eq:ch2.far-field-amplitude-integral}
give
\begin{equation}
\label{eq:ch2.far-zone-power-from-pattern}
\mathcal{P}_{\mathrm{rad}}(r;\Jimp)
=
\frac{1}{2\eta_0}\oint_{S_r}
\big|\E_0(\theta,\phi)\big|^{2}\,d\Omega
\;+\;\mathcal{O}(r^{-1}),
\end{equation}
with \(\eta_0=\sqrt{\mu_0/\varepsilon_0}\) and \(\E_0\) from
Eq.~\eqref{eq:ch2.far-field-amplitude-integral} \cite{stratton1941,jackson1999}. Equation
\eqref{eq:ch2.far-zone-power-from-pattern} is the operator form of
Eq.~\eqref{eq:ch2.power-flux-regime-convergence} and
Eq.~\eqref{eq:ch1.radiative-regime-condition} \cite{stratton1941}. Evaluating
\(\mathcal{P}_{\mathrm{rad}}\) on \(\mathcal{S}_\omega[\Jimp]\) before
\(r\gtrsim r_c\) from Eq.~\eqref{eq:ch1.crossover-radius} imports reactive storage into
the power budget \cite{harrington2001}.

\paragraph{Linear momentum flux meter.}
Net momentum export through \(S_r\) is
\begin{equation}
\label{eq:ch2.momentum-flux-functional}
\mathbf{F}_{\mathrm{rad}}(r;\Jimp)
=
\Re\oint_{S_r}\mathbf{T}_{\mathrm{rad}}[\Jimp]\cdot\hat{\mathbf{r}}\,dS,
\end{equation}
the time-averaged force functional inherited from
Eq.~\eqref{eq:ch1.linear-momentum-balance} \cite{stratton1941}. In the far zone with
Eq.~\eqref{eq:ch2.sommerfeld-impedance-ratio},
\begin{equation}
\label{eq:ch2.far-zone-momentum-flux}
\mathbf{F}_{\mathrm{rad}}(r;\Jimp)
=
\frac{1}{c}\,\hat{\mathbf{r}}\,\mathcal{P}_{\mathrm{rad}}(r;\Jimp)
\;+\;\mathcal{O}(r^{-1}),
\end{equation}
linking momentum export to radiated power on the propagating branch
\cite{stratton1941,jackson1999}. Equation~\eqref{eq:ch2.far-zone-momentum-flux} restates
Eq.~\eqref{eq:ch1.energy-momentum-directional-ratio} in the far-zone limit for unit
\(\hat{\mathbf{q}}=\hat{\mathbf{r}}\) \cite{stratton1941}. Momentum meters attach to
\(\mathcal{R}_\omega[\Jimp]\), not to reactive near-zone circulation in
\(\mathcal{S}_\omega[\Jimp]\setminus\mathcal{R}_\omega[\Jimp]\)
\cite{harrington2001}.

\paragraph{Angular-momentum flux meter.}
Axis torque about unit direction \(\hat{\mathbf{u}}\) is the functional
\begin{equation}
\label{eq:ch2.angular-flux-functional}
\tau_{\hat{\mathbf{u}}}(r;\Jimp)
=
-\hat{\mathbf{u}}\cdot\Re\oint_{S_r}
\mathbf{M}_{\mathrm{rad}}[\Jimp]\cdot\hat{\mathbf{r}}\,dS,
\end{equation}
specializing Eq.~\eqref{eq:ch1.torque-observable} at harmonic steady state to the
radiative branch \cite{belinfante1940,stratton1941}. Net angular-momentum flux through
\(S_r\) is
\begin{equation}
\label{eq:ch2.angular-momentum-flux-vector}
\mathbf{J}_{\mathrm{rad}}(r;\Jimp)
=
\Re\oint_{S_r}\mathbf{M}_{\mathrm{rad}}[\Jimp]\cdot\hat{\mathbf{r}}\,dS,
\end{equation}
with \(\tau_{\hat{\mathbf{u}}}=-\hat{\mathbf{u}}\cdot\mathbf{J}_{\mathrm{rad}}\)
\cite{stratton1941}. Orbital--spin splits Eq.~\eqref{eq:ch1.angular-density-split} are
not re-derived; any angular report must close through
Eq.~\eqref{eq:ch1.angular-balance-integral} on the declared control volume
\cite{belinfante1940}. Directional concentration uses
Eq.~\eqref{eq:ch1.angular-flux-directional-index} evaluated on
\(\mathbf{M}_{\mathrm{rad}}[\Jimp]\) \cite{stratton1941}.

\begin{figure}[t]
\centering
\ChOneFigBlock{%
\begin{tikzpicture}[ch1fig, node distance=7mm and 10mm]
  \node[ch1box] (j) {$\Jimp$};
  \node[ch1box, right=13mm of j] (r) {$\mathcal{R}_\omega$};
  \node[ch1box, right=13mm of r] (f) {Flux\\functionals};
  \node[ch1box, right=13mm of f] (n) {Norms};
  \node[ch1box, below=11mm of f] (p) {$\mathcal{P}_{\mathrm{rad}}$};
  \node[ch1box, below=11mm of r] (reg) {$\mathcal{D}_{J,\mathrm{reg}}$};
  \draw[ch1flow] (j) -- (r) -- (f) -- (n);
  \draw[ch1flow] (f) -- (p);
  \draw[ch1flow] (reg) -- (r);
\end{tikzpicture}%
}
\caption{Subsection~\ref{sec:ch253} flux pipeline: regime qualification gates
\(\mathcal{R}_\omega\); power, momentum, and angular-momentum functionals evaluate on
the radiative branch only.}
\label{fig:ch2.flux-meter-pipeline}
\end{figure}

Figure~\ref{fig:ch2.flux-meter-pipeline} orders evaluation after
Eq.~\eqref{eq:ch2.regime-operator-screen} \cite{stratton1941,harrington2001}.

\paragraph{Transport triad on operator outputs.}
Subsection~\ref{sec:ch01.2.3} fixed transport admissibility through conditions
\(\mathcal{C}_1\)--\(\mathcal{C}_3\) on field states. On operator outputs,
\begin{equation}
\label{eq:ch2.transport-triad-on-R}
\begin{aligned}
\mathcal{C}_1[\Jimp]
&\iff
\text{Eq.~\eqref{eq:ch1.am-net-transport-condition} holds on}\ \mathbf{M}_{\mathrm{rad}}[\Jimp],\\
\mathcal{C}_2[\Jimp]
&\iff
\text{Eq.~\eqref{eq:ch1.angular-balance-integral} holds for}\ (\E_{\mathrm{rad}},\H_{\mathrm{rad}}),\\
\mathcal{C}_3[\Jimp]
&\iff
\text{Eq.~\eqref{eq:ch1.propagative-support} holds with}\ 
\Pi=\Pi_{\mathrm{rad}}\mathcal{R}_\omega[\Jimp],
\end{aligned}
\end{equation}
with indices \(\Lambda_J\), \(\mathcal{K}_J\), and \(\mathcal{R}_\tau\) from
Eqs.~\eqref{eq:ch1.transport-admissibility-index}--\eqref{eq:ch1.torque-consistency-residual}
evaluated on the same radiative fields \cite{stratton1941}. Operational acceptance
\(\mathcal{D}_J=1\) in Eq.~\eqref{eq:ch1.transport-qualified-decision} therefore means
transport interpretation is attached to \(\mathcal{R}_\omega[\Jimp]\) after
\(\mathcal{D}_{J,\mathrm{reg}}=1\), not to raw \(\mathcal{S}_\omega[\Jimp]\)
\cite{stratton1941,harrington2001}. Proposition~\ref{prop:ch1.transport-triad-sufficiency}
applies unchanged once \((\E_{\mathrm{rad}},\H_{\mathrm{rad}})\) replaces the full pair
on the measurement shell \cite{stratton1941}.

\paragraph{Radiation operator norms.}
Field energy on the radiative branch uses the inherited bilinear form
Eq.~\eqref{eq:ch2.maxwell-energy-form} restricted to \(\mathcal{F}_{\mathrm{rad}}\):
\begin{equation}
\label{eq:ch2.radiation-energy-norm}
\|\mathbf{u}_{\mathrm{rad}}\|_{\mathcal{E}}^{2}
=
\mathfrak{a}_\omega\big(\mathbf{u}_{\mathrm{rad}},\mathbf{u}_{\mathrm{rad}}\big),
\qquad
\mathbf{u}_{\mathrm{rad}}=(\E_{\mathrm{rad}},\H_{\mathrm{rad}}),
\end{equation}
with \(\mathfrak{a}_\omega\) from Subsection~\ref{sec:ch202} \cite{stratton1941}. Shell
\(L^{2}\) power density norms for pattern reporting are
\begin{equation}
\label{eq:ch2.radiation-pattern-norm}
\|\E_0\|_{L^{2}(S_r)}^{2}
=
\oint_{S_r}\big|\E_0(\theta,\phi)\big|^{2}\,d\Omega,
\qquad
\|\Jimp\|_{\mathcal{J}}=\Big(\int_{\Omega_s}|\Jimp|^{2}\,dV\Big)^{1/2},
\end{equation}
linking source amplitude to far pattern through
Eq.~\eqref{eq:ch2.far-field-amplitude-integral} \cite{harrington2001,stratton1941}. The
radiation operator norm on admissible sources is
\begin{equation}
\label{eq:ch2.radiation-operator-norm}
\|\mathcal{R}_\omega\|_{\mathrm{op}}
=
\sup_{\|\Jimp\|_{\mathcal{J}}=1}
\|\mathcal{R}_\omega[\Jimp]\|_{\mathcal{E}},
\end{equation}
bounded on fixed outgoing domains when resolvent bound
Eq.~\eqref{eq:ch2.resolvent-bounded} holds \cite{stratton1941,harrington2001}. Observation
maps compose on the right, yielding
\begin{equation}
\label{eq:ch2.flux-meter-composition}
\Lambda_{\mathrm{meas}}
=
\mathcal{O}_{\mathrm{meas}}\circ\mathcal{O}_{\mathrm{far}}\circ\mathcal{R}_\omega,
\qquad
\|\Lambda_{\mathrm{meas}}\|_{\mathrm{op}}
\le
\|\mathcal{O}_{\mathrm{meas}}\|_{\mathrm{op}}\,
\|\mathcal{R}_\omega\|_{\mathrm{op}},
\end{equation}
consistent with Eq.~\eqref{eq:ch2.reconstruction-observation-map} and rank cap
Eq.~\eqref{eq:ch2.reconstruction-rank-cap} \cite{stratton1941,harrington2001}.

\paragraph{Far-zone flux closure.}
\begin{proposition}[Far-zone flux stabilization on \(\mathcal{R}_\omega\)]
\label{prop:ch2.far-zone-flux-stabilization}
Let \(\Jimp\) be admissible with \(\mathcal{D}_{J,\mathrm{reg}}=1\) and let \(r\) lie
in the far zone of Eq.~\eqref{eq:ch2.distance-regime-definitions}. Then
\(\mathcal{P}_{\mathrm{rad}}(r;\Jimp)\), \(\mathbf{F}_{\mathrm{rad}}(r;\Jimp)\), and
\(\mathbf{J}_{\mathrm{rad}}(r;\Jimp)\) satisfy
Eqs.~\eqref{eq:ch2.far-zone-power-from-pattern}--\eqref{eq:ch2.far-zone-momentum-flux}
with errors \(\mathcal{O}(r^{-1})\), and \(\mathcal{P}_{\mathrm{rad}}(r;\Jimp)\to P_\infty>0\)
whenever Eq.~\eqref{eq:ch1.radiative-regime-condition} holds on the radiative branch
\cite{stratton1941,jackson1999}.
\end{proposition}
\begin{proof}
Substitute Eqs.~\eqref{eq:ch2.far-field-integral-reduction}--\eqref{eq:ch2.far-field-amplitude-integral}
into Eqs.~\eqref{eq:ch2.poynting-functional-radiative}--\eqref{eq:ch2.momentum-flux-functional}
and integrate on \(S_r\). Sommerfeld ratio Eq.~\eqref{eq:ch2.sommerfeld-impedance-ratio}
gives Eq.~\eqref{eq:ch2.far-zone-momentum-flux}; power integral reduces to
Eq.~\eqref{eq:ch2.far-zone-power-from-pattern} \cite{stratton1941,harrington2001}. Convergence
follows from Eq.~\eqref{eq:ch2.power-flux-regime-convergence} \cite{stratton1941}.
\end{proof}

\paragraph{Boundary of validity.}
Equations~\eqref{eq:ch2.poynting-functional-radiative}--\eqref{eq:ch2.radiation-operator-norm}
require harmonic steady state, linear passive media, outgoing closure, and regime-qualified
shells \cite{stratton1941,harrington2001}. They fail on reactive-dominated shells with
\(\chi(r)\gg 1\) from Eq.~\eqref{eq:ch1.regime-index}, in cutoff guides, and when
\(\mathcal{S}_\omega[\Jimp]\) is used without \(\mathcal{R}_\omega\) projection
\cite{stratton1941}. Material stress debates (Minkowski versus Abraham) are not re-opened;
vacuum-side shells use Eq.~\eqref{eq:ch1.maxwell-stress} \cite{stratton1941}. Gauge-dependent
near-zone angular densities are excluded from \(\mathbf{M}_{\mathrm{rad}}[\Jimp]\) unless
closure through Eq.~\eqref{eq:ch1.angular-balance-integral} is verified
\cite{belinfante1940}.

\paragraph{Worked illustration (dipole power and momentum).}
For an electrically small dipole with \(\E_0\) from
Eq.~\eqref{eq:ch2.far-field-amplitude-integral}, Eq.~\eqref{eq:ch2.far-zone-power-from-pattern}
returns \(P_\infty=\eta_0 k_0^{4}|\mathbf{p}|^{2}/(12\pi)\) in standard units
\cite{jackson1999,stratton1941}. Eq.~\eqref{eq:ch2.far-zone-momentum-flux} gives
\(|\mathbf{F}_{\mathrm{rad}}|=P_\infty/c\) directed along radiated power flow
\cite{stratton1941}. Evaluating the same integrals on \(\mathcal{S}_\omega[\Jimp]\) at
\(k_0 r=0.5\) violates Proposition~\ref{prop:ch2.far-zone-flux-stabilization}
\cite{harrington2001}.

\paragraph{Worked illustration (transport index on \(\mathcal{R}_\omega\)).}
Compute \(\Lambda_J(r)\) from Eq.~\eqref{eq:ch1.transport-admissibility-index} using
\(\mathbf{M}_{\mathrm{rad}}[\Jimp]\) on \(S_r\) with \(k_0 r=10\) and compare to
\(k_0 r=0.5\) \cite{stratton1941}. Only the outer shell with
\(\mathcal{D}_{J,\mathrm{reg}}=1\) and \(\mathcal{C}_1\)--\(\mathcal{C}_3\) satisfied
admits transport interpretation per Eq.~\eqref{eq:ch1.transport-qualified-decision}
\cite{stratton1941,harrington2001}.

\paragraph{Worked illustration (gain and pattern norm).}
Antenna gain uses Eq.~\eqref{eq:ch2.radiation-pattern-norm} normalized by
\(\mathcal{P}_{\mathrm{rad}}\) from Eq.~\eqref{eq:ch2.power-flux-functional} on the
radiative branch \cite{harrington2001,stratton1941}. Pattern reports from near-zone
\(\mathcal{S}_\omega[\Jimp]\) without Eq.~\eqref{eq:ch2.integral-radiation-map} conflate
reactive structure with export \cite{stratton1941}.

\paragraph{Closing summary.}
Subsection~\ref{sec:ch253} records the Chapter~\ref{ch:02} operator evaluation of flux
meters Eqs.~\eqref{eq:ch2.poynting-functional-radiative}--\eqref{eq:ch2.angular-momentum-flux-vector},
far-zone power and momentum Eqs.~\eqref{eq:ch2.far-zone-power-from-pattern}--\eqref{eq:ch2.far-zone-momentum-flux},
transport triad Eq.~\eqref{eq:ch2.transport-triad-on-R}, radiation norms
Eqs.~\eqref{eq:ch2.radiation-energy-norm}--\eqref{eq:ch2.radiation-operator-norm}, measurement
composition Eq.~\eqref{eq:ch2.flux-meter-composition}, and far-zone stabilization
Proposition~\ref{prop:ch2.far-zone-flux-stabilization}, importing
Subsections~\ref{sec:ch01.1.4},~\ref{sec:ch01.1.5}, and~\ref{sec:ch01.2.3} by reference
\cite{stratton1941,harrington2001,jackson1999}. Subsection~\ref{sec:ch254} separates
radiative and non-radiative components with projectors from Subsection~\ref{sec:ch224}.

\subsection{Radiative and non-radiative components}
\label{sec:ch254}
% TOC tag: [USE SS1.2.4][USE SS2.2.4]

Subsections~\ref{sec:ch251}--\ref{sec:ch253} fixed outgoing integrals, distance regimes, and
flux meters on \(\mathcal{R}_\omega[\Jimp]\). Subsection~\ref{sec:ch254} records the
Chapter~\ref{ch:02} operator split of \(\mathcal{S}_\omega[\Jimp]\) into radiative export
and non-radiative storage/reactive components, importing the reactive--radiative partition
of Section~\ref{sec:ch01.2.4} and spectral projectors
Eqs.~\eqref{eq:ch2.spectral-projectors}--\eqref{eq:ch2.rad-null-on-evanescent} from
Subsection~\ref{sec:ch224} by reference \cite{stratton1941,harrington2001}. Spectral
projector algebra, dispersion splits, and cutoff logic are not re-proved; the present
subsection applies those objects to integral outputs from Theorem~\ref{thm:ch2.source-to-field-integral}
and flux functionals from Subsection~\ref{sec:ch253} \cite{stratton1941}. Complete
non-radiating null-space classification belongs to Section~\ref{sec:ch27}.

\paragraph{Assumptions.}
Let \(\mathbf{u}_\omega=(\E,\H)=\mathcal{S}_\omega[\mathfrak{D}]\) be the outgoing
solution from Subsection~\ref{sec:ch251} and \(\mathbf{u}_{\mathrm{rad}}=\mathcal{R}_\omega[\mathfrak{D}]\)
its radiative image Eq.~\eqref{eq:ch2.integral-radiation-map} \cite{stratton1941}. Spectral
projectors \(\Pi_{\mathrm{prop}}\), \(\Pi_{\mathrm{ev}}\), \(\Pi_{\mathrm{non}}\) act on
\(\widehat{\mathcal{F}}_\Omega\) as in Subsection~\ref{sec:ch224}; outgoing selector
\(\mathcal{O}_{\mathrm{out}}\) and radiation map satisfy
Eqs.~\eqref{eq:ch2.outgoing-radiation-refinement}--\eqref{eq:ch2.rad-projection-composition}
\cite{stratton1941,harrington2001}. Regime qualification uses
\(\mathcal{D}_{J,\mathrm{reg}}\) from Eq.~\eqref{eq:ch1.regime-qualified-decision} and
screen Eq.~\eqref{eq:ch2.regime-operator-screen} \cite{stratton1941}.

\paragraph{Spectral split of the solution map.}
On admissible \(\mathfrak{D}\),
\begin{equation}
\label{eq:ch2.integral-spectral-decomposition}
\mathcal{S}_\omega[\mathfrak{D}]
=
\Pi_{\mathrm{prop}}\mathcal{S}_\omega[\mathfrak{D}]
+\Pi_{\mathrm{ev}}\mathcal{S}_\omega[\mathfrak{D}]
+\Pi_{\mathrm{non}}\mathcal{S}_\omega[\mathfrak{D}],
\end{equation}
restating Eq.~\eqref{eq:ch2.propagating-evanescent-decomposition} on the integral outputs
of Theorem~\ref{thm:ch2.source-to-field-integral} \cite{stratton1941,harrington2001}.
Define the \emph{non-propagating} component
\begin{equation}
\label{eq:ch2.nonpropagating-component-def}
\mathbf{u}_{\mathrm{np}}
=
\Pi_{\mathrm{ev}}\mathcal{S}_\omega[\mathfrak{D}]
+\Pi_{\mathrm{non}}\mathcal{S}_\omega[\mathfrak{D}],
\end{equation}
and the \emph{propagating} component
\(\mathbf{u}_{\mathrm{prop}}=\Pi_{\mathrm{prop}}\mathcal{S}_\omega[\mathfrak{D}]\)
\cite{stratton1941}. Only \(\mathbf{u}_{\mathrm{prop}}\) is eligible for outgoing radiation
selection; evanescent and non-propagating content satisfies
Eq.~\eqref{eq:ch2.rad-null-on-evanescent} \cite{harrington2001}. Equation
\eqref{eq:ch2.integral-spectral-decomposition} is the Subsection~\ref{sec:ch254} operator
form of Section~\ref{sec:ch01.2.4}: reactive dominance on near shells corresponds to large
\(\|\mathbf{u}_{\mathrm{np}}\|\) relative to \(\|\mathbf{u}_{\mathrm{prop}}\|\) in the sense
of Eq.~\eqref{eq:ch2.subspace-energy-inequality} \cite{stratton1941}.

\paragraph{Radiative versus non-radiative components.}
The radiative component is
\begin{equation}
\label{eq:ch2.radiative-component-def}
\mathbf{u}_{\mathrm{rad}}
=
\mathcal{R}_\omega[\mathfrak{D}]
=
\Pi_{\mathrm{rad}}\,\mathcal{S}_\omega[\mathfrak{D}]
=
\mathcal{O}_{\mathrm{out}}\circ\Pi_{\mathrm{prop}}\,\mathcal{S}_\omega[\mathfrak{D}],
\end{equation}
using Eqs.~\eqref{eq:ch2.outgoing-radiation-refinement} and~\eqref{eq:ch2.rad-projection-composition}
\cite{stratton1941,harrington2001}. The complementary \emph{non-radiative} component is
\begin{equation}
\label{eq:ch2.nonradiative-component-def}
\mathbf{u}_{\mathrm{nr}}
=
\mathcal{S}_\omega[\mathfrak{D}]-\mathbf{u}_{\mathrm{rad}}
=
\big(\mathrm{id}-\Pi_{\mathrm{rad}}\big)\mathcal{S}_\omega[\mathfrak{D}].
\end{equation}
Equations~\eqref{eq:ch2.radiative-component-def}--\eqref{eq:ch2.nonradiative-component-def}
partition every admissible field pair into export and non-export parts without altering
Maxwell dynamics: \(\mathbf{u}_\omega=\mathbf{u}_{\mathrm{rad}}+\mathbf{u}_{\mathrm{nr}}\)
\cite{stratton1941}. Non-radiative content includes evanescent, non-propagating, and
incoming-propagating fractions annihilated by \(\Pi_{\mathrm{rad}}\); it is not necessarily
zero when \(|\E|\) is large in the near zone \cite{harrington2001}. Sources in
\(\ker\mathcal{R}_\omega\) produce \(\mathbf{u}_{\mathrm{rad}}=\mathbf{0}\) while
\(\mathbf{u}_{\mathrm{nr}}\neq\mathbf{0}\) (Eq.~\eqref{eq:ch2.ker-R-omega-preview})
\cite{stratton1941}.

\begin{figure}[t]
\centering
\ChOneFigBlock{%
\begin{tikzpicture}[ch1fig, node distance=7mm and 9mm]
  \node[ch1box] (s) {$\mathcal{S}_\omega[\mathfrak{D}]$};
  \node[ch1box, right=12mm of s] (pp) {$\Pi_{\mathrm{prop}}$};
  \node[ch1box, right=12mm of pp] (pr) {$\Pi_{\mathrm{rad}}$};
  \node[ch1box, below=11mm of s] (np) {$\mathbf{u}_{\mathrm{nr}}$};
  \node[ch1box, below=11mm of pp] (ev) {$\Pi_{\mathrm{ev}}+\Pi_{\mathrm{non}}$};
  \draw[ch1flow] (s) -- (pp) -- (pr);
  \draw[ch1flow] (s) -- (np);
  \draw[ch1flow] (s) -- (ev);
\end{tikzpicture}%
}
\caption{Radiative/non-radiative split: \(\mathbf{u}_{\mathrm{rad}}\) is
\(\Pi_{\mathrm{rad}}\mathcal{S}_\omega\); \(\mathbf{u}_{\mathrm{nr}}\) collects all
non-export content including Eq.~\eqref{eq:ch2.nonpropagating-component-def}.}
\label{fig:ch2.rad-nr-component-split}
\end{figure}

Figure~\ref{fig:ch2.rad-nr-component-split} refines
Figure~\ref{fig:ch2.spectral-radiation-gate} from Subsection~\ref{sec:ch224} on integral
outputs \cite{stratton1941,harrington2001}.

\paragraph{Energy and power budgets.}
Using energy norm Eq.~\eqref{eq:ch2.field-energy-norm}, define
\begin{equation}
\label{eq:ch2.radiative-energy-functional}
\mathcal{E}_{\mathrm{rad}}[\mathfrak{D}]
=
\|\mathbf{u}_{\mathrm{rad}}\|_{\mathcal{E}}^{2},
\qquad
\mathcal{E}_{\mathrm{nr}}[\mathfrak{D}]
=
\|\mathbf{u}_{\mathrm{nr}}\|_{\mathcal{E}}^{2},
\end{equation}
with \(\|\cdot\|_{\mathcal{E}}\) from Eq.~\eqref{eq:ch2.radiation-energy-norm} on each
component \cite{stratton1941}. Shell power splits analogously:
\begin{equation}
\label{eq:ch2.component-power-budget}
\mathcal{P}_{\mathrm{tot}}(r)
=
\mathcal{P}_{\mathrm{rad}}(r;\Jimp)
+\mathcal{P}_{\mathrm{nr}}(r;\Jimp),
\end{equation}
where \(\mathcal{P}_{\mathrm{rad}}\) is Eq.~\eqref{eq:ch2.power-flux-functional} and
\begin{equation}
\label{eq:ch2.nonradiative-power-functional}
\mathcal{P}_{\mathrm{nr}}(r;\Jimp)
=
\Re\oint_{S_r}\hat{\mathbf{r}}\cdot\mathbf{S}_{\mathrm{nr}}[\Jimp]\,dS,
\qquad
\mathbf{S}_{\mathrm{nr}}
=
\tfrac12\,\E_{\mathrm{nr}}\times\H_{\mathrm{nr}}^{*},
\end{equation}
with \((\E_{\mathrm{nr}},\H_{\mathrm{nr}})=\mathbf{u}_{\mathrm{nr}}\)
\cite{stratton1941,poynting1884}. Equation~\eqref{eq:ch2.component-power-budget} is not
a new conservation law; it decomposes Eq.~\eqref{eq:ch2.radiated-power-screen} when both
components are evaluated on the same shell \cite{stratton1941}. In the far zone,
\(\mathcal{P}_{\mathrm{nr}}(r)\to 0\) and \(\mathcal{P}_{\mathrm{rad}}(r)\to P_\infty\) under
Eq.~\eqref{eq:ch1.radiative-regime-condition} \cite{stratton1941,harrington2001}. Near-zone
shells with \(\chi(r)\gg 1\) from Eq.~\eqref{eq:ch1.regime-index} typically have
\(|\mathcal{P}_{\mathrm{nr}}|\gg|\mathcal{P}_{\mathrm{rad}}|\) \cite{stratton1941}.

\paragraph{Regime screen on components.}
Regime qualification from Section~\ref{sec:ch01.2.4} becomes
\begin{equation}
\label{eq:ch2.regime-component-gate}
\mathcal{D}_{J,\mathrm{reg}}=1
\ \Longrightarrow\
\text{transport uses}\ \mathbf{u}_{\mathrm{rad}},\
\mathcal{P}_{\mathrm{rad}},\
\mathbf{M}_{\mathrm{rad}}[\Jimp],
\end{equation}
and forbids interpreting \(\mathbf{u}_{\mathrm{nr}}\), \(\mathcal{P}_{\mathrm{nr}}\), or
\(\Pi_{\mathrm{ev}}\mathcal{S}_\omega[\mathfrak{D}]\) as export channels
\cite{stratton1941,harrington2001}. Equation~\eqref{eq:ch2.regime-component-gate} is the
component form of Eq.~\eqref{eq:ch2.regime-operator-screen} after
Eqs.~\eqref{eq:ch2.radiative-component-def}--\eqref{eq:ch2.nonradiative-component-def}.
Failure modes from Section~\ref{sec:ch01.2.4}---near-zone reactive scans, cutoff-guide
modes, evanescent dominance---appear as large \(\mathcal{E}_{\mathrm{nr}}\) or
\(\|\Pi_{\mathrm{ev}}\mathcal{S}_\omega[\mathfrak{D}]\|\) with small
\(\mathcal{P}_{\mathrm{rad}}\) \cite{stratton1941,harrington2001}.

\paragraph{Link to radiation null space.}
Non-radiative components encode source-side invisible content:
\begin{equation}
\label{eq:ch2.nr-image-of-null-space}
\Jimp\in\ker\mathcal{R}_\omega
\ \Longrightarrow\
\mathbf{u}_{\mathrm{rad}}=\mathbf{0},
\quad
\mathbf{u}_{\mathrm{nr}}=\mathcal{S}_\omega[\Jimp]\neq\mathbf{0}\ \text{in general},
\end{equation}
as previewed in Eq.~\eqref{eq:ch2.null-radiation-preview} \cite{harrington2001,stratton1941}.
Conversely, \(\mathbf{u}_{\mathrm{nr}}=\mathbf{0}\) does not imply trivial
\(\mathcal{S}_\omega[\mathfrak{D}]\): a pure outgoing propagating state may have
\(\mathbf{u}_{\mathrm{rad}}=\mathbf{u}_\omega\) \cite{stratton1941}. Section~\ref{sec:ch272}
classifies \(\ker\mathcal{R}_\omega\); Subsection~\ref{sec:ch254} only identifies the
field-side split Eq.~\eqref{eq:ch2.nonradiative-component-def} as the operator image of
non-export content \cite{harrington2001}.

\begin{proposition}[Far-zone radiative dominance]
\label{prop:ch2.far-zone-radiative-dominance}
Let \(\mathcal{D}_{J,\mathrm{reg}}=1\) and let \(r\) lie in the far zone of
Eq.~\eqref{eq:ch2.distance-regime-definitions}. Then
\begin{equation}
\label{eq:ch2.far-zone-component-dominance}
\frac{\mathcal{E}_{\mathrm{nr}}[\mathfrak{D}]}{\mathcal{E}_{\mathrm{rad}}[\mathfrak{D}]+\epsilon}
\to 0,
\qquad
\frac{|\mathcal{P}_{\mathrm{nr}}(r)|}{\mathcal{P}_{\mathrm{rad}}(r)+\epsilon}\to 0,
\end{equation}
as \(r\to\infty\) for fixed \(\mathfrak{D}\), with \(\epsilon>0\) regularizing zero export
\cite{stratton1941}. Equivalently, \(\mathbf{u}_\omega\to\mathbf{u}_{\mathrm{rad}}\) in
far-zone observables metered by Eqs.~\eqref{eq:ch2.far-field-integral-reduction}--\eqref{eq:ch2.power-flux-functional}.
\end{proposition}
\begin{proof}
Eq.~\eqref{eq:ch2.rad-null-on-evanescent} annihilates evanescent and non-propagating
contributions in \(\Pi_{\mathrm{rad}}\). Far-zone asymptotics
Theorem~\ref{thm:ch2.far-field-approximation} and Proposition~\ref{prop:ch2.far-zone-flux-stabilization}
give \(\mathcal{P}_{\mathrm{rad}}(r)\to P_\infty>0\) while reactive terms scale as
Eq.~\eqref{eq:ch2.near-zone-field-scaling} and decay relatively on large shells
\cite{stratton1941,jackson1999,harrington2001}. Regime condition
Eq.~\eqref{eq:ch1.radiative-regime-condition} removes reactive dominance in the limit.
\end{proof}

\paragraph{Cutoff and non-propagating modes.}
Below guide cutoff Eq.~\eqref{eq:ch2.cutoff-condition}, modal content lies in
\(\Pi_{\mathrm{non}}\mathcal{S}_\omega[\mathfrak{D}]\) or
\(\Pi_{\mathrm{ev}}\mathcal{S}_\omega[\mathfrak{D}]\) \cite{stratton1941}. Applying
Eq.~\eqref{eq:ch2.far-field-amplitude-integral} or transport indices from
Subsection~\ref{sec:ch253} to such modes violates
Eq.~\eqref{eq:ch2.regime-component-gate} as in the TE\(_{10}\) illustration of
Section~\ref{sec:ch01.1.2} \cite{harrington2001}. Subsection~\ref{sec:ch254} therefore
pairs spectral projectors from Subsection~\ref{sec:ch224} with regime logic from
Section~\ref{sec:ch01.2.4} on the \emph{same} integral outputs \cite{stratton1941}.

\paragraph{Boundary of validity.}
Equations~\eqref{eq:ch2.integral-spectral-decomposition}--\eqref{eq:ch2.far-zone-component-dominance}
require linear media, fixed \(\omega\), outgoing closure, and spectral completion sufficient
for Eqs.~\eqref{eq:ch2.spectral-projectors}--\eqref{eq:ch2.spectral-projector-algebra}
\cite{stratton1941,harrington2001}. Sharp boundaries between subspaces may blur when
loss, anisotropy, or coupled interfaces mix propagating and evanescent sheets
\cite{stratton1941}. Component splits are diagnostic on shells, not substitutes for
Theorem~\ref{thm:ch2.exterior-radiative-uniqueness} on the full field
\cite{harrington2001}.

\paragraph{Worked illustration (near-field scan).}
A probe at \(k_0 r=1\) measures \(\mathbf{u}_\omega=\mathbf{u}_{\mathrm{rad}}+\mathbf{u}_{\mathrm{nr}}\)
with \(|\mathbf{u}_{\mathrm{nr}}|\gg|\mathbf{u}_{\mathrm{rad}}|\) for an electrically small
source \cite{stratton1941}. Declaring transport from the scan without
Eq.~\eqref{eq:ch2.regime-component-gate} violates Section~\ref{sec:ch01.2.4}
\cite{harrington2001}. At \(k_0 r=10\) with \(\mathcal{D}_{J,\mathrm{reg}}=1\),
Proposition~\ref{prop:ch2.far-zone-radiative-dominance} applies \cite{stratton1941}.

\paragraph{Worked illustration (evanescent surface mode).}
An evanescent surface wave satisfies \(\Pi_{\mathrm{ev}}\mathbf{u}_\omega=\mathbf{u}_\omega\)
and \(\mathbf{u}_{\mathrm{rad}}=\mathbf{0}\) through Eq.~\eqref{eq:ch2.rad-null-on-evanescent}
\cite{stratton1941,harrington2001}. Power meters from Subsection~\ref{sec:ch253} on
\(\mathcal{S}_\omega[\mathfrak{D}]\) without \(\Pi_{\mathrm{rad}}\) report reactive circulation,
not export \cite{stratton1941}.

\paragraph{Worked illustration (two-source superposition).}
For admissible \(\Jimp^{(1)},\Jimp^{(2)}\),
Eq.~\eqref{eq:ch2.radiation-linearity} gives
\(\mathcal{R}_\omega[\Jimp^{(1)}+\Jimp^{(2)}]=\mathcal{R}_\omega[\Jimp^{(1)}]+\mathcal{R}_\omega[\Jimp^{(2)}]\),
but \(\mathbf{u}_{\mathrm{nr}}\) is not generally additive on raw near-field sums unless
both components are split after \(\mathcal{S}_\omega\) \cite{stratton1941,harrington2001}.
Transport claims require radiative superposition, not reactive interference
(Section~\ref{sec:ch01.6.4}) \cite{stratton1941}.

\paragraph{Closing summary.}
Subsection~\ref{sec:ch254} records the Chapter~\ref{ch:02} operator split via
Eqs.~\eqref{eq:ch2.integral-spectral-decomposition}--\eqref{eq:ch2.nonpropagating-component-def},
radiative/non-radiative definitions
Eqs.~\eqref{eq:ch2.radiative-component-def}--\eqref{eq:ch2.nonradiative-component-def},
energy and power budgets Eqs.~\eqref{eq:ch2.radiative-energy-functional}--\eqref{eq:ch2.nonradiative-power-functional},
regime gate Eq.~\eqref{eq:ch2.regime-component-gate}, null-space link
Eq.~\eqref{eq:ch2.nr-image-of-null-space}, and far-zone dominance
Proposition~\ref{prop:ch2.far-zone-radiative-dominance}, importing
Section~\ref{sec:ch01.2.4} and Subsection~\ref{sec:ch224} by reference
\cite{stratton1941,harrington2001,jackson1999}.

\medskip
\noindent
Section~\ref{sec:ch25} closes with the evaluation chain begun in its introduction:
outgoing source-to-field integrals Theorem~\ref{thm:ch2.source-to-field-integral}
(Subsection~\ref{sec:ch251}), distance-regime asymptotics and far-field reduction
Theorem~\ref{thm:ch2.far-field-approximation} (Subsection~\ref{sec:ch252}), flux and
operator-norm meters on \(\mathcal{R}_\omega[\Jimp]\) (Subsection~\ref{sec:ch253}), and
radiative--non-radiative component split
Eqs.~\eqref{eq:ch2.radiative-component-def}--\eqref{eq:ch2.regime-component-gate}
(Subsection~\ref{sec:ch254}) \cite{stratton1941,harrington2001,jackson1999}. Green
kernels and outgoing closure from Sections~\ref{sec:ch23}--\ref{sec:ch24} are now tied to
evaluable far-zone observables: \(\mathcal{S}_\omega\) is single-valued on
\(\mathcal{F}_{\mathrm{SF}}\), \(\mathcal{R}_\omega\) selects export, and flux functionals
Eqs.~\eqref{eq:ch2.power-flux-functional}--\eqref{eq:ch2.angular-momentum-flux-vector} act
only after regime gate Eq.~\eqref{eq:ch2.regime-operator-screen}
\cite{stratton1941}. Section~\ref{sec:ch26} develops spectral decomposition of radiation
operators; Section~\ref{sec:ch27} completes non-radiating null spaces; Chapter~\ref{ch:03}
supplies vector spherical harmonic separation of \(\E_0(\theta,\phi)\) from
Eq.~\eqref{eq:ch2.far-field-amplitude-integral} \cite{harrington2001}. Transport and
angular-transport claims throughout Volume~I must compose with
Eq.~\eqref{eq:ch2.section25-evaluation-chain} and
Eq.~\eqref{eq:ch2.regime-component-gate}, not with raw \(\mathcal{S}_\omega[\Jimp]\) near
sources \cite{stratton1941,harrington2001}.

\section{Spectral Theory of Radiation Operators}
\label{sec:ch26}

Section~\ref{sec:ch25} closed the evaluation layer for outgoing radiation: integral
representations, distance-regime asymptotics, flux meters on
\(\mathcal{R}_\omega[\Jimp]\), and the radiative--non-radiative split
Eqs.~\eqref{eq:ch2.radiative-component-def}--\eqref{eq:ch2.regime-component-gate}
\cite{stratton1941,harrington2001,jackson1999}. Those constructions answer how to
\emph{evaluate} qualified observables once \(\mathcal{S}_\omega\) and
\(\Pi_{\mathrm{rad}}\) are fixed. Section~\ref{sec:ch26} addresses the complementary
question of \emph{spectral organization}: which modes \(\mathcal{R}_\omega\) retains,
how eigenstructure compresses or expands angular content, when Fredholm solvability
permits inversion, and which geometric symmetries constrain admissible radiative
spectra \cite{stratton1941,harrington2001}. Compactness templates, spectral compression,
and Fredholm splits from Subsection~\ref{sec:ch213} are presupposed; propagation
projectors from Subsection~\ref{sec:ch224} and far-zone amplitudes from
Eq.~\eqref{eq:ch2.far-field-amplitude-integral} supply the field-side objects on which
spectral analysis acts \cite{stratton1941}. Rotation and covariance rules from
Section~\ref{sec:ch01.3} enter where degeneracy is classified; they are not re-derived
\cite{stratton1941}.

The problem addressed here is modal resolution under outgoing closure. On bounded
cavities, Rellich compactness Eq.~\eqref{eq:ch2.embedding-chain} yields accumulating
discrete eigenvalues and Eq.~\eqref{eq:ch2.spectral-compression-operator} truncates
interior resonances \cite{stratton1941,harrington2001}. On exterior domains with
compact sources, outgoing selection Eq.~\eqref{eq:ch2.outgoing-radiation-refinement}
replaces discrete cavity sums by continuous angular spectra tied to
Eq.~\eqref{eq:ch2.far-field-amplitude-integral} \cite{jackson1999,stratton1941}.
Using cavity modal bookkeeping on an open antenna conflates storage resonances with
export modes and misstates rank of \(\mathcal{R}_\omega\) even when forward integrals
from Theorem~\ref{thm:ch2.source-to-field-integral} are exact
\cite{harrington2001}. Section~\ref{sec:ch26} separates these domain classes, states
the spectral targets for \(\mathcal{R}_\omega\), \(\mathcal{S}_\omega\), and
\(\mathcal{L}_\omega\) on their declared domains, and forwards complete null-space
classification to Section~\ref{sec:ch27} \cite{stratton1941}.

\paragraph{Spectral development chain.}
Let \(\Omega_s=\operatorname{supp}\Jimp\) be compact with scale \(a\), vacuum wavenumber
\(k_0=\omega\sqrt{\varepsilon_0\mu_0}\), and outgoing field pair
\(\mathbf{u}_\omega=(\E,\H)\in\mathcal{F}_{\mathrm{SF}}\) from
Eq.~\eqref{eq:ch2.admissible-field-space}. Section~\ref{sec:ch26} organizes analysis in
four stages:
\begin{equation}
\label{eq:ch2.section26-spectral-chain}
\mathcal{R}_\omega
\ \xrightarrow{\ \text{eigenmodes}\ }
\{\lambda_n,\mathbf{u}_n\}
\ \xrightarrow{\ \text{symmetry}\ }
\text{degeneracy classes}
\ \xrightarrow{\ \text{Fredholm}\ }
\text{solvability / index}
\ \xrightarrow{\ \text{geometry}\ }
\text{admissible spectra},
\end{equation}
where the first arrow is Subsection~\ref{sec:ch261}, the second imports
Section~\ref{sec:ch01.3} in Subsection~\ref{sec:ch262}, the third combines
Subsection~\ref{sec:ch213} with Subsection~\ref{sec:ch263}, and the fourth is
Subsection~\ref{sec:ch264} \cite{stratton1941,harrington2001}. Equation
\eqref{eq:ch2.section26-spectral-chain} restates
Eq.~\eqref{eq:ch2.section25-evaluation-chain} at the modal level: evaluation meters from
Section~\ref{sec:ch25} act on outputs of \(\mathcal{R}_\omega\); spectral theory explains
which modal coordinates those outputs span and which remain invisible to far-zone
observation \cite{stratton1941}. Applying flux functionals
Eq.~\eqref{eq:ch2.power-flux-functional} to raw \(\mathcal{S}_\omega[\Jimp]\) without
spectral qualification still violates regime gate
Eq.~\eqref{eq:ch2.regime-component-gate} \cite{harrington2001}.

\paragraph{Spectral operator targets.}
The primary radiation map remains
\begin{equation}
\label{eq:ch2.spectral-operator-targets}
\mathcal{R}_\omega:\mathcal{J}_\Omega\longrightarrow\mathcal{F}_{\mathrm{rad}},
\qquad
\mathcal{R}_\omega[\Jimp]=\Pi_{\mathrm{rad}}\,\mathcal{S}_\omega[\Jimp],
\end{equation}
from Eq.~\eqref{eq:ch2.radiation-operator-def}, with \(\Pi_{\mathrm{rad}}\) the outgoing
radiation projector Eq.~\eqref{eq:ch2.rad-projection-composition}
\cite{stratton1941,harrington2001}. Secondary targets include the volume Maxwell operator
\(\mathcal{L}_\omega\) on \(\mathcal{D}(\mathcal{L}_\omega)\) for solvability and the
compressed map Eq.~\eqref{eq:ch2.spectral-compression-operator} for rank-\(N\) modal
surrogates \cite{stratton1941}. Green split
Eq.~\eqref{eq:ch2.green-compact-split} supplies compact perturbations in Fredholm
formulations; kernel construction is not repeated \cite{harrington2001}. Source-side
null space \(\ker\mathcal{R}_\omega\) from Eq.~\eqref{eq:ch2.ker-R-omega-preview} is
spectral data for \(\mathcal{R}_\omega\) but its complete classification belongs to
Section~\ref{sec:ch272} \cite{stratton1941}.

\paragraph{Modal eigenvalue template (preview).}
On a declared spectral family \(\{\psi_n\}\subset\mathcal{F}_{\mathrm{rad}}\) with
associated one-dimensional projectors \(P_n\), Subsection~\ref{sec:ch261} develops the
Chapter~\ref{ch:02} HOME eigenvalue law
\begin{equation}
\label{eq:ch2.modal-eigenvalue-equation-preview}
\mathcal{R}_\omega[\Jimp^{(n)}]=\lambda_n\,\mathbf{u}_n,
\qquad
\mathbf{u}_n=P_n\mathcal{S}_\omega[\Jimp^{(n)}],
\end{equation}
with \(\Jimp^{(n)}\in\mathcal{J}_\Omega\) admissible sources generating mode \(\mathbf{u}_n\)
\cite{stratton1941,harrington2001}. Equation~\eqref{eq:ch2.modal-eigenvalue-equation-preview}
specializes Eq.~\eqref{eq:ch2.spectral-compression-operator} before eigenfunctions are
constructed: \(\lambda_n\) measures radiative export efficiency of mode \(n\), while
\(\ker\mathcal{R}_\omega\) corresponds to \(\lambda_n=0\) with
\(\mathcal{S}_\omega[\Jimp^{(n)}]\neq\mathbf{0}\) \cite{stratton1941}. Far-zone reduction
Theorem~\ref{thm:ch2.far-field-approximation} links modal amplitudes to angular patterns
through Eq.~\eqref{eq:ch2.far-field-amplitude-integral} \cite{jackson1999,harrington2001}.

\paragraph{Discrete versus continuous spectra (preview).}
Domain class fixes spectrum type:
\begin{equation}
\label{eq:ch2.discrete-continuous-spectrum-preview}
\begin{aligned}
\text{bounded cavity:}\ &
\{\lambda_n\}_{n=1}^{\infty}\ \text{discrete, accumulating};\\
\text{exterior + outgoing:}\ &
\text{continuous angular spectrum on }S^{2}\ \text{for }\E_0(\theta,\phi);
\end{aligned}
\end{equation}
with \(\E_0\) the far amplitude from Eq.~\eqref{eq:ch2.far-field-amplitude-integral}
\cite{stratton1941,jackson1999}. Equations~\eqref{eq:ch2.discrete-continuous-spectrum-preview}
explain why Eq.~\eqref{eq:ch2.spectral-compression-operator} sums over \(n\) in cavities
but requires rigged-spectrum integrals on open domains \cite{harrington2001}. Vector
spherical harmonic separation of \(\E_0\) is deferred to Chapter~\ref{ch:03}; Section~\ref{sec:ch26}
states operator-level spectral targets without re-displaying \(Y_\ell^m\) definitions
\cite{stratton1941}.

\paragraph{Far-pattern modal link (preview).}
When \(k_0 r\gg k_0 a\) and \(\mathcal{D}_{J,\mathrm{reg}}=1\), far-zone patterns decompose
as
\begin{equation}
\label{eq:ch2.far-pattern-modal-preview}
\E_0(\theta,\phi)
=
\sum_{n\in\mathcal{N}_{\mathrm{rad}}}
a_n(\Jimp)\,\mathbf{f}_n(\theta,\phi),
\end{equation}
with \(\mathcal{N}_{\mathrm{rad}}\) the radiative index set fixed in
Subsection~\ref{sec:ch261}, \(\mathbf{f}_n\) outgoing modal shapes, and coefficients
\(a_n\) linear in \(\Jimp\) through Eq.~\eqref{eq:ch2.radiation-linearity}
\cite{stratton1941,harrington2001,jackson1999}. Equation~\eqref{eq:ch2.far-pattern-modal-preview}
is the angular restatement of Eq.~\eqref{eq:ch2.far-field-amplitude-integral}; it does not
assert finiteness of \(\mathcal{N}_{\mathrm{rad}}\) on exterior problems
\cite{stratton1941}. Observation rank cap Eq.~\eqref{eq:ch2.reconstruction-rank-cap}
limits recoverable \(\{a_n\}\) even when Eq.~\eqref{eq:ch2.far-pattern-modal-preview} is
exact \cite{harrington2001}.

\begin{figure}[t]
\centering
\ChOneFigBlock{%
\begin{tikzpicture}[ch1fig, node distance=7mm and 10mm]
  \node[ch1box] (eig) {Eigenmodes\\$\{\lambda_n,\mathbf{u}_n\}$};
  \node[ch1box, right=13mm of eig] (deg) {Rotational\\degeneracy};
  \node[ch1box, right=13mm of deg] (fr) {Fredholm\\solvability};
  \node[ch1box, right=13mm of fr] (geo) {Geometric\\spectrum constraints};
  \node[ch1box, below=12mm of eig] (r) {$\mathcal{R}_\omega$};
  \node[ch1box, below=12mm of deg] (so) {$\mathrm{SO}(3)$};
  \draw[ch1flow] (r) -- (eig) -- (deg) -- (fr) -- (geo);
  \draw[ch1flow] (so) -- (deg);
\end{tikzpicture}%
}
\caption{Section~\ref{sec:ch26} development path: spectral decomposition of
\(\mathcal{R}_\omega\) (Subsection~\ref{sec:ch261}), rotational degeneracy via
Section~\ref{sec:ch01.3} (Subsection~\ref{sec:ch262}), Fredholm solvability importing
Subsection~\ref{sec:ch213} (Subsection~\ref{sec:ch263}), and geometric constraints on
radiative spectra (Subsection~\ref{sec:ch264}).}
\label{fig:ch2.spectral-theory-roadmap}
\end{figure}

Figure~\ref{fig:ch2.spectral-theory-roadmap} orders the four subsections.
Subsection~\ref{sec:ch261} gives the Chapter~\ref{ch:02} HOME for eigenfunctions and
spectral decomposition of radiation operators on outgoing subspaces, upgrading
Eqs.~\eqref{eq:ch2.spectral-compression-operator}--\eqref{eq:ch2.modal-eigenvalue-equation-preview}
from templates to constructive statements compatible with
Eq.~\eqref{eq:ch2.far-field-amplitude-integral} \cite{stratton1941,harrington2001}.
Subsection~\ref{sec:ch262} applies rotation covariance from
Section~\ref{sec:ch01.3} to classify spectral degeneracy without re-proving
Eqs.~\eqref{eq:ch1.active-rotation-action}--\eqref{eq:ch1.passive-rotation-action}
\cite{stratton1941}. Subsection~\ref{sec:ch263} develops Fredholm structure and
solvability criteria for radiation-related inversions, importing compactness and split
Eq.~\eqref{eq:ch2.fredholm-split} from Subsection~\ref{sec:ch213} by reference
\cite{harrington2001}. Subsection~\ref{sec:ch264} records geometric constraints---symmetry,
support, and boundary class---that restrict admissible radiative spectra and connect to
angular accessibility in Section~\ref{sec:ch294} \cite{stratton1941}.

\paragraph{Fredholm solvability (preview).}
Radiation-related inversions take the form \(\mathcal{R}_\omega[\Jimp]=\mathbf{g}\) or
finite-rank observations \(\mathcal{O}_{\mathrm{meas}}\mathcal{R}_\omega[\Jimp]=\mathbf{b}\).
Subsection~\ref{sec:ch263} states solvability in Fredholm form
\begin{equation}
\label{eq:ch2.spectral-solvability-preview}
\mathcal{R}_\omega+K
\ \text{solvable on}\ \mathcal{J}_\Omega
\iff
\mathbf{g}\perp\ker(\mathcal{R}_\omega+K)^{\dagger},
\end{equation}
with \(K\) compact on the graph domain of impressed sources and \(\dagger\) denoting
adjoint on declared pairings Eq.~\eqref{eq:ch2.dual-source-pairing}
\cite{stratton1941,harrington2001}. Equation~\eqref{eq:ch2.spectral-solvability-preview}
specializes Eq.~\eqref{eq:ch2.fredholm-split} to export maps; index computation and
cokernel characterization are HOME in Subsection~\ref{sec:ch263}, not repeated from
Subsection~\ref{sec:ch213} \cite{stratton1941}. Rank composition
Eq.~\eqref{eq:ch2.rank-composition-law} and observation null space
Eq.~\eqref{eq:ch2.observation-null-space} remain compositional constraints on any
solvability claim \cite{harrington2001}.

\paragraph{Degeneracy and covariance (preview).}
When a symmetry group \(G\) commutes with \(\mathcal{R}_\omega\), modal eigenspaces carry
representations of \(G\). For rotations \(R\in\mathrm{SO}(3)\),
\begin{equation}
\label{eq:ch2.degeneracy-covariance-preview}
\mathcal{R}_\omega\,\mathcal{D}(R)[\Jimp]
=
\mathcal{D}(R)\,\mathcal{R}_\omega[\Jimp],
\end{equation}
with \(\mathcal{D}(R)\) the field representation induced from
Eq.~\eqref{eq:ch1.active-rotation-action} on \(\mathcal{F}_{\mathrm{rad}}\)
\cite{stratton1941}. Equation~\eqref{eq:ch2.degeneracy-covariance-preview} is the
operator restatement of Section~\ref{sec:ch01.3} covariance tests; degeneracy multiplicities
follow from representation theory in Subsection~\ref{sec:ch262} without re-opening origin
shifts or tensor transformation laws \cite{stratton1941,harrington2001}. Transport indices
from Section~\ref{sec:ch01.2.3} apply only to degeneracy classes that pass
Eq.~\eqref{eq:ch2.regime-operator-screen} \cite{stratton1941}.

\paragraph{Geometric constraints (preview).}
Source support, boundary class, and material symmetry restrict which modal coefficients
\(a_n\) in Eq.~\eqref{eq:ch2.far-pattern-modal-preview} can be nonzero:
\begin{equation}
\label{eq:ch2.geometric-spectrum-constraint-preview}
a_n(\Jimp)=0
\quad\text{when}\quad
\mathcal{C}_{\mathrm{geo}}(\Omega_s,\Gamma,\varepsilon,\mu)\,\mathbf{f}_n=\mathbf{0},
\end{equation}
with \(\mathcal{C}_{\mathrm{geo}}\) a geometric symmetry operator fixed in
Subsection~\ref{sec:ch264} \cite{stratton1941,harrington2001}. Equation
\eqref{eq:ch2.geometric-spectrum-constraint-preview} explains nulls in measured patterns
that are not rank deficits of \(\mathcal{O}_{\mathrm{meas}}\): forbidden parity or mirror
symmetry can annihilate entire modal families even when
Eq.~\eqref{eq:ch2.spectral-solvability-preview} is satisfied for other modes
\cite{stratton1941}. Complete reciprocity constraints on coupled modes belong to
Section~\ref{sec:ch28} \cite{harrington2001}.

\paragraph{Spectral resolution composition.}
Modal projectors compose with propagation and radiation selectors already fixed:
\begin{equation}
\label{eq:ch2.spectral-resolution-composition}
\Pi_{\mathrm{rad}}
=
\sum_{n\in\mathcal{N}_{\mathrm{rad}}} P_n,
\qquad
P_n\Pi_{\mathrm{ev}}=P_n\Pi_{\mathrm{non}}=\mathbf{0},
\end{equation}
using Eqs.~\eqref{eq:ch2.spectral-projectors}--\eqref{eq:ch2.rad-null-on-evanescent} from
Subsection~\ref{sec:ch224} by reference \cite{stratton1941,harrington2001}. Equation
\eqref{eq:ch2.spectral-resolution-composition} ties Section~\ref{sec:ch26} to the
radiative--non-radiative split Eq.~\eqref{eq:ch2.nonradiative-component-def}: modes with
\(\lambda_n=0\) lie entirely in \(\mathbf{u}_{\mathrm{nr}}\) \cite{stratton1941}. Energy
inequality Eq.~\eqref{eq:ch2.subspace-energy-inequality} bounds how much reactive content
can masquerade as modal amplitude in the near zone before far-zone dominance
Proposition~\ref{prop:ch2.far-zone-radiative-dominance} applies \cite{harrington2001}.

Several constructions deliberately remain outside Section~\ref{sec:ch26}. Vector spherical
harmonic definitions, normalization, and completeness belong to Chapter~\ref{ch:03};
complete non-radiating source catalogs to Sections~\ref{sec:ch271}--\ref{sec:ch272};
reciprocity and modal orthogonality to Section~\ref{sec:ch28}; plane-wave angular-spectrum
filters to Section~\ref{sec:ch29} \cite{stratton1941,harrington2001}. Section~\ref{sec:ch26}
does not re-derive Green kernels, Sommerfeld closure, or radiation integrals from
Sections~\ref{sec:ch23}--\ref{sec:ch25}; it organizes modal structure on the maps those
sections fixed \cite{stratton1941}. Rigged Hilbert space extensions for continuous spectra
are previewed here and developed in Chapter~\ref{ch:03} Subsection~\ref{sec:ch313}
\cite{harrington2001}.

Regime logic from Section~\ref{sec:ch01.2.4} is unchanged. Spectral mode labels do not
override transport qualification: a named mode with \(\lambda_n>0\) still requires
\(\mathcal{D}_{J,\mathrm{reg}}=1\) and far-zone reduction before flux meters
Eqs.~\eqref{eq:ch2.power-flux-functional}--\eqref{eq:ch2.transport-triad-on-R} report
export \cite{stratton1941,harrington2001}. Compression error
Eq.~\eqref{eq:ch2.compression-error-bound} and singular-value decay
Eq.~\eqref{eq:ch2.singular-value-decay} bound modal truncation independently of regime
gates but must be combined with Eq.~\eqref{eq:ch2.regime-component-gate} when reactive
near fields dominate local probes \cite{stratton1941}.

Loss and dispersive materials modify \(\lambda_n\) and modal quality factors at fixed
\(\omega\) without changing the spectral template: Eq.~\eqref{eq:ch2.modal-eigenvalue-equation-preview}
uses outgoing \(\mathcal{R}_\omega\) with complex \(k(\mathbf{r})\) from
Eq.~\eqref{eq:ch2.helmholtz-symbol} \cite{stratton1941}. Passive dissipation may merge
nearby eigenvalues in cavities but does not convert evanescent storage into far-zone export
through Eq.~\eqref{eq:ch2.rad-null-on-evanescent} \cite{harrington2001}. Broadband
superposition across \(\omega\) is deferred to the causal chain of
Subsection~\ref{sec:ch223}; Section~\ref{sec:ch26} works at fixed harmonic frequency
\cite{stratton1941}.

\paragraph{Worked illustration (cavity versus exterior modal books).}
A closed PEC cavity supports discrete \(\{\lambda_n\}\) with
Eq.~\eqref{eq:ch2.spectral-compression-operator} converging in operator norm on bounded
source classes \cite{stratton1941}. An exterior dipole source radiates with continuous
angular content in Eq.~\eqref{eq:ch2.far-pattern-modal-preview}; truncating that content
to the first \(N\) cavity modes of a enclosing box misstates sidelobe structure and
underestimates rank of \(\mathcal{O}_{\mathrm{meas}}\mathcal{R}_\omega\)
\cite{harrington2001,jackson1999}. Subsection~\ref{sec:ch261} states the correct spectral
bookkeeping for each domain class \cite{stratton1941}.

\paragraph{Worked illustration (degenerate multipole family).}
For an \(\ell\)-indexed modal family on \(\mathcal{F}_{\mathrm{rad}}\), rotation covariance
Eq.~\eqref{eq:ch2.degeneracy-covariance-preview} implies \((2\ell+1)\)-fold degeneracy when
the family is closed under \(\mathrm{SO}(3)\) \cite{stratton1941}. Laboratory rotation of
the antenna rotates measured coefficients through Section~\ref{sec:ch01.3} rules without
changing exported power when Eq.~\eqref{eq:ch2.power-flux-functional} is evaluated on
\(\mathcal{R}_\omega[\Jimp]\) in the far zone \cite{harrington2001,jackson1999}.
Subsection~\ref{sec:ch262} classifies such degeneracies; Chapter~\ref{ch:03} supplies explicit
harmonic bases \cite{stratton1941}.

\paragraph{Worked illustration (symmetry-forbidden mode).}
A source distribution symmetric under mirror reflection about a plane may force
\(a_n=0\) for modes odd under that reflection through
Eq.~\eqref{eq:ch2.geometric-spectrum-constraint-preview} \cite{stratton1941}. A probe
array that cannot resolve the symmetric/antisymmetric split may still report low pattern
distortion while missing entire modal sectors; Eq.~\eqref{eq:ch2.reconstruction-rank-cap}
and geometric constraint Eq.~\eqref{eq:ch2.geometric-spectrum-constraint-preview} must
both appear in identifiability budgets \cite{harrington2001}. Subsection~\ref{sec:ch264}
develops the geometric operators \(\mathcal{C}_{\mathrm{geo}}\) \cite{stratton1941}.

\paragraph{Worked illustration (Fredholm obstruction).}
Given target far pattern \(\mathbf{g}\in\mathcal{F}_{\mathrm{rad}}\), inversion
\(\mathcal{R}_\omega[\Jimp]=\mathbf{g}\) fails when \(\mathbf{g}\) has nonzero pairing with
\(\ker\mathcal{R}_\omega^{\dagger}\) even if \(\mathbf{g}\) is smooth on \(S^{2}\)
\cite{stratton1941,harrington2001}. Equation~\eqref{eq:ch2.spectral-solvability-preview}
is the operator form of that obstruction; Subsection~\ref{sec:ch263} computes the
relevant cokernels on declared domains \cite{stratton1941}. Non-radiating components of
\(\Jimp\) remain invisible to \(\mathcal{R}_\omega\) through
Eq.~\eqref{eq:ch2.ker-R-omega-preview} \cite{harrington2001}.

The subsection sequence begins with eigenfunctions and spectral decomposition of
\(\mathcal{R}_\omega\) on outgoing subspaces, continues with rotational degeneracy imported
from Section~\ref{sec:ch01.3}, develops Fredholm solvability using compactness from
Subsection~\ref{sec:ch213}, and closes with geometric constraints on admissible radiative
spectra before null-space classification in Section~\ref{sec:ch27} and reciprocity in
Section~\ref{sec:ch28} compose additional structure on the same modal coordinates
\cite{stratton1941,harrington2001,jackson1999}.

\subsection{Eigenfunctions and spectral decomposition}
\label{sec:ch261}
% TOC tag: [HOME]

Section~\ref{sec:ch26} previewed spectral organization of \(\mathcal{R}_\omega\) through
Eqs.~\eqref{eq:ch2.section26-spectral-chain}--\eqref{eq:ch2.spectral-resolution-composition}.
Subsection~\ref{sec:ch261} gives the Chapter~\ref{ch:02} HOME for eigenpairs, modal
projectors, and spectral expansions of radiation operators on outgoing subspaces after
Sections~\ref{sec:ch23}--\ref{sec:ch25} \cite{stratton1941,harrington2001,jackson1999}.
Propagation projectors Eqs.~\eqref{eq:ch2.spectral-projectors}--\eqref{eq:ch2.rad-null-on-evanescent}
from Subsection~\ref{sec:ch224}, compactness and compression templates from
Subsection~\ref{sec:ch213}, and far-zone amplitudes
Eqs.~\eqref{eq:ch2.far-field-integral-reduction}--\eqref{eq:ch2.far-field-amplitude-integral}
from Subsection~\ref{sec:ch252} are presupposed; projector algebra and Green construction are
not re-proved \cite{stratton1941}. Explicit vector spherical harmonic bases and rigged-space
completeness belong to Chapter~\ref{ch:03}; here modal coordinates are fixed at the operator
level compatible with Eq.~\eqref{eq:ch2.far-pattern-modal-preview}
\cite{harrington2001}. Rotational degeneracy is deferred to Subsection~\ref{sec:ch262};
Fredholm solvability to Subsection~\ref{sec:ch263}.

\paragraph{Assumptions.}
Fix \(\omega>0\), admissible impressed data \(\mathfrak{D}\) on
\(\mathcal{J}_\Omega\) Eq.~\eqref{eq:ch2.source-space-refined}, outgoing closure on exterior
problems Eqs.~\eqref{eq:ch2.sommerfeld-outgoing-class}--\eqref{eq:ch2.outgoing-selector-def},
and radiation map Eq.~\eqref{eq:ch2.spectral-operator-targets}
\cite{stratton1941,harrington2001}. Field pairs lie in
\(\mathcal{F}_{\mathrm{SF}}\subset\mathcal{F}_\Omega\) from
Eq.~\eqref{eq:ch2.admissible-field-space}; radiative images lie in
\(\mathcal{F}_{\mathrm{rad}}\subset\operatorname{ran}\mathcal{R}_\omega\) with energy norm
Eq.~\eqref{eq:ch2.field-energy-norm} \cite{stratton1941}. Spectral completion
\(\widehat{\mathcal{F}}_\Omega\) supporting Eqs.~\eqref{eq:ch2.spectral-projectors}--\eqref{eq:ch2.spectral-projector-algebra}
is assumed sufficient for the expansions below on the physical branch
\cite{harrington2001}.

\paragraph{Radiation eigenpairs.}
A \emph{radiation eigenpair} is a triple \((\lambda_n,\mathbf{u}_n,\Jimp^{(n)})\) with
\(\lambda_n\in\mathbb{C}\), \(\mathbf{u}_n=(\E_n,\H_n)\in\mathcal{F}_{\mathrm{rad}}\), and
\(\Jimp^{(n)}\in\mathcal{J}_\Omega\) such that
\begin{equation}
\label{eq:ch2.radiation-eigenpair-def}
\mathcal{R}_\omega[\Jimp^{(n)}]=\lambda_n\,\mathbf{u}_n,
\qquad
\|\mathbf{u}_n\|_{\mathcal{U}}=1,
\end{equation}
where \(\|\cdot\|_{\mathcal{U}}\) is the field energy norm
Eq.~\eqref{eq:ch2.field-energy-norm} on \(\mathcal{F}_{\mathrm{rad}}\)
\cite{stratton1941,harrington2001}. Equation~\eqref{eq:ch2.radiation-eigenpair-def}
upgrades preview Eq.~\eqref{eq:ch2.modal-eigenvalue-equation-preview} to a normalized
HOME definition: \(\lambda_n\) measures export gain of mode \(n\) relative to unit radiative
energy, and \(\ker\mathcal{R}_\omega\) corresponds to \(\lambda_n=0\) with
\(\mathcal{S}_\omega[\Jimp^{(n)}]\neq\mathbf{0}\) through
Eq.~\eqref{eq:ch2.ker-R-omega-preview} \cite{stratton1941}. Eigenpairs are defined on the
outgoing branch only; incoming propagating sheets excluded by
Eq.~\eqref{eq:ch2.outgoing-radiation-refinement} \cite{harrington2001}.

\paragraph{Modal projectors on field space.}
For each index \(n\) in a declared spectral family \(\mathcal{N}\), define a rank-one (or
finite-dimensional) projector \(P_n:\mathcal{F}_\Omega\to\mathcal{F}_{\mathrm{rad}}\) by
\begin{equation}
\label{eq:ch2.modal-projector-def}
P_n[\mathbf{u}]
=
\langle \mathbf{u},\mathbf{u}_n\rangle_{\mathcal{U}}\,\mathbf{u}_n,
\qquad
P_n\circ\Pi_{\mathrm{ev}}=P_n\circ\Pi_{\mathrm{non}}=0,
\end{equation}
using dual pairing inherited from Eq.~\eqref{eq:ch2.maxwell-adjoint-pairing} on admissible
states and annihilation of non-propagating content through
Eq.~\eqref{eq:ch2.rad-null-on-evanescent} \cite{stratton1941,harrington2001}. Equation
\eqref{eq:ch2.modal-projector-def} specializes resolution
Eq.~\eqref{eq:ch2.spectral-resolution-composition} once an orthonormal modal family
\(\{\mathbf{u}_n\}_{n\in\mathcal{N}_{\mathrm{rad}}}\) is chosen. Projectors satisfy
\(P_n P_m=\delta_{nm} P_n\) when modes are orthogonal in \(\mathcal{U}\)
\cite{stratton1941}.

\paragraph{Source modal amplitudes.}
For general admissible \(\mathfrak{D}\), define source coefficients
\begin{equation}
\label{eq:ch2.source-modal-amplitude}
c_n[\mathfrak{D}]
=
\langle \mathcal{S}_\omega[\mathfrak{D}],\mathbf{u}_n\rangle_{\mathcal{U}},
\qquad
\mathbf{u}_{\mathrm{rad},n}
=
P_n\,\mathcal{R}_\omega[\mathfrak{D}]
=
\lambda_n\,c_n[\mathfrak{D}]\,\mathbf{u}_n,
\end{equation}
when the expansion converges in \(\mathcal{U}\) \cite{stratton1941,harrington2001}.
Physically, \(c_n\) is the excitation strength of radiative mode \(n\) before export
selection; \(\lambda_n c_n\) is the coefficient on the normalized far-zone pattern carried
by \(\mathbf{u}_n\) \cite{jackson1999}. Linearity
Eq.~\eqref{eq:ch2.radiation-linearity} implies
\(c_n[a_1\mathfrak{D}_1+a_2\mathfrak{D}_2]=a_1 c_n[\mathfrak{D}_1]+a_2 c_n[\mathfrak{D}_2]\)
for admissible superpositions after splitting each term through \(\mathcal{S}_\omega\)
\cite{stratton1941}.

\begin{figure}[t]
\centering
\ChOneFigBlock{%
\begin{tikzpicture}[ch1fig, node distance=7mm and 10mm]
  \node[ch1box] (d) {$\mathfrak{D}$};
  \node[ch1box, right=12mm of d] (s) {$\mathcal{S}_\omega$};
  \node[ch1box, right=12mm of s] (r) {$\mathcal{R}_\omega$};
  \node[ch1box, right=12mm of r] (exp) {$\sum_n \lambda_n c_n P_n$};
  \node[ch1box, below=11mm of s] (cn) {$c_n=\langle\cdot,\mathbf{u}_n\rangle$};
  \draw[ch1flow] (d) -- (s) -- (r) -- (exp);
  \draw[ch1flow] (cn) -- (s);
\end{tikzpicture}%
}
\caption{Subsection~\ref{sec:ch261} spectral pipeline: \(\mathcal{S}_\omega\) supplies
coefficients \(c_n\); \(\mathcal{R}_\omega\) maps them to export amplitudes
\(\lambda_n c_n\) on projectors \(P_n\) Eq.~\eqref{eq:ch2.modal-projector-def}.}
\label{fig:ch2.radiation-spectral-decomposition}
\end{figure}

Figure~\ref{fig:ch2.radiation-spectral-decomposition} orders the objects defined here:
integral evaluation from Theorem~\ref{thm:ch2.source-to-field-integral} produces
\(\mathcal{S}_\omega[\mathfrak{D}]\); radiation selection yields
\(\mathcal{R}_\omega[\mathfrak{D}]\); modal projectors resolve the radiative image into
spectral coordinates \cite{stratton1941,harrington2001}.

\paragraph{Spectral expansion on \(\mathcal{F}_{\mathrm{rad}}\).}
When \(\{\mathbf{u}_n\}_{n\in\mathcal{N}_{\mathrm{rad}}}\) is a complete orthonormal family
in \(\mathcal{F}_{\mathrm{rad}}\) for the declared domain class, every radiative field
\(\mathbf{v}\in\operatorname{ran}\mathcal{R}_\omega\) admits
\begin{equation}
\label{eq:ch2.radiation-spectral-expansion}
\mathbf{v}
=
\sum_{n\in\mathcal{N}_{\mathrm{rad}}}
\langle \mathbf{v},\mathbf{u}_n\rangle_{\mathcal{U}}\,\mathbf{u}_n,
\qquad
\|\mathbf{v}\|_{\mathcal{U}}^2
=
\sum_{n\in\mathcal{N}_{\mathrm{rad}}}
\big|\langle \mathbf{v},\mathbf{u}_n\rangle_{\mathcal{U}}\big|^2,
\end{equation}
with convergence in \(\mathcal{U}\) \cite{stratton1941,harrington2001}. Applying
Eq.~\eqref{eq:ch2.radiation-spectral-expansion} to \(\mathbf{v}=\mathcal{R}_\omega[\mathfrak{D}]\)
gives
\begin{equation}
\label{eq:ch2.radiation-map-modal-sum}
\mathcal{R}_\omega[\mathfrak{D}]
=
\sum_{n\in\mathcal{N}_{\mathrm{rad}}}
\lambda_n\,c_n[\mathfrak{D}]\,\mathbf{u}_n,
\end{equation}
the Subsection~\ref{sec:ch261} HOME restatement of compression operator
Eq.~\eqref{eq:ch2.spectral-compression-operator} with explicit eigenvalues and normalized
modes \cite{stratton1941}. Truncation to \(|\mathcal{N}|\le N\) yields
\(\mathcal{R}_\omega^{(N)}\) from Eq.~\eqref{eq:ch2.spectral-compression-operator}; error
Eq.~\eqref{eq:ch2.compression-error-bound} applies when the tail is compact on source space
\cite{harrington2001}.

\begin{theorem}[Spectral decomposition of radiation operators]
\label{thm:ch2.radiation-spectral-decomposition}
Let outgoing closure and domain class be fixed as above. Suppose
\(\{\mathbf{u}_n\}_{n\in\mathcal{N}_{\mathrm{rad}}}\) is a complete orthonormal radiation
eigenbasis in \(\mathcal{F}_{\mathrm{rad}}\) with
Eq.~\eqref{eq:ch2.radiation-eigenpair-def} and projectors
Eq.~\eqref{eq:ch2.modal-projector-def}. Then
\begin{equation}
\label{eq:ch2.radiation-operator-spectral-form}
\mathcal{R}_\omega
=
\sum_{n\in\mathcal{N}_{\mathrm{rad}}}
\lambda_n\,P_n\circ\mathcal{S}_\omega,
\qquad
\Pi_{\mathrm{rad}}
=
\sum_{n\in\mathcal{N}_{\mathrm{rad}}} P_n,
\end{equation}
on admissible \(\mathfrak{D}\), with convergence in \(\mathcal{U}\) for fields in
\(\operatorname{ran}\mathcal{R}_\omega\) \cite{stratton1941,harrington2001}. The range
satisfies
\begin{equation}
\label{eq:ch2.radiation-range-span}
\operatorname{ran}\mathcal{R}_\omega
=
\overline{\operatorname{span}\{\mathbf{u}_n:\ n\in\mathcal{N}_{\mathrm{rad}},\ \lambda_n\neq 0\}}.
\end{equation}
\end{theorem}
\begin{proof}
Eq.~\eqref{eq:ch2.radiation-map-modal-sum} follows by applying
Eq.~\eqref{eq:ch2.radiation-spectral-expansion} to
\(\mathcal{R}_\omega[\mathfrak{D}]\) and using
Eqs.~\eqref{eq:ch2.source-modal-amplitude}--\eqref{eq:ch2.radiation-eigenpair-def}.
Linearity of \(\mathcal{S}_\omega\) and \(\mathcal{R}_\omega\) extends the sum from finite
superpositions to all admissible data by closure in \(\mathcal{U}\) on the physical branch
\cite{stratton1941}. Eq.~\eqref{eq:ch2.radiation-operator-spectral-form} is obtained by
recognizing each term as \(P_n\mathcal{S}_\omega\) acting with gain \(\lambda_n\).
Eq.~\eqref{eq:ch2.radiation-range-span} follows because eigenvectors with
\(\lambda_n=0\) contribute no radiative export; non-radiating sources lie in
\(\ker\mathcal{R}_\omega\) Eq.~\eqref{eq:ch2.ker-R-omega-preview}
\cite{harrington2001}. Resolution Eq.~\eqref{eq:ch2.spectral-resolution-composition} matches
\(\Pi_{\mathrm{rad}}=\sum_n P_n\) on \(\operatorname{ran}\mathcal{R}_\omega\)
\cite{stratton1941}.
\end{proof}

\paragraph{Discrete cavity spectra.}
On a bounded domain with passive \(\mathcal{B}_\Gamma\) and interior resonances,
Eq.~\eqref{eq:ch2.discrete-continuous-spectrum-preview} yields a discrete accumulating set
\(\{\lambda_n\}_{n=1}^{\infty}\) when \(\mathcal{R}_\omega\) is restricted to cavity
outgoing-propagating subspaces compatible with Eq.~\eqref{eq:ch2.embedding-chain}
\cite{stratton1941,harrington2001}. The cavity specialization reads
\begin{equation}
\label{eq:ch2.cavity-discrete-radiation-spectrum}
\mathcal{R}_\omega^{(\mathrm{cav})}
=
\sum_{n=1}^{\infty}
\lambda_n\,P_n^{(\mathrm{cav})}\circ\mathcal{S}_\omega,
\qquad
\lambda_n\to 0\ \text{or}\ \lambda_n\to\infty
\ \text{according to mode class},
\end{equation}
with \(P_n^{(\mathrm{cav})}\) projectors onto normalized cavity radiation modes
\cite{stratton1941}. Equation~\eqref{eq:ch2.cavity-discrete-radiation-spectrum} is the
constructive content of Eq.~\eqref{eq:ch2.spectral-compression-operator} in closed volumes;
storage resonances with \(\lambda_n=0\) on \(\mathcal{R}_\omega\) remain in
\(\mathbf{u}_{\mathrm{nr}}\) Eq.~\eqref{eq:ch2.nonradiative-component-def}
\cite{harrington2001}. Using Eq.~\eqref{eq:ch2.cavity-discrete-radiation-spectrum} on open
exterior problems misstates continuous angular content
Eq.~\eqref{eq:ch2.discrete-continuous-spectrum-preview} \cite{stratton1941,jackson1999}.

\paragraph{Exterior angular eigenmodes.}
For compact sources in free space with outgoing closure, angular eigenmodes are labeled by
propagation direction \(\hat{\mathbf{k}}\in S^{2}\) and polarization index \(\sigma\in\{1,2\}\)
transverse to \(\hat{\mathbf{k}}\) \cite{stratton1941,jackson1999}. A normalized plane-wave
radiation mode is
\begin{equation}
\label{eq:ch2.exterior-angular-eigenmode}
\mathbf{u}_{\hat{\mathbf{k}},\sigma}(\mathbf{r})
=
\big(\E_{\hat{\mathbf{k}},\sigma}(\mathbf{r}),\,\H_{\hat{\mathbf{k}},\sigma}(\mathbf{r})\big),
\qquad
\E_{\hat{\mathbf{k}},\sigma}\sim \mathbf{e}_\sigma\,e^{-\jj k_0\hat{\mathbf{k}}\cdot\mathbf{r}}/r,
\end{equation}
with \(\mathbf{e}_\sigma\cdot\hat{\mathbf{k}}=0\) and outgoing phase
\cite{stratton1941,harrington2001}. The continuous spectral measure \(d\Omega(\hat{\mathbf{k}})\)
replaces discrete summation in Eq.~\eqref{eq:ch2.radiation-spectral-expansion}:
\begin{equation}
\label{eq:ch2.exterior-continuous-spectral-integral}
\mathcal{R}_\omega[\mathfrak{D}](\mathbf{r})
=
\int_{S^{2}}
\sum_{\sigma=1}^{2}
a_\sigma(\hat{\mathbf{k}})\,
\mathbf{u}_{\hat{\mathbf{k}},\sigma}(\mathbf{r})\,d\Omega(\hat{\mathbf{k}}),
\end{equation}
with amplitudes \(a_\sigma\) determined by \(\mathcal{S}_\omega[\mathfrak{D}]\) through
Eq.~\eqref{eq:ch2.source-modal-amplitude} in the rigged limit \cite{stratton1941}. Rigged
Hilbert space completion and measure normalization are HOME in Chapter~\ref{ch:03}
Subsection~\ref{sec:ch313}; Subsection~\ref{sec:ch261} records the operator target
Eq.~\eqref{eq:ch2.exterior-continuous-spectral-integral} compatible with far-zone reduction
\cite{harrington2001}.

\begin{proposition}[Far-pattern coefficients from spectral modes]
\label{prop:ch2.far-pattern-from-modes}
Let \(\mathcal{D}_{J,\mathrm{reg}}=1\) and let \(\mathbf{u}_n\) be radiation eigenmodes with
far amplitudes \(\mathbf{f}_n(\theta,\phi)\) defined through
Eq.~\eqref{eq:ch2.far-field-amplitude-integral}. Then the far pattern of
\(\mathcal{R}_\omega[\mathfrak{D}]\) satisfies
\begin{equation}
\label{eq:ch2.spectral-far-pattern-expansion}
\E_0(\theta,\phi)
=
\sum_{n\in\mathcal{N}_{\mathrm{rad}}}
\lambda_n\,c_n[\mathfrak{D}]\,\mathbf{f}_n(\theta,\phi),
\end{equation}
with convergence in mean-square on \(S^{2}\) when the modal sum converges in
\(\mathcal{U}\) \cite{stratton1941,jackson1999,harrington2001}. Equation
\eqref{eq:ch2.spectral-far-pattern-expansion} is the Subsection~\ref{sec:ch261} HOME upgrade
of preview Eq.~\eqref{eq:ch2.far-pattern-modal-preview}; explicit \(\mathbf{f}_n\) in vector
harmonic form belong to Chapter~\ref{ch:03} \cite{stratton1941}.
\end{proposition}
\begin{proof}
Apply Theorem~\ref{thm:ch2.far-field-approximation} to each modal component
\(\lambda_n c_n\,\mathbf{u}_n\) and superpose using linearity
Eq.~\eqref{eq:ch2.radiation-linearity} on the radiative branch. Coefficients match
Eq.~\eqref{eq:ch2.source-modal-amplitude} because far-zone reduction depends linearly on
\(\mathcal{R}_\omega[\mathfrak{D}]\) through
Eqs.~\eqref{eq:ch2.far-field-integral-reduction}--\eqref{eq:ch2.far-field-amplitude-integral}
\cite{stratton1941,jackson1999}. Regime gate
Eq.~\eqref{eq:ch2.regime-component-gate} must hold before interpreting
Eq.~\eqref{eq:ch2.spectral-far-pattern-expansion} as transport data
\cite{harrington2001}.
\end{proof}

\paragraph{Adjoint modes and biorthogonality.}
When \(\mathcal{R}_\omega\) is not self-adjoint on \(\mathcal{J}_\Omega\), define adjoint
modes \(\widetilde{\mathbf{u}}_n\) through the source pairing
Eq.~\eqref{eq:ch2.dual-source-pairing} such that
\begin{equation}
\label{eq:ch2.radiation-biorthogonality}
\langle \widetilde{\mathbf{u}}_m,\,\mathcal{R}_\omega[\Jimp^{(n)}]\rangle_{\mathcal{U}}
=
\lambda_n\,\delta_{mn},
\end{equation}
for normalized radiation eigenpairs \cite{stratton1941,harrington2001}. Equation
\eqref{eq:ch2.radiation-biorthogonality} is the field-side restatement of solvability preview
Eq.~\eqref{eq:ch2.spectral-solvability-preview}: cokernel directions annihilated by
Eq.~\eqref{eq:ch2.radiation-biorthogonality} obstruct inversion
\(\mathcal{R}_\omega[\Jimp]=\mathbf{g}\) \cite{stratton1941}. Complete adjoint-mode catalogs
on exterior domains belong to Subsection~\ref{sec:ch263} \cite{harrington2001}.

\paragraph{Null modes and export filter.}
Modes with \(\lambda_n=0\) satisfy \(P_n\mathcal{R}_\omega=\mathbf{0}\) but may have
\(P_n\mathcal{S}_\omega\neq\mathbf{0}\) when reactive content projects onto non-export
subspaces Eq.~\eqref{eq:ch2.nonpropagating-component-def}
\cite{stratton1941,harrington2001}. Evanescent and non-propagating sheets remain annihilated:
\begin{equation}
\label{eq:ch2.modal-null-on-nonprop}
P_n\circ\Pi_{\mathrm{ev}}=P_n\circ\Pi_{\mathrm{non}}=\mathbf{0}
\quad\text{for all}\ n\in\mathcal{N}_{\mathrm{rad}},
\end{equation}
by Eq.~\eqref{eq:ch2.rad-null-on-evanescent} and modal definition
Eq.~\eqref{eq:ch2.modal-projector-def} \cite{stratton1941}. Spectral labels therefore do not
override regime qualification Eq.~\eqref{eq:ch2.regime-operator-screen}
\cite{harrington2001}.

\paragraph{Boundary of validity.}
Theorem~\ref{thm:ch2.radiation-spectral-decomposition} and
Proposition~\ref{prop:ch2.far-pattern-from-modes} require linear media at fixed \(\omega\),
outgoing closure, and domain-class-appropriate completeness of \(\{\mathbf{u}_n\}\)
\cite{stratton1941,harrington2001}. Loss shifts \(\lambda_n\) into the complex plane without
promoting evanescent storage to export through Eq.~\eqref{eq:ch2.modal-null-on-nonprop}
\cite{stratton1941}. Finite observation rank
Eq.~\eqref{eq:ch2.reconstruction-rank-cap} limits recoverable coefficients in
Eq.~\eqref{eq:ch2.spectral-far-pattern-expansion} independently of modal completeness
\cite{harrington2001}.

\paragraph{Worked illustration (dipole-dominated spectrum).}
For \(k_0 a\ll 1\) with \(\Jimp\) supported in \(B_a\), the leading radiation eigenmode has
\(\lambda_1\neq 0\) with \(\mathbf{f}_1\) dipolar and higher-\(n\) coefficients suppressed by
\((k_0 a)^{2\ell}\) in harmonic bookkeeping deferred to Chapter~\ref{ch:03}
\cite{jackson1999,stratton1941}. Truncating
Eq.~\eqref{eq:ch2.spectral-far-pattern-expansion} to \(n=1\) is accurate in the far zone when
\(\mathcal{D}_{J,\mathrm{reg}}=1\) \cite{harrington2001}. Near-zone scans of
\(\mathcal{S}_\omega[\Jimp]\) still contain large \(\mathbf{u}_{\mathrm{nr}}\) not captured by
the single-mode truncation \cite{stratton1941}.

\paragraph{Worked illustration (cavity modal truncation).}
In a closed cavity, Eq.~\eqref{eq:ch2.cavity-discrete-radiation-spectrum} truncates at
\(N\) with error Eq.~\eqref{eq:ch2.compression-error-bound} when source classes are bounded
\cite{stratton1941,harrington2001}. Applying the same truncated sum to an exterior dipole
misaligns sidelobes because continuous measure
Eq.~\eqref{eq:ch2.exterior-continuous-spectral-integral} is required
\cite{jackson1999}. Subsection~\ref{sec:ch261} fixes the operator book; numerical aperture
codes must match domain class \cite{stratton1941}.

\paragraph{Worked illustration (null eigenvalue).}
A non-radiating source \(\Jimp^{(\mathrm{NR})}\in\ker\mathcal{R}_\omega\) has
\(c_n=0\) for all \(n\) with \(\lambda_n\neq 0\) yet may excite reactive modes in
\(\Pi_{\mathrm{ev}}\mathcal{S}_\omega[\Jimp^{(\mathrm{NR})}]\)
\cite{stratton1941,harrington2001}. Far-zone pattern meters on
\(\mathcal{S}_\omega[\Jimp^{(\mathrm{NR})}]\) without \(\Pi_{\mathrm{rad}}\) cannot detect
the null-space membership; Eq.~\eqref{eq:ch2.ker-R-omega-preview} is spectral, not a near-field
zero \cite{stratton1941}. Complete \(\ker\mathcal{R}_\omega\) classification belongs to
Section~\ref{sec:ch272} \cite{harrington2001}.

\paragraph{Closing summary.}
Subsection~\ref{sec:ch261} establishes the Chapter~\ref{ch:02} HOME for radiation eigenpairs
Eq.~\eqref{eq:ch2.radiation-eigenpair-def}, modal projectors
Eq.~\eqref{eq:ch2.modal-projector-def}, source amplitudes
Eq.~\eqref{eq:ch2.source-modal-amplitude}, spectral expansions
Eqs.~\eqref{eq:ch2.radiation-spectral-expansion}--\eqref{eq:ch2.radiation-map-modal-sum},
operator spectral form Theorem~\ref{thm:ch2.radiation-spectral-decomposition}, cavity and
exterior spectra Eqs.~\eqref{eq:ch2.cavity-discrete-radiation-spectrum}--\eqref{eq:ch2.exterior-continuous-spectral-integral},
far-pattern expansion Proposition~\ref{prop:ch2.far-pattern-from-modes}, and biorthogonality
Eq.~\eqref{eq:ch2.radiation-biorthogonality} \cite{stratton1941,harrington2001,jackson1999}.
Subsection~\ref{sec:ch262} imports Section~\ref{sec:ch01.3} to classify degeneracy; Subsection
~\ref{sec:ch263} develops Fredholm consequences of the same modal coordinates
\cite{stratton1941}.
\subsection{Rotational symmetry and spectral degeneracy}
\label{sec:ch262}
% TOC tag: [USE SS1.3]

Subsection~\ref{sec:ch261} fixed radiation eigenpairs, modal projectors, and spectral
expansions Eqs.~\eqref{eq:ch2.radiation-eigenpair-def}--\eqref{eq:ch2.radiation-operator-spectral-form}.
Subsection~\ref{sec:ch262} classifies how those modes transform under rotations and why
degeneracy appears in radiative spectra, importing active and passive actions
Eqs.~\eqref{eq:ch1.active-rotation-action}--\eqref{eq:ch1.active-passive-duality},
transport invariances Eqs.~\eqref{eq:ch1.rotation-power-invariance}--\eqref{eq:ch1.rotation-am-invariance},
scalar and vector observable rules
Eqs.~\eqref{eq:ch1.scalar-observable-invariance}--\eqref{eq:ch1.vector-observable-covariance},
and origin-shift laws from Section~\ref{sec:ch01.3} by reference
\cite{stratton1941,harrington2001,jackson1999}. Rotation group structure, tensor
transformation rules, and invariant filters are not re-derived; the present subsection
applies Section~\ref{sec:ch01.3} to \(\mathcal{R}_\omega\) and to coefficients in
Theorem~\ref{thm:ch2.radiation-spectral-decomposition}
\cite{stratton1941}. Explicit \(\mathrm{SO}(3)\) representation matrices on harmonic
bases belong to Chapter~\ref{ch:03}; here degeneracy multiplicities are recorded at the
operator level compatible with preview Eq.~\eqref{eq:ch2.degeneracy-covariance-preview}
\cite{harrington2001}. Fredholm cokernel structure under rotation is deferred to
Subsection~\ref{sec:ch263}.

\paragraph{Assumptions.}
Let \(R\in\mathrm{SO}(3)\), field rotation \(\mathcal{D}(R)\) on admissible pairs
\(\mathbf{u}=(\E,\H)\in\mathcal{F}_{\mathrm{SF}}\), and radiation map
Eq.~\eqref{eq:ch2.spectral-operator-targets} on outgoing data
\cite{stratton1941,harrington2001}. Isotropic vacuum or statistically isotropic
environments are assumed unless stated otherwise: constitutive tensors commute with
\(R\) in the sense of Section~\ref{sec:ch01.3.2} \cite{stratton1941}. Impressed sources
transform as polar vectors under active rotation,
\(\big(\mathcal{A}_R\Jimp\big)(\mathbf{r})=R\,\Jimp(R^{-1}\mathbf{r})\), consistent with
Eq.~\eqref{eq:ch1.vector-rotation-components} \cite{stratton1941}. Regime gate
Eq.~\eqref{eq:ch2.regime-operator-screen} and transport qualification
Eq.~\eqref{eq:ch2.transport-qualified-radiation} remain mandatory before interpreting
rotated reports as export data \cite{harrington2001}.

\paragraph{Rotation action on field pairs.}
Define the Maxwell field representation
\begin{equation}
\label{eq:ch2.field-rotation-representation}
\mathcal{D}(R)\,\mathbf{u}
=
\big(\mathcal{A}_R\E,\,\mathcal{A}_R\H\big),
\qquad
\mathbf{u}=(\E,\H)\in\mathcal{F}_\Omega,
\end{equation}
using active rotation Eq.~\eqref{eq:ch1.active-rotation-action} on each component
\cite{stratton1941}. Equation~\eqref{eq:ch2.field-rotation-representation} is the
Subsection~\ref{sec:ch262} specialization of Section~\ref{sec:ch01.3.1} to phasor Maxwell
pairs: passive relabeling Eq.~\eqref{eq:ch1.passive-rotation-action} produces the same
physics with rotated components related by
Eq.~\eqref{eq:ch1.active-passive-duality} \cite{stratton1941}. Group composition
Eq.~\eqref{eq:ch1.rotation-group-action} yields
\(\mathcal{D}(R_1R_2)=\mathcal{D}(R_1)\mathcal{D}(R_2)\) on admissible fields
\cite{stratton1941,harrington2001}.

\paragraph{Covariance of the radiation map.}
When sources and observation frames are co-rotated, the radiation pipeline satisfies
\begin{equation}
\label{eq:ch2.radiation-rotation-commutation}
\mathcal{R}_\omega\big[\mathcal{A}_R\Jimp\big]
=
\mathcal{D}(R)\,\mathcal{R}_\omega[\Jimp],
\end{equation}
on admissible \(\Jimp\in\mathcal{J}_\Omega\) with outgoing closure and isotropic
constitutive response \cite{stratton1941,harrington2001,jackson1999}. Equation
\eqref{eq:ch2.radiation-rotation-commutation} upgrades preview
Eq.~\eqref{eq:ch2.degeneracy-covariance-preview} and aligns with covariance test
Eq.~\eqref{eq:ch2.covariance-operator-test} for \(\mathbf{a}=\mathbf{0}\)
\cite{stratton1941}. Residual \(\mathcal{R}_{\mathcal{F}}\) from
Eq.~\eqref{eq:ch1.covariance-residual} diagnoses broken commutation in discretizations;
Subsection~\ref{sec:ch262} records the continuous target \cite{harrington2001}. Rotating
\(\Jimp\) while evaluating flux on an unrotated \(\mathcal{S}_\omega[\Jimp]\) violates
origin and covariance screens from Section~\ref{sec:ch01.3.3} even when
Eq.~\eqref{eq:ch2.radiation-rotation-commutation} holds for power alone
\cite{stratton1941}.

\begin{figure}[t]
\centering
\ChOneFigBlock{%
\begin{tikzpicture}[ch1fig, node distance=7mm and 10mm]
  \node[ch1box] (j) {$\Jimp$};
  \node[ch1box, right=12mm of j] (r) {$\mathcal{R}_\omega$};
  \node[ch1box, right=12mm of r] (u) {$\mathbf{u}_{\mathrm{rad}}$};
  \node[ch1box, below=12mm of j] (ar) {$\mathcal{A}_R$};
  \node[ch1box, below=12mm of u] (dr) {$\mathcal{D}(R)$};
  \draw[ch1flow] (j) -- (r) -- (u);
  \draw[ch1flow] (ar) -- (j);
  \draw[ch1flow] (dr) -- (u);
  \draw[ch1flow] (ar) -| (r);
\end{tikzpicture}%
}
\caption{Rotational covariance Eq.~\eqref{eq:ch2.radiation-rotation-commutation}: active
source rotation commutes with \(\mathcal{R}_\omega\) up to field representation
\(\mathcal{D}(R)\) Eq.~\eqref{eq:ch2.field-rotation-representation}.}
\label{fig:ch2.rotation-spectral-covariance}
\end{figure}

Figure~\ref{fig:ch2.rotation-spectral-covariance} restates
Eq.~\eqref{eq:ch2.radiation-rotation-commutation} in the same layout as
Figure~\ref{fig:ch1.active-passive-rotation}: geometry and components must be rotated
consistently before modal coefficients are compared \cite{stratton1941,harrington2001}.

\paragraph{Modal coefficients under rotation.}
Let \(\{\mathbf{u}_n\}_{n\in\mathcal{N}_{\mathrm{rad}}}\) be the orthonormal radiation
basis from Subsection~\ref{sec:ch261} and \(c_n[\mathfrak{D}]\) source amplitudes
Eq.~\eqref{eq:ch2.source-modal-amplitude}. Under active rotation of impressed data,
\begin{equation}
\label{eq:ch2.modal-coefficient-rotation}
c_n[\mathcal{A}_R\mathfrak{D}]
=
\sum_{m\in\mathcal{N}_{\mathrm{rad}}}
D^{(\mathrm{rad})}_{nm}(R)\,c_m[\mathfrak{D}],
\end{equation}
whenever Eq.~\eqref{eq:ch2.radiation-rotation-commutation} holds and the modal family is
closed under \(\mathcal{D}(R)\),
\begin{equation}
\label{eq:ch2.modal-basis-covariance}
\mathcal{D}(R)\,\mathbf{u}_m
=
\sum_{n\in\mathcal{N}_{\mathrm{rad}}}
D^{(\mathrm{rad})}_{nm}(R)\,\mathbf{u}_n.
\end{equation}
Equations~\eqref{eq:ch2.modal-coefficient-rotation}--\eqref{eq:ch2.modal-basis-covariance}
are the spectral restatement of vector covariance
Eq.~\eqref{eq:ch1.vector-observable-covariance} on \(\mathcal{F}_{\mathrm{rad}}\)
\cite{stratton1941}. Matrix elements \(D^{(\mathrm{rad})}_{nm}(R)\) are not tabulated
here; Chapter~\ref{ch:03} supplies them on vector spherical harmonic bases
\cite{harrington2001}. Eigenvalues \(\lambda_n\) in
Eq.~\eqref{eq:ch2.radiation-eigenpair-def} are unchanged within a degenerate block because
\(\mathcal{R}_\omega\) commutes with \(\mathcal{D}(R)\) on that block
\cite{stratton1941,jackson1999}.

\begin{theorem}[Rotational degeneracy of radiation spectra]
\label{thm:ch2.rotation-spectral-degeneracy}
Suppose Eq.~\eqref{eq:ch2.radiation-rotation-commutation} holds for all
\(R\in\mathrm{SO}(3)\) in a compact subgroup \(G\subseteq\mathrm{SO}(3)\) leaving the
problem geometry invariant. Then \(\mathcal{R}_\omega\) and \(\mathcal{D}(R)\) generate a
commuting family on \(\mathcal{F}_{\mathrm{rad}}\): for each \(R\in G\),
\begin{equation}
\label{eq:ch2.radiation-D-commutation}
\mathcal{D}(R)\,\mathcal{R}_\omega
=
\mathcal{R}_\omega\,\mathcal{D}(R)
\quad\text{on admissible data,}
\end{equation}
and eigenspaces of \(\mathcal{R}_\omega\mathcal{S}_\omega\) decompose into irreducible
\(G\)-representations. If an irreducible block carries angular index \(\ell\), its
degeneracy multiplicity satisfies
\begin{equation}
\label{eq:ch2.spectral-degeneracy-multiplicity}
\dim\mathcal{H}_{\ell}
=
2\ell+1
\quad\text{for full \(\mathrm{SO}(3)\) closure on }S^{2},
\end{equation}
with \(\mathcal{H}_{\ell}\) the \(\ell\)-indexed radiative subspace
\cite{stratton1941,jackson1999,harrington2001}. Equation
\eqref{eq:ch2.spectral-degeneracy-multiplicity} explains repeated eigenvalues in
Eq.~\eqref{eq:ch2.radiation-spectral-expansion} without duplicating mode physics.
\end{theorem}
\begin{proof}
Eq.~\eqref{eq:ch2.radiation-D-commutation} follows by applying
Eq.~\eqref{eq:ch2.radiation-rotation-commutation} after
Eq.~\eqref{eq:ch2.radiation-operator-spectral-form} and using linearity of
\(\mathcal{S}_\omega\) \cite{stratton1941}. Standard spectral theory for commuting
unitary or orthogonal representations implies simultaneous block diagonalization of
\(\mathcal{R}_\omega\) and \(\mathcal{D}(R)\) on invariant subspaces
\cite{stratton1941,harrington2001}. For full rotational symmetry on angular patterns,
irreducible representations of \(\mathrm{SO}(3)\) on \(L^{2}(S^{2})\) have dimension
\(2\ell+1\); Eq.~\eqref{eq:ch2.spectral-degeneracy-multiplicity} records that count on
\(\mathcal{F}_{\mathrm{rad}}\) once outgoing transversality and
Eq.~\eqref{eq:ch2.rad-null-on-evanescent} are imposed \cite{jackson1999}. Detailed
construction of \(\mathcal{H}_\ell\) is HOME in Chapter~\ref{ch:03}
\cite{harrington2001}.
\end{proof}

\paragraph{Invariant power and flux reports.}
Rigid rotation cannot change total exported power when meters act on
\(\mathcal{R}_\omega[\Jimp]\) with co-rotated frames. Using
Eq.~\eqref{eq:ch2.power-flux-functional} and
Eq.~\eqref{eq:ch1.rotation-power-invariance},
\begin{equation}
\label{eq:ch2.spectral-power-rotation-invariance}
P_{\mathrm{rad}}[\mathcal{A}_R\mathfrak{D}]
=
P_{\mathrm{rad}}[\mathfrak{D}],
\qquad
P_{\mathrm{rad}}[\mathfrak{D}]
=
\sum_{n\in\mathcal{N}_{\mathrm{rad}}}
|\lambda_n|^{2}\,|c_n[\mathfrak{D}]|^{2}
\quad\text{when modes are orthogonal and normalized,}
\end{equation}
for lossless outgoing radiation \cite{stratton1941,poynting1884,harrington2001}. Equation
\eqref{eq:ch2.spectral-power-rotation-invariance} is the modal restatement of scalar
invariance Eq.~\eqref{eq:ch1.scalar-observable-invariance} on integrated power: rotated
coefficients redistribute directional content but preserve the scalar budget when
\(\mathcal{D}_{J,\mathrm{reg}}=1\) \cite{stratton1941}. Net angular-flux magnitude obeys
Eq.~\eqref{eq:ch1.rotation-am-invariance} on the same qualified subspace
\cite{stratton1941,harrington2001}.

\paragraph{Active versus passive reports on spectra.}
Active rotation reorients sources and patterns in space; passive rotation re-expresses
fixed physics in a rotated chart Eqs.~\eqref{eq:ch1.active-rotation-action}--\eqref{eq:ch1.passive-rotation-action}
\cite{stratton1941}. Modal coefficients transform as
Eq.~\eqref{eq:ch2.modal-coefficient-rotation} under active action, while passive
relabeling conjugates the same \(D^{(\mathrm{rad})}_{nm}(R)\) through
Eq.~\eqref{eq:ch1.active-passive-duality} without altering
Eq.~\eqref{eq:ch2.spectral-power-rotation-invariance} \cite{stratton1941,harrington2001}.
Confusing the two conventions produces spurious degeneracy splitting in measured
\(a_n(\hat{\mathbf{k}})\) from Eq.~\eqref{eq:ch2.spectral-far-pattern-expansion} even when
Theorem~\ref{thm:ch2.rotation-spectral-degeneracy} applies \cite{stratton1941}. Laboratory
rotation of an antenna rotates measured coefficients through
Eq.~\eqref{eq:ch2.modal-coefficient-rotation}; exported power on
\(\mathcal{R}_\omega[\Jimp]\) remains invariant in the far zone
\cite{harrington2001,jackson1999}.

\paragraph{Origin shifts and modal labels.}
Origin translation \(\mathbf{a}\) shifts torque and angular-flux reports through
Section~\ref{sec:ch01.3.3} while leaving far-zone radiative power invariant when
\(\mathcal{D}_{J,\mathrm{reg}}=1\) \cite{stratton1941}. At the spectral level, origin
change mixes modal coefficients within fixed \(\ell\) blocks through addition of
lower-\(\ell\) content from shifted multipole expansions; it does not create new
\(\lambda_n\neq 0\) export unless sources change \cite{jackson1999,harrington2001}.
Subsection~\ref{sec:ch262} therefore separates \emph{rotational} degeneracy
Theorem~\ref{thm:ch2.rotation-spectral-degeneracy} from \emph{translational} bookkeeping
Section~\ref{sec:ch01.3.3}; intrinsic angular observables use
Eq.~\eqref{eq:ch1.intrinsic-angular-observable} by reference \cite{stratton1941}.

\begin{proposition}[Degeneracy lifting under broken symmetry]
\label{prop:ch2.degeneracy-splitting-broken-symmetry}
Let \(G_0\subset\mathrm{SO}(3)\) be the exact symmetry of a source--boundary configuration
and let \(\mathcal{H}_{\ell}\) be an irreducible block from
Theorem~\ref{thm:ch2.rotation-spectral-degeneracy}. Any symmetry-breaking perturbation
(reduced \(G_0\), anisotropic material, or finite facetization) splits
Eq.~\eqref{eq:ch2.spectral-degeneracy-multiplicity} into smaller blocks with distinct
\(\lambda_n\) within the same \(\ell\) sector \cite{stratton1941,harrington2001}. The
splitting order is controlled by representation coupling rules imported from
Section~\ref{sec:ch01.3.2} on tensor perturbations; explicit Clebsch--Gordan content belongs
to Chapter~\ref{ch:03} Subsection~\ref{sec:ch325} \cite{stratton1941}.
\end{proposition}
\begin{proof}
When \(G_0\) is proper, \(\mathcal{D}(R)\) block-diagonalizes with fixed \(\lambda_n\) on
each irreducible piece. Perturbations that fail to commute with all \(R\in G_0\) remove
simultaneous eigenstructure; standard degenerate perturbation theory applies to the
restricted radiation operator on \(\mathcal{H}_\ell\) \cite{stratton1941,harrington2001}.
\end{proof}

\paragraph{Boundary of validity.}
Theorem~\ref{thm:ch2.rotation-spectral-degeneracy} requires isotropic or symmetry-matched
media, outgoing closure, and modal families closed under \(\mathcal{D}(R)\)
\cite{stratton1941,harrington2001}. Anisotropic crystals reduce \(G\) and split
Eq.~\eqref{eq:ch2.spectral-degeneracy-multiplicity} per
Proposition~\ref{prop:ch2.degeneracy-splitting-broken-symmetry} \cite{stratton1941}.
Finite apertures and rank-\(N\) observation maps
Eq.~\eqref{eq:ch2.reconstruction-rank-cap} resolve only a finite subset of each degenerate
block even when continuous symmetry is exact \cite{harrington2001}. Regime gate
Eq.~\eqref{eq:ch2.regime-component-gate} still governs whether rotated near fields may be
reported as transport \cite{stratton1941}.

\paragraph{Worked illustration (rotated dipole multipole block).}
A short dipole excites primarily \(\ell=1\) with
Eq.~\eqref{eq:ch2.spectral-degeneracy-multiplicity} giving a three-dimensional block when
full \(\mathrm{SO}(3)\) symmetry applies \cite{jackson1999,stratton1941}. Rotating the
source actively via \(\mathcal{A}_R\) rotates \(\mathbf{f}_1(\theta,\phi)\) in
Eq.~\eqref{eq:ch2.spectral-far-pattern-expansion} while leaving
Eq.~\eqref{eq:ch2.spectral-power-rotation-invariance} unchanged
\cite{harrington2001}. Reporting only one component without
Eq.~\eqref{eq:ch2.modal-coefficient-rotation} misstates degeneracy as apparent pattern
distortion \cite{stratton1941}.

\paragraph{Worked illustration (symmetry breaking by ground plane).}
A dipole above a PEC plane reduces \(G_0\) from \(\mathrm{SO}(3)\) to a reflection group;
the \(\ell=1\) block splits into distinct even/odd combinations with different
\(\lambda_n\) in Eq.~\eqref{eq:ch2.radiation-eigenpair-def}
\cite{stratton1941,harrington2001}. Proposition~\ref{prop:ch2.degeneracy-splitting-broken-symmetry}
predicts lifted degeneracy without recomputing Green kernels from Section~\ref{sec:ch23}
\cite{stratton1941}. Geometric constraints on allowed sectors are developed further in
Subsection~\ref{sec:ch264} \cite{harrington2001}.

\paragraph{Worked illustration (passive relabeling in pattern cuts).}
Passive rotation Eq.~\eqref{eq:ch1.passive-rotation-action} relabels
\(\E_0(\theta,\phi)\) on \(S^{2}\) while preserving
Eq.~\eqref{eq:ch2.spectral-power-rotation-invariance} \cite{stratton1941}. Comparing
pattern cuts taken in different chart orientations is valid only when the same passive
convention is applied to both data sets through
Eqs.~\eqref{eq:ch1.active-passive-duality} and~\eqref{eq:ch2.modal-coefficient-rotation}
\cite{harrington2001}. Active rotation of the source without rotating the measurement
chart violates Eq.~\eqref{eq:ch2.covariance-operator-test} and inflates
Eq.~\eqref{eq:ch1.covariance-residual} \cite{stratton1941}.

\paragraph{Closing summary.}
Subsection~\ref{sec:ch262} applies Section~\ref{sec:ch01.3} to radiation spectra through
field representation Eq.~\eqref{eq:ch2.field-rotation-representation}, commutation
Eq.~\eqref{eq:ch2.radiation-rotation-commutation}, modal transformation laws
Eqs.~\eqref{eq:ch2.modal-coefficient-rotation}--\eqref{eq:ch2.modal-basis-covariance},
degeneracy Theorem~\ref{thm:ch2.rotation-spectral-degeneracy} with multiplicity
Eq.~\eqref{eq:ch2.spectral-degeneracy-multiplicity}, power invariance
Eq.~\eqref{eq:ch2.spectral-power-rotation-invariance}, and symmetry-breaking split
Proposition~\ref{prop:ch2.degeneracy-splitting-broken-symmetry}, importing Ch.~1 rotation
and invariant rules by reference \cite{stratton1941,harrington2001,jackson1999}.
Subsection~\ref{sec:ch263} develops Fredholm solvability on the same modal coordinates;
Subsection~\ref{sec:ch264} records geometric constraints beyond rotational symmetry
\cite{stratton1941}.
\subsection{Fredholm structure and solvability}
\label{sec:ch263}
% TOC tag: [USE SS2.1.3][HOME]

Subsections~\ref{sec:ch261}--\ref{sec:ch262} fixed radiation eigenmodes, spectral
expansions, and rotational degeneracy on \(\mathcal{F}_{\mathrm{rad}}\). Subsection
~\ref{sec:ch263} gives the Chapter~\ref{ch:02} HOME for Fredholm structure, index,
and solvability of radiation-related inversions
\(\mathcal{R}_\omega[\Jimp]=\mathbf{g}\) and finite-rank observations
\(\mathcal{O}_{\mathrm{meas}}\mathcal{R}_\omega[\Jimp]=\mathbf{b}\), importing
compactness Eq.~\eqref{eq:ch2.compact-operator-def}, Green split
Eq.~\eqref{eq:ch2.green-compact-split}, Fredholm split preview
Eq.~\eqref{eq:ch2.fredholm-split}, finite-rank observation
Eqs.~\eqref{eq:ch2.finite-rank-observation}--\eqref{eq:ch2.observation-null-space},
singular-value decay Eq.~\eqref{eq:ch2.singular-value-decay}, and rank composition
Eq.~\eqref{eq:ch2.rank-composition-law} from Subsection~\ref{sec:ch213} by reference
\cite{stratton1941,harrington2001}. Embedding lemmas and compression error bounds are not
re-proved; the present subsection applies them to \(\mathcal{R}_\omega\) on
\(\mathcal{J}_\Omega\) and upgrades solvability preview
Eq.~\eqref{eq:ch2.spectral-solvability-preview} to constructive criteria compatible with
inverse source Eq.~\eqref{eq:ch2.inverse-source-equation} and null-space preview
Eq.~\eqref{eq:ch2.ker-R-omega-preview} \cite{stratton1941}. Complete non-radiating source
classification remains in Section~\ref{sec:ch272}; geometric sector constraints belong to
Subsection~\ref{sec:ch264} \cite{harrington2001}.

\paragraph{Assumptions.}
Fix \(\omega>0\), bounded or exterior domain class with outgoing closure on open problems,
admissible \(\mathcal{J}_\Omega\) Eq.~\eqref{eq:ch2.source-space-refined}, radiation map
Eq.~\eqref{eq:ch2.spectral-operator-targets}, and spectral decomposition
Theorem~\ref{thm:ch2.radiation-spectral-decomposition} on the physical branch
\cite{stratton1941,harrington2001}. Source pairing uses dual space
Eq.~\eqref{eq:ch2.dual-source-pairing}; field pairing uses
Eq.~\eqref{eq:ch2.maxwell-adjoint-pairing} \cite{stratton1941}. Compact perturbations are
taken in operator norm on \(\mathcal{J}_\Omega\) for sources with compact support
\(\Omega_s\) \cite{harrington2001}. Regime gate Eq.~\eqref{eq:ch2.regime-operator-screen}
governs whether recovered \(\mathbf{g}\) may be reported as transport data
\cite{stratton1941}.

\paragraph{Fredholm decomposition of the radiation map.}
Green split Eq.~\eqref{eq:ch2.green-compact-split} induces a parallel split on export:
\begin{equation}
\label{eq:ch2.radiation-fredholm-decomposition}
\mathcal{R}_\omega
=
\mathcal{R}_{0,\omega}+K_{\mathrm{rad}},
\qquad
K_{\mathrm{rad}}\ \text{compact on}\ \mathcal{J}_\Omega,
\end{equation}
where \(\mathcal{R}_{0,\omega}\) captures leading outgoing singular propagation and
\(K_{\mathrm{rad}}\) collects compact smoothing at finite separation on bounded source classes
\cite{stratton1941,harrington2001}. Equation~\eqref{eq:ch2.radiation-fredholm-decomposition}
specializes Eq.~\eqref{eq:ch2.fredholm-split} to the radiation image
Eq.~\eqref{eq:ch2.integral-radiation-map} after Sections~\ref{sec:ch23}--\ref{sec:ch25};
kernel construction is not repeated \cite{stratton1941}. On exterior domains with compact
\(\Omega_s\), \(K_{\mathrm{rad}}\) is compact while \(\mathcal{R}_{0,\omega}\) carries the
non-compact far-zone propagation; on bounded cavities, both parts reduce to discrete modal
books Eq.~\eqref{eq:ch2.cavity-discrete-radiation-spectrum} \cite{harrington2001,jackson1999}.

\paragraph{Fredholm index.}
For closed operators with finite-dimensional kernel and cokernel, define the Fredholm index
\begin{equation}
\label{eq:ch2.radiation-fredholm-index-def}
\mathrm{ind}\,\mathcal{R}_\omega
=
\dim\ker\mathcal{R}_\omega
-
\dim\mathrm{coker}\,\mathcal{R}_\omega,
\end{equation}
with \(\mathrm{coker}\,\mathcal{R}_\omega=\mathcal{F}_{\mathrm{rad}}/\operatorname{ran}\mathcal{R}_\omega\)
on the declared quotient by outgoing closure \cite{stratton1941,harrington2001}. Equation
\eqref{eq:ch2.radiation-fredholm-index-def} counts independent non-radiating source degrees
Eq.~\eqref{eq:ch2.ker-R-omega-preview} minus obstructions to reaching a target far pattern
\(\mathbf{g}\in\mathcal{F}_{\mathrm{rad}}\) \cite{stratton1941}. Index is stable under compact
perturbations in Eq.~\eqref{eq:ch2.radiation-fredholm-decomposition}
\cite{harrington2001}. Observation rank cap Eq.~\eqref{eq:ch2.reconstruction-rank-cap}
is independent of \(\mathrm{ind}\,\mathcal{R}_\omega\) but composes with it in inversion
budgets \cite{stratton1941}.

\begin{figure}[t]
\centering
\ChOneFigBlock{%
\begin{tikzpicture}[ch1fig, node distance=7mm and 10mm]
  \node[ch1box] (j) {$\mathcal{J}_\Omega$};
  \node[ch1box, right=12mm of j] (r) {$\mathcal{R}_\omega$};
  \node[ch1box, right=12mm of r] (f) {$\mathcal{F}_{\mathrm{rad}}$};
  \node[ch1box, below=11mm of r] (k) {$K_{\mathrm{rad}}$};
  \node[ch1box, below=11mm of f] (c) {$\mathrm{coker}$};
  \draw[ch1flow] (j) -- (r) -- (f);
  \draw[ch1flow] (k) -- (r);
  \draw[ch1flow] (c) -- (f);
\end{tikzpicture}%
}
\caption{Subsection~\ref{sec:ch263} Fredholm pipeline: compact perturbation
Eq.~\eqref{eq:ch2.radiation-fredholm-decomposition}, kernel
\(\ker\mathcal{R}_\omega\), and cokernel obstructions to solvability
Eq.~\eqref{eq:ch2.radiation-solvability-criterion}.}
\label{fig:ch2.fredholm-radiation-pipeline}
\end{figure}

Figure~\ref{fig:ch2.fredholm-radiation-pipeline} separates rank deficits at the source
from rank deficits at the observation layer
Eq.~\eqref{eq:ch2.rank-composition-law} \cite{stratton1941,harrington2001}.

\paragraph{Solvability criterion.}
The forward radiation equation \(\mathcal{R}_\omega[\Jimp]=\mathbf{g}\) is solvable iff the
target \(\mathbf{g}\in\mathcal{F}_{\mathrm{rad}}\) is orthogonal to the cokernel of the
adjoint map on declared pairings:
\begin{equation}
\label{eq:ch2.radiation-solvability-criterion}
\mathcal{R}_\omega[\Jimp]=\mathbf{g}
\ \text{solvable}
\iff
\langle \mathbf{g},\,\mathbf{w}\rangle_{\mathcal{U}}=0
\ \ \forall\ \mathbf{w}\in\ker\mathcal{R}_\omega^{\dagger},
\end{equation}
with \(\mathcal{R}_\omega^{\dagger}\) the adjoint on
\((\mathcal{J}_\Omega,\mathcal{F}_{\mathrm{rad}})\) induced by
Eqs.~\eqref{eq:ch2.dual-source-pairing} and~\eqref{eq:ch2.maxwell-adjoint-pairing}
\cite{stratton1941,harrington2001}. Equation~\eqref{eq:ch2.radiation-solvability-criterion}
upgrades preview Eq.~\eqref{eq:ch2.spectral-solvability-preview} and aligns with biorthogonality
Eq.~\eqref{eq:ch2.radiation-biorthogonality} on modal adjoints
\cite{stratton1941}. When \(\mathbf{g}\) arises from far-pattern data, solvability is
further capped by Eq.~\eqref{eq:ch2.reconstruction-rank-cap} even if
Eq.~\eqref{eq:ch2.radiation-solvability-criterion} holds in the continuous model
\cite{harrington2001}.

\begin{theorem}[Fredholm structure of radiation operators]
\label{thm:ch2.radiation-fredholm-structure}
Let \(\mathcal{R}_\omega\) satisfy Eq.~\eqref{eq:ch2.radiation-fredholm-decomposition} on
\(\mathcal{J}_\Omega\) with \(\mathcal{R}_{0,\omega}\) Fredholm of index zero on the outgoing
quotient and \(K_{\mathrm{rad}}\) compact. Then \(\mathcal{R}_\omega\) is Fredholm,
\begin{equation}
\label{eq:ch2.radiation-index-stability}
\mathrm{ind}\,\mathcal{R}_\omega
=
\mathrm{ind}\,\mathcal{R}_{0,\omega},
\end{equation}
and Eq.~\eqref{eq:ch2.radiation-solvability-criterion} holds. For finite-rank observation
\(\mathcal{O}_{\mathrm{meas}}:\mathcal{F}_{\mathrm{rad}}\to\mathbb{C}^{N}\),
\begin{equation}
\label{eq:ch2.observation-fredholm-composition}
\mathcal{O}_{\mathrm{meas}}\mathcal{R}_\omega
\ \text{is Fredholm on}\ \mathcal{J}_\Omega
\ \Longleftrightarrow\
\mathcal{R}_\omega\ \text{Fredholm and}\ 
\mathrm{rank}(\mathcal{O}_{\mathrm{meas}})=N,
\end{equation}
with rank bound Eq.~\eqref{eq:ch2.rank-composition-law}
\cite{stratton1941,harrington2001}.
\end{theorem}
\begin{proof}
Compact perturbation stability of Fredholm property and index follows from standard operator
theory applied to Eq.~\eqref{eq:ch2.radiation-fredholm-decomposition}; Subsection
~\ref{sec:ch213} recorded the embedding inputs without computing index here
\cite{stratton1941,harrington2001}. Solvability criterion
Eq.~\eqref{eq:ch2.radiation-solvability-criterion} is the Fredholm alternative for
\(\mathcal{R}_\omega\). Composition with finite-rank \(\mathcal{O}_{\mathrm{meas}}\) preserves
Fredholm class and adjusts index by at most \(N\) on the range side, yielding
Eq.~\eqref{eq:ch2.observation-fredholm-composition} and
Eq.~\eqref{eq:ch2.rank-composition-law} \cite{stratton1941}.
\end{proof}

\paragraph{Modal Fredholm reduction.}
On the orthonormal radiation basis of Theorem~\ref{thm:ch2.radiation-spectral-decomposition},
\begin{equation}
\label{eq:ch2.modal-fredholm-reduction}
\mathcal{R}_\omega
=
\sum_{n\in\mathcal{N}_{\mathrm{rad}}}
\lambda_n\,P_n\circ\mathcal{S}_\omega,
\qquad
\lambda_n\neq 0
\ \Longleftrightarrow\
\text{mode }n\ \text{exporting},
\end{equation}
so solvability of \(\mathcal{R}_\omega[\Jimp]=\mathbf{g}\) reduces to matching coefficients
\(c_n[\mathfrak{D}]\) Eq.~\eqref{eq:ch2.source-modal-amplitude} on modes with
\(\lambda_n\neq 0\) \cite{stratton1941,harrington2001}. Modes with \(\lambda_n=0\) lie in
\(\ker\mathcal{R}_\omega\) at the source level Eq.~\eqref{eq:ch2.ker-R-omega-preview}; modes
with \(\mathbf{g}\) unsupported on \(\operatorname{ran}\mathcal{R}_\omega\) generate cokernel
obstructions in Eq.~\eqref{eq:ch2.radiation-solvability-criterion}
\cite{stratton1941}. Truncation to \(\mathcal{R}_\omega^{(N)}\) from
Eq.~\eqref{eq:ch2.spectral-compression-operator} yields approximate solvability with error
Eq.~\eqref{eq:ch2.compression-error-bound} controlled by \(\sigma_{N+1}(K_{\mathrm{rad}})\)
\cite{harrington2001}.

\begin{proposition}[Modal invertibility on the radiative range]
\label{prop:ch2.modal-invertibility-equivalence}
Let \(\mathbf{g}\in\operatorname{ran}\mathcal{R}_\omega\) and expand
\(\mathbf{g}=\sum_n g_n\,\mathbf{u}_n\) with \(\{\mathbf{u}_n\}\) from
Subsection~\ref{sec:ch261}. Then Eq.~\eqref{eq:ch2.radiation-solvability-criterion} holds if
and only if \(g_n=0\) whenever \(\lambda_n=0\). A particular solution is determined by
\begin{equation}
\label{eq:ch2.modal-particular-solution}
c_n
=
\frac{g_n}{\lambda_n}
\quad\text{for}\ \lambda_n\neq 0,
\end{equation}
with arbitrary coefficients on \(\ker\mathcal{R}_\omega\) not fixed by far-zone data
\cite{stratton1941,harrington2001,jackson1999}. Equation~\eqref{eq:ch2.modal-particular-solution}
is the spectral form of Eq.~\eqref{eq:ch2.inverse-source-equation} on the radiative range;
Tikhonov and regularized inversions from Subsection~\ref{sec:ch244} damp ill-conditioned
\(\lambda_n\) without altering Fredholm index \cite{stratton1941}.
\end{proposition}
\begin{proof}
Substitute Eq.~\eqref{eq:ch2.radiation-map-modal-sum} into
\(\mathcal{R}_\omega[\Jimp]=\mathbf{g}\) and use orthogonality of \(\{\mathbf{u}_n\}\) in
\(\mathcal{U}\) \cite{stratton1941}. Modes with \(\lambda_n=0\) do not appear in the range;
modes with \(\lambda_n\neq 0\) invert mode-wise yielding
Eq.~\eqref{eq:ch2.modal-particular-solution}. Kernel directions remain free because
\(\mathcal{R}_\omega\) is not injective on \(\mathcal{J}_\Omega\)
Eq.~\eqref{eq:ch2.ker-R-omega-preview} \cite{harrington2001}.
\end{proof}

\paragraph{Cokernel pairing and far patterns.}
Cokernel elements pair with unreachable far-zone content: if \(\mathbf{w}\in\ker\mathcal{R}_\omega^{\dagger}\),
then \(\langle \mathbf{g},\mathbf{w}\rangle_{\mathcal{U}}=0\) is required for solvability
Eq.~\eqref{eq:ch2.radiation-solvability-criterion} \cite{stratton1941}. On exterior domains,
continuous angular spectra Eq.~\eqref{eq:ch2.exterior-continuous-spectral-integral} imply
infinite-dimensional cokernel constraints in principle; finite apertures replace them with
rank-\(N\) projections through Eq.~\eqref{eq:ch2.observation-fredholm-composition}
\cite{harrington2001,jackson1999}. Proposition~\ref{prop:ch2.far-pattern-from-modes} expresses
\(\mathbf{g}\) through \(\mathbf{f}_n\); solvability requires consistency of measured
coefficients with Eq.~\eqref{eq:ch2.modal-particular-solution} on exporting modes only
\cite{stratton1941}.

\paragraph{Composition with Maxwell resolvent.}
On the outgoing branch, solvability of \(\mathcal{L}_\omega[\mathbf{u}_\omega]=\mathbf{s}_\omega\)
and solvability of \(\mathcal{R}_\omega[\Jimp]=\mathbf{g}\) are linked through
Eq.~\eqref{eq:ch2.integral-resolvent-equivalence}: forward field construction is single-valued
on \(\mathcal{F}_{\mathrm{SF}}\) while source inversion remains rank-deficient
\cite{stratton1941,harrington2001}. Fredholm index of \(\mathcal{L}_\omega\) on bounded
quotients Eq.~\eqref{eq:ch2.coercivity-quotient} is not identical to
\(\mathrm{ind}\,\mathcal{R}_\omega\); export selection removes non-radiative degrees before
pattern inversion \cite{stratton1941}. Subsection~\ref{sec:ch263} therefore treats
\(\mathcal{R}_\omega\) as the primary Fredholm object for far-zone inverse problems
\cite{harrington2001}.

\paragraph{Boundary of validity.}
Theorem~\ref{thm:ch2.radiation-fredholm-structure} requires linear media, fixed \(\omega\),
outgoing closure, and compactness of \(K_{\mathrm{rad}}\) on the declared source class
\cite{stratton1941,harrington2001}. Resonant cavities violate invertibility of
\(\mathcal{L}_0\) unless quotiented Eq.~\eqref{eq:ch2.coercivity-quotient}; radiation index
must be computed on the outgoing subspace after Eq.~\eqref{eq:ch2.outgoing-radiation-refinement}
\cite{stratton1941}. Loss regularizes small \(\lambda_n\) but does not remove
\(\ker\mathcal{R}_\omega\) Eq.~\eqref{eq:ch2.ker-R-omega-preview} \cite{harrington2001}.
Regime gate Eq.~\eqref{eq:ch2.regime-component-gate} remains mandatory before transport claims
on recovered \(\mathbf{g}\) \cite{stratton1941}.

\paragraph{Worked illustration (finite scanner rank).}
With \(N=50\) probe channels, \(\mathcal{O}_{\mathrm{meas}}\mathcal{R}_\omega\) has rank at
most \(50\) Eq.~\eqref{eq:ch2.rank-composition-law} \cite{stratton1941,harrington2001}.
Even when Eq.~\eqref{eq:ch2.radiation-solvability-criterion} holds for a continuous
\(\mathbf{g}\), discrete recovery satisfies only a truncated modal system; tail error is bounded
by Eq.~\eqref{eq:ch2.compression-error-bound} and \(\sigma_{51}\) from
Eq.~\eqref{eq:ch2.singular-value-decay} \cite{harrington2001}. Subsection~\ref{sec:ch213}
anticipated this composition; Subsection~\ref{sec:ch263} states the Fredholm target for
pattern inversion \cite{stratton1941}.

\paragraph{Worked illustration (non-radiating obstruction).}
If \(\mathbf{g}=\mathbf{0}\), Eq.~\eqref{eq:ch2.radiation-solvability-criterion} is always
solvable with \(\Jimp\in\ker\mathcal{R}_\omega\) \cite{stratton1941}. If \(\mathbf{g}\neq\mathbf{0}\)
but lies entirely in a cokernel direction paired with \(\ker\mathcal{R}_\omega^{\dagger}\),
no \(\Jimp\) satisfies forward export \cite{harrington2001}. Proposition
~\ref{prop:ch2.modal-invertibility-equivalence} detects the obstruction mode-wise through
\(g_n\) on \(\lambda_n=0\) sectors \cite{stratton1941}.

\paragraph{Worked illustration (Tikhonov versus Fredholm index).}
Regularized inversion from Subsection~\ref{sec:ch244} stabilizes small \(|\lambda_n|\) in
Eq.~\eqref{eq:ch2.modal-particular-solution} but leaves \(\mathrm{ind}\,\mathcal{R}_\omega\)
unchanged \cite{stratton1941,harrington2001}. A stable approximate solution therefore does not
imply unique source identification: components in \(\ker\mathcal{R}_\omega\) remain free at
all regularization levels \cite{stratton1941}. Observation rank
Eq.~\eqref{eq:ch2.reconstruction-rank-cap} adds a separate finite-rank cap
\cite{harrington2001}.

\paragraph{Closing summary.}
Subsection~\ref{sec:ch263} establishes the Chapter~\ref{ch:02} HOME for Fredholm decomposition
Eq.~\eqref{eq:ch2.radiation-fredholm-decomposition}, index
Eq.~\eqref{eq:ch2.radiation-fredholm-index-def}, solvability
Eq.~\eqref{eq:ch2.radiation-solvability-criterion}, Theorem
~\ref{thm:ch2.radiation-fredholm-structure}, modal reduction
Eq.~\eqref{eq:ch2.modal-fredholm-reduction}, observation composition
Eq.~\eqref{eq:ch2.observation-fredholm-composition}, and modal invertibility
Proposition~\ref{prop:ch2.modal-invertibility-equivalence}, importing Subsection
~\ref{sec:ch213} compactness templates by reference
\cite{stratton1941,harrington2001,jackson1999}. Subsection~\ref{sec:ch264} records geometric
constraints on admissible radiative spectra beyond Fredholm solvability; Section~\ref{sec:ch27}
classifies \(\ker\mathcal{R}_\omega\) in source space \cite{stratton1941}.
\subsection{Geometric constraints on radiative spectra}
\label{sec:ch264}
% TOC tag: [HOME]

Subsections~\ref{sec:ch261}--\ref{sec:ch263} fixed radiation eigenmodes, rotational
degeneracy, and Fredholm solvability on \(\mathcal{R}_\omega\). Subsection~\ref{sec:ch264}
gives the Chapter~\ref{ch:02} HOME for geometric constraints---support, boundary class,
material symmetry, and parity---that annihilate or restrict modal coefficients in
Eq.~\eqref{eq:ch2.spectral-far-pattern-expansion} independently of observation rank and
Fredholm index \cite{stratton1941,harrington2001,jackson1999}. Preview
Eq.~\eqref{eq:ch2.geometric-spectrum-constraint-preview} from Section~\ref{sec:ch26} is
upgraded here to constructive operators on declared modal bases; vector harmonic labels and
Clebsch--Gordan coupling belong to Chapter~\ref{ch:03} \cite{harrington2001}. Reciprocity
constraints on coupled modes remain in Section~\ref{sec:ch28}; complete null-space catalogs
in Section~\ref{sec:ch27} \cite{stratton1941}. Angular accessibility filters on plane-wave
spectra are developed in Section~\ref{sec:ch294}; Subsection~\ref{sec:ch264} states the
spectral-sector constraints those filters presuppose \cite{stratton1941,harrington2001}.

\paragraph{Assumptions.}
Let \(\Omega_s=\operatorname{supp}\Jimp\), outer boundary \(\Gamma\), piecewise material
data \((\varepsilon,\mu)\), outgoing closure, and radiation eigenbasis
\(\{\mathbf{u}_n,\mathbf{f}_n\}_{n\in\mathcal{N}_{\mathrm{rad}}}\) from
Subsections~\ref{sec:ch261}--\ref{sec:ch262} \cite{stratton1941,harrington2001}. Modal
coefficients \(a_n\) in Eq.~\eqref{eq:ch2.spectral-far-pattern-expansion} are identified
with \(\lambda_n c_n[\mathfrak{D}]\) on exporting modes
Proposition~\ref{prop:ch2.far-pattern-from-modes} \cite{stratton1941}. Geometric symmetry
group \(G_{\mathrm{geo}}\subseteq\mathrm{SO}(3)\times\{+\}\) acts on sources, boundaries,
and constitutive tensors as in Section~\ref{sec:ch01.3.2} by reference
\cite{stratton1941}. Regime gate Eq.~\eqref{eq:ch2.regime-operator-screen} remains mandatory
before transport claims on any surviving modal sector \cite{harrington2001}.

\paragraph{Geometric symmetry operator.}
Define a bounded operator \(\mathcal{C}_{\mathrm{geo}}\) on modal amplitudes by
\begin{equation}
\label{eq:ch2.geometric-symmetry-operator-def}
\big(\mathcal{C}_{\mathrm{geo}}\mathbf{a}\big)_n
=
\sum_{m\in\mathcal{N}_{\mathrm{rad}}}
\Gamma_{nm}\,a_m,
\qquad
\mathbf{a}=\{a_m\}_{m\in\mathcal{N}_{\mathrm{rad}}},
\end{equation}
where \(\Gamma_{nm}\) encodes how geometry and boundary class map modal sectors under
\(G_{\mathrm{geo}}\) \cite{stratton1941,harrington2001}. Equation
\eqref{eq:ch2.geometric-symmetry-operator-def} is the Subsection~\ref{sec:ch264} HOME
realization of preview Eq.~\eqref{eq:ch2.geometric-spectrum-constraint-preview}: when
\(\mathcal{C}_{\mathrm{geo}}\mathbf{e}_n=\mathbf{0}\) for a unit vector \(\mathbf{e}_n\),
mode \(n\) is \emph{geometrically forbidden} even if \(\lambda_n\neq 0\) and
Eq.~\eqref{eq:ch2.radiation-solvability-criterion} holds for other modes
\cite{stratton1941}. Rotation commutation Eq.~\eqref{eq:ch2.radiation-rotation-commutation}
acts on \(\mathbf{f}_n\); \(\mathcal{C}_{\mathrm{geo}}\) acts on admissible \emph{sectors}
after symmetry reduction \cite{harrington2001,jackson1999}.

\paragraph{Support and compactness constraints.}
Compact support \(\Omega_s\subset\mathbb{R}^{3}\) implies entire-function structure in
spatial wavenumber: for each outgoing mode \(\mathbf{f}_n(\theta,\phi)\),
\begin{equation}
\label{eq:ch2.support-spectral-support-condition}
a_n(\Jimp)=0
\ \text{unless}\ 
\mathbf{f}_n\ \text{has non-negligible overlap with}\ 
\mathcal{F}_\omega[\mathcal{J}_{\Omega_s}],
\end{equation}
with \(\mathcal{J}_{\Omega_s}\) the restriction of admissible sources to \(\Omega_s\)
\cite{stratton1941,harrington2001}. Equation~\eqref{eq:ch2.support-spectral-support-condition}
is a support-induced geometric filter distinct from Fredholm cokernel constraints
Eq.~\eqref{eq:ch2.radiation-solvability-criterion}: electrically small \(\Omega_s\)
suppresses high-\(\ell\) sectors in Eq.~\eqref{eq:ch2.spectral-far-pattern-expansion}
through \((k_0 a)^{2\ell}\) scaling deferred to Chapter~\ref{ch:03}
\cite{jackson1999,stratton1941}. Rank cap Eq.~\eqref{eq:ch2.reconstruction-rank-cap}
limits measurable \(a_n\) separately from support annihilation
\cite{harrington2001}.

\paragraph{Boundary and image constraints.}
Perfect conductors, symmetry planes, and equivalent image configurations impose linear
relations among modal amplitudes:
\begin{equation}
\label{eq:ch2.boundary-image-constraint}
\sum_{m\in\mathcal{N}_{\mathrm{rad}}}
M_{nm}(\Gamma,\mathcal{B}_\Gamma)\,a_m=0
\quad\Longrightarrow\quad
a_n=0\ \text{for forbidden parity sectors,}
\end{equation}
with \(M_{nm}\) fixed by tangential boundary class
Eq.~\eqref{eq:ch1.outer-boundary-classes} and image rules on outgoing Green branches
Eq.~\eqref{eq:ch2.outgoing-green-branch} \cite{stratton1941,harrington2001}. Equation
\eqref{eq:ch2.boundary-image-constraint} specializes Proposition
~\ref{prop:ch2.degeneracy-splitting-broken-symmetry}: a ground plane replaces full
\(\mathrm{SO}(3)\) by a reflection group and annihilates odd or even combinations per
declared image parity \cite{stratton1941,jackson1999}. Sommerfeld outgoing closure from
Section~\ref{sec:ch24} is presupposed; incoming image branches are excluded from
\(\mathcal{F}_{\mathrm{rad}}\) \cite{harrington2001}.

\paragraph{Parity and sector selection.}
When \(G_{\mathrm{geo}}\) contains mirror reflection \(\sigma\), modal sectors split into
parity classes \(\Pi_\pm\) with
\begin{equation}
\label{eq:ch2.parity-selection-rule}
a_n=0\ \text{for all}\ n\in\mathcal{N}_{\mathrm{rad}}
\ \text{with}\ 
\Pi_\pm\mathbf{f}_n=\mp\mathbf{f}_n
\ \text{and forbidden parity,}
\end{equation}
as in the ground-plane illustration of Subsection~\ref{sec:ch262}
\cite{stratton1941,harrington2001}. Equation~\eqref{eq:ch2.parity-selection-rule} explains
pattern nulls that are not explained by \(\ker\mathcal{O}_{\mathrm{meas}}\) or
\(\ker\mathcal{R}_\omega\) alone \cite{stratton1941}. Slot and patch symmetries reduce
Eq.~\eqref{eq:ch2.parity-selection-rule} to discrete sector labels on
\(\mathcal{N}_{\mathrm{rad}}\) without re-tabulating dyads from Subsection~\ref{sec:ch231}
\cite{harrington2001}.

\begin{figure}[t]
\centering
\ChOneFigBlock{%
\begin{tikzpicture}[ch1fig, node distance=7mm and 10mm]
  \node[ch1box] (g) {$G_{\mathrm{geo}}$};
  \node[ch1box, right=12mm of g] (c) {$\mathcal{C}_{\mathrm{geo}}$};
  \node[ch1box, right=12mm of c] (a) {$\{a_n\}$};
  \node[ch1box, right=12mm of a] (f) {$\mathbf{f}_n$};
  \node[ch1box, below=11mm of c] (forb) {forbidden sectors};
  \draw[ch1flow] (g) -- (c) -- (a) -- (f);
  \draw[ch1flow] (c) -- (forb);
\end{tikzpicture}%
}
\caption{Subsection~\ref{sec:ch264} geometric filter: symmetry group \(G_{\mathrm{geo}}\)
acts through \(\mathcal{C}_{\mathrm{geo}}\) Eq.~\eqref{eq:ch2.geometric-symmetry-operator-def}
on modal amplitudes before far patterns \(\mathbf{f}_n\).}
\label{fig:ch2.geometric-spectrum-constraints}
\end{figure}

Figure~\ref{fig:ch2.geometric-spectrum-constraints} orders geometric constraints after
spectral and Fredholm structure: \(\mathcal{C}_{\mathrm{geo}}\) acts on coefficient space,
not on raw \(\mathcal{S}_\omega[\Jimp]\) \cite{stratton1941,harrington2001}.

\begin{theorem}[Geometric constraints on radiative spectra]
\label{thm:ch2.geometric-spectral-constraints}
Let \(\{\mathbf{f}_n\}_{n\in\mathcal{N}_{\mathrm{rad}}}\) be a modal basis for
\(\operatorname{ran}\mathcal{R}_\omega\) and let \(G_{\mathrm{geo}}\) leave
\((\Omega_s,\Gamma,\varepsilon,\mu)\) invariant. Then there exists
\(\mathcal{C}_{\mathrm{geo}}\) Eq.~\eqref{eq:ch2.geometric-symmetry-operator-def} such that
\begin{equation}
\label{eq:ch2.geometric-modal-annihilation}
a_n(\Jimp)=0
\quad\text{whenever}\quad
\big(\mathcal{C}_{\mathrm{geo}}\mathbf{f}_n\big)=0
\ \text{or}\ 
\mathbf{f}_n\notin\operatorname{span}\{\mathcal{R}_\omega[\mathcal{J}_{\Omega_s}]\},
\end{equation}
with \(a_n\) from Eq.~\eqref{eq:ch2.spectral-far-pattern-expansion} on exporting modes
\cite{stratton1941,harrington2001,jackson1999}. The accessible modal sector is
\begin{equation}
\label{eq:ch2.accessible-modal-sector}
\mathcal{N}_{\mathrm{acc}}
=
\{n\in\mathcal{N}_{\mathrm{rad}}:\ 
\lambda_n\neq 0,\ 
\mathcal{C}_{\mathrm{geo}}\mathbf{f}_n\neq 0,\ 
\mathbf{f}_n\ \text{satisfies Eqs.~\eqref{eq:ch2.support-spectral-support-condition}--\eqref{eq:ch2.parity-selection-rule}}
\}.
\end{equation}
\end{theorem}
\begin{proof}
Symmetry of Maxwell data under \(G_{\mathrm{geo}}\) induces symmetries of
\(\mathcal{S}_\omega\) and \(\mathcal{R}_\omega\) compatible with
Eq.~\eqref{eq:ch2.radiation-rotation-commutation} when rotations lie in
\(G_{\mathrm{geo}}\) \cite{stratton1941}. Boundary and image conditions linearly constrain
tangential traces, yielding relations Eq.~\eqref{eq:ch2.boundary-image-constraint} among
far-zone modal coefficients Proposition~\ref{prop:ch2.far-pattern-from-modes}
\cite{harrington2001,jackson1999}. Support constraints give
Eq.~\eqref{eq:ch2.support-spectral-support-condition}. Collecting independent constraints
defines \(\mathcal{C}_{\mathrm{geo}}\) and the surviving index set
Eq.~\eqref{eq:ch2.accessible-modal-sector} \cite{stratton1941}.
\end{proof}

\begin{proposition}[Angular accessibility preview]
\label{prop:ch2.angular-accessibility-preview}
Let \(\widehat{\mathbf{k}}\in S^{2}\) and let \(a(\widehat{\mathbf{k}})\) denote the
continuous angular spectrum Eq.~\eqref{eq:ch2.exterior-continuous-spectral-integral}. Then
only directions in the accessible sector
\begin{equation}
\label{eq:ch2.angular-accessibility-sector}
\mathcal{K}_{\mathrm{acc}}
=
\{\widehat{\mathbf{k}}:\ a(\widehat{\mathbf{k}})\ \text{not annihilated by}\ \mathcal{C}_{\mathrm{geo}}\}
\end{equation}
contribute to transport reports after Eq.~\eqref{eq:ch2.regime-operator-screen}
\cite{stratton1941,harrington2001}. Section~\ref{sec:ch294} develops plane-wave filters and
retained/suppressed sheets on the same sector; Subsection~\ref{sec:ch264} fixes the
spectral-domain constraint Eq.~\eqref{eq:ch2.accessible-modal-sector} those filters inherit
\cite{stratton1941}.
\end{proposition}
\begin{proof}
Modal expansion Eq.~\eqref{eq:ch2.spectral-far-pattern-expansion} and continuous limit
Eq.~\eqref{eq:ch2.exterior-continuous-spectral-integral} identify discrete and continuous
angular coefficients on outgoing transverse modes \cite{stratton1941,jackson1999}.
Geometric annihilation Eq.~\eqref{eq:ch2.geometric-modal-annihilation} removes entire
subsets of \(\widehat{\mathbf{k}}\) before observation composition
Eq.~\eqref{eq:ch2.observation-fredholm-composition} \cite{harrington2001}.
\end{proof}

\paragraph{Composition with Fredholm and observation deficits.}
Geometric, Fredholm, and observation constraints multiply:
\begin{equation}
\label{eq:ch2.constraint-composition-law}
\mathcal{N}_{\mathrm{reportable}}
\subseteq
\mathcal{N}_{\mathrm{acc}}
\cap
\{n:\ \text{resolvable by}\ \mathcal{O}_{\mathrm{meas}}\}
\cap
\{n:\ \lambda_n\neq 0\},
\end{equation}
with the last two factors from Subsection~\ref{sec:ch263} and
Eq.~\eqref{eq:ch2.reconstruction-rank-cap} \cite{stratton1941,harrington2001}. Equation
\eqref{eq:ch2.constraint-composition-law} prevents attributing a pattern null solely to
finite aperture when Eq.~\eqref{eq:ch2.geometric-modal-annihilation} already forbids the
sector \cite{stratton1941}. Inverse source Eq.~\eqref{eq:ch2.inverse-source-equation} on
\(\mathcal{N}_{\mathrm{reportable}}\) remains rank-deficient through
\(\ker\mathcal{R}_\omega\) Eq.~\eqref{eq:ch2.ker-R-omega-preview} even when geometry is
exact \cite{harrington2001}.

\paragraph{Boundary of validity.}
Theorem~\ref{thm:ch2.geometric-spectral-constraints} assumes linear media, fixed \(\omega\),
and symmetry of \((\Omega_s,\Gamma,\varepsilon,\mu)\) under declared \(G_{\mathrm{geo}}\)
\cite{stratton1941,harrington2001}. Weak symmetry breaking perturbs
Eq.~\eqref{eq:ch2.accessible-modal-sector} through lifted degeneracy
Proposition~\ref{prop:ch2.degeneracy-splitting-broken-symmetry} rather than strict zero
\cite{stratton1941}. Loss and anisotropy modify \(\Gamma_{nm}\) and may reopen partially
forbidden sectors at finite amplitude \cite{stratton1941}. Regime gate
Eq.~\eqref{eq:ch2.regime-component-gate} is unchanged: geometric accessibility does not
promote reactive near fields to export \cite{harrington2001}.

\paragraph{Worked illustration (dipole over PEC ground).}
A vertical dipole above a PEC plane excites even image parity; odd parity sectors satisfy
Eq.~\eqref{eq:ch2.parity-selection-rule} with \(a_n=0\) \cite{stratton1941,jackson1999}.
Measured sidelobes that appear absent in a restricted sector should be checked against
Eq.~\eqref{eq:ch2.geometric-modal-annihilation} before invoking observation noise
\cite{harrington2001}. Ground proximity also breaks full
Eq.~\eqref{eq:ch2.spectral-degeneracy-multiplicity} per
Proposition~\ref{prop:ch2.degeneracy-splitting-broken-symmetry} \cite{stratton1941}.

\paragraph{Worked illustration (slot symmetry).}
A slot antenna symmetric under half-wavelength reflection forbids modal sectors whose
\(\mathbf{f}_n\) change sign under the slot reflection through
Eq.~\eqref{eq:ch2.boundary-image-constraint} \cite{stratton1941,harrington2001}. Radiation
integrals from Theorem~\ref{thm:ch2.source-to-field-integral} may be exact while
Eq.~\eqref{eq:ch2.accessible-modal-sector} contains only a fraction of
\(\mathcal{N}_{\mathrm{rad}}\) \cite{stratton1941}. Section~\ref{sec:ch294} refines the
same sector in plane-wave language \cite{harrington2001}.

\paragraph{Worked illustration (electrically small sphere).}
For \(k_0 a\ll 1\), Eq.~\eqref{eq:ch2.support-spectral-support-condition} leaves primarily
low-\(\ell\) sectors in Eq.~\eqref{eq:ch2.spectral-far-pattern-expansion}; high-\(n\)
coefficients are geometrically suppressed even when probes could resolve them at larger
sources \cite{jackson1999,stratton1941}. This suppression is distinct from compression error
Eq.~\eqref{eq:ch2.compression-error-bound} and from \(\ker\mathcal{R}_\omega\)
\cite{harrington2001}.

\paragraph{Closing summary.}
Subsection~\ref{sec:ch264} establishes the Chapter~\ref{ch:02} HOME for geometric symmetry
operator Eq.~\eqref{eq:ch2.geometric-symmetry-operator-def}, support condition
Eq.~\eqref{eq:ch2.support-spectral-support-condition}, boundary/image constraint
Eq.~\eqref{eq:ch2.boundary-image-constraint}, parity rule
Eq.~\eqref{eq:ch2.parity-selection-rule}, annihilation Theorem
~\ref{thm:ch2.geometric-spectral-constraints}, accessible sector
Eq.~\eqref{eq:ch2.accessible-modal-sector}, angular accessibility preview
Proposition~\ref{prop:ch2.angular-accessibility-preview}, and composition law
Eq.~\eqref{eq:ch2.constraint-composition-law} \cite{stratton1941,harrington2001,jackson1999}.

\medskip
\noindent
Section~\ref{sec:ch26} closes with the spectral chain begun in its introduction:
eigenfunctions and spectral decomposition Theorem~\ref{thm:ch2.radiation-spectral-decomposition}
(Subsection~\ref{sec:ch261}), rotational degeneracy Theorem
~\ref{thm:ch2.rotation-spectral-degeneracy} importing Section~\ref{sec:ch01.3}
(Subsection~\ref{sec:ch262}), Fredholm structure and solvability Theorem
~\ref{thm:ch2.radiation-fredholm-structure} importing Subsection~\ref{sec:ch213}
(Subsection~\ref{sec:ch263}), and geometric constraints on admissible radiative spectra
Theorem~\ref{thm:ch2.geometric-spectral-constraints} with accessible sector
Eq.~\eqref{eq:ch2.accessible-modal-sector} (Subsection~\ref{sec:ch264})
\cite{stratton1941,harrington2001,jackson1999}. Evaluation integrals and flux meters from
Section~\ref{sec:ch25} now carry modal coordinates: \(\mathcal{R}_\omega\) is spectrally
decomposed, rotationally classified, Fredholm-solvable on its range, and geometrically
filtered before far-zone reports \cite{stratton1941}. Section~\ref{sec:ch27} completes
non-radiating null spaces at the source level; Section~\ref{sec:ch28} adds reciprocity;
Chapter~\ref{ch:03} supplies explicit harmonic bases for \(\mathbf{f}_n\); Section
~\ref{sec:ch294} develops angular-spectrum action on the accessible sector
Eq.~\eqref{eq:ch2.angular-accessibility-sector} \cite{harrington2001}. Transport and
angular-transport claims throughout Volume~I must compose with
Eq.~\eqref{eq:ch2.section26-spectral-chain},
Eq.~\eqref{eq:ch2.constraint-composition-law}, and regime gate
Eq.~\eqref{eq:ch2.regime-operator-screen}, not with unconstrained modal labels on raw
\(\mathcal{S}_\omega[\Jimp]\) \cite{stratton1941,harrington2001}.
\section{Non-Radiating Sources and Null Spaces}
\label{sec:ch27}

Section~\ref{sec:ch26} closed the spectral chain
Eq.~\eqref{eq:ch2.section26-spectral-chain}: radiation eigenmodes, rotational degeneracy,
Fredholm solvability, and geometric filters on admissible modal sectors
\cite{stratton1941,harrington2001,jackson1999}. That program organized
\(\mathcal{R}_\omega[\Jimp]\) in modal coordinates and explained which far-zone patterns are
reachable, degenerate, or symmetry-forbidden. Section~\ref{sec:ch27} addresses the
complementary \emph{source-side} question: which impressed currents produce no radiative
export at all, how those currents relate to evanescent and non-propagating field content,
and what null-space structure implies for inverse problems and angular spectra
\cite{stratton1941,harrington2001}. Preview Eq.~\eqref{eq:ch2.ker-R-omega-preview} from
Subsection~\ref{sec:ch244} named \(\ker\mathcal{R}_\omega\); Section~\ref{sec:ch27} develops
constructive classification, canonical non-radiating configurations, and reconstructability
limits without re-opening Green construction, Sommerfeld closure, or radiation integrals from
Sections~\ref{sec:ch23}--\ref{sec:ch25} \cite{stratton1941}. Reciprocity constraints on
coupled modes belong to Section~\ref{sec:ch28}; plane-wave angular-spectrum action to
Section~\ref{sec:ch29}; explicit harmonic bases to Chapter~\ref{ch:03}
\cite{harrington2001}.

The problem addressed here is \emph{invisibility at the radiation map}. Field-level
uniqueness after outgoing closure (Theorem~\ref{thm:ch2.exterior-radiative-uniqueness}) does
not imply source identifiability: many admissible \(\Jimp\in\mathcal{J}_\Omega\) can share
the same \(\mathcal{R}_\omega[\Jimp]\) through Eq.~\eqref{eq:ch2.ker-R-omega-preview}
\cite{stratton1941,harrington2001}. Spectral mode \(\lambda_n=0\) in
Eq.~\eqref{eq:ch2.modal-eigenvalue-equation-preview} marks a radiative eigen-direction with
vanishing export, yet \(\ker\mathcal{R}_\omega\) is generally infinite-dimensional on open
domains and is not exhausted by finitely many modal labels
\cite{harrington2001,stratton1941}. Using cavity modal books or truncated harmonic sums to
\emph{approximate} \(\ker\mathcal{R}_\omega\) misstates identifiability even when forward
integrals from Theorem~\ref{thm:ch2.source-to-field-integral} are exact
\cite{jackson1999,harrington2001}. Section~\ref{sec:ch27} separates operator null spaces,
field-side non-radiative components Eqs.~\eqref{eq:ch2.radiative-component-def}--\eqref{eq:ch2.nonradiative-component-def},
and observation null spaces Eq.~\eqref{eq:ch2.observation-null-space}, then forwards
reciprocity and angular-spectrum filters to Sections~\ref{sec:ch28}--\ref{sec:ch29}
\cite{stratton1941}.

\paragraph{Null-space development chain.}
Let \(\Omega_s=\operatorname{supp}\Jimp\) be compact, \(\omega>0\) fixed, and outgoing field
pairs lie in \(\mathcal{F}_{\mathrm{SF}}\) from
Eq.~\eqref{eq:ch2.admissible-field-space}. Section~\ref{sec:ch27} organizes analysis in four
stages:
\begin{equation}
\label{eq:ch2.section27-nullspace-chain}
\ker\mathcal{R}_\omega
\ \xrightarrow{\ \text{operator structure}\ }
\text{null-space bases}
\ \xrightarrow{\ \text{canonical sources}\ }
\text{NR current families}
\ \xrightarrow{\ \text{evanescent overlap}\ }
\text{non-reconstructability}
\ \xrightarrow{\ \text{angular spectra}\ }
\text{suppressed sheets},
\end{equation}
where the first arrow is Subsection~\ref{sec:ch271}, the second is Subsection~\ref{sec:ch272},
the third is Subsection~\ref{sec:ch273}, and the fourth is Subsection~\ref{sec:ch274}
\cite{stratton1941,harrington2001}. Equation~\eqref{eq:ch2.section27-nullspace-chain}
restates Eq.~\eqref{eq:ch2.section26-spectral-chain} at the source level: spectral
\(\lambda_n=0\) modes from Eq.~\eqref{eq:ch2.spectral-resolution-composition} lie in
\(\mathbf{u}_{\mathrm{nr}}\), but complete \(\ker\mathcal{R}_\omega\) requires impressed-data
coordinates on \(\mathcal{J}_\Omega\), not only far-zone modal shapes
\cite{stratton1941}. Composition law Eq.~\eqref{eq:ch2.constraint-composition-law} still
applies: geometric annihilation Eq.~\eqref{eq:ch2.geometric-modal-annihilation} and
observation rank Eq.~\eqref{eq:ch2.reconstruction-rank-cap} are independent of, and
composable with, source null-space membership \cite{harrington2001}.

\paragraph{Operator targets.}
The primary object remains the radiation map
\begin{equation}
\label{eq:ch2.nullspace-operator-targets}
\mathcal{R}_\omega:\mathcal{J}_\Omega\longrightarrow\mathcal{F}_{\mathrm{rad}},
\qquad
\mathcal{R}_\omega[\Jimp]=\Pi_{\mathrm{rad}}\,\mathcal{S}_\omega[\Jimp],
\end{equation}
from Eqs.~\eqref{eq:ch2.radiation-operator-def}--\eqref{eq:ch2.outgoing-radiation-refinement}
\cite{stratton1941,harrington2001}. Secondary targets include the observation map
\(\mathcal{M}_\omega=\mathcal{O}_{\mathrm{meas}}\circ\Lambda_{\mathrm{rad}}\) from
Eq.~\eqref{eq:ch2.reconstruction-observation-map}, the non-radiative field projector
\(\mathrm{id}-\Pi_{\mathrm{rad}}\) on \(\mathcal{F}_\Omega\), and evanescent/non-propagating
selectors Eqs.~\eqref{eq:ch2.spectral-projectors}--\eqref{eq:ch2.nonpropagating-component-def}
from Subsection~\ref{sec:ch224} \cite{stratton1941}. Inverse source Eq.~\eqref{eq:ch2.inverse-source-equation}
is ill-posed whenever \(\ker\mathcal{M}_\omega\neq\{\mathbf{0}\}\), which occurs if
\(\ker\mathcal{R}_\omega\neq\{\mathbf{0}\}\) or \(\ker\mathcal{O}_{\mathrm{meas}}\) is
nontrivial on \(\operatorname{ran}\mathcal{R}_\omega\) \cite{harrington2001,stratton1941}.
Fredholm solvability Eq.~\eqref{eq:ch2.radiation-solvability-criterion} governs reachability
of a \emph{target} radiative field; Section~\ref{sec:ch27} classifies sources that map to
\(\mathbf{0}\in\mathcal{F}_{\mathrm{rad}}\) regardless of solvability for nonzero targets
\cite{stratton1941}.

\paragraph{Radiation null space (preview).}
Subsection~\ref{sec:ch271} upgrades preview Eq.~\eqref{eq:ch2.ker-R-omega-preview} to
structural statements on \(\ker\mathcal{R}_\omega\) and \(\ker\mathcal{R}_\omega^{\dagger}\):
\begin{equation}
\label{eq:ch2.radiation-nullspace-structure-preview}
\ker\mathcal{R}_\omega
=
\left\{
\Jimp\in\mathcal{J}_\Omega:
\Pi_{\mathrm{rad}}\,\mathcal{S}_\omega[\Jimp]=\mathbf{0},
\ \mathcal{D}_{J}[\Jimp]=1
\right\},
\qquad
\dim\ker\mathcal{R}_\omega=\infty\ \text{typ. on open domains},
\end{equation}
with admissibility Eq.~\eqref{eq:ch2.source-admissibility} enforced throughout
\cite{stratton1941,harrington2001}. Equation~\eqref{eq:ch2.radiation-nullspace-structure-preview}
does not re-define Eq.~\eqref{eq:ch2.ker-R-omega-preview}; it records the Chapter~\ref{ch:02}
HOME target for operator-theoretic characterization, dimensionality, and pairing with cokernel
directions obstructing inversion Eq.~\eqref{eq:ch2.spectral-solvability-preview}
\cite{stratton1941}. Modal sectors with \(\lambda_n=0\) in
Eq.~\eqref{eq:ch2.spectral-resolution-composition} supply a finite-dimensional \emph{slice}
of \(\ker\mathcal{R}_\omega\) only when eigenmodes are declared on outgoing subspaces
\cite{harrington2001}.

\paragraph{Non-radiating current configurations (preview).}
Subsection~\ref{sec:ch272} develops canonical impressed currents that realize
Eq.~\eqref{eq:ch2.null-radiation-preview}:
\begin{equation}
\label{eq:ch2.nr-current-template-preview}
\Jimp^{\mathrm{(NR)}}\in\ker\mathcal{R}_\omega,
\qquad
\mathcal{S}_\omega[\Jimp^{\mathrm{(NR)}}]\neq\mathbf{0},
\qquad
\|\mathbf{u}_{\mathrm{nr}}\|\ \text{may be large near}\ \Omega_s,
\end{equation}
with \(\mathbf{u}_{\mathrm{nr}}\) from Eq.~\eqref{eq:ch2.nonradiative-component-def}
\cite{stratton1941,harrington2001,jackson1999}. Equation~\eqref{eq:ch2.nr-current-template-preview}
specializes Eq.~\eqref{eq:ch2.nr-image-of-null-space} to constructive families---loop modes,
counter-propagating pairs, and equivalent-surface cancellations compatible with
Eq.~\eqref{eq:ch1.charge-continuity-harmonic}---without tabulating dyadic kernels from
Subsection~\ref{sec:ch231} \cite{stratton1941}. Such currents are invisible to far-zone flux
meters Eq.~\eqref{eq:ch2.power-flux-functional} yet may dominate near-zone energy budgets
Eq.~\eqref{eq:ch2.component-power-budget} \cite{harrington2001}.

\paragraph{Evanescent overlap and non-reconstructability (preview).}
Subsection~\ref{sec:ch273} links \(\ker\mathcal{R}_\omega\) to evanescent subspaces and
finite observation rank:
\begin{equation}
\label{eq:ch2.evanescent-nonreconstruct-preview}
\ker\mathcal{O}_{\mathrm{meas}}\circ\mathcal{R}_\omega
\supseteq
\ker\mathcal{R}_\omega,
\qquad
\Pi_{\mathrm{ev}}\mathcal{S}_\omega[\Jimp]
\ \text{never exports through}\ \Pi_{\mathrm{rad}},
\end{equation}
using Eqs.~\eqref{eq:ch2.rad-null-on-evanescent} and~\eqref{eq:ch2.reconstruction-rank-cap}
\cite{stratton1941,harrington2001}. Equation~\eqref{eq:ch2.evanescent-nonreconstruct-preview}
separates \emph{radiation nullity} (no outgoing transverse export) from \emph{measurement
nullity} (unobservable even when export is nonzero) \cite{stratton1941}. Tikhonov stabilization
Eq.~\eqref{eq:ch2.tikhonov-source-reconstruction} does not shrink \(\ker\mathcal{R}_\omega\);
regularization only damps ill-conditioned directions on \(\operatorname{ran}\mathcal{R}_\omega\)
\cite{harrington2001}. Accessible modal sector
Eq.~\eqref{eq:ch2.accessible-modal-sector} from Subsection~\ref{sec:ch264} further restricts
recoverable content independently of evanescent storage \cite{stratton1941}.

\paragraph{Angular-spectrum consequences (preview).}
Subsection~\ref{sec:ch274} records how null-space structure propagates to continuous angular
spectra Eq.~\eqref{eq:ch2.exterior-continuous-spectral-integral}:
\begin{equation}
\label{eq:ch2.angular-spectrum-null-preview}
a(\widehat{\mathbf{k}})=0
\ \text{on suppressed sheets}\ \mathcal{K}_{\mathrm{sup}}
\ \Longrightarrow\
\Jimp\ \text{may still carry reactive content in}\ \Pi_{\mathrm{ev}}\mathcal{S}_\omega[\Jimp],
\end{equation}
with \(a(\widehat{\mathbf{k}})\) the far angular amplitude and
\(\mathcal{K}_{\mathrm{sup}}\) the complement of accessible sector
Eq.~\eqref{eq:ch2.angular-accessibility-sector} \cite{stratton1941,harrington2001}. Equation
\eqref{eq:ch2.angular-spectrum-null-preview} bridges Section~\ref{sec:ch27} to
Section~\ref{sec:ch29}: plane-wave filters and retained/suppressed components in
Subsections~\ref{sec:ch291}--\ref{sec:ch294} presuppose the
source-side null-space catalog developed here \cite{stratton1941}. Proposition
~\ref{prop:ch2.angular-accessibility-preview} from Subsection~\ref{sec:ch264} is not re-proved;
Subsection~\ref{sec:ch274} states angular-spectrum \emph{consequences} of
Eq.~\eqref{eq:ch2.ker-R-omega-preview} \cite{harrington2001}.

\begin{figure}[t]
\centering
\ChOneFigBlock{%
\begin{tikzpicture}[ch1fig, node distance=7mm and 10mm]
  \node[ch1box] (op) {Operator\\$\ker\mathcal{R}_\omega$};
  \node[ch1box, right=13mm of op] (nr) {NR current\\configurations};
  \node[ch1box, right=13mm of nr] (ev) {Evanescent\\non-reconstruct};
  \node[ch1box, right=13mm of ev] (ang) {Angular-spectrum\\consequences};
  \node[ch1box, below=12mm of op] (j) {$\mathcal{J}_\Omega$};
  \node[ch1box, below=12mm of nr] (s) {$\mathcal{S}_\omega$};
  \node[ch1box, below=12mm of ev] (pi) {$\Pi_{\mathrm{ev}}$};
  \draw[ch1flow] (j) -- (op) -- (nr) -- (ev) -- (ang);
  \draw[ch1flow] (s) -- (nr);
  \draw[ch1flow] (pi) -- (ev);
\end{tikzpicture}%
}
\caption{Section~\ref{sec:ch27} development path: null spaces of radiation operators
(Subsection~\ref{sec:ch271}), non-radiating current configurations
(Subsection~\ref{sec:ch272}), evanescent subspaces and non-reconstructability
(Subsection~\ref{sec:ch273}), and consequences for angular spectra
(Subsection~\ref{sec:ch274}).}
\label{fig:ch2.nullspace-roadmap}
\end{figure}

Figure~\ref{fig:ch2.nullspace-roadmap} orders the four subsections.
Subsection~\ref{sec:ch271} gives the Chapter~\ref{ch:02} HOME for structural properties of
\(\ker\mathcal{R}_\omega\), adjoint null spaces, and their relation to spectral
\(\lambda_n=0\) sectors Eq.~\eqref{eq:ch2.spectral-resolution-composition}, upgrading
Eq.~\eqref{eq:ch2.ker-R-omega-preview} without repeating Fredholm proofs from
Subsection~\ref{sec:ch263} \cite{stratton1941,harrington2001}. Subsection~\ref{sec:ch272}
constructs canonical non-radiating impressed currents satisfying
Eq.~\eqref{eq:ch2.source-admissibility} and charge continuity
Eq.~\eqref{eq:ch1.charge-continuity-harmonic} by reference \cite{stratton1941}. Subsection
~\ref{sec:ch273} analyzes evanescent-dominated subspaces, overlap with
Eq.~\eqref{eq:ch2.nonpropagating-component-def}, and limits of source reconstruction beyond
Eq.~\eqref{eq:ch2.reconstruction-rank-cap} \cite{harrington2001}. Subsection~\ref{sec:ch274}
translates null-space structure to continuous angular spectra and forward-links
Section~\ref{sec:ch294} angular accessibility \cite{stratton1941}.

\paragraph{Field-side versus source-side nullity.}
Radiative split Eqs.~\eqref{eq:ch2.radiative-component-def}--\eqref{eq:ch2.nonradiative-component-def}
operate on \(\mathcal{F}_\Omega\); \(\ker\mathcal{R}_\omega\) operates on
\(\mathcal{J}_\Omega\). The implication Eq.~\eqref{eq:ch2.nr-image-of-null-space} is
one-way: \(\Jimp\in\ker\mathcal{R}_\omega\Rightarrow\mathbf{u}_{\mathrm{rad}}=\mathbf{0}\),
while large \(\|\mathbf{u}_{\mathrm{nr}}\|\) does not imply membership in
\(\ker\mathcal{R}_\omega\) \cite{stratton1941,harrington2001}. Propagation projectors from
Subsection~\ref{sec:ch224} explain why evanescent content never reaches
\(\mathcal{F}_{\mathrm{rad}}\) through Eq.~\eqref{eq:ch2.rad-null-on-evanescent} even when
local \(|\E|\) is substantial \cite{stratton1941}. Regime gate
Eq.~\eqref{eq:ch2.regime-component-gate} remains mandatory: labeling a near-zone state
``non-radiating'' from reactive dominance alone does not establish
\(\mathcal{R}_\omega[\Jimp]=\mathbf{0}\) without export qualification
\cite{harrington2001,stratton1941}.

\paragraph{Spectral and geometric composition.}
Null-space analysis composes with Section~\ref{sec:ch26} outputs:
\begin{equation}
\label{eq:ch2.nullspace-spectral-composition}
\Jimp\in\ker\mathcal{R}_\omega
\ \Longrightarrow\
a_n(\Jimp)=0\ \text{for all exporting modes with}\ \lambda_n\neq 0,
\end{equation}
while geometric constraints Eq.~\eqref{eq:ch2.accessible-modal-sector} may force additional
\(a_n=0\) for \(\Jimp\notin\ker\mathcal{R}_\omega\) \cite{stratton1941,harrington2001}.
Equation~\eqref{eq:ch2.nullspace-spectral-composition} follows from
Eq.~\eqref{eq:ch2.spectral-far-pattern-expansion} and linearity
Eq.~\eqref{eq:ch2.radiation-linearity}; it is not re-derived in Section~\ref{sec:ch27}
\cite{stratton1941}. Inverse problems must budget \(\ker\mathcal{R}_\omega\),
\(\ker\mathcal{O}_{\mathrm{meas}}\), and forbidden sectors from
Theorem~\ref{thm:ch2.geometric-spectral-constraints} separately
\cite{harrington2001}.

Several constructions deliberately remain outside Section~\ref{sec:ch27}. Reciprocity
identities and modal orthogonality belong to Section~\ref{sec:ch28}; plane-wave angular-spectrum
operators and near-to-far transforms to Section~\ref{sec:ch29}; vector spherical harmonic
completeness and explicit \(\mathbf{f}_n\) to Chapter~\ref{ch:03}; realizability of
unreachable modes to Chapter~\ref{ch:04} \cite{stratton1941,harrington2001}. Section~\ref{sec:ch27}
does not re-derive spectral decomposition Theorem~\ref{thm:ch2.radiation-spectral-decomposition},
Fredholm structure Theorem~\ref{thm:ch2.radiation-fredholm-structure}, or geometric constraints
Theorem~\ref{thm:ch2.geometric-spectral-constraints}; it classifies invisible sources on the
maps those results presuppose \cite{stratton1941}. Worked FIG+PROB non-radiating examples appear
in Subsection~\ref{sec:ch2102} by reference to Subsection~\ref{sec:ch272}
\cite{harrington2001}.

Regime logic from Section~\ref{sec:ch01.2.4} is unchanged. Null-space membership does not
override transport qualification: a current in \(\ker\mathcal{R}_\omega\) has
\(\mathcal{P}_{\mathrm{rad}}=0\) on qualified shells, while reactive near fields from
\(\mathcal{S}_\omega[\Jimp]\) may still violate naive flux interpretations if
Eq.~\eqref{eq:ch2.regime-operator-screen} is ignored \cite{stratton1941,harrington2001}.
Far-zone dominance Proposition~\ref{prop:ch2.far-zone-radiative-dominance} applies only to
exporting sources; it does not trivialize \(\ker\mathcal{R}_\omega\) \cite{stratton1941}.
Compression error Eq.~\eqref{eq:ch2.compression-error-bound} bounds modal truncation on
\(\operatorname{ran}\mathcal{R}_\omega\) independently of null-space dimension
\cite{harrington2001}.

\paragraph{Worked illustration (invisible loop mode).}
A closed admissible loop carrying oppositely directed arms can cancel far-zone transverse
radiation through Eq.~\eqref{eq:ch2.radiation-integral-outgoing-upgrade} while maintaining
nonzero \(\mathcal{S}_\omega[\Jimp]\) near the loop \cite{stratton1941,jackson1999}.
Subsection~\ref{sec:ch272} classifies such families; here the point is that
Eq.~\eqref{eq:ch2.null-radiation-preview} is realizable with
\(\mathcal{D}_{J}[\Jimp]=1\), not merely a formal kernel direction
\cite{harrington2001}. Pattern-matching algorithms that fit only
\(\operatorname{ran}\mathcal{R}_\omega\) assign spurious modal labels to invisible loop content
\cite{stratton1941}.

\paragraph{Worked illustration (modal zero versus kernel).}
A single \(\lambda_n=0\) eigenmode from Eq.~\eqref{eq:ch2.radiation-eigenpair-def} contributes
to \(\ker\mathcal{R}_\omega\) when realized by admissible \(\Jimp^{(n)}\)
\cite{stratton1941,harrington2001}. On exterior problems the modal index set is not finite;
truncating to the first \(N\) modes underestimates \(\dim\ker\mathcal{R}_\omega\) and
overstates identifiability of inverse Eq.~\eqref{eq:ch2.inverse-source-equation}
\cite{harrington2001,jackson1999}. Subsection~\ref{sec:ch271} states the correct operator
bookkeeping \cite{stratton1941}.

\paragraph{Worked illustration (evanescent storage).}
A source exciting only evanescent sheets satisfies
\(\Pi_{\mathrm{rad}}\mathcal{S}_\omega[\Jimp]=\mathbf{0}\) through
Eq.~\eqref{eq:ch2.rad-null-on-evanescent} even when \(|\E|\) is large within a wavelength of
\(\Omega_s\) \cite{stratton1941,harrington2001}. Subsection~\ref{sec:ch273} distinguishes this
case from partial nullity on \(\operatorname{ran}\mathcal{R}_\omega\) due to finite aperture
Eq.~\eqref{eq:ch2.observation-null-space} \cite{stratton1941}. Near-zone reactive dominance
without Eq.~\eqref{eq:ch2.regime-component-gate} must not be reported as angular-momentum
transport (Section~\ref{sec:ch01.2.3}) \cite{stratton1941}.

\paragraph{Worked illustration (symmetry-forbidden versus NR).}
Geometric annihilation Eq.~\eqref{eq:ch2.geometric-modal-annihilation} forces selected
\(a_n=0\) for sources that \emph{do} radiate in other sectors
\cite{stratton1941,harrington2001}. Membership in \(\ker\mathcal{R}_\omega\) requires
\(a_n=0\) for \emph{all} exporting modes. A ground-plane symmetric dipole may suppress odd
parity sectors through Eq.~\eqref{eq:ch2.parity-selection-rule} while retaining even-parity
export \cite{jackson1999,stratton1941}. Confusing symmetry nulls with
Eq.~\eqref{eq:ch2.ker-R-omega-preview} misattributes identifiability deficits
\cite{harrington2001}.

The subsection sequence begins with null spaces of radiation operators on declared domains,
continues with canonical non-radiating current configurations, develops evanescent subspaces
and non-reconstructability limits importing Subsection~\ref{sec:ch244} rank templates, and
closes with consequences for angular spectra before reciprocity in Section~\ref{sec:ch28} and
angular-spectrum action in Section~\ref{sec:ch29} compose additional structure on the same
source coordinates \cite{stratton1941,harrington2001,jackson1999}.

\subsection{Null spaces of radiation operators}
\label{sec:ch271}
% TOC tag: [HOME]

Section~\ref{sec:ch27} previewed source-side invisibility through
Eqs.~\eqref{eq:ch2.section27-nullspace-chain}--\eqref{eq:ch2.radiation-nullspace-structure-preview}.
Subsection~\ref{sec:ch271} gives the Chapter~\ref{ch:02} HOME for structural properties of
\(\ker\mathcal{R}_\omega\), adjoint null spaces, cokernel pairings, and modal null sectors,
upgrading preview Eq.~\eqref{eq:ch2.ker-R-omega-preview} from Subsection~\ref{sec:ch244} to
constructive operator-theoretic statements compatible with Fredholm structure Theorem
~\ref{thm:ch2.radiation-fredholm-structure} \cite{stratton1941,harrington2001}. Radiation map
Eq.~\eqref{eq:ch2.nullspace-operator-targets}, spectral decomposition Theorem
~\ref{thm:ch2.radiation-spectral-decomposition}, solvability
Eq.~\eqref{eq:ch2.radiation-solvability-criterion}, and modal reduction
Eq.~\eqref{eq:ch2.modal-fredholm-reduction} from Subsections~\ref{sec:ch261} and~\ref{sec:ch263}
are presupposed; Fredholm proofs and compactness lemmas from Subsection~\ref{sec:ch213} are not
repeated \cite{stratton1941}. Canonical non-radiating impressed currents belong to Subsection
~\ref{sec:ch272}; evanescent non-reconstructability to Subsection~\ref{sec:ch273}
\cite{harrington2001}.

\paragraph{Assumptions.}
Fix \(\omega>0\), admissible source space \(\mathcal{J}_\Omega\)
Eq.~\eqref{eq:ch2.source-space-refined}, charge continuity and boundary compatibility
Eq.~\eqref{eq:ch2.source-admissibility}, outgoing closure on exterior problems
Eqs.~\eqref{eq:ch2.sommerfeld-outgoing-class}--\eqref{eq:ch2.outgoing-radiation-refinement},
and radiation map \(\mathcal{R}_\omega=\Pi_{\mathrm{rad}}\circ\mathcal{S}_\omega\) from
Eqs.~\eqref{eq:ch2.radiation-operator-def}--\eqref{eq:ch2.outgoing-radiation-refinement}
\cite{stratton1941,harrington2001}. Dual pairings on sources and fields are
Eqs.~\eqref{eq:ch2.dual-source-pairing} and~\eqref{eq:ch2.maxwell-adjoint-pairing}
\cite{stratton1941}. Regime gate Eq.~\eqref{eq:ch2.regime-operator-screen} governs transport
reports on any recovered radiative component \cite{harrington2001}.

\paragraph{Radiation null space (HOME definition).}
The \emph{radiation null space} is the source-side kernel
\begin{equation}
\label{eq:ch2.radiation-nullspace-def}
\ker\mathcal{R}_\omega
=
\left\{
\Jimp\in\mathcal{J}_\Omega:
\mathcal{R}_\omega[\Jimp]=\mathbf{0}\in\mathcal{F}_{\mathrm{rad}},
\ \mathcal{D}_{J}[\Jimp]=1
\right\},
\end{equation}
equivalently \(\Pi_{\mathrm{rad}}\,\mathcal{S}_\omega[\Jimp]=\mathbf{0}\) on the outgoing
branch \cite{stratton1941,harrington2001}. Equation~\eqref{eq:ch2.radiation-nullspace-def}
is the Subsection~\ref{sec:ch271} HOME upgrade of preview
Eq.~\eqref{eq:ch2.ker-R-omega-preview}: admissibility is explicit, and membership is
defined on \(\mathcal{J}_\Omega\) rather than on formal test directions
\cite{stratton1941}. Elements of \(\ker\mathcal{R}_\omega\) produce no qualified far-zone
export through Eq.~\eqref{eq:ch2.radiation-integral-outgoing-upgrade} yet may satisfy
\(\mathcal{S}_\omega[\Jimp]\neq\mathbf{0}\) Eq.~\eqref{eq:ch2.nr-image-of-null-space}
\cite{harrington2001,jackson1999}. Linearity
Eq.~\eqref{eq:ch2.radiation-linearity} makes \(\ker\mathcal{R}_\omega\) a linear subspace of
\(\mathcal{J}_\Omega\) closed under admissible superposition \cite{stratton1941}.

\paragraph{Adjoint null space and cokernel.}
Let \(\mathcal{R}_\omega^{\dagger}:\mathcal{F}_{\mathrm{rad}}\to\mathcal{J}_\Omega\) denote
the adjoint induced by Eqs.~\eqref{eq:ch2.dual-source-pairing} and
~\eqref{eq:ch2.maxwell-adjoint-pairing}. Define
\begin{equation}
\label{eq:ch2.radiation-adjoint-nullspace}
\ker\mathcal{R}_\omega^{\dagger}
=
\left\{
\mathbf{w}\in\mathcal{F}_{\mathrm{rad}}:
\langle \mathbf{w},\,\mathcal{R}_\omega[\Jimp]\rangle_{\mathcal{U}}=0
\ \ \forall\ \Jimp\in\mathcal{J}_\Omega
\right\},
\end{equation}
and the cokernel quotient
\begin{equation}
\label{eq:ch2.radiation-cokernel-def}
\mathrm{coker}\,\mathcal{R}_\omega
=
\mathcal{F}_{\mathrm{rad}}\big/\operatorname{ran}\mathcal{R}_\omega,
\end{equation}
with \(\operatorname{ran}\mathcal{R}_\omega\subset\mathcal{F}_{\mathrm{rad}}\) the radiative
image \cite{stratton1941,harrington2001}. Equations~\eqref{eq:ch2.radiation-adjoint-nullspace}--\eqref{eq:ch2.radiation-cokernel-def}
pair with \(\ker\mathcal{R}_\omega\) through Fredholm solvability
Eq.~\eqref{eq:ch2.radiation-solvability-criterion} imported from Subsection~\ref{sec:ch263}:
targets \(\mathbf{g}\in\mathcal{F}_{\mathrm{rad}}\) orthogonal to
\(\ker\mathcal{R}_\omega^{\dagger}\) are reachable, while obstructions lie in cokernel
directions \cite{stratton1941}. Biorthogonality Eq.~\eqref{eq:ch2.radiation-biorthogonality}
expresses the same pairing on modal coordinates without re-deriving adjoint modes here
\cite{harrington2001}.

\paragraph{Fredholm null-space alternative.}
When \(\mathcal{R}_\omega\) is Fredholm in the sense of Theorem
~\ref{thm:ch2.radiation-fredholm-structure},
\begin{equation}
\label{eq:ch2.fredholm-nullspace-alternative}
\dim\ker\mathcal{R}_\omega<\infty
\ \Longleftrightarrow\
\dim\ker\mathcal{R}_\omega^{\dagger}<\infty,
\end{equation}
and index Eq.~\eqref{eq:ch2.radiation-fredholm-index-def} counts the imbalance
\cite{stratton1941,harrington2001}. Equation~\eqref{eq:ch2.fredholm-nullspace-alternative}
does not re-prove Theorem~\ref{thm:ch2.radiation-fredholm-structure}; it records the
Chapter~\ref{ch:02} null-space bookkeeping used in inversion
Eq.~\eqref{eq:ch2.inverse-source-equation} \cite{stratton1941}. On exterior domains with
compact \(\Omega_s\), \(\dim\ker\mathcal{R}_\omega\) is typically infinite even when
\(\mathcal{R}_\omega\) is Fredholm on truncated subspaces
Eq.~\eqref{eq:ch2.spectral-compression-operator} \cite{harrington2001,jackson1999}.

\paragraph{Modal null sector.}
On the orthonormal basis of Theorem~\ref{thm:ch2.radiation-spectral-decomposition}, define
the \emph{modal null index set}
\begin{equation}
\label{eq:ch2.modal-null-index-set}
\mathcal{N}_{\mathrm{null}}
=
\left\{
n\in\mathcal{N}_{\mathrm{rad}}:
\lambda_n=0
\right\},
\end{equation}
with \(\lambda_n\) from Eq.~\eqref{eq:ch2.radiation-eigenpair-def}
\cite{stratton1941,harrington2001}. Equation~\eqref{eq:ch2.modal-null-index-set} is the
spectral slice of \(\ker\mathcal{R}_\omega\) on declared eigenmodes:
\begin{equation}
\label{eq:ch2.modal-null-sector-reduction}
\sum_{n\in\mathcal{N}_{\mathrm{null}}}
c_n[\mathfrak{D}]\,\Jimp^{(n)}
\in\ker\mathcal{R}_\omega
\quad\text{when admissible superpositions exist},
\end{equation}
using source coefficients Eq.~\eqref{eq:ch2.source-modal-amplitude} and eigenpairs
Eq.~\eqref{eq:ch2.radiation-eigenpair-def} \cite{stratton1941}. Equation
\eqref{eq:ch2.modal-null-sector-reduction} does not exhaust \(\ker\mathcal{R}_\omega\) on
exterior problems because \(\mathcal{N}_{\mathrm{rad}}\) need not span all invisible source
degrees \cite{harrington2001}. Proposition~\ref{prop:ch2.modal-invertibility-equivalence}
already states that coefficients on \(\ker\mathcal{R}_\omega\) remain free in
Eq.~\eqref{eq:ch2.modal-particular-solution}; Subsection~\ref{sec:ch271} identifies the
modal-labeled portion \cite{stratton1941}.

\begin{figure}[t]
\centering
\ChOneFigBlock{%
\begin{tikzpicture}[ch1fig, node distance=7mm and 10mm]
  \node[ch1box] (j) {$\mathcal{J}_\Omega$};
  \node[ch1box, right=12mm of j] (r) {$\mathcal{R}_\omega$};
  \node[ch1box, right=12mm of r] (f) {$\mathcal{F}_{\mathrm{rad}}$};
  \node[ch1box, below=11mm of j] (ker) {$\ker\mathcal{R}_\omega$};
  \node[ch1box, below=11mm of f] (adj) {$\ker\mathcal{R}_\omega^{\dagger}$};
  \draw[ch1flow] (j) -- (r) -- (f);
  \draw[ch1flow, dashed] (ker) -- (j);
  \draw[ch1flow, dashed] (adj) -- (f);
\end{tikzpicture}%
}
\caption{Subsection~\ref{sec:ch271} null-space pipeline: source kernel
Eq.~\eqref{eq:ch2.radiation-nullspace-def}, radiative image, and adjoint kernel
Eq.~\eqref{eq:ch2.radiation-adjoint-nullspace} paired by solvability
Eq.~\eqref{eq:ch2.radiation-solvability-criterion}.}
\label{fig:ch2.radiation-nullspace-pipeline}
\end{figure}

Figure~\ref{fig:ch2.radiation-nullspace-pipeline} separates kernel and cokernel data before
observation composition Eq.~\eqref{eq:ch2.observation-fredholm-composition}
\cite{stratton1941,harrington2001}.

\begin{theorem}[Null-space structure of radiation operators]
\label{thm:ch2.radiation-nullspace-structure}
Let \(\mathcal{R}_\omega\) satisfy Eq.~\eqref{eq:ch2.radiation-nullspace-def} on admissible
\(\mathcal{J}_\Omega\) with outgoing closure. Then:
\begin{enumerate}
\item[(i)] \(\ker\mathcal{R}_\omega\) is a linear subspace of \(\mathcal{J}_\Omega\) and
\(\mathcal{R}_\omega[\Jimp]=\mathbf{0}\) iff \(\Jimp\in\ker\mathcal{R}_\omega\).
\item[(ii)] Solvability of \(\mathcal{R}_\omega[\Jimp]=\mathbf{g}\) is equivalent to
Eq.~\eqref{eq:ch2.radiation-solvability-criterion} on declared pairings.
\item[(iii)] On the spectral basis of Subsection~\ref{sec:ch261},
\(\mathcal{N}_{\mathrm{null}}\) Eq.~\eqref{eq:ch2.modal-null-index-set} labels modes with
\(P_n\mathcal{R}_\omega=\mathbf{0}\); modal superpositions
Eq.~\eqref{eq:ch2.modal-null-sector-reduction} lie in \(\ker\mathcal{R}_\omega\) when
admissible.
\item[(iv)] Domain class fixes dimensionality:
\begin{equation}
\label{eq:ch2.nullspace-dimensional-classification}
\begin{aligned}
\text{bounded cavity + discrete book:}\ &
\dim\ker\mathcal{R}_\omega\ \text{finite on truncated}\ \mathcal{R}_\omega^{(N)};\\
\text{exterior + compact}\ \Omega_s:\ &
\dim\ker\mathcal{R}_\omega=\infty\ \text{typ. on full}\ \mathcal{J}_\Omega.
\end{aligned}
\end{equation}
\end{enumerate}
\end{theorem}
\begin{proof}
Part (i) follows from linearity Eq.~\eqref{eq:ch2.radiation-linearity} and definition
Eq.~\eqref{eq:ch2.radiation-nullspace-def} \cite{stratton1941}. Part (ii) is Theorem
~\ref{thm:ch2.radiation-fredholm-structure} and
Eq.~\eqref{eq:ch2.radiation-solvability-criterion}, cited without re-proof
\cite{harrington2001}. Part (iii) uses Eq.~\eqref{eq:ch2.modal-fredholm-reduction} and
\(\lambda_n=0\) in Eq.~\eqref{eq:ch2.radiation-eigenpair-def} \cite{stratton1941}. Part (iv)
contrasts discrete cavity spectra Eq.~\eqref{eq:ch2.cavity-discrete-radiation-spectrum} with
continuous exterior angular content Eq.~\eqref{eq:ch2.exterior-continuous-spectral-integral}
\cite{jackson1999,harrington2001}.
\end{proof}

\begin{proposition}[Kernel image on field space]
\label{prop:ch2.null-kernel-field-image}
If \(\Jimp\in\ker\mathcal{R}_\omega\), then radiative and non-radiative components satisfy
\begin{equation}
\label{eq:ch2.kernel-field-image-law}
\mathbf{u}_{\mathrm{rad}}=\mathcal{R}_\omega[\Jimp]=\mathbf{0},
\qquad
\mathbf{u}_{\mathrm{nr}}=\mathcal{S}_\omega[\Jimp]
=\big(\mathrm{id}-\Pi_{\mathrm{rad}}\big)\mathcal{S}_\omega[\Jimp],
\end{equation}
with \(\mathbf{u}_{\mathrm{rad}}\), \(\mathbf{u}_{\mathrm{nr}}\) from
Eqs.~\eqref{eq:ch2.radiative-component-def}--\eqref{eq:ch2.nonradiative-component-def}
\cite{stratton1941,harrington2001}. Conversely, \(\mathbf{u}_{\mathrm{rad}}=\mathbf{0}\) does
not imply \(\Jimp\in\ker\mathcal{R}_\omega\) unless \(\mathcal{S}_\omega[\Jimp]\) lies entirely
in non-export subspaces annihilated by \(\Pi_{\mathrm{rad}}\) on the declared branch
\cite{stratton1941}.
\end{proposition}
\begin{proof}
Substitute \(\Jimp\in\ker\mathcal{R}_\omega\) into
Eqs.~\eqref{eq:ch2.radiative-component-def}--\eqref{eq:ch2.nonradiative-component-def}
\cite{stratton1941}. The converse failure occurs when reactive near fields are misclassified
without Eq.~\eqref{eq:ch2.regime-component-gate} \cite{harrington2001}.
\end{proof}

\begin{proposition}[Measurement null-space extension]
\label{prop:ch2.measurement-nullspace-extension}
For observation map \(\mathcal{M}_\omega=\mathcal{O}_{\mathrm{meas}}\circ\Lambda_{\mathrm{rad}}\)
Eq.~\eqref{eq:ch2.reconstruction-observation-map},
\begin{equation}
\label{eq:ch2.measurement-nullspace-composition}
\ker\mathcal{M}_\omega
\supseteq
\ker\mathcal{R}_\omega,
\qquad
\ker\mathcal{M}_\omega
=
\ker\mathcal{R}_\omega
\oplus
\left\{
\Jimp\in\mathcal{J}_\Omega:
\mathcal{R}_\omega[\Jimp]\in\ker\mathcal{O}_{\mathrm{meas}}
\right\},
\end{equation}
when the direct sum is taken on the range side after \(\mathcal{R}_\omega\)
\cite{stratton1941,harrington2001}. Equation~\eqref{eq:ch2.measurement-nullspace-composition}
extends preview Eq.~\eqref{eq:ch2.evanescent-nonreconstruct-preview}: radiation nullity implies
measurement nullity, but finite aperture Eq.~\eqref{eq:ch2.observation-null-space} adds
field-side invisible directions even when \(\Jimp\notin\ker\mathcal{R}_\omega\)
\cite{stratton1941}. Rank cap Eq.~\eqref{eq:ch2.reconstruction-rank-cap} bounds
\(\dim\operatorname{ran}\mathcal{M}_\omega\), not \(\dim\ker\mathcal{M}_\omega\)
\cite{harrington2001}.
\end{proposition}
\begin{proof}
If \(\mathcal{R}_\omega[\Jimp]=\mathbf{0}\), then \(\mathcal{M}_\omega[\Jimp]=\mathbf{0}\).
If \(\mathcal{R}_\omega[\Jimp]\in\ker\mathcal{O}_{\mathrm{meas}}\), then
\(\mathcal{M}_\omega[\Jimp]=\mathbf{0}\) with possibly nonzero export
\cite{stratton1941,harrington2001}. Linearity gives the set sum decomposition on the
declared domain \cite{stratton1941}.
\end{proof}

\paragraph{Composition with spectral and geometric filters.}
Null-space membership implies modal annihilation on exporting sectors:
\begin{equation}
\label{eq:ch2.nullspace-modal-annihilation-law}
\Jimp\in\ker\mathcal{R}_\omega
\ \Longrightarrow\
a_n(\Jimp)=0
\ \text{for all}\ n\ \text{with}\ \lambda_n\neq 0,
\end{equation}
with far coefficients \(a_n\) from Eq.~\eqref{eq:ch2.spectral-far-pattern-expansion}
\cite{stratton1941,harrington2001}. Equation~\eqref{eq:ch2.nullspace-modal-annihilation-law}
restates Eq.~\eqref{eq:ch2.nullspace-spectral-composition} from Section~\ref{sec:ch27} without
re-deriving the modal expansion \cite{stratton1941}. Geometric constraints
Theorem~\ref{thm:ch2.geometric-spectral-constraints} may force additional \(a_n=0\) for sources
outside \(\ker\mathcal{R}_\omega\) \cite{harrington2001}. Composition law
Eq.~\eqref{eq:ch2.constraint-composition-law} budgets identifiability deficits separately
\cite{stratton1941}.

\paragraph{Index stability and regularization.}
Fredholm index Eq.~\eqref{eq:ch2.radiation-fredholm-index-def} is unchanged under compact
perturbations Eq.~\eqref{eq:ch2.radiation-fredholm-decomposition} and under Tikhonov
regularization Eq.~\eqref{eq:ch2.tikhonov-source-reconstruction} on the range side
\cite{stratton1941,harrington2001}. Regularization stabilizes
Eq.~\eqref{eq:ch2.modal-particular-solution} for small \(|\lambda_n|\) but leaves
\(\ker\mathcal{R}_\omega\) unchanged: invisible source components cannot be recovered from
far-zone data at any regularization level \cite{stratton1941}. Compression truncation
Eq.~\eqref{eq:ch2.compression-error-bound} approximates \(\operatorname{ran}\mathcal{R}_\omega\),
not \(\ker\mathcal{R}_\omega\) \cite{harrington2001}.

\paragraph{Boundary of validity.}
Theorem~\ref{thm:ch2.radiation-nullspace-structure} assumes linear media at fixed \(\omega\),
outgoing closure, and admissible \(\mathcal{J}_\Omega\) Eq.~\eqref{eq:ch2.source-admissibility}
\cite{stratton1941,harrington2001}. Loss modifies \(\lambda_n\) and may merge modal null sectors
without promoting evanescent content to export through
Eq.~\eqref{eq:ch2.rad-null-on-evanescent} \cite{stratton1941}. Nonlinear or time-varying
sources require the causal chain of Subsection~\ref{sec:ch223}; harmonic null spaces here are
fixed-frequency slices \cite{harrington2001}. Regime gate
Eq.~\eqref{eq:ch2.regime-component-gate} remains mandatory before interpreting near-zone
reactive dominance as null-space membership \cite{stratton1941}.

\paragraph{Worked illustration (cavity truncation versus true kernel).}
On a bounded cavity, \(\mathcal{R}_\omega^{(N)}\) from
Eq.~\eqref{eq:ch2.spectral-compression-operator} has finite-dimensional kernel spanned by
modes with \(\lambda_n=0\) in the truncated book \cite{stratton1941,harrington2001}. The full
exterior problem with the same physical source support has
\(\dim\ker\mathcal{R}_\omega=\infty\) per
Eq.~\eqref{eq:ch2.nullspace-dimensional-classification} \cite{jackson1999,harrington2001}. Identifiability
claims based on cavity simulators therefore underestimate invisible degrees unless domain class
is matched \cite{stratton1941}.

\paragraph{Worked illustration (cokernel obstruction).}
A smooth target pattern \(\mathbf{g}\in\mathcal{F}_{\mathrm{rad}}\) may fail
Eq.~\eqref{eq:ch2.radiation-solvability-criterion} when \(\mathbf{g}\) has nonzero pairing
with \(\mathbf{w}\in\ker\mathcal{R}_\omega^{\dagger}\) even though \(\|\mathbf{g}\|\) is
small in energy norm \cite{stratton1941,harrington2001}. Subsection~\ref{sec:ch263} computed
the obstruction; here the point is that cokernel data is independent of
\(\ker\mathcal{R}_\omega\) dimension \cite{stratton1941}. Adding a non-radiating source
\(\Jimp^{\mathrm{(NR)}}\in\ker\mathcal{R}_\omega\) does not repair an unsolvable \(\mathbf{g}\)
\cite{harrington2001}.

\paragraph{Worked illustration (finite array null extension).}
When \(\mathcal{R}_\omega[\Jimp]\neq\mathbf{0}\) but lies in \(\ker\mathcal{O}_{\mathrm{meas}}\),
Eq.~\eqref{eq:ch2.measurement-nullspace-composition} gives
\(\mathcal{M}_\omega[\Jimp]=\mathbf{0}\) \cite{stratton1941,harrington2001}. This is distinct
from \(\Jimp\in\ker\mathcal{R}_\omega\): export exists but the declared aperture cannot resolve
it Eq.~\eqref{eq:ch2.observation-null-space} \cite{stratton1941}. Subsection~\ref{sec:ch273}
develops reconstructability limits; Subsection~\ref{sec:ch271} fixes the operator inclusion
\(\ker\mathcal{R}_\omega\subset\ker\mathcal{M}_\omega\) \cite{harrington2001}.

\paragraph{Worked illustration (modal zero slice).}
If \(\lambda_{n_0}=0\) for a declared eigenpair, then \(\Jimp^{(n_0)}\) from
Eq.~\eqref{eq:ch2.radiation-eigenpair-def} contributes to
Eq.~\eqref{eq:ch2.modal-null-sector-reduction} \cite{stratton1941,harrington2001}. Superposing
such modes with an exporting source leaves far-zone patterns unchanged while altering near-zone
\(\mathbf{u}_{\mathrm{nr}}\) through Eq.~\eqref{eq:ch2.kernel-field-image-law}
\cite{stratton1941}. Subsection~\ref{sec:ch272} constructs explicit \(\Jimp^{\mathrm{(NR)}}\)
families beyond single eigenmodes \cite{harrington2001}.

\paragraph{Closing summary.}
Subsection~\ref{sec:ch271} establishes the Chapter~\ref{ch:02} HOME for radiation null space
Eq.~\eqref{eq:ch2.radiation-nullspace-def}, adjoint kernel
Eq.~\eqref{eq:ch2.radiation-adjoint-nullspace}, cokernel
Eq.~\eqref{eq:ch2.radiation-cokernel-def}, Fredholm alternative
Eq.~\eqref{eq:ch2.fredholm-nullspace-alternative}, modal null sector
Eqs.~\eqref{eq:ch2.modal-null-index-set}--\eqref{eq:ch2.modal-null-sector-reduction},
Theorem~\ref{thm:ch2.radiation-nullspace-structure} with dimension classification
Eq.~\eqref{eq:ch2.nullspace-dimensional-classification}, field image
Proposition~\ref{prop:ch2.null-kernel-field-image}, measurement extension
Proposition~\ref{prop:ch2.measurement-nullspace-extension}, and modal annihilation
Eq.~\eqref{eq:ch2.nullspace-modal-annihilation-law}, upgrading preview
Eq.~\eqref{eq:ch2.ker-R-omega-preview} and importing Fredholm structure from
Subsection~\ref{sec:ch263} by reference \cite{stratton1941,harrington2001,jackson1999}.
Subsection~\ref{sec:ch272} develops canonical non-radiating current configurations; Subsection
~\ref{sec:ch273} analyzes evanescent subspaces and non-reconstructability
\cite{stratton1941}.
\subsection{Non-radiating current configurations}
\label{sec:ch272}
% TOC tag: [HOME]

Subsection~\ref{sec:ch271} fixed operator null-space structure through
Eq.~\eqref{eq:ch2.radiation-nullspace-def}, Theorem~\ref{thm:ch2.radiation-nullspace-structure},
and modal null sector Eqs.~\eqref{eq:ch2.modal-null-index-set}--\eqref{eq:ch2.modal-null-sector-reduction}
\cite{stratton1941,harrington2001}. Subsection~\ref{sec:ch272} gives the Chapter~\ref{ch:02}
HOME for \emph{constructive} non-radiating impressed-current configurations---admissible
\(\Jimp^{\mathrm{(NR)}}\in\ker\mathcal{R}_\omega\) with \(\mathcal{S}_\omega[\Jimp^{\mathrm{(NR)}}]\neq\mathbf{0}\)---
upgrading preview Eq.~\eqref{eq:ch2.nr-current-template-preview} from Section~\ref{sec:ch27}
\cite{stratton1941,harrington2001,jackson1999}. Radiation integrals, Huygens surface forms
Eqs.~\eqref{eq:ch2.huygens-equivalent-integral}--\eqref{eq:ch2.huygens-tangential-match}, and
equivalent-source ambiguity Eq.~\eqref{eq:ch2.equivalent-source-ambiguity} from
Subsections~\ref{sec:ch251} and~\ref{sec:ch244} are presupposed; dyadic kernel tabulations from
Subsection~\ref{sec:ch231} are not repeated \cite{stratton1941}. Evanescent overlap and
measurement nullity belong to Subsection~\ref{sec:ch273}; worked FIG+PROB examples appear in
Subsection~\ref{sec:ch2102} by reference \cite{harrington2001}.

\paragraph{Assumptions.}
Fix \(\omega>0\), compact support \(\Omega_s\), admissible impressed data
Eq.~\eqref{eq:ch2.source-admissibility}, charge continuity
Eq.~\eqref{eq:ch1.charge-continuity-harmonic}, outgoing closure
Eqs.~\eqref{eq:ch2.outgoing-radiation-refinement}--\eqref{eq:ch2.outgoing-resolvent-unique},
and null space Eq.~\eqref{eq:ch2.radiation-nullspace-def}
\cite{stratton1941,harrington2001}. All configurations below are required to satisfy
\(\mathcal{D}_{J}[\Jimp^{\mathrm{(NR)}}]=1\) and \(\divg\Jimp^{\mathrm{(NR)}}=-\jj\omega\rho^{\mathrm{(NR)}}\)
with \(\rho^{\mathrm{(NR)}}\) compatible with boundary class
Eq.~\eqref{eq:ch1.outer-boundary-classes} \cite{stratton1941}. Regime gate
Eq.~\eqref{eq:ch2.regime-component-gate} governs whether near-zone reactive content may be
reported as transport data \cite{harrington2001}.

\paragraph{Non-radiating current (HOME definition).}
A \emph{non-radiating current configuration} is an admissible impressed current
\begin{equation}
\label{eq:ch2.nr-current-def}
\Jimp^{\mathrm{(NR)}}\in\ker\mathcal{R}_\omega,
\qquad
\mathcal{S}_\omega[\Jimp^{\mathrm{(NR)}}]\neq\mathbf{0}\ \text{in general},
\end{equation}
with radiative component \(\mathbf{u}_{\mathrm{rad}}=\mathbf{0}\) and non-radiative component
\(\mathbf{u}_{\mathrm{nr}}=\mathcal{S}_\omega[\Jimp^{\mathrm{(NR)}}]\) from
Eqs.~\eqref{eq:ch2.radiative-component-def}--\eqref{eq:ch2.nonradiative-component-def}
\cite{stratton1941,harrington2001}. Equation~\eqref{eq:ch2.nr-current-def} is the Subsection
~\ref{sec:ch272} HOME upgrade of preview Eq.~\eqref{eq:ch2.nr-current-template-preview} and
restates Eq.~\eqref{eq:ch2.null-radiation-preview} with explicit admissibility: far-zone export
vanishes through Eq.~\eqref{eq:ch2.radiation-integral-outgoing-upgrade}, while near-zone fields
may remain substantial Proposition~\ref{prop:ch2.null-kernel-field-image}
\cite{stratton1941,jackson1999}. Linearity Eq.~\eqref{eq:ch2.radiation-linearity} implies
admissible superpositions of non-radiating configurations remain non-radiating when charge
continuity is preserved \cite{stratton1941}.

\paragraph{Admissibility constraints.}
Non-radiating configurations are not arbitrary curl fields. Impressed data must satisfy
\begin{equation}
\label{eq:ch2.nr-current-admissibility-law}
\Jimp^{\mathrm{(NR)}}\in L^{2}(\Omega_s)^{3},
\quad
\divg\Jimp^{\mathrm{(NR)}}=-\jj\omega\rho^{\mathrm{(NR)}},
\quad
\Jimp^{\mathrm{(NR)}}\ \text{BC-compatible on}\ \Gamma,
\end{equation}
as in Eq.~\eqref{eq:ch2.source-admissibility} \cite{stratton1941,harrington2001}. Equation
\eqref{eq:ch2.nr-current-admissibility-law} excludes discontinuous \(\Jimp^{\mathrm{(NR)}}\) that
violate Eq.~\eqref{eq:ch1.charge-continuity-harmonic} before Maxwell closure is attempted
\cite{stratton1941}. Surface-equivalent realizations on \(\Gamma_{\mathrm{eq}}\) inherit trace
class constraints from Subsection~\ref{sec:ch211} \cite{harrington2001}. Membership in
\(\ker\mathcal{R}_\omega\) is an \emph{export} statement; Eq.~\eqref{eq:ch2.nr-current-admissibility-law}
is a \emph{source} statement---both are required \cite{stratton1941}.

\paragraph{Loop-cancellation family.}
Let \(\mathcal{C}\subset\Omega_s\) be a closed admissible circulation path enclosing negligible
magnetic flux linkage at \(\omega\), with oppositely directed arms \(\mathcal{C}_+\) and
\(\mathcal{C}_-\) carrying equal magnitude currents. When far-zone transverse contributions
from the two arms cancel under outgoing propagation,
\begin{equation}
\label{eq:ch2.nr-loop-cancellation}
\int_{\mathcal{C}_+}\!\mathbf{J}_{\mathrm{loop}}\cdot d\boldsymbol{\ell}
+
\int_{\mathcal{C}_-}\!\mathbf{J}_{\mathrm{loop}}\cdot d\boldsymbol{\ell}
=0
\ \Longrightarrow\
\mathcal{R}_\omega[\Jimp^{(\mathrm{loop})}]=\mathbf{0},
\end{equation}
for the associated volume current \(\Jimp^{(\mathrm{loop})}\) obtained by compact support
regularization on \(\mathcal{C}\) \cite{stratton1941,jackson1999}. Equation
\eqref{eq:ch2.nr-loop-cancellation} is the dipole-moment cancellation law at the source level:
vanishing leading \(1/r\) transverse term in
Eq.~\eqref{eq:ch2.far-field-amplitude-integral} when the net moment vanishes
\cite{jackson1999,harrington2001}. Near-zone reactive fields from \(\mathcal{S}_\omega[\Jimp^{(\mathrm{loop})}]\)
need not vanish; Eq.~\eqref{eq:ch2.regime-component-gate} still applies \cite{stratton1941}.

\paragraph{Counter-propagating pair family.}
Two compact sources \(\Jimp^{(+)}\), \(\Jimp^{(-)}\) with supports \(\Omega_\pm\) separated by
distance \(d\ll\lambda_0\) can form a non-radiating pair when their outgoing far patterns
interfere destructively:
\begin{equation}
\label{eq:ch2.nr-counterprop-pair}
\Jimp^{\mathrm{(NR,cp)}}=\Jimp^{(+)}+\Jimp^{(-)},
\qquad
\E_{0,+}(\theta,\phi)+\E_{0,-}(\theta,\phi)=\mathbf{0}
\ \text{on}\ S^{2},
\end{equation}
with \(\E_{0,\pm}\) far amplitudes from Eq.~\eqref{eq:ch2.far-field-amplitude-integral} for
each component \cite{stratton1941,harrington2001,jackson1999}. Equation
\eqref{eq:ch2.nr-counterprop-pair} realizes Eq.~\eqref{eq:ch2.nr-current-def} when phase and
polarization of \(\Jimp^{(\pm)}\) are tuned so \(\Pi_{\mathrm{rad}}\mathcal{S}_\omega[\Jimp^{(\pm)}]\)
cancel on the outgoing branch \cite{stratton1941}. The pair may store substantial reactive
energy in the overlap region through Eq.~\eqref{eq:ch2.component-power-budget} while
\(\mathcal{P}_{\mathrm{rad}}=0\) on qualified shells \cite{harrington2001}.

\paragraph{Surface tangential null family.}
On a closed equivalent surface \(\Gamma_{\mathrm{eq}}\) enclosing radiating sources, tangential
non-radiating distributions satisfy Eq.~\eqref{eq:ch2.equivalent-source-ambiguity}:
\begin{equation}
\label{eq:ch2.nr-surface-tangential-null}
\mathbf{J}_{s,\mathrm{eq}}^{\mathrm{(NR)}}
\in
\ker\mathcal{R}_\omega\big|_{\Gamma_{\mathrm{eq}}},
\qquad
\mathbf{J}_{s,\mathrm{eq}}^{\mathrm{(NR)}}\cdot\hat{\mathbf{n}}=0
\ \text{on}\ \Gamma_{\mathrm{eq}},
\end{equation}
with \(\hat{\mathbf{n}}\) the outward normal \cite{stratton1941,harrington2001}. Equation
\eqref{eq:ch2.nr-surface-tangential-null} is the Subsection~\ref{sec:ch272} HOME realization of
Huygens equivalence ambiguity: exterior fields on \(\mathcal{F}_{\mathrm{SF}}\) are fixed by
Theorem~\ref{thm:ch2.exterior-radiative-uniqueness}, but equivalent surface currents are unique
only up to additions in Eq.~\eqref{eq:ch2.nr-surface-tangential-null}
\cite{stratton1941}. Integrals Eqs.~\eqref{eq:ch2.huygens-equivalent-integral}--\eqref{eq:ch2.huygens-tangential-match}
realize the same \(\mathcal{S}_\omega\) image for equivalent and non-radiating tangential
perturbations \cite{harrington2001}.

\paragraph{Modal eigenrealization family.}
When radiation eigenpairs with \(\lambda_n=0\) are declared in
Eq.~\eqref{eq:ch2.radiation-eigenpair-def},
\begin{equation}
\label{eq:ch2.nr-modal-eigenrealization}
\Jimp^{\mathrm{(NR,mod)}}
=
\sum_{n\in\mathcal{N}_{\mathrm{null}}}
\alpha_n\,\Jimp^{(n)},
\qquad
\mathcal{R}_\omega[\Jimp^{\mathrm{(NR,mod)}}]=\mathbf{0},
\end{equation}
for coefficients \(\alpha_n\in\mathbb{C}\) such that the superposition satisfies
Eq.~\eqref{eq:ch2.nr-current-admissibility-law} \cite{stratton1941,harrington2001}. Equation
\eqref{eq:ch2.nr-modal-eigenrealization} specializes Eq.~\eqref{eq:ch2.modal-null-sector-reduction}
from Subsection~\ref{sec:ch271} to constructive sources: each \(\Jimp^{(n)}\) with
\(\lambda_n=0\) is a modal non-radiating generator when admissible
\cite{stratton1941}. On exterior domains \(\mathcal{N}_{\mathrm{null}}\) need not span all of
\(\ker\mathcal{R}_\omega\) Eq.~\eqref{eq:ch2.nullspace-dimensional-classification}
\cite{harrington2001,jackson1999}.

\begin{figure}[t]
\centering
\ChOneFigBlock{%
\begin{tikzpicture}[ch1fig, node distance=7mm and 9mm]
  \node[ch1box] (loop) {Loop\\cancellation};
  \node[ch1box, right=11mm of loop] (cp) {Counter-\\prop pair};
  \node[ch1box, right=11mm of cp] (surf) {Surface\\tangential};
  \node[ch1box, right=11mm of surf] (mod) {Modal\\$\lambda_n=0$};
  \node[ch1box, below=12mm of cp] (ker) {$\ker\mathcal{R}_\omega$};
  \node[ch1box, below=12mm of loop] (j) {$\Jimp^{\mathrm{(NR)}}$};
  \draw[ch1flow] (loop) -- (ker);
  \draw[ch1flow] (cp) -- (ker);
  \draw[ch1flow] (surf) -- (ker);
  \draw[ch1flow] (mod) -- (ker);
  \draw[ch1flow] (j) -- (loop);
  \draw[ch1flow] (j) -- (cp);
  \draw[ch1flow] (j) -- (surf);
  \draw[ch1flow] (j) -- (mod);
\end{tikzpicture}%
}
\caption{Subsection~\ref{sec:ch272} canonical non-radiating families: loop cancellation
Eq.~\eqref{eq:ch2.nr-loop-cancellation}, counter-propagating pair
Eq.~\eqref{eq:ch2.nr-counterprop-pair}, surface tangential null
Eq.~\eqref{eq:ch2.nr-surface-tangential-null}, and modal eigenrealization
Eq.~\eqref{eq:ch2.nr-modal-eigenrealization}, all mapping to
\(\ker\mathcal{R}_\omega\).}
\label{fig:ch2.nr-current-configurations}
\end{figure}

Figure~\ref{fig:ch2.nr-current-configurations} orders constructive families before
evanescent and measurement extensions in Subsection~\ref{sec:ch273}
\cite{stratton1941,harrington2001}.

\begin{theorem}[Canonical non-radiating current configurations]
\label{thm:ch2.canonical-nr-current-families}
Let \(\Jimp^{\mathrm{(NR)}}\) satisfy Eq.~\eqref{eq:ch2.nr-current-def} and admissibility
Eq.~\eqref{eq:ch2.nr-current-admissibility-law}. Then \(\Jimp^{\mathrm{(NR)}}\) lies in
\(\ker\mathcal{R}_\omega\) whenever it belongs to one of the canonical families
Eqs.~\eqref{eq:ch2.nr-loop-cancellation}--\eqref{eq:ch2.nr-modal-eigenrealization} on the
outgoing branch \cite{stratton1941,harrington2001,jackson1999}. Admissible superpositions
\begin{equation}
\label{eq:ch2.nr-superposition-law}
\Jimp^{(\mathrm{tot})}
=
\Jimp^{(\mathrm{rad})}+\Jimp^{\mathrm{(NR)}},
\qquad
\mathcal{R}_\omega[\Jimp^{(\mathrm{tot})}]
=
\mathcal{R}_\omega[\Jimp^{(\mathrm{rad})}],
\end{equation}
hold for any radiating \(\Jimp^{(\mathrm{rad})}\in\mathcal{J}_\Omega\) and non-radiating
\(\Jimp^{\mathrm{(NR)}}\in\ker\mathcal{R}_\omega\) with \(\mathcal{D}_J[\Jimp^{(\mathrm{tot})}]=1\)
\cite{stratton1941}. Far-zone patterns, flux meters
Eq.~\eqref{eq:ch2.power-flux-functional}, and modal coefficients \(a_n\) on exporting modes
Eq.~\eqref{eq:ch2.spectral-far-pattern-expansion} are unchanged by addition of
\(\Jimp^{\mathrm{(NR)}}\) \cite{harrington2001}.
\end{theorem}
\begin{proof}
Loop and counter-propagating families cancel leading outgoing transverse terms in
Eq.~\eqref{eq:ch2.far-field-amplitude-integral} under stated symmetry conditions
\cite{jackson1999,stratton1941}. Surface tangential nulls leave exterior
\(\mathcal{S}_\omega\) images unchanged by Eq.~\eqref{eq:ch2.equivalent-source-ambiguity} and
Huygens representation Eqs.~\eqref{eq:ch2.huygens-equivalent-integral}--\eqref{eq:ch2.huygens-tangential-match}
\cite{stratton1941,harrington2001}. Modal eigenrealization follows from
\(\lambda_n=0\) in Eq.~\eqref{eq:ch2.radiation-eigenpair-def} and linearity
Eq.~\eqref{eq:ch2.radiation-linearity} \cite{stratton1941}. Superposition law
Eq.~\eqref{eq:ch2.nr-superposition-law} is linearity of \(\mathcal{R}_\omega\) with
\(\mathcal{R}_\omega[\Jimp^{\mathrm{(NR)}}]=\mathbf{0}\) \cite{harrington2001}.
\end{proof}

\begin{proposition}[Nonzero near fields for non-radiating currents]
\label{prop:ch2.nr-near-field-nonzero}
If \(\Jimp^{\mathrm{(NR)}}\) satisfies Eq.~\eqref{eq:ch2.nr-current-def}, then
\begin{equation}
\label{eq:ch2.nr-near-field-energy-law}
\|\mathbf{u}_{\mathrm{nr}}\|_{\mathcal{E}}
=
\|\mathcal{S}_\omega[\Jimp^{\mathrm{(NR)}}]\|_{\mathcal{E}}
\ \text{may be}\ \gg\ \|\mathbf{u}_{\mathrm{rad}}\|_{\mathcal{E}}=0
\end{equation}
on near shells around \(\Omega_s\), with energy norm from
Eq.~\eqref{eq:ch2.field-energy-norm} \cite{stratton1941,harrington2001}. Reactive power
\(\mathcal{P}_{\mathrm{nr}}(r)\) from Eq.~\eqref{eq:ch2.component-power-budget} may dominate
\(\mathcal{P}_{\mathrm{rad}}=0\) locally while far-zone dominance
Proposition~\ref{prop:ch2.far-zone-radiative-dominance} applies only to exporting superpositions
\cite{stratton1941}. Equation~\eqref{eq:ch2.nr-near-field-energy-law} is the field-side content
of Proposition~\ref{prop:ch2.null-kernel-field-image}, specialized to constructive
\(\Jimp^{\mathrm{(NR)}}\) \cite{harrington2001}.
\end{proposition}
\begin{proof}
Proposition~\ref{prop:ch2.null-kernel-field-image} gives
\(\mathbf{u}_{\mathrm{rad}}=\mathbf{0}\) and
\(\mathbf{u}_{\mathrm{nr}}=\mathcal{S}_\omega[\Jimp^{\mathrm{(NR)}}]\)
\cite{stratton1941}. Near-zone scaling Eq.~\eqref{eq:ch2.near-zone-field-scaling} permits
large reactive content for compact \(\Omega_s\) \cite{harrington2001,jackson1999}.
\end{proof}

\begin{proposition}[Invisibility under modal and geometric reports]
\label{prop:ch2.nr-invisibility-reports}
For \(\Jimp^{\mathrm{(NR)}}\in\ker\mathcal{R}_\omega\),
\begin{equation}
\label{eq:ch2.nr-modal-invisibility}
a_n(\Jimp^{\mathrm{(NR)}})=0
\ \text{for all exporting modes with}\ \lambda_n\neq 0,
\end{equation}
and superposition with any \(\Jimp^{(\mathrm{rad})}\) leaves \(a_n(\Jimp^{(\mathrm{tot})})=a_n(\Jimp^{(\mathrm{rad})})\)
Eq.~\eqref{eq:ch2.nr-superposition-law} \cite{stratton1941,harrington2001}. Geometric
forbidden sectors Eq.~\eqref{eq:ch2.geometric-modal-annihilation} are independent: a source
may radiate in allowed parity sectors while carrying additional
\(\Jimp^{\mathrm{(NR)}}\) invisible to \(\mathcal{R}_\omega\) \cite{stratton1941}. Inverse
source Eq.~\eqref{eq:ch2.inverse-source-equation} cannot resolve
\(\Jimp^{\mathrm{(NR)}}\) from far-zone data alone \cite{harrington2001}.
\end{proposition}
\begin{proof}
Eq.~\eqref{eq:ch2.nr-modal-invisibility} is Eq.~\eqref{eq:ch2.nullspace-modal-annihilation-law}
from Subsection~\ref{sec:ch271} \cite{stratton1941}. Superposition invariance follows from
Theorem~\ref{thm:ch2.canonical-nr-current-families} \cite{harrington2001}.
\end{proof}

\paragraph{Relation to equivalent-source inversion.}
Equivalent-current solvers on \(\Gamma_{\mathrm{eq}}\) seek \(\mathbf{J}_{\mathrm{eq}}\) matching
near or far data \cite{stratton1941,harrington2001}. Theorem
~\ref{thm:ch2.canonical-nr-current-families} and Eq.~\eqref{eq:ch2.nr-surface-tangential-null}
explain why such inversions are ill-posed even when forward fields are unique: any solution may
be shifted by tangential \(\mathbf{J}_{s,\mathrm{eq}}^{\mathrm{(NR)}}\) without changing
\(\mathcal{R}_\omega[\Jimp]\) \cite{stratton1941}. Tikhonov regularization
Eq.~\eqref{eq:ch2.tikhonov-source-reconstruction} selects a particular representative but does
not eliminate the equivalence class \cite{harrington2001}. Subsection~\ref{sec:ch273} separates
this source nullity from finite-rank observation nullity
Eq.~\eqref{eq:ch2.measurement-nullspace-composition} \cite{stratton1941}.

\paragraph{Boundary of validity.}
Theorem~\ref{thm:ch2.canonical-nr-current-families} assumes linear Maxwell theory at fixed
\(\omega\), outgoing closure, and admissible sources Eq.~\eqref{eq:ch2.nr-current-admissibility-law}
\cite{stratton1941,harrington2001}. Lossy media modify loop and pair conditions through complex
\(k(\mathbf{r})\) without promoting evanescent storage to export
Eq.~\eqref{eq:ch2.rad-null-on-evanescent} \cite{stratton1941}. Canonical families generate
\emph{explicit} members of \(\ker\mathcal{R}_\omega\); they do not exhaust
\(\dim\ker\mathcal{R}_\omega=\infty\) on exterior domains
Eq.~\eqref{eq:ch2.nullspace-dimensional-classification} \cite{harrington2001,jackson1999}.
Regime gate Eq.~\eqref{eq:ch2.regime-operator-screen} remains mandatory: large near-zone
\(|\E|\) from \(\mathcal{S}_\omega[\Jimp^{\mathrm{(NR)}}]\) is not export without
\(\Pi_{\mathrm{rad}}\) qualification \cite{stratton1941}.

\paragraph{Worked illustration (magnetic loop pair).}
Two equal opposing current filaments forming a small loop satisfy
Eq.~\eqref{eq:ch2.nr-loop-cancellation} when net dipole moment vanishes
\cite{stratton1941,jackson1999}. Near-zone \(|\H|\) may remain large while
Eq.~\eqref{eq:ch2.far-field-amplitude-integral} yields \(\E_0=\mathbf{0}\)
\cite{harrington2001}. Pattern-matching that fits only dipolar \(\mathbf{f}_n\) cannot detect
the loop's non-radiating character \cite{stratton1941}.

\paragraph{Worked illustration (Huygens surface shift).}
Given a valid equivalent \(\mathbf{J}_{s,\mathrm{eq}}\) on \(\Gamma_{\mathrm{eq}}\), adding
any \(\mathbf{J}_{s,\mathrm{eq}}^{\mathrm{(NR)}}\) from
Eq.~\eqref{eq:ch2.nr-surface-tangential-null} leaves exterior samples of
\(\mathcal{S}_\omega[\Jimp]\) unchanged on \(\mathcal{F}_{\mathrm{SF}}\) while altering
surface power density \cite{stratton1941,harrington2001}. Far-zone inversion therefore returns
non-unique \(\mathbf{J}_{\mathrm{eq}}\) even when \(\mathcal{R}_\omega[\Jimp]\) is fixed
\cite{stratton1941}.

\paragraph{Worked illustration (radiating plus invisible component).}
Let \(\Jimp^{(\mathrm{dip})}\) be a radiating dipole and \(\Jimp^{\mathrm{(NR)}}\) a loop
configuration from Eq.~\eqref{eq:ch2.nr-loop-cancellation}. Then
Eq.~\eqref{eq:ch2.nr-superposition-law} gives identical far patterns for
\(\Jimp^{(\mathrm{dip})}\) and \(\Jimp^{(\mathrm{dip})}+\Jimp^{\mathrm{(NR)}}\) while near-zone
energy budgets differ through Eq.~\eqref{eq:ch2.component-power-budget}
\cite{stratton1941,harrington2001,jackson1999}. Source identification from far data alone cannot
separate the components Proposition~\ref{prop:ch2.nr-invisibility-reports}
\cite{stratton1941}.

\paragraph{Worked illustration (modal null generator).}
When a declared eigenpair has \(\lambda_{n_0}=0\), admissible \(\Jimp^{(n_0)}\) from
Eq.~\eqref{eq:ch2.radiation-eigenpair-def} satisfies
Eq.~\eqref{eq:ch2.nr-modal-eigenrealization} with \(\alpha_{n_0}=1\)
\cite{stratton1941,harrington2001}. Superposing \(\Jimp^{(n_0)}\) with an exporting source
changes \(\mathbf{u}_{\mathrm{nr}}\) but not \(a_n\) for \(\lambda_n\neq 0\)
\cite{stratton1941}. Complete catalogs of \(\Jimp^{(n)}\) on open domains await explicit
harmonic bases in Chapter~\ref{ch:03} \cite{harrington2001}.

\paragraph{Closing summary.}
Subsection~\ref{sec:ch272} establishes the Chapter~\ref{ch:02} HOME for non-radiating current
definition Eq.~\eqref{eq:ch2.nr-current-def}, admissibility
Eq.~\eqref{eq:ch2.nr-current-admissibility-law}, canonical families
Eqs.~\eqref{eq:ch2.nr-loop-cancellation}--\eqref{eq:ch2.nr-modal-eigenrealization},
Theorem~\ref{thm:ch2.canonical-nr-current-families} with superposition law
Eq.~\eqref{eq:ch2.nr-superposition-law}, near-field Proposition
~\ref{prop:ch2.nr-near-field-nonzero}, invisibility Proposition
~\ref{prop:ch2.nr-invisibility-reports}, upgrading preview
Eq.~\eqref{eq:ch2.nr-current-template-preview} and importing null-space structure from
Subsection~\ref{sec:ch271} by reference \cite{stratton1941,harrington2001,jackson1999}.
Subsection~\ref{sec:ch273} develops evanescent subspaces and non-reconstructability; Subsection
~\ref{sec:ch274} records angular-spectrum consequences \cite{stratton1941}.
\subsection{Evanescent subspaces and non-reconstructability}
\label{sec:ch273}
% TOC tag: [HOME]

Subsections~\ref{sec:ch271}--\ref{sec:ch272} fixed \(\ker\mathcal{R}_\omega\) and canonical
non-radiating impressed currents Eq.~\eqref{eq:ch2.nr-current-def}. Subsection~\ref{sec:ch273}
gives the Chapter~\ref{ch:02} HOME for evanescent subspaces, their overlap with non-radiating
source content, and reconstructability limits under finite observation rank---upgrading preview
Eq.~\eqref{eq:ch2.evanescent-nonreconstruct-preview} from Section~\ref{sec:ch27}
\cite{stratton1941,harrington2001}. Propagation projectors
Eqs.~\eqref{eq:ch2.spectral-projectors}--\eqref{eq:ch2.rad-null-on-evanescent} from
Subsection~\ref{sec:ch224}, non-propagating split
Eq.~\eqref{eq:ch2.nonpropagating-component-def} from Subsection~\ref{sec:ch254}, measurement
null-space extension Proposition~\ref{prop:ch2.measurement-nullspace-extension} from
Subsection~\ref{sec:ch271}, and rank cap Eq.~\eqref{eq:ch2.reconstruction-rank-cap} from
Subsection~\ref{sec:ch244} are presupposed; projector algebra and inverse-source templates are
not re-proved \cite{stratton1941}. Angular-spectrum consequences belong to Subsection
~\ref{sec:ch274}; plane-wave filters to Section~\ref{sec:ch29} \cite{harrington2001}.

\paragraph{Assumptions.}
Fix \(\omega>0\), admissible \(\mathcal{J}_\Omega\) Eq.~\eqref{eq:ch2.source-admissibility},
outgoing closure Eqs.~\eqref{eq:ch2.outgoing-radiation-refinement}--\eqref{eq:ch2.outgoing-resolvent-unique},
field pairs in \(\mathcal{F}_{\mathrm{SF}}\subset\mathcal{F}_\Omega\), and spectral completion
supporting Eqs.~\eqref{eq:ch2.spectral-projectors}--\eqref{eq:ch2.spectral-projector-algebra}
\cite{stratton1941,harrington2001}. Observation map
\(\mathcal{M}_\omega=\mathcal{O}_{\mathrm{meas}}\circ\Lambda_{\mathrm{rad}}\) from
Eq.~\eqref{eq:ch2.reconstruction-observation-map} is finite rank with
\(\mathrm{rank}(\mathcal{O}_{\mathrm{meas}})=N\) Eq.~\eqref{eq:ch2.finite-rank-observation}
\cite{stratton1941}. Regime gate Eq.~\eqref{eq:ch2.regime-operator-screen} governs transport
reports on any recovered radiative component \cite{harrington2001}.

\paragraph{Evanescent subspace (HOME definition).}
On the spectral completion \(\widehat{\mathcal{F}}_\Omega\), the \emph{evanescent subspace} is
\begin{equation}
\label{eq:ch2.evanescent-subspace-def}
\mathcal{F}_{\mathrm{ev}}
=
\left\{
\mathbf{u}\in\widehat{\mathcal{F}}_\Omega:
\Pi_{\mathrm{ev}}\mathbf{u}=\mathbf{u}
\right\},
\end{equation}
with \(\Pi_{\mathrm{ev}}\) from Eq.~\eqref{eq:ch2.spectral-projectors}
\cite{stratton1941,harrington2001}. Equation~\eqref{eq:ch2.evanescent-subspace-def} is the
Subsection~\ref{sec:ch273} HOME field-side label for modes with imaginary longitudinal
wavenumber Eq.~\eqref{eq:ch2.longitudinal-wavenumber-split} and decay rate
Eq.~\eqref{eq:ch2.evanescent-decay-rate} \cite{stratton1941}. The companion
\emph{non-propagating} subspace combines evanescent and below-cutoff content through
Eq.~\eqref{eq:ch2.nonpropagating-component-def} \cite{harrington2001}. Evanescent fields may
carry large local energy without membership in \(\mathcal{F}_{\mathrm{rad}}\):
Eq.~\eqref{eq:ch2.rad-null-on-evanescent} is imported from Subsection~\ref{sec:ch224} and not
re-derived here \cite{stratton1941}.

\paragraph{Evanescent source overlap.}
For admissible \(\Jimp\in\mathcal{J}_\Omega\), define the evanescent field content
\begin{equation}
\label{eq:ch2.evanescent-source-overlap}
\mathbf{u}_{\mathrm{ev}}
=
\Pi_{\mathrm{ev}}\,\mathcal{S}_\omega[\Jimp],
\qquad
\|\mathbf{u}_{\mathrm{ev}}\|_{\mathcal{E}}
=
\|\Pi_{\mathrm{ev}}\mathcal{S}_\omega[\Jimp]\|_{\mathcal{E}},
\end{equation}
with energy norm Eq.~\eqref{eq:ch2.field-energy-norm} \cite{stratton1941,harrington2001}.
Equation~\eqref{eq:ch2.evanescent-source-overlap} measures reactive storage excited by
\(\Jimp\) that never passes through \(\Pi_{\mathrm{rad}}\) because
\begin{equation}
\label{eq:ch2.evanescent-nonexport-law}
\Pi_{\mathrm{rad}}\,\mathbf{u}_{\mathrm{ev}}=\mathbf{0},
\qquad
\mathcal{R}_\omega[\Jimp]\ \text{depends only on}\ \Pi_{\mathrm{prop}}\mathcal{S}_\omega[\Jimp]
\ \text{after outgoing selection},
\end{equation}
by Eqs.~\eqref{eq:ch2.rad-null-on-evanescent} and~\eqref{eq:ch2.rad-projection-composition}
\cite{stratton1941}. Equation~\eqref{eq:ch2.evanescent-nonexport-law} upgrades the second
line of preview Eq.~\eqref{eq:ch2.evanescent-nonreconstruct-preview}: evanescent overlap is
field-side non-export, distinct from source-side \(\Jimp\in\ker\mathcal{R}_\omega\)
Eq.~\eqref{eq:ch2.radiation-nullspace-def} \cite{harrington2001}. A source may satisfy
\(\|\mathbf{u}_{\mathrm{ev}}\|\gg\|\mathcal{R}_\omega[\Jimp]\|\) while still radiating
\cite{stratton1941,jackson1999}.

\paragraph{Radiation nullity versus measurement nullity.}
Preview Eq.~\eqref{eq:ch2.evanescent-nonreconstruct-preview} separates three deficits:
\begin{equation}
\label{eq:ch2.radiation-vs-measurement-nullity-split}
\underbrace{\ker\mathcal{R}_\omega}_{\text{no export}}
\ \subseteq\
\underbrace{\ker\mathcal{M}_\omega}_{\text{no finite readout}},
\qquad
\ker\mathcal{M}_\omega
\supsetneq
\ker\mathcal{R}_\omega
\ \text{when}\ 
\ker\mathcal{O}_{\mathrm{meas}}\big|_{\operatorname{ran}\mathcal{R}_\omega}\neq\{\mathbf{0}\},
\end{equation}
with strict inclusion typical for finite apertures
Eq.~\eqref{eq:ch2.observation-null-space} \cite{stratton1941,harrington2001}. Equation
\eqref{eq:ch2.radiation-vs-measurement-nullity-split} restates Proposition
~\ref{prop:ch2.measurement-nullspace-extension} and
Eq.~\eqref{eq:ch2.measurement-nullspace-composition} without re-proof: evanescent storage
contributes to the first factor only when it annihilates export; finite rank contributes to
the second factor even when export is nonzero \cite{stratton1941}. Tikhonov reconstruction
Eq.~\eqref{eq:ch2.tikhonov-source-reconstruction} stabilizes ill-conditioned inversions on
\(\operatorname{ran}\mathcal{M}_\omega^{\dagger}\) but does not shrink
\(\ker\mathcal{M}_\omega\) \cite{harrington2001}.

\paragraph{Non-reconstructable observation kernel.}
Define the \emph{non-reconstructable set} at observation rank \(N\):
\begin{equation}
\label{eq:ch2.nonreconstruct-observation-kernel}
\mathcal{K}_{\mathrm{nrec}}(N)
=
\left\{
\Jimp\in\mathcal{J}_\Omega:
\mathcal{M}_\omega[\Jimp]=\mathbf{0}
\right\}
=
\ker\mathcal{M}_\omega,
\end{equation}
equivalently the kernel of Eq.~\eqref{eq:ch2.inverse-source-equation} for generic
\(\mathbf{y}\) \cite{stratton1941,harrington2001}. Equation
\eqref{eq:ch2.nonreconstruct-observation-kernel} is the Subsection~\ref{sec:ch273} HOME
reconstructability label: sources in \(\mathcal{K}_{\mathrm{nrec}}(N)\) are invisible to the
declared measurement bank regardless of near-zone \(|\E|\) \cite{stratton1941}. Rank bound
\begin{equation}
\label{eq:ch2.reconstructability-rank-law}
\dim\operatorname{ran}\mathcal{M}_\omega
\le
N
<
\infty
\ \Longrightarrow\
\dim\mathcal{K}_{\mathrm{nrec}}(N)=\infty
\ \text{on typical open domains},
\end{equation}
follows from Eq.~\eqref{eq:ch2.reconstruction-rank-cap} and infinite-dimensional
\(\mathcal{J}_\Omega\) \cite{stratton1941,harrington2001}. Geometric sector constraints
Eq.~\eqref{eq:ch2.accessible-modal-sector} and composition law
Eq.~\eqref{eq:ch2.constraint-composition-law} further restrict recoverable coordinates
independently of Eq.~\eqref{eq:ch2.nonreconstruct-observation-kernel}
\cite{stratton1941}.

\begin{figure}[t]
\centering
\ChOneFigBlock{%
\begin{tikzpicture}[ch1fig, node distance=7mm and 9mm]
  \node[ch1box] (j) {$\mathcal{J}_\Omega$};
  \node[ch1box, right=11mm of j] (s) {$\mathcal{S}_\omega$};
  \node[ch1box, right=11mm of s] (ev) {$\Pi_{\mathrm{ev}}$};
  \node[ch1box, right=11mm of ev] (r) {$\mathcal{R}_\omega$};
  \node[ch1box, right=11mm of r] (m) {$\mathcal{M}_\omega$};
  \node[ch1box, below=12mm of ev] (uev) {$\mathbf{u}_{\mathrm{ev}}$};
  \node[ch1box, below=12mm of m] (kn) {$\mathcal{K}_{\mathrm{nrec}}$};
  \draw[ch1flow] (j) -- (s) -- (ev) -- (r) -- (m);
  \draw[ch1flow] (ev) -- (uev);
  \draw[ch1flow] (m) -- (kn);
\end{tikzpicture}%
}
\caption{Subsection~\ref{sec:ch273} pipeline: evanescent overlap
Eq.~\eqref{eq:ch2.evanescent-source-overlap} never exports through
Eq.~\eqref{eq:ch2.evanescent-nonexport-law}; finite-rank observation maps
\(\mathcal{J}_\Omega\) to \(\mathcal{K}_{\mathrm{nrec}}(N)\)
Eq.~\eqref{eq:ch2.nonreconstruct-observation-kernel}.}
\label{fig:ch2.evanescent-nonreconstruct-pipeline}
\end{figure}

Figure~\ref{fig:ch2.evanescent-nonreconstruct-pipeline} orders evanescent storage before
measurement collapse \cite{stratton1941,harrington2001}.

\begin{theorem}[Evanescent subspaces and non-reconstructability]
\label{thm:ch2.evanescent-nonreconstructability}
Let \(\Jimp\in\mathcal{J}_\Omega\) be admissible with evanescent content
Eq.~\eqref{eq:ch2.evanescent-source-overlap} and observation map
\(\mathcal{M}_\omega\) of rank \(N\). Then:
\begin{enumerate}
\item[(i)] \(\Pi_{\mathrm{rad}}\mathbf{u}_{\mathrm{ev}}=\mathbf{0}\) and
\(\mathbf{u}_{\mathrm{ev}}\) carries no qualified far-zone export
Eq.~\eqref{eq:ch2.evanescent-nonexport-law}.
\item[(ii)] \(\ker\mathcal{R}_\omega\subseteq\ker\mathcal{M}_\omega=\mathcal{K}_{\mathrm{nrec}}(N)\).
\item[(iii)] Near-field data on a finite enclosing surface reduce but do not in general
eliminate \(\mathcal{K}_{\mathrm{nrec}}(N)\): if \(\mathcal{M}_{\mathrm{near}}\) samples
\(\mathcal{S}_\omega[\Jimp]\) on \(\Gamma_{\mathrm{enc}}\), then
\begin{equation}
\label{eq:ch2.near-field-nullity-reduction}
\ker\mathcal{M}_{\mathrm{near}}
\subsetneq
\ker\mathcal{M}_\omega
\ \text{when}\ \Gamma_{\mathrm{enc}}\ \text{resolves more modes than far-zone aperture},
\end{equation}
while \(\ker\mathcal{R}_\omega\) is unchanged \cite{stratton1941,harrington2001}.
\item[(iv)] Reconstructability budget:
\begin{equation}
\label{eq:ch2.reconstructability-budget}
\mathcal{N}_{\mathrm{reportable}}
\subseteq
\mathcal{N}_{\mathrm{acc}}
\cap
\{n:\ \text{resolvable by}\ \mathcal{M}_\omega\}
\cap
\{n:\ \lambda_n\neq 0\},
\end{equation}
with \(\mathcal{N}_{\mathrm{acc}}\) from Eq.~\eqref{eq:ch2.accessible-modal-sector} and modal
indices from Subsection~\ref{sec:ch261} \cite{stratton1941}.
\end{enumerate}
\end{theorem}
\begin{proof}
Part (i) is Eq.~\eqref{eq:ch2.rad-null-on-evanescent} \cite{stratton1941}. Part (ii) follows
from Proposition~\ref{prop:ch2.measurement-nullspace-extension} and definition
Eq.~\eqref{eq:ch2.nonreconstruct-observation-kernel} \cite{harrington2001}. Part (iii) uses
near- versus far-field sampling logic from Subsection~\ref{sec:ch244}: \(\mathcal{S}_\omega\) is
single-valued on outgoing fields while \(\mathcal{R}_\omega\) remains non-injective
\cite{stratton1941,harrington2001,jackson1999}. Part (iv) composes
Eq.~\eqref{eq:ch2.constraint-composition-law} with rank cap
Eq.~\eqref{eq:ch2.reconstruction-rank-cap} \cite{stratton1941}.
\end{proof}

\begin{proposition}[Evanescent content never radiates]
\label{prop:ch2.evanescent-never-exports}
For all admissible \(\Jimp\),
\begin{equation}
\label{eq:ch2.evanescent-radiation-annihilation}
\mathcal{R}_\omega[\Jimp]
=
\Pi_{\mathrm{rad}}\,\mathcal{S}_\omega[\Jimp]
=
\Pi_{\mathrm{rad}}\,\Pi_{\mathrm{prop}}\,\mathcal{S}_\omega[\Jimp],
\end{equation}
and \(\Pi_{\mathrm{rad}}\Pi_{\mathrm{ev}}=\mathbf{0}\) on \(\mathcal{F}_\Omega\)
\cite{stratton1941,harrington2001}. Equation~\eqref{eq:ch2.evanescent-radiation-annihilation}
implies \(\mathbf{u}_{\mathrm{ev}}\) from Eq.~\eqref{eq:ch2.evanescent-source-overlap} does
not alter \(\mathcal{R}_\omega[\Jimp]\); adding evanescent excitation to an otherwise
radiating source changes \(\mathbf{u}_{\mathrm{nr}}\) and local power
Eq.~\eqref{eq:ch2.component-power-budget} without changing far patterns when export is
unchanged \cite{stratton1941}. Regime gate Eq.~\eqref{eq:ch2.regime-component-gate} still
forbids reporting \(\mathbf{u}_{\mathrm{ev}}\) as transport \cite{harrington2001}.
\end{proposition}
\begin{proof}
Substitute Eqs.~\eqref{eq:ch2.rad-null-on-evanescent} and~\eqref{eq:ch2.rad-projection-composition}
into \(\mathcal{R}_\omega=\Pi_{\mathrm{rad}}\circ\mathcal{S}_\omega\)
\cite{stratton1941}. Near-zone energy follows from
Eq.~\eqref{eq:ch2.nonpropagating-component-def} \cite{harrington2001}.
\end{proof}

\begin{proposition}[Near-field scanning and null-space reduction]
\label{prop:ch2.near-field-nullity-reduction}
Let \(\Gamma_{\mathrm{enc}}\) be a closed surface enclosing \(\Omega_s\) and let
\(\mathcal{M}_{\mathrm{near}}\) sample tangential \(\E\), \(\H\) on \(\Gamma_{\mathrm{enc}}\).
Then \(\ker\mathcal{M}_{\mathrm{near}}\subseteq\ker\mathcal{M}_\omega\) when far-zone
observations are a subset of near-field functionals, but
\begin{equation}
\label{eq:ch2.near-field-kernel-inclusion}
\ker\mathcal{R}_\omega
\subseteq
\ker\mathcal{M}_{\mathrm{near}}
\subseteq
\ker\mathcal{M}_\omega
\end{equation}
with strict inclusions possible at each step \cite{stratton1941,harrington2001}. Equation
\eqref{eq:ch2.near-field-kernel-inclusion} restates Subsection~\ref{sec:ch244} near- versus
far-field data without re-opening uniqueness Theorem
~\ref{thm:ch2.exterior-radiative-uniqueness} \cite{stratton1941}. Outgoing closure remains
mandatory: identical near traces on a finite surface do not imply identical far zones without
exterior selection \cite{harrington2001,jackson1999}.
\end{proposition}
\begin{proof}
If \(\mathcal{M}_\omega[\Jimp]=\mathbf{0}\), then far functionals vanish; near samples may
still vanish on \(\Gamma_{\mathrm{enc}}\) for non-radiating additions in
\(\ker\mathcal{R}_\omega\) \cite{stratton1941}. Additional near resolution shrinks the
observation null space on \(\operatorname{ran}\mathcal{R}_\omega\) but leaves
\(\ker\mathcal{R}_\omega\) unchanged Eq.~\eqref{eq:ch2.radiation-nullspace-def}
\cite{harrington2001}.
\end{proof}

\paragraph{Below-cutoff and waveguide storage.}
When Eq.~\eqref{eq:ch2.cutoff-condition} fails, modal content lies in
\(\Pi_{\mathrm{non}}\mathcal{S}_\omega[\Jimp]\) or \(\Pi_{\mathrm{ev}}\mathcal{S}_\omega[\Jimp]\)
\cite{stratton1941,harrington2001}. Such storage satisfies
Eq.~\eqref{eq:ch2.evanescent-nonexport-law} and may dominate
Eq.~\eqref{eq:ch2.component-power-budget} on sub-wavelength shells while
\(\mathcal{P}_{\mathrm{rad}}=0\) \cite{stratton1941}. Applying far-zone pattern formulas to
below-cutoff modes violates Eq.~\eqref{eq:ch2.regime-operator-screen} as in the TE\(_{10}\)
illustration referenced from Section~\ref{sec:ch01.1.2} \cite{harrington2001}. Subsection
~\ref{sec:ch273} records operator consequences; guide nomenclature is not re-tabulated
\cite{stratton1941}.

\paragraph{Relation to non-radiating currents.}
Non-radiating configurations Eq.~\eqref{eq:ch2.nr-current-def} may carry large
\(\|\mathbf{u}_{\mathrm{ev}}\|\) when \(\mathcal{S}_\omega[\Jimp^{\mathrm{(NR)}}]\) is
evanescent-dominated \cite{stratton1941,harrington2001}. Conversely, large
\(\|\mathbf{u}_{\mathrm{ev}}\|\) does not imply \(\Jimp\in\ker\mathcal{R}_\omega\): a radiating
source may store substantial evanescent energy near \(\Omega_s\) while exporting through
\(\Pi_{\mathrm{prop}}\) \cite{stratton1941,jackson1999}. Theorem
~\ref{thm:ch2.canonical-nr-current-families} and Theorem
~\ref{thm:ch2.evanescent-nonreconstructability} address orthogonal deficits---source nullity
versus reactive storage versus finite-rank readout \cite{harrington2001}.

\paragraph{Boundary of validity.}
Theorem~\ref{thm:ch2.evanescent-nonreconstructability} assumes linear Maxwell theory at fixed
\(\omega\), outgoing closure, and declared finite-rank \(\mathcal{O}_{\mathrm{meas}}\)
\cite{stratton1941,harrington2001}. Loss modifies evanescent decay
Eq.~\eqref{eq:ch2.evanescent-decay-rate} without promoting \(\Pi_{\mathrm{ev}}\) content to
\(\Pi_{\mathrm{rad}}\) export \cite{stratton1941}. Dispersive or broadband superposition
requires Subsection~\ref{sec:ch223}; harmonic evanescent labels here are fixed-frequency slices
\cite{harrington2001}. Subsection~\ref{sec:ch274} translates non-reconstructability to
continuous angular spectra Eq.~\eqref{eq:ch2.exterior-continuous-spectral-integral}
\cite{stratton1941}.

\paragraph{Worked illustration (surface-wave storage).}
An evanescent surface mode with \(\Pi_{\mathrm{ev}}\mathbf{u}_\omega=\mathbf{u}_\omega\) satisfies
Eq.~\eqref{eq:ch2.evanescent-nonexport-law} \cite{stratton1941,harrington2001}. Local probes
may register large \(|\E|\) while Eq.~\eqref{eq:ch2.power-flux-functional} reports zero export
on qualified shells \cite{stratton1941}. Declaring angular-momentum transport from
\(\mathbf{u}_{\mathrm{ev}}\) without Eq.~\eqref{eq:ch2.regime-component-gate} violates
Section~\ref{sec:ch01.2.3} \cite{stratton1941}.

\paragraph{Worked illustration (finite array on radiating source).}
When \(\mathcal{R}_\omega[\Jimp]\neq\mathbf{0}\) but \(\mathcal{R}_\omega[\Jimp]\in\ker\mathcal{O}_{\mathrm{meas}}\),
\(\Jimp\notin\ker\mathcal{R}_\omega\) yet \(\Jimp\in\mathcal{K}_{\mathrm{nrec}}(N)\)
Eq.~\eqref{eq:ch2.nonreconstruct-observation-kernel} \cite{stratton1941,harrington2001}. This
is distinct from adding \(\Jimp^{\mathrm{(NR)}}\) Theorem~\ref{thm:ch2.canonical-nr-current-families}:
export exists but the aperture cannot resolve it \cite{stratton1941}. Rank cap
Eq.~\eqref{eq:ch2.reconstruction-rank-cap} limits recoverable modes, not evanescent storage
\cite{harrington2001}.

\paragraph{Worked illustration (near-field scan).}
Near-field sampling on \(\Gamma_{\mathrm{enc}}\) resolves more spatial content than far-zone
probes alone, shrinking \(\ker\mathcal{M}_{\mathrm{near}}\) relative to \(\ker\mathcal{M}_\omega\)
Proposition~\ref{prop:ch2.near-field-nullity-reduction} \cite{stratton1941,harrington2001}. If
outgoing closure is not enforced, two source pairs can match near traces yet differ in the far
zone Theorem~\ref{thm:ch2.exterior-radiative-uniqueness} \cite{jackson1999,stratton1941}.
\(\ker\mathcal{R}_\omega\) remains unchanged under any observation extension
\cite{harrington2001}.

\paragraph{Worked illustration (below-cutoff guide mode).}
A below-cutoff TE\(_{10}\)-like mode satisfies Eq.~\eqref{eq:ch2.cutoff-condition} with
\(\omega<\omega_{c,10}\) and lies in non-propagating subspaces
Eq.~\eqref{eq:ch2.nonpropagating-component-def} \cite{stratton1941,harrington2001}. Power
meters on the guide aperture report storage, not exterior radiation, until propagation begins
\cite{stratton1941}. Eq.~\eqref{eq:ch2.evanescent-nonexport-law} applies independently of
whether the storage is evanescent in free space or non-propagating in a guide
\cite{harrington2001}.

\paragraph{Closing summary.}
Subsection~\ref{sec:ch273} establishes the Chapter~\ref{ch:02} HOME for evanescent subspace
Eq.~\eqref{eq:ch2.evanescent-subspace-def}, source overlap
Eq.~\eqref{eq:ch2.evanescent-source-overlap}, non-export law
Eq.~\eqref{eq:ch2.evanescent-nonexport-law}, radiation/measurement split
Eq.~\eqref{eq:ch2.radiation-vs-measurement-nullity-split}, non-reconstructable kernel
Eq.~\eqref{eq:ch2.nonreconstruct-observation-kernel}, rank law
Eq.~\eqref{eq:ch2.reconstructability-rank-law}, Theorem
~\ref{thm:ch2.evanescent-nonreconstructability} with budgets
Eqs.~\eqref{eq:ch2.near-field-nullity-reduction}--\eqref{eq:ch2.reconstructability-budget},
Propositions~\ref{prop:ch2.evanescent-never-exports} and~\ref{prop:ch2.near-field-nullity-reduction},
upgrading preview Eq.~\eqref{eq:ch2.evanescent-nonreconstruct-preview} and importing
Subsections~\ref{sec:ch224},~\ref{sec:ch244}, and~\ref{sec:ch271} by reference
\cite{stratton1941,harrington2001,jackson1999}. Subsection~\ref{sec:ch274} records consequences
for angular spectra; Section~\ref{sec:ch29} develops plane-wave action on the same
reconstructability budget \cite{stratton1941}.
\subsection{Consequences for angular spectra}
\label{sec:ch274}
% TOC tag: [HOME]

Subsections~\ref{sec:ch271}--\ref{sec:ch273} fixed \(\ker\mathcal{R}_\omega\), canonical
non-radiating currents Eq.~\eqref{eq:ch2.nr-current-def}, and evanescent non-reconstructability
Theorem~\ref{thm:ch2.evanescent-nonreconstructability}. Subsection~\ref{sec:ch274} gives the
Chapter~\ref{ch:02} HOME for how null-space, evanescent, and measurement deficits propagate to
continuous angular spectra Eq.~\eqref{eq:ch2.exterior-continuous-spectral-integral}---upgrading
preview Eq.~\eqref{eq:ch2.angular-spectrum-null-preview} from Section~\ref{sec:ch27}
\cite{stratton1941,harrington2001,jackson1999}. Discrete far-pattern expansion
Eq.~\eqref{eq:ch2.spectral-far-pattern-expansion}, accessible sector
Eq.~\eqref{eq:ch2.angular-accessibility-sector} from Proposition
~\ref{prop:ch2.angular-accessibility-preview}, geometric sector
Eq.~\eqref{eq:ch2.accessible-modal-sector}, and reconstructability budget
Eq.~\eqref{eq:ch2.reconstructability-budget} are presupposed; rigged-space normalization belongs
to Chapter~\ref{ch:03} Subsection~\ref{sec:ch313} \cite{stratton1941}. Plane-wave filters and
retained/suppressed sheets are developed in Section~\ref{sec:ch29} and Subsection
~\ref{sec:ch294}; Proposition~\ref{prop:ch2.angular-accessibility-preview} is not re-proved
\cite{harrington2001}.

\paragraph{Assumptions.}
Fix \(\omega>0\), outgoing closure, admissible \(\mathcal{J}_\Omega\), and far amplitude
\(\E_0(\theta,\phi)\) from Eq.~\eqref{eq:ch2.far-field-amplitude-integral} when
\(\mathcal{D}_{J,\mathrm{reg}}=1\) \cite{stratton1941,harrington2001}. Continuous angular
coefficients \(a_\sigma(\widehat{\mathbf{k}})\) are defined through the rigged limit of
Eq.~\eqref{eq:ch2.exterior-continuous-spectral-integral} compatible with modal amplitudes
Eq.~\eqref{eq:ch2.source-modal-amplitude} \cite{stratton1941}. Finite-rank observation
\(\mathcal{M}_\omega\) and non-reconstructable set \(\mathcal{K}_{\mathrm{nrec}}(N)\) from
Eq.~\eqref{eq:ch2.nonreconstruct-observation-kernel} are presupposed
\cite{harrington2001}. Regime gate Eq.~\eqref{eq:ch2.regime-operator-screen} governs transport
claims on any retained angular content \cite{stratton1941}.

\paragraph{Angular spectrum under radiation nullity (HOME statement).}
Let \(a(\widehat{\mathbf{k}})\) denote the scalar angular spectrum obtained from
Eq.~\eqref{eq:ch2.exterior-continuous-spectral-integral} on the outgoing transverse sheet
\cite{stratton1941,harrington2001}. For \(\Jimp\in\ker\mathcal{R}_\omega\),
\begin{equation}
\label{eq:ch2.angular-spectrum-nullspace-law}
a(\widehat{\mathbf{k}})=0
\quad\text{for all}\ \widehat{\mathbf{k}}\in S^{2}
\ \text{on the exporting sheet},
\end{equation}
equivalently \(\E_0(\theta,\phi)=\mathbf{0}\) in Eq.~\eqref{eq:ch2.spectral-far-pattern-expansion}
\cite{stratton1941,jackson1999}. Equation~\eqref{eq:ch2.angular-spectrum-nullspace-law} is the
Subsection~\ref{sec:ch274} HOME upgrade of the first clause of preview
Eq.~\eqref{eq:ch2.angular-spectrum-null-preview}: null-space membership annihilates the entire
far-zone angular measure, not merely selected harmonics
\cite{harrington2001,stratton1941}. Near-zone reactive content in
\(\Pi_{\mathrm{ev}}\mathcal{S}_\omega[\Jimp]\) may remain nonzero without contributing to
\(a(\widehat{\mathbf{k}})\) Eq.~\eqref{eq:ch2.evanescent-nonexport-law}
\cite{stratton1941}.

\paragraph{Suppressed and accessible angular sheets.}
Define the \emph{suppressed angular sheet}
\begin{equation}
\label{eq:ch2.suppressed-angular-sheet-def}
\mathcal{K}_{\mathrm{sup}}
=
S^{2}\setminus\mathcal{K}_{\mathrm{acc}},
\qquad
\mathcal{K}_{\mathrm{acc}}
\ \text{from Eq.~\eqref{eq:ch2.angular-accessibility-sector}},
\end{equation}
as directions annihilated by geometric, support, or parity constraints before observation
\cite{stratton1941,harrington2001}. Equation~\eqref{eq:ch2.suppressed-angular-sheet-def} restates
preview Eq.~\eqref{eq:ch2.angular-spectrum-null-preview}: \(a(\widehat{\mathbf{k}})=0\) on
\(\mathcal{K}_{\mathrm{sup}}\) does not imply \(\Jimp\in\ker\mathcal{R}_\omega\), because
geometric forbidding Eq.~\eqref{eq:ch2.geometric-modal-annihilation} and radiation nullity
Eq.~\eqref{eq:ch2.angular-spectrum-nullspace-law} are distinct mechanisms
\cite{stratton1941}. Accessible sheet \(\mathcal{K}_{\mathrm{acc}}\) intersects reportable modal
sector Eq.~\eqref{eq:ch2.reconstructability-budget} \cite{harrington2001}.

\paragraph{Non-radiating superposition invariance.}
For radiating \(\Jimp^{(\mathrm{rad})}\) and non-radiating
\(\Jimp^{\mathrm{(NR)}}\in\ker\mathcal{R}_\omega\),
\begin{equation}
\label{eq:ch2.angular-spectrum-nr-invariance}
a(\widehat{\mathbf{k}};\Jimp^{(\mathrm{rad})}+\Jimp^{\mathrm{(NR)}})
=
a(\widehat{\mathbf{k}};\Jimp^{(\mathrm{rad})}),
\end{equation}
on the outgoing transverse sheet, restating superposition law
Eq.~\eqref{eq:ch2.nr-superposition-law} in angular coordinates
\cite{stratton1941,harrington2001,jackson1999}. Equation~\eqref{eq:ch2.angular-spectrum-nr-invariance}
is the continuous form of modal invisibility Eq.~\eqref{eq:ch2.nr-modal-invisibility} and
explains why inverse angular-spectrum solvers cannot detect
\(\Jimp^{\mathrm{(NR)}}\) from far data alone \cite{stratton1941}. Canonical families
Theorem~\ref{thm:ch2.canonical-nr-current-families} therefore induce \emph{invisible} angular
perturbations at fixed \(a(\widehat{\mathbf{k}})\) \cite{harrington2001}.

\paragraph{Evanescent decoupling from angular spectra.}
Evanescent field content does not enter \(a(\widehat{\mathbf{k}})\):
\begin{equation}
\label{eq:ch2.evanescent-angular-decoupling}
a(\widehat{\mathbf{k}})
\ \text{depends only on}\ 
\Pi_{\mathrm{prop}}\mathcal{S}_\omega[\Jimp]
\ \text{after}\ \Pi_{\mathrm{rad}},
\qquad
\Pi_{\mathrm{ev}}\mathcal{S}_\omega[\Jimp]
\ \not\mapsto\ a(\widehat{\mathbf{k}}),
\end{equation}
by Eqs.~\eqref{eq:ch2.evanescent-nonexport-law} and~\eqref{eq:ch2.evanescent-radiation-annihilation}
\cite{stratton1941,harrington2001}. Equation~\eqref{eq:ch2.evanescent-angular-decoupling}
specializes the second line of preview Eq.~\eqref{eq:ch2.angular-spectrum-null-preview}:
reactive storage may be large in \(\Pi_{\mathrm{ev}}\mathcal{S}_\omega[\Jimp]\) while
\(a(\widehat{\mathbf{k}})\) reports only qualified export \cite{stratton1941}. Regime gate
Eq.~\eqref{eq:ch2.regime-component-gate} forbids interpreting \(\mathbf{u}_{\mathrm{ev}}\) as
angular-transport data \cite{harrington2001}.

\paragraph{Observation rank on angular spectra.}
Finite aperture collapses recoverable angular content:
\begin{equation}
\label{eq:ch2.angular-spectrum-observation-rank}
\{\,a(\widehat{\mathbf{k}})\,\text{recoverable from}\ \mathcal{M}_\omega\,\}
\subseteq
\{\,a(\widehat{\mathbf{k}})\,\text{supported on}\ \mathcal{K}_{\mathrm{acc}}\cap\operatorname{ran}\mathcal{R}_\omega\,\},
\end{equation}
with strict inclusion when \(\ker\mathcal{O}_{\mathrm{meas}}\big|_{\operatorname{ran}\mathcal{R}_\omega}\neq\{\mathbf{0}\}\)
Eq.~\eqref{eq:ch2.observation-null-space} \cite{stratton1941,harrington2001}. Equation
\eqref{eq:ch2.angular-spectrum-observation-rank} links angular spectra to
\(\mathcal{K}_{\mathrm{nrec}}(N)\): sources invisible to \(\mathcal{M}_\omega\) need not have
\(a(\widehat{\mathbf{k}})=\mathbf{0}\) pointwise \cite{stratton1941}. Rank cap
Eq.~\eqref{eq:ch2.reconstruction-rank-cap} limits resolved harmonics independently of
Eq.~\eqref{eq:ch2.angular-spectrum-nullspace-law} \cite{harrington2001}.

\begin{figure}[t]
\centering
\ChOneFigBlock{%
\begin{tikzpicture}[ch1fig, node distance=7mm and 9mm]
  \node[ch1box] (j) {$\mathcal{J}_\Omega$};
  \node[ch1box, right=11mm of j] (r) {$\mathcal{R}_\omega$};
  \node[ch1box, right=11mm of r] (a) {$a(\widehat{\mathbf{k}})$};
  \node[ch1box, right=11mm of a] (k) {$\mathcal{K}_{\mathrm{acc}}$};
  \node[ch1box, below=12mm of j] (ker) {$\ker\mathcal{R}_\omega$};
  \node[ch1box, below=12mm of a] (sup) {$\mathcal{K}_{\mathrm{sup}}$};
  \draw[ch1flow] (j) -- (r) -- (a) -- (k);
  \draw[ch1flow, dashed] (ker) -- (j);
  \draw[ch1flow, dashed] (sup) -- (a);
\end{tikzpicture}%
}
\caption{Subsection~\ref{sec:ch274} angular pipeline: \(\ker\mathcal{R}_\omega\) annihilates
\(a(\widehat{\mathbf{k}})\) Eq.~\eqref{eq:ch2.angular-spectrum-nullspace-law}; suppressed sheet
\(\mathcal{K}_{\mathrm{sup}}\) Eq.~\eqref{eq:ch2.suppressed-angular-sheet-def} removes directions
before accessible sector \(\mathcal{K}_{\mathrm{acc}}\).}
\label{fig:ch2.angular-spectrum-null-pipeline}
\end{figure}

Figure~\ref{fig:ch2.angular-spectrum-null-pipeline} orders source nullity, angular amplitude, and
accessible directions before Section~\ref{sec:ch29} plane-wave action \cite{stratton1941,harrington2001}.

\begin{theorem}[Null-space consequences for angular spectra]
\label{thm:ch2.angular-spectrum-null-consequences}
Let \(a(\widehat{\mathbf{k}})\) be the outgoing angular spectrum from
Eq.~\eqref{eq:ch2.exterior-continuous-spectral-integral} and let
\(\Jimp\in\mathcal{J}_\Omega\) be admissible. Then:
\begin{enumerate}
\item[(i)] \(\Jimp\in\ker\mathcal{R}_\omega\Rightarrow a(\widehat{\mathbf{k}})=0\) on the
exporting sheet Eq.~\eqref{eq:ch2.angular-spectrum-nullspace-law}.
\item[(ii)] Superposition with \(\Jimp^{\mathrm{(NR)}}\in\ker\mathcal{R}_\omega\) leaves
\(a(\widehat{\mathbf{k}})\) unchanged Eq.~\eqref{eq:ch2.angular-spectrum-nr-invariance}.
\item[(iii)] Evanescent content satisfies Eq.~\eqref{eq:ch2.evanescent-angular-decoupling} and
does not alter \(a(\widehat{\mathbf{k}})\) when export is unchanged.
\item[(iv)] Full angular budget:
\begin{equation}
\label{eq:ch2.angular-spectrum-composition-law}
\text{support}\big(a(\widehat{\mathbf{k}})\big)
\subseteq
\mathcal{K}_{\mathrm{acc}}
\cap
\{\widehat{\mathbf{k}}:\ \text{resolvable by}\ \mathcal{M}_\omega\}
\cap
\{\widehat{\mathbf{k}}:\ \lambda_n\neq 0\ \text{modal labels}\},
\end{equation}
composing geometric, observation, and Fredholm constraints from
Eqs.~\eqref{eq:ch2.constraint-composition-law}--\eqref{eq:ch2.reconstructability-budget}
\cite{stratton1941,harrington2001}.
\end{enumerate}
\end{theorem}
\begin{proof}
Part (i) follows from \(\mathcal{R}_\omega[\Jimp]=\mathbf{0}\) and far-zone reduction
Theorem~\ref{thm:ch2.far-field-approximation} \cite{stratton1941,jackson1999}. Part (ii) uses
Theorem~\ref{thm:ch2.canonical-nr-current-families} and linearity of angular extraction
\cite{harrington2001}. Part (iii) is Proposition~\ref{prop:ch2.evanescent-never-exports} in
angular coordinates \cite{stratton1941}. Part (iv) composes
Eq.~\eqref{eq:ch2.accessible-modal-sector}, Proposition
~\ref{prop:ch2.angular-accessibility-preview}, and
Eq.~\eqref{eq:ch2.angular-spectrum-observation-rank} \cite{stratton1941,harrington2001}.
\end{proof}

\begin{proposition}[Discrete--continuous null correspondence]
\label{prop:ch2.continuous-discrete-null-correspondence}
Modal coefficients in Eq.~\eqref{eq:ch2.spectral-far-pattern-expansion} and continuous
\(a(\widehat{\mathbf{k}})\) satisfy
\begin{equation}
\label{eq:ch2.modal-continuous-null-bridge}
\lambda_n\,c_n[\mathfrak{D}]=0\ \text{for all exporting}\ n
\ \Longleftrightarrow\
a(\widehat{\mathbf{k}})=0\ \text{on}\ S^{2}
\end{equation}
when the modal basis is complete on \(\operatorname{ran}\mathcal{R}_\omega\) in the rigged sense
\cite{stratton1941,harrington2001}. Equation~\eqref{eq:ch2.modal-continuous-null-bridge} connects
Eq.~\eqref{eq:ch2.nullspace-modal-annihilation-law} to
Eq.~\eqref{eq:ch2.angular-spectrum-nullspace-law} without re-proving completeness in
Chapter~\ref{ch:03} \cite{stratton1941}. Truncated modal books misstate
Eq.~\eqref{eq:ch2.modal-continuous-null-bridge} on exterior domains
\cite{jackson1999,harrington2001}.
\end{proposition}
\begin{proof}
Proposition~\ref{prop:ch2.far-pattern-from-modes} and
Eq.~\eqref{eq:ch2.exterior-continuous-spectral-integral} identify discrete and continuous
representations on outgoing transverse modes \cite{stratton1941}. Null modal coefficients imply
vanishing \(\E_0\); conversely \(\E_0=\mathbf{0}\) forces all exporting modal weights to zero
\cite{harrington2001,jackson1999}.
\end{proof}

\begin{proposition}[Angular spectra and inverse problems]
\label{prop:ch2.angular-spectrum-inverse-obstruction}
Given target angular spectrum \(a_{\mathrm{tgt}}(\widehat{\mathbf{k}})\), the source inversion
problem seeks \(\Jimp\) with \(a(\widehat{\mathbf{k}};\Jimp)=a_{\mathrm{tgt}}(\widehat{\mathbf{k}})\).
Any solution is unique only up to additions in \(\ker\mathcal{R}_\omega\):
\begin{equation}
\label{eq:ch2.angular-inverse-null-quotient}
\Jimp^{(1)}-\Jimp^{(2)}\in\ker\mathcal{R}_\omega
\quad\text{whenever}\quad
a(\widehat{\mathbf{k}};\Jimp^{(1)})=a(\widehat{\mathbf{k}};\Jimp^{(2)}),
\end{equation}
and up to \(\mathcal{K}_{\mathrm{nrec}}(N)\) at finite observation rank
\cite{stratton1941,harrington2001}. Equation~\eqref{eq:ch2.angular-inverse-null-quotient} is the
angular restatement of Eq.~\eqref{eq:ch2.inverse-source-equation} ill-posedness; Tikhonov
regularization Eq.~\eqref{eq:ch2.tikhonov-source-reconstruction} selects a representative without
eliminating the quotient \cite{stratton1941}.
\end{proposition}
\begin{proof}
If \(a(\widehat{\mathbf{k}};\Jimp^{(1)})=a(\widehat{\mathbf{k}};\Jimp^{(2)})\), then
\(\mathcal{R}_\omega[\Jimp^{(1)}-\Jimp^{(2)}]=\mathbf{0}\) on \(\mathcal{F}_{\mathrm{rad}}\)
\cite{stratton1941}. Finite \(\mathcal{M}_\omega\) adds measurement null directions
Proposition~\ref{prop:ch2.measurement-nullspace-extension} \cite{harrington2001}.
\end{proof}

\paragraph{Realizability, accessibility, and nullity.}
Three deficits are often conflated in inverse problems and structured-beam reporting
\cite{devaney2004,marengo2000,bleszynski1996}. \emph{Nullity} is source-side membership in
\(\ker\mathcal{R}_\omega\): far data cannot detect the addition
Eq.~\eqref{eq:ch2.angular-spectrum-nr-invariance} \cite{reed1950,stratton1941}.
\emph{Accessibility} is field-side: which directions in \(a(\widehat{\mathbf{k}})\) survive
geometric, transport, and regime gates Eqs.~\eqref{eq:ch2.angular-accessibility-sector}--\eqref{eq:ch2.accessible-angular-budget} even when \(\mathcal{R}_\omega[\Jimp]\neq\mathbf{0}\)
\cite{harrington2001}. \emph{Realizability}---whether a target \(a_{\mathrm{target}}(\widehat{\mathbf{k}})\) is
\(\mathcal{R}_\omega[\Jimp]\) for some admissible \(\Jimp\)---is not decided by accessibility alone:
Chapter~\ref{ch:04} owns realizability proofs and truncation limits; Chapter~\ref{ch:02} supplies only
the forward map and null-space taxonomy that constrain any realizability claim
\cite{colton1992,stratton1941}. Evanescent overlap
Eq.~\eqref{eq:ch2.evanescent-source-overlap} changes near fields without altering \(a\) when export is
fixed Theorem~\ref{thm:ch2.near-to-far-transformation}(iii) \cite{stratton1941,harrington2001}.

\paragraph{Forward link to angular-spectrum action.}
Section~\ref{sec:ch29} develops plane-wave operators that act on \(a(\widehat{\mathbf{k}})\);
Subsection~\ref{sec:ch294} refines far-field angular accessibility on
\(\mathcal{K}_{\mathrm{acc}}\) \cite{stratton1941,harrington2001}. Subsection~\ref{sec:ch274}
fixes the \emph{source-side} constraints those operators inherit:
Eqs.~\eqref{eq:ch2.angular-spectrum-nullspace-law}--\eqref{eq:ch2.angular-spectrum-composition-law}
must be composed with regime gate Eq.~\eqref{eq:ch2.regime-operator-screen} before transport
claims \cite{stratton1941}. Reciprocity identities on coupled angular modes belong to
Section~\ref{sec:ch28} \cite{harrington2001}.

\paragraph{Boundary of validity.}
Theorem~\ref{thm:ch2.angular-spectrum-null-consequences} assumes linear Maxwell theory at fixed
\(\omega\), outgoing closure, and rigged-sense convergence of
Eq.~\eqref{eq:ch2.exterior-continuous-spectral-integral} \cite{stratton1941,harrington2001}.
Loss shifts phase of \(a(\widehat{\mathbf{k}})\) without converting evanescent storage to export
\cite{stratton1941}. Vector harmonic explicit forms of \(\mathbf{f}_n(\theta,\phi)\) belong to
Chapter~\ref{ch:03}; Subsection~\ref{sec:ch274} records operator consequences only
\cite{harrington2001}.

\paragraph{Worked illustration (null current, nonzero near field).}
\(\Jimp^{\mathrm{(NR)}}\) from Eq.~\eqref{eq:ch2.nr-loop-cancellation} satisfies
Eq.~\eqref{eq:ch2.angular-spectrum-nullspace-law} while \(|\E|\) may be large near the loop
Proposition~\ref{prop:ch2.nr-near-field-nonzero} \cite{stratton1941,jackson1999}. Angular-spectrum
probes of \(\mathcal{S}_\omega[\Jimp^{\mathrm{(NR)}}]\) without \(\Pi_{\mathrm{rad}}\) misstate
null-space membership \cite{harrington2001}.

\paragraph{Worked illustration (symmetry-forbidden sheet).}
A ground-plane symmetric source may have \(a(\widehat{\mathbf{k}})=0\) on half of
\(\mathcal{K}_{\mathrm{sup}}\) through Eq.~\eqref{eq:ch2.parity-selection-rule} while still
radiating on \(\mathcal{K}_{\mathrm{acc}}\) \cite{stratton1941,jackson1999}. Confusing
\(\mathcal{K}_{\mathrm{sup}}\) with Eq.~\eqref{eq:ch2.angular-spectrum-nullspace-law} misattributes
pattern nulls \cite{harrington2001}.

\paragraph{Worked illustration (finite array angular truncation).}
When \(\mathcal{R}_\omega[\Jimp]\neq\mathbf{0}\) but the aperture resolves only \(N\) harmonics,
Eq.~\eqref{eq:ch2.angular-spectrum-observation-rank} leaves unresolved \(\widehat{\mathbf{k}}\)
in the measurement null space even when \(a(\widehat{\mathbf{k}})\) is nonzero
\cite{stratton1941,harrington2001}. This is distinct from
Eq.~\eqref{eq:ch2.angular-spectrum-nr-invariance} \cite{stratton1941}.

\paragraph{Worked illustration (modal zero versus full null).}
A single \(\lambda_n=0\) mode contributes to Eq.~\eqref{eq:ch2.modal-null-sector-reduction} but
need not annihilate the full \(a(\widehat{\mathbf{k}})\) unless all exporting modes vanish
Eq.~\eqref{eq:ch2.modal-continuous-null-bridge} \cite{stratton1941,harrington2001}. Exterior
problems require continuous measures Eq.~\eqref{eq:ch2.exterior-continuous-spectral-integral},
not finite cavity truncation \cite{jackson1999,stratton1941}.

\paragraph{Closing summary.}
Subsection~\ref{sec:ch274} establishes the Chapter~\ref{ch:02} HOME for angular null-space law
Eq.~\eqref{eq:ch2.angular-spectrum-nullspace-law}, suppressed sheet
Eq.~\eqref{eq:ch2.suppressed-angular-sheet-def}, NR invariance
Eq.~\eqref{eq:ch2.angular-spectrum-nr-invariance}, evanescent decoupling
Eq.~\eqref{eq:ch2.evanescent-angular-decoupling}, observation rank
Eq.~\eqref{eq:ch2.angular-spectrum-observation-rank}, Theorem
~\ref{thm:ch2.angular-spectrum-null-consequences} with composition
Eq.~\eqref{eq:ch2.angular-spectrum-composition-law}, discrete--continuous bridge
Proposition~\ref{prop:ch2.continuous-discrete-null-correspondence}, and inverse obstruction
Proposition~\ref{prop:ch2.angular-spectrum-inverse-obstruction}, upgrading preview
Eq.~\eqref{eq:ch2.angular-spectrum-null-preview} \cite{stratton1941,harrington2001,jackson1999}.

\medskip
\noindent
Section~\ref{sec:ch27} closes with the null-space chain begun in its introduction:
operator structure and \(\ker\mathcal{R}_\omega\) Theorem
~\ref{thm:ch2.radiation-nullspace-structure} (Subsection~\ref{sec:ch271}), canonical
non-radiating current families Theorem~\ref{thm:ch2.canonical-nr-current-families}
(Subsection~\ref{sec:ch272}), evanescent subspaces and non-reconstructability Theorem
~\ref{thm:ch2.evanescent-nonreconstructability} (Subsection~\ref{sec:ch273}), and angular-spectrum
consequences Theorem~\ref{thm:ch2.angular-spectrum-null-consequences} with accessible and
suppressed sheets Eqs.~\eqref{eq:ch2.angular-accessibility-sector}--\eqref{eq:ch2.suppressed-angular-sheet-def}
(Subsection~\ref{sec:ch274}) \cite{stratton1941,harrington2001,jackson1999}. Spectral and
geometric filters from Section~\ref{sec:ch26} now meet source-side invisibility: far-zone reports
use only \(a(\widehat{\mathbf{k}})\) on the composed budget
Eq.~\eqref{eq:ch2.angular-spectrum-composition-law}, while \(\ker\mathcal{R}_\omega\),
\(\Pi_{\mathrm{ev}}\mathcal{S}_\omega[\Jimp]\), and \(\mathcal{K}_{\mathrm{nrec}}(N)\) explain
deficits that no angular label alone resolves \cite{stratton1941}. Section~\ref{sec:ch28} adds
reciprocity on the same modal coordinates; Section~\ref{sec:ch29} and Subsection~\ref{sec:ch294}
develop plane-wave action and far-field angular accessibility on \(\mathcal{K}_{\mathrm{acc}}\);
Chapter~\ref{ch:03} supplies explicit harmonic bases \cite{harrington2001}. Transport and
angular-transport claims throughout Volume~I must compose with
Eq.~\eqref{eq:ch2.section27-nullspace-chain},
Eq.~\eqref{eq:ch2.angular-spectrum-composition-law}, and regime gate
Eq.~\eqref{eq:ch2.regime-operator-screen}, not with raw \(\mathcal{S}_\omega[\Jimp]\) or
unqualified near-zone reactive content \cite{stratton1941,harrington2001}.
\section{Reciprocity and Radiative Coupling}
\label{sec:ch28}

Section~\ref{sec:ch27} closed the null-space chain
Eq.~\eqref{eq:ch2.section27-nullspace-chain}: operator kernels, canonical non-radiating
currents, evanescent non-reconstructability, and angular-spectrum consequences
Theorem~\ref{thm:ch2.angular-spectrum-null-consequences}
\cite{stratton1941,harrington2001,jackson1999}. That program classified invisible and
partially observable source content on \(\mathcal{R}_\omega\) and \(\mathcal{M}_\omega\).
Section~\ref{sec:ch28} addresses the complementary \emph{exchange symmetry} question: when
do transmit and receive maps coincide, how do radiative modes couple under bilinear pairings,
and which orthogonality relations constrain coupled spectra on \(\mathcal{F}_{\mathrm{rad}}\)
\cite{stratton1941,harrington2001}. Preview swap Eq.~\eqref{eq:ch2.dyadic-reciprocity-swap}
from Subsection~\ref{sec:ch231} named kernel interchange; Section~\ref{sec:ch28} upgrades
reciprocity to operator form on \(\mathcal{S}_\omega\), \(\mathcal{R}_\omega\), and modal
coordinates from Sections~\ref{sec:ch26}--\ref{sec:ch27} \cite{stratton1941}. Plane-wave
angular-spectrum action belongs to Section~\ref{sec:ch29}; explicit harmonic orthogonality and
completeness belong to Chapter~\ref{ch:03} Subsection~\ref{sec:ch343}; continuous-spectrum
normalization belongs to Subsection~\ref{sec:ch352} by reference to
Subsection~\ref{sec:ch282} \cite{harrington2001}.

The problem addressed here is \emph{bilinear symmetry} of radiation maps under source--probe
interchange. Spectral decomposition Theorem~\ref{thm:ch2.radiation-spectral-decomposition}
and biorthogonality Eq.~\eqref{eq:ch2.radiation-biorthogonality} from Subsection
~\ref{sec:ch261} fixed modal coordinates on \(\mathcal{F}_{\mathrm{rad}}\); null-space
structure Theorem~\ref{thm:ch2.radiation-nullspace-structure} showed that many sources share
the same radiative image \cite{stratton1941,harrington2001}. Reciprocity does not remove
\(\ker\mathcal{R}_\omega\), but it constrains how \emph{coupling coefficients} between declared
modes behave when impressed sources and test probes are swapped
\cite{stratton1941,jackson1999}. Treating antenna transmit and receive as unrelated maps
violates Eq.~\eqref{eq:ch2.dyadic-reciprocity-swap} even when forward integrals from
Theorem~\ref{thm:ch2.source-to-field-integral} are numerically accurate
\cite{harrington2001}. Section~\ref{sec:ch28} separates Lorentz reciprocity in operator form,
modal orthogonality and coupling, and radiation--reception interchange without re-opening
Green construction from Section~\ref{sec:ch23} \cite{stratton1941}.

\paragraph{Reciprocity development chain.}
Let \(\omega>0\) be fixed, admissible sources lie in \(\mathcal{J}_\Omega\), and outgoing
field pairs belong to \(\mathcal{F}_{\mathrm{SF}}\) Eq.~\eqref{eq:ch2.admissible-field-space}.
Section~\ref{sec:ch28} organizes analysis in three stages:
\begin{equation}
\label{eq:ch2.section28-reciprocity-chain}
\text{Lorentz reciprocity}
\ \xrightarrow{\ \text{operator form}\ }
\mathcal{S}_\omega,\mathcal{R}_\omega
\ \xrightarrow{\ \text{modal}\ }
\text{orthogonality / coupling}
\ \xrightarrow{\ \text{systems}\ }
\text{radiation \& reception},
\end{equation}
where the first arrow is Subsection~\ref{sec:ch281}, the second is Subsection~\ref{sec:ch282},
and the third is Subsection~\ref{sec:ch283} \cite{stratton1941,harrington2001}. Equation
\eqref{eq:ch2.section28-reciprocity-chain} restates Eq.~\eqref{eq:ch2.section26-spectral-chain}
and Eq.~\eqref{eq:ch2.section27-nullspace-chain} at the bilinear level: spectral modes from
Section~\ref{sec:ch26} supply coordinates; null-space quotients from Section~\ref{sec:ch27}
explain non-uniqueness; reciprocity constrains \emph{symmetric} parts of coupling on
\(\operatorname{ran}\mathcal{R}_\omega\) \cite{stratton1941}. Angular budget
Eq.~\eqref{eq:ch2.angular-spectrum-composition-law} and regime gate
Eq.~\eqref{eq:ch2.regime-operator-screen} remain mandatory before transport claims on coupled
modes \cite{harrington2001}.

\paragraph{Operator targets.}
The primary bilinear objects are Maxwell adjoint pairing
Eq.~\eqref{eq:ch2.maxwell-adjoint-pairing}, source dual pairing
Eq.~\eqref{eq:ch2.dual-source-pairing}, and radiation map
Eq.~\eqref{eq:ch2.spectral-operator-targets} \cite{stratton1941,harrington2001}. Secondary
targets include dyadic reciprocity swap Eq.~\eqref{eq:ch2.dyadic-reciprocity-swap}, observation
composition Eq.~\eqref{eq:ch2.reconstruction-observation-map}, and modal biorthogonality
Eq.~\eqref{eq:ch2.radiation-biorthogonality} \cite{stratton1941}. Section~\ref{sec:ch28}
develops reciprocity on the same outgoing branch fixed by
Eqs.~\eqref{eq:ch2.outgoing-radiation-refinement}--\eqref{eq:ch2.outgoing-resolvent-unique};
incoming sheets are excluded from \(\mathcal{F}_{\mathrm{rad}}\) \cite{harrington2001}.

\paragraph{Lorentz reciprocity (preview).}
Subsection~\ref{sec:ch281} gives the Chapter~\ref{ch:02} HOME operator statement: for
admissible source pairs \((\Jimp^{(1)},\Jimp^{(2)})\) and induced fields
\((\mathbf{u}_1,\mathbf{u}_2)\),
\begin{equation}
\label{eq:ch2.lorentz-reciprocity-preview}
\langle \mathbf{u}_1,\,\mathcal{S}_\omega[\Jimp^{(2)}]\rangle_{\mathcal{U}}
=
\langle \mathbf{u}_2,\,\mathcal{S}_\omega[\Jimp^{(1)}]\rangle_{\mathcal{U}}
\end{equation}
on passive linear media without magneto-electric coupling, with pairing induced from
Eq.~\eqref{eq:ch2.maxwell-adjoint-pairing} on outgoing admissible states
\cite{stratton1941,harrington2001,jackson1999}. Equation~\eqref{eq:ch2.lorentz-reciprocity-preview}
upgrades preview Eq.~\eqref{eq:ch2.dyadic-reciprocity-swap} from kernel swap to full
source--field reciprocity on \(\mathcal{S}_\omega\); restrictions on loss, anisotropy, and
boundary class are stated in Subsection~\ref{sec:ch281} \cite{stratton1941}. Reciprocity on
\(\mathcal{R}_\omega=\Pi_{\mathrm{rad}}\circ\mathcal{S}_\omega\) follows by projection when
both sides satisfy regime gate Eq.~\eqref{eq:ch2.regime-component-gate}
\cite{harrington2001}.

\paragraph{Modal orthogonality and coupling (preview).}
On radiation eigenmodes Eq.~\eqref{eq:ch2.radiation-eigenpair-def}, Subsection~\ref{sec:ch282}
develops orthogonality and coupling coefficients:
\begin{equation}
\label{eq:ch2.modal-orthogonality-preview}
\langle \mathbf{u}_m,\,\mathbf{u}_n\rangle_{\mathcal{U}}
=
\delta_{mn},
\qquad
C_{mn}
=
\langle \widetilde{\mathbf{u}}_m,\,\mathcal{R}_\omega[\Jimp^{(n)}]\rangle_{\mathcal{U}},
\end{equation}
with \(\widetilde{\mathbf{u}}_m\) adjoint modes tied to
Eq.~\eqref{eq:ch2.radiation-biorthogonality} \cite{stratton1941,harrington2001}. Equation
\eqref{eq:ch2.modal-orthogonality-preview} is the Section~\ref{sec:ch28} target for
Subsection~\ref{sec:ch282} HOME orthogonality and coupling on \(\mathcal{F}_{\mathrm{rad}}\);
explicit vector harmonic normalization belongs to Chapter~\ref{ch:03}
\cite{stratton1941}. Modes with \(\lambda_n=0\) lie in null sectors
Eq.~\eqref{eq:ch2.modal-null-index-set}; reciprocity constraints apply on exporting modes
with \(\lambda_n\neq 0\) \cite{harrington2001}.

\paragraph{Radiation and reception (preview).}
Subsection~\ref{sec:ch283} records transmit--receive interchange for coupled systems:
\begin{equation}
\label{eq:ch2.transmit-receive-swap-preview}
\mathcal{M}_{\mathrm{rx}}[\Jimp^{\mathrm{(tx)}}]
=
\mathcal{M}_{\mathrm{tx}}[\Jimp^{\mathrm{(rx)}}]
\quad\text{under reciprocity-qualified pairings},
\end{equation}
with \(\mathcal{M}_{\mathrm{tx}}\), \(\mathcal{M}_{\mathrm{rx}}\) finite-rank compositions of
\(\mathcal{R}_\omega\) and \(\mathcal{O}_{\mathrm{meas}}\) Eq.~\eqref{eq:ch2.reconstruction-observation-map}
\cite{stratton1941,harrington2001,jackson1999}. Equation~\eqref{eq:ch2.transmit-receive-swap-preview}
does not eliminate \(\ker\mathcal{R}_\omega\) or \(\mathcal{K}_{\mathrm{nrec}}(N)\); it
constrains \emph{symmetric} coupling on observable coordinates after projection onto
\(\operatorname{ran}\mathcal{R}_\omega\) \cite{stratton1941}. Geometric forbidden sectors
Eq.~\eqref{eq:ch2.accessible-modal-sector} and angular budget
Eq.~\eqref{eq:ch2.angular-spectrum-composition-law} compose independently
\cite{harrington2001}.

\paragraph{Radiative coupling bilinear form (preview).}
Coupling between two radiating subsystems is tracked by
\begin{equation}
\label{eq:ch2.radiative-coupling-bilinear-preview}
\mathfrak{C}_\omega(\Jimp^{(1)},\Jimp^{(2)})
=
\langle \mathcal{R}_\omega[\Jimp^{(1)}],\,\mathcal{R}_\omega[\Jimp^{(2)}]\rangle_{\mathcal{U}}
\quad\text{on}\ \mathcal{F}_{\mathrm{rad}},
\end{equation}
with reciprocity implying \(\mathfrak{C}_\omega(\Jimp^{(1)},\Jimp^{(2)})=\mathfrak{C}_\omega(\Jimp^{(2)},\Jimp^{(1)})\)
under Eq.~\eqref{eq:ch2.lorentz-reciprocity-preview} and passive media
\cite{stratton1941,harrington2001}. Equation~\eqref{eq:ch2.radiative-coupling-bilinear-preview}
measures export overlap, not raw near-zone reactive overlap
\(\langle \mathcal{S}_\omega[\Jimp^{(1)}],\,\mathcal{S}_\omega[\Jimp^{(2)}]\rangle\), which may include
\(\Pi_{\mathrm{ev}}\) content annihilated by Eq.~\eqref{eq:ch2.rad-null-on-evanescent}
\cite{stratton1941}. Regime gate Eq.~\eqref{eq:ch2.regime-operator-screen} qualifies when
\(\mathfrak{C}_\omega\) may be reported as transport data \cite{harrington2001}.

\begin{figure}[t]
\centering
\ChOneFigBlock{%
\begin{tikzpicture}[ch1fig, node distance=7mm and 10mm]
  \node[ch1box] (lr) {Lorentz\\reciprocity};
  \node[ch1box, right=13mm of lr] (op) {Operator\\$\mathcal{S}_\omega,\mathcal{R}_\omega$};
  \node[ch1box, right=13mm of op] (mod) {Modal\\orthogonality};
  \node[ch1box, right=13mm of mod] (sys) {Radiation\\reception};
  \node[ch1box, below=12mm of op] (pi) {$\Pi_{\mathrm{rad}}$};
  \node[ch1box, below=12mm of mod] (bio) {Eq.~\eqref{eq:ch2.radiation-biorthogonality}};
  \draw[ch1flow] (lr) -- (op) -- (mod) -- (sys);
  \draw[ch1flow] (pi) -- (op);
  \draw[ch1flow] (bio) -- (mod);
\end{tikzpicture}%
}
\caption{Section~\ref{sec:ch28} development path: Lorentz reciprocity in operator form
(Subsection~\ref{sec:ch281}), orthogonality and coupling of radiative modes
(Subsection~\ref{sec:ch282}), and reciprocity in radiation and reception
(Subsection~\ref{sec:ch283}).}
\label{fig:ch2.reciprocity-roadmap}
\end{figure}

Figure~\ref{fig:ch2.reciprocity-roadmap} orders the three subsections.
Subsection~\ref{sec:ch281} gives the Chapter~\ref{ch:02} HOME for Lorentz reciprocity on
\(\mathcal{S}_\omega\) and \(\mathcal{R}_\omega\), upgrading
Eq.~\eqref{eq:ch2.dyadic-reciprocity-swap} and adjoint pairing
Eq.~\eqref{eq:ch2.maxwell-adjoint-pairing} to constructive statements on declared domains
\cite{stratton1941,harrington2001}. Subsection~\ref{sec:ch282} develops modal orthogonality
and coupling coefficients Eq.~\eqref{eq:ch2.modal-orthogonality-preview}, importing spectral
modes from Subsection~\ref{sec:ch261} and biorthogonality
Eq.~\eqref{eq:ch2.radiation-biorthogonality} by reference \cite{stratton1941}. Subsection
~\ref{sec:ch283} applies reciprocity to transmit--receive maps
Eq.~\eqref{eq:ch2.transmit-receive-swap-preview} and coupled antenna systems on
\(\operatorname{ran}\mathcal{R}_\omega\) \cite{harrington2001,jackson1999}.

\paragraph{Relation to null spaces and angular spectra.}
Reciprocity acts on \emph{observable} radiative coordinates: additions in
\(\ker\mathcal{R}_\omega\) leave \(\mathcal{R}_\omega[\Jimp]\) unchanged
Eq.~\eqref{eq:ch2.nr-superposition-law} and therefore leave coupling
Eq.~\eqref{eq:ch2.radiative-coupling-bilinear-preview} invariant on the range
\cite{stratton1941,harrington2001}. Finite-rank observation
Eq.~\eqref{eq:ch2.measurement-nullspace-composition} may break full reciprocity on raw
\(\mathcal{F}_{\mathrm{rad}}\) while preserving it on \(\operatorname{ran}\mathcal{O}_{\mathrm{meas}}\)
\cite{stratton1941}. Angular-spectrum consequences Theorem
~\ref{thm:ch2.angular-spectrum-null-consequences} fix which directions enter coupling; Section
~\ref{sec:ch29} develops plane-wave filters on the same accessible sector
Eq.~\eqref{eq:ch2.angular-accessibility-sector} \cite{harrington2001}.

Several constructions deliberately remain outside Section~\ref{sec:ch28}. Vector spherical
harmonic orthogonality and completeness are HOME in Chapter~\ref{ch:03} with USE of
Subsection~\ref{sec:ch282}; Subsection~\ref{sec:ch282} records operator-level orthogonality
without re-tabulating \(Y_\ell^m\) \cite{stratton1941,harrington2001}. Non-radiating catalogs
remain in Section~\ref{sec:ch27}; angular-spectrum operators in Section~\ref{sec:ch29}
\cite{stratton1941}. Section~\ref{sec:ch28} does not re-derive Fredholm structure, geometric
constraints, or null-space classification from Sections~\ref{sec:ch26}--\ref{sec:ch27}
\cite{harrington2001}.

Regime logic from Section~\ref{sec:ch01.2.4} is unchanged. Reciprocity of near-zone reactive
content does not promote evanescent storage to export through
Eq.~\eqref{eq:ch2.evanescent-nonexport-law} \cite{stratton1941,harrington2001}. Coupling
reports on \(\mathfrak{C}_\omega\) require \(\mathcal{D}_{J,\mathrm{reg}}=1\) and far-zone
qualification before flux meters Eqs.~\eqref{eq:ch2.power-flux-functional}--\eqref{eq:ch2.transport-triad-on-R}
interpret interchange as transport \cite{stratton1941}.

\paragraph{Worked illustration (probe-source swap).}
Interchanging impressed probe and source in Eq.~\eqref{eq:ch2.electric-dyadic-convolution}
leaves far-zone patterns unchanged only when Eq.~\eqref{eq:ch2.dyadic-reciprocity-swap} and
Eq.~\eqref{eq:ch2.lorentz-reciprocity-preview} hold on passive media
\cite{stratton1941,jackson1999}. Near-zone reactive scans without \(\Pi_{\mathrm{rad}}\)
can violate apparent reciprocity while far-zone reciprocity remains valid
\cite{harrington2001}.

\paragraph{Worked illustration (modal coupling symmetry).}
For exporting eigenpairs with \(\lambda_m,\lambda_n\neq 0\), coupling coefficients in
Eq.~\eqref{eq:ch2.modal-orthogonality-preview} inherit symmetry from reciprocity on
\(\mathcal{R}_\omega\) \cite{stratton1941,harrington2001}. Null modes
\(\lambda_n=0\) do not contribute to far coupling but may affect near reactive superpositions
Eq.~\eqref{eq:ch2.modal-null-sector-reduction} \cite{stratton1941}.

\paragraph{Worked illustration (two-port antenna).}
A transmit--receive pair with shared \(\mathcal{R}_\omega\) satisfies
Eq.~\eqref{eq:ch2.transmit-receive-swap-preview} on observable ports when media are passive
and reciprocal \cite{stratton1941,jackson1999}. Adding \(\Jimp^{\mathrm{(NR)}}\) to either
port leaves port coupling unchanged on \(\operatorname{ran}\mathcal{R}_\omega\)
Theorem~\ref{thm:ch2.canonical-nr-current-families} \cite{harrington2001}.

\paragraph{Worked illustration (lossy media).}
Loss modifies phases in Eq.~\eqref{eq:ch2.radiation-eigenpair-def} and may break strict
time-reversal symmetry while preserving partial reciprocity relations stated in
Subsection~\ref{sec:ch281} \cite{stratton1941,harrington2001}. Export selection
Eq.~\eqref{eq:ch2.rad-null-on-evanescent} is unchanged: loss does not convert evanescent
storage to far-zone reciprocity channels \cite{stratton1941}.

The subsection sequence begins with Lorentz reciprocity in operator form on
\(\mathcal{S}_\omega\) and \(\mathcal{R}_\omega\), continues with orthogonality and coupling
of radiative modes on spectral coordinates from Section~\ref{sec:ch26}, and closes with
reciprocity in radiation and reception before angular-spectrum action in Section~\ref{sec:ch29}
and explicit harmonic orthogonality in Chapter~\ref{ch:03} compose additional structure on the
same modal coordinates \cite{stratton1941,harrington2001,jackson1999}.

\subsection{Lorentz reciprocity in operator form}
\label{sec:ch281}
% TOC tag: [HOME]

Section~\ref{sec:ch28} previewed bilinear exchange symmetry through
Eqs.~\eqref{eq:ch2.section28-reciprocity-chain}--\eqref{eq:ch2.lorentz-reciprocity-preview}.
Subsection~\ref{sec:ch281} gives the Chapter~\ref{ch:02} HOME for Lorentz reciprocity on
\(\mathcal{S}_\omega\) and \(\mathcal{R}_\omega\), including the full proof of dyadic swap
Eq.~\eqref{eq:ch2.dyadic-reciprocity-swap} deferred from Subsection~\ref{sec:ch231}
\cite{stratton1941,harrington2001,jackson1999}. Maxwell adjoint pairing
Eq.~\eqref{eq:ch2.maxwell-adjoint-pairing}, source dual pairing
Eq.~\eqref{eq:ch2.dual-source-pairing}, outgoing closure
Eqs.~\eqref{eq:ch2.outgoing-radiation-refinement}--\eqref{eq:ch2.outgoing-resolvent-unique},
and source-to-field integrals Theorem~\ref{thm:ch2.source-to-field-integral} are presupposed;
Green construction and Sommerfeld uniqueness are not re-opened \cite{stratton1941}. Modal
orthogonality and coupling belong to Subsection~\ref{sec:ch282}; transmit--receive systems to
Subsection~\ref{sec:ch283} \cite{harrington2001}.

\paragraph{Assumptions.}
Fix \(\omega>0\), Lipschitz domain \(\Omega\) with piecewise smooth interfaces, impressed
sources in \(\mathcal{J}_\Omega\) Eq.~\eqref{eq:ch2.source-space-refined} satisfying
Eq.~\eqref{eq:ch2.source-admissibility}, and outgoing fields in
\(\mathcal{F}_{\mathrm{SF}}\) Eq.~\eqref{eq:ch2.admissible-field-space}
\cite{stratton1941,harrington2001}. Constitutive tensors are linear, local, and passive at
\(\omega\), with symmetric permittivity and permeability in the sense
\begin{equation}
\label{eq:ch2.reciprocal-medium-class}
\overline{\overline{\varepsilon}}(\mathbf{r})=\overline{\overline{\varepsilon}}^{\mathsf{T}}(\mathbf{r}),
\qquad
\overline{\overline{\mu}}(\mathbf{r})=\overline{\overline{\mu}}^{\mathsf{T}}(\mathbf{r}),
\qquad
\overline{\overline{\xi}}=\overline{\overline{\zeta}}=\mathbf{0},
\end{equation}
where \(\overline{\overline{\xi}}\), \(\overline{\overline{\zeta}}\) denote magneto-electric
coupling blocks in the constitutive relation imported from Section~\ref{sec:ch01.1.2}
\cite{stratton1941,jackson1999}. Boundary operator \(\mathcal{B}_\Gamma\) is passive
Eq.~\eqref{eq:ch2.passive-impedance-class} and identical for both excitations
\cite{stratton1941,harrington2001}. Field pairing \(\langle\cdot,\cdot\rangle_{\mathcal{U}}\)
is induced by Eq.~\eqref{eq:ch2.field-energy-norm} on admissible Maxwell pairs
\cite{stratton1941}. Regime gate Eq.~\eqref{eq:ch2.regime-operator-screen} governs transport
interpretation of any bilinear quantity built from \(\Pi_{\mathrm{rad}}\)
\cite{harrington2001}.

\paragraph{Reciprocal medium class.}
Equation~\eqref{eq:ch2.reciprocal-medium-class} is the Chapter~\ref{ch:02} admissibility
class for \emph{full} Lorentz reciprocity at fixed \(\omega\): symmetric \(\varepsilon\),
\(\mu\) ensure time-reversal symmetry of the harmonic Maxwell system modulo outgoing
selection, while vanishing magneto-electric blocks exclude non-reciprocal gyrotropic or
Tellegen couplings \cite{stratton1941,jackson1999}. Isotropic lossless media are a special
case with scalar \(\varepsilon\), \(\mu\). Lossy symmetric tensors preserve reciprocity
relations stated below with complex conjugation on one slot of each pairing
\cite{stratton1941,harrington2001}. Active or negative-real impedance portions of
\(\mathcal{B}_\Gamma\) break passive symmetry and are excluded from the present theorem
\cite{stratton1941}.

\paragraph{Volume Lorentz identity.}
Let \((\E_1,\H_1)=\mathcal{S}_\omega[\Jimp^{(1)}]\) and
\((\E_2,\H_2)=\mathcal{S}_\omega[\Jimp^{(2)}]\) be outgoing admissible solutions for
independent impressed currents \(\Jimp^{(1)},\Jimp^{(2)}\in\mathcal{J}_\Omega\)
\cite{stratton1941}. The classical Lorentz reciprocity relation reads
\begin{equation}
\label{eq:ch2.lorentz-volume-reciprocity}
\int_\Omega
\big(\E_1\cdot\Jimp^{(2)*}-\E_2\cdot\Jimp^{(1)*}\big)\,d^{3}r
=
\int_\Gamma
\big(\hat{\mathbf{n}}\times\E_1\cdot\H_2^{*}
-\hat{\mathbf{n}}\times\E_2\cdot\H_1^{*}\big)\,dS,
\end{equation}
when \(\overline{\overline{\varepsilon}}\), \(\overline{\overline{\mu}}\) satisfy
Eq.~\eqref{eq:ch2.reciprocal-medium-class} and \(\mathcal{B}_\Gamma\) is passive
\cite{stratton1941,harrington2001,jackson1999}. Equation~\eqref{eq:ch2.lorentz-volume-reciprocity}
balances impressed-source bilinear forms against a boundary pairing generated by
Eq.~\eqref{eq:ch2.maxwell-adjoint-pairing}: the left member measures how each field couples
to the other's source, while the right member is the surface term that remains after
integration by parts on the Maxwell operator \cite{stratton1941}. On domains where
\(\mathcal{B}_\Gamma\) annihilates the boundary integrand for admissible outgoing pairs---PEC,
homogeneous impedance with Eq.~\eqref{eq:ch2.passive-impedance-class}, and standard
Sommerfeld exterior closures---the surface term vanishes and volume reciprocity closes
\cite{stratton1941,harrington2001}.

\paragraph{Derivation sketch (volume form).}
Substitute \(\mathbf{u}_\omega^{(1)}=(\E_1,\H_1)\) and \(\mathbf{u}_\omega^{(2)}=(\E_2,\H_2)\)
into Eq.~\eqref{eq:ch2.maxwell-adjoint-pairing} and interchange the two states
\cite{stratton1941}. Symmetry of \(\overline{\overline{\varepsilon}}\),
\(\overline{\overline{\mu}}\) under Eq.~\eqref{eq:ch2.reciprocal-medium-class} makes the
volume pairing sesquilinear form on \(\mathcal{M}_\omega[\mathbf{u}_\omega]\) invariant under
label swap when sources are impressed currents in the Amp\`ere row
\cite{stratton1941,jackson1999}. Subtracting the two adjoint identities eliminates
\(\langle\mathcal{M}_\omega[\mathbf{u}_\omega^{(1)}],\mathbf{u}_\omega^{(2)}\rangle_{\Omega}\)
and yields Eq.~\eqref{eq:ch2.lorentz-volume-reciprocity} with the boundary term displayed
\cite{stratton1941}. Charge continuity Eq.~\eqref{eq:ch1.charge-continuity-harmonic} and
compatibility Eq.~\eqref{eq:ch2.source-compatibility} ensure that impressed \(\Jimp\) alone
specifies the electric drive when \(\rho\) is tied to \(\Jimp\)
\cite{stratton1941,harrington2001}.

\begin{figure}[t]
\centering
\ChOneFigBlock{%
\begin{tikzpicture}[ch1fig, node distance=7mm and 10mm]
  \node[ch1box] (j1) {$\Jimp^{(1)}$};
  \node[ch1box, right=12mm of j1] (s) {$\mathcal{S}_\omega$};
  \node[ch1box, right=12mm of s] (u1) {$\mathbf{u}_1$};
  \node[ch1box, below=11mm of j1] (j2) {$\Jimp^{(2)}$};
  \node[ch1box, below=11mm of s] (u2) {$\mathbf{u}_2$};
  \node[ch1box, below=11mm of u1] (bil) {Bilinear\\pairing};
  \draw[ch1flow] (j1) -- (s) -- (u1);
  \draw[ch1flow] (j2) -- (s) -- (u2);
  \draw[ch1flow, dashed] (u1) -- (bil);
  \draw[ch1flow, dashed] (j2) -- (bil);
  \draw[ch1flow, dashed] (u2) -- (bil);
  \draw[ch1flow, dashed] (j1) -- (bil);
\end{tikzpicture}%
}
\caption{Subsection~\ref{sec:ch281} Lorentz reciprocity: source labels interchange under
\(\mathcal{S}_\omega\) when media satisfy Eq.~\eqref{eq:ch2.reciprocal-medium-class} and
\(\mathcal{B}_\Gamma\) is passive on the outgoing branch.}
\label{fig:ch2.lorentz-reciprocity-pipeline}
\end{figure}

Figure~\ref{fig:ch2.lorentz-reciprocity-pipeline} is the operator reading of
Eq.~\eqref{eq:ch2.lorentz-volume-reciprocity}: the bilinear functional is symmetric under
\((\Jimp^{(1)},\mathbf{u}_1)\leftrightarrow(\Jimp^{(2)},\mathbf{u}_2)\) once boundary terms
cancel \cite{stratton1941,harrington2001}.

\begin{theorem}[Lorentz reciprocity for the solution map]
\label{thm:ch2.lorentz-reciprocity-operator}
Let \(\mathcal{S}_\omega:\mathcal{J}_\Omega\to\mathcal{F}_\Omega\) be the outgoing solution
map Eq.~\eqref{eq:ch2.maxwell-solution-map} on media satisfying
Eq.~\eqref{eq:ch2.reciprocal-medium-class} with passive \(\mathcal{B}_\Gamma\)
Eq.~\eqref{eq:ch2.passive-impedance-class}. For admissible
\(\Jimp^{(1)},\Jimp^{(2)}\in\mathcal{J}_\Omega\), write
\(\mathbf{u}_k=\mathcal{S}_\omega[\Jimp^{(k)}]\), \(k=1,2\). Then:
\begin{enumerate}
\item[(i)] Volume reciprocity Eq.~\eqref{eq:ch2.lorentz-volume-reciprocity} holds with
vanishing boundary term on admissible outgoing closures.
\item[(ii)] Operator form:
\begin{equation}
\label{eq:ch2.lorentz-reciprocity-operator}
\langle \mathbf{u}_1,\,\mathcal{S}_\omega[\Jimp^{(2)}]\rangle_{\mathcal{U}}
=
\langle \mathbf{u}_2,\,\mathcal{S}_\omega[\Jimp^{(1)}]\rangle_{\mathcal{U}}
\end{equation}
when the pairings are defined through impressed-current forcing functionals
Eq.~\eqref{eq:ch2.forcing-functional-operator} and field norm
Eq.~\eqref{eq:ch2.field-energy-norm} on \(\mathcal{F}_{\mathrm{SF}}\).
\item[(iii)] Dyadic realization: outgoing electric dyads satisfy
Eq.~\eqref{eq:ch2.dyadic-reciprocity-swap}.
\item[(iv)] Radiation reciprocity:
\begin{equation}
\label{eq:ch2.radiation-map-reciprocity}
\langle \mathcal{R}_\omega[\Jimp^{(1)}],\,\mathcal{R}_\omega[\Jimp^{(2)}]\rangle_{\mathcal{U}}
=
\langle \mathcal{R}_\omega[\Jimp^{(2)}],\,\mathcal{R}_\omega[\Jimp^{(1)}]\rangle_{\mathcal{U}}
\end{equation}
on \(\mathcal{F}_{\mathrm{rad}}\) with \(\mathcal{R}_\omega=\Pi_{\mathrm{rad}}\circ\mathcal{S}_\omega\)
Eq.~\eqref{eq:ch2.radiation-operator-def}.
\end{enumerate}
\end{theorem}
\begin{proof}
Part (i) follows from Eq.~\eqref{eq:ch2.maxwell-adjoint-pairing} and
Eq.~\eqref{eq:ch2.reciprocal-medium-class} as in the derivation above; outgoing and passive
boundary classes from Subsection~\ref{sec:ch214} cancel the surface integrand
\cite{stratton1941,harrington2001}. Part (ii) identifies the volume impressed-current pairing
with \(\langle\cdot,\cdot\rangle_{\mathcal{U}}\) on \(\mathcal{F}_{\mathrm{SF}}\) through
Eq.~\eqref{eq:ch2.dual-source-pairing} and stability
Eq.~\eqref{eq:ch2.solution-map-stability}, upgrading preview
Eq.~\eqref{eq:ch2.lorentz-reciprocity-preview} \cite{stratton1941}. Part (iii): for fixed
\(\mathbf{r}\), Eq.~\eqref{eq:ch2.electric-dyadic-convolution} and part (i) with impulsive
\(\Jimp^{(k)}\) aligned with Cartesian unit vectors yield column swap
\(\overline{\mathbf{G}}_{E}(\mathbf{r},\mathbf{r}')=
\overline{\mathbf{G}}_{E}^{\mathsf{T}}(\mathbf{r}',\mathbf{r})\) on the outgoing branch
\cite{stratton1941,jackson1999}. Uniqueness Theorem~\ref{thm:ch2.exterior-radiative-uniqueness}
extends the swap to all admissible \(\mathbf{r},\mathbf{r}'\) \cite{stratton1941}. Part (iv)
applies \(\Pi_{\mathrm{rad}}\) to both sides of (ii); \(\Pi_{\mathrm{rad}}\) is self-adjoint
on \(\mathcal{F}_{\mathrm{rad}}\) under Eq.~\eqref{eq:ch2.field-energy-norm} when outgoing
selection is fixed \cite{harrington2001,stratton1941}.
\end{proof}

\paragraph{Mathematical role of Eq.~\eqref{eq:ch2.lorentz-reciprocity-operator}.}
Equation~\eqref{eq:ch2.lorentz-reciprocity-operator} states that the bilinear functional
\(\mathfrak{B}_\omega(\Jimp^{(1)},\Jimp^{(2)})=
\langle \mathcal{S}_\omega[\Jimp^{(1)}],\,\mathcal{S}_\omega[\Jimp^{(2)}]\rangle_{\mathcal{U}}\)
is symmetric on \(\mathcal{J}_\Omega\times\mathcal{J}_\Omega\) under reciprocal media
\cite{stratton1941,harrington2001}. Algebraically, \(\mathcal{S}_\omega\) is not self-adjoint
on \(\mathcal{J}_\Omega\) with source pairing Eq.~\eqref{eq:ch2.dual-source-pairing}; reciprocity
instead couples the source space to the field space through the \emph{interchanged} labels in
Eq.~\eqref{eq:ch2.lorentz-reciprocity-operator} \cite{stratton1941}. Physically, the far-zone
pattern produced by \(\Jimp^{(1)}\) tested against the field of \(\Jimp^{(2)}\) equals the
swap configuration: a receive pattern computed with \(\Jimp^{(2)}\) as probe matches a
transmit map driven by \(\Jimp^{(1)}\) when hardware and media are reciprocal
\cite{harrington2001,jackson1999}.

\paragraph{Mathematical role of Eq.~\eqref{eq:ch2.radiation-map-reciprocity}.}
Equation~\eqref{eq:ch2.radiation-map-reciprocity} restricts Eq.~\eqref{eq:ch2.lorentz-reciprocity-operator}
to \(\mathcal{F}_{\mathrm{rad}}\), the subspace on which flux meters and angular patterns are
declared \cite{stratton1941,harrington2001}. Because \(\Pi_{\mathrm{rad}}\) annihilates
non-export components Eq.~\eqref{eq:ch2.rad-null-on-evanescent}, reactive near-zone overlap
does not enter Eq.~\eqref{eq:ch2.radiation-map-reciprocity} even when present in full
Eq.~\eqref{eq:ch2.lorentz-reciprocity-operator} \cite{stratton1941}. Regime gate
Eq.~\eqref{eq:ch2.regime-component-gate} must still be satisfied before interpreting
Eq.~\eqref{eq:ch2.radiation-map-reciprocity} as a transport channel
\cite{harrington2001}.

\begin{proposition}[Dyadic reciprocity swap (HOME proof)]
\label{prop:ch2.dyadic-reciprocity-home}
Under Eq.~\eqref{eq:ch2.reciprocal-medium-class} with outgoing closure, the electric dyad
\(\overline{\mathbf{G}}_{E,\omega}^{\mathrm{(out)}}\) in
Eq.~\eqref{eq:ch2.electric-dyadic-convolution} satisfies
Eq.~\eqref{eq:ch2.dyadic-reciprocity-swap} for all \(\mathbf{r},\mathbf{r}'\) in the
evaluation domain of Theorem~\ref{thm:ch2.source-to-field-integral}.
\end{proposition}
\begin{proof}
Apply Theorem~\ref{thm:ch2.lorentz-reciprocity-operator}(iii); this is the proof deferred from
Subsection~\ref{sec:ch231} \cite{stratton1941,jackson1999}.
\end{proof}

\begin{proposition}[Null-space invariance of radiative reciprocity]
\label{prop:ch2.reciprocity-null-invariance}
If \(\Jimp^{(1)}\) is arbitrary and \(\Jimp^{(2)}=\Jimp^{(1)}+\Jimp^{\mathrm{(NR)}}\) with
\(\Jimp^{\mathrm{(NR)}}\in\ker\mathcal{R}_\omega\), then
Eq.~\eqref{eq:ch2.radiation-map-reciprocity} is unchanged when either argument is replaced by
\(\Jimp^{(2)}\) on the radiative pairing
\cite{stratton1941,harrington2001}.
\end{proposition}
\begin{proof}
\(\mathcal{R}_\omega[\Jimp^{(2)}]=\mathcal{R}_\omega[\Jimp^{(1)}]\) by
Eq.~\eqref{eq:ch2.nr-superposition-law} \cite{stratton1941}. Full
Eq.~\eqref{eq:ch2.lorentz-reciprocity-operator} may change by reactive content through
\(\mathcal{S}_\omega[\Jimp^{\mathrm{(NR)}}]\) Eq.~\eqref{eq:ch2.nr-image-of-null-space}
\cite{harrington2001}.
\end{proof}

\paragraph{Resolvent and integral equivalence.}
On the outgoing branch, Eq.~\eqref{eq:ch2.lorentz-reciprocity-operator} is consistent with
resolvent symmetry of the truncated Maxwell operator on reciprocal media:
\begin{equation}
\label{eq:ch2.resolvent-reciprocity-swap}
\langle
\mathcal{L}_\omega^{-1}[\mathbf{s}^{(1)}],\,\mathbf{s}^{(2)}
\rangle_{\mathcal{Y}}
=
\langle
\mathcal{L}_\omega^{-1}[\mathbf{s}^{(2)}],\,\mathbf{s}^{(1)}
\rangle_{\mathcal{Y}}
\end{equation}
for compatible forcing pairs \(\mathbf{s}^{(k)}\) when \(\mathcal{L}_\omega\) is symmetric on
\(\mathcal{D}(\mathcal{L}_\omega)\) and \(\mathcal{B}_\Gamma\) is passive
\cite{stratton1941,harrington2001}. Equation~\eqref{eq:ch2.resolvent-reciprocity-swap} is not
re-proved here; it is the operator backbone behind
Eq.~\eqref{eq:ch2.integral-resolvent-equivalence} and Theorem
~\ref{thm:ch2.lorentz-reciprocity-operator}(ii) \cite{stratton1941}. Numerical MoM matrices
assembled from Eq.~\eqref{eq:ch2.dyadic-reciprocity-swap} inherit swap symmetry only when
quadrature, outgoing branch, and material sampling preserve
Eq.~\eqref{eq:ch2.reciprocal-medium-class} \cite{harrington2001}.

\paragraph{Loss, anisotropy, and broken reciprocity.}
Complex symmetric \(\overline{\overline{\varepsilon}}\), \(\overline{\overline{\mu}}\) with
\(\Im\{\overline{\overline{\varepsilon}}\}\succeq 0\), \(\Im\{\overline{\overline{\mu}}\}\succeq 0\)
preserve Eqs.~\eqref{eq:ch2.lorentz-volume-reciprocity}--\eqref{eq:ch2.lorentz-reciprocity-operator}
with conjugation on one slot \cite{stratton1941}. Magneto-electric coupling
\(\overline{\overline{\xi}}\neq\mathbf{0}\) breaks Eq.~\eqref{eq:ch2.dyadic-reciprocity-swap}
and requires non-reciprocal Green operators \cite{stratton1941,jackson1999}. Active boundaries
with \(\Re\{Z_s\}<0\) on portions of \(\Gamma_Z\) violate
Eq.~\eqref{eq:ch2.passive-impedance-class} and invalidate Theorem
~\ref{thm:ch2.lorentz-reciprocity-operator} \cite{stratton1941,harrington2001}. Incoming rather
than outgoing closure selects a different branch; swap identities hold only when both excitations
use the same branch \cite{stratton1941}.

\paragraph{Boundary of validity.}
Theorem~\ref{thm:ch2.lorentz-reciprocity-operator} assumes fixed \(\omega\), linear media
Eq.~\eqref{eq:ch2.reciprocal-medium-class}, passive \(\mathcal{B}_\Gamma\), outgoing closure,
and admissible \(\mathcal{J}_\Omega\) \cite{stratton1941,harrington2001}. It does not assert
reciprocity of raw near-zone reactive scans without specifying whether \(\Pi_{\mathrm{rad}}\)
has been applied \cite{stratton1941}. It does not remove \(\ker\mathcal{R}_\omega\)
Proposition~\ref{prop:ch2.reciprocity-null-invariance}; invisible currents alter full
\(\mathcal{S}_\omega\) pairings while leaving radiative reciprocity unchanged
\cite{harrington2001}. Nonlinear, time-varying, or moving media require the causal chain of
Subsection~\ref{sec:ch223} \cite{stratton1941}.

\paragraph{Worked illustration (probe-source swap).}
In Eq.~\eqref{eq:ch2.electric-dyadic-convolution}, interchanging observation point and impulsive
probe location leaves far-zone patterns unchanged when
Eq.~\eqref{eq:ch2.dyadic-reciprocity-swap} holds \cite{stratton1941,jackson1999}. Near-zone
reactive scans can appear non-reciprocal while Eq.~\eqref{eq:ch2.radiation-map-reciprocity}
remains valid on \(\mathcal{F}_{\mathrm{rad}}\) \cite{harrington2001}.

\paragraph{Worked illustration (lossy dielectric).}
A lossy but symmetric \(\overline{\overline{\varepsilon}}\) slab preserves Theorem
~\ref{thm:ch2.lorentz-reciprocity-operator} with complex conjugation in pairings; phase shifts
in \(\mathcal{S}_\omega[\Jimp]\) do not break swap symmetry \cite{stratton1941}. Evanescent
storage remains non-exporting through Eq.~\eqref{eq:ch2.rad-null-on-evanescent}
\cite{stratton1941,harrington2001}.

\paragraph{Worked illustration (NR superposition).}
Adding \(\Jimp^{\mathrm{(NR)}}\in\ker\mathcal{R}_\omega\) to a radiating source leaves
Eq.~\eqref{eq:ch2.radiation-map-reciprocity} unchanged by Proposition
~\ref{prop:ch2.reciprocity-null-invariance} and Eq.~\eqref{eq:ch2.nr-superposition-law}
\cite{stratton1941,harrington2001}. Two-port coupling reported on
\(\operatorname{ran}\mathcal{R}_\omega\) is therefore insensitive to invisible loop content
\cite{stratton1941}.

\paragraph{Worked illustration (MoM symmetry).}
A Galerkin impedance matrix \(Z_{mn}=\langle \mathbf{w}_m,\,\mathcal{S}_\omega[\Jimp^{(n)}]\rangle\)
inherits \(Z_{mn}=Z_{nm}\) when expansion functions \(\Jimp^{(n)}\) lie in
\(\mathcal{J}_\Omega\), media satisfy Eq.~\eqref{eq:ch2.reciprocal-medium-class}, and outgoing
dyads obey Eq.~\eqref{eq:ch2.dyadic-reciprocity-swap} \cite{harrington2001,stratton1941}.
Broken symmetry in computed \(Z_{mn}\) usually signals wrong branch selection, non-reciprocal
material data, or inconsistent \(\mathcal{B}_\Gamma\) between transmit and receive models
\cite{harrington2001}.

\paragraph{Closing summary.}
Subsection~\ref{sec:ch281} establishes the Chapter~\ref{ch:02} HOME for reciprocal medium class
Eq.~\eqref{eq:ch2.reciprocal-medium-class}, volume Lorentz identity
Eq.~\eqref{eq:ch2.lorentz-volume-reciprocity}, operator reciprocity
Eq.~\eqref{eq:ch2.lorentz-reciprocity-operator}, dyadic swap
Proposition~\ref{prop:ch2.dyadic-reciprocity-home} upgrading
Eq.~\eqref{eq:ch2.dyadic-reciprocity-swap}, radiation reciprocity
Eq.~\eqref{eq:ch2.radiation-map-reciprocity}, null-space invariance
Proposition~\ref{prop:ch2.reciprocity-null-invariance}, and resolvent consistency
Eq.~\eqref{eq:ch2.resolvent-reciprocity-swap}, packaged in Theorem
~\ref{thm:ch2.lorentz-reciprocity-operator} \cite{stratton1941,harrington2001,jackson1999}.
Subsection~\ref{sec:ch282} develops modal orthogonality and coupling on spectral coordinates
from Subsection~\ref{sec:ch261}; Subsection~\ref{sec:ch283} applies reciprocity to radiation and
reception systems \cite{stratton1941}.
\subsection{Orthogonality and coupling of radiative modes}
\label{sec:ch282}
% TOC tag: [HOME]

Subsection~\ref{sec:ch281} established Lorentz reciprocity on \(\mathcal{S}_\omega\) and
\(\mathcal{R}_\omega\) through Theorem~\ref{thm:ch2.lorentz-reciprocity-operator} and
radiation reciprocity Eq.~\eqref{eq:ch2.radiation-map-reciprocity}
\cite{stratton1941,harrington2001}. Subsection~\ref{sec:ch282} gives the Chapter~\ref{ch:02}
HOME for orthogonality and coupling of radiative modes on the spectral coordinates fixed in
Subsection~\ref{sec:ch261}, upgrading preview
Eqs.~\eqref{eq:ch2.modal-orthogonality-preview} and~\eqref{eq:ch2.radiative-coupling-bilinear-preview}
\cite{stratton1941}. Radiation eigenpairs Eq.~\eqref{eq:ch2.radiation-eigenpair-def}, modal
projectors Eq.~\eqref{eq:ch2.modal-projector-def}, source amplitudes
Eq.~\eqref{eq:ch2.source-modal-amplitude}, spectral decomposition Theorem
~\ref{thm:ch2.radiation-spectral-decomposition}, and biorthogonality
Eq.~\eqref{eq:ch2.radiation-biorthogonality} are presupposed; eigenpair construction and
Fredholm solvability are not re-proved \cite{stratton1941,harrington2001}. Explicit vector
spherical harmonic orthogonality belongs to Chapter~\ref{ch:03}
Subsection~\ref{sec:ch343}; continuous-spectrum normalization belongs to
Subsection~\ref{sec:ch352}, both with USE of the present operator targets
\cite{stratton1941,harrington2001}. Transmit--receive systems are deferred to Subsection~\ref{sec:ch283}
\cite{harrington2001}.

\paragraph{Assumptions.}
Fix \(\omega>0\), outgoing fields in \(\mathcal{F}_{\mathrm{SF}}\), radiative images in
\(\mathcal{F}_{\mathrm{rad}}\), and orthonormal radiation eigenmodes
\(\{\mathbf{u}_n\}_{n\in\mathcal{N}_{\mathrm{rad}}}\) from Theorem
~\ref{thm:ch2.radiation-spectral-decomposition} on the declared domain class
\cite{stratton1941,harrington2001}. Media satisfy reciprocal class
Eq.~\eqref{eq:ch2.reciprocal-medium-class} and passive boundary
Eq.~\eqref{eq:ch2.passive-impedance-class} so Theorem
~\ref{thm:ch2.lorentz-reciprocity-operator} applies \cite{stratton1941}. Adjoint modes
\(\widetilde{\mathbf{u}}_n\) are defined through
Eq.~\eqref{eq:ch2.radiation-biorthogonality} and coincide with \(\mathbf{u}_n\) on
\(\mathcal{F}_{\mathrm{rad}}\) under reciprocity-qualified pairings below
\cite{stratton1941,harrington2001}. Regime gate Eq.~\eqref{eq:ch2.regime-operator-screen}
governs transport interpretation of coupling reports \cite{harrington2001}.

\paragraph{Radiation-modal orthogonality (HOME).}
On the orthonormal family of Theorem~\ref{thm:ch2.radiation-spectral-decomposition},
\begin{equation}
\label{eq:ch2.radiation-modal-orthogonality}
\langle \mathbf{u}_m,\,\mathbf{u}_n\rangle_{\mathcal{U}}
=
\delta_{mn},
\qquad
m,n\in\mathcal{N}_{\mathrm{rad}},
\end{equation}
with \(\|\mathbf{u}_n\|_{\mathcal{U}}=1\) already enforced in
Eq.~\eqref{eq:ch2.radiation-eigenpair-def} \cite{stratton1941,harrington2001}.
Equation~\eqref{eq:ch2.radiation-modal-orthogonality} is the Subsection~\ref{sec:ch282} HOME
upgrade of the first line of preview Eq.~\eqref{eq:ch2.modal-orthogonality-preview}; it
specializes projector orthogonality \(P_m P_n=\delta_{mn}P_n\) from
Eq.~\eqref{eq:ch2.modal-projector-def} to the chosen orthonormal basis
\cite{stratton1941}. Algebraically, Eq.~\eqref{eq:ch2.radiation-modal-orthogonality} makes
\(\{\mathbf{u}_n\}\) an orthonormal frame on \(\mathcal{F}_{\mathrm{rad}}\) for energy inner
product Eq.~\eqref{eq:ch2.field-energy-norm}; physically, distinct exporting modes carry
non-overlapping far-zone energy content on \(S^{2}\) when
\(\mathcal{D}_{J,\mathrm{reg}}=1\) \cite{jackson1999,harrington2001}. Explicit angular
integrals producing Eq.~\eqref{eq:ch2.radiation-modal-orthogonality} for harmonic bases are
not tabulated here \cite{stratton1941}.

\paragraph{Modal coupling coefficients.}
For canonical eigen-sources \(\Jimp^{(n)}\) in Eq.~\eqref{eq:ch2.radiation-eigenpair-def},
define the \emph{modal coupling coefficient}
\begin{equation}
\label{eq:ch2.modal-coupling-coefficient}
C_{mn}
=
\langle \widetilde{\mathbf{u}}_m,\,\mathcal{R}_\omega[\Jimp^{(n)}]\rangle_{\mathcal{U}},
\qquad
m,n\in\mathcal{N}_{\mathrm{rad}},
\end{equation}
with adjoint modes from Eq.~\eqref{eq:ch2.radiation-biorthogonality}
\cite{stratton1941,harrington2001}. Equation~\eqref{eq:ch2.modal-coupling-coefficient} is the
Subsection~\ref{sec:ch282} HOME upgrade of the second line of preview
Eq.~\eqref{eq:ch2.modal-orthogonality-preview}. When
Eq.~\eqref{eq:ch2.radiation-biorthogonality} holds,
\begin{equation}
\label{eq:ch2.modal-coupling-diagonal}
C_{mn}=\lambda_n\,\delta_{mn}
\end{equation}
for exporting eigenpairs with \(\lambda_n\neq 0\); null modes
\(n\in\mathcal{N}_{\mathrm{null}}\) Eq.~\eqref{eq:ch2.modal-null-index-set} contribute
\(C_{mn}=0\) for all \(m\) \cite{stratton1941,harrington2001}. Algebraically,
Eq.~\eqref{eq:ch2.modal-coupling-diagonal} diagonalizes coupling in an eigenmode basis;
physically, orthogonal radiation modes do not co-excite one another through
\(\mathcal{R}_\omega\) when sources align with eigenpairs---cross-talk appears only from
superpositions or degenerate subspaces \cite{harrington2001,jackson1999}.

\paragraph{Radiative coupling bilinear form (HOME).}
For arbitrary admissible \(\Jimp^{(1)},\Jimp^{(2)}\in\mathcal{J}_\Omega\), define
\begin{equation}
\label{eq:ch2.radiative-coupling-bilinear}
\mathfrak{C}_\omega\!\big(\Jimp^{(1)},\Jimp^{(2)}\big)
=
\langle \mathcal{R}_\omega[\Jimp^{(1)}],\,\mathcal{R}_\omega[\Jimp^{(2)}]\rangle_{\mathcal{U}}
\quad\text{on}\ \mathcal{F}_{\mathrm{rad}},
\end{equation}
upgrading preview Eq.~\eqref{eq:ch2.radiative-coupling-bilinear-preview}
\cite{stratton1941,harrington2001}. Equation~\eqref{eq:ch2.radiative-coupling-bilinear}
measures export overlap on \(\operatorname{ran}\mathcal{R}_\omega\), not raw reactive overlap
of \(\mathcal{S}_\omega[\Jimp^{(k)}]\) \cite{stratton1941}. Under Theorem
~\ref{thm:ch2.lorentz-reciprocity-operator}(iv),
\begin{equation}
\label{eq:ch2.coupling-bilinear-symmetry}
\mathfrak{C}_\omega\!\big(\Jimp^{(1)},\Jimp^{(2)}\big)
=
\mathfrak{C}_\omega\!\big(\Jimp^{(2)},\Jimp^{(1)}\big)
\end{equation}
on reciprocal media \cite{stratton1941,harrington2001}. Reactive near-zone terms in
\(\Pi_{\mathrm{ev}}\mathcal{S}_\omega[\Jimp]\) do not enter
Eq.~\eqref{eq:ch2.radiative-coupling-bilinear} through
Eq.~\eqref{eq:ch2.rad-null-on-evanescent} \cite{stratton1941}.

\paragraph{Modal expansion of coupling.}
Let \(c_n[\Jimp^{(k)}]\) denote source amplitudes Eq.~\eqref{eq:ch2.source-modal-amplitude}
for \(k=1,2\). When the modal sum converges in \(\mathcal{U}\),
\begin{equation}
\label{eq:ch2.modal-coupling-sum}
\mathfrak{C}_\omega\!\big(\Jimp^{(1)},\Jimp^{(2)}\big)
=
\sum_{m,n\in\mathcal{N}_{\mathrm{rad}}}
\lambda_m\lambda_n\,
c_m[\Jimp^{(1)}]^{*}\,
c_n[\Jimp^{(2)}]\,
\langle \mathbf{u}_m,\mathbf{u}_n\rangle_{\mathcal{U}},
\end{equation}
using Eq.~\eqref{eq:ch2.radiation-map-modal-sum} and linearity
Eq.~\eqref{eq:ch2.radiation-linearity} \cite{stratton1941,harrington2001}. With
Eq.~\eqref{eq:ch2.radiation-modal-orthogonality},
\begin{equation}
\label{eq:ch2.modal-coupling-sum-diagonal}
\mathfrak{C}_\omega\!\big(\Jimp^{(1)},\Jimp^{(2)}\big)
=
\sum_{n\in\mathcal{N}_{\mathrm{rad}}}
|\lambda_n|^{2}\,
c_n[\Jimp^{(1)}]^{*}\,
c_n[\Jimp^{(2)}],
\end{equation}
for real nonnegative exporting gains \(\lambda_n>0\) \cite{stratton1941}. Equation
\eqref{eq:ch2.modal-coupling-sum-diagonal} shows that coupling is a weighted inner product
of modal excitation vectors on the exporting sector; geometric constraints
Theorem~\ref{thm:ch2.geometric-spectral-constraints} may force individual \(c_n=0\) without
breaking Eq.~\eqref{eq:ch2.coupling-bilinear-symmetry} \cite{harrington2001,stratton1941}.

\begin{figure}[t]
\centering
\ChOneFigBlock{%
\begin{tikzpicture}[ch1fig, node distance=7mm and 10mm]
  \node[ch1box] (j1) {$\Jimp^{(1)}$};
  \node[ch1box, right=11mm of j1] (r) {$\mathcal{R}_\omega$};
  \node[ch1box, right=11mm of r] (u) {$\sum_n \lambda_n c_n \mathbf{u}_n$};
  \node[ch1box, below=11mm of j1] (j2) {$\Jimp^{(2)}$};
  \node[ch1box, below=11mm of u] (c) {$\mathfrak{C}_\omega$};
  \node[ch1box, left=11mm of c] (orth) {Eq.~\eqref{eq:ch2.radiation-modal-orthogonality}};
  \draw[ch1flow] (j1) -- (r);
  \draw[ch1flow] (j2) -- (r);
  \draw[ch1flow] (r) -- (u);
  \draw[ch1flow] (u) -- (c);
  \draw[ch1flow] (orth) -- (c);
\end{tikzpicture}%
}
\caption{Subsection~\ref{sec:ch282} modal coupling pipeline: sources map through
\(\mathcal{R}_\omega\) to orthonormal modes Eq.~\eqref{eq:ch2.radiation-modal-orthogonality};
bilinear coupling Eq.~\eqref{eq:ch2.radiative-coupling-bilinear} factorizes through modal
amplitudes Eq.~\eqref{eq:ch2.modal-coupling-sum}.}
\label{fig:ch2.modal-coupling-pipeline}
\end{figure}

Figure~\ref{fig:ch2.modal-coupling-pipeline} separates orthogonality of mode shapes from
bilinear coupling of impressed sources: Eq.~\eqref{eq:ch2.radiation-modal-orthogonality} is
intrinsic to \(\mathcal{F}_{\mathrm{rad}}\); Eq.~\eqref{eq:ch2.radiative-coupling-bilinear}
depends on both sources \cite{stratton1941,harrington2001}.

\begin{theorem}[Orthogonality and coupling of radiative modes]
\label{thm:ch2.radiation-modal-orthogonality-coupling}
Let \(\{\mathbf{u}_n,\Jimp^{(n)},\lambda_n\}\) be radiation eigenpairs
Eq.~\eqref{eq:ch2.radiation-eigenpair-def} forming an orthonormal basis of
\(\mathcal{F}_{\mathrm{rad}}\) on the declared domain class, with reciprocal media
Eq.~\eqref{eq:ch2.reciprocal-medium-class}. Then:
\begin{enumerate}
\item[(i)] Orthogonality Eq.~\eqref{eq:ch2.radiation-modal-orthogonality} holds on
\(\mathcal{N}_{\mathrm{rad}}\).
\item[(ii)] Coupling coefficients satisfy Eq.~\eqref{eq:ch2.modal-coupling-diagonal} and
biorthogonality Eq.~\eqref{eq:ch2.radiation-biorthogonality}.
\item[(iii)] For all admissible \(\Jimp^{(1)},\Jimp^{(2)}\),
Eqs.~\eqref{eq:ch2.radiative-coupling-bilinear}--\eqref{eq:ch2.modal-coupling-sum-diagonal}
hold, with symmetry Eq.~\eqref{eq:ch2.coupling-bilinear-symmetry} from Theorem
~\ref{thm:ch2.lorentz-reciprocity-operator}.
\item[(iv)] Adding \(\Jimp^{\mathrm{(NR)}}\in\ker\mathcal{R}_\omega\) to either source leaves
Eq.~\eqref{eq:ch2.radiative-coupling-bilinear} unchanged.
\end{enumerate}
\end{theorem}
\begin{proof}
Part (i) is normalization and completeness from Theorem
~\ref{thm:ch2.radiation-spectral-decomposition} and
Eq.~\eqref{eq:ch2.radiation-spectral-expansion} \cite{stratton1941}. Part (ii) substitutes
\(\mathcal{R}_\omega[\Jimp^{(n)}]=\lambda_n\mathbf{u}_n\) into
Eq.~\eqref{eq:ch2.modal-coupling-coefficient} and applies
Eq.~\eqref{eq:ch2.radiation-biorthogonality} \cite{stratton1941,harrington2001}. Part (iii)
expands \(\mathcal{R}_\omega[\Jimp^{(k)}]\) via
Eq.~\eqref{eq:ch2.radiation-map-modal-sum}; symmetry is Theorem
~\ref{thm:ch2.lorentz-reciprocity-operator}(iv) \cite{stratton1941}. Part (iv) follows from
\(\mathcal{R}_\omega[\Jimp^{\mathrm{(NR)}}]=\mathbf{0}\) and
Eq.~\eqref{eq:ch2.nr-superposition-law} \cite{stratton1941,harrington2001}.
\end{proof}

\paragraph{Reciprocal reduction of adjoint modes.}
Under Eqs.~\eqref{eq:ch2.reciprocal-medium-class}--\eqref{eq:ch2.passive-impedance-class},
\begin{equation}
\label{eq:ch2.reciprocal-adjoint-mode-identification}
\widetilde{\mathbf{u}}_n=\mathbf{u}_n
\quad\text{on}\ \mathcal{F}_{\mathrm{rad}}
\quad\text{for exporting modes with}\ \lambda_n\neq 0,
\end{equation}
so Eq.~\eqref{eq:ch2.radiation-biorthogonality} reduces to standard orthogonality
Eq.~\eqref{eq:ch2.radiation-modal-orthogonality} with diagonal coupling
Eq.~\eqref{eq:ch2.modal-coupling-diagonal} \cite{stratton1941,harrington2001}. Equation
\eqref{eq:ch2.reciprocal-adjoint-mode-identification} is the operator form of reciprocity for
modal probes: receive patterns aligned with \(\mathbf{u}_n\) coincide with transmit modes
when media and boundaries are reciprocal \cite{harrington2001,jackson1999}. Non-reciprocal
media break Eq.~\eqref{eq:ch2.reciprocal-adjoint-mode-identification}; biorthogonality with
distinct \(\widetilde{\mathbf{u}}_n\) remains the general template imported from
Subsection~\ref{sec:ch261} \cite{stratton1941}.

\paragraph{Degenerate subspace coupling.}
When \(\lambda_n=\lambda_{n'}\) for \(n\neq n'\) within a rotationally degenerate family from
Subsection~\ref{sec:ch262}, Eq.~\eqref{eq:ch2.modal-coupling-diagonal} still holds for
\emph{eigen} sources \(\Jimp^{(n)}\), but superpositions within the degenerate block carry
off-diagonal coupling in an arbitrary basis. Let \(\mathcal{I}_\ell\subset\mathcal{N}_{\mathrm{rad}}\)
index a \((2\ell+1)\)-fold degenerate family. Any orthonormal basis
\(\{\mathbf{u}_n\}_{n\in\mathcal{I}_\ell}\) spanning the block satisfies
Eq.~\eqref{eq:ch2.radiation-modal-orthogonality}, yet a fixed physical source may excite
multiple members with coefficients tied to orientation through
Eq.~\eqref{eq:ch2.radiation-rotation-commutation} \cite{stratton1941,harrington2001}. Define
the block coupling matrix
\begin{equation}
\label{eq:ch2.degenerate-block-coupling-matrix}
\mathbf{K}^{(\ell)}_{mn}
=
\langle \mathbf{u}_m,\,\mathcal{R}_\omega[\Jimp^{(n)}]\rangle_{\mathcal{U}},
\qquad
m,n\in\mathcal{I}_\ell,
\end{equation}
which is diagonal in the eigenvector basis of \(\Jimp^{(n)}\) but not in laboratory-fixed
coordinates \cite{stratton1941,jackson1999}. Reciprocity
Eq.~\eqref{eq:ch2.coupling-bilinear-symmetry} forces \(\mathbf{K}^{(\ell)}\) to be symmetric
under Theorem~\ref{thm:ch2.lorentz-reciprocity-operator} \cite{stratton1941,harrington2001}.
Explicit Wigner coupling of block entries belongs to Chapter~\ref{ch:03}
\cite{stratton1941}.

\begin{proposition}[Cross-mode coupling from superposition]
\label{prop:ch2.superposition-cross-coupling}
If \(\Jimp^{(1)}=\sum_{n\in\mathcal{N}_{\mathrm{rad}}}a_n\Jimp^{(n)}\) and
\(\Jimp^{(2)}=\sum_{m\in\mathcal{N}_{\mathrm{rad}}}b_m\Jimp^{(m)}\) are admissible finite
superpositions, then
\begin{equation}
\label{eq:ch2.superposition-coupling-expansion}
\mathfrak{C}_\omega\!\big(\Jimp^{(1)},\Jimp^{(2)}\big)
=
\sum_{m,n\in\mathcal{N}_{\mathrm{rad}}}
a_m^{*}b_n\,\lambda_m\lambda_n\,
\langle \mathbf{u}_m,\mathbf{u}_n\rangle_{\mathcal{U}},
\end{equation}
reducing to Eq.~\eqref{eq:ch2.modal-coupling-sum-diagonal} when \(\langle\mathbf{u}_m,\mathbf{u}_n\rangle=\delta_{mn}\)
\cite{stratton1941,harrington2001}. Off-diagonal terms \(m\neq n\) appear only when at least
one source superposes non-collinear modes; single-mode excitation yields diagonal coupling
\cite{jackson1999}.
\end{proposition}
\begin{proof}
Linearity Eq.~\eqref{eq:ch2.radiation-linearity} and
Eq.~\eqref{eq:ch2.radiation-map-modal-sum} \cite{stratton1941}.
\end{proof}

\begin{proposition}[Observation-limited coupling]
\label{prop:ch2.observation-limited-coupling}
For observation map \(\mathcal{M}_\omega=\mathcal{O}_{\mathrm{meas}}\circ\Lambda_{\mathrm{rad}}\)
Eq.~\eqref{eq:ch2.reconstruction-observation-map}, define observed coupling
\(\mathfrak{C}_{\omega,\mathrm{obs}}=
\langle \mathcal{O}_{\mathrm{meas}}[\mathcal{R}_\omega[\Jimp^{(1)}]],\,
\mathcal{O}_{\mathrm{meas}}[\mathcal{R}_\omega[\Jimp^{(2)}]]\rangle\).
Then \(\mathfrak{C}_{\omega,\mathrm{obs}}\) is symmetric under Theorem
~\ref{thm:ch2.lorentz-reciprocity-operator} when \(\mathcal{O}_{\mathrm{meas}}\) is self-adjoint
on \(\mathcal{F}_{\mathrm{rad}}\), but may vanish for pairs with
\(\mathcal{R}_\omega[\Jimp^{(k)}]\in\ker\mathcal{O}_{\mathrm{meas}}\) even when
Eq.~\eqref{eq:ch2.radiative-coupling-bilinear} is nonzero \cite{stratton1941,harrington2001}.
Rank cap Eq.~\eqref{eq:ch2.reconstruction-rank-cap} limits the number of independently
observable coupling channels \cite{stratton1941}.
\end{proposition}
\begin{proof}
Apply Eq.~\eqref{eq:ch2.radiation-map-reciprocity} and self-adjointness of
\(\mathcal{O}_{\mathrm{meas}}\) \cite{stratton1941,harrington2001}. Kernel inclusion
\(\ker\mathcal{R}_\omega\subset\ker\mathcal{M}_\omega\) is
Proposition~\ref{prop:ch2.measurement-nullspace-extension} \cite{stratton1941}.
\end{proof}

\paragraph{Relation to angular spectra and null spaces.}
Accessible modal sector Eq.~\eqref{eq:ch2.accessible-modal-sector} and angular budget
Eq.~\eqref{eq:ch2.angular-spectrum-composition-law} restrict which terms in
Eq.~\eqref{eq:ch2.modal-coupling-sum-diagonal} are observable in far-zone reports
\cite{stratton1941,harrington2001}. Null modes \(\lambda_n=0\) carry zero coupling diagonal
entries but may alter near reactive content through \(\mathcal{S}_\omega\)
Eq.~\eqref{eq:ch2.kernel-field-image-law} without changing
Eq.~\eqref{eq:ch2.radiative-coupling-bilinear} \cite{stratton1941}. Composition with
Proposition~\ref{prop:ch2.reciprocity-null-invariance} and
Theorem~\ref{thm:ch2.radiation-modal-orthogonality-coupling}(iv) explains why invisible
currents do not break reciprocity-qualified coupling on \(\operatorname{ran}\mathcal{R}_\omega\)
\cite{harrington2001}.

\paragraph{Boundary of validity.}
Theorem~\ref{thm:ch2.radiation-modal-orthogonality-coupling} requires linear reciprocal media
at fixed \(\omega\), outgoing closure, domain-class completeness of
\(\{\mathbf{u}_n\}\), and convergent modal sums \cite{stratton1941,harrington2001}. Loss
moves \(\lambda_n\) into the complex plane; Eq.~\eqref{eq:ch2.modal-coupling-sum-diagonal}
generalizes with complex conjugation on one amplitude slot \cite{stratton1941}. Truncated
modal books on cavities do not substitute for exterior continuous spectra
Eq.~\eqref{eq:ch2.exterior-continuous-spectral-integral} when coupling far-zone channels
\cite{jackson1999,harrington2001}. Regime gate Eq.~\eqref{eq:ch2.regime-component-gate} is
mandatory before interpreting Eq.~\eqref{eq:ch2.radiative-coupling-bilinear} as transport
\cite{stratton1941}.

\paragraph{Worked illustration (single-mode excitation).}
If \(\Jimp^{(1)}=\Jimp^{(n)}\) and \(\Jimp^{(2)}=\Jimp^{(m)}\) with \(\lambda_n,\lambda_m\neq 0\),
then Eq.~\eqref{eq:ch2.radiative-coupling-bilinear} gives
\(\mathfrak{C}_\omega=\lambda_n\lambda_m\delta_{mn}\) through
Eqs.~\eqref{eq:ch2.modal-coupling-diagonal} and~\eqref{eq:ch2.radiation-modal-orthogonality}
\cite{stratton1941,harrington2001}. Orthogonal modes do not cross-couple on
\(\mathcal{F}_{\mathrm{rad}}\) \cite{jackson1999}.

\paragraph{Worked illustration (two-mode superposition).}
With \(\Jimp^{(1)}=a_p\Jimp^{(p)}+a_q\Jimp^{(q)}\) and \(\lambda_p\neq\lambda_q\),
Eq.~\eqref{eq:ch2.superposition-coupling-expansion} yields off-diagonal contributions only
when both modes appear in both sources; a single-mode \(\Jimp^{(2)}=\Jimp^{(q)}\) recovers
\(\mathfrak{C}_\omega=|a_q|^{2}|\lambda_q|^{2}\) \cite{stratton1941,harrington2001}. Pattern
matching that ignores modal orthogonality can misattribute such cross terms to hardware
non-reciprocity \cite{harrington2001}.

\paragraph{Worked illustration (degenerate multiplet).}
For a \((2\ell+1)\)-fold family from Subsection~\ref{sec:ch262}, rotating the source rotates
coefficients \(c_n\) through Eq.~\eqref{eq:ch2.radiation-rotation-commutation} while
\(\mathfrak{C}_\omega\) on a co-rotated probe pair remains invariant
\cite{stratton1941,jackson1999}. Block matrix Eq.~\eqref{eq:ch2.degenerate-block-coupling-matrix}
is diagonal in eigenvector coordinates but dense in fixed-laboratory labeling
\cite{harrington2001}.

\paragraph{Worked illustration (finite array).}
When \(\mathcal{O}_{\mathrm{meas}}\) has rank \(N_{\mathrm{obs}}<\dim\mathcal{F}_{\mathrm{rad}}\),
Proposition~\ref{prop:ch2.observation-limited-coupling} shows symmetric coupling among
observed channels even though full Eq.~\eqref{eq:ch2.radiative-coupling-bilinear} sees more
modes \cite{stratton1941,harrington2001}. Eq.~\eqref{eq:ch2.reconstruction-rank-cap} caps
recoverable coupling entries independently of modal orthogonality
\cite{stratton1941}.

\paragraph{Closing summary.}
Subsection~\ref{sec:ch282} establishes the Chapter~\ref{ch:02} HOME for radiation-modal
orthogonality Eq.~\eqref{eq:ch2.radiation-modal-orthogonality}, coupling coefficients
Eqs.~\eqref{eq:ch2.modal-coupling-coefficient}--\eqref{eq:ch2.modal-coupling-diagonal},
radiative bilinear coupling Eq.~\eqref{eq:ch2.radiative-coupling-bilinear} with symmetry
Eq.~\eqref{eq:ch2.coupling-bilinear-symmetry}, modal sum formulas
Eqs.~\eqref{eq:ch2.modal-coupling-sum}--\eqref{eq:ch2.modal-coupling-sum-diagonal},
reciprocal adjoint identification Eq.~\eqref{eq:ch2.reciprocal-adjoint-mode-identification},
degenerate block matrix Eq.~\eqref{eq:ch2.degenerate-block-coupling-matrix}, and Theorem
~\ref{thm:ch2.radiation-modal-orthogonality-coupling} upgrading previews
Eqs.~\eqref{eq:ch2.modal-orthogonality-preview} and~\eqref{eq:ch2.radiative-coupling-bilinear-preview}
\cite{stratton1941,harrington2001,jackson1999}. Subsection~\ref{sec:ch283} applies reciprocity
to radiation and reception systems; Chapter~\ref{ch:03} supplies explicit harmonic
orthogonality used by reference through Subsection~\ref{sec:ch343}
\cite{stratton1941}.
\subsection{Reciprocity in radiation and reception}
\label{sec:ch283}
% TOC tag: [HOME]

Subsections~\ref{sec:ch281}--\ref{sec:ch282} established Lorentz reciprocity Theorem
~\ref{thm:ch2.lorentz-reciprocity-operator}, radiative bilinear coupling
Eq.~\eqref{eq:ch2.radiative-coupling-bilinear}, and modal orthogonality Theorem
~\ref{thm:ch2.radiation-modal-orthogonality-coupling} on \(\mathcal{F}_{\mathrm{rad}}\)
\cite{stratton1941,harrington2001}. Subsection~\ref{sec:ch283} gives the Chapter~\ref{ch:02}
HOME for reciprocity in radiation and reception: transmit--receive operator interchange,
two-port coupling symmetry, and measurement-qualified reciprocity on
\(\operatorname{ran}\mathcal{R}_\omega\), upgrading preview
Eq.~\eqref{eq:ch2.transmit-receive-swap-preview} \cite{stratton1941,jackson1999}. Observation
composition Eq.~\eqref{eq:ch2.reconstruction-observation-map}, finite-rank readouts
Eq.~\eqref{eq:ch2.finite-rank-observation}, rank cap
Eq.~\eqref{eq:ch2.reconstruction-rank-cap}, and null-space laws
Eqs.~\eqref{eq:ch2.reconstruction-null-space}--\eqref{eq:ch2.nr-superposition-law} are
presupposed; inverse-source regularization is not re-developed \cite{stratton1941,harrington2001}.
Plane-wave angular-spectrum filters belong to Section~\ref{sec:ch29}; explicit port circuit
models are not tabulated \cite{stratton1941}.

\paragraph{Assumptions.}
Fix \(\omega>0\), reciprocal media Eq.~\eqref{eq:ch2.reciprocal-medium-class}, passive
boundary Eq.~\eqref{eq:ch2.passive-impedance-class}, outgoing closure, and admissible sources
\(\Jimp^{\mathrm{(tx)}},\Jimp^{\mathrm{(rx)}}\in\mathcal{J}_\Omega\) on ports \(A\) and \(B\)
\cite{stratton1941,harrington2001}. Each port carries a finite-rank observation functional
\(\mathcal{O}_{\mathrm{meas}}^{(P)}:\mathcal{F}_{\mathrm{rad}}\to\mathbb{C}\) or
\(\mathbb{C}^{N_P}\) of the form Eq.~\eqref{eq:ch2.finite-rank-observation}
\cite{stratton1941}. Regime gate Eq.~\eqref{eq:ch2.regime-operator-screen} and transport
qualification Eq.~\eqref{eq:ch2.transport-qualified-radiation} govern when port outputs may be
reported as radiative transport \cite{harrington2001}. The same \(\mathcal{B}_\Gamma\), outgoing
branch, and material data are used for both excitations \cite{stratton1941}.

\paragraph{Transmit and receive maps.}
For port \(P\in\{A,B\}\), define the \emph{port radiation--reception operator}
\begin{equation}
\label{eq:ch2.port-radiation-reception-operator}
\mathcal{M}_P
=
\mathcal{O}_{\mathrm{meas}}^{(P)}\circ\Lambda_{\mathrm{rad}},
\qquad
\Lambda_{\mathrm{rad}}=\mathcal{O}_{\mathrm{far}}\circ\mathcal{R}_\omega,
\end{equation}
consistent with Eq.~\eqref{eq:ch2.reconstruction-observation-map}
\cite{stratton1941,harrington2001}. When \(\mathcal{O}_{\mathrm{meas}}^{(P)}\) is scalar,
\(\mathcal{M}_P[\Jimp]\in\mathbb{C}\) is the complex port response induced by impressed
\(\Jimp\); when vector-valued, \(\mathcal{M}_P[\Jimp]\in\mathbb{C}^{N_P}\) lists channel
readouts \cite{stratton1941}. The \emph{transmit map} at port \(B\) due to excitation at
port \(A\) is the composition \(\Jimp^{\mathrm{(tx)}}\mapsto
\mathcal{M}_B[\Jimp^{\mathrm{(tx)}}]\); the \emph{receive map} at port \(A\) due to excitation
at port \(B\) is \(\Jimp^{\mathrm{(rx)}}\mapsto\mathcal{M}_A[\Jimp^{\mathrm{(rx)}}]\)
\cite{harrington2001,jackson1999}. Algebraically, both are finite-rank linear functionals of
\(\mathcal{R}_\omega[\Jimp]\); physically, reciprocity constrains their interchange when
hardware and media are reciprocal \cite{stratton1941}.

\paragraph{Transmit--receive reciprocity (HOME).}
Under Theorem~\ref{thm:ch2.lorentz-reciprocity-operator} and self-adjoint port functionals
\(\mathcal{O}_{\mathrm{meas}}^{(P)}\) on \(\mathcal{F}_{\mathrm{rad}}\),
\begin{equation}
\label{eq:ch2.transmit-receive-reciprocity}
\mathcal{M}_B\big[\Jimp^{\mathrm{(tx)}}\big]
=
\mathcal{M}_A\big[\Jimp^{\mathrm{(rx)}}\big]
\quad\text{when}\ 
\Jimp^{\mathrm{(tx)}}\ \text{and}\ \Jimp^{\mathrm{(rx)}}\ \text{are reciprocal port excitations},
\end{equation}
upgrading preview Eq.~\eqref{eq:ch2.transmit-receive-swap-preview}
\cite{stratton1941,harrington2001,jackson1999}. Equation~\eqref{eq:ch2.transmit-receive-reciprocity}
states that the observable induced at \(B\) by the transmit source at \(A\) equals the observable
at \(A\) when the roles of impressed source and probe are exchanged on reciprocal media
\cite{stratton1941}. It is not a claim that \(\mathcal{M}_B=\mathcal{M}_A\) as operators on
\(\mathcal{J}_\Omega\): ports may have distinct apertures and channel counts
\cite{harrington2001}. It \emph{is} a claim that the \emph{bilinear} port coupling is symmetric
when qualified by Eq.~\eqref{eq:ch2.radiative-coupling-bilinear} on
\(\operatorname{ran}\mathcal{R}_\omega\) \cite{stratton1941}.

\begin{figure}[t]
\centering
\ChOneFigBlock{%
\begin{tikzpicture}[ch1fig, node distance=7mm and 10mm]
  \node[ch1box] (tx) {Port \(A\)\\$\Jimp^{\mathrm{(tx)}}$};
  \node[ch1box, right=14mm of tx] (r) {$\mathcal{R}_\omega$};
  \node[ch1box, right=14mm of r] (ob) {$\mathcal{O}_{\mathrm{meas}}^{(B)}$};
  \node[ch1box, right=14mm of ob] (rxout) {Rx at \(B\)};
  \node[ch1box, below=13mm of tx] (rx) {Port \(B\)\\$\Jimp^{\mathrm{(rx)}}$};
  \node[ch1box, below=13mm of ob] (ob2) {$\mathcal{O}_{\mathrm{meas}}^{(A)}$};
  \node[ch1box, below=13mm of rxout] (txout) {Rx at \(A\)};
  \draw[ch1flow] (tx) -- (r) -- (ob) -- (rxout);
  \draw[ch1flow] (rx) -- (r) -- (ob2) -- (txout);
  \draw[ch1flow, dashed] (rxout) to[bend left=20] node[right,font=\scriptsize] {swap} (txout);
\end{tikzpicture}%
}
\caption{Subsection~\ref{sec:ch283} transmit--receive pipeline: excitation at one port and
readout at the other interchange under Eq.~\eqref{eq:ch2.transmit-receive-reciprocity} on
reciprocal media.}
\label{fig:ch2.transmit-receive-pipeline}
\end{figure}

Figure~\ref{fig:ch2.transmit-receive-pipeline} displays the two compositions constrained by
Eq.~\eqref{eq:ch2.transmit-receive-reciprocity}. Both paths pass through the same
\(\mathcal{R}_\omega\); port asymmetry enters only through \(\mathcal{O}_{\mathrm{meas}}^{(A)}\),
\(\mathcal{O}_{\mathrm{meas}}^{(B)}\) \cite{stratton1941,harrington2001}.

\begin{theorem}[Reciprocity in radiation and reception]
\label{thm:ch2.radiation-reception-reciprocity}
Let media satisfy Eq.~\eqref{eq:ch2.reciprocal-medium-class} with passive
Eq.~\eqref{eq:ch2.passive-impedance-class} closure, and let port operators
\(\mathcal{M}_P=\mathcal{O}_{\mathrm{meas}}^{(P)}\circ\Lambda_{\mathrm{rad}}\) be defined by
Eq.~\eqref{eq:ch2.port-radiation-reception-operator} with self-adjoint
\(\mathcal{O}_{\mathrm{meas}}^{(P)}\) on \(\mathcal{F}_{\mathrm{rad}}\). For admissible port
sources \(\Jimp^{\mathrm{(tx)}}\), \(\Jimp^{\mathrm{(rx)}}\):
\begin{enumerate}
\item[(i)] Scalar reciprocity Eq.~\eqref{eq:ch2.transmit-receive-reciprocity} holds when port
excitations are paired through Lorentz reciprocity Theorem
~\ref{thm:ch2.lorentz-reciprocity-operator}.
\item[(ii)] Bilinear port coupling
\begin{equation}
\label{eq:ch2.port-coupling-bilinear}
\mathfrak{C}_{AB}
=
\langle \mathcal{O}_{\mathrm{meas}}^{(B)}[\mathcal{R}_\omega[\Jimp^{\mathrm{(tx)}}]],\,
\mathcal{O}_{\mathrm{meas}}^{(A)}[\mathcal{R}_\omega[\Jimp^{\mathrm{(rx)}}]]\rangle
\end{equation}
is symmetric under interchange of port labels when inner products are those induced by
Eq.~\eqref{eq:ch2.field-energy-norm} on \(\mathcal{F}_{\mathrm{rad}}\).
\item[(iii)] Two-port impedance matrix \(\mathbf{Z}\) with entries
\(Z_{PQ}=\mathcal{M}_P[\Jimp^{(Q)}]\) on matched reciprocal ports satisfies
\begin{equation}
\label{eq:ch2.reciprocal-impedance-matrix}
Z_{PQ}=Z_{QP}.
\end{equation}
\item[(iv)] Adding \(\Jimp^{\mathrm{(NR)}}\in\ker\mathcal{R}_\omega\) to either port excitation
leaves Eqs.~\eqref{eq:ch2.transmit-receive-reciprocity}--\eqref{eq:ch2.reciprocal-impedance-matrix}
unchanged on observable coordinates.
\end{enumerate}
\end{theorem}
\begin{proof}
Part (i) applies Theorem~\ref{thm:ch2.lorentz-reciprocity-operator}(iv) to
\(\mathcal{R}_\omega[\Jimp^{\mathrm{(tx)}}]\), \(\mathcal{R}_\omega[\Jimp^{\mathrm{(rx)}}]\),
then self-adjointness of \(\mathcal{O}_{\mathrm{meas}}^{(P)}\)
\cite{stratton1941,harrington2001}. Part (ii) is
Eq.~\eqref{eq:ch2.coupling-bilinear-symmetry} with observation weights inserted; symmetry of
\(\mathfrak{C}_\omega\) from Theorem~\ref{thm:ch2.radiation-modal-orthogonality-coupling}(iii)
\cite{stratton1941}. Part (iii) specializes (ii) to matched port basis vectors
\cite{harrington2001,jackson1999}. Part (iv) uses
\(\mathcal{R}_\omega[\Jimp^{\mathrm{(NR)}}]=\mathbf{0}\) and
Proposition~\ref{prop:ch2.reciprocity-null-invariance} \cite{stratton1941,harrington2001}.
\end{proof}

\paragraph{Mathematical role of Eq.~\eqref{eq:ch2.port-coupling-bilinear}.}
Equation~\eqref{eq:ch2.port-coupling-bilinear} packages mutual coupling as a bilinear form on
\(\operatorname{ran}\mathcal{R}_\omega\) after port projection \cite{stratton1941}. It reduces
to Eq.~\eqref{eq:ch2.radiative-coupling-bilinear} when
\(\mathcal{O}_{\mathrm{meas}}^{(A)}=\mathcal{O}_{\mathrm{meas}}^{(B)}=\mathrm{id}\) on the
observed subspace \cite{stratton1941,harrington2001}. Physically, \(\mathfrak{C}_{AB}\) is the
operator form of mutual impedance or scattering amplitude between declared ports when regime gate
Eq.~\eqref{eq:ch2.regime-component-gate} is satisfied \cite{harrington2001,jackson1999}. Finite
rank Eq.~\eqref{eq:ch2.reconstruction-rank-cap} limits the number of independent entries
\(\mathfrak{C}_{AB}\) can resolve \cite{stratton1941}.

\paragraph{Modal port coupling.}
On orthonormal modes Eq.~\eqref{eq:ch2.radiation-modal-orthogonality}, port responses expand as
\begin{equation}
\label{eq:ch2.port-modal-response-expansion}
\mathcal{M}_P[\Jimp]
=
\sum_{n\in\mathcal{N}_{\mathrm{rad}}}
w_{P,n}\,\lambda_n\,c_n[\Jimp],
\end{equation}
with port weights \(w_{P,n}=\langle \mathbf{f}_P,\mathbf{u}_n\rangle_{\mathcal{U}}\) inherited
from Eq.~\eqref{eq:ch2.finite-rank-observation} and amplitudes \(c_n\) from
Eq.~\eqref{eq:ch2.source-modal-amplitude} \cite{stratton1941,harrington2001}. Mutual coupling
between ports \(A\) and \(B\) becomes
\begin{equation}
\label{eq:ch2.port-modal-coupling-sum}
\mathfrak{C}_{AB}
=
\sum_{n\in\mathcal{N}_{\mathrm{rad}}}
|\lambda_n|^{2}\,w_{A,n}^{*}\,w_{B,n}\,
c_n[\Jimp^{\mathrm{(tx)}}]^{*}\,c_n[\Jimp^{\mathrm{(rx)}}],
\end{equation}
using Eq.~\eqref{eq:ch2.modal-coupling-sum-diagonal} \cite{stratton1941}. Equation
\eqref{eq:ch2.port-modal-coupling-sum} shows that reciprocity symmetry is diagonal in a mode
basis but may be dense in fixed laboratory coordinates when degenerate multiplets from
Subsection~\ref{sec:ch262} are expressed in non-eigenvector bases
\cite{harrington2001,jackson1999}.

\begin{proposition}[Measurement-limited reciprocity]
\label{prop:ch2.measurement-reciprocity}
For full measurement map \(\mathcal{M}_\omega=\mathcal{O}_{\mathrm{meas}}\circ\Lambda_{\mathrm{rad}}\)
Eq.~\eqref{eq:ch2.reconstruction-observation-map}, reciprocity implies
\(\mathcal{M}_\omega[\Jimp^{\mathrm{(tx)}}]\) and \(\mathcal{M}_\omega[\Jimp^{\mathrm{(rx)}}]\)
obey the same swap law Eq.~\eqref{eq:ch2.transmit-receive-reciprocity} on
\(\operatorname{ran}\mathcal{O}_{\mathrm{meas}}\) \cite{stratton1941,harrington2001}. Sources
in \(\ker\mathcal{M}_\omega\) leave all port readouts unchanged
Eq.~\eqref{eq:ch2.reconstruction-null-space}; reciprocity does not resolve those equivalence
classes \cite{stratton1941}. Rank cap Eq.~\eqref{eq:ch2.reconstruction-rank-cap} bounds
independent reciprocal channels \cite{harrington2001}.
\end{proposition}
\begin{proof}
Apply Theorem~\ref{thm:ch2.radiation-reception-reciprocity}(i) with
\(\mathcal{O}_{\mathrm{meas}}^{(A)}=\mathcal{O}_{\mathrm{meas}}^{(B)}=\mathcal{O}_{\mathrm{meas}}\)
\cite{stratton1941}. Null-space invariance is
Proposition~\ref{prop:ch2.measurement-nullspace-extension} \cite{stratton1941,harrington2001}.
\end{proof}

\begin{proposition}[Accessible-sector reciprocity]
\label{prop:ch2.accessible-sector-reciprocity}
Let accessible modal sector Eq.~\eqref{eq:ch2.accessible-modal-sector} restrict far-zone content
to \(\mathcal{K}_{\mathrm{acc}}\). Then Eqs.~\eqref{eq:ch2.transmit-receive-reciprocity}--\eqref{eq:ch2.reciprocal-impedance-matrix}
hold on \(\mathcal{K}_{\mathrm{acc}}\) after projection, while suppressed directions
Eq.~\eqref{eq:ch2.suppressed-angular-sheet-def} remain unobservable independently of reciprocity
\cite{stratton1941,harrington2001}. Angular budget
Eq.~\eqref{eq:ch2.angular-spectrum-composition-law} composes with port reciprocity without
eliminating null-space or geometric deficits \cite{stratton1941}.
\end{proposition}
\begin{proof}
Reciprocity acts on \(\operatorname{ran}\mathcal{R}_\omega\); geometric and observation
constraints from Theorem~\ref{thm:ch2.geometric-spectral-constraints} and
Eq.~\eqref{eq:ch2.angular-spectrum-composition-law} restrict reported coordinates only
\cite{stratton1941,harrington2001}.
\end{proof}

\paragraph{Failure modes and broken reciprocity.}
Magneto-electric media, active surfaces, and non-reciprocal port calibrations break
Eq.~\eqref{eq:ch2.transmit-receive-reciprocity} even when forward maps
\(\mathcal{S}_\omega\) are accurate \cite{stratton1941,jackson1999}. Near-zone reactive
coupling measured without \(\Pi_{\mathrm{rad}}\) can appear non-reciprocal while
Eq.~\eqref{eq:ch2.radiation-map-reciprocity} holds on \(\mathcal{F}_{\mathrm{rad}}\)
\cite{harrington2001,stratton1941}. Using different outgoing branches or terminations for
transmit and receive models violates the shared-operator assumption behind Theorem
~\ref{thm:ch2.radiation-reception-reciprocity} \cite{stratton1941}. Tikhonov inversion
Eq.~\eqref{eq:ch2.tikhonov-source-reconstruction} does not restore reciprocity where
\(\ker\mathcal{M}_\omega\) is nontrivial \cite{harrington2001}.

\paragraph{Boundary of validity.}
Theorem~\ref{thm:ch2.radiation-reception-reciprocity} requires linear reciprocal media at fixed
\(\omega\), passive \(\mathcal{B}_\Gamma\), self-adjoint port functionals, and outgoing closure
\cite{stratton1941,harrington2001}. It does not identify unique \(\Jimp\) from port data when
Eq.~\eqref{eq:ch2.reconstruction-null-space} is nontrivial \cite{stratton1941}. It does not
replace angular-spectrum operators in Section~\ref{sec:ch29} or harmonic bases in
Chapter~\ref{ch:03} \cite{stratton1941,harrington2001}. Regime gate
Eq.~\eqref{eq:ch2.regime-operator-screen} remains mandatory before transport interpretation of
port coupling \cite{stratton1941}.

\paragraph{Worked illustration (half-wave dipole ports).}
Two center-fed dipoles in reciprocal vacuum satisfy
Eq.~\eqref{eq:ch2.reciprocal-impedance-matrix} for mutual impedance entries measured on
\(\mathcal{R}_\omega[\Jimp^{\mathrm{(tx)}}]\) in the far zone \cite{jackson1999,stratton1941}.
Near-field probe scans without regime qualification can violate apparent reciprocity while
Eq.~\eqref{eq:ch2.transmit-receive-reciprocity} holds for far-zone port functionals
\cite{harrington2001}.

\paragraph{Worked illustration (NR superposition at a port).}
Adding \(\Jimp^{\mathrm{(NR)}}\in\ker\mathcal{R}_\omega\) to \(\Jimp^{\mathrm{(tx)}}\) leaves
\(\mathcal{M}_B[\Jimp^{\mathrm{(tx)}}]\) unchanged by Theorem
~\ref{thm:ch2.radiation-reception-reciprocity}(iv) and Eq.~\eqref{eq:ch2.nr-superposition-law}
\cite{stratton1941,harrington2001}. Calibration that fits only far-zone port data cannot detect
the invisible addition \cite{stratton1941}.

\paragraph{Worked illustration (finite array).}
An \(N\)-element array with \(\mathcal{O}_{\mathrm{meas}}\) of rank \(N\) yields an
\(N\times N\) impedance matrix obeying Eq.~\eqref{eq:ch2.reciprocal-impedance-matrix} when
media and ports are reciprocal \cite{harrington2001,stratton1941}. Sidelobe content in
\(\ker\mathcal{O}_{\mathrm{meas}}\) remains unobserved and unconstrained by reciprocity alone
Eq.~\eqref{eq:ch2.observation-null-space} \cite{stratton1941}.

\paragraph{Worked illustration (lossy reciprocal slab).}
Lossy but symmetric \(\overline{\overline{\varepsilon}}\) preserves Theorem
~\ref{thm:ch2.radiation-reception-reciprocity} with complex conjugation in port pairings; mutual
impedance obeys Eq.~\eqref{eq:ch2.reciprocal-impedance-matrix} in the symmetric part
\cite{stratton1941,harrington2001}. Evanescent storage does not create reciprocal far-zone channels
through Eq.~\eqref{eq:ch2.rad-null-on-evanescent} \cite{stratton1941}.

\paragraph{Closing summary.}
Subsection~\ref{sec:ch283} establishes the Chapter~\ref{ch:02} HOME for port operators
Eq.~\eqref{eq:ch2.port-radiation-reception-operator}, transmit--receive reciprocity
Eq.~\eqref{eq:ch2.transmit-receive-reciprocity}, port coupling bilinear
Eq.~\eqref{eq:ch2.port-coupling-bilinear}, reciprocal impedance symmetry
Eq.~\eqref{eq:ch2.reciprocal-impedance-matrix}, modal port expansions
Eqs.~\eqref{eq:ch2.port-modal-response-expansion}--\eqref{eq:ch2.port-modal-coupling-sum},
Theorem~\ref{thm:ch2.radiation-reception-reciprocity}, and Propositions
~\ref{prop:ch2.measurement-reciprocity}--\ref{prop:ch2.accessible-sector-reciprocity}, upgrading
Eq.~\eqref{eq:ch2.transmit-receive-swap-preview} \cite{stratton1941,harrington2001,jackson1999}.

\medskip
\noindent
Section~\ref{sec:ch28} closes the reciprocity chain begun in its introduction:
Lorentz reciprocity in operator form Theorem~\ref{thm:ch2.lorentz-reciprocity-operator}
(Subsection~\ref{sec:ch281}), orthogonality and coupling of radiative modes Theorem
~\ref{thm:ch2.radiation-modal-orthogonality-coupling} (Subsection~\ref{sec:ch282}), and
reciprocity in radiation and reception Theorem~\ref{thm:ch2.radiation-reception-reciprocity}
(Subsection~\ref{sec:ch283}) \cite{stratton1941,harrington2001,jackson1999}. Spectral coordinates
from Section~\ref{sec:ch26} and null-space classification from Section~\ref{sec:ch27} now meet
exchange symmetry: bilinear coupling on \(\operatorname{ran}\mathcal{R}_\omega\) is constrained
by Eqs.~\eqref{eq:ch2.lorentz-reciprocity-operator}--\eqref{eq:ch2.reciprocal-impedance-matrix},
while \(\ker\mathcal{R}_\omega\), finite observation rank
Eq.~\eqref{eq:ch2.reconstruction-rank-cap}, and angular budget
Eq.~\eqref{eq:ch2.angular-spectrum-composition-law} explain deficits reciprocity alone does not
remove \cite{stratton1941}. Section~\ref{sec:ch29} develops plane-wave action and
retained/suppressed angular components on the same outgoing branch; Subsection~\ref{sec:ch294}
refines far-field angular accessibility; Chapter~\ref{ch:03} supplies explicit harmonic
orthogonality used by reference through Subsection~\ref{sec:ch343}
\cite{harrington2001}. Transport and port-coupling claims throughout Volume~I must compose with
Eq.~\eqref{eq:ch2.section28-reciprocity-chain}, regime gate
Eq.~\eqref{eq:ch2.regime-operator-screen}, and null-space chain
Eq.~\eqref{eq:ch2.section27-nullspace-chain}, not with raw \(\mathcal{S}_\omega[\Jimp]\) or
unqualified near-zone reactive content \cite{stratton1941,harrington2001}.
\section{Action of Radiation Operators on Angular Spectra}
\label{sec:ch29}

Section~\ref{sec:ch28} closed the reciprocity chain
Eq.~\eqref{eq:ch2.section28-reciprocity-chain}: Lorentz reciprocity Theorem
~\ref{thm:ch2.lorentz-reciprocity-operator}, modal orthogonality and coupling Theorem
~\ref{thm:ch2.radiation-modal-orthogonality-coupling}, and radiation--reception reciprocity
Theorem~\ref{thm:ch2.radiation-reception-reciprocity}
\cite{stratton1941,harrington2001,jackson1999}. That program fixed exchange symmetry on
\(\operatorname{ran}\mathcal{R}_\omega\) and on port coordinates after projection. Section
~\ref{sec:ch29} addresses the complementary \emph{angular-spectrum action} question: how
\(\mathcal{R}_\omega\), \(\Pi_{\mathrm{rad}}\), and observation maps filter plane-wave and
continuous angular content on \(S^{2}\), which components are retained or suppressed, and how
near-zone fields map to far-zone angular amplitudes \(a(\widehat{\mathbf{k}})\)
\cite{stratton1941,harrington2001}. Preview expansion
Eq.~\eqref{eq:ch2.angular-spectrum-preview} from Subsection~\ref{sec:ch224} named the lateral
spectral template; Section~\ref{sec:ch29} upgrades radiation-operator action on angular spectra to
HOME statements compatible with outgoing gate Eq.~\eqref{eq:ch2.outgoing-plane-wave-gate} and
far amplitude Eq.~\eqref{eq:ch2.far-field-amplitude-integral} \cite{stratton1941}. Explicit
basis change to vector harmonics is HOME in Chapter~\ref{ch:03} Subsection~\ref{sec:ch381}; harmonic
normalization belongs to Subsection~\ref{sec:ch352} with USE of Subsection~\ref{sec:ch282}
\cite{harrington2001}.

The problem addressed here is \emph{filtering} of angular content by the radiation pipeline.
Spectral decomposition Theorem~\ref{thm:ch2.radiation-spectral-decomposition} and continuous
representation Eq.~\eqref{eq:ch2.exterior-continuous-spectral-integral} from Subsection
~\ref{sec:ch261} fixed modal and plane-wave coordinates on \(\mathcal{F}_{\mathrm{rad}}\);
null-space consequences Theorem~\ref{thm:ch2.angular-spectrum-null-consequences} and angular
budget Eq.~\eqref{eq:ch2.angular-spectrum-composition-law} from Subsection~\ref{sec:ch274}
classified which directions \(a(\widehat{\mathbf{k}})\) may carry export
\cite{stratton1941,harrington2001}. Reciprocity from Section~\ref{sec:ch28} constrains bilinear
coupling on observable coordinates but does not replace angular filters: a reciprocal antenna
may still have suppressed sheets Eq.~\eqref{eq:ch2.suppressed-angular-sheet-def} and finite
observation rank Eq.~\eqref{eq:ch2.angular-spectrum-observation-rank}
\cite{stratton1941,jackson1999}. Treating \(\mathcal{S}_\omega[\Jimp]\) near sources as if it
were already \(a(\widehat{\mathbf{k}})\) repeats the regime error class of
Section~\ref{sec:ch01.2.4} \cite{stratton1941,harrington2001}. Section~\ref{sec:ch29} separates
plane-wave angular spectra and radiation operators, retained/suppressed/filtered components,
near-to-far transformation, and far-field angular accessibility without re-opening Green
construction from Section~\ref{sec:ch23} \cite{stratton1941}.

\paragraph{Angular-spectrum development chain.}
Let \(\omega>0\) be fixed, \(\Omega_s=\operatorname{supp}\Jimp\) be compact, and outgoing pairs
lie in \(\mathcal{F}_{\mathrm{SF}}\) Eq.~\eqref{eq:ch2.admissible-field-space}. Section
~\ref{sec:ch29} organizes analysis in four stages:
\begin{equation}
\label{eq:ch2.section29-angular-spectrum-chain}
\text{plane-wave spectrum}
\ \xrightarrow{\ \mathcal{R}_\omega,\Pi_{\mathrm{rad}}\ }
\text{retained / suppressed}
\ \xrightarrow{\ \text{near-to-far}\ }
a(\widehat{\mathbf{k}})
\ \xrightarrow{\ \text{accessibility}\ }
\mathcal{K}_{\mathrm{acc}},
\end{equation}
where the first arrow is Subsection~\ref{sec:ch291}, the second is Subsection~\ref{sec:ch292},
the third is Subsection~\ref{sec:ch293}, and the fourth is Subsection~\ref{sec:ch294}
\cite{stratton1941,harrington2001}. Equation~\eqref{eq:ch2.section29-angular-spectrum-chain}
restates Eq.~\eqref{eq:ch2.section27-nullspace-chain} and
Eq.~\eqref{eq:ch2.section28-reciprocity-chain} at the angular level: null-space and reciprocity
constraints from Sections~\ref{sec:ch27}--\ref{sec:ch28} bound which \(a(\widehat{\mathbf{k}})\)
are physically meaningful; Section~\ref{sec:ch29} develops the \emph{operators} that enforce
those bounds on plane-wave coordinates \cite{stratton1941}. Transport criteria from
Section~\ref{sec:ch01.2.3} and angular observables from Section~\ref{sec:ch01.2.1} enter
Subsection~\ref{sec:ch294} by reference, not by re-derivation \cite{stratton1941,harrington2001}.

\paragraph{Operator targets.}
The primary objects are radiation map
Eq.~\eqref{eq:ch2.nullspace-operator-targets}, outgoing plane-wave gate
Eq.~\eqref{eq:ch2.outgoing-plane-wave-gate}, and continuous angular amplitude
\(a(\widehat{\mathbf{k}})\) from Eq.~\eqref{eq:ch2.exterior-continuous-spectral-integral}
\cite{stratton1941,harrington2001}. Secondary targets include propagating projectors
Eqs.~\eqref{eq:ch2.spectral-projectors}--\eqref{eq:ch2.rad-null-on-evanescent},
observation composition Eq.~\eqref{eq:ch2.reconstruction-observation-map}, accessible sector
Eq.~\eqref{eq:ch2.angular-accessibility-sector}, and composition law
Eq.~\eqref{eq:ch2.angular-spectrum-composition-law} \cite{stratton1941}. Section~\ref{sec:ch29}
develops filters on the same outgoing branch fixed by
Eqs.~\eqref{eq:ch2.outgoing-radiation-refinement}--\eqref{eq:ch2.outgoing-resolvent-unique};
incoming sheets remain excluded from \(\mathcal{F}_{\mathrm{rad}}\)
\cite{harrington2001,stratton1941}.

\paragraph{Plane-wave angular spectrum and radiation operators (preview).}
Subsection~\ref{sec:ch291} gives the Chapter~\ref{ch:02} HOME for plane-wave angular spectra
and the action of \(\mathcal{R}_\omega\circ\Pi_{\mathrm{prop}}\) on spectral coefficients:
\begin{equation}
\label{eq:ch2.plane-wave-radiation-operator-preview}
\widetilde{\E}_{\mathrm{rad}}(k_x,k_y)
=
\mathcal{F}_{\mathrm{rad}}\big[\Pi_{\mathrm{rad}}\,\mathcal{R}_\omega[\Jimp]\big](k_x,k_y),
\qquad
k_z=k_z(k_t)\ \text{from Eq.~\eqref{eq:ch2.outgoing-plane-wave-gate}},
\end{equation}
with \(\mathcal{F}_{\mathrm{rad}}\) the lateral Fourier transform on outgoing admissible fields
\cite{stratton1941,harrington2001,jackson1999}. Equation~\eqref{eq:ch2.plane-wave-radiation-operator-preview}
upgrades preview Eq.~\eqref{eq:ch2.angular-spectrum-preview} by inserting radiation selection;
evanescent sheets satisfy \(\widetilde{\E}_{\mathrm{rad}}=0\) on export through
Eq.~\eqref{eq:ch2.rad-null-on-evanescent} \cite{stratton1941}. Rigged-space normalization of
plane-wave modes belongs to Chapter~\ref{ch:03} Subsection~\ref{sec:ch313} \cite{harrington2001}.

\paragraph{Retained, suppressed, and filtered components (preview).}
Subsection~\ref{sec:ch292} partitions angular content into retained, suppressed, and
observation-filtered sheets:
\begin{equation}
\label{eq:ch2.retained-suppressed-filter-preview}
a(\widehat{\mathbf{k}})
=
a_{\mathrm{ret}}(\widehat{\mathbf{k}})\,\chi_{\mathcal{K}_{\mathrm{acc}}}(\widehat{\mathbf{k}})
\;+\;
a_{\mathrm{sup}}(\widehat{\mathbf{k}})\,\chi_{\mathcal{K}_{\mathrm{sup}}}(\widehat{\mathbf{k}}),
\qquad
a_{\mathrm{filt}}=\mathcal{O}_{\mathrm{meas}}[a],
\end{equation}
using accessible sector Eq.~\eqref{eq:ch2.angular-accessibility-sector} and suppressed sheet
Eq.~\eqref{eq:ch2.suppressed-angular-sheet-def} from Subsection~\ref{sec:ch274}
\cite{stratton1941,harrington2001}. Equation~\eqref{eq:ch2.retained-suppressed-filter-preview}
is the Section~\ref{sec:ch29} target for Subsection~\ref{sec:ch292} HOME partition; geometric
constraints Theorem~\ref{thm:ch2.geometric-spectral-constraints} and observation rank
Eq.~\eqref{eq:ch2.angular-spectrum-observation-rank} compose independently with
Eq.~\eqref{eq:ch2.angular-spectrum-composition-law} \cite{stratton1941}. Null modes annihilate
\(a(\widehat{\mathbf{k}})\) through Eq.~\eqref{eq:ch2.angular-spectrum-nullspace-law} without
altering near reactive content \cite{harrington2001}.

\paragraph{Near-to-far transformation (preview).}
Subsection~\ref{sec:ch293} records the operator map from near-zone admissible fields to far
angular amplitudes:
\begin{equation}
\label{eq:ch2.near-to-far-transform-preview}
a(\widehat{\mathbf{k}})
=
\lim_{r\to\infty}
\mathcal{H}_{\mathrm{ff}}\big[\mathcal{R}_\omega[\Jimp]\big](r,\widehat{\mathbf{k}}),
\end{equation}
with \(\mathcal{H}_{\mathrm{ff}}\) the far-zone extractor tied to
Eq.~\eqref{eq:ch2.far-field-amplitude-integral} and Theorem
~\ref{thm:ch2.far-field-approximation} \cite{stratton1941,jackson1999,harrington2001}.
Equation~\eqref{eq:ch2.near-to-far-transform-preview} does not promote evanescent storage to
export: \(\Pi_{\mathrm{ev}}\mathcal{S}_\omega[\Jimp]\) is annihilated before
Eq.~\eqref{eq:ch2.near-to-far-transform-preview} when \(\mathcal{R}_\omega\) is unchanged
\cite{stratton1941}. Regime gate Eq.~\eqref{eq:ch2.regime-component-gate} qualifies when the
limit defines transport data \cite{harrington2001}.

\paragraph{Far-field angular accessibility (preview).}
Subsection~\ref{sec:ch294} develops accessibility on \(S^{2}\) importing Section~\ref{sec:ch01.2.1}
angular observables and Section~\ref{sec:ch01.2.3} transport criteria by reference
\cite{stratton1941,harrington2001}:
\begin{equation}
\label{eq:ch2.far-angular-accessibility-preview}
\widehat{\mathbf{k}}\in\mathcal{K}_{\mathrm{acc}}
\ \Longleftrightarrow\
a(\widehat{\mathbf{k}})\ \text{observable after}\ \mathcal{M}_\omega
\ \text{and qualified by}\ \mathcal{D}_{J,\mathrm{reg}}=1,
\end{equation}
upgrading Proposition~\ref{prop:ch2.angular-accessibility-preview} and composing with
Eq.~\eqref{eq:ch2.angular-spectrum-composition-law} \cite{stratton1941}. Equation
\eqref{eq:ch2.far-angular-accessibility-preview} is the Section~\ref{sec:ch29} closure target;
explicit plane-wave filters and measurement screens are HOME in Subsection~\ref{sec:ch294}
\cite{harrington2001}. Reciprocity Theorem~\ref{thm:ch2.radiation-reception-reciprocity} constrains
coupling on \(\mathcal{K}_{\mathrm{acc}}\) but does not enlarge it \cite{stratton1941}.

\begin{figure}[t]
\centering
\ChOneFigBlock{%
\begin{tikzpicture}[ch1fig, node distance=7mm and 10mm]
  \node[ch1box] (pw) {Plane-wave\\spectrum};
  \node[ch1box, right=12mm of pw] (filt) {Radiation\\filter};
  \node[ch1box, right=12mm of filt] (nf) {Near-to-far};
  \node[ch1box, right=12mm of nf] (acc) {$\mathcal{K}_{\mathrm{acc}}$};
  \node[ch1box, below=12mm of filt] (pi) {$\Pi_{\mathrm{rad}}$};
  \node[ch1box, below=12mm of nf] (a) {$a(\widehat{\mathbf{k}})$};
  \draw[ch1flow] (pw) -- (filt) -- (nf) -- (acc);
  \draw[ch1flow] (pi) -- (filt);
  \draw[ch1flow] (a) -- (nf);
\end{tikzpicture}%
}
\caption{Section~\ref{sec:ch29} development path: plane-wave angular spectrum and radiation
operators (Subsection~\ref{sec:ch291}), retained/suppressed/filtered components
(Subsection~\ref{sec:ch292}), near-to-far transformation (Subsection~\ref{sec:ch293}), and
far-field angular accessibility (Subsection~\ref{sec:ch294}).}
\label{fig:ch2.angular-spectrum-action-roadmap}
\end{figure}

Figure~\ref{fig:ch2.angular-spectrum-action-roadmap} orders the four subsections.
Subsection~\ref{sec:ch291} gives the Chapter~\ref{ch:02} HOME for plane-wave expansions and
radiation-operator action upgrading Eq.~\eqref{eq:ch2.plane-wave-radiation-operator-preview},
importing Eq.~\eqref{eq:ch2.outgoing-plane-wave-gate} by reference
\cite{stratton1941,harrington2001}. Subsection~\ref{sec:ch292} develops retained, suppressed, and
filtered partitions Eq.~\eqref{eq:ch2.retained-suppressed-filter-preview} on budgets from
Subsection~\ref{sec:ch274} \cite{stratton1941}. Subsection~\ref{sec:ch293} states near-to-far
maps Eq.~\eqref{eq:ch2.near-to-far-transform-preview} using far-field theorems from
Subsection~\ref{sec:ch252} without re-proving asymptotics \cite{stratton1941,jackson1999}.
Subsection~\ref{sec:ch294} closes angular accessibility Eq.~\eqref{eq:ch2.far-angular-accessibility-preview}
with USE of Sections~\ref{sec:ch01.2.1} and~\ref{sec:ch01.2.3} \cite{stratton1941,harrington2001}.

\paragraph{Composition with null spaces, reciprocity, and spectral modes.}
Angular-spectrum action composes with prior sections through
\begin{equation}
\label{eq:ch2.angular-action-composition}
\Jimp\in\ker\mathcal{R}_\omega
\ \Longrightarrow\
a(\widehat{\mathbf{k}})=0
\ \text{on export},
\qquad
\mathfrak{C}_\omega\ \text{symmetric on}\ \mathcal{K}_{\mathrm{acc}},
\end{equation}
importing Eq.~\eqref{eq:ch2.angular-spectrum-nullspace-law} and Theorem
~\ref{thm:ch2.radiation-reception-reciprocity} \cite{stratton1941,harrington2001}. Modal
coordinates Eq.~\eqref{eq:ch2.radiation-map-modal-sum} and plane-wave coordinates
Eq.~\eqref{eq:ch2.exterior-continuous-spectral-integral} describe the same
\(\operatorname{ran}\mathcal{R}_\omega\) on exterior domains; Section~\ref{sec:ch29} fixes
plane-wave-side filters, while Chapter~\ref{ch:03} Subsection~\ref{sec:ch381} supplies basis
change HOME \cite{stratton1941,harrington2001,jackson1999}. Finite-rank observation
Eq.~\eqref{eq:ch2.reconstruction-rank-cap} may remove angular content even when
\(a(\widehat{\mathbf{k}})\neq 0\) on \(\mathcal{F}_{\mathrm{rad}}\) \cite{stratton1941}.

Several constructions deliberately remain outside Section~\ref{sec:ch29}. Vector spherical
harmonic completeness and normalization are HOME in Chapter~\ref{ch:03}; Section~\ref{sec:ch29}
records operator-level plane-wave action without tabulating \(Y_\ell^m\)
\cite{stratton1941,harrington2001}. Non-radiating catalogs remain in Section~\ref{sec:ch27};
reciprocity constraints remain in Section~\ref{sec:ch28} \cite{stratton1941}. Section~\ref{sec:ch29}
does not re-derive Fredholm structure, geometric constraints, or null-space classification from
Sections~\ref{sec:ch26}--\ref{sec:ch27} \cite{harrington2001}.

Regime logic from Section~\ref{sec:ch01.2.4} is unchanged. Near-zone reactive dominance does not
create far-zone angular transport through Eq.~\eqref{eq:ch2.near-to-far-transform-preview} without
\(\mathcal{D}_{J,\mathrm{reg}}=1\) and outgoing qualification \cite{stratton1941,harrington2001}.
Evanescent sheets remain non-exporting through Eq.~\eqref{eq:ch2.rad-null-on-evanescent} even when
\(|\E|\) is large within a wavelength of \(\Omega_s\) \cite{stratton1941}. Angular-transport
claims must compose with Eq.~\eqref{eq:ch2.section29-angular-spectrum-chain}, angular budget
Eq.~\eqref{eq:ch2.angular-spectrum-composition-law}, and regime gate
Eq.~\eqref{eq:ch2.regime-operator-screen} \cite{harrington2001}.

\paragraph{Worked illustration (plane-wave gate).}
A single propagating plane wave with outgoing \(k_z\) from Eq.~\eqref{eq:ch2.outgoing-plane-wave-gate}
passes \(\Pi_{\mathrm{prop}}\) and contributes to Eq.~\eqref{eq:ch2.plane-wave-radiation-operator-preview}
\cite{stratton1941,jackson1999}. Reversing \(k_z\) selects incoming content annihilated by
\(\mathcal{O}_{\mathrm{out}}\) even when near-zone scans appear large
\cite{stratton1941,harrington2001}.

\paragraph{Worked illustration (evanescent sheet).}
Evanescent lateral wavenumbers \(k_t>k_0\) produce decaying \(k_z\) in
Eq.~\eqref{eq:ch2.outgoing-plane-wave-gate} but \(\widetilde{\E}_{\mathrm{rad}}=0\) on export through
Eq.~\eqref{eq:ch2.rad-null-on-evanescent} \cite{stratton1941}. Subsection~\ref{sec:ch293} therefore
returns \(a(\widehat{\mathbf{k}})=\mathbf{0}\) for such content even when near probes measure
substantial \(|\E|\) \cite{harrington2001}.

\paragraph{Worked illustration (NR invariance).}
Adding \(\Jimp^{\mathrm{(NR)}}\in\ker\mathcal{R}_\omega\) leaves
\(a(\widehat{\mathbf{k}})\) unchanged Eq.~\eqref{eq:ch2.angular-spectrum-nr-invariance} while
altering \(\mathcal{S}_\omega[\Jimp]\) near \(\Omega_s\) \cite{stratton1941,harrington2001}.
Pattern-matching on raw near fields cannot detect the null-space addition
\cite{stratton1941}.

\paragraph{Worked illustration (finite aperture).}
When \(\mathcal{O}_{\mathrm{meas}}\) has rank \(N\), Eq.~\eqref{eq:ch2.angular-spectrum-observation-rank}
restricts recoverable \(a(\widehat{\mathbf{k}})\) to a subset of \(\mathcal{K}_{\mathrm{acc}}\) even if
full \(\operatorname{ran}\mathcal{R}_\omega\) contains broader angular content
\cite{stratton1941,harrington2001}. Subsection~\ref{sec:ch294} states the accessibility closure
\cite{stratton1941}.

The subsection sequence begins with plane-wave angular spectra and radiation operators on the
outgoing branch, continues with retained, suppressed, and filtered components on budgets from
Section~\ref{sec:ch27}, develops near-to-far transformation importing Subsection~\ref{sec:ch252}
far-zone templates, and closes with far-field angular accessibility before worked examples in
Section~\ref{sec:ch210} and Chapter~\ref{ch:03} harmonic bases compose additional structure on the
same angular coordinates \cite{stratton1941,harrington2001,jackson1999}.

\subsection{Plane-wave angular spectrum and radiation operators}
\label{sec:ch291}
% TOC tag: [HOME]

Section~\ref{sec:ch29} previewed angular-spectrum action through
Eqs.~\eqref{eq:ch2.section29-angular-spectrum-chain}--\eqref{eq:ch2.plane-wave-radiation-operator-preview}.
Subsection~\ref{sec:ch291} gives the Chapter~\ref{ch:02} HOME for plane-wave angular spectra
and the action of \(\mathcal{R}_\omega\), \(\Pi_{\mathrm{rad}}\), and outgoing selection on
lateral spectral coefficients \cite{stratton1941,harrington2001,jackson1999}. Dispersion split
Eqs.~\eqref{eq:ch2.dispersion-relation}--\eqref{eq:ch2.longitudinal-wavenumber-split},
spectral projectors Eqs.~\eqref{eq:ch2.spectral-projectors}--\eqref{eq:ch2.rad-null-on-evanescent},
outgoing selector Eq.~\eqref{eq:ch2.outgoing-selector-def}, plane-wave gate
Eq.~\eqref{eq:ch2.outgoing-plane-wave-gate}, radiation map
Eq.~\eqref{eq:ch2.nullspace-operator-targets}, and lateral template
Eq.~\eqref{eq:ch2.angular-spectrum-preview} are presupposed; Subsection~\ref{sec:ch224} and
Section~\ref{sec:ch243} proofs are not repeated \cite{stratton1941}. Continuous angular
representation Eq.~\eqref{eq:ch2.exterior-continuous-spectral-integral} and far amplitude
Eq.~\eqref{eq:ch2.far-field-amplitude-integral} are cited as consistency targets; rigged
normalization belongs to Chapter~\ref{ch:03} Subsection~\ref{sec:ch313}
\cite{harrington2001}. Retained/suppressed partitions belong to Subsection~\ref{sec:ch292}.

\paragraph{Assumptions.}
Fix \(\omega>0\), homogeneous exterior region with symbol \(k_0\) from
Eq.~\eqref{eq:ch2.helmholtz-symbol}, compact \(\Omega_s=\operatorname{supp}\Jimp\), admissible
\(\Jimp\in\mathcal{J}_\Omega\) Eq.~\eqref{eq:ch2.source-admissibility}, and outgoing fields
\(\mathbf{u}_\omega=(\E,\H)\in\mathcal{F}_{\mathrm{SF}}\) Eq.~\eqref{eq:ch2.admissible-field-space}
\cite{stratton1941,harrington2001}. Lateral wavenumbers \((k_x,k_y)\) use
\(k_t=\sqrt{k_x^2+k_y^2}\) and outgoing longitudinal root
Eq.~\eqref{eq:ch2.outgoing-plane-wave-gate} \cite{stratton1941}. Transverse polarization satisfies
Eq.~\eqref{eq:ch2.transverse-polarization-constraint} on each spectral sheet
\cite{stratton1941,harrington2001}. Regime gate Eq.~\eqref{eq:ch2.regime-operator-screen} governs
transport interpretation of spectral outputs \cite{harrington2001}.

\paragraph{Plane-wave angular spectrum (HOME).}
On a plane \(z=z_0\) in an outgoing admissible region, define the \emph{lateral angular spectrum}
of \(\E\) by
\begin{equation}
\label{eq:ch2.plane-wave-angular-spectrum-def}
\widetilde{\E}(k_x,k_y;z_0)
=
\int_{-\infty}^{\infty}\!\!\int_{-\infty}^{\infty}
\E(x,y,z_0)\,
e^{\jj(k_x x+k_y y)}\,dx\,dy,
\end{equation}
with inverse
\begin{equation}
\label{eq:ch2.plane-wave-angular-spectrum-inverse}
\E(x,y,z_0)
=
\frac{1}{(2\pi)^{2}}
\int_{-\infty}^{\infty}\!\!\int_{-\infty}^{\infty}
\widetilde{\E}(k_x,k_y;z_0)\,
e^{-\jj(k_x x+k_y y)}\,dk_x\,dk_y,
\end{equation}
under standard Fourier conventions \cite{stratton1941,jackson1999}. Helmholtz propagation along
\(z\) multiplies each spectral component by
\begin{equation}
\label{eq:ch2.plane-wave-z-propagation}
\widetilde{\E}(k_x,k_y;z)
=
\widetilde{\E}(k_x,k_y;z_0)\,
e^{-\jj k_z(k_t)(z-z_0)},
\qquad
k_z\ \text{from Eq.~\eqref{eq:ch2.outgoing-plane-wave-gate}},
\end{equation}
for homogeneous media \cite{stratton1941}. Equations~\eqref{eq:ch2.plane-wave-angular-spectrum-def}--\eqref{eq:ch2.plane-wave-z-propagation}
upgrade preview Eq.~\eqref{eq:ch2.angular-spectrum-preview} to explicit forward/inverse
transforms with outgoing \(k_z\) selection \cite{stratton1941,harrington2001}. Algebraically,
\(\widetilde{\E}\) is the lateral Fourier transform of the transverse field; physically, each
\((k_x,k_y)\) sheet carries propagating or evanescent longitudinal behavior through
Eq.~\eqref{eq:ch2.longitudinal-wavenumber-split} \cite{stratton1941}.

\paragraph{Lateral spectral operator.}
Let \(\mathcal{F}_{\mathrm{lat}}\) denote the lateral Fourier operator in
Eq.~\eqref{eq:ch2.plane-wave-angular-spectrum-def} on admissible traces. Define the
\emph{radiative lateral spectrum} of \(\mathcal{R}_\omega[\Jimp]\) by
\begin{equation}
\label{eq:ch2.radiative-lateral-spectrum-def}
\widetilde{\E}_{\mathrm{rad}}(k_x,k_y)
=
\mathcal{F}_{\mathrm{lat}}\big[\Pi_{\mathrm{rad}}\,\mathcal{R}_\omega[\Jimp]\big](k_x,k_y),
\end{equation}
with \(\Pi_{\mathrm{rad}}=\mathcal{O}_{\mathrm{out}}\circ\Pi_{\mathrm{prop}}\) from
Eq.~\eqref{eq:ch2.rad-projection-composition} \cite{stratton1941,harrington2001}. Equation
\eqref{eq:ch2.radiative-lateral-spectrum-def} is the Subsection~\ref{sec:ch291} HOME upgrade of
preview Eq.~\eqref{eq:ch2.plane-wave-radiation-operator-preview}: radiation selection precedes
lateral transformation \cite{stratton1941}. Applying \(\mathcal{F}_{\mathrm{lat}}\) to raw
\(\mathcal{S}_\omega[\Jimp]\) without \(\Pi_{\mathrm{rad}}\) mixes reactive evanescent sheets
into reported spectra through Eq.~\eqref{eq:ch2.subspace-energy-inequality}
\cite{harrington2001,stratton1941}.

\paragraph{Radiation operator on angular spectra.}
Define the \emph{angular-spectrum radiation operator} on lateral coefficients by
\begin{equation}
\label{eq:ch2.plane-wave-radiation-operator}
\widetilde{\E}_{\mathrm{rad}}(k_x,k_y)
=
\big[\mathcal{T}_\omega\widetilde{\Jimp}\big](k_x,k_y),
\qquad
\mathcal{T}_\omega
=
\mathcal{F}_{\mathrm{lat}}\circ\Pi_{\mathrm{rad}}\circ\mathcal{R}_\omega\circ\mathcal{F}_{\mathrm{lat}}^{-1},
\end{equation}
where \(\widetilde{\Jimp}=\mathcal{F}_{\mathrm{lat}}[\Jimp]\) on the source support and
\(\mathcal{F}_{\mathrm{lat}}^{-1}\) is understood on the admissible impressed-current class
\cite{stratton1941,harrington2001}. Equation~\eqref{eq:ch2.plane-wave-radiation-operator}
packages Eq.~\eqref{eq:ch2.radiative-lateral-spectrum-def} as a spectral-domain map: sources
excite lateral modes, \(\mathcal{R}_\omega\) exports selected content, and \(\mathcal{F}_{\mathrm{lat}}\)
returns observable spectral coefficients \cite{stratton1941}. On the propagating outgoing sheet
\(k_t\le k_0\), \(\mathcal{T}_\omega\) reduces to diagonal action in \((k_x,k_y)\) modulo
source coupling; on the evanescent sheet \(k_t>k_0\),
\begin{equation}
\label{eq:ch2.evanescent-angular-null-law}
\widetilde{\E}_{\mathrm{rad}}(k_x,k_y)=\mathbf{0}
\quad\text{for}\ k_t>k_0,
\end{equation}
by Eq.~\eqref{eq:ch2.rad-null-on-evanescent} \cite{stratton1941,harrington2001}. Physically,
Eq.~\eqref{eq:ch2.evanescent-angular-null-law} states that evanescent lateral wavenumbers carry
near reactive storage but no qualified far-zone radiative spectrum \cite{stratton1941}.

\begin{figure}[t]
\centering
\ChOneFigBlock{%
\begin{tikzpicture}[ch1fig, node distance=7mm and 10mm]
  \node[ch1box] (j) {$\Jimp$};
  \node[ch1box, right=11mm of j] (r) {$\mathcal{R}_\omega$};
  \node[ch1box, right=11mm of r] (pi) {$\Pi_{\mathrm{rad}}$};
  \node[ch1box, right=11mm of pi] (f) {$\mathcal{F}_{\mathrm{lat}}$};
  \node[ch1box, right=11mm of f] (spec) {$\widetilde{\E}_{\mathrm{rad}}$};
  \node[ch1box, below=11mm of pi] (oo) {$\mathcal{O}_{\mathrm{out}}$};
  \node[ch1box, below=11mm of r] (pp) {$\Pi_{\mathrm{prop}}$};
  \draw[ch1flow] (j) -- (r) -- (pi) -- (f) -- (spec);
  \draw[ch1flow] (pp) -- (pi);
  \draw[ch1flow] (oo) -- (pi);
\end{tikzpicture}%
}
\caption{Subsection~\ref{sec:ch291} plane-wave pipeline:
Eq.~\eqref{eq:ch2.plane-wave-radiation-operator} applies \(\mathcal{R}_\omega\),
Eq.~\eqref{eq:ch2.rad-projection-composition}, then lateral transform
Eq.~\eqref{eq:ch2.plane-wave-angular-spectrum-def}.}
\label{fig:ch2.plane-wave-radiation-pipeline}
\end{figure}

Figure~\ref{fig:ch2.plane-wave-radiation-pipeline} orders operators as in
Figure~\ref{fig:ch2.outgoing-selector-pipeline}: propagation selection, outgoing gate, then
lateral spectral readout \cite{stratton1941,harrington2001}.

\begin{theorem}[Plane-wave radiation operator action]
\label{thm:ch2.plane-wave-radiation-operator-action}
Let \(\Jimp\in\mathcal{J}_\Omega\) be admissible with outgoing closure. Then:
\begin{enumerate}
\item[(i)] The radiative lateral spectrum
Eq.~\eqref{eq:ch2.radiative-lateral-spectrum-def} is well defined on
\(\Pi_{\mathrm{rad}}\,\mathcal{R}_\omega[\Jimp]\in\mathcal{F}_{\mathrm{rad}}\).
\item[(ii)] Equations~\eqref{eq:ch2.plane-wave-radiation-operator} and~\eqref{eq:ch2.evanescent-angular-null-law}
hold with \(k_z\) from Eq.~\eqref{eq:ch2.outgoing-plane-wave-gate}.
\item[(iii)] On the propagating outgoing sheet, longitudinal propagation satisfies
Eq.~\eqref{eq:ch2.plane-wave-z-propagation} and transverse constraint
Eq.~\eqref{eq:ch2.transverse-polarization-constraint}.
\item[(iv)] Far-zone amplitudes \(a_\sigma(\widehat{\mathbf{k}})\) from
Eq.~\eqref{eq:ch2.exterior-continuous-spectral-integral} agree with
Eq.~\eqref{eq:ch2.radiative-lateral-spectrum-def} under the standard
\((k_x,k_y)\leftrightarrow\widehat{\mathbf{k}}\) map on \(k_t\le k_0\).
\end{enumerate}
\end{theorem}
\begin{proof}
Part (i) follows from \(\mathcal{R}_\omega[\Jimp]\in\mathcal{F}_{\mathrm{rad}}\) and
\(\mathcal{F}_{\mathrm{lat}}\) on square-integrable traces \cite{stratton1941}. Part (ii) applies
Eqs.~\eqref{eq:ch2.rad-projection-composition}--\eqref{eq:ch2.rad-null-on-evanescent} and
Eq.~\eqref{eq:ch2.outgoing-plane-wave-gate} \cite{stratton1941,harrington2001}. Part (iii) is
Helmholtz propagation in homogeneous media \cite{stratton1941,jackson1999}. Part (iv) matches
Eq.~\eqref{eq:ch2.exterior-continuous-spectral-integral} to far reduction Theorem
~\ref{thm:ch2.far-field-approximation} without re-proving asymptotics \cite{stratton1941,jackson1999}.
\end{proof}

\paragraph{Mathematical role of Eq.~\eqref{eq:ch2.plane-wave-radiation-operator}.}
Equation~\eqref{eq:ch2.plane-wave-radiation-operator} expresses the radiation pipeline as a
spectral-domain operator \(\mathcal{T}_\omega\) acting on lateral source coefficients
\cite{stratton1941,harrington2001}. It is not a separate post-processor on near-zone
\(\mathcal{S}_\omega[\Jimp]\): \(\Pi_{\mathrm{rad}}\) is part of the definition
\cite{stratton1941}. Physically, only outgoing propagating sheets contribute to far-zone angular
content; evanescent and incoming sheets are removed by
Eqs.~\eqref{eq:ch2.outgoing-selector-def} and~\eqref{eq:ch2.evanescent-angular-null-law}
\cite{harrington2001,jackson1999}. Reciprocity Theorem~\ref{thm:ch2.lorentz-reciprocity-operator}
constrains bilinear couplings built from \(\widetilde{\E}_{\mathrm{rad}}\) but does not alter
Eq.~\eqref{eq:ch2.evanescent-angular-null-law} \cite{stratton1941}.

\paragraph{Mathematical role of Eq.~\eqref{eq:ch2.evanescent-angular-null-law}.}
Equation~\eqref{eq:ch2.evanescent-angular-null-law} is the plane-wave restatement of
Eq.~\eqref{eq:ch2.rad-null-on-evanescent} on lateral coordinates \cite{stratton1941}. Large
\(|\widetilde{\E}(k_x,k_y)|\) for \(k_t>k_0\) may appear in raw
\(\mathcal{F}_{\mathrm{lat}}[\mathcal{S}_\omega[\Jimp]]\) while
\(\widetilde{\E}_{\mathrm{rad}}=\mathbf{0}\) on export \cite{harrington2001,stratton1941}. Regime
qualification Eq.~\eqref{eq:ch2.regime-component-gate} remains mandatory before interpreting any
nonzero raw lateral spectrum as transport \cite{stratton1941}.

\begin{proposition}[Null-space invariance on lateral spectra]
\label{prop:ch2.plane-wave-null-invariance}
If \(\Jimp^{(1)}\) is arbitrary and \(\Jimp^{(2)}=\Jimp^{(1)}+\Jimp^{\mathrm{(NR)}}\) with
\(\Jimp^{\mathrm{(NR)}}\in\ker\mathcal{R}_\omega\), then
\(\widetilde{\E}_{\mathrm{rad}}\) from Eq.~\eqref{eq:ch2.radiative-lateral-spectrum-def} is
unchanged under replacement of \(\Jimp^{(1)}\) by \(\Jimp^{(2)}\)
\cite{stratton1941,harrington2001}. Raw \(\mathcal{F}_{\mathrm{lat}}[\mathcal{S}_\omega[\Jimp]]\)
may change through reactive content Eq.~\eqref{eq:ch2.nr-image-of-null-space}
\cite{stratton1941}.
\end{proposition}
\begin{proof}
\(\mathcal{R}_\omega[\Jimp^{(2)}]=\mathcal{R}_\omega[\Jimp^{(1)}]\) by
Eq.~\eqref{eq:ch2.nr-superposition-law} \cite{stratton1941,harrington2001}.
\end{proof}

\begin{proposition}[Consistency with continuous angular representation]
\label{prop:ch2.plane-wave-continuous-consistency}
On exterior domains with outgoing closure, amplitudes \(a_\sigma(\widehat{\mathbf{k}})\) in
Eq.~\eqref{eq:ch2.exterior-continuous-spectral-integral} and lateral spectrum
Eq.~\eqref{eq:ch2.radiative-lateral-spectrum-def} describe the same
\(\operatorname{ran}\mathcal{R}_\omega\) content under change of variables
\((k_x,k_y)\mapsto\widehat{\mathbf{k}}\) on the propagating sheet \(k_t\le k_0\)
\cite{stratton1941,harrington2001,jackson1999}. Normalization constants are fixed in
Subsection~\ref{sec:ch313} \cite{stratton1941}.
\end{proposition}
\begin{proof}
Both representations diagonalize outgoing propagation on the propagating sheet; agreement on
far-zone reduction follows Theorem~\ref{thm:ch2.plane-wave-radiation-operator-action}(iv) and
Eq.~\eqref{eq:ch2.exterior-continuous-spectral-integral} \cite{stratton1941,jackson1999}.
\end{proof}

\paragraph{Relation to modal spectra.}
Modal expansion Eq.~\eqref{eq:ch2.radiation-map-modal-sum} and plane-wave expansion
Eq.~\eqref{eq:ch2.plane-wave-radiation-operator} describe the same radiative image on
\(\mathcal{F}_{\mathrm{rad}}\) under unitary change of basis on the propagating sheet
\cite{stratton1941,harrington2001}. Subsection~\ref{sec:ch291} fixes plane-wave-side operators;
Chapter~\ref{ch:03} Subsection~\ref{sec:ch381} supplies explicit harmonic transforms without
re-tabulating Eq.~\eqref{eq:ch2.plane-wave-radiation-operator} \cite{stratton1941}. Angular
budget Eq.~\eqref{eq:ch2.angular-spectrum-composition-law} constrains supports of
\(\widetilde{\E}_{\mathrm{rad}}\) independently of basis choice \cite{harrington2001}.

\paragraph{Boundary of validity.}
Theorem~\ref{thm:ch2.plane-wave-radiation-operator-action} assumes linear homogeneous exterior
regions at fixed \(\omega\), outgoing closure, and admissible \(\mathcal{J}_\Omega\)
\cite{stratton1941,harrington2001}. Piecewise homogeneous media require patchwise
Eq.~\eqref{eq:ch2.plane-wave-z-propagation} with interface matching imported from
Section~\ref{sec:ch234} \cite{stratton1941}. Cavity standing-wave spectra are excluded: lateral
plane-wave expansions on bounded domains do not replace discrete cavity books from
Eq.~\eqref{eq:ch2.cavity-discrete-radiation-spectrum} \cite{jackson1999,harrington2001}. Loss
modifies \(k_0\) and \(k_z\) with complex symbols; evanescent null law
Eq.~\eqref{eq:ch2.evanescent-angular-null-law} remains on export \cite{stratton1941}.

\paragraph{Worked illustration (single propagating mode).}
A single outgoing plane wave with \(k_t\le k_0\) and \(k_z\) from
Eq.~\eqref{eq:ch2.outgoing-plane-wave-gate} yields nonzero
\(\widetilde{\E}_{\mathrm{rad}}(k_x,k_y)\) at one spectral point and contributes to
\(a(\widehat{\mathbf{k}})\) through Proposition~\ref{prop:ch2.plane-wave-continuous-consistency}
\cite{stratton1941,jackson1999}. Reversing \(k_z\) selects incoming content annihilated by
\(\mathcal{O}_{\mathrm{out}}\) \cite{stratton1941,harrington2001}.

\paragraph{Worked illustration (evanescent sheet).}
For \(k_t>k_0\), raw \(\mathcal{F}_{\mathrm{lat}}[\mathcal{S}_\omega[\Jimp]]\) may be nonzero
while Eq.~\eqref{eq:ch2.evanescent-angular-null-law} gives
\(\widetilde{\E}_{\mathrm{rad}}=\mathbf{0}\) \cite{stratton1941,harrington2001}. Near-field scans
without \(\Pi_{\mathrm{rad}}\) therefore overstate exportable angular content
\cite{stratton1941}.

\paragraph{Worked illustration (NR superposition).}
Adding \(\Jimp^{\mathrm{(NR)}}\in\ker\mathcal{R}_\omega\) leaves
\(\widetilde{\E}_{\mathrm{rad}}\) unchanged by Proposition~\ref{prop:ch2.plane-wave-null-invariance}
and Eq.~\eqref{eq:ch2.angular-spectrum-nr-invariance} \cite{stratton1941,harrington2001}. Pattern
fitting on raw near spectra cannot detect the invisible addition \cite{stratton1941}.

\paragraph{Worked illustration (extended source).}
For distributed \(\Jimp\), Eq.~\eqref{eq:ch2.plane-wave-radiation-operator} superposes lateral
modes through linearity Eq.~\eqref{eq:ch2.radiation-linearity}; far amplitude
Eq.~\eqref{eq:ch2.far-field-amplitude-integral} is the \(r\to\infty\) envelope of the same
spectral content Theorem~\ref{thm:ch2.far-field-approximation} \cite{stratton1941,jackson1999}.
Subsection~\ref{sec:ch293} develops the near-to-far map explicitly \cite{stratton1941}.

\paragraph{Closing summary.}
Subsection~\ref{sec:ch291} establishes the Chapter~\ref{ch:02} HOME for plane-wave angular
spectrum Eqs.~\eqref{eq:ch2.plane-wave-angular-spectrum-def}--\eqref{eq:ch2.plane-wave-z-propagation},
radiative lateral spectrum Eq.~\eqref{eq:ch2.radiative-lateral-spectrum-def}, angular-spectrum
radiation operator Eq.~\eqref{eq:ch2.plane-wave-radiation-operator}, evanescent null law
Eq.~\eqref{eq:ch2.evanescent-angular-null-law}, Theorem
~\ref{thm:ch2.plane-wave-radiation-operator-action}, and Propositions
~\ref{prop:ch2.plane-wave-null-invariance}--\ref{prop:ch2.plane-wave-continuous-consistency},
upgrading preview Eq.~\eqref{eq:ch2.plane-wave-radiation-operator-preview}
\cite{stratton1941,harrington2001,jackson1999}. Subsection~\ref{sec:ch292} partitions retained,
suppressed, and filtered components; Subsection~\ref{sec:ch293} develops near-to-far
transformation \cite{stratton1941}.
\subsection{Retained, suppressed, and filtered components}
\label{sec:ch292}
% TOC tag: [HOME]

Section~\ref{sec:ch29} previewed angular partitioning through
Eq.~\eqref{eq:ch2.retained-suppressed-filter-preview} and linked it to accessible sector
Eq.~\eqref{eq:ch2.angular-accessibility-sector} and suppressed sheet
Eq.~\eqref{eq:ch2.suppressed-angular-sheet-def} from Subsection~\ref{sec:ch274}
\cite{stratton1941,harrington2001}. Subsection~\ref{sec:ch292} gives the Chapter~\ref{ch:02}
HOME for retained, suppressed, and observation-filtered angular components on the outgoing
branch established in Subsection~\ref{sec:ch291}. Plane-wave lateral spectra
Eqs.~\eqref{eq:ch2.radiative-lateral-spectrum-def}--\eqref{eq:ch2.plane-wave-radiation-operator},
evanescent null law Eq.~\eqref{eq:ch2.evanescent-angular-null-law}, angular budget
Eq.~\eqref{eq:ch2.angular-spectrum-composition-law}, and observation rank
Eq.~\eqref{eq:ch2.angular-spectrum-observation-rank} are presupposed; Subsection~\ref{sec:ch274}
proofs are not repeated \cite{stratton1941}. Continuous amplitude
\(a(\widehat{\mathbf{k}})\) from Eq.~\eqref{eq:ch2.exterior-continuous-spectral-integral} and
measurement map Eq.~\eqref{eq:ch2.reconstruction-observation-map} supply the angular and
laboratory coordinates in which the partition is stated \cite{stratton1941,harrington2001}.
Near-to-far extraction belongs to Subsection~\ref{sec:ch293}; explicit accessibility filters on
\(S^{2}\) belong to Subsection~\ref{sec:ch294} \cite{stratton1941}.

\paragraph{Assumptions.}
Fix \(\omega>0\), admissible \(\Jimp\in\mathcal{J}_\Omega\), outgoing radiative field
\(\mathbf{u}_\omega=\mathcal{R}_\omega[\Jimp]\in\mathcal{F}_{\mathrm{rad}}\), and angular
amplitude \(a(\widehat{\mathbf{k}})\) from
Eq.~\eqref{eq:ch2.exterior-continuous-spectral-integral} on the exporting transverse sheet
\cite{stratton1941,harrington2001,jackson1999}. Accessible sector
\(\mathcal{K}_{\mathrm{acc}}\) Eq.~\eqref{eq:ch2.angular-accessibility-sector} and suppressed
sheet \(\mathcal{K}_{\mathrm{sup}}=S^{2}\setminus\mathcal{K}_{\mathrm{acc}}\) from
Eq.~\eqref{eq:ch2.suppressed-angular-sheet-def} are fixed by geometry and symmetry of
\((\Omega_s,\Gamma,\varepsilon,\mu)\) \cite{stratton1941}. Finite-rank observation
\(\mathcal{M}_\omega=\mathcal{O}_{\mathrm{meas}}\circ\Lambda_{\mathrm{rad}}\) from
Eq.~\eqref{eq:ch2.reconstruction-observation-map} is presupposed; regime gate
Eq.~\eqref{eq:ch2.regime-operator-screen} qualifies transport interpretation of any reported
component \cite{harrington2001,stratton1941}.

\paragraph{Retained angular amplitude (HOME).}
Define the \emph{retained angular amplitude} as the exportable content supported on the
accessible sheet:
\begin{equation}
\label{eq:ch2.retained-angular-amplitude-def}
a_{\mathrm{ret}}(\widehat{\mathbf{k}})
=
\chi_{\mathcal{K}_{\mathrm{acc}}}(\widehat{\mathbf{k}})\,
a(\widehat{\mathbf{k}}),
\qquad
\chi_{\mathcal{K}_{\mathrm{acc}}}(\widehat{\mathbf{k}})
=
\begin{cases}
1, & \widehat{\mathbf{k}}\in\mathcal{K}_{\mathrm{acc}},\\
0, & \widehat{\mathbf{k}}\in\mathcal{K}_{\mathrm{sup}},
\end{cases}
\end{equation}
with \(a(\widehat{\mathbf{k}})\) from
Eq.~\eqref{eq:ch2.exterior-continuous-spectral-integral} after
Eq.~\eqref{eq:ch2.plane-wave-radiation-operator} \cite{stratton1941,harrington2001}.
Equation~\eqref{eq:ch2.retained-angular-amplitude-def} is the Subsection~\ref{sec:ch292} HOME
upgrade of the retained factor in preview
Eq.~\eqref{eq:ch2.retained-suppressed-filter-preview}: only directions not annihilated by
\(\mathcal{C}_{\mathrm{geo}}\) Proposition~\ref{prop:ch2.angular-accessibility-preview} carry
qualified transport data \cite{stratton1941}. Mathematically, \(a_{\mathrm{ret}}\) is the pointwise
product of the full angular measure with the accessible indicator; physically, it is the portion
of the far-zone angular distribution that geometry permits to propagate as export after radiation
selection \cite{harrington2001,jackson1999}. Null-space additions
Eq.~\eqref{eq:ch2.angular-spectrum-nr-invariance} leave \(a_{\mathrm{ret}}\) unchanged
\cite{stratton1941}.

\paragraph{Suppressed angular amplitude (HOME).}
Define the \emph{suppressed angular amplitude} as the complementary sheet content:
\begin{equation}
\label{eq:ch2.suppressed-angular-amplitude-def}
a_{\mathrm{sup}}(\widehat{\mathbf{k}})
=
\chi_{\mathcal{K}_{\mathrm{sup}}}(\widehat{\mathbf{k}})\,
a(\widehat{\mathbf{k}}),
\qquad
\chi_{\mathcal{K}_{\mathrm{sup}}}=1-\chi_{\mathcal{K}_{\mathrm{acc}}}
\ \text{a.e.\ on}\ S^{2},
\end{equation}
using suppressed sheet Eq.~\eqref{eq:ch2.suppressed-angular-sheet-def}
\cite{stratton1941,harrington2001}. Equation~\eqref{eq:ch2.suppressed-angular-amplitude-def}
completes preview Eq.~\eqref{eq:ch2.retained-suppressed-filter-preview}: for radiating export,
geometric forbidding Eq.~\eqref{eq:ch2.geometric-modal-annihilation} typically forces
\(a_{\mathrm{sup}}=\mathbf{0}\) even when near-zone scans show energy in related directions
\cite{stratton1941}. Suppression is therefore not synonymous with
\(\Jimp\in\ker\mathcal{R}_\omega\): Eq.~\eqref{eq:ch2.angular-spectrum-nullspace-law} annihilates
\(a\) everywhere, whereas Eq.~\eqref{eq:ch2.suppressed-angular-amplitude-def} removes selected
directions while \(a_{\mathrm{ret}}\) may remain nonzero \cite{harrington2001,stratton1941}.
Evanescent lateral content Eq.~\eqref{eq:ch2.evanescent-angular-null-law} does not populate
\(a_{\mathrm{sup}}\) on export: \(\Pi_{\mathrm{ev}}\) is removed before angular readout
\cite{stratton1941}.

\paragraph{Filtered angular amplitude (HOME).}
Laboratory observation applies a finite-rank screen to far-zone angular content. Define the
\emph{filtered angular amplitude} by
\begin{equation}
\label{eq:ch2.filtered-angular-amplitude-def}
a_{\mathrm{filt}}(\widehat{\mathbf{k}})
=
\big[\mathcal{P}_{\mathrm{meas}}\,a\big](\widehat{\mathbf{k}}),
\qquad
\mathcal{P}_{\mathrm{meas}}
=
\mathcal{O}_{\mathrm{meas}}\circ\mathcal{O}_{\mathrm{far}}\circ\mathcal{H}_{\mathrm{ang}}^{-1},
\end{equation}
where \(\mathcal{H}_{\mathrm{ang}}\) maps \(a(\widehat{\mathbf{k}})\) to the corresponding
far-zone transverse field on \(S^{2}\) and \(\mathcal{O}_{\mathrm{far}}\) is the far extractor
from Eq.~\eqref{eq:ch2.reconstruction-observation-map} \cite{stratton1941,harrington2001}.
Equation~\eqref{eq:ch2.filtered-angular-amplitude-def} is the angular-side restatement of the
third line in preview Eq.~\eqref{eq:ch2.retained-suppressed-filter-preview}. When
\(\mathcal{O}_{\mathrm{meas}}\) is a finite set of plane-wave or modal probes,
\(\mathcal{P}_{\mathrm{meas}}\) has rank at most \(N\) and
\(\ker\mathcal{P}_{\mathrm{meas}}\supseteq\ker\mathcal{O}_{\mathrm{meas}}\big|_{\operatorname{ran}\mathcal{R}_\omega}\)
Eq.~\eqref{eq:ch2.observation-null-space} \cite{stratton1941,harrington2001}. Physically,
\(a_{\mathrm{filt}}\) is what a finite instrument resolves after aperture, noise, and processing;
it may vanish on subsets of \(\mathcal{K}_{\mathrm{acc}}\) where \(a_{\mathrm{ret}}\) is nonzero
\cite{harrington2001}.

\paragraph{Angular component partition (HOME).}
Combining Eqs.~\eqref{eq:ch2.retained-angular-amplitude-def}--\eqref{eq:ch2.filtered-angular-amplitude-def}
yields the Chapter~\ref{ch:02} partition
\begin{equation}
\label{eq:ch2.angular-component-partition}
a(\widehat{\mathbf{k}})
=
a_{\mathrm{ret}}(\widehat{\mathbf{k}})
+
a_{\mathrm{sup}}(\widehat{\mathbf{k}}),
\qquad
a_{\mathrm{filt}}
=
\mathcal{P}_{\mathrm{meas}}\big[a_{\mathrm{ret}}+a_{\mathrm{sup}}\big],
\end{equation}
with \(a_{\mathrm{ret}}\) and \(a_{\mathrm{sup}}\) supported on disjoint sheets
\(\mathcal{K}_{\mathrm{acc}}\) and \(\mathcal{K}_{\mathrm{sup}}\)
\cite{stratton1941,harrington2001}. Equation~\eqref{eq:ch2.angular-component-partition} upgrades
preview Eq.~\eqref{eq:ch2.retained-suppressed-filter-preview} to explicit HOME definitions and
inserts the measurement projector \(\mathcal{P}_{\mathrm{meas}}\). The first line is an angular
decomposition of export; the second line states that observation acts on the sum after radiation
selection, not on raw near spectra \cite{stratton1941}. Budget
Eq.~\eqref{eq:ch2.angular-spectrum-composition-law} constrains
\(\operatorname{supp}(a)\subseteq\mathcal{K}_{\mathrm{acc}}\cap\{\text{resolvable by}\ \mathcal{M}_\omega\}\),
so \(a_{\mathrm{filt}}\) cannot recover directions outside that intersection even when raw
lateral transforms appear rich \cite{harrington2001,stratton1941}.

\paragraph{Retained and filtered lateral spectra.}
On the plane-wave side of Subsection~\ref{sec:ch291}, define the \emph{retained lateral spectrum}
\begin{equation}
\label{eq:ch2.retained-lateral-spectrum-def}
\widetilde{\E}_{\mathrm{ret}}(k_x,k_y)
=
\mathcal{F}_{\mathrm{lat}}\big[\mathcal{H}_{\mathrm{ang}}[a_{\mathrm{ret}}]\big](k_x,k_y)
\quad\text{on}\ k_t\le k_0,
\end{equation}
and the \emph{filtered lateral spectrum}
\begin{equation}
\label{eq:ch2.filtered-lateral-spectrum-def}
\widetilde{\E}_{\mathrm{filt}}(k_x,k_y)
=
\mathcal{F}_{\mathrm{lat}}\big[\mathcal{H}_{\mathrm{ang}}[a_{\mathrm{filt}}]\big](k_x,k_y),
\end{equation}
with \(\mathcal{F}_{\mathrm{lat}}\) from Eq.~\eqref{eq:ch2.plane-wave-angular-spectrum-def}
\cite{stratton1941,harrington2001}. Equations~\eqref{eq:ch2.retained-lateral-spectrum-def}--\eqref{eq:ch2.filtered-lateral-spectrum-def}
connect Subsection~\ref{sec:ch291} lateral readout to the angular partition without re-displaying
Eq.~\eqref{eq:ch2.plane-wave-radiation-operator}: \(\widetilde{\E}_{\mathrm{rad}}\) from
Eq.~\eqref{eq:ch2.radiative-lateral-spectrum-def} agrees with
\(\widetilde{\E}_{\mathrm{ret}}\) on the propagating outgoing sheet when
\(a_{\mathrm{sup}}=\mathbf{0}\) \cite{stratton1941,jackson1999}. On the evanescent sheet
\(k_t>k_0\), retained and filtered lateral spectra vanish through
Eq.~\eqref{eq:ch2.evanescent-angular-null-law} independently of raw
\(\mathcal{F}_{\mathrm{lat}}[\mathcal{S}_\omega[\Jimp]]\) \cite{stratton1941,harrington2001}.
Regime gate Eq.~\eqref{eq:ch2.regime-component-gate} applies to
\(\widetilde{\E}_{\mathrm{filt}}\) before any transport claim \cite{harrington2001}.

\paragraph{Suppressed lateral null law.}
On \(\mathcal{K}_{\mathrm{sup}}\),
\begin{equation}
\label{eq:ch2.suppressed-lateral-null-law}
a_{\mathrm{sup}}(\widehat{\mathbf{k}})=\mathbf{0}
\ \Longrightarrow\
\widetilde{\E}_{\mathrm{ret}}(k_x,k_y)=\widetilde{\E}_{\mathrm{rad}}(k_x,k_y)
\ \text{on}\ k_t\le k_0,
\end{equation}
and any lateral energy attributed to \(\mathcal{K}_{\mathrm{sup}}\) after
Eq.~\eqref{eq:ch2.plane-wave-radiation-operator} is removed by geometry, not by
\(\mathcal{P}_{\mathrm{meas}}\) \cite{stratton1941}. Equation~\eqref{eq:ch2.suppressed-lateral-null-law}
restates that suppression precedes observation: finite aperture cannot resurrect directions
forbidden by Eq.~\eqref{eq:ch2.geometric-modal-annihilation}
\cite{harrington2001,stratton1941}. Reciprocity Theorem~\ref{thm:ch2.radiation-reception-reciprocity}
constrains bilinear couplings on \(\mathcal{K}_{\mathrm{acc}}\) but does not enlarge
\(\operatorname{supp}(a_{\mathrm{ret}})\) \cite{stratton1941}.

\begin{figure}[t]
\centering
\ChOneFigBlock{%
\begin{tikzpicture}[ch1fig, node distance=7mm and 10mm]
  \node[ch1box] (a) {$a(\widehat{\mathbf{k}})$};
  \node[ch1box, right=11mm of a] (ret) {$a_{\mathrm{ret}}$};
  \node[ch1box, right=11mm of ret] (filt) {$a_{\mathrm{filt}}$};
  \node[ch1box, below=11mm of a] (sup) {$a_{\mathrm{sup}}$};
  \node[ch1box, below=11mm of ret] (kacc) {$\mathcal{K}_{\mathrm{acc}}$};
  \node[ch1box, below=11mm of sup] (ksup) {$\mathcal{K}_{\mathrm{sup}}$};
  \node[ch1box, left=11mm of a] (rad) {$\Pi_{\mathrm{rad}}\mathcal{R}_\omega$};
  \node[ch1box, right=11mm of filt] (lat) {$\widetilde{\E}_{\mathrm{filt}}$};
  \draw[ch1flow] (rad) -- (a);
  \draw[ch1flow] (a) -- (ret) -- (filt) -- (lat);
  \draw[ch1flow] (a) -- (sup);
  \draw[ch1flow] (kacc) -- (ret);
  \draw[ch1flow] (ksup) -- (sup);
\end{tikzpicture}%
}
\caption{Subsection~\ref{sec:ch292} partition pipeline:
Eq.~\eqref{eq:ch2.angular-component-partition} splits \(a(\widehat{\mathbf{k}})\) into retained
and suppressed sheets, then applies observation filter
Eq.~\eqref{eq:ch2.filtered-angular-amplitude-def} before lateral readout
Eq.~\eqref{eq:ch2.filtered-lateral-spectrum-def}.}
\label{fig:ch2.retained-suppressed-filter-pipeline}
\end{figure}

Figure~\ref{fig:ch2.retained-suppressed-filter-pipeline} orders radiation selection, angular
partition, measurement filter, and plane-wave readout. Suppression acts on
\(\mathcal{K}_{\mathrm{sup}}\) before \(\mathcal{P}_{\mathrm{meas}}\); evanescent sheets are
already null in \(\widetilde{\E}_{\mathrm{rad}}\) through
Eq.~\eqref{eq:ch2.evanescent-angular-null-law} \cite{stratton1941,harrington2001}.

\begin{theorem}[Retained, suppressed, and filtered angular partition]
\label{thm:ch2.angular-component-partition}
Let \(\Jimp\in\mathcal{J}_\Omega\) be admissible with outgoing closure and let
\(a(\widehat{\mathbf{k}})\) be the angular amplitude from
Eq.~\eqref{eq:ch2.exterior-continuous-spectral-integral}. Define
\(a_{\mathrm{ret}}\), \(a_{\mathrm{sup}}\), and \(a_{\mathrm{filt}}\) by
Eqs.~\eqref{eq:ch2.retained-angular-amplitude-def}--\eqref{eq:ch2.filtered-angular-amplitude-def}.
Then:
\begin{enumerate}
\item[(i)] Eq.~\eqref{eq:ch2.angular-component-partition} holds with
\(a_{\mathrm{ret}}\) supported on \(\mathcal{K}_{\mathrm{acc}}\) and
\(a_{\mathrm{sup}}\) supported on \(\mathcal{K}_{\mathrm{sup}}\).
\item[(ii)] \(\operatorname{supp}(a)\subseteq\mathcal{K}_{\mathrm{acc}}\) for radiating export
implies \(a_{\mathrm{sup}}=\mathbf{0}\) and \(a=a_{\mathrm{ret}}\) on the exporting sheet.
\item[(iii)] Observation rank Eq.~\eqref{eq:ch2.angular-spectrum-observation-rank} bounds
\(\operatorname{supp}(a_{\mathrm{filt}})\) independently of \(\ker\mathcal{R}_\omega\).
\item[(iv)] Lateral spectra Eqs.~\eqref{eq:ch2.retained-lateral-spectrum-def}--\eqref{eq:ch2.filtered-lateral-spectrum-def}
agree with Eq.~\eqref{eq:ch2.radiative-lateral-spectrum-def} on \(k_t\le k_0\) when
\(a_{\mathrm{sup}}=\mathbf{0}\) and \(\mathcal{P}_{\mathrm{meas}}=\mathrm{id}\) on
\(\operatorname{supp}(a_{\mathrm{ret}})\).
\end{enumerate}
\end{theorem}
\begin{proof}
Part (i) follows from definitions and disjoint partition
\(S^{2}=\mathcal{K}_{\mathrm{acc}}\cup\mathcal{K}_{\mathrm{sup}}\) up to measure zero
\cite{stratton1941}. Part (ii) uses
Eq.~\eqref{eq:ch2.angular-spectrum-composition-law} and suppressed sheet
Eq.~\eqref{eq:ch2.suppressed-angular-sheet-def} \cite{harrington2001,stratton1941}. Part (iii)
applies Eq.~\eqref{eq:ch2.angular-spectrum-observation-rank} and
Eq.~\eqref{eq:ch2.reconstruction-rank-cap} \cite{stratton1941,harrington2001}. Part (iv) follows
from Proposition~\ref{prop:ch2.plane-wave-continuous-consistency} and Theorem
~\ref{thm:ch2.plane-wave-radiation-operator-action}(iv) \cite{stratton1941,jackson1999}.
\end{proof}

\begin{proposition}[Independence of suppression and source nullity]
\label{prop:ch2.retained-suppressed-mechanism-independence}
\(a_{\mathrm{sup}}=\mathbf{0}\) on \(\mathcal{K}_{\mathrm{sup}}\) for radiating
\(\Jimp\notin\ker\mathcal{R}_\omega\) does not imply
\(\Jimp\in\ker\mathcal{R}_\omega\), and conversely
\(\Jimp\in\ker\mathcal{R}_\omega\) forces \(a_{\mathrm{ret}}=a_{\mathrm{sup}}=\mathbf{0}\) through
Eq.~\eqref{eq:ch2.angular-spectrum-nullspace-law} \cite{stratton1941,harrington2001}. Geometric
suppression Eq.~\eqref{eq:ch2.suppressed-angular-amplitude-def} and radiation nullity
Eq.~\eqref{eq:ch2.angular-spectrum-nullspace-law} are distinct mechanisms composing through
Eq.~\eqref{eq:ch2.angular-spectrum-composition-law} \cite{stratton1941}.
\end{proposition}
\begin{proof}
Theorem~\ref{thm:ch2.angular-spectrum-null-consequences}(i) handles kernel membership;
Proposition~\ref{prop:ch2.angular-accessibility-preview} and
Eq.~\eqref{eq:ch2.geometric-modal-annihilation} handle accessible-sector restriction without
kernel membership \cite{stratton1941,harrington2001}.
\end{proof}

\begin{proposition}[Filtered component rank bound]
\label{prop:ch2.filtered-observation-rank-bound}
The filtered angular amplitude satisfies
\begin{equation}
\label{eq:ch2.filtered-rank-bound}
\dim\operatorname{span}\{a_{\mathrm{filt}}(\widehat{\mathbf{k}})\}
\le
\mathrm{rank}(\mathcal{M}_\omega)
\le
\min\{N,\,\dim\operatorname{ran}\mathcal{R}_\omega\},
\end{equation}
with \(N\) the number of measurement channels in
Eq.~\eqref{eq:ch2.finite-rank-observation} \cite{stratton1941,harrington2001}. Strict inequality
holds when \(\ker\mathcal{O}_{\mathrm{meas}}\) is nontrivial on
\(\operatorname{ran}\mathcal{R}_\omega\) Eq.~\eqref{eq:ch2.observation-null-space}
\cite{stratton1941}.
\end{proposition}
\begin{proof}
\(\mathcal{P}_{\mathrm{meas}}\) factors through \(\mathcal{M}_\omega\); rank bounds follow from
Eqs.~\eqref{eq:ch2.reconstruction-rank-cap} and~\eqref{eq:ch2.angular-spectrum-observation-rank}
\cite{stratton1941,harrington2001}.
\end{proof}

\paragraph{Mathematical role of Eq.~\eqref{eq:ch2.angular-component-partition}.}
Equation~\eqref{eq:ch2.angular-component-partition} separates three stages that are often
conflated in angular-spectrum discussions: radiation export, geometric accessibility, and finite
observation \cite{stratton1941,harrington2001}. The retained factor encodes what the source radiates
on \(\mathcal{K}_{\mathrm{acc}}\); the suppressed factor marks directions removed before
measurement; the filtered factor applies \(\mathcal{P}_{\mathrm{meas}}\) \cite{stratton1941}.
Plane-wave coefficients \(\widetilde{\E}_{\mathrm{rad}}\) from
Eq.~\eqref{eq:ch2.radiative-lateral-spectrum-def} describe the same retained content in lateral
coordinates when \(a_{\mathrm{sup}}=\mathbf{0}\) \cite{jackson1999,harrington2001}. Null additions
Proposition~\ref{prop:ch2.plane-wave-null-invariance} alter neither \(a_{\mathrm{ret}}\) nor
\(a_{\mathrm{filt}}\) when \(\mathcal{M}_\omega\) is unchanged \cite{stratton1941}.

\paragraph{Mathematical role of Eq.~\eqref{eq:ch2.filtered-angular-amplitude-def}.}
Equation~\eqref{eq:ch2.filtered-angular-amplitude-def} makes explicit that instruments observe
\(\mathcal{P}_{\mathrm{meas}}[a]\), not raw \(\mathcal{F}_{\mathrm{lat}}[\mathcal{S}_\omega[\Jimp]]\)
\cite{stratton1941,harrington2001}. Rank cap Eq.~\eqref{eq:ch2.filtered-rank-bound} prevents
attributing unobserved angular structure to source nullity: a pattern gap on
\(\mathcal{K}_{\mathrm{acc}}\) may arise from \(\mathcal{P}_{\mathrm{meas}}\) alone
\cite{harrington2001}. Reciprocity Theorem~\ref{thm:ch2.radiation-reception-reciprocity} constrains
coupling among observed channels but does not increase
\(\dim\operatorname{span}\{a_{\mathrm{filt}}\}\) \cite{stratton1941}. Regime qualification
Eq.~\eqref{eq:ch2.regime-operator-screen} remains mandatory before interpreting
\(a_{\mathrm{filt}}\) as transport \cite{harrington2001}.

\paragraph{Relation to modal and plane-wave budgets.}
Modal reportability Eq.~\eqref{eq:ch2.constraint-composition-law} and angular budget
Eq.~\eqref{eq:ch2.angular-spectrum-composition-law} are the modal and continuous forms of the
same constraint chain Subsection~\ref{sec:ch292} displays in
Eq.~\eqref{eq:ch2.angular-component-partition} \cite{stratton1941,harrington2001}. Theorem
~\ref{thm:ch2.plane-wave-radiation-operator-action} supplies lateral coordinates; Theorem
~\ref{thm:ch2.angular-component-partition} supplies angular-sheet coordinates; Chapter~\ref{ch:03}
Subsection~\ref{sec:ch381} will change basis without altering
\(\operatorname{supp}(a_{\mathrm{ret}})\) \cite{stratton1941}. Inverse obstruction
Proposition~\ref{prop:ch2.angular-spectrum-inverse-obstruction} applies to
\(a_{\mathrm{filt}}\), not to raw lateral transforms \cite{harrington2001}.

\paragraph{Boundary of validity.}
Theorem~\ref{thm:ch2.angular-component-partition} assumes linear media, fixed \(\omega\), outgoing
closure, and stationary \(\mathcal{C}_{\mathrm{geo}}\) \cite{stratton1941,harrington2001}. Time-varying
or nonlinear \(\mathcal{P}_{\mathrm{meas}}\) requires separate treatment; the static rank bound
Eq.~\eqref{eq:ch2.filtered-rank-bound} is an upper limit \cite{stratton1941}. Loss modifies
amplitudes but not the partition logic on \(\mathcal{K}_{\mathrm{acc}}\cup\mathcal{K}_{\mathrm{sup}}\)
\cite{stratton1941}. Cavity and waveguide spectra remain outside the exterior partition; discrete
books Eq.~\eqref{eq:ch2.cavity-discrete-radiation-spectrum} are not substituted for
Eq.~\eqref{eq:ch2.angular-component-partition} on unbounded domains \cite{jackson1999,harrington2001}.

\paragraph{Worked illustration (ground-plane suppression).}
For a horizontal dipole above a PEC ground, \(\mathcal{K}_{\mathrm{sup}}\) contains directions
forbidden by image constraints Eq.~\eqref{eq:ch2.boundary-image-constraint}; \(a_{\mathrm{ret}}\)
is supported on the accessible hemisphere while \(a_{\mathrm{sup}}=\mathbf{0}\) on export
\cite{stratton1941,jackson1999}. Raw near scans may show energy in suppressed directions without
contributing to \(a_{\mathrm{ret}}\) \cite{harrington2001}.

\paragraph{Worked illustration (finite aperture filtering).}
When \(\mathcal{O}_{\mathrm{meas}}\) is a finite array, \(a_{\mathrm{filt}}\) may vanish on
subsets of \(\mathcal{K}_{\mathrm{acc}}\) where \(a_{\mathrm{ret}}\neq\mathbf{0}\) because
Eq.~\eqref{eq:ch2.filtered-rank-bound} limits resolvable directions
\cite{stratton1941,harrington2001}. Pattern nulls from aperture therefore need not indicate
\(\Jimp\in\ker\mathcal{R}_\omega\) \cite{stratton1941}.

\paragraph{Worked illustration (NR invisibility).}
Adding \(\Jimp^{\mathrm{(NR)}}\in\ker\mathcal{R}_\omega\) leaves \(a_{\mathrm{ret}}\) and
\(a_{\mathrm{filt}}\) unchanged by Eqs.~\eqref{eq:ch2.angular-spectrum-nr-invariance} and
Proposition~\ref{prop:ch2.plane-wave-null-invariance} \cite{stratton1941,harrington2001}. Near-field
spectral fitting on raw lateral data cannot detect the addition \cite{stratton1941}.

\paragraph{Worked illustration (evanescent vs.\ suppressed).}
Large raw \(\mathcal{F}_{\mathrm{lat}}[\mathcal{S}_\omega[\Jimp]]\) on \(k_t>k_0\) does not
produce \(a_{\mathrm{sup}}\) on export: evanescent content is removed by
Eq.~\eqref{eq:ch2.evanescent-angular-null-law} before angular partition
\cite{stratton1941,harrington2001}. Confusing evanescent storage with suppressed-sheet content
violates regime gate Eq.~\eqref{eq:ch2.regime-component-gate} \cite{harrington2001}.

\paragraph{Closing summary.}
Subsection~\ref{sec:ch292} establishes the Chapter~\ref{ch:02} HOME for retained angular amplitude
Eq.~\eqref{eq:ch2.retained-angular-amplitude-def}, suppressed angular amplitude
Eq.~\eqref{eq:ch2.suppressed-angular-amplitude-def}, filtered angular amplitude
Eq.~\eqref{eq:ch2.filtered-angular-amplitude-def}, full partition
Eq.~\eqref{eq:ch2.angular-component-partition}, lateral counterparts
Eqs.~\eqref{eq:ch2.retained-lateral-spectrum-def}--\eqref{eq:ch2.filtered-lateral-spectrum-def},
Theorem~\ref{thm:ch2.angular-component-partition}, and Propositions
~\ref{prop:ch2.retained-suppressed-mechanism-independence}--\ref{prop:ch2.filtered-observation-rank-bound},
upgrading preview Eq.~\eqref{eq:ch2.retained-suppressed-filter-preview}
\cite{stratton1941,harrington2001,jackson1999}. Subsection~\ref{sec:ch293} develops the near-to-far
operator map Eq.~\eqref{eq:ch2.near-to-far-transform-preview}; Subsection~\ref{sec:ch294} fixes
far-field angular accessibility on \(S^{2}\) \cite{stratton1941}.
\subsection{Near-to-far transformation}
\label{sec:ch293}
% TOC tag: [HOME]

Section~\ref{sec:ch29} previewed the near-to-far map through
Eq.~\eqref{eq:ch2.near-to-far-transform-preview} and tied it to far-zone reduction
Eqs.~\eqref{eq:ch2.far-field-integral-reduction}--\eqref{eq:ch2.far-field-amplitude-integral}
and Theorem~\ref{thm:ch2.far-field-approximation} from Subsection~\ref{sec:ch252}
\cite{stratton1941,jackson1999,harrington2001}. Subsection~\ref{sec:ch293} gives the
Chapter~\ref{ch:02} HOME for the operator that extracts far angular amplitudes
\(a(\widehat{\mathbf{k}})\) from radiative fields after
Subsections~\ref{sec:ch291}--\ref{sec:ch292}. Plane-wave lateral spectra
Eqs.~\eqref{eq:ch2.radiative-lateral-spectrum-def}--\eqref{eq:ch2.plane-wave-radiation-operator},
angular partition Eq.~\eqref{eq:ch2.angular-component-partition}, evanescent null law
Eq.~\eqref{eq:ch2.evanescent-angular-null-law}, and continuous representation
Eq.~\eqref{eq:ch2.exterior-continuous-spectral-integral} are presupposed; asymptotic proofs in
Subsection~\ref{sec:ch252} are not repeated \cite{stratton1941}. Observation-shell regime map
Eq.~\eqref{eq:ch2.observation-shell-regime-map} and regime gate
Eq.~\eqref{eq:ch2.regime-component-gate} qualify when the limit defines transport data
\cite{harrington2001,stratton1941}. Far-field angular accessibility belongs to
Subsection~\ref{sec:ch294} \cite{stratton1941}.

\paragraph{Assumptions.}
Fix \(\omega>0\), admissible \(\Jimp\in\mathcal{J}_\Omega\), outgoing radiative field
\(\mathbf{u}_\omega=\mathcal{R}_\omega[\Jimp]\in\mathcal{F}_{\mathrm{rad}}\), and distance
thresholds from Eq.~\eqref{eq:ch2.distance-regime-definitions} with \(r>r_{\min}(a,\lambda_0)\)
for source support radius \(a\) \cite{stratton1941,jackson1999,harrington2001}. Angular amplitude
\(a(\widehat{\mathbf{k}})\) is measured on the unit sphere with outgoing transverse polarization
\cite{stratton1941}. Regime qualification \(\mathcal{D}_{J,\mathrm{reg}}=1\) from
Eq.~\eqref{eq:ch2.regime-operator-screen} is presupposed before interpreting any near-to-far
output as transport \cite{harrington2001}.

\paragraph{Far-zone extractor (HOME).}
On observation shells \(S_r=\{|\mathbf{r}|=r\}\), define the \emph{far-zone extractor}
\(\mathcal{H}_{\mathrm{ff}}\) by the transverse \(e^{-\jj k_0 r}/r\) coefficient in direction
\(\widehat{\mathbf{k}}\):
\begin{equation}
\label{eq:ch2.far-zone-extractor-def}
\mathcal{H}_{\mathrm{ff}}[\mathbf{u}_\omega](r,\widehat{\mathbf{k}})
=
r\,e^{\jj k_0 r}\,
\widehat{\mathbf{k}}\times\Big(\widehat{\mathbf{k}}\times\E(\mathbf{r})\Big)\Big|_{|\mathbf{r}|=r},
\qquad
\mathbf{r}=r\,\widehat{\mathbf{k}},
\end{equation}
with \(\E\) the electric phasor in \(\mathbf{u}_\omega=(\E,\H)\) and transversality enforced by
the double cross product \cite{stratton1941,jackson1999,harrington2001}. Equation
\eqref{eq:ch2.far-zone-extractor-def} is the Subsection~\ref{sec:ch293} HOME specification of the
extractor previewed in Eq.~\eqref{eq:ch2.near-to-far-transform-preview}: it isolates the
outgoing radiative envelope compatible with Eq.~\eqref{eq:ch2.far-field-integral-reduction}
\cite{stratton1941}. Mathematically, \(\mathcal{H}_{\mathrm{ff}}\) is a linear functional on
shell traces; physically, it is the standard far-field post-processor that removes reactive
\(r^{-2}\) and \(r^{-3}\) contamination when \(k_0 r\gg 1\)
\cite{harrington2001,jackson1999}. The extractor acts on \(\mathcal{R}_\omega[\Jimp]\), not on raw
\(\mathcal{S}_\omega[\Jimp]\) Eq.~\eqref{eq:ch2.observation-shell-regime-map}
\cite{stratton1941}.

\paragraph{Near-to-far transformation (HOME).}
Define the \emph{near-to-far transformation} on admissible sources by
\begin{equation}
\label{eq:ch2.near-to-far-transform-def}
a(\widehat{\mathbf{k}})
=
\lim_{r\to\infty}
\mathcal{H}_{\mathrm{ff}}\big[\mathcal{R}_\omega[\Jimp]\big](r,\widehat{\mathbf{k}}),
\end{equation}
with \(a(\widehat{\mathbf{k}})\) the angular amplitude entering
Eq.~\eqref{eq:ch2.exterior-continuous-spectral-integral} \cite{stratton1941,harrington2001}.
Equation~\eqref{eq:ch2.near-to-far-transform-def} upgrades preview
Eq.~\eqref{eq:ch2.near-to-far-transform-preview} to explicit extractor
Eq.~\eqref{eq:ch2.far-zone-extractor-def} and fixes the radiation map input
\cite{stratton1941,jackson1999}. The limit exists and is independent of the particular
\(r\)-path when Theorem~\ref{thm:ch2.far-field-approximation} applies; error on finite shells
scales as in Eq.~\eqref{eq:ch2.power-flux-regime-convergence}
\cite{stratton1941,harrington2001}. Physically, Eq.~\eqref{eq:ch2.near-to-far-transform-def}
states that far angular content is determined by the radiative branch alone after outgoing
selection \cite{stratton1941}.

\paragraph{Near-to-far operator composition.}
Package Eq.~\eqref{eq:ch2.near-to-far-transform-def} as the composed operator
\begin{equation}
\label{eq:ch2.near-to-far-operator-composition}
\Lambda_{\mathrm{nf}}
=
\mathcal{H}_{\mathrm{ang}}\circ\mathcal{H}_{\mathrm{ff}}^{\infty}\circ\mathcal{R}_\omega,
\qquad
a(\widehat{\mathbf{k}})=\Lambda_{\mathrm{nf}}[\Jimp](\widehat{\mathbf{k}}),
\end{equation}
where \(\mathcal{H}_{\mathrm{ff}}^{\infty}\) denotes the \(r\to\infty\) limit in
Eq.~\eqref{eq:ch2.near-to-far-transform-def} and \(\mathcal{H}_{\mathrm{ang}}\) maps the
transverse far vector in Eq.~\eqref{eq:ch2.far-zone-extractor-def} to scalar polarization
amplitudes \(a_\sigma(\widehat{\mathbf{k}})\) in
Eq.~\eqref{eq:ch2.exterior-continuous-spectral-integral}
\cite{stratton1941,harrington2001}. Equation~\eqref{eq:ch2.near-to-far-operator-composition}
aligns with far readout
\(\Lambda_{\mathrm{rad}}=\mathcal{O}_{\mathrm{far}}\circ\mathcal{R}_\omega\) from
Eq.~\eqref{eq:ch2.reconstruction-observation-map} when \(\mathcal{O}_{\mathrm{far}}\) implements
Eq.~\eqref{eq:ch2.far-zone-extractor-def} on the angular side \cite{stratton1941}. Measurement
\(\mathcal{M}_\omega=\mathcal{O}_{\mathrm{meas}}\circ\Lambda_{\mathrm{rad}}\) therefore factors
near-to-far extraction before finite-rank observation
\cite{harrington2001,stratton1941}. Null additions \(\Jimp^{\mathrm{(NR)}}\in\ker\mathcal{R}_\omega\)
leave \(\Lambda_{\mathrm{nf}}[\Jimp]\) unchanged through
Eq.~\eqref{eq:ch2.angular-spectrum-nr-invariance} \cite{stratton1941}.

\paragraph{Link to plane-wave and retained spectra.}
On the propagating outgoing sheet \(k_t\le k_0\), the near-to-far output agrees with retained
lateral content from Subsection~\ref{sec:ch292}:
\begin{equation}
\label{eq:ch2.near-to-far-lateral-link}
a(\widehat{\mathbf{k}})
\ \longleftrightarrow\
\widetilde{\E}_{\mathrm{ret}}(k_x,k_y)
\ \text{under}\ (k_x,k_y)\mapsto\widehat{\mathbf{k}},
\end{equation}
when \(a_{\mathrm{sup}}=\mathbf{0}\) and Proposition
~\ref{prop:ch2.plane-wave-continuous-consistency} applies \cite{stratton1941,jackson1999}.
Equation~\eqref{eq:ch2.near-to-far-lateral-link} connects Subsection~\ref{sec:ch293} angular
limit to Subsection~\ref{sec:ch291} lateral spectra without re-displaying
Eq.~\eqref{eq:ch2.plane-wave-radiation-operator}: both describe
\(\Pi_{\mathrm{rad}}\mathcal{R}_\omega[\Jimp]\) in different coordinates
\cite{stratton1941,harrington2001}. Filtered content satisfies
\(a_{\mathrm{filt}}=\mathcal{P}_{\mathrm{meas}}[a]\) with \(a\) from
Eq.~\eqref{eq:ch2.near-to-far-transform-def} \cite{stratton1941}.

\paragraph{Evanescent annihilation before near-to-far.}
Evanescent near storage does not enter Eq.~\eqref{eq:ch2.near-to-far-transform-def}:
\begin{equation}
\label{eq:ch2.near-to-far-evanescent-null}
\Pi_{\mathrm{ev}}\mathcal{S}_\omega[\Jimp]
\ \not\mapsto\
a(\widehat{\mathbf{k}})
\quad\text{when}\quad
\mathcal{R}_\omega[\Jimp]\ \text{is unchanged},
\end{equation}
by Eqs.~\eqref{eq:ch2.evanescent-angular-null-law},~\eqref{eq:ch2.evanescent-angular-decoupling},
and~\eqref{eq:ch2.rad-null-on-evanescent} \cite{stratton1941,harrington2001}. Equation
\eqref{eq:ch2.near-to-far-evanescent-null} restates preview
Eq.~\eqref{eq:ch2.near-to-far-transform-preview}: reactive \(\Pi_{\mathrm{ev}}\) content may be
large in the near zone Eq.~\eqref{eq:ch2.near-zone-field-scaling} while the far limit reports
only qualified export \cite{stratton1941,jackson1999}. Regime gate
Eq.~\eqref{eq:ch2.regime-component-gate} forbids promoting \(\Pi_{\mathrm{ev}}\) storage to
\(a(\widehat{\mathbf{k}})\) \cite{harrington2001}.

\begin{figure}[t]
\centering
\ChOneFigBlock{%
\begin{tikzpicture}[ch1fig, node distance=7mm and 10mm]
  \node[ch1box] (j) {$\Jimp$};
  \node[ch1box, right=11mm of j] (r) {$\mathcal{R}_\omega$};
  \node[ch1box, right=11mm of r] (near) {Near shell\\$r\sim r_c$};
  \node[ch1box, right=11mm of near] (hff) {$\mathcal{H}_{\mathrm{ff}}$};
  \node[ch1box, right=11mm of hff] (a) {$a(\widehat{\mathbf{k}})$};
  \node[ch1box, below=11mm of r] (pi) {$\Pi_{\mathrm{rad}}$};
  \node[ch1box, below=11mm of near] (ev) {$\Pi_{\mathrm{ev}}=\mathbf{0}$};
  \draw[ch1flow] (j) -- (r) -- (near) -- (hff) -- (a);
  \draw[ch1flow] (pi) -- (r);
  \draw[ch1flow, dashed] (ev) -- (near);
\end{tikzpicture}%
}
\caption{Subsection~\ref{sec:ch293} near-to-far pipeline:
Eq.~\eqref{eq:ch2.near-to-far-transform-def} applies \(\mathcal{H}_{\mathrm{ff}}\) on
\(\mathcal{R}_\omega[\Jimp]\) after \(r\to\infty\); evanescent content does not feed
Eq.~\eqref{eq:ch2.near-to-far-evanescent-null}.}
\label{fig:ch2.near-to-far-pipeline}
\end{figure}

Figure~\ref{fig:ch2.near-to-far-pipeline} orders radiation selection, finite-shell observation,
far-zone extraction, and angular amplitude. Near reactive content is dashed: it may appear on
inner shells but is annihilated in the limit when export is unchanged
\cite{stratton1941,harrington2001}.

\begin{theorem}[Near-to-far transformation of radiative fields]
\label{thm:ch2.near-to-far-transformation}
Let \(\Jimp\in\mathcal{J}_\Omega\) be admissible with outgoing closure and support in
\(B_a(0)\). Define \(a(\widehat{\mathbf{k}})\) by Eq.~\eqref{eq:ch2.near-to-far-transform-def}.
Then:
\begin{enumerate}
\item[(i)] The limit in Eq.~\eqref{eq:ch2.near-to-far-transform-def} exists for
\(r>r_{\min}(a,\lambda_0)\) and agrees with \(\E_0(\theta,\phi)\) from
Eq.~\eqref{eq:ch2.far-field-amplitude-integral} under Theorem
~\ref{thm:ch2.far-field-approximation}.
\item[(ii)] Eq.~\eqref{eq:ch2.near-to-far-operator-composition} satisfies
\(\Lambda_{\mathrm{nf}}[\Jimp]=a\) on the exporting transverse sheet and is linear on
\(\mathcal{J}_\Omega\).
\item[(iii)] Evanescent content satisfies Eq.~\eqref{eq:ch2.near-to-far-evanescent-null}; raw
\(\mathcal{S}_\omega[\Jimp]\) near traces may differ from \(\Lambda_{\mathrm{nf}}[\Jimp]\).
\item[(iv)] Retained amplitude Eq.~\eqref{eq:ch2.retained-angular-amplitude-def} equals
\(a\) from Eq.~\eqref{eq:ch2.near-to-far-transform-def} when
\(\operatorname{supp}(a)\subseteq\mathcal{K}_{\mathrm{acc}}\).
\end{enumerate}
\end{theorem}
\begin{proof}
Part (i) applies Theorem~\ref{thm:ch2.far-field-approximation} to
Eq.~\eqref{eq:ch2.far-zone-extractor-def} and identifies the limit with
Eq.~\eqref{eq:ch2.far-field-amplitude-integral} \cite{stratton1941,jackson1999}. Part (ii) follows
from linearity of \(\mathcal{R}_\omega\), \(\mathcal{H}_{\mathrm{ff}}\), and
\(\mathcal{H}_{\mathrm{ang}}\) on admissible classes \cite{stratton1941,harrington2001}. Part (iii)
uses Proposition~\ref{prop:ch2.evanescent-never-exports} and
Eq.~\eqref{eq:ch2.evanescent-angular-decoupling} \cite{stratton1941}. Part (iv) is Theorem
~\ref{thm:ch2.angular-component-partition}(ii) with Eq.~\eqref{eq:ch2.near-to-far-lateral-link}
\cite{stratton1941,harrington2001}.
\end{proof}

\begin{proposition}[Consistency with the far-field amplitude integral]
\label{prop:ch2.near-to-far-integral-consistency}
Under Theorem~\ref{thm:ch2.far-field-approximation}, amplitudes from
Eq.~\eqref{eq:ch2.near-to-far-transform-def} and
Eq.~\eqref{eq:ch2.far-field-amplitude-integral} satisfy
\begin{equation}
\label{eq:ch2.near-to-far-integral-agreement}
a(\widehat{\mathbf{k}})
=
-\jj\omega\mu_0\,\frac{1}{4\pi}
\int_{\Omega_s}
\Jimp(\mathbf{r}')\,
e^{-\jj k_0\widehat{\mathbf{k}}\cdot\mathbf{r}'}\,d^{3}r'
\end{equation}
up to transverse polarization components in
Eq.~\eqref{eq:ch2.exterior-continuous-spectral-integral}
\cite{stratton1941,jackson1999,harrington2001}. Equation
\eqref{eq:ch2.near-to-far-integral-agreement} is the Subsection~\ref{sec:ch293} angular restatement
of Eq.~\eqref{eq:ch2.far-field-amplitude-integral}; asymptotics are not re-proved here
\cite{stratton1941}.
\end{proposition}
\begin{proof}
Theorem~\ref{thm:ch2.near-to-far-transformation}(i) and
Eq.~\eqref{eq:ch2.far-field-integral-reduction} identify the limit extractor with the radiation
integral \cite{stratton1941,jackson1999}.
\end{proof}

\begin{proposition}[Null-space invariance of near-to-far maps]
\label{prop:ch2.near-to-far-null-invariance}
If \(\Jimp^{(2)}=\Jimp^{(1)}+\Jimp^{\mathrm{(NR)}}\) with
\(\Jimp^{\mathrm{(NR)}}\in\ker\mathcal{R}_\omega\), then
\(\Lambda_{\mathrm{nf}}[\Jimp^{(2)}]=\Lambda_{\mathrm{nf}}[\Jimp^{(1)}]\)
\cite{stratton1941,harrington2001}. Near-zone fields
\(\mathcal{S}_\omega[\Jimp^{(2)}]\) and \(\mathcal{S}_\omega[\Jimp^{(1)}]\) may differ while
Eq.~\eqref{eq:ch2.near-to-far-transform-def} is unchanged \cite{stratton1941}.
\end{proposition}
\begin{proof}
\(\mathcal{R}_\omega[\Jimp^{(2)}]=\mathcal{R}_\omega[\Jimp^{(1)}]\) implies identical
\(\mathcal{H}_{\mathrm{ff}}\) limits \cite{stratton1941,harrington2001}.
\end{proof}

\begin{proposition}[Finite-shell error bound]
\label{prop:ch2.near-to-far-finite-shell-error}
On shells with \(k_0 r\gg 1\) and \(r\gg r_{\mathrm{F}}\) from
Eq.~\eqref{eq:ch2.fresnel-parameter},
\begin{equation}
\label{eq:ch2.near-to-far-shell-error}
\big|\mathcal{H}_{\mathrm{ff}}[\mathcal{R}_\omega[\Jimp]](r,\widehat{\mathbf{k}})
-a(\widehat{\mathbf{k}})\big|
=
\mathcal{O}\big((k_0 a)^{2}/(k_0 r)\big)
\quad\text{for extended sources},
\end{equation}
with the dipole scaling \(\mathcal{O}((k_0 r)^{-1})\) when \(k_0 a\ll 1\), as in Theorem
~\ref{thm:ch2.far-field-approximation} \cite{stratton1941,jackson1999,harrington2001}. Equation
\eqref{eq:ch2.near-to-far-shell-error} quantifies when finite-radius near-to-far post-processing
approximates Eq.~\eqref{eq:ch2.near-to-far-transform-def} \cite{stratton1941}.
\end{proposition}
\begin{proof}
Apply Theorem~\ref{thm:ch2.far-field-approximation} to
Eq.~\eqref{eq:ch2.far-zone-extractor-def} on \(S_r\) \cite{stratton1941,jackson1999}.
\end{proof}

\paragraph{Mathematical role of Eq.~\eqref{eq:ch2.near-to-far-transform-def}.}
Equation~\eqref{eq:ch2.near-to-far-transform-def} is the angular-spectrum-side closure of the
far-zone template Eq.~\eqref{eq:ch2.far-field-integral-reduction}: it converts radiative fields
to continuous amplitudes \(a(\widehat{\mathbf{k}})\) used in
Eq.~\eqref{eq:ch2.exterior-continuous-spectral-integral} \cite{stratton1941,harrington2001}. The
map is not an optional post-processor on raw near data: input is
\(\mathcal{R}_\omega[\Jimp]\), consistent with observation-shell map
Eq.~\eqref{eq:ch2.observation-shell-regime-map} \cite{stratton1941}. Plane-wave lateral spectra
Eq.~\eqref{eq:ch2.radiative-lateral-spectrum-def} and angular limit
Eq.~\eqref{eq:ch2.near-to-far-lateral-link} are equivalent coordinate descriptions of the same
export \cite{jackson1999,harrington2001}. Reciprocity Theorem
~\ref{thm:ch2.radiation-reception-reciprocity} constrains bilinear couplings built from
\(a(\widehat{\mathbf{k}})\) but does not modify Eq.~\eqref{eq:ch2.near-to-far-evanescent-null}
\cite{stratton1941}.

\paragraph{Mathematical role of Eq.~\eqref{eq:ch2.far-zone-extractor-def}.}
Equation~\eqref{eq:ch2.far-zone-extractor-def} makes explicit how laboratory far-field
algorithms isolate \(\E_0(\theta,\phi)\) before pattern integration or modal projection
\cite{stratton1941,harrington2001}. Applying the extractor to
\(\mathcal{S}_\omega[\Jimp]\) on near shells confuses reactive and radiative scaling
Eq.~\eqref{eq:ch2.near-zone-field-scaling}; applying it to \(\mathcal{R}_\omega[\Jimp]\) aligns
with power-flux convergence Eq.~\eqref{eq:ch2.power-flux-regime-convergence}
\cite{stratton1941,jackson1999}. Filtered observation
Eq.~\eqref{eq:ch2.filtered-angular-amplitude-def} composes after
Eq.~\eqref{eq:ch2.near-to-far-transform-def}, not before radiation selection
\cite{harrington2001,stratton1941}.

\paragraph{Relation to spectral modes and angular partition.}
Modal far pattern Eq.~\eqref{eq:ch2.spectral-far-pattern-expansion} and angular limit
Eq.~\eqref{eq:ch2.near-to-far-transform-def} describe the same \(\operatorname{ran}\mathcal{R}_\omega\)
content under Proposition~\ref{prop:ch2.far-pattern-from-modes}
\cite{stratton1941,jackson1999}. Retained/suppressed partition Theorem
~\ref{thm:ch2.angular-component-partition} constrains \(\operatorname{supp}(a)\) independently of
the near-to-far limit procedure \cite{stratton1941,harrington2001}. Chapter~\ref{ch:03} harmonic
bases change coordinates without altering Eq.~\eqref{eq:ch2.near-to-far-operator-composition}
\cite{stratton1941}. Example~\ref{sec:ch2103} will apply
Eq.~\eqref{eq:ch2.near-to-far-evanescent-null} in a worked setting \cite{stratton1941}.

\paragraph{Boundary of validity.}
Theorem~\ref{thm:ch2.near-to-far-transformation} assumes linear homogeneous or piecewise-homogeneous
media, fixed \(\omega\), outgoing closure, and observation outside \(\Omega_s\)
\cite{stratton1941,harrington2001,jackson1999}. It fails in cutoff guides, cavity standing waves,
and lossless open resonances without outgoing selection \cite{stratton1941}. Surface waves and
strong anisotropy require modified extractors beyond
Eq.~\eqref{eq:ch2.far-zone-extractor-def} \cite{harrington2001}. The theorem does not assert
source uniqueness---only field asymptotics and angular amplitudes on the outgoing branch
\cite{stratton1941}. Cavity discrete spectra Eq.~\eqref{eq:ch2.cavity-discrete-radiation-spectrum}
do not replace Eq.~\eqref{eq:ch2.near-to-far-transform-def} on exterior domains
\cite{jackson1999,harrington2001}.

\paragraph{Worked illustration (dipole shells).}
For \(k_0 a\ll 1\), Eq.~\eqref{eq:ch2.near-zone-field-scaling} dominates on \(k_0 r\lesssim 1\)
while Eq.~\eqref{eq:ch2.near-to-far-transform-def} stabilizes for \(k_0 r\gg 1\)
\cite{stratton1941,jackson1999}. Finite-shell error Eq.~\eqref{eq:ch2.near-to-far-shell-error}
decreases as \(r^{-1}\) once \(\chi(r)\ll 1\) from Section~\ref{sec:ch01.2.4}
\cite{harrington2001,stratton1941}.

\paragraph{Worked illustration (extended aperture).}
When \(k_0 a\gg 1\), Fraunhofer distance \(r_{\mathrm{F}}\) from
Eq.~\eqref{eq:ch2.fresnel-parameter} sets the shell where
Eq.~\eqref{eq:ch2.far-zone-extractor-def} approximates
Eq.~\eqref{eq:ch2.near-to-far-transform-def} \cite{stratton1941,jackson1999}. Intermediate-zone
patterns may differ from \(a(\widehat{\mathbf{k}})\) even when \(\mathcal{D}_{J,\mathrm{reg}}=1\)
\cite{harrington2001}.

\paragraph{Worked illustration (evanescent near storage).}
Large \(\Pi_{\mathrm{ev}}\mathcal{S}_\omega[\Jimp]\) near the source does not alter
Eq.~\eqref{eq:ch2.near-to-far-transform-def} when \(\mathcal{R}_\omega[\Jimp]\) is fixed, by
Eq.~\eqref{eq:ch2.near-to-far-evanescent-null} \cite{stratton1941,harrington2001}. Confusing
near reactive scans with far angular export violates regime gate
Eq.~\eqref{eq:ch2.regime-component-gate} \cite{harrington2001}.

\paragraph{Worked illustration (NR superposition).}
Adding \(\Jimp^{\mathrm{(NR)}}\in\ker\mathcal{R}_\omega\) changes near traces but leaves
\(a(\widehat{\mathbf{k}})\) unchanged by Proposition~\ref{prop:ch2.near-to-far-null-invariance}
\cite{stratton1941,harrington2001}. Near-to-far fitting on raw \(\mathcal{S}_\omega[\Jimp]\)
without \(\mathcal{R}_\omega\) cannot detect the invisible addition \cite{stratton1941}.

\paragraph{Closing summary.}
Subsection~\ref{sec:ch293} establishes the Chapter~\ref{ch:02} HOME for far-zone extractor
Eq.~\eqref{eq:ch2.far-zone-extractor-def}, near-to-far transformation
Eq.~\eqref{eq:ch2.near-to-far-transform-def}, operator composition
Eq.~\eqref{eq:ch2.near-to-far-operator-composition}, lateral link
Eq.~\eqref{eq:ch2.near-to-far-lateral-link}, evanescent null
Eq.~\eqref{eq:ch2.near-to-far-evanescent-null}, Theorem
~\ref{thm:ch2.near-to-far-transformation}, and Propositions
~\ref{prop:ch2.near-to-far-integral-consistency}--\ref{prop:ch2.near-to-far-finite-shell-error},
upgrading preview Eq.~\eqref{eq:ch2.near-to-far-transform-preview}
\cite{stratton1941,jackson1999,harrington2001}. Subsection~\ref{sec:ch294} develops far-field
angular accessibility Eq.~\eqref{eq:ch2.far-angular-accessibility-preview} on \(S^{2}\)
\cite{stratton1941}.
\subsection{Far-field angular accessibility}
\label{sec:ch294}
% TOC tag: [HOME][USE SS1.2.1][USE SS1.2.3]

Section~\ref{sec:ch29} previewed far-field angular accessibility through
Eq.~\eqref{eq:ch2.far-angular-accessibility-preview} and required USE of directional
observables from Section~\ref{sec:ch01.2.1} and transport criteria from
Section~\ref{sec:ch01.2.3} \cite{stratton1941,harrington2001}. Subsection~\ref{sec:ch294}
gives the Chapter~\ref{ch:02} HOME closure for which directions on \(S^{2}\) carry
qualified, observable far-zone angular content after Subsections~\ref{sec:ch291}--\ref{sec:ch293}.
Near-to-far map Eq.~\eqref{eq:ch2.near-to-far-transform-def}, angular partition
Eq.~\eqref{eq:ch2.angular-component-partition}, accessible sector
Eq.~\eqref{eq:ch2.angular-accessibility-sector}, composition law
Eq.~\eqref{eq:ch2.angular-spectrum-composition-law}, and observation rank
Eq.~\eqref{eq:ch2.angular-spectrum-observation-rank} are presupposed; Proposition
~\ref{prop:ch2.angular-accessibility-preview} is upgraded, not re-proved
\cite{stratton1941}. Directional flux maps Eqs.~\eqref{eq:ch1.directional-power-density}--\eqref{eq:ch1.directional-am-flux}
and transport triad Eq.~\eqref{eq:ch1.transport-qualified-decision} enter by reference only
\cite{stratton1941,belinfante1940}. Explicit harmonic filters belong to Chapter~\ref{ch:03};
Section~\ref{sec:ch210} supplies worked examples \cite{stratton1941,harrington2001}.

\paragraph{Assumptions.}
Fix \(\omega>0\), admissible \(\Jimp\in\mathcal{J}_\Omega\), outgoing radiative field
\(\mathbf{u}_\omega=\mathcal{R}_\omega[\Jimp]\), angular amplitude \(a(\widehat{\mathbf{k}})\) from
Eq.~\eqref{eq:ch2.near-to-far-transform-def}, and regime qualification
\(\mathcal{D}_{J,\mathrm{reg}}=1\) from Eq.~\eqref{eq:ch2.regime-operator-screen}
\cite{stratton1941,harrington2001}. Accessible sector \(\mathcal{K}_{\mathrm{acc}}\) and
suppressed sheet \(\mathcal{K}_{\mathrm{sup}}\) are fixed by
Eqs.~\eqref{eq:ch2.angular-accessibility-sector}--\eqref{eq:ch2.suppressed-angular-sheet-def}
\cite{stratton1941}. Finite-rank measurement \(\mathcal{M}_\omega\) from
Eq.~\eqref{eq:ch2.reconstruction-observation-map} is presupposed \cite{harrington2001}.

\paragraph{Far-field accessibility criterion (HOME).}
Define \emph{far-field angular accessibility} by
\begin{equation}
\label{eq:ch2.far-angular-accessibility-def}
\widehat{\mathbf{k}}\in\mathcal{K}_{\mathrm{acc}}^{\mathrm{ff}}
\ \Longleftrightarrow\
\Big[\,
a(\widehat{\mathbf{k}})\neq 0
\ \text{after}\ \Lambda_{\mathrm{nf}}
\ \text{and}\ \mathcal{D}_{J,\mathrm{reg}}=1
\,\Big]
\ \text{and}\ 
\widehat{\mathbf{k}}\in\mathcal{K}_{\mathrm{acc}},
\end{equation}
with \(\Lambda_{\mathrm{nf}}\) from Eq.~\eqref{eq:ch2.near-to-far-operator-composition} and
\(\mathcal{K}_{\mathrm{acc}}\) from Eq.~\eqref{eq:ch2.angular-accessibility-sector}
\cite{stratton1941,harrington2001,jackson1999}. Equation~\eqref{eq:ch2.far-angular-accessibility-def}
upgrades preview Eq.~\eqref{eq:ch2.far-angular-accessibility-preview} by inserting explicit
near-to-far qualification and geometric sector membership \cite{stratton1941}. Mathematically,
\(\mathcal{K}_{\mathrm{acc}}^{\mathrm{ff}}\subseteq\mathcal{K}_{\mathrm{acc}}\) is the
far-field reportable subset after radiation and regime gates; physically, it marks directions
where angular amplitude survives export and lies in the geometrically accessible hemisphere
\cite{harrington2001,stratton1941}. Directions in \(\mathcal{K}_{\mathrm{sup}}\) are excluded
before observation composition \cite{stratton1941}.

\paragraph{Angular accessibility operator.}
Package Eq.~\eqref{eq:ch2.far-angular-accessibility-def} as the \emph{accessibility projector}
\begin{equation}
\label{eq:ch2.angular-accessibility-operator}
\mathcal{A}_{\mathrm{ff}}
=
\chi_{\mathcal{K}_{\mathrm{acc}}}\circ\mathcal{D}_{J,\mathrm{reg}}\circ\Lambda_{\mathrm{nf}},
\qquad
a_{\mathrm{acc}}(\widehat{\mathbf{k}})
=
\big[\mathcal{A}_{\mathrm{ff}}[\Jimp]\big](\widehat{\mathbf{k}}),
\end{equation}
with \(a_{\mathrm{acc}}\) the accessible angular amplitude before finite measurement
\cite{stratton1941,harrington2001}. Equation~\eqref{eq:ch2.angular-accessibility-operator}
aligns retained amplitude Eq.~\eqref{eq:ch2.retained-angular-amplitude-def} with near-to-far
output when \(\operatorname{supp}(a)\subseteq\mathcal{K}_{\mathrm{acc}}\) Theorem
~\ref{thm:ch2.near-to-far-transformation}(iv) \cite{stratton1941}. Filtered observation
Eq.~\eqref{eq:ch2.filtered-angular-amplitude-def} applies after
Eq.~\eqref{eq:ch2.angular-accessibility-operator}:
\(a_{\mathrm{filt}}=\mathcal{P}_{\mathrm{meas}}[a_{\mathrm{acc}}]\)
\cite{harrington2001,stratton1941}. Null additions leave \(a_{\mathrm{acc}}\) unchanged through
Eq.~\eqref{eq:ch2.angular-spectrum-nr-invariance} \cite{stratton1941}.

\paragraph{Plane-wave accessibility filter (HOME).}
On lateral coordinates, define the \emph{plane-wave accessibility filter} by
\begin{equation}
\label{eq:ch2.plane-wave-accessibility-filter}
\widetilde{\E}_{\mathrm{acc}}(k_x,k_y)
=
\chi_{\mathcal{K}_{\mathrm{acc}}^{\mathrm{ff}}}(\widehat{\mathbf{k}})\,
\widetilde{\E}_{\mathrm{rad}}(k_x,k_y),
\qquad
(k_x,k_y)\mapsto\widehat{\mathbf{k}}\ \text{on}\ k_t\le k_0,
\end{equation}
with \(\widetilde{\E}_{\mathrm{rad}}\) from Eq.~\eqref{eq:ch2.radiative-lateral-spectrum-def}
\cite{stratton1941,harrington2001,jackson1999}. Equation~\eqref{eq:ch2.plane-wave-accessibility-filter}
is the Subsection~\ref{sec:ch294} HOME plane-wave screen promised in
Eq.~\eqref{eq:ch2.far-angular-accessibility-preview}: it removes geometrically forbidden
directions on the propagating outgoing sheet without re-tabulating
Eq.~\eqref{eq:ch2.plane-wave-radiation-operator} \cite{stratton1941}. Evanescent sheets satisfy
\(\widetilde{\E}_{\mathrm{acc}}=\mathbf{0}\) for \(k_t>k_0\) through
Eq.~\eqref{eq:ch2.evanescent-angular-null-law} \cite{stratton1941,harrington2001}. Physically,
Eq.~\eqref{eq:ch2.plane-wave-accessibility-filter} is the angular-spectrum-side statement of which
plane-wave directions participate in far-zone reports \cite{jackson1999,harrington2001}.

\paragraph{Transport qualification gate (USE).}
Far-field accessibility composes with Chapter~\ref{ch:01} transport screening by reference.
A direction \(\widehat{\mathbf{k}}\in\mathcal{K}_{\mathrm{acc}}^{\mathrm{ff}}\) supports a
qualified transport report on solid-angle sector \(\Delta\Omega(\widehat{\mathbf{k}})\) only when
Section~\ref{sec:ch01.2.3} propagative support Eq.~\eqref{eq:ch1.propagative-support} and
transport decision Eq.~\eqref{eq:ch1.transport-qualified-decision} hold for the corresponding
directional channels Eqs.~\eqref{eq:ch1.directional-power-density}--\eqref{eq:ch1.directional-am-flux}
\cite{stratton1941,belinfante1940}. Restated in operator form,
\begin{equation}
\label{eq:ch2.far-angular-transport-gate}
\mathcal{T}_{\mathrm{acc}}(\widehat{\mathbf{k}})=1
\ \Longrightarrow\
\int_{\Delta\Omega(\widehat{\mathbf{k}})}\Pi(\theta,\phi;r)\,d\Omega>0
\ \text{and}\ 
\mathcal{Q}_{J}=1,
\end{equation}
with \(\Pi\) from Eq.~\eqref{eq:ch1.directional-power-density},
\(\mathcal{Q}_{J}\) from Eq.~\eqref{eq:ch1.torque-transport-qualified-criterion}, and
\(r\) in the far zone where Eq.~\eqref{eq:ch2.near-to-far-transform-def} applies
\cite{stratton1941,harrington2001}. Equation~\eqref{eq:ch2.far-angular-transport-gate} imports
Section~\ref{sec:ch01.2.3} criteria without re-deriving the transport triad: Chapter~\ref{ch:02}
records when operator-level accessibility implies Chapter~\ref{ch:01} transport admissibility on
the same sector \cite{stratton1941}. Regime gate Eq.~\eqref{eq:ch2.regime-component-gate} remains
mandatory: \(\mathcal{T}_{\mathrm{acc}}=1\) does not promote reactive near content to export
\cite{harrington2001}.

\paragraph{Angular measurement screen (HOME).}
After accessibility projection, finite instruments apply
\begin{equation}
\label{eq:ch2.angular-measurement-screen}
a_{\mathrm{obs}}(\widehat{\mathbf{k}})
=
\big[\mathcal{P}_{\mathrm{meas}}\,a_{\mathrm{acc}}\big](\widehat{\mathbf{k}}),
\qquad
\operatorname{supp}(a_{\mathrm{obs}})
\subseteq
\{\widehat{\mathbf{k}}:\ \text{resolvable by}\ \mathcal{M}_\omega\}
\cap
\mathcal{K}_{\mathrm{acc}}^{\mathrm{ff}},
\end{equation}
with \(\mathcal{P}_{\mathrm{meas}}\) from Eq.~\eqref{eq:ch2.filtered-angular-amplitude-def}
\cite{stratton1941,harrington2001}. Equation~\eqref{eq:ch2.angular-measurement-screen} is the
Subsection~\ref{sec:ch294} HOME measurement closure: observation rank
Eq.~\eqref{eq:ch2.angular-spectrum-observation-rank} may strictly shrink
\(\operatorname{supp}(a_{\mathrm{obs}})\) inside \(\mathcal{K}_{\mathrm{acc}}^{\mathrm{ff}}\)
\cite{stratton1941}. Rank cap Eq.~\eqref{eq:ch2.filtered-rank-bound} limits resolvable directions
independently of source nullity \cite{harrington2001}. Physically, \(a_{\mathrm{obs}}\) is the
laboratory-accessible angular spectrum after geometry, regime, near-to-far extraction, and
aperture \cite{stratton1941,jackson1999}.

\paragraph{Accessible angular budget (HOME).}
Collecting Subsection~\ref{sec:ch294} constraints,
\begin{equation}
\label{eq:ch2.accessible-angular-budget}
\operatorname{supp}(a_{\mathrm{obs}})
\subseteq
\mathcal{K}_{\mathrm{acc}}^{\mathrm{ff}}
\cap
\{\widehat{\mathbf{k}}:\ \mathcal{T}_{\mathrm{acc}}(\widehat{\mathbf{k}})=1\}
\cap
\{\widehat{\mathbf{k}}:\ \text{resolvable by}\ \mathcal{M}_\omega\},
\end{equation}
composing geometric, transport, and observation factors with angular budget
Eq.~\eqref{eq:ch2.angular-spectrum-composition-law} \cite{stratton1941,harrington2001}. Equation
\eqref{eq:ch2.accessible-angular-budget} is the Section~\ref{sec:ch29} closure of
Eq.~\eqref{eq:ch2.section29-angular-spectrum-chain}: each factor removes a distinct failure mode
--- radiation nullity, geometric forbidding, regime violation, transport disqualification, and
finite aperture --- without conflating them \cite{stratton1941}. Reciprocity Theorem
~\ref{thm:ch2.radiation-reception-reciprocity} constrains coupling among observed channels on
\(\operatorname{supp}(a_{\mathrm{obs}})\) but does not enlarge
Eq.~\eqref{eq:ch2.accessible-angular-budget} \cite{stratton1941,harrington2001}.

\begin{figure}[t]
\centering
\ChOneFigBlock{%
\begin{tikzpicture}[ch1fig, node distance=7mm and 10mm]
  \node[ch1box] (nf) {$\Lambda_{\mathrm{nf}}$};
  \node[ch1box, right=11mm of nf] (acc) {$\mathcal{A}_{\mathrm{ff}}$};
  \node[ch1box, right=11mm of acc] (t) {$\mathcal{T}_{\mathrm{acc}}$};
  \node[ch1box, right=11mm of t] (m) {$\mathcal{P}_{\mathrm{meas}}$};
  \node[ch1box, right=11mm of m] (obs) {$a_{\mathrm{obs}}$};
  \node[ch1box, below=11mm of acc] (k) {$\mathcal{K}_{\mathrm{acc}}$};
  \node[ch1box, below=11mm of t] (c1) {USE §1.2.3};
  \draw[ch1flow] (nf) -- (acc) -- (t) -- (m) -- (obs);
  \draw[ch1flow] (k) -- (acc);
  \draw[ch1flow, dashed] (c1) -- (t);
\end{tikzpicture}%
}
\caption{Subsection~\ref{sec:ch294} accessibility pipeline:
Eq.~\eqref{eq:ch2.angular-accessibility-operator} applies near-to-far extraction, geometric
sector \(\mathcal{K}_{\mathrm{acc}}\), transport gate Eq.~\eqref{eq:ch2.far-angular-transport-gate}
(USE Section~\ref{sec:ch01.2.3}), and measurement screen
Eq.~\eqref{eq:ch2.angular-measurement-screen}.}
\label{fig:ch2.far-angular-accessibility-pipeline}
\end{figure}

Figure~\ref{fig:ch2.far-angular-accessibility-pipeline} orders near-to-far extraction, geometric
accessibility, transport qualification imported from Chapter~\ref{ch:01}, and finite observation.
Suppressed directions never enter \(a_{\mathrm{obs}}\); evanescent content is removed before
\(\Lambda_{\mathrm{nf}}\) Eq.~\eqref{eq:ch2.near-to-far-evanescent-null}
\cite{stratton1941,harrington2001}.

\begin{theorem}[Far-field angular accessibility]
\label{thm:ch2.far-angular-accessibility}
Let \(\Jimp\in\mathcal{J}_\Omega\) be admissible with outgoing closure and define
\(a_{\mathrm{acc}}\), \(a_{\mathrm{obs}}\), and \(\mathcal{K}_{\mathrm{acc}}^{\mathrm{ff}}\) by
Eqs.~\eqref{eq:ch2.far-angular-accessibility-def}--\eqref{eq:ch2.angular-measurement-screen}. Then:
\begin{enumerate}
\item[(i)] Eq.~\eqref{eq:ch2.accessible-angular-budget} holds with equality of supports up to
measure-zero sets on \(S^{2}\).
\item[(ii)] \(\mathcal{K}_{\mathrm{acc}}^{\mathrm{ff}}\subseteq\mathcal{K}_{\mathrm{acc}}\) and
\(a_{\mathrm{acc}}(\widehat{\mathbf{k}})=\mathbf{0}\) on \(\mathcal{K}_{\mathrm{sup}}\).
\item[(iii)] Transport gate Eq.~\eqref{eq:ch2.far-angular-transport-gate} is necessary for
qualified angular reports importing Section~\ref{sec:ch01.2.1}--\ref{sec:ch01.2.3} by reference.
\item[(iv)] Plane-wave filter Eq.~\eqref{eq:ch2.plane-wave-accessibility-filter} agrees with
\(a_{\mathrm{acc}}\) under Eq.~\eqref{eq:ch2.near-to-far-lateral-link} on \(k_t\le k_0\).
\end{enumerate}
\end{theorem}
\begin{proof}
Part (i) composes Eqs.~\eqref{eq:ch2.angular-spectrum-composition-law},
~\eqref{eq:ch2.angular-accessibility-operator}, and~\eqref{eq:ch2.angular-measurement-screen}
\cite{stratton1941,harrington2001}. Part (ii) follows from
Eq.~\eqref{eq:ch2.suppressed-angular-sheet-def} and Theorem
~\ref{thm:ch2.angular-component-partition} \cite{stratton1941}. Part (iii) cites
Eq.~\eqref{eq:ch1.transport-qualified-decision} and
Eq.~\eqref{eq:ch1.propagative-support} without re-proof \cite{stratton1941,belinfante1940}. Part (iv)
uses Proposition~\ref{prop:ch2.plane-wave-continuous-consistency} and Theorem
~\ref{thm:ch2.near-to-far-transformation} \cite{stratton1941,jackson1999}.
\end{proof}

\begin{proposition}[Reciprocity constraint on accessible directions]
\label{prop:ch2.accessibility-reciprocity-constraint}
On \(\operatorname{supp}(a_{\mathrm{obs}})\), port reciprocity
Eq.~\eqref{eq:ch2.transmit-receive-reciprocity} and modal coupling
Eq.~\eqref{eq:ch2.radiative-coupling-bilinear} hold after projection, but reciprocity does not
 enlarge \(\mathcal{K}_{\mathrm{acc}}^{\mathrm{ff}}\) or restore directions in
\(\mathcal{K}_{\mathrm{sup}}\) \cite{stratton1941,harrington2001}. Equation
\eqref{eq:ch2.accessible-angular-budget} is unchanged under reciprocal media when geometry and
\(\mathcal{M}_\omega\) are fixed \cite{stratton1941}.
\end{proposition}
\begin{proof}
Theorem~\ref{thm:ch2.radiation-reception-reciprocity} constrains bilinear forms on observed
channels; geometric sector Eq.~\eqref{eq:ch2.angular-accessibility-sector} is independent of
reciprocity \cite{stratton1941,harrington2001}.
\end{proof}

\begin{proposition}[Transport consistency with Chapter 1 observables]
\label{prop:ch2.accessibility-transport-consistency}
If \(\widehat{\mathbf{k}}\in\operatorname{supp}(a_{\mathrm{obs}})\) and
\(\mathcal{T}_{\mathrm{acc}}(\widehat{\mathbf{k}})=1\), then directional power and angular-flux
maps Eqs.~\eqref{eq:ch1.directional-power-density}--\eqref{eq:ch1.directional-am-flux} support
nontrivial export on \(\Delta\Omega(\widehat{\mathbf{k}})\) in the far zone, consistent with
Section~\ref{sec:ch01.2.3} transport triad Eq.~\eqref{eq:ch1.transport-qualified-decision}
\cite{stratton1941,belinfante1940}. Conversely, failure of
Eq.~\eqref{eq:ch1.propagative-support} on a sector excludes that sector from qualified
\(\operatorname{supp}(a_{\mathrm{obs}})\) even when near scans show strong local fields
\cite{stratton1941,harrington2001}.
\end{proposition}
\begin{proof}
Eq.~\eqref{eq:ch2.far-angular-transport-gate} restates Section~\ref{sec:ch01.2.3} criteria on
sectors selected by \(a_{\mathrm{acc}}\) \cite{stratton1941}. Regime gate
Eq.~\eqref{eq:ch2.regime-component-gate} blocks reactive-only sectors \cite{harrington2001}.
\end{proof}

\begin{proposition}[Observation-limited accessibility]
\label{prop:ch2.accessibility-observation-rank}
The observed accessible spectrum satisfies
\begin{equation}
\label{eq:ch2.observed-accessibility-rank}
\dim\operatorname{span}\{a_{\mathrm{obs}}(\widehat{\mathbf{k}})\}
\le
\mathrm{rank}(\mathcal{M}_\omega)
\le
\min\{N,\,\dim\operatorname{ran}\mathcal{R}_\omega\},
\end{equation}
with strict inequality possible on subsets of \(\mathcal{K}_{\mathrm{acc}}^{\mathrm{ff}}\) where
\(a_{\mathrm{acc}}\neq\mathbf{0}\) \cite{stratton1941,harrington2001}. Equation
\eqref{eq:ch2.observed-accessibility-rank} specializes
Eq.~\eqref{eq:ch2.filtered-rank-bound} to the accessibility pipeline
\cite{stratton1941}.
\end{proposition}
\begin{proof}
Apply Proposition~\ref{prop:ch2.filtered-observation-rank-bound} to
Eq.~\eqref{eq:ch2.angular-measurement-screen} \cite{stratton1941,harrington2001}.
\end{proof}

\paragraph{Mathematical role of Eq.~\eqref{eq:ch2.far-angular-accessibility-def}.}
Equation~\eqref{eq:ch2.far-angular-accessibility-def} closes Section~\ref{sec:ch29} by stating
which directions on \(S^{2}\) remain after radiation selection, geometric constraints, and
regime qualification \cite{stratton1941,harrington2001}. It upgrades preview
Eq.~\eqref{eq:ch2.far-angular-accessibility-preview} with explicit dependence on
\(\Lambda_{\mathrm{nf}}\) and \(\mathcal{K}_{\mathrm{acc}}\) \cite{stratton1941}. Operator form
Eq.~\eqref{eq:ch2.angular-accessibility-operator} packages the same logic for source-to-amplitude
maps \cite{harrington2001}. Chapter~\ref{ch:03} harmonic bases relabel
\(\mathcal{K}_{\mathrm{acc}}^{\mathrm{ff}}\) without altering the budget
Eq.~\eqref{eq:ch2.accessible-angular-budget} \cite{stratton1941}.

\paragraph{Mathematical role of Eq.~\eqref{eq:ch2.plane-wave-accessibility-filter}.}
Equation~\eqref{eq:ch2.plane-wave-accessibility-filter} is the lateral-coordinate screen
completing the plane-wave chain begun in Subsection~\ref{sec:ch291} \cite{stratton1941,jackson1999}.
It removes forbidden directions before near-to-far comparison and measurement
\cite{harrington2001}. Agreement with \(a_{\mathrm{acc}}\) on \(k_t\le k_0\) follows from
Eq.~\eqref{eq:ch2.near-to-far-lateral-link} \cite{stratton1941}. Evanescent sheets remain null
through Eq.~\eqref{eq:ch2.evanescent-angular-null-law} independently of raw lateral scans
\cite{stratton1941,harrington2001}.

\paragraph{Relation to the Section~\ref{sec:ch29} chain.}
Subsections~\ref{sec:ch291}--\ref{sec:ch294} complete
Eq.~\eqref{eq:ch2.section29-angular-spectrum-chain}:
plane-wave radiation operators, retained/suppressed/filtered partition, near-to-far extraction,
and far-field accessibility \cite{stratton1941,harrington2001}. Figure
~\ref{fig:ch2.angular-spectrum-action-roadmap} orders the same sequence. Modal coordinates
Eq.~\eqref{eq:ch2.spectral-far-pattern-expansion} and plane-wave coordinates
Eq.~\eqref{eq:ch2.exterior-continuous-spectral-integral} describe identical export on
\(\operatorname{ran}\mathcal{R}_\omega\); Section~\ref{sec:ch29} fixes operator-level filters on
the plane-wave branch \cite{stratton1941,jackson1999}. Inverse obstruction
Proposition~\ref{prop:ch2.angular-spectrum-inverse-obstruction} applies to \(a_{\mathrm{obs}}\),
not to pre-accessibility transforms \cite{harrington2001,stratton1941}.

\paragraph{Boundary of validity.}
Theorem~\ref{thm:ch2.far-angular-accessibility} assumes linear media, fixed \(\omega\), outgoing
closure, and stationary geometry \(\mathcal{C}_{\mathrm{geo}}\) \cite{stratton1941,harrington2001}.
Time-varying \(\mathcal{M}_\omega\) requires separate treatment; static rank bounds remain upper
limits \cite{stratton1941}. Loss modifies amplitudes but not the inclusion chain in
Eq.~\eqref{eq:ch2.accessible-angular-budget} \cite{stratton1941}. Cavity and guide spectra remain
outside exterior accessibility; discrete books do not substitute for
Eq.~\eqref{eq:ch2.far-angular-accessibility-def} on unbounded domains
\cite{jackson1999,harrington2001}. Section~\ref{sec:ch01.2.3} transport criteria are imported, not
extended to near reactive zones without regime qualification \cite{stratton1941,harrington2001}.

\paragraph{Worked illustration (ground-plane sector).}
For a horizontal source above a PEC ground, \(\mathcal{K}_{\mathrm{acc}}^{\mathrm{ff}}\) is a
hemisphere while \(\mathcal{K}_{\mathrm{sup}}\) contains image-forbidden directions
Eq.~\eqref{eq:ch2.boundary-image-constraint} \cite{stratton1941,jackson1999}. Near scans may show
energy in suppressed directions without contributing to \(a_{\mathrm{obs}}\) \cite{harrington2001}.

\paragraph{Worked illustration (transport gate).}
A direction with \(a_{\mathrm{acc}}\neq\mathbf{0}\) but failing
Eq.~\eqref{eq:ch1.propagative-support} on \(\Delta\Omega(\widehat{\mathbf{k}})\) does not satisfy
\(\mathcal{T}_{\mathrm{acc}}=1\) and is excluded from qualified reports by
Eq.~\eqref{eq:ch2.accessible-angular-budget}, consistent with Section~\ref{sec:ch01.2.3}
\cite{stratton1941,belinfante1940}.

\paragraph{Worked illustration (finite array).}
When \(\mathcal{O}_{\mathrm{meas}}\) has rank \(N\), \(a_{\mathrm{obs}}\) may vanish on portions
of \(\mathcal{K}_{\mathrm{acc}}^{\mathrm{ff}}\) where \(a_{\mathrm{acc}}\neq\mathbf{0}\) because of
Eq.~\eqref{eq:ch2.observed-accessibility-rank} \cite{stratton1941,harrington2001}. Pattern nulls from
aperture therefore need not indicate \(\Jimp\in\ker\mathcal{R}_\omega\) \cite{stratton1941}.

\paragraph{Worked illustration (NR invisibility).}
Adding \(\Jimp^{\mathrm{(NR)}}\in\ker\mathcal{R}_\omega\) leaves \(a_{\mathrm{acc}}\) and
\(a_{\mathrm{obs}}\) unchanged while altering near fields \cite{stratton1941,harrington2001}.
Accessibility reports therefore cannot detect NR additions from far data alone
\cite{stratton1941}.

\paragraph{Closing summary.}
Subsection~\ref{sec:ch294} establishes the Chapter~\ref{ch:02} HOME for far-field accessibility
Eq.~\eqref{eq:ch2.far-angular-accessibility-def}, accessibility operator
Eq.~\eqref{eq:ch2.angular-accessibility-operator}, plane-wave filter
Eq.~\eqref{eq:ch2.plane-wave-accessibility-filter}, transport gate
Eq.~\eqref{eq:ch2.far-angular-transport-gate} (USE Sections~\ref{sec:ch01.2.1}--\ref{sec:ch01.2.3}),
measurement screen Eq.~\eqref{eq:ch2.angular-measurement-screen}, budget
Eq.~\eqref{eq:ch2.accessible-angular-budget}, Theorem~\ref{thm:ch2.far-angular-accessibility}, and
Propositions~\ref{prop:ch2.accessibility-reciprocity-constraint}--\ref{prop:ch2.accessibility-observation-rank},
upgrading preview Eq.~\eqref{eq:ch2.far-angular-accessibility-preview}
\cite{stratton1941,harrington2001,jackson1999}. Section~\ref{sec:ch29} is complete; Section~\ref{sec:ch210}
develops worked examples and Section~\ref{sec:ch02.11} records the Chapter~\ref{ch:02} forward map
\cite{stratton1941}.
\section{Examples}
\label{sec:ch210}

Section~\ref{sec:ch210} is a \emph{degrees-of-freedom diagnostic suite}: four worked problems
test distinct failure modes of the radiation pipeline developed in Sections~\ref{sec:ch20}--\ref{sec:ch29}.
Subsections~\ref{sec:ch2101}--\ref{sec:ch2104} isolate compact truncation, source nullity,
evanescent near-to-far contamination, and structured-source topology--transport separation,
respectively \cite{stratton1941,harrington2001,devaney2004,mandel1995}. Every numerical ratio
quoted below is \textbf{illustrative}: values follow author-calibrated discretizations of the stated
geometries and are not certified measurement data; they are included only to show order-of-magnitude
separation among mechanisms \cite{hansen1988}.

\paragraph{Master DOF partition.}
Radiation reports can fail for distinct reasons that must not be merged in software or laboratory
interpretation \cite{marengo2000,bleszynski1996}:
\begin{equation}
\label{eq:ch2.dof-partition-master}
\underbrace{\ker\mathcal{R}_\omega}_{\text{source nullity (NR)}}
\ \oplus\
\underbrace{\ker\mathcal{O}_{\mathrm{meas}}\big|_{\operatorname{ran}\mathcal{R}_\omega}}_{\text{aperture rank}}
\ \oplus\
\underbrace{\Pi_{\mathrm{ev}}\mathcal{S}_\omega[\mathcal{J}_\Omega]}_{\text{evanescent storage}}
\ \oplus\
\underbrace{\text{tail of compact/ modal truncation}}_{\text{spectral compression}}
\ \oplus\
\underbrace{\mathcal{K}_{\mathrm{sup}}}_{\text{geometric suppression}},
\end{equation}
with accessibility \(\mathcal{K}_{\mathrm{acc}}^{\mathrm{ff}}\) from
Eq.~\eqref{eq:ch2.accessible-angular-budget} acting after Eq.~\eqref{eq:ch2.dof-partition-master}
\cite{stratton1941,harrington2001}. Subsections~\ref{sec:ch2101}--\ref{sec:ch2104} each activate a
different term in Eq.~\eqref{eq:ch2.dof-partition-master}; cross-references replace repeated
local split displays \cite{stratton1941}.

\subsection{Compactness of radiation operators}
\label{sec:ch2101}
% TOC tag: [FIG+PROB][USE SS2.1.3]

Subsection~\ref{sec:ch2101} illustrates compactness and rank compression for radiation-related
maps, importing Subsection~\ref{sec:ch213} by reference only
\cite{stratton1941,harrington2001,kato1995}. Compact operator definition
Eq.~\eqref{eq:ch2.compact-operator-def}, Green split Eq.~\eqref{eq:ch2.green-compact-split},
singular-value decay Eq.~\eqref{eq:ch2.singular-value-decay}, compression bound
Eq.~\eqref{eq:ch2.compression-error-bound}, and rank composition
Eq.~\eqref{eq:ch2.rank-composition-law} are presupposed; proofs are not repeated
\cite{stratton1941}. Radiation map Eq.~\eqref{eq:ch2.radiation-operator-composition},
measurement chain Eq.~\eqref{eq:ch2.operational-measurement-chain}, and null-space structure
Section~\ref{sec:ch27} supply context \cite{harrington2001,stratton1941}.

\paragraph{Problem.}
A compact impressed-current ball \(\Omega_s=B_a(0)\) carries an electric dipole
\(\Jimp(\mathbf{r})=\mathbf{J}_0\,\delta(\mathbf{r})\) in the distributional limit with
\(\int_{\Omega_s}\Jimp\,d^{3}r=\mathbf{p}\) and \(k_0 a=1.5\) at fixed \(\omega\)
\cite{stratton1941,jackson1999}. A near-field scanner records \(N=36\) tangential electric
samples \(\mathbf{y}\in\mathbb{C}^{36}\) on a closed surface \(S_b\) with radius \(b=2a\)
through measurement map \(\mathcal{M}_\omega\) from
Eq.~\eqref{eq:ch2.reconstruction-observation-map} \cite{harrington2001}. The task is to:
\begin{enumerate}
\item[(a)] identify the compact and non-compact parts of the probe-to-far-field map implied by
Subsection~\ref{sec:ch213};
\item[(b)] bound truncation error when only \(M=20\) retained modes are used in a spherical-harmonic
surrogate of \(\mathcal{R}_\omega[\Jimp]\);
\item[(c)] state whether unresolved modal content is attributable to compact compression, to
\(\ker\mathcal{R}_\omega\), or to finite rank of \(\mathcal{O}_{\mathrm{meas}}\).
\end{enumerate}

\begin{figure}[t]
\centering
\ChOneFigBlock{%
\begin{tikzpicture}[ch1fig, node distance=7mm and 10mm]
  \node[ch1box] (j) {$\Jimp$\\$k_0a=1.5$};
  \node[ch1box, right=12mm of j] (s) {$\mathcal{S}_\omega$};
  \node[ch1box, right=12mm of s] (r) {$\mathcal{R}_\omega$};
  \node[ch1box, right=12mm of r] (m) {$\mathcal{M}_\omega$};
  \node[ch1box, right=12mm of m] (y) {$\mathbf{y}\in\mathbb{C}^{36}$};
  \node[ch1box, below=11mm of s] (split) {$G_{\mathrm{sing}}+K_{\mathrm{cpt}}$};
  \node[ch1box, below=11mm of m] (rank) {$\mathrm{rank}\le 36$};
  \draw[ch1flow] (j) -- (s) -- (r) -- (m) -- (y);
  \draw[ch1flow] (split) -- (s);
  \draw[ch1flow] (rank) -- (m);
\end{tikzpicture}%
}
\caption{Worked problem~\ref{sec:ch2101}: compact source, Green split on \(\mathcal{S}_\omega\),
radiation map, and rank-\(36\) near-field probe bank. Compact part \(K_{\mathrm{cpt}}\) controls
modal truncation error; \(\mathcal{R}_\omega\) is not compact on \(\mathcal{J}_\Omega\).}
\label{fig:ch2.wp-compact-radiation-probes}
\end{figure}

Figure~\ref{fig:ch2.wp-compact-radiation-probes} orders the operator chain. The schematic
separates compact smoothing on the source-to-field map from non-compact radiation selection and
finite-rank readout \cite{stratton1941,harrington2001}.

\paragraph{Step 1 — Operator split (USE §2.1.3).}
On bounded domains with compact \(\Omega_s\), Subsection~\ref{sec:ch213} records
\begin{equation}
\label{eq:ch2.ex-compact-green-split}
\mathcal{S}_\omega
=
G_{\mathrm{sing}}+K_{\mathrm{cpt}},
\end{equation}
with \(K_{\mathrm{cpt}}\) compact on \(\mathcal{J}_\Omega\) Eq.~\eqref{eq:ch2.green-compact-split}
\cite{stratton1941}. Equation~\eqref{eq:ch2.ex-compact-green-split} restates the split without
re-proving kernel construction from Section~\ref{sec:ch23} \cite{stratton1941,harrington2001}.
Mathematically, \(G_{\mathrm{sing}}\) captures non-compact singular behavior as probe and source
points coincide; \(K_{\mathrm{cpt}}\) maps bounded source sets to relatively compact field sets in
the \(L^{2}\) topology on \(S_b\) \cite{stratton1941}. Physically, smooth interpolation between
widely separated samples is a compact smoothing operation, whereas pointwise singular self-terms
are not \cite{harrington2001,jackson1999}.

\paragraph{Step 2 — Radiation map is not compact.}
Radiation selection defines
\(\mathcal{R}_\omega=\Pi_{\mathrm{rad}}\circ\mathcal{S}_\omega\) from
Eq.~\eqref{eq:ch2.radiation-operator-composition} \cite{stratton1941}. The map
\(\mathcal{R}_\omega:\mathcal{J}_\Omega\to\mathcal{F}_{\mathrm{rad}}\) is \emph{not} compact:
\(\ker\mathcal{R}_\omega\neq\{\mathbf{0}\}\) is infinite-dimensional
Theorem~\ref{thm:ch2.radiation-nullspace-structure}, so bounded source balls do not map to
relatively compact sets in \(\mathcal{F}_{\mathrm{rad}}\) \cite{stratton1941,harrington2001}.
Compactness enters through \(K_{\mathrm{cpt}}\) in Eq.~\eqref{eq:ch2.ex-compact-green-split} and
through finite observation, not through \(\mathcal{R}_\omega\) itself \cite{stratton1941}. Confusing
non-radiating additions with modal truncation error violates the distinction established in
Subsection~\ref{sec:ch213} \cite{harrington2001}.

\paragraph{Step 3 — Finite-rank measurement.}
The probe bank implements \(\mathcal{O}_{\mathrm{meas}}:\mathcal{F}_{\mathrm{rad}}\to\mathbb{C}^{36}\)
with \(\mathrm{rank}(\mathcal{O}_{\mathrm{meas}})\le 36\) from
Eq.~\eqref{eq:ch2.finite-rank-observation} \cite{stratton1941,harrington2001}. Composed map
\(\mathcal{M}_\omega=\mathcal{O}_{\mathrm{meas}}\circ\mathcal{O}_{\mathrm{far}}\circ\mathcal{R}_\omega\)
therefore satisfies
\begin{equation}
\label{eq:ch2.ex-compact-rank-cap}
\mathrm{rank}(\mathcal{M}_\omega)
\le
\min\{36,\,\dim\mathcal{R}(\mathcal{R}_\omega)\}
\end{equation}
by Eq.~\eqref{eq:ch2.rank-composition-law} \cite{stratton1941}. Mathematically, at most \(36\)
independent complex numbers are available to infer far-zone content; physically, the scanner cannot
resolve more than \(36\) radiation degrees of freedom regardless of Maxwell residual in the volume
\cite{harrington2001,stratton1941}. Null-space directions in
\(\ker\mathcal{O}_{\mathrm{meas}}\big|_{\operatorname{ran}\mathcal{R}_\omega}\) are
infinite-dimensional Eq.~\eqref{eq:ch2.observation-null-space} \cite{stratton1941}.

\paragraph{Step 4 — Singular values of the compact part.}
Let \(\mathcal{K}\) denote the probe-to-modal map obtained by composing sample interpolation with
compact \(K_{\mathrm{cpt}}\) on the dipole-supported subspace. Subsection~\ref{sec:ch213} gives
\begin{equation}
\label{eq:ch2.ex-compact-svd-decay}
\|\mathcal{K}-\mathcal{K}_M\|_{\mathrm{op}}
=
\inf_{\mathrm{rank}\,\mathcal{K}_M\le M}\|\mathcal{K}-\mathcal{K}_M\|_{\mathrm{op}}
\sim
\sigma_{M+1}(\mathcal{K}),
\end{equation}
restating Eq.~\eqref{eq:ch2.singular-value-decay} on the worked geometry \cite{stratton1941}.
For \(k_0 a=1.5\) and \(b=2a\), numerical evaluation of the discretized \(\mathcal{K}\) yields
\(\sigma_{20}/\sigma_1\approx 0.12\) and \(\sigma_{37}/\sigma_1\approx 0.04\): the twenty-first
mode is already below \(12\%\) of the dominant singular value \cite{harrington2001,stratton1941}.
Physically, retaining \(M=20\) modes captures most radiative energy resolvable on \(S_b\); modes
\(21\)–\(36\) lie in the tail controlled by compact smoothing, not in
\(\ker\mathcal{R}_\omega\) \cite{stratton1941,jackson1999}.

\paragraph{Step 5 — Truncation error bound.}
Subsection~\ref{sec:ch213} bounds compression error by
\begin{equation}
\label{eq:ch2.ex-compact-truncation-bound}
\|\mathcal{R}_\omega[\Jimp]-\mathcal{R}_\omega^{(M)}[\Jimp]\|_{\mathcal{U}}
\le
C\,\sigma_{M+1}(K_{\mathrm{cpt}})\,\|\Jimp\|_{\mathcal{J}},
\end{equation}
importing Eq.~\eqref{eq:ch2.compression-error-bound} with \(\mathcal{R}_\omega^{(M)}\) the
\(M\)-mode surrogate Eq.~\eqref{eq:ch2.spectral-compression-operator}
\cite{stratton1941,harrington2001}. With \(\sigma_{21}\approx 0.12\,\sigma_1\) and dipole
normalization \(\|\Jimp\|_{\mathcal{J}}\) fixed by \(|\mathbf{p}|\), the bound states that no
rank-\(20\) operator can reduce far-field pattern error below order \(0.12\,|\mathbf{E}_0|\) on
the outgoing branch \cite{stratton1941,jackson1999}. Mathematically, the tail is controlled by
compact \(K_{\mathrm{cpt}}\); physically, increasing \(M\) beyond the singular-value knee yields
diminishing returns unless \(N\) or \(b/a\) increases \cite{harrington2001}.

\paragraph{Step 6 — Compact truncation term in Eq.~\eqref{eq:ch2.dof-partition-master}.}
For the dipole problem, the active mechanism is the compact tail of
\(\sigma_n(K_{\mathrm{cpt}})\) in Eq.~\eqref{eq:ch2.dof-partition-master}, not
\(\ker\mathcal{R}_\omega\) \cite{stratton1941,harrington2001,kato1995}. Adding
\(\Jimp^{\mathrm{(NR)}}\in\ker\mathcal{R}_\omega\) leaves \(\mathbf{y}\) unchanged;
modes beyond rank \(36\) lie in \(\ker\mathcal{O}_{\mathrm{meas}}\); modes \(21\)–\(36\) sit in
the compact tail with error \(\sim\sigma_{21}\) (illustrative ratio \(\sigma_{21}/\sigma_1\approx 0.12\))
\cite{stratton1941}. Pattern-fitting software must identify which term in
Eq.~\eqref{eq:ch2.dof-partition-master} removed unresolved content
\cite{harrington2001,stratton1941,jackson1999}.

\paragraph{Physical meaning (summary).}
Subsection~\ref{sec:ch2101} shows that compactness is a property of Green smoothing on bounded
sources, not of \(\mathcal{R}_\omega\) itself \cite{stratton1941,harrington2001,kato1995}.
Regime gate Eq.~\eqref{eq:ch2.regime-operator-screen} remains mandatory before interpreting
truncated modal reports as transport \cite{harrington2001}.

\subsection{Non-radiating source configurations}
\label{sec:ch2102}
% TOC tag: [FIG+PROB][USE SS2.7.2]

Subsection~\ref{sec:ch2102} activates the \(\ker\mathcal{R}_\omega\) term of
Eq.~\eqref{eq:ch2.dof-partition-master}, importing Subsection~\ref{sec:ch272} by reference only
\cite{stratton1941,harrington2001,jackson1999,reed1950,bleszynski1996}. Non-radiating current definition Eq.~\eqref{eq:ch2.nr-current-def}, canonical families
Eqs.~\eqref{eq:ch2.nr-loop-cancellation}--\eqref{eq:ch2.nr-modal-eigenrealization},
Theorem~\ref{thm:ch2.canonical-nr-current-families}, superposition law
Eq.~\eqref{eq:ch2.nr-superposition-law}, near-field law
Eq.~\eqref{eq:ch2.nr-near-field-energy-law}, and modal invisibility
Eq.~\eqref{eq:ch2.nr-modal-invisibility} are presupposed; proofs are not repeated
\cite{stratton1941}. Null-space operator structure Theorem~\ref{thm:ch2.radiation-nullspace-structure},
angular invariance Eq.~\eqref{eq:ch2.angular-spectrum-nr-invariance}, and inverse obstruction
Proposition~\ref{prop:ch2.angular-spectrum-inverse-obstruction} supply context
\cite{harrington2001,stratton1941}.

\paragraph{Problem.}
A compact electric dipole \(\Jimp^{(\mathrm{rad})}\) with moment \(|\mathbf{p}|=p_0\) and
\(k_0 a=1.2\) radiates in free space with outgoing closure \cite{stratton1941,jackson1999}. An
engineer adds a small magnetic loop \(\Jimp^{(\mathrm{loop})}\) of radius \(a_{\mathrm{loop}}=0.3a\)
with equal opposing filaments so net dipole moment vanishes and loop cancellation
Eq.~\eqref{eq:ch2.nr-loop-cancellation} applies \cite{stratton1941,harrington2001}. Define
\(\Jimp^{(\mathrm{tot})}=\Jimp^{(\mathrm{rad})}+\Jimp^{(\mathrm{loop})}\) with charge continuity
preserved on \(\Omega_s\) Eq.~\eqref{eq:ch2.nr-current-admissibility-law}
\cite{stratton1941}. Near-field probes on \(S_b\) at \(b=2a\) record \(\mathbf{y}_{\mathrm{near}}\);
a far-zone array estimates \(\E_0(\theta,\phi)\) from Eq.~\eqref{eq:ch2.far-field-amplitude-integral}.
The task is to:
\begin{enumerate}
\item[(a)] verify \(\Jimp^{(\mathrm{loop})}\in\ker\mathcal{R}_\omega\) using Subsection~\ref{sec:ch272}
without re-deriving canonical families;
\item[(b)] quantify near-field change versus far-field invariance under superposition;
\item[(c)] explain why source inversion Eq.~\eqref{eq:ch2.inverse-source-equation} cannot detect
\(\Jimp^{(\mathrm{loop})}\) from far data alone.
\end{enumerate}

\begin{figure}[t]
\centering
\ChOneFigBlock{%
\begin{tikzpicture}[ch1fig, node distance=7mm and 10mm]
  \node[ch1box] (rad) {$\Jimp^{(\mathrm{rad})}$\\dipole};
  \node[ch1box, right=11mm of rad] (loop) {$\Jimp^{(\mathrm{loop})}$\\NR loop};
  \node[ch1box, right=11mm of loop] (tot) {$\Jimp^{(\mathrm{tot})}$};
  \node[ch1box, right=11mm of tot] (r) {$\mathcal{R}_\omega$};
  \node[ch1box, right=11mm of r] (ff) {$\E_0$\\unchanged};
  \node[ch1box, below=11mm of tot] (s) {$\mathcal{S}_\omega$};
  \node[ch1box, below=11mm of s] (near) {Near\\$\mathbf{y}$ changes};
  \draw[ch1flow] (rad) -- (tot);
  \draw[ch1flow] (loop) -- (tot);
  \draw[ch1flow] (tot) -- (r) -- (ff);
  \draw[ch1flow] (tot) -- (s) -- (near);
  \draw[ch1flow, dashed] (loop) -- node[below,font=\scriptsize] {$\ker\mathcal{R}_\omega$} (r);
\end{tikzpicture}%
}
\caption{Worked problem~\ref{sec:ch2102}: superposition of radiating dipole and loop-cancellation
non-radiating current. Far export \(\E_0\) is invariant; near probes detect \(\mathcal{S}_\omega[\Jimp^{(\mathrm{tot})}]\neq\mathcal{S}_\omega[\Jimp^{(\mathrm{rad})}]\).}
\label{fig:ch2.wp-nr-superposition-invisibility}
\end{figure}

Figure~\ref{fig:ch2.wp-nr-superposition-invisibility} separates radiation map invariance from
source-field map sensitivity \cite{stratton1941,harrington2001}.

\paragraph{Step 1 — Loop membership in \(\ker\mathcal{R}_\omega\) (USE §2.7.2).}
Subsection~\ref{sec:ch272} records that equal opposing arms on \(\mathcal{C}_+\), \(\mathcal{C}_-\)
with vanishing net circulation imply
\begin{equation}
\label{eq:ch2.ex-nr-loop-test}
\int_{\mathcal{C}_+}\!\mathbf{J}_{\mathrm{loop}}\cdot d\boldsymbol{\ell}
+
\int_{\mathcal{C}_-}\!\mathbf{J}_{\mathrm{loop}}\cdot d\boldsymbol{\ell}
=0
\ \Longrightarrow\
\mathcal{R}_\omega[\Jimp^{(\mathrm{loop})}]=\mathbf{0},
\end{equation}
restating Eq.~\eqref{eq:ch2.nr-loop-cancellation} on the worked geometry \cite{stratton1941,jackson1999}.
Mathematically, vanishing leading transverse moment in
Eq.~\eqref{eq:ch2.far-field-amplitude-integral} annihilates \(\E_0\); physically, the loop is
magnetically dominated near \(\Omega_s\) but electrically invisible at infinity on the outgoing
branch \cite{harrington2001,stratton1941}. Theorem~\ref{thm:ch2.canonical-nr-current-families}
certifies \(\Jimp^{(\mathrm{loop})}\) as admissible non-radiating configuration without reopening
the proof \cite{stratton1941}.

\paragraph{Step 2 — Superposition invariance on export.}
For \(\Jimp^{(\mathrm{tot})}=\Jimp^{(\mathrm{rad})}+\Jimp^{(\mathrm{loop})}\), Subsection~\ref{sec:ch272}
gives
\begin{equation}
\label{eq:ch2.ex-nr-superposition-test}
\mathcal{R}_\omega[\Jimp^{(\mathrm{tot})}]
=
\mathcal{R}_\omega[\Jimp^{(\mathrm{rad})}]+\mathcal{R}_\omega[\Jimp^{(\mathrm{loop})}]
=
\mathcal{R}_\omega[\Jimp^{(\mathrm{rad})}],
\end{equation}
by linearity Eq.~\eqref{eq:ch2.radiation-linearity} and
\(\mathcal{R}_\omega[\Jimp^{(\mathrm{loop})}]=\mathbf{0}\) \cite{stratton1941,harrington2001}. Far
pattern amplitudes therefore satisfy
\begin{equation}
\label{eq:ch2.ex-nr-far-invariance}
\E_0(\theta,\phi;\Jimp^{(\mathrm{tot})})
=
\E_0(\theta,\phi;\Jimp^{(\mathrm{rad})}),
\qquad
a(\widehat{\mathbf{k}};\Jimp^{(\mathrm{tot})})
=
a(\widehat{\mathbf{k}};\Jimp^{(\mathrm{rad})}),
\end{equation}
importing Eq.~\eqref{eq:ch2.nr-superposition-law} and angular invariance
Eq.~\eqref{eq:ch2.angular-spectrum-nr-invariance} \cite{stratton1941,jackson1999}. Mathematically,
\(\ker\mathcal{R}_\omega\) is invisible to all far-zone readouts including
Eq.~\eqref{eq:ch2.near-to-far-transform-def}; physically, antenna pattern measurements cannot
distinguish \(\Jimp^{(\mathrm{tot})}\) from \(\Jimp^{(\mathrm{rad})}\) on qualified shells
\cite{harrington2001,stratton1941}. Modal coefficients obey
Eq.~\eqref{eq:ch2.nr-modal-invisibility} for \(\Jimp^{(\mathrm{loop})}\) alone
\cite{stratton1941}.

\paragraph{Step 3 — Near-field detectability versus far invisibility.}
Source-to-field map \(\mathcal{S}_\omega\) is generally \emph{not} invariant:
\begin{equation}
\label{eq:ch2.ex-nr-near-detectability}
\mathcal{S}_\omega[\Jimp^{(\mathrm{tot})}]
=
\mathcal{S}_\omega[\Jimp^{(\mathrm{rad})}]
+
\mathcal{S}_\omega[\Jimp^{(\mathrm{loop})}],
\qquad
\|\mathcal{S}_\omega[\Jimp^{(\mathrm{loop})}]\|_{\mathcal{E}}
\ \text{may be}\ \gg\ 0
\end{equation}
on near shells, restating Proposition~\ref{prop:ch2.nr-near-field-nonzero} and
Eq.~\eqref{eq:ch2.nr-near-field-energy-law} \cite{stratton1941,harrington2001}. For
\(k_0 b=2.4\), numerical evaluation gives \(|\Delta\mathbf{y}|/|\mathbf{y}_{\mathrm{rad}}|\approx 0.35\)
on tangential probes: the loop is visible near the source while far \(\E_0\) is unchanged
\cite{stratton1941,jackson1999}. Regime gate Eq.~\eqref{eq:ch2.regime-component-gate} forbids
interpreting the near increment as angular transport without \(\mathcal{D}_{J,\mathrm{reg}}=1\)
\cite{harrington2001}. This step is distinct from evanescent storage
Eq.~\eqref{eq:ch2.near-to-far-evanescent-null}: the loop contributes to
\(\mathcal{S}_\omega[\Jimp^{(\mathrm{tot})}]\) but not to \(\mathcal{R}_\omega[\Jimp^{(\mathrm{tot})}]\)
\cite{stratton1941}.

\paragraph{Step 4 — Power and flux meters.}
Radiated power on qualified far shells satisfies
\(\mathcal{P}_{\mathrm{rad}}(\Jimp^{(\mathrm{tot})})=\mathcal{P}_{\mathrm{rad}}(\Jimp^{(\mathrm{rad})})\)
from Eq.~\eqref{eq:ch2.power-flux-regime-convergence} and
Eq.~\eqref{eq:ch2.ex-nr-far-invariance} \cite{stratton1941,harrington2001}. Reactive budget
Eq.~\eqref{eq:ch2.component-power-budget} may increase locally because
\(\mathcal{P}_{\mathrm{nr}}\) from \(\mathcal{S}_\omega[\Jimp^{(\mathrm{loop})}]\) adds near
\(\Omega_s\) while export remains unchanged \cite{stratton1941}. Physically, integrating Poynting
flux through a near sphere confuses storage with radiation unless regime index
Eq.~\eqref{eq:ch1.regime-index} from Section~\ref{sec:ch01.2.4} qualifies the shell
\cite{harrington2001,stratton1941}.

\paragraph{Step 5 — Inverse source ill-posedness.}
Far-zone data fix \(\mathcal{R}_\omega[\Jimp]\) but not \(\Jimp\). Proposition
~\ref{prop:ch2.angular-spectrum-inverse-obstruction} and Eq.~\eqref{eq:ch2.angular-inverse-null-quotient}
imply
\begin{equation}
\label{eq:ch2.ex-nr-inverse-quotient}
\Jimp^{(1)}-\Jimp^{(2)}\in\ker\mathcal{R}_\omega
\quad\text{whenever}\quad
\E_0(\Jimp^{(1)})=\E_0(\Jimp^{(2)}),
\end{equation}
so \(\Jimp^{(\mathrm{tot})}\) and \(\Jimp^{(\mathrm{rad})}\) belong to the same equivalence class
\cite{stratton1941,harrington2001}. Tikhonov reconstruction
Eq.~\eqref{eq:ch2.tikhonov-source-reconstruction} may return either representative depending on
regularization without contradiction \cite{stratton1941}. Mathematically, \(\ker\mathcal{R}_\omega\)
has infinite dimension on exterior domains
Eq.~\eqref{eq:ch2.nullspace-dimensional-classification}; physically, adding loops, surface
tangential nulls Eq.~\eqref{eq:ch2.nr-surface-tangential-null}, or modal \(\lambda_n=0\)
realizations Eq.~\eqref{eq:ch2.nr-modal-eigenrealization} leaves patterns unchanged
\cite{harrington2001,jackson1999}.

\paragraph{Step 6 — Source nullity term in Eq.~\eqref{eq:ch2.dof-partition-master}.}
Far invariance Eq.~\eqref{eq:ch2.ex-nr-far-invariance} arises because
\(\Jimp^{(\mathrm{loop})}\in\ker\mathcal{R}_\omega\)---the active term is source nullity in
Eq.~\eqref{eq:ch2.dof-partition-master}, not finite probe rank or compact truncation
\cite{stratton1941,reed1950,bleszynski1996}. Geometric suppression
Eq.~\eqref{eq:ch2.suppressed-angular-sheet-def} remains a separate term: the dipole radiates on
\(\mathcal{K}_{\mathrm{acc}}\) while the loop adds invisible near content
\cite{harrington2001,stratton1941}.

\paragraph{Physical meaning (summary).}
Subsection~\ref{sec:ch2102} demonstrates that non-radiating configurations are physically realizable
and operationally consequential: near diagnostics change while far observables do not
\cite{stratton1941,jackson1999,harrington2001,devaney2004}. Pattern-matching based solely on
\(\E_0\) or \(a(\widehat{\mathbf{k}})\) cannot detect \(\Jimp^{(\mathrm{loop})}\); near measurements
and source-domain constraints Eq.~\eqref{eq:ch2.nr-current-admissibility-law} are required
\cite{stratton1941}.

\subsection{Evanescent suppression in the far field}
\label{sec:ch2103}
% TOC tag: [FIG+PROB][USE SS2.9.3]

Subsection~\ref{sec:ch2103} activates the evanescent-storage term of
Eq.~\eqref{eq:ch2.dof-partition-master}, importing Subsection~\ref{sec:ch293} by reference only
\cite{stratton1941,harrington2001,jackson1999,mandel1995}. Near-to-far definition
Eq.~\eqref{eq:ch2.near-to-far-transform-def}, far-zone extractor
Eq.~\eqref{eq:ch2.far-zone-extractor-def}, operator composition
Eq.~\eqref{eq:ch2.near-to-far-operator-composition}, evanescent annihilation
Eq.~\eqref{eq:ch2.near-to-far-evanescent-null}, angular null law
Eq.~\eqref{eq:ch2.evanescent-angular-null-law}, Theorem~\ref{thm:ch2.near-to-far-transformation},
and Propositions~\ref{prop:ch2.near-to-far-null-invariance}--\ref{prop:ch2.near-to-far-finite-shell-error}
are presupposed; asymptotic proofs in Subsection~\ref{sec:ch252} and evanescent HOME content in
Subsection~\ref{sec:ch273} are not repeated \cite{stratton1941}. Evanescent overlap
Eq.~\eqref{eq:ch2.evanescent-source-overlap}, non-export law
Eq.~\eqref{eq:ch2.evanescent-nonexport-law}, regime gate
Eq.~\eqref{eq:ch2.regime-component-gate}, and near-zone scaling
Eq.~\eqref{eq:ch2.near-zone-field-scaling} supply context \cite{harrington2001,stratton1941}.

\paragraph{Problem.}
A compact impressed-current ball \(\Omega_s=B_a(0)\) carries an electric dipole
\(\Jimp^{(\mathrm{d})}(\mathbf{r})=\mathbf{J}_0\,\delta(\mathbf{r})\) in the distributional limit
with \(\int_{\Omega_s}\Jimp^{(\mathrm{d})}\,d^{3}r=\mathbf{p}\) and \(k_0 a=2.0\) at fixed
\(\omega\) \cite{stratton1941,jackson1999}. A second admissible current
\(\Jimp^{(\mathrm{ev})}\) is chosen so that
\(\mathcal{R}_\omega[\Jimp^{(\mathrm{ev})}]=\mathbf{0}\) while
\(\|\Pi_{\mathrm{ev}}\mathcal{S}_\omega[\Jimp^{(\mathrm{ev})}]\|_{\mathcal{E}}\) is comparable to
the radiative energy of \(\Jimp^{(\mathrm{d})}\) \cite{stratton1941,harrington2001}. This is
\emph{field-side evanescent excitation with vanishing export}, not membership of
\(\Jimp^{(\mathrm{ev})}\) in the canonical non-radiating \emph{source} catalog of
Subsection~\ref{sec:ch272}: the same far invariance can arise from \(\Jimp^{\mathrm{(NR)}}\in\ker\mathcal{R}_\omega\)
with a different near signature \cite{reed1950,bleszynski1996,stratton1941}. Define the
superposition \(\Jimp=\Jimp^{(\mathrm{d})}+\Jimp^{(\mathrm{ev})}\). A near-field probe records
tangential \(\E\) on shell \(S_b\) with \(b=1.25a\), while a far-field post-processor applies
extractor Eq.~\eqref{eq:ch2.far-zone-extractor-def} on shells \(S_r\) with increasing \(r\)
\cite{harrington2001,stratton1941}. The task is to:
\begin{enumerate}
\item[(a)] quantify near-shell contrast between \(\mathcal{S}_\omega[\Jimp]\) and
\(\mathcal{R}_\omega[\Jimp]\) when evanescent storage is large;
\item[(b)] verify that \(a(\widehat{\mathbf{k}})\) from Eq.~\eqref{eq:ch2.near-to-far-transform-def}
is unchanged by \(\Jimp^{(\mathrm{ev})}\) and that lateral spectra on \(k_t>k_0\) vanish under
Eq.~\eqref{eq:ch2.evanescent-angular-null-law};
\item[(c)] diagnose failure modes when the near-to-far limit is applied to raw
\(\mathcal{S}_\omega[\Jimp]\) instead of \(\mathcal{R}_\omega[\Jimp]\).
\end{enumerate}

\begin{figure}[t]
\centering
\ChOneFigBlock{%
\begin{tikzpicture}[ch1fig, node distance=7mm and 10mm]
  \node[ch1box] (j) {$\Jimp^{(\mathrm{d})}+\Jimp^{(\mathrm{ev})}$};
  \node[ch1box, right=11mm of j] (s) {$\mathcal{S}_\omega$};
  \node[ch1box, right=11mm of s] (r) {$\mathcal{R}_\omega$};
  \node[ch1box, right=11mm of r] (hff) {$\mathcal{H}_{\mathrm{ff}}$};
  \node[ch1box, right=11mm of hff] (a) {$a(\widehat{\mathbf{k}})$};
  \node[ch1box, below=11mm of s] (near) {Near shell\\$b=1.25a$};
  \node[ch1box, below=11mm of r] (ev) {$\Pi_{\mathrm{ev}}\neq\mathbf{0}$};
  \node[ch1box, below=11mm of hff] (lim) {$r\to\infty$};
  \draw[ch1flow] (j) -- (s) -- (r) -- (hff) -- (a);
  \draw[ch1flow, dashed] (s) -- (near);
  \draw[ch1flow, dashed] (ev) -- (s);
  \draw[ch1flow] (lim) -- (hff);
\end{tikzpicture}%
}
\caption{Worked problem~\ref{sec:ch2103}: evanescent current adds large near reactive content
through \(\mathcal{S}_\omega\) but does not alter \(\mathcal{R}_\omega[\Jimp]\) or far amplitude
\(a(\widehat{\mathbf{k}})\) from Eq.~\eqref{eq:ch2.near-to-far-transform-def}. Dashed path:
near scans; solid path: qualified near-to-far export (USE Subsection~\ref{sec:ch293}).}
\label{fig:ch2.wp-evanescent-far-suppression}
\end{figure}

Figure~\ref{fig:ch2.wp-evanescent-far-suppression} orders source superposition, full field map,
radiation selection, far-zone extraction, and angular amplitude. Evanescent overlap appears only on
the dashed near branch; the far limit reports radiative export alone
\cite{stratton1941,harrington2001}.

\paragraph{Step 1 — Spectral split (USE §2.9.3).}
On admissible \(\Jimp\), Subsection~\ref{sec:ch293} fixes the near-to-far input as
\(\mathcal{R}_\omega[\Jimp]\), not \(\mathcal{S}_\omega[\Jimp]\)
Eq.~\eqref{eq:ch2.near-to-far-transform-def} \cite{stratton1941}. Decompose the impressed
current into radiative and evanescent-overlap parts through
\begin{equation}
\label{eq:ch2.ex-ev-spectral-decomposition}
\Jimp
=
\Jimp^{(\mathrm{d})}+\Jimp^{(\mathrm{ev})},
\qquad
\mathcal{R}_\omega[\Jimp^{(\mathrm{ev})}]=\mathbf{0},
\qquad
\|\Pi_{\mathrm{ev}}\mathcal{S}_\omega[\Jimp^{(\mathrm{ev})}]\|_{\mathcal{E}}
=
\beta\,\|\mathcal{R}_\omega[\Jimp^{(\mathrm{d})}]\|_{\mathcal{E}},
\end{equation}
with \(\beta=1.8\) in the worked normalization \cite{stratton1941,harrington2001}. Equation
\eqref{eq:ch2.ex-ev-spectral-decomposition} restates evanescent source overlap
Eq.~\eqref{eq:ch2.evanescent-source-overlap} on the dipole-plus-evanescent geometry without
re-proving Proposition~\ref{prop:ch2.evanescent-never-exports}
\cite{stratton1941}. Mathematically, \(\Jimp^{(\mathrm{ev})}\) lies in a class that excites
\(\mathcal{F}_{\mathrm{ev}}\) Eq.~\eqref{eq:ch2.evanescent-subspace-def} while annihilating under
\(\Pi_{\mathrm{rad}}\) Eq.~\eqref{eq:ch2.evanescent-nonexport-law}; physically, it models reactive
near storage---for example, a superposition of forbidden-angle plane-wave trials with imaginary
longitudinal wavenumber Eq.~\eqref{eq:ch2.evanescent-decay-rate}---that does not add outgoing
export \cite{harrington2001,jackson1999}.

\paragraph{Step 2 — Near-shell contrast.}
On \(S_b\) with \(b=1.25a\), near-zone scaling Eq.~\eqref{eq:ch2.near-zone-field-scaling}
dominates: reactive \(r^{-3}\) and inductive \(r^{-2}\) terms can exceed the radiative \(r^{-1}\)
envelope \cite{stratton1941,jackson1999}. Tangential energy norms on \(S_b\) yield
\begin{equation}
\label{eq:ch2.ex-ev-near-amplitude-ratio}
\frac{\|\E_{\mathrm{near}}^{(\mathrm{full})}\|_{L^{2}(S_b)}}
     {\|\E_{\mathrm{near}}^{(\mathrm{rad})}\|_{L^{2}(S_b)}}
=
\frac{\|\Pi_{\mathrm{near}}\mathcal{S}_\omega[\Jimp^{(\mathrm{d})}+\Jimp^{(\mathrm{ev})}]\|}
     {\|\Pi_{\mathrm{near}}\mathcal{R}_\omega[\Jimp^{(\mathrm{d})}]\|}
\approx
\sqrt{1+\beta^{2}}
\approx
2.16,
\end{equation}
where \(\Pi_{\mathrm{near}}\) denotes tangential shell restriction and \(\E_{\mathrm{near}}^{(\mathrm{rad})}\)
is computed from \(\mathcal{R}_\omega[\Jimp^{(\mathrm{d})}]\) alone \cite{stratton1941,harrington2001}.
Equation~\eqref{eq:ch2.ex-ev-near-amplitude-ratio} quantifies task (a): a scanner that records raw
\(\mathcal{S}_\omega[\Jimp]\) sees roughly twice the tangential magnitude of the radiative
component at \(b=1.25a\) \cite{stratton1941}. Mathematically, the ratio follows from orthogonal
energy split Eq.~\eqref{eq:ch2.subspace-energy-inequality} on evanescent and radiative subspaces;
physically, it warns that near-field ``hot spots'' need not imply additional far-zone power or
modified angular patterns \cite{harrington2001,jackson1999}. Regime map
Eq.~\eqref{eq:ch2.observation-shell-regime-map} classifies \(b=1.25a\) as reactive-dominated for
\(k_0 a=2\) \cite{stratton1941}.

\paragraph{Step 3 — Far amplitude invariance (USE §2.9.3).}
Near-to-far transformation Eq.~\eqref{eq:ch2.near-to-far-transform-def} gives
\begin{equation}
\label{eq:ch2.ex-ev-far-invariance}
a(\widehat{\mathbf{k}})
=
\lim_{r\to\infty}
\mathcal{H}_{\mathrm{ff}}\big[\mathcal{R}_\omega[\Jimp^{(\mathrm{d})}+\Jimp^{(\mathrm{ev})}]\big]
=
\lim_{r\to\infty}
\mathcal{H}_{\mathrm{ff}}\big[\mathcal{R}_\omega[\Jimp^{(\mathrm{d})}]\big],
\end{equation}
because \(\mathcal{R}_\omega[\Jimp^{(\mathrm{ev})}]=\mathbf{0}\) and \(\Lambda_{\mathrm{nf}}\) is
linear Eq.~\eqref{eq:ch2.near-to-far-operator-composition} \cite{stratton1941,harrington2001}.
Theorem~\ref{thm:ch2.near-to-far-transformation}(iii) restates the same conclusion through
Eq.~\eqref{eq:ch2.near-to-far-evanescent-null}: evanescent storage in
\(\Pi_{\mathrm{ev}}\mathcal{S}_\omega[\Jimp]\) does not map to \(a(\widehat{\mathbf{k}})\) when
\(\mathcal{R}_\omega[\Jimp]\) is unchanged \cite{stratton1941}. Equation
\eqref{eq:ch2.ex-ev-far-invariance} answers task (b): dipole angular pattern, directional gain, and
plane-wave lateral content on \(k_t\le k_0\) are identical with and without
\(\Jimp^{(\mathrm{ev})}\) \cite{jackson1999,harrington2001}. Proposition
~\ref{prop:ch2.near-to-far-null-invariance} further guarantees that any
\(\Jimp^{(\mathrm{ev})}\in\ker\mathcal{R}_\omega\) leaves \(a\) unchanged, distinct from the
evanescent-overlap case where \(\Jimp^{(\mathrm{ev})}\) is not necessarily non-radiating as a
source but still annihilates under \(\mathcal{R}_\omega\) \cite{stratton1941,harrington2001}.

\paragraph{Step 4 — Lateral spectrum null on the evanescent sheet.}
Plane-wave radiation operator Eq.~\eqref{eq:ch2.plane-wave-radiation-operator} applied to the
superposition gives, on \(k_t>k_0\),
\begin{equation}
\label{eq:ch2.ex-ev-lateral-null}
\widetilde{\E}_{\mathrm{rad}}(k_x,k_y)
=
\big[\mathcal{T}_\omega\,\widetilde{\Jimp}\big](k_x,k_y)
=
\mathbf{0},
\qquad
k_t>k_0,
\end{equation}
importing Eq.~\eqref{eq:ch2.evanescent-angular-null-law} without re-displaying
Eq.~\eqref{eq:ch2.rad-null-on-evanescent} \cite{stratton1941,harrington2001}. Raw lateral transform
\(\mathcal{F}_{\mathrm{lat}}[\mathcal{S}_\omega[\Jimp]]\) can show large coefficients on the
evanescent sheet; radiative lateral spectrum
Eq.~\eqref{eq:ch2.radiative-lateral-spectrum-def} does not \cite{stratton1941}. Mathematically,
\(\mathcal{T}_\omega\) factors \(\Pi_{\mathrm{rad}}\circ\mathcal{R}_\omega\) before
\(\mathcal{F}_{\mathrm{lat}}\); physically, a near-field Fourier scan that omits radiation selection
misreports ``angular content'' on forbidden \(k_t\) \cite{harrington2001,jackson1999}. Link
Eq.~\eqref{eq:ch2.near-to-far-lateral-link} identifies \(a(\widehat{\mathbf{k}})\) with
\(\widetilde{\E}_{\mathrm{ret}}\) only on \(k_t\le k_0\) when retained content vanishes
\cite{stratton1941}.

\paragraph{Step 5 — Wrong versus correct near-to-far pipelines.}
Task (c) compares two post-processing chains on the same near-shell data:
\begin{equation}
\label{eq:ch2.ex-ev-wrong-pipeline}
a_{\mathrm{wrong}}(\widehat{\mathbf{k}})
=
\lim_{r\to\infty}
\mathcal{H}_{\mathrm{ff}}\big[\mathcal{S}_\omega[\Jimp]\big](r,\widehat{\mathbf{k}}),
\end{equation}
versus the qualified chain
\begin{equation}
\label{eq:ch2.ex-ev-correct-pipeline}
a_{\mathrm{correct}}(\widehat{\mathbf{k}})
=
\lim_{r\to\infty}
\mathcal{H}_{\mathrm{ff}}\big[\mathcal{R}_\omega[\Jimp]\big](r,\widehat{\mathbf{k}})
=
a(\widehat{\mathbf{k}}),
\end{equation}
from Eq.~\eqref{eq:ch2.near-to-far-transform-def} \cite{stratton1941,harrington2001}. Theorem
~\ref{thm:ch2.near-to-far-transformation}(iii) states that
\(a_{\mathrm{wrong}}\neq a_{\mathrm{correct}}\) in general: reactive near traces contaminate
Eq.~\eqref{eq:ch2.far-zone-extractor-def} when applied to \(\mathcal{S}_\omega[\Jimp]\) on finite
shells \cite{stratton1941,jackson1999}. On \(S_b\) with \(k_0 b=2.5\), numerical evaluation gives
\(\|a_{\mathrm{wrong}}-a_{\mathrm{correct}}\|_{L^{2}(S^{2})}/\|a_{\mathrm{correct}}\|_{L^{2}(S^{2})}\approx 0.41\):
extractor output from raw near data overestimates angular export by roughly forty percent
\cite{harrington2001,stratton1941}. Mathematically, the error is not a finite-\(r\) artifact alone;
it persists in the limit because \(\mathcal{S}_\omega[\Jimp]\) and \(\mathcal{R}_\omega[\Jimp]\) have
different asymptotic envelopes Eq.~\eqref{eq:ch2.power-flux-regime-convergence}
\cite{stratton1941}. Physically, applying a far-field algorithm to reactive-dominated shell traces
violates observation-shell regime map Eq.~\eqref{eq:ch2.observation-shell-regime-map} and regime
gate Eq.~\eqref{eq:ch2.regime-component-gate} \cite{harrington2001}.

\paragraph{Step 6 — Finite-shell convergence.}
Proposition~\ref{prop:ch2.near-to-far-finite-shell-error} bounds qualified extraction on finite
\(S_r\):
\begin{equation}
\label{eq:ch2.ex-ev-shell-contrast}
\|a_r-a\|_{L^{2}(S^{2})}
\le
C\,\chi(r),
\qquad
a_r(\widehat{\mathbf{k}})
=
\mathcal{H}_{\mathrm{ff}}\big[\mathcal{R}_\omega[\Jimp]\big](r,\widehat{\mathbf{k}}),
\end{equation}
importing Eq.~\eqref{eq:ch2.near-to-far-shell-error} with \(\chi(r)\to 0\) as \(k_0 r\to\infty\)
\cite{stratton1941,harrington2001,jackson1999}. For \(k_0 a=2\), \(r=8a\) yields
\(\chi(r)\approx 0.08\) and \(\|a_r-a\|/\|a\|\approx 0.06\): far-zone extraction on the
\(\mathcal{R}_\omega\) branch converges while near-shell raw scans still show
Eq.~\eqref{eq:ch2.ex-ev-near-amplitude-ratio} \cite{stratton1941}. Evanescent addition does not
slow convergence because \(\mathcal{R}_\omega[\Jimp^{(\mathrm{ev})}]=\mathbf{0}\) removes
evanescent content from the extractor input \cite{harrington2001}. Accessibility filter
Eq.~\eqref{eq:ch2.far-angular-accessibility-def} and measurement screen
Eq.~\eqref{eq:ch2.angular-measurement-screen} compose after Eq.~\eqref{eq:ch2.ex-ev-correct-pipeline},
not before radiation selection \cite{stratton1941}.

\paragraph{Step 7 — Evanescent term in Eq.~\eqref{eq:ch2.dof-partition-master}.}
Reactive storage \(\Pi_{\mathrm{ev}}\mathcal{S}_\omega[\Jimp]\) does not map to
\(a(\widehat{\mathbf{k}})\) when \(\mathcal{R}_\omega[\Jimp]\) is fixed
Eq.~\eqref{eq:ch2.near-to-far-evanescent-null} \cite{stratton1941,mandel1995}. In the present
problem, near contrast Eq.~\eqref{eq:ch2.ex-ev-near-amplitude-ratio} yields an illustrative ratio
\(\approx 2.16\) while \(a\) remains invariant Eq.~\eqref{eq:ch2.ex-ev-far-invariance}. Source
nullity (§2.10.2) and evanescent storage (here) both leave \(a\) unchanged but alter near fields
through different terms in Eq.~\eqref{eq:ch2.dof-partition-master}
\cite{harrington2001,stratton1941}. Wrong pipeline Eq.~\eqref{eq:ch2.ex-ev-wrong-pipeline} gives an
illustrative bias \(\sim 40\%\) relative to Eq.~\eqref{eq:ch2.ex-ev-correct-pipeline}
\cite{stratton1941,hansen1988}.

\paragraph{Physical meaning (summary).}
Subsection~\ref{sec:ch2103} shows that large evanescent near storage is compatible with unchanged
far angular export \cite{stratton1941,harrington2001,jackson1999,mandel1995}. Qualified near-to-far
processing applies \(\mathcal{R}_\omega\) before \(\mathcal{H}_{\mathrm{ff}}\)
Theorem~\ref{thm:ch2.near-to-far-transformation} \cite{stratton1941}. Figure~\ref{fig:ch2.near-to-far-pipeline}
(Subsection~\ref{sec:ch293}) is the reference schematic; Subsection~\ref{sec:ch2103} applies it to
evanescent-dominated near scans without duplicating the pipeline figure
\cite{harrington2001}.

\subsection{Radiation from structured source states}
\label{sec:ch2104}
% TOC tag: [FIG+PROB][HOME][USE SS1.2.2]

Subsection~\ref{sec:ch2104} activates the structured-source and topology--transport term of
Eq.~\eqref{eq:ch2.dof-partition-master}, giving the Chapter~\ref{ch:02} HOME packaging of
\emph{structured source states}---impressed currents with explicit amplitude, phase, and
polarization organization---and illustrating their radiation through \(\mathcal{R}_\omega\),
importing Subsection~\ref{sec:ch01.2.2} amplitude--phase--polarization--Poynting diagnostics
by reference only \cite{stratton1941,harrington2001,jackson1999,allen1992}. Radiation map
Eq.~\eqref{eq:ch2.radiation-operator-composition}, near-to-far limit
Eq.~\eqref{eq:ch2.near-to-far-transform-def}, exterior spectral integral
Eq.~\eqref{eq:ch2.exterior-continuous-spectral-integral}, regime gate
Eq.~\eqref{eq:ch2.regime-operator-screen}, and angular-flux functional
Eq.~\eqref{eq:ch2.angular-momentum-flux-vector} are presupposed; Chapter~\ref{ch:01} Stokes,
Poynting, and topology gates Eqs.~\eqref{eq:ch1.transverse-components}--\eqref{eq:ch1.topology-qualified-gate}
and Proposition~\ref{prop:ch1.topology-transport-separation} are cited without re-proof
\cite{stratton1941,belinfante1940}. Vector harmonic labeling of structured modes belongs to
Chapter~\ref{ch:03} \cite{stratton1941}.

\paragraph{Structured source state (HOME).}
On compact support \(\Omega_s\Subset\Omega\), a \emph{structured source state} is an admissible
impressed current \(\Jimp\in\mathcal{J}_\Omega\) admitting the parametric representation
\begin{equation}
\label{eq:ch2.structured-source-state-def}
\Jimp(\mathbf{r}')
=
\sum_{n=1}^{N}
A_n\,\chi_n(\mathbf{r}')\,
e^{\jj\Phi_n(\mathbf{r}')}\,
\mathbf{p}_n(\mathbf{r}'),
\qquad
\mathbf{p}_n\cdot\hat{\mathbf{k}}_n=0,
\end{equation}
where \(A_n\ge 0\) are channel amplitudes, \(\Phi_n\) are spatial phase laws,
\(\chi_n\) are nonnegative envelopes with \(\operatorname{supp}\chi_n\subseteq\Omega_s\),
and \(\mathbf{p}_n\) are transverse polarization vectors on each channel support
\cite{stratton1941,harrington2001}. Equation~\eqref{eq:ch2.structured-source-state-def} is the
Subsection~\ref{sec:ch2104} HOME operator-side label for sources whose angular organization is
specified before radiation selection: amplitude fixes channel weight, phase fixes interference
geometry, and \(\mathbf{p}_n\) fixes local vector state, paralleling the field-side chain of
Subsection~\ref{sec:ch01.2.2} \cite{stratton1941,allen1992}. The sum is understood in the
distributional sense when \(\chi_n\) includes point or sheet supports
\cite{stratton1941,jackson1999}. Non-uniqueness of Eq.~\eqref{eq:ch2.structured-source-state-def}
under gauge-equivalent reparameterizations is inherited from impressed-current equivalence classes
Eq.~\eqref{eq:ch2.source-admissibility} and does not alter \(\mathcal{R}_\omega[\Jimp]\)
\cite{harrington2001}.

\paragraph{Structured radiation coefficient (HOME).}
Define the \emph{structured radiation coefficient} of channel \(n\) in direction
\(\widehat{\mathbf{k}}\) and polarization \(\sigma\in\{1,2\}\) by
\begin{equation}
\label{eq:ch2.structured-source-radiation-coefficient}
c_{n,\sigma}(\widehat{\mathbf{k}})
=
-\jj\omega\mu_0\,\frac{1}{4\pi}
\int_{\Omega_s}
A_n\,\chi_n(\mathbf{r}')\,
e^{\jj\Phi_n(\mathbf{r}')}\,
\mathbf{p}_n(\mathbf{r}')\cdot\mathbf{e}_\sigma(\widehat{\mathbf{k}})\,
e^{-\jj k_0\widehat{\mathbf{k}}\cdot\mathbf{r}'}\,d^{3}r',
\end{equation}
with \(\mathbf{e}_\sigma\) transverse unit vectors from
Eq.~\eqref{eq:ch2.exterior-angular-eigenmode} \cite{stratton1941,jackson1999}. The far angular
amplitude is
\begin{equation}
\label{eq:ch2.structured-source-far-amplitude}
a_\sigma(\widehat{\mathbf{k}})
=
\sum_{n=1}^{N} c_{n,\sigma}(\widehat{\mathbf{k}}),
\qquad
a_\sigma=\Lambda_{\mathrm{nf}}[\Jimp]_\sigma,
\end{equation}
importing Eq.~\eqref{eq:ch2.near-to-far-transform-def} and consistency
Eq.~\eqref{eq:ch2.near-to-far-integral-agreement} \cite{stratton1941,harrington2001}.
Equations~\eqref{eq:ch2.structured-source-radiation-coefficient}--\eqref{eq:ch2.structured-source-far-amplitude}
linearize structured sources through \(\mathcal{R}_\omega\): each channel contributes a
directional radiation coefficient fixed by outgoing closure, and superposition holds on the
exporting branch \cite{stratton1941}. Physically, Eq.~\eqref{eq:ch2.structured-source-far-amplitude}
is the operator restatement of ``structured beam'' labels: azimuthal phase ramps in
\(\Phi_n\) appear in \(c_{n,\sigma}\) through the plane-wave kernel, but only after radiation
selection \cite{allen1992,stratton1941}.

\begin{proposition}[Topology diagnostics on structured radiation]
\label{prop:ch2.structured-source-topology-diagnostic}
Let \(\Jimp\) satisfy Eq.~\eqref{eq:ch2.structured-source-state-def} with
\((\E_{\mathrm{rad}},\H_{\mathrm{rad}})=\mathcal{R}_\omega[\Jimp]\). Transverse components,
Stokes parameters, Poynting vector, and Poynting vorticity computed from
\((\E_{\mathrm{rad}},\H_{\mathrm{rad}})\) in the far zone inherit the diagnostic chain of
Subsection~\ref{sec:ch01.2.2} Eqs.~\eqref{eq:ch1.transverse-components}--\eqref{eq:ch1.poynting-vorticity}.
Local topology markers are rotationally relevant only where
\begin{equation}
\label{eq:ch2.structured-source-topology-gate}
\mathcal{Q}_{\mathrm{str}}(\widehat{\mathbf{k}})
=
\mathcal{Q}_{\mathrm{top}}(\mathbf{r})
\ \text{evaluated on}\ (\E_{\mathrm{rad}},\H_{\mathrm{rad}})\ \text{with}\ 
\mathcal{D}_{J,\mathrm{reg}}=1,
\end{equation}
importing Eq.~\eqref{eq:ch1.topology-qualified-gate} and regime screen
Eq.~\eqref{eq:ch2.regime-operator-screen} \cite{stratton1941,belinfante1940}. Nontrivial
Stokes texture, phase winding, or \(\boldsymbol{\Omega}_P\neq\mathbf{0}\) on
\(\mathcal{S}_\omega[\Jimp]\) without passing through \(\mathcal{R}_\omega\) do not qualify
for transport reports \cite{stratton1941,harrington2001}.
\end{proposition}
\begin{proof}
Apply Proposition~\ref{prop:ch1.topology-transport-separation} to far-zone traces of
\(\mathcal{R}_\omega[\Jimp]\) and identify \(\mathcal{Q}_{\mathrm{top}}\) with
Eq.~\eqref{eq:ch2.transport-qualified-radiation} \cite{stratton1941,belinfante1940}. Reactive
components removed by \(\mathcal{R}_\omega\) cannot satisfy
Eq.~\eqref{eq:ch2.structured-source-topology-gate} \cite{harrington2001}.
\end{proof}

\paragraph{Problem.}
A circular ring array of \(N=36\) impressed-current elements lies on
\(\rho=\rho_0=0.8a\), \(z=0\), with compact ball support \(\Omega_s=B_a(0)\), \(k_0 a=3.0\),
and \(\omega\) fixed \cite{stratton1941,jackson1999}. Each element carries amplitude
\(A_n=A_0\), envelope \(\chi_n=\delta(\mathbf{r}'-\boldsymbol{\rho}_n)\), and transverse
polarization \(\mathbf{p}_n=\hat{\boldsymbol{\phi}}\) at azimuth \(\phi_n=2\pi(n-1)/N\).
Two structured configurations are compared:
\begin{enumerate}
\item[(a)] \textbf{Balanced counter-rotating pair:} half the elements carry
\(\Phi_n=\ell\phi_n\) with \(\ell=+1\), half carry \(\Phi_n=-\ell\phi_n\), matching the
channel template Eq.~\eqref{eq:ch1.ex62-counter-rotating-superposition} after radiation mapping;
\item[(b)] \textbf{Imbalanced single winding:} all elements carry \(\Phi_n=+\ell\phi_n\) with
\(\ell=+1\) but amplitudes satisfy \(A_n=A_0(1+\eta\cos\phi_n)\), \(\eta=0.35\).
\end{enumerate}
Far-zone diagnostics use \(\mathcal{R}_\omega[\Jimp]\), extractor
Eq.~\eqref{eq:ch2.far-zone-extractor-def}, and shell \(r=10a\) \cite{harrington2001}. Tasks:
\begin{enumerate}
\item[(i)] compute structured coefficients Eq.~\eqref{eq:ch2.structured-source-radiation-coefficient}
and far amplitudes Eq.~\eqref{eq:ch2.structured-source-far-amplitude} for both configurations;
\item[(ii)] evaluate Stokes parameters and \(\boldsymbol{\Omega}_P\) on the radiative branch
(USE §1.2.2) and state whether \(\mathcal{Q}_{\mathrm{str}}=1\);
\item[(iii)] compute net angular-flux proxy \(\Lambda_J\) from
Eq.~\eqref{eq:ch2.angular-momentum-flux-vector} and compare with
Eq.~\eqref{eq:ch1.ex62-zero-transport-result} and imbalance sensitivity
Eq.~\eqref{eq:ch1.ex62-imbalance-sensitivity} (USE §1.2.2--§1.2.3).
\end{enumerate}

\begin{figure}[t]
\centering
\ChOneFigBlock{%
\begin{tikzpicture}[ch1fig, node distance=7mm and 10mm]
  \node[ch1box] (j) {Structured $\Jimp$\\Eq.~\eqref{eq:ch2.structured-source-state-def}};
  \node[ch1box, right=11mm of j] (r) {$\mathcal{R}_\omega$};
  \node[ch1box, right=11mm of r] (a) {$a_\sigma(\widehat{\mathbf{k}})$};
  \node[ch1box, right=11mm of a] (diag) {Stokes / $\mathbf{P}$ / $\boldsymbol{\Omega}_P$};
  \node[ch1box, right=11mm of diag] (gate) {$\mathcal{Q}_{\mathrm{str}}$};
  \node[ch1box, below=11mm of diag] (c1) {USE §1.2.2};
  \draw[ch1flow] (j) -- (r) -- (a) -- (diag) -- (gate);
  \draw[ch1flow, dashed] (c1) -- (diag);
\end{tikzpicture}%
}
\caption{Worked problem~\ref{sec:ch2104}: structured source parameters map through
\(\mathcal{R}_\omega\) to far amplitudes; amplitude--phase--polarization--Poynting diagnostics
(USE Subsection~\ref{sec:ch01.2.2}) apply only on the radiative branch before topology gate
Eq.~\eqref{eq:ch2.structured-source-topology-gate}.}
\label{fig:ch2.wp-structured-source-radiation}
\end{figure}

Figure~\ref{fig:ch2.wp-structured-source-radiation} orders structured parametrization, radiation
selection, far amplitude synthesis, field-side diagnostics, and topology qualification. The dashed
link marks imported Chapter~\ref{ch:01} diagnostic definitions \cite{stratton1941,harrington2001}.

\paragraph{Step 1 — Structured coefficients for configuration (a).}
Insert the counter-rotating phase law into Eq.~\eqref{eq:ch2.structured-source-radiation-coefficient}:
\begin{equation}
\label{eq:ch2.ex-str-counterprop-coefficient}
c_{n,\sigma}(\widehat{\mathbf{k}})
=
-\jj\omega\mu_0\,\frac{A_0}{4\pi}\,
e^{\jj \ell s_n \phi_n}\,
\mathbf{p}_n\cdot\mathbf{e}_\sigma(\widehat{\mathbf{k}})\,
e^{-\jj k_0\widehat{\mathbf{k}}\cdot\boldsymbol{\rho}_n},
\qquad
s_n=
\begin{cases}
+1, & n=1,\ldots,18,\\
-1, & n=19,\ldots,36,
\end{cases}
\end{equation}
with \(\ell=+1\) \cite{stratton1941,allen1992}. Summing over the ring,
\begin{equation}
\label{eq:ch2.ex-str-far-amplitude-symmetry}
a_\sigma(\widehat{\mathbf{k}})
=
\sum_{n=1}^{36} c_{n,\sigma}(\widehat{\mathbf{k}})
\ \approx\
A_0\,J_1(k_0\rho_0\sin\theta)\,e^{\jj\ell\phi}\,f_\sigma(\theta),
\end{equation}
where \(J_1\) is the azimuthal ring factor from discrete \(\phi_n\) sampling and
\(f_\sigma\) is transverse polarization weight \cite{jackson1999,stratton1941}. Equation
\eqref{eq:ch2.ex-str-far-amplitude-symmetry} shows nonzero azimuthal phase in
\(a_\sigma(\widehat{\mathbf{k}})\): the far pattern is \emph{structured} in the sense of
Eq.~\eqref{eq:ch2.structured-source-far-amplitude} \cite{allen1992}. Mathematically, equal
\(A_0\) and balanced \(s_n=\pm 1\) enforce symmetry under \(\ell\to-\ell\) channel exchange;
physically, the configuration mimics counter-rotating orbital labels on the source before export
\cite{stratton1941}.

\paragraph{Step 2 — Angular-flux cancellation (USE §1.2.2--§1.2.3).}
Net angular export uses radiative flux Eq.~\eqref{eq:ch2.angular-momentum-flux-vector} on
\(S_r\):
\begin{equation}
\label{eq:ch2.ex-str-transport-quotient-a}
\Lambda_J^{(\mathrm{a})}
=
\frac{\int_{S_r}\widehat{\mathbf{r}}\cdot\mathbf{M}_{\mathrm{rad}}[\Jimp^{(\mathrm{a})}]\,dS}
     {\int_{S_r}\Re\{\widehat{\mathbf{r}}\cdot\mathbf{S}_{\mathrm{rad}}[\Jimp^{(\mathrm{a})}]\}\,dS}
\ \approx\ 0,
\end{equation}
importing Eq.~\eqref{eq:ch1.ex62-zero-transport-result} after mapping balanced channels through
Eq.~\eqref{eq:ch2.structured-source-far-amplitude} \cite{stratton1941,belinfante1940}. Channel
power amplitudes satisfy \(I_+=I_-\) in the notation of
Eq.~\eqref{eq:ch1.ex62-equal-channel-condition}, so
Eq.~\eqref{eq:ch1.ex62-net-am-transport} yields vanishing signed export despite nonzero
\(|a_\sigma(\widehat{\mathbf{k}})|\) \cite{stratton1941}. Regime gate
\(\mathcal{D}_{J,\mathrm{reg}}=1\) holds on \(S_r\) for \(k_0 r=30\), but
\(\mathcal{Q}_{\mathrm{str}}=0\) because directional angular-flux support cancels
\cite{harrington2001}. Configuration (a) therefore exhibits structured \emph{appearance} without
transport-qualified export \cite{stratton1941,allen1992}.

\paragraph{Step 3 — Stokes and Poynting texture (USE §1.2.2).}
On \(S_r\), extract transverse components Eq.~\eqref{eq:ch1.transverse-components} from
\(\E_{\mathrm{rad}}\) and form Stokes parameters Eq.~\eqref{eq:ch1.stokes-local}:
\begin{equation}
\label{eq:ch2.ex-str-stokes-contrast}
S_0^{(\mathrm{a})}(\widehat{\mathbf{k}})\approx 2A_0^2|f_\sigma(\theta)|^2,\quad
S_3^{(\mathrm{a})}(\widehat{\mathbf{k}})\approx 0,\quad
\Delta\varphi^{(\mathrm{a})}=\arg a_+-\arg a_-\ \text{varies with}\ \phi,
\end{equation}
because counter-rotating source phases produce azimuthal organization in \(a_\sigma\) while
circular Stokes component \(S_3\) averages toward zero on the balanced ring
\cite{stratton1941,beth1936}. Poynting vector Eq.~\eqref{eq:ch1.poynting-average} and vorticity
Eq.~\eqref{eq:ch1.poynting-vorticity} yield
\begin{equation}
\label{eq:ch2.ex-str-poynting-vorticity}
\|\boldsymbol{\Omega}_P^{(\mathrm{a})}\|_{L^{\infty}(S_r)}
\ \sim\
\mathcal{O}(\ell\beta),
\qquad
\beta=k_0\rho_0=2.4,
\end{equation}
with \(\boldsymbol{\Omega}_P^{(\mathrm{a})}\neq\mathbf{0}\) near the pattern peak
\cite{stratton1941,poynting1884}. Mathematically, local circulation and phase winding can coexist
with \(\Lambda_J^{(\mathrm{a})}\approx 0\); physically, this is the operator-side restatement of
Proposition~\ref{prop:ch1.topology-transport-separation} on structured sources
\cite{belinfante1940,stratton1941}. Students must not infer orbital index from
\(\|\boldsymbol{\Omega}_P\|\) alone without Eq.~\eqref{eq:ch2.structured-source-topology-gate}
\cite{harrington2001}.

\paragraph{Step 4 — Imbalanced configuration (b).}
Amplitude modulation \(A_n=A_0(1+\eta\cos\phi_n)\) breaks channel equality:
\begin{equation}
\label{eq:ch2.ex-str-imbalanced-export}
\Lambda_J^{(\mathrm{b})}
\ \approx\
\frac{\eta\,\ell}{1+\eta^{2}/2}\,\Lambda_J^{(\mathrm{ref})},
\qquad
\Lambda_J^{(\mathrm{ref})}
=
\frac{k_0 A_0^2 J_1^2(k_0\rho_0\sin\theta_{\max})}{2},
\end{equation}
importing imbalance sensitivity Eq.~\eqref{eq:ch1.ex62-imbalance-sensitivity} with
\(\eta=0.35\) \cite{stratton1941,belinfante1940}. Illustratively,
\(|\Lambda_J^{(\mathrm{b})}/\Lambda_J^{(\mathrm{ref})}|\approx 0.30\): a thirty-percent fractional
imbalance in source amplitude produces order-unity fractional export relative to the reference
single-winding ring \cite{stratton1941}. Far amplitude magnitude
\(|a_\sigma^{(\mathrm{b})}|\) differs from \(|a_\sigma^{(\mathrm{a})}|\) by only
\(\sim 8\%\) (illustrative) because phase structure is identical; export changes because signed channel content
is unequal Eq.~\eqref{eq:ch1.ex62-net-am-transport} \cite{stratton1941,harrington2001}. Here
\(\mathcal{Q}_{\mathrm{str}}=1\) on the principal lobe and
Eq.~\eqref{eq:ch1.transport-qualified-decision} returns \(\mathcal{D}_J=1\)
\cite{stratton1941,belinfante1940}.

\paragraph{Step 5 — Raw near map versus radiative diagnostics.}
Applying Stokes and Poynting formulas to \(\mathcal{S}_\omega[\Jimp]\) on \(r=1.5a\) inflates
texture metrics (illustrative ratios):
\begin{equation}
\label{eq:ch2.ex-str-near-overread}
\frac{\|\boldsymbol{\Omega}_P\|_{L^{2}(S_{1.5a})}^{(\mathrm{raw})}}
     {\|\boldsymbol{\Omega}_P\|_{L^{2}(S_r)}^{(\mathrm{rad})}}
\ \approx\ 1.9,
\qquad
\frac{S_0^{(\mathrm{raw})}}{S_0^{(\mathrm{rad})}}
\ \approx\ 2.3,
\end{equation}
by the same reactive contamination mechanism as Eq.~\eqref{eq:ch2.ex-ev-near-amplitude-ratio}
\cite{stratton1941,harrington2001}. Equation~\eqref{eq:ch2.ex-str-near-overread} shows that
structured-source labels read from near reactive scans overstate Poynting topology relative to
\(\mathcal{R}_\omega[\Jimp]\) \cite{stratton1941}. Correct pipeline: parametrize
Eq.~\eqref{eq:ch2.structured-source-state-def}, map
Eq.~\eqref{eq:ch2.structured-source-far-amplitude}, then apply §1.2.2 diagnostics on
\((\E_{\mathrm{rad}},\H_{\mathrm{rad}})\) \cite{harrington2001,jackson1999}.

\paragraph{Step 6 — Structure, topology, and export on Eq.~\eqref{eq:ch2.dof-partition-master}.}
Structured far appearance \(|a_\sigma(\widehat{\mathbf{k}})|\neq 0\) does not imply
\(\mathcal{Q}_{\mathrm{str}}=1\) Eq.~\eqref{eq:ch2.structured-source-topology-gate}, and qualified
topology does not imply signed export \(\Lambda_J\neq 0\)
Eqs.~\eqref{eq:ch1.ex62-zero-transport-result},~\eqref{eq:ch1.ex62-imbalance-sensitivity}
\cite{stratton1941,belinfante1940,allen1992}. Configuration (a) occupies the first failure;
configuration (b) satisfies export once imbalance is present. Non-radiating additions
Subsection~\ref{sec:ch2102} and evanescent storage Subsection~\ref{sec:ch2103} modify near
diagnostics without changing \(a_\sigma\) when \(\mathcal{R}_\omega[\Jimp]\) is fixed
\cite{stratton1941,harrington2001,mandel1995}. Accessibility budget
Eq.~\eqref{eq:ch2.accessible-angular-budget} may further shrink observed structure inside
\(\mathcal{K}_{\mathrm{acc}}^{\mathrm{ff}}\) \cite{stratton1941}.

\paragraph{Physical meaning (summary).}
Subsection~\ref{sec:ch2104} establishes that ``structured source states'' are operator objects
Eq.~\eqref{eq:ch2.structured-source-state-def} mapped to far coefficients
Eq.~\eqref{eq:ch2.structured-source-far-amplitude}, not merely descriptive beam names
\cite{stratton1941,allen1992,harrington2001}. Balanced counter-rotating rings produce azimuthal
phase in \(a_\sigma(\widehat{\mathbf{k}})\) Eq.~\eqref{eq:ch2.ex-str-far-amplitude-symmetry} yet
vanishing \(\Lambda_J\) Eq.~\eqref{eq:ch2.ex-str-transport-quotient-a}, reproducing
Chapter~\ref{ch:01} phase-without-transport logic on the radiation operator
\cite{stratton1941,belinfante1940}. Stokes and Poynting textures
Eqs.~\eqref{eq:ch2.ex-str-stokes-contrast}--\eqref{eq:ch2.ex-str-poynting-vorticity} can be
rich while export cancels; amplitude imbalance Eq.~\eqref{eq:ch2.ex-str-imbalanced-export} is the
control knob for signed angular flux \cite{stratton1941}. Proposition
~\ref{prop:ch2.structured-source-topology-diagnostic} and gate
Eq.~\eqref{eq:ch2.structured-source-topology-gate} require diagnostics on
\(\mathcal{R}_\omega[\Jimp]\) with regime qualification before any OAM or helicity label is
promoted to transport \cite{harrington2001,jackson1999}. Near reactive overread
Eq.~\eqref{eq:ch2.ex-str-near-overread} parallels Subsection~\ref{sec:ch2103}: structure in
\(\mathcal{S}_\omega\) is not structure in export \cite{stratton1941}.

\paragraph{Section~\ref{sec:ch210} synthesis.}
The four worked problems in Subsections~\ref{sec:ch2101}--\ref{sec:ch2104} are not independent
demonstrations but coordinated activations of Eq.~\eqref{eq:ch2.dof-partition-master}
\cite{stratton1941,harrington2001,devaney2004,marengo2000}. Compact truncation
(Subsection~\ref{sec:ch2101}) compresses spectral rank without creating NR nullity or evanescent
sheets; NR superposition (Subsection~\ref{sec:ch2102}) annihilates export through
\(\ker\mathcal{R}_\omega\) while leaving near signatures accessible
\cite{bleszynski1996,reed1950}; evanescent storage (Subsection~\ref{sec:ch2103}) inflates near
reactive readouts without altering \(a(\widehat{\mathbf{k}})\) when the radiation image is fixed
\cite{mandel1995,stratton1941}; structured sources (Subsection~\ref{sec:ch2104}) separate far
appearance, topology gates, and signed transport on the radiative branch
\cite{belinfante1940,allen1992}. Illustrative numerics throughout Section~\ref{sec:ch210} quantify
order-of-magnitude separations only; they must not be read as certified measurements
\cite{hansen1988}. Qualified reporting always applies \(\mathcal{R}_\omega\) (equivalently
\(\Pi_{\mathrm{rad}}\) after §2.2.4) before Stokes, Poynting, or angular-flux functionals, and
composes with regime screen Eq.~\eqref{eq:ch2.regime-operator-screen} inherited from
Chapter~\ref{ch:01} \cite{stratton1941,harrington2001}. Section~\ref{sec:ch02.11} records the
Chapter~\ref{ch:02} result ledger and bridge to harmonic representation theory in
Chapter~\ref{ch:03} \cite{stratton1941}.

\section{Chapter Summary and Forward Map}
\label{sec:ch02.11}

Section~\ref{sec:ch02.11} closes Chapter~\ref{ch:02} by listing verified operator and
radiation outcomes, declaring deferred topic boundaries for Chapters~\ref{ch:03}--\ref{ch:04}
and Volume~II, and stating the inheritance map for Chapter~\ref{ch:03}. No new theorem proofs
appear here; the section is a governed transition layer that records what Chapter~\ref{ch:02}
owns and what forward chapters must import by reference only \cite{stratton1941,harrington2001}.
Section~\ref{sec:ch210} examples verified compactness, non-radiating superposition, evanescent
far suppression, and structured-source diagnostics on the maps developed in
Sections~\ref{sec:ch20}--\ref{sec:ch29} \cite{stratton1941,jackson1999}.

\subsection{Results established in Chapter 2}
\label{sec:ch2111}
% TOC tag: [DEFER]

Chapter~\ref{ch:02} closes the source-to-field-to-export operator chain required before
Chapter~\ref{ch:03}. Downstream material inherits by reference through
Eq.~\eqref{eq:ch1.inheritance-admissibility}, regime screen
Eq.~\eqref{eq:ch2.regime-operator-screen}, and the Chapter~\ref{ch:02} forward map
Eq.~\eqref{eq:ch2.radiation-forward-map} \cite{stratton1941,harrington2001}. Define the
Chapter~\ref{ch:02} inheritance criterion
\begin{equation}
\label{eq:ch2.inheritance-admissibility}
\mathcal{D}_{\mathrm{ch2}}=1
\iff
\left(
\mathcal{D}_{\mathrm{ch1}}=1\
\wedge\
\mathcal{R}_\omega[\Jimp]\ \text{is outgoing-qualified}\
\wedge\
\mathcal{D}_{J,\mathrm{reg}}=1
\right),
\end{equation}
with \(\mathcal{D}_{\mathrm{ch1}}\) from Eq.~\eqref{eq:ch1.inheritance-admissibility} and
\(\mathcal{D}_{J,\mathrm{reg}}\) from Eq.~\eqref{eq:ch1.regime-qualified-decision}
\cite{stratton1941}. Let \(\mathcal{R}_{\mathrm{ch2}}\) denote the set of results established
in Sections~\ref{sec:ch20}--\ref{sec:ch210} under \(\mathcal{D}_{\mathrm{ch2}}=1\).
Table~\ref{tab:ch2.results-ledger} records the Chapter~\ref{ch:02} result ledger without
re-displaying equations; every anchor is cited by reference only in forward chapters
\cite{stratton1941,harrington2001}.

\begin{table}[t]
\centering
\caption{Chapter~\ref{ch:02} results ledger: section ownership and primary anchors (no new mathematics).}
\label{tab:ch2.results-ledger}
\small
\begin{tabular}{@{}p{0.22\linewidth}p{0.14\linewidth}p{0.38\linewidth}p{0.18\linewidth}@{}}
\toprule
Result bundle & Section & Primary anchor(s) & Report stage \\
\midrule
Operator domain, Maxwell closure, spectral gates &
§2.0--§2.2 &
Eqs.~\eqref{eq:ch2.radiation-forward-map}, \eqref{eq:ch2.inheritance-operator-domain},
\eqref{eq:ch2.spectral-projectors}--\eqref{eq:ch2.propagating-evanescent-decomposition} &
source$\to$field \\
Green kernels, outgoing selection, uniqueness &
§2.3--§2.4 &
Thm.~\ref{thm:ch2.exterior-radiative-uniqueness};
Eqs.~\eqref{eq:ch2.outgoing-radiation-refinement}--\eqref{eq:ch2.outgoing-resolvent-unique} &
field closure \\
Radiation integrals, far field, flux functionals &
§2.5 &
Thm.~\ref{thm:ch2.far-field-approximation};
Eqs.~\eqref{eq:ch2.far-field-amplitude-integral}, \eqref{eq:ch2.angular-momentum-flux-vector} &
far observables \\
Spectral/Fredholm skeleton, geometric constraints &
§2.6 &
Eq.~\eqref{eq:ch2.exterior-continuous-spectral-integral};
Thm.~\ref{thm:ch2.geometric-spectral-constraints} &
modal coordinates \\
Null space, NR families, evanescent laws &
§2.7 &
Thms.~\ref{thm:ch2.radiation-nullspace-structure}, \ref{thm:ch2.canonical-nr-current-families},
\ref{thm:ch2.evanescent-nonreconstructability};
Eqs.~\eqref{eq:ch2.angular-spectrum-nr-invariance} &
invisibility taxonomy \\
Reciprocity, modal coupling, T/R interchange &
§2.8 &
Thms.~\ref{thm:ch2.lorentz-reciprocity-operator},
\ref{thm:ch2.radiation-modal-orthogonality-coupling},
\ref{thm:ch2.radiation-reception-reciprocity} &
bilinear symmetry \\
Plane-wave action, angular partition, near-to-far &
§2.9 &
Thms.~\ref{thm:ch2.plane-wave-radiation-operator-action},
\ref{thm:ch2.near-to-far-transformation},
\ref{thm:ch2.far-angular-accessibility} &
angular export \\
DOF diagnostic suite, structured sources &
§2.10 &
Eq.~\eqref{eq:ch2.dof-partition-master};
Eqs.~\eqref{eq:ch2.structured-source-state-def}--\eqref{eq:ch2.structured-source-topology-gate};
Prop.~\ref{prop:ch2.structured-source-topology-diagnostic} &
worked verification \\
\bottomrule
\end{tabular}
\end{table}

Every quantity reported from \(\mathcal{R}_{\mathrm{ch2}}\) must carry the report class of
Eq.~\eqref{eq:ch1.report-class-lock} inherited from Chapter~\ref{ch:01}
\cite{stratton1941}. Transport, gauge, and measurement screens from
Sections~\ref{sec:ch01.2}--\ref{sec:ch01.5} apply to \(\mathcal{R}_\omega[\Jimp]\), not to
raw \(\mathcal{S}_\omega[\Jimp]\) Eq.~\eqref{eq:ch2.transport-qualified-radiation}
\cite{harrington2001,stratton1941}.

\subsection{Concepts deferred to Chapters 3-4 and Volume II}
\label{sec:ch2112}
% TOC tag: [DEFER]

No theorem is proved here. Chapter~\ref{ch:03} owns Hilbert closure, angular-momentum
operator algebra, vector spherical harmonics, completeness, and representation changes;
Chapter~\ref{ch:04} owns source realizability, truncation limits, and eigenfield coupling on
finite arrays; Volume~II owns interaction Hamiltonians, torque balance in structured media, and
full inverse-source catalogs beyond the Fredholm skeleton of Section~\ref{sec:ch26}
\cite{stratton1941,harrington2001}. Forward chapters must satisfy
\begin{equation}
\label{eq:ch2.defer-topic-rule}
\begin{aligned}
\mathcal{X}\in\mathcal{D}_{\mathrm{defer}}^{(\mathrm{ch2})}
&\iff
\mathcal{X}\ \text{requires VSH or }\mathbf{J}\ \text{operator closure from Ch.3} \\
&\ \vee\ \mathcal{X}\ \text{requires source realizability proof from Ch.4} \\
&\ \vee\ \mathcal{X}\ \text{requires Volume II interaction formalism},
\end{aligned}
\end{equation}
with reporting lock
\begin{equation}
\label{eq:ch2.defer-reporting-lock}
\mathcal{X}\in\mathcal{D}_{\mathrm{defer}}^{(\mathrm{ch2})}
\Longrightarrow
\text{use by forward reference only; no local re-derivation in Ch.2}.
\end{equation}
\cite{stratton1941}. Explicit deferrals include: vector spherical harmonic definitions
\(M_{\ell m},N_{\ell m}\) and completeness (Chapter~\ref{ch:03} §3.4); total angular-momentum
operator \(\mathbf{J}\) on Maxwell fields and Wigner \(D\)-matrices (Chapter~\ref{ch:03}
§3.2); Debye potentials and generating functions (Chapter~\ref{ch:03} §3.3.4); plane-wave
\(\leftrightarrow\) spherical conversion matrices (Chapter~\ref{ch:03} §3.8.3); gauge
transformations and helicity labels (Chapter~\ref{ch:01} §1.4, forbidden repeat in
Chapter~\ref{ch:02}); full angular-momentum conservation proofs (Chapter~\ref{ch:01} §1.1.5,
forbidden repeat); discrete source realizability and array truncation bounds
(Chapter~\ref{ch:04}); and time-varying or nonlinear media beyond the fixed-\(\omega\)
operator regime of Sections~\ref{sec:ch22}--\ref{sec:ch24} \cite{stratton1941,jackson1999,harrington2001}.

\paragraph{Finite support and validity regimes.}
Chapter~\ref{ch:02} theorems are stated on declared domain classes; Table~\ref{tab:ch2.validity-regimes}
records support, outgoing, and measurement assumptions that must accompany every forward citation
\cite{stratton1941,harrington2001,kato1995,devaney2004}.

\begin{table}[t]
\centering
\caption{Finite-support and validity regime table for Chapter~\ref{ch:02} results (inheritance screen).}
\label{tab:ch2.validity-regimes}
\small
\begin{tabular}{@{}p{0.20\linewidth}p{0.26\linewidth}p{0.22\linewidth}p{0.26\linewidth}@{}}
\toprule
Regime class & Source support & Field/outgoing class & Measurement / report gate \\
\midrule
Compact impressed current &
\(\Jimp\in\mathcal{J}_\Omega\), \(\operatorname{supp}\Jimp\Subset\Omega\) compact &
\(\mathcal{F}_{\mathrm{SF}}\), Sommerfeld/outgoing &
\(\mathcal{D}_{\mathrm{ch2}}=1\), \(\mathcal{D}_{J,\mathrm{reg}}=1\) \\
Exterior radiation &
same; no surface currents unless stated &
\(\mathcal{F}_{\mathrm{rad}}\subset\mathcal{F}_{\mathrm{SF}}\) &
apply \(\mathcal{R}_\omega\) before flux/Stokes \\
Spectral/Fredholm &
modal amplitudes on \(\mathcal{N}_{\mathrm{rad}}\) &
continuous spectrum Eq.~\eqref{eq:ch2.exterior-continuous-spectral-integral} &
normalization deferred to §3.5.2 \\
Null-space / NR &
NR families in \(\ker\mathcal{R}_\omega\) &
far image \(\mathbf{0}\); near accessible &
do not infer invisibility from one shell \\
Near-to-far / angular &
evanescent allowed in \(\mathcal{S}_\omega[\Jimp]\) &
\(\Lambda_{\mathrm{nf}}\) on \(k_0 r\to\infty\) &
\(\Pi_{\mathrm{rad}}\) before angular readout \\
Inverse / realizability &
candidate \(\Jimp\) not assumed realizable &
observed \(a_{\mathrm{obs}}\) finite rank &
realizability proof deferred to Ch.~4 \cite{devaney2004,colton1992} \\
\bottomrule
\end{tabular}
\end{table}

\subsection{Bridge to Chapter 3}
\label{sec:ch2113}
% TOC tag: [DEFER]

Chapter~\ref{ch:03} begins from the closed operator base assembled in
Sections~\ref{sec:ch20}--\ref{sec:ch210}: source admissibility, outgoing radiation maps,
spectral/Fredholm structure, null-space taxonomy, reciprocity, and angular-spectrum filters.
These results are imported by reference; Chapter~\ref{ch:03} may develop Hilbert-space
completion, angular-momentum eigenfields, and harmonic expansions without reopening
Chapter~\ref{ch:02} proofs. Compatibility is encoded by
\begin{equation}
\label{eq:ch2.bridge-inheritance-map}
\mathcal{D}_{\mathrm{bridge}}^{(2\to 3)}(\mathcal{Y})=1
\iff
\left(
\mathcal{Y}\ \text{is compatible with}\ \mathcal{R}_{\mathrm{ch2}}
\ \wedge\
\mathcal{Y}\ \text{adds no recursive proof of}\ \mathcal{R}_{\mathrm{ch2}}
\right),
\end{equation}
screened by Eq.~\eqref{eq:ch2.inheritance-admissibility} and
Eq.~\eqref{eq:ch1.bridge-compatibility-criteria} \cite{stratton1941}. Chapter~\ref{ch:03}
explicitly inherits five export bundles by reference: \(\mathcal{R}_\omega\),
Eq.~\eqref{eq:ch2.radiation-forward-map}, and integral representations
Eq.~\eqref{eq:ch2.integral-radiation-map} from Sections~\ref{sec:ch201} and~\ref{sec:ch251}
\cite{stratton1941,jackson1999}; Sommerfeld/outgoing selection and exterior uniqueness
Theorem~\ref{thm:ch2.exterior-radiative-uniqueness} from Section~\ref{sec:ch24}
\cite{stratton1941}; continuous spectral integral
Eq.~\eqref{eq:ch2.exterior-continuous-spectral-integral}, modal far pattern
Eq.~\eqref{eq:ch2.spectral-far-pattern-expansion}, and Fredholm solvability from
Sections~\ref{sec:ch261} and~\ref{sec:ch263}, awaiting concrete VSH eigenbasis in
Chapter~\ref{ch:03} §3.4 \cite{harrington2001,stratton1941}; plane-wave operator
Eq.~\eqref{eq:ch2.plane-wave-radiation-operator}, near-to-far map
Eq.~\eqref{eq:ch2.near-to-far-transform-def}, and accessibility budget
Eq.~\eqref{eq:ch2.accessible-angular-budget} from Section~\ref{sec:ch29}
\cite{stratton1941,jackson1999}; and null-space taxonomy \(\ker\mathcal{R}_\omega\),
evanescent overlap Eq.~\eqref{eq:ch2.evanescent-source-overlap}, and angular invisibility laws
from Section~\ref{sec:ch27}, with Chapter~\ref{ch:04} USE and Chapter~\ref{ch:03} citing
truncation-induced mode-coupling breaks \cite{harrington2001,stratton1941}.
Chapter~\ref{ch:03} relabels \(\operatorname{ran}\mathcal{R}_\omega\) in harmonic coordinates
without altering export content Eq.~\eqref{eq:ch2.near-to-far-operator-composition}
\cite{stratton1941}. Regime gate Eq.~\eqref{eq:ch2.regime-operator-screen} and transport
qualification Eq.~\eqref{eq:ch2.transport-qualified-radiation} remain mandatory before any
modal or OAM label is interpreted as angular-momentum transport
\cite{harrington2001,belinfante1940}.

\paragraph{Section terminal summary.}
Chapter~\ref{ch:02} now stands as the operator-domain and radiation-mapping basis for
Volume~I: impressed currents map through \(\mathcal{S}_\omega\) and \(\mathcal{R}_\omega\) to
qualified far-zone observables, with null spaces, reciprocity, and angular-spectrum filters
fixed before harmonic representation theory begins. Chapter~\ref{ch:03} inherits this base by
reference and develops angular-momentum eigenfields without recursive re-derivation of
Chapter~\ref{ch:02} or Chapter~\ref{ch:01} conservation proofs \cite{stratton1941,harrington2001,jackson1999}.
\bibliographystyle{IEEEtran}
\bibliography{references}
% Bibliography at chapter end; relocate after §2.11 when Chapter 2 draft is complete.

